% IRdrone user manual
% Compilation: pdflatex irdrone_help.tex (two passes for the table of contents and the references)
% Each button has a \label{btn:...}: these labels are used to open the PDF at the right page.
\documentclass[a4paper,11pt]{report}
\usepackage[utf8]{inputenc}
\usepackage[T1]{fontenc}
\usepackage[english]{babel}
\usepackage[margin=2.5cm]{geometry}
\usepackage{graphicx}
\usepackage{float}
\usepackage{caption}
\usepackage{subcaption}
\usepackage{xcolor}
\usepackage[colorlinks=true,linkcolor=blue!50!black]{hyperref}
\usepackage{refcount}

\usepackage{epsfig}
\usepackage{mathrsfs}
\usepackage{amsfonts}
\usepackage{amssymb}
\usepackage{shadow}
\usepackage{amsmath}
\usepackage{amsmath}
\usepackage{multirow}
\usepackage{makeidx}
\usepackage{color}
\usepackage{stmaryrd}

% Light version: compile with  pdflatex "\def\IRDlight{1}% Manuel utilisateur IRdrone
% Compilation : pdflatex irdrone_help.tex (deux passes pour la table des matieres et les references)
% Chaque bouton a un \label{btn:...} : ces labels serviront a ouvrir le PDF a la bonne page.
\documentclass[a4paper,11pt]{report}
\usepackage[utf8]{inputenc}
\usepackage[T1]{fontenc}
\usepackage[french]{babel}
\usepackage[margin=2.5cm]{geometry}
\usepackage{graphicx}
\usepackage{float}
\usepackage{caption}
\usepackage{subcaption}
\usepackage{xcolor}
\usepackage[colorlinks=true,linkcolor=blue!50!black]{hyperref}
\usepackage{refcount}

\usepackage{epsfig}
\usepackage{mathrsfs}
\usepackage{amsfonts}
\usepackage{amssymb}
\usepackage{shadow}
\usepackage{amsmath}
\usepackage{amsmath}
\usepackage{multirow}
\usepackage{makeidx}
\usepackage{color}
\usepackage{stmaryrd}

% Version legere : compiler avec  pdflatex "\def\IRDlight{1}% Manuel utilisateur IRdrone
% Compilation : pdflatex irdrone_help.tex (deux passes pour la table des matieres et les references)
% Chaque bouton a un \label{btn:...} : ces labels serviront a ouvrir le PDF a la bonne page.
\documentclass[a4paper,11pt]{report}
\usepackage[utf8]{inputenc}
\usepackage[T1]{fontenc}
\usepackage[french]{babel}
\usepackage[margin=2.5cm]{geometry}
\usepackage{graphicx}
\usepackage{float}
\usepackage{caption}
\usepackage{subcaption}
\usepackage{xcolor}
\usepackage[colorlinks=true,linkcolor=blue!50!black]{hyperref}
\usepackage{refcount}

\usepackage{epsfig}
\usepackage{mathrsfs}
\usepackage{amsfonts}
\usepackage{amssymb}
\usepackage{shadow}
\usepackage{amsmath}
\usepackage{amsmath}
\usepackage{multirow}
\usepackage{makeidx}
\usepackage{color}
\usepackage{stmaryrd}

% Version legere : compiler avec  pdflatex "\def\IRDlight{1}% Manuel utilisateur IRdrone
% Compilation : pdflatex irdrone_help.tex (deux passes pour la table des matieres et les references)
% Chaque bouton a un \label{btn:...} : ces labels serviront a ouvrir le PDF a la bonne page.
\documentclass[a4paper,11pt]{report}
\usepackage[utf8]{inputenc}
\usepackage[T1]{fontenc}
\usepackage[french]{babel}
\usepackage[margin=2.5cm]{geometry}
\usepackage{graphicx}
\usepackage{float}
\usepackage{caption}
\usepackage{subcaption}
\usepackage{xcolor}
\usepackage[colorlinks=true,linkcolor=blue!50!black]{hyperref}
\usepackage{refcount}

\usepackage{epsfig}
\usepackage{mathrsfs}
\usepackage{amsfonts}
\usepackage{amssymb}
\usepackage{shadow}
\usepackage{amsmath}
\usepackage{amsmath}
\usepackage{multirow}
\usepackage{makeidx}
\usepackage{color}
\usepackage{stmaryrd}

% Version legere : compiler avec  pdflatex "\def\IRDlight{1}\input{irdrone_help.tex}"
% (voir build_help.bat). Les images viennent alors de images_light/ au lieu de images/.
\ifdefined\IRDlight
  \newcommand{\IRDimgdir}{images_light}
  \graphicspath{{images_light/}}
\else
  \newcommand{\IRDimgdir}{images}
  \graphicspath{{images/}}
\fi

% Symboles des boutons du logiciel (pdflatex ne les connait pas)
\DeclareUnicodeCharacter{25C0}{\ensuremath{\blacktriangleleft}}
\DeclareUnicodeCharacter{25B6}{\ensuremath{\blacktriangleright}}
\DeclareUnicodeCharacter{2699}{\ensuremath{\circledast}}
\DeclareUnicodeCharacter{2922}{\ensuremath{\nearrow}}
\DeclareUnicodeCharacter{21BA}{\ensuremath{\circlearrowleft}}
\DeclareUnicodeCharacter{21BB}{\ensuremath{\circlearrowright}}
\DeclareUnicodeCharacter{25B2}{\ensuremath{\blacktriangle}}
\DeclareUnicodeCharacter{25BC}{\ensuremath{\blacktriangledown}}
% Bouton du logiciel
\DeclareRobustCommand{\bouton}[1]{\texorpdfstring{\fbox{\small\textsf{#1}}}{#1}}
% A completer
\newcommand{\todo}[1]{\textcolor{red}{[\`A r\'ediger : #1]}}
% Sections historiques : repliees dans le manuel web, texte normal dans le PDF
\newcommand{\legacybegin}[1]{}
\newcommand{\legacyend}{}
% Figure tolerante : affiche un cadre si l'image n'existe pas encore.
% Usage : \figureIRD[label]{fichier}{legende}{largeur}
%   - label optionnel (ex. fig:img_VIR_1), utilisable avec \ref{...} / \autoref{...}
%   - fichier avec ou sans extension ; sans extension, essaie .png, .jpg, .jpeg puis .pdf
\makeatletter
\newcommand{\ird@file}{}
\newcommand{\ird@find}[1]{%
  \def\ird@file{}%
  \IfFileExists{\IRDimgdir/#1}{\def\ird@file{#1}}{%
  \IfFileExists{\IRDimgdir/#1.png}{\def\ird@file{#1.png}}{%
  \IfFileExists{\IRDimgdir/#1.jpg}{\def\ird@file{#1.jpg}}{%
  \IfFileExists{\IRDimgdir/#1.jpeg}{\def\ird@file{#1.jpeg}}{%
  \IfFileExists{\IRDimgdir/#1.pdf}{\def\ird@file{#1.pdf}}{}}}}}}
\newcommand{\figureIRD}[4][]{%
  \ird@find{#2}%
  \begin{figure}[htbp]\centering
    \ifx\ird@file\@empty
      \fbox{\parbox{0.6\linewidth}{\centering\small\textcolor{gray}{Image manquante : \detokenize{#2}}}}%
    \else
      \includegraphics[width=#4\linewidth]{\ird@file}%
    \fi
    \caption{#3}%
    \if\relax\detokenize{#1}\relax\else\label{#1}\fi
  \end{figure}}
% Paire de figures cote a cote, placee a l'endroit exact ou elle est ecrite ([H]).
% Usage : \figurePaireIRD{label}{fichier gauche}{legende gauche}{fichier droit}{legende droite}{legende globale}
%   - le label designe la paire ; les sous-figures ont pour labels <label>_avant et <label>_apres
\newcommand{\ird@incsub}[1]{%
  \ird@find{#1}%
  \ifx\ird@file\@empty
    \fbox{\parbox{0.8\linewidth}{\centering\small\textcolor{gray}{Image manquante : \detokenize{#1}}}}%
  \else
    \includegraphics[width=\linewidth]{\ird@file}%
  \fi}
\newcommand{\figurePaireIRD}[6]{%
  \begin{figure}[H]\centering
    \begin{subfigure}[t]{0.48\linewidth}\centering
      \ird@incsub{#2}%
      \caption{#3}\label{#1_avant}
    \end{subfigure}\hfill
    \begin{subfigure}[t]{0.48\linewidth}\centering
      \ird@incsub{#4}%
      \caption{#5}\label{#1_apres}
    \end{subfigure}
    \caption{#6}\label{#1}
  \end{figure}}
\makeatother

\title{IRdrone\\Manuel de l'utilisateur}
\author{}
\date{\today}


\begin{document}
	
	% --------------------------------------------------
	% Cover page
	% --------------------------------------------------
	\begin{titlepage}
		\centering
						
		\vspace*{2cm}		
		%{\Huge\bfseries IRdrone\par}		
		%\vspace{0.5cm}		
		{\LARGE Manuel de l'utilisateur\par}
		\vspace{6cm}		
		\includegraphics[width=0.95\textwidth]{Logo-IRdrone.png}		
		
		\vfill		
		{\large Version du \today\par}		
	\end{titlepage}
	
	% --------------------------------------------------
	% Table of contents
	% --------------------------------------------------
	\tableofcontents



%==========================================================
\chapter{Le projet IRDrone}\label{ch:overview}


\section{Observer le couvert végétal.}

\subsection{Contexte général.}
La cartographie de la végétation est aujourd'hui classiquement employée dans de nombreux types de diagnostics environnementaux (site Natura 2000, ZNIEFF, études d'impact environnemental etc.). 
Cette cartographie est basée sur une segmentation (détourage d'ensembles homogènes au plan de la végétation) du couvert végétal en vue plane aérienne. 
Cette approche est généralement complétée par des prospections au sol (relevés botaniques, phytosociologiques etc.). 
La segmentation d'ensembles homogènes repose sur l'analyse des photographies aériennes en couleur naturelle, notamment sur leur couleur, leur texture, leur structure apparente, leur opacité etc.

\figurePaireIRD{fig:etude_de_cas}
{step0_carto_aerienne.jpg}{Délimitation des habitats sur fond image aérienne.}
{step0_habitats.jpg}{Carte restituée des habitats sur fond IGN topo 25.}
{Exemple d'exploitation de données terrain à l'aide de cartes et photos aériennes.}

Il est toutefois délicat de distinguer les essences (arbres) et les espèces végétales (espèces herbacées) sur des photographies aériennes classiques. Les dynamiques et les éléments fins du paysage sont également difficiles à cerner. 
Les photographies aériennes IR disponibles sont souvent anciennes, donc difficilement utilisables dans le cadre de diagnostics actuels. 
D'autre part la couverture nationale est lacunaire, avec des régions mal couvertes. 
Ajoutons que toutes les photos existantes sont réalisées en altitude (3-4000 m) et manquent donc de précision pour des diagnostics réalisés couramment à l'échelle du 1/5000, voire moins. \\




L'utilisation d'un drone pour acquérir des photographies présente les avantages de centrer les clichés sur les secteurs d'études, d'obtenir des images d'une précision largement supérieure (altitude de vol de 120 m), d'obtenir autant d'images récentes qu'on le souhaite, de pouvoir les réaliser à différentes saisons. 


\subsection{Intérêt des images multispectrales.}

Les photographies aériennes dans le proche infrarouge
\footnote{On ne confondra pas ce rayonnement infrarouge proche qui en journée provient du rayonnement solaire et est réfléchi par les objets avec le rayonnement infrarouge lointain qui lui est émis par tous les objets.} NIR sont plus riches en informations écologiques que les photographies en couleur naturelle. 
L'objectif principal de la photographie aérienne NIR est donc de faciliter l'interprétation des photographies aériennes en se basant sur les différences de facteur de réflexion du rayonnement infrarouge entre les types de végétations (en raison de différences structurelles entre les tissus foliaires des espèces végétales).

\figureIRD[fig:img_NDVI_1]{img_NDVI_1.png}{Carte de l'indice NDVI de développement des végétaux.}{0.7}
\figureIRD[fig:img_VIR_1]{img_VIR_1.png}{Exemple d'image composite de type VIR. Les couches R/G/B sont remplac\'ees par les couches NIR/R/G. Ce type de représentation est très utile, entre autres, pour distinguer entre résineux et feuillus.}{0.7} 


\subsection{Que cherchons-nous à observer ?}

L'emploi des photographies aériennes se justifie pour plusieurs raisons, qui ont trait essentiellement à la précision et à la rapidité d'analyse. 
D'une manière générale, sur les images aériennes NIR les contrastes entre les différents types d'habitats sont exacerbés, ce qui permet une lecture plus fine du paysage. 
Dans les peuplements forestiers denses, il est difficile d'estimer les proportions de recouvrement résineux/feuillus en se basant sur une prospection de terrain en raison des grandes surfaces à parcourir et de la hauteur des arbres. 
La différenciation entre résineux/feuillus sur les photographies en couleur naturelle est souvent ambiguë. 
Au contraire, la photographie aérienne NIR permet de discerner précisément et rapidement les contributions de chacune des essences dans les peuplements mélangés couvrant des superficies considérables, et ainsi d'avoir accès à une information concernant la gestion spécifique des peuplements (types de traitements sylvicoles, classes d'âge etc.).  D'autre part, certains peuplements arborescents mêlant feuillus et résineux en proportions variables peuvent, selon l'importance prise par les feuillus, relever ou non d'une protection européenne (directive \og Habitats \fg). 
L'évaluation des proportions des essences devient objective et précisément quantifiable à l'aide des photographies aériennes NIR. 
L'analyse des photographies aériennes NIR permet en outre une meilleure appréhension des phénomènes de dynamique de la végétation.
 
La détection des très jeunes ligneux devient possible en raison de leur vive coloration et de leur position au sein de milieux ouverts (fronts de colonisation ligneux au sein d'espaces ouverts) alors que leur présence est difficilement détectable en couleur naturelle en raison de la petite taille de leur houppier et de leur couleur comparable à celle des habitats ouverts colonisés. 
Une autre application importante concerne la localisation instantanée et exhaustive des masses d'eau (qui peuvent être discrètes, de petite taille et donc passer inaperçues en couleur naturelle) qui apparaissent comme des taches sombres  en NIR. 
Cette application est particulièrement importante en raison du statut réglementaire des zones humides, qui constituent généralement un des enjeux principaux des diagnostics environnementaux. 


\subsection{Principe de l'observation.}

Les végétaux produisent de la chlorophylle en absorbant la lumière visible du soleil (VIS). Seule une petite partie de ce rayonnement est utilisée dans ce mécanisme. Or le rayonnement solaire arrivant au sol comporte aussi des ultra-violets (UV) et des infrarouges proches (NIR) \footnote{Les rayonnements UV et NIR ne sont pas visibles par l'\oe il humain.} 

\figureIRD[Facteur_reflexion_chloro-carot]{Spectres-absorption-pigments-photosynthetiques.pdf}{Facteur d'absorption des pigments photosynthétiques dans la bande spectrale visible.}{1}

C'est essentiellement le rayonnement visible dans les bandes \og bleu \fg{} ($[400\,nm ,\, 450\,nm]$) et \og rouge \fg{} ($[630\,nm ,\, 680\,nm]$) qui contribue à la production chlorophyllienne.  
La partie non absorbée du rayonnement par la plante est transmise (la plante est \og transparente \fg{}) ou bien réfléchie (renvoyée vers l'environnement). 
Le rendement chlorophyllien étant nul dans la bande spectrale \og jaune-vert \fg ($[500\,nm ,\, 600\,nm]$) celle-ci est réfléchie par la plante (\ref{Facteur_reflexion_chloro-carot}). 
C'est de là que provient la couleur verte de ces végétaux. De même les rayonnements UV et IR ne participent pas à la production de chlorophylle. 
En revanche ils produiraient son échauffement s'ils étaient absorbés. 
Ainsi pour se protéger les plantes n'absorbent pas les rayonnements UV et IR émis par le soleil, qui sont alors transmis ou réfléchis. 
Les plantes sont totalement transparentes aux rayonnements ultraviolets.
En revanche elles réfléchissent une partie du rayonnement infrarouge.
C'est cette propriété qui va être exploitée dans le projet. 
Les résineux et les feuillus n'ont pas le même facteur de réflexion du rayonnement infrarouge. 
La figure \ref{Facteur_reflexion_vegetaux} montre les courbes typiques du facteur de réflexion spectral des résineux et des feuillus. 
Notons que les feuillus produisent aussi de la bêta-carotène et de la fucoxanthine qui absorbent le bleu et une partie du vert mais réfléchissent le rouge. 
C'est pour cette raison qu'à l'automne, lorsque les feuillus ne produisent plus de chlorophylle mais que la bêta-carotène est encore présente, ces espèces arborent des teintes allant du marron orangé à l'ocre rouge.

\figureIRD[Facteur_reflexion_vegetaux]{Facteur_reflexion_vegetaux.pdf}{Facteur de réflexion des résineux et des feuillus. Dans la bande spectrale visible $[400\,nm ,\,700\,nm]$ le comportement des deux espèces est très similaire. En revanche dans la bande spectrale infrarouge IR-A $[700\,nm ,\, 1.4\,\mu m]$ bien présente dans le spectre solaire les deux espèces se distinguent nettement.}{1}


On remarque que les feuilles des arbres réfléchissent le rayonnement autour de 550 nm qui correspond à la couleur verte perçue par l'\oe il humain. Mais les courbes étant identiques dans la partie visible du spectre (400 nm à 700 nm) on ne peut pas distinguer entre les différentes essences. En revanche les facteurs de réflexion dans le domaine infrarouge sont bien distincts entre résineux et feuillus. L'idée est d'exploiter cette propriété en produisant une image infrarouge  du couvert végétal
\footnote{Éventuellement combinée à une image dans le spectre visible.}. Sur cette image les feuillus devraient apparaître  \og plus clairs\fg{}  que les résineux. On pourra ensuite tracer une carte du couvert végétal. Signalons toutefois que selon l'époque de l'année et l'âge des sujets ces courbes peuvent évoluer rendant plus délicate l'interprétation des images.

\figureIRD[Facteur_absorption_vegetaux]{Par_action_spectrum.jpg}{Facteur d'absorption des végétaux et rendement de photosynthèse.}{0.5}


\subsection{Le matériel pour l'observation.}
Pour tracer une carte il faut des images \og vues du ciel \fg{}. Les professionnels utilisent des photos satellites (alt. 300km) ou prises d'avion (alt. 2000-3000 m). Mais comme expliqué plus haut, la résolution spatiale et temporelle n'est pas adaptée à des études fines de terrain. Dans le cadre de ce projet on envisage donc d'utiliser un petit drone qui permet d'obtenir avec une très grande souplesse des vues aériennes des zones à étudier (alt. 120 m) avec une très bonne résolution spatiale et une maîtrise totale des périodes de prise de vue.

La plupart des \og petits \fg{} drones existants ne permettent pas l'obtention simultanée de clichés en couleur naturelle et en IR. Or il est indispensable de disposer des deux types de clichés géoréférencés (IR pour interprétation de la végétation et couleur naturelle pour la réalisation de cartes standardisées). \\ 

Il existe des drones et des caméras multi-spectrales pour drones
\footnote{Par exemple DJI Phantom 4 multispectral, Parrot Sequoia, Micasense RedEdge.} . 
Il s'agit de caméras multi canaux  qui permettent d'acquérir, outre l'image visible classique, des bandes spectrales étroites dans le rouge, le vert, le bleu et dans une ou deux bandes dans la proche infrarouge. Il existe même des versions équipées de caméra thermiques \footnote{Par exemple la caméra Micasense Altum  possède un capteur Flir IR. Son poids est de 400g. Cependant il s'agit ici de rayonnement infrarouge thermique (LWIR) dont les longueurs d'onde sont dix fois plus grandes que celles du  proche infrarouge solaire.  Soulignons que ce matériel est très onéreux.}.
D'autres drones
sont équipés de deux caméras l'une dans le visible et l'autre dans l'infrarouge. 
Enfin la caméra de certains drones
\footnote{Par exemple le DJI Mavic Pro.2 équipé d'une très performante caméra Hasselblad. Son filtre \og hot mirror\fg{} est fixé directement sur le capteur. Des drones \og lourds\fg{} munis de nacelles peuvent embarquer des reflex numériques qui eux aussi peuvent être \og défiltrés \fg{}. }  
peut être rendue sensible dans le proche infrarouge. Cela nécessite de démonter
\footnote{Cette opération de \og défiltrage \fg est délicate. Elle peut être réalisée par des sociétés spécialisées.} le filtre interne du capteur numérique qui coupe le rayonnement infrarouge et de munir le système optique d'un filtre adapté
\footnote{Avec un jeu de filtres externes on peut choisir le type de bande spectrale IR. Avec un filtre passe haut IR $760\,\mathrm{nm}$ on obtiendra une bande large $[760, 1100]$. Il existe aussi des filtres à bande étroite, par exemple $850\,\mathrm{nm}\pm 25\mathrm{nm}$). On peut même remettre un filtre \og hot mirror\fg{} qui rétablira l'imagerie dans le spectre visible.}.  
De plus le drone n'enregistre alors qu'un seul type d'image selon le filtre dont il est équipé. \\  


L'idée du projet \textsf{IRDrone} est d'équiper un petit  drone de coût abordable 
\footnote{Ici un DJI Mavic Air.2  dont le poids n'excède pas 600 g.} d'une seconde  caméra. La caméra équipant le drone fournira des images en couleur dans le spectre visible (RGB). La deuxième caméra embarquée, que nous désignerons par caméra IR, devra produire des images dans la bande spectrale du proche infrarouge \og solaire \fg 
\footnote{Concrètement avec les capteurs CMOS disponibles la bande spectrale accessible s'étend de $720 \,\mathrm{nm}$ à $1100\,\mathrm{nm}$. Elle englobe les bandes étroites RedEdge et NIR (pour Near Infra Red)}.\\

La caméra NIR devra être légère pour ne pas perturber le vol du drone et ne pas diminuer son autonomie. Elle devra être d'un coût réduit ce qui conduit naturellement à se tourner soit vers des caméras de surveillance (\og vision de nuit\fg{}) soit vers des \og action cam \fg{} grand public. Les premières ne conviennent pas car elles offrent des performances médiocres en résolution. Les secondes sont conçues pour donner de belles images dans le domaine visible. Il faudra donc adapter l'optique par ablation du filtre \og hot mirror\fg{} et mise en place d'un filtre IR.  \\


Après chaque vol il sera nécessaire de faire un traitement informatique des images collectées (corrections géométriques, colorimétriques, calage et fusion d'images visibles et IR ...). Ce sera le rôle dévolu au code Python  IRDrone.py associé au projet \textsf{IRDrone} dont les fonctionnalités sont  exposées ci-après.


\section{Organisation d'une mission IRDrone}

\subsection{Préparation au bureau}

Avant le début de la partie active de la mission le pilote doit préparer soigneusement son vol. 
D'abord au bureau, en repérant les chemins d'accès, la topographie du terrain, les obstacles et la dimension de la zone à étudier. 
Cette dimension conditionne l'altitude de vol. On peut imaginer ne faire qu'un vol unique à haute altitude pour couvrir en une fois la zone, ou bien plusieurs vols successifs. Il faut garder à l'esprit que le drone DJI chargé de sa caméra infrarouge a une durée de vol limitée à 25 minutes. 
En pratique la phase de synchronisation prend 2 minutes. Avec une vitesse de $3\,\mathrm{m/s}$ le drone peut parcourir environ 4\,km, ce qui n'est pas énorme compte tenu de la nécessité de \og balayer \fg~la zone de survol afin que les clichés se superposent (si on souhaite faire un assemblage des images). 
On suppose aussi que l'opérateur a sa qualification de pilote de drone et les autorisations de survol.

\subsection{Préparatifs du vol sur le terrain}
Une fois sur le terrain le pilote installera la zone de décollage. Après balisage il positionne le stand de décollage et le système de synchronisation par mire ArUco avec son alimentation
\footnote{Que faire si on a oublié le dispositif de synchronisation ? On peut \og sauver \fg~la mission en déplaçant lentement un objet sous le drone. Il faudra ensuite estimer manuellement l'offset de synchronisation des deux caméras. Ça peut marcher\dots{} des fois !}. 
Un siège pliant est utile pour piloter en regardant le retour caméra afin de suivre au mieux le plan de vol. 
De même une petite manche à air pour apprécier la direction et la vitesse du vent est utile. 
Enfin, mesurer précisément la position du point de décollage avec un GPS en position fixe est appréciable\footnote{Voir les possibilités d'un GPS Garmin 750, qui a ce type de fonction.}. 
Avant chaque mission sur le terrain \textbf{le GPS du drone doit être calibré et sa caméra étalonnée} (direction de visée). 
Ces deux actions sont indispensables pour la réussite de l'exploitation des images par le logiciel IRDrone. 
Le pilote vérifie que les batteries sont chargées et que les cartes SD (vides) sont bien présentes sur le drone DJI et la caméra SJCam M20. 
Il est judicieux de prérégler les paramètres de la caméra du DJI (ceux qui peuvent l'être) avant le décollage. 
À ce stade le pilote peut déclencher l'hyperlapse de la caméra infrarouge (SJCam M20), positionner le drone sur son stand de décollage et enfin mettre sous tension le dispositif de synchronisation (dans la dernière version du dispositif, la mire ne se met à tourner que lorsque le drone a quitté le stand de décollage).

\subsection{Le vol de la mission}
Une fois ces préparatifs faits, la mission comporte 4 phases successives :
\begin{enumerate}
	\item \emph{Takeoff Point -- Synchronization} : synchronisation à proximité immédiate du point de décollage, à une altitude d'environ 2\,m. 
	Au début de cette phase la caméra NIR est déjà déclenchée en mode hyperlapse (fréquence de prise de vue 3\,s). 
	Pendant cette phase le pilote stabilise le drone au-dessus du dispositif de synchronisation, règle les paramètres de la caméra VIS et déclenche la prise de vue VIS en mode hyperlapse. Il reste au point fixe une dizaine de secondes.
	\item \emph{Approach} : déplacement vers la zone à survoler pour l'étude. 
	Le pilote prend rapidement de l'altitude. 
	Il se déplace à la vitesse maximum permise par le mode hyperlapse du DJI. 
	Une valeur de $3\,\mathrm{m/s}$ est un bon choix. Lors de la préparation de la mission il est fortement conseillé de choisir un point de décollage au plus près de la zone de survol.
	\item \emph{Hyperlapse Survey} : acquisition des images multispectrales au-dessus de la zone de survol (\emph{survey area}). 
	Le drone se déplace en gardant une orientation constante. 
	Il avance puis se déplace latéralement puis recule, etc. (voir figure~\ref{fig:4_step_of_the_mission}). 
	Lors du décollage il est conseillé d'orienter le drone dans le sens de déplacement principal au-dessus de la zone de survol. 
	À la sortie de la zone de survol le pilote peut interrompre l'enregistrement hyperlapse de la caméra VIS. 
	Ceci a deux avantages : cela évite de saturer la carte SD, les images étant en haute résolution, mais surtout cela libère la contrainte de vitesse du drone imposée par le mode hyperlapse.
	\item \emph{Return to Home} (RTH) : c'est la phase de retour au point de décollage. Ce retour peut se faire à vitesse maximum. Le mode RTH peut être enclenché pour un vol automatique.
\end{enumerate}

\figureIRD[fig:4_step_of_the_mission]{4_step_of_the_mission.pdf}{Les phases de vol d'une mission IRDrone sur le terrain.}{0.8} 

\subsection{Exploitation des images}

Une fois l'acquisition terminée, la phase suivante consiste à traiter numériquement les images des deux caméras à l'aide du logiciel IRDrone. 


\section{Objectif du logiciel IRDrone}
IRDrone.py est une application interactive et conviviale qui permet de générer des images multispectrales pour des études d'impact environnemental ou de botanique. 
Elle utilise pour cela les images prises lors des campagnes sur le terrain par deux caméras, l'une dans le spectre visible (\og VIS \fg) et l'autre dans le spectre du rayonnement dans le proche infrarouge (\og NIR \fg). 
Les deux caméras sont fixées sur un drone. 
Actuellement le système supporte un drone DJI Mavic Air 2 muni de sa caméra VIS montée sur un gimbal et une caméra SJCam M20 modifiée pour capturer des images dans le spectre du proche infrarouge (typiquement au-delà de 810 nm). 
L'application synchronise les horloges des deux caméras puis aligne les images VIS et NIR.
Le processus principal produit des images multispectrales alignées (4 bandes red, green, blue, nir au format tif). 
Un post-traitement permet d'orthorectifier et géoréférencer les images, soit unitairement, soit sous forme d'assemblage (mapping avec le code ODM ; figure~\ref{fig:img_mapping_1}).
\figureIRD[fig:img_mapping_1]{img_mapping_sample.png}{Assemblage à l'aide d'Open Drone Map (ODM) d'images dans le spectre VIS. La visualisation est faite avec le logiciel QGIS. Les courbes de niveau ont une équidistance de 1\,m. Les points rouges indiquent l'emplacement des clich\'es et le numéro de la prise de vue.}{0.8}
L'exploitation peut être réalisée avec le logiciel QGIS afin d'obtenir des images spécialisées (indice NDVI, figure~\ref{fig:img_NDVI_1} ; image VIR nir/red/green,
figure~\ref{fig:img_VIR_1}). 


\section{Fonctionnalités principales du logiciel IRDrone}
Au lancement, une petite fenêtre \og Loading, please wait \fg{} s'affiche pendant le chargement (quelques dizaines de secondes), puis la fenêtre principale s'ouvre (chapitre~\ref{ch:home}). Elle liste les missions du dossier de travail et, une fois une mission ouverte, ses sept étapes avec leur avancement. Un clic sur une étape la lance ; les étapes 1 à 4 s'enchaînent d'elles-mêmes.

L'exploitation d'une mission comporte 7 étapes :
\begin{enumerate}
	\item \textbf{Mission} : les images des deux caméras sont trouvées sur les cartes SD, copiées dans un nouveau dossier de mission, et les paramètres de la mission sont enregistrés.
	Le nom du dossier de mission est généré automatiquement en fonction de la localisation et de la date
	\footnote{\label{fn:MissionFolderName} Par défaut le nom du dossier est FLY-YYYYMMDD-hhmm-location. Il peut être modifié une fois pour toutes à la création de la mission dans le champ \og Location \fg. Par exemple si le champ \og Location \fg~a la valeur Blassac on peut préciser \og Blassac Bourg \fg.}.
	\item \textbf{Ligne de temps du vol} : les images de synchronisation (mire ArUco au décollage) et celles du relevé au-dessus de la zone d'étude sont détectées automatiquement ; l'utilisateur vérifie.
	\item \textbf{Synchronisation} : le décalage entre les horloges des deux caméras est mesuré et les images VIS et NIR sont appariées.
	\item \textbf{Alignement}, c\oe ur d'IRDrone : chaque paire d'images est alignée pour obtenir 4 couches spectrales superposées.
	\item \textbf{Revue} des images alignées (chapitre~\ref{ch:review}) : chaque image est vérifiée sur une ligne de temps, les mal alignées sont signalées (écartées de la carte) ou réalignées à la main.
	\item \textbf{Cartographie et géoréférencement} : géoréférencement, orthorectification, assemblage en orthophoto (mapping), modèle de terrain et courbes de niveau.
	\item \textbf{Archivage} de la mission.
\end{enumerate}

\figureIRD[fig:Air_Mission_folder]{Air_Mission_folder.png}{Le dossier principal \texttt{Air-Mission} avec deux missions en cours de traitement (ici Chalivoy-Milon et Labout) et une mission archivée (ici Blassac). Sous Windows les icônes sont attribués automatiquement.}{0.6}

\subsection*{Aide en ligne}\label{sec:help}
\begin{itemize}
	\item \textbf{Info-bulles.} Laisser le pointeur de la souris quelques instants sur un bouton (ou une case à cocher) affiche une info-bulle qui résume son action.
	\item \textbf{Touche \textsf{F1}.} La touche \textsf{F1}, le pointeur de la souris sur un bouton, ouvre ce manuel à la section qui décrit ce bouton. Sur une étape de la page d'accueil, elle ouvre le chapitre de l'étape. Ailleurs dans une fenêtre, elle ouvre le chapitre de la fenêtre, comme le bouton \bouton{?}.
	\item \textbf{Bouton \bouton{?}.} Chaque fenêtre a un bouton \bouton{?} qui ouvre ce manuel au chapitre de la fenêtre. Dans les fenêtres qui se pilotent au clavier (ligne de temps, revue), son info-bulle rappelle les touches. Il est en bas à gauche ; le bouton \bouton{Close} des fenêtres d'étape et de la revue est en bas à droite ; aux étapes 2 à 5, le bouton qui passe à l'étape suivante (\bouton{Continue with \ldots ▶}) est à sa droite et ne lance jamais un calcul tout seul. Les actions d'une étape (\bouton{Run alignment}, \bouton{Build the map\ldots}, \ldots) sont dans la fenêtre, au-dessus de la barre de progression ; la principale est bleu foncé, les autres blanches, et celle qui est irréversible (l'archivage) rouge. Les étapes~3 à~6 affichent en haut le nom de l'étape et, dessous, la mission. Dans la fenêtre principale aussi, en bas à gauche, il ouvre le manuel entier.
	\item \textbf{Boutons \bouton{↺} et \bouton{Close} des fenêtres qui modifient quelque chose} (détails de la mission, ligne de temps, réglage du décalage, réglages de l'étape~4). Fermer garde ce qui a été changé : dès qu'une valeur change, \bouton{Close} devient \bouton{Save and close} ; fermer la fenêtre par sa barre de titre ou avec \textsf{Échap} fait de même. Le bouton \bouton{↺}, à côté de \bouton{?}, n'apparaît que s'il y a quelque chose à remettre : il remet les valeurs et la fenêtre reste ouverte.
\end{itemize}
Le manuel s'ouvre dans le lecteur PDF par défaut (navigateur, Acrobat), à la bonne page quand le lecteur le permet.

\section{Organisation numérique d'une mission IRDrone}

%\todo{Dossier \texttt{FLY-YYYYMMDD-HHmm-localit\'e}, sous-dossiers \texttt{ImgIRdrone}, \texttt{mapping\_MULTI}, etc.}
On suppose ici que votre ordinateur sous Windows comporte un dossier \texttt{C:/Air-Mission}.
Ce dossier de travail est choisi au premier lancement de l'application ; il est rappel\'e en bas de la fen\^etre
principale et peut \^etre chang\'e avec le bouton \og engrenage \fg{} plac\'e \`a c\^ot\'e.

Chaque mission est placée dans un dossier \texttt{C:/Air-Mission/FLY-YYYYMMDD-hhmm-location}. 

Ce dossier comporte les sous-dossiers :
\begin{enumerate}
	\item \texttt{AerialPhotography}
	      \begin{itemize}
	      	\item \texttt{VIS}
	      	  \begin{itemize}
	      	  	\item \texttt{VIS\_XXXX.dng}. Images de la caméra DJI transférées,
	      	  	\item \texttt{VIS\_XXXX.exif}. Fichiers \og compagnons \fg~des images contenant des données exif,
	      	  	\item \texttt{time\_line.json}. Fichier qui donne la chronologie des images par rapport à l'horloge de la caméra VIS,
	      	  	\item \texttt{transfer\_info\_VIS\_dng.json}. Fichier qui résume les choix de l'étape~2,
	      	  	\item \texttt{VIS\_XXXX\_RawTherapee.tif}. Images\footnote{\label{fn:rawtherapee_1}Ces fichiers très volumineux sont indispensables au traitement numérique des images. Ils sont générés une seule fois à la première utilisation de l'étape~4. Lorsque la mission est entièrement traitée et doit être archivée, ils peuvent être supprimés.} créées à l'étape~4.
	      	  \end{itemize}
	      		
	      	\item \texttt{NIR}
	      	\begin{itemize}
	      		\item \texttt{NIR\_XXXX.dng}. Images de la caméra infrarouge transférées et converties (RAW propriétaire vers dng), 
	      		\item \texttt{NIR\_XXXX.exif}. Fichiers \og compagnons \fg~des images contenant des données exif,   
	      		\item \texttt{time\_line.json}. Fichier qui donne la chronologie des images par rapport à l'horloge de la caméra NIR,
	      		\item \texttt{transfer\_info\_NIR\_dng.json}. Fichier qui résume les choix de l'étape~2,
	      		\item \texttt{NIR\_XXXX\_RawTherapee.tif}. Images\textsuperscript{\ref{fn:rawtherapee_1}} créées à l'étape~4.
	      		\end{itemize}
      		\item \texttt{shoot\_points.json}. Dictionnaire des \emph{shoot points} (prise de vue avec comme référence l'image VIS et l'image NIR la plus proche de l'instant de la prise de vue VIS). Il est créé à l'étape~3.
	    \end{itemize}
	\item \texttt{cameras}
	         \begin{itemize}
	         	\item Non utilisé actuellement.
	         \end{itemize}
	\item \texttt{FlightAnalytics}. C'est un dossier technique.
			\begin{itemize}
				\item \texttt{mission\_parameters.json}. Fichier contenant les paramètres de la mission (modifiable avec la fenêtre \og Mission details \fg, §\ref{btn:step1_details}).
				\item \texttt{flight\_timeline.json}. Phases du vol confirmées à l'étape~2.
				\item \texttt{fly\_vertical\_profil.png}. Graphe qui permet de visualiser les altitudes de prise de vue par rapport au sol.
				\item \texttt{Angles analysis.png},  \texttt{Angles analysis dispersion.png}, \texttt{Angles analysis spectrum.png}. Graphes d'analyse de la qualité du vol. Le premier de ces trois graphes permet par exemple de juger si l'attitude du drone pendant la mission est restée dans les spécifications de vol.
			\end{itemize}
	\item \texttt{ImgIRdrone}
		\begin{itemize}
		\item Pour chaque shoot point de numéro XXXX on trouve : 
			\begin{itemize}
				\item \texttt{VIS\_XXXX\_REDO.bat},
				\item \texttt{VIS\_XXXX.jpg},
				\item \texttt{VIR\_XXXX\_local.jpg},
				\item \texttt{VIS\_XXXX\_motion\_model.json},
				\item \texttt{VIS\_XXXX\_motion\_model.npy}.
			\end{itemize}
		\item \texttt{review.json}. Mesures et points vérifiés de la revue des images alignées (chapitre~\ref{ch:review}).
		\item \texttt{shoot\_points.json}. La liste des shoot points traités avec un indicateur de réussite. Attention : cet indicateur suppose par défaut que l'alignement est correct. L'utilisateur dispose d'une procédure pour examiner les alignements et les déclarer, le cas échéant, défaillants.
		\item \texttt{imgMultispectral}. Dossier contenant les quadruplets d'images multispectrales au format tif. Ici ces images ne sont pas géoréférencées. Pour chaque shoot point XXXX :
			\begin{itemize}
				\item \texttt{IRD\_XXXX\_1.tif}. Image de la couche rouge. 
				\item \texttt{IRD\_XXXX\_2.tif}. Image de la couche verte.
				\item \texttt{IRD\_XXXX\_3.tif}. Image de la couche bleue.
				\item \texttt{IRD\_XXXX\_4.tif}. Image de la couche infrarouge proche.
				\item \texttt{IRD\_XXXX\_1.exif}. Fichier \og compagnon \fg~contenant des données exif\footnote{Les données exif sont aussi présentes dans les images \texttt{IRD\_XXXX\_Y.tif}. Toutefois l'extraction de certaines de ces données exif nécessite l'utilisation d'un logiciel tiers spécifique comme ExifTool. C'est une opération coûteuse qu'on ne réalise qu'une fois. On peut donc considérer les fichiers exif compagnons comme des données qui ont été mises en cache.}.
			\end{itemize}
		\item \texttt{VIS\_XXXX\_align\_traces\_manual}. Dossier contenant toutes les informations et images suite à une opération d'alignement manuel des images VIS et NIR du shoot point XXXX.
		\end{itemize}
	\item \texttt{mapping\_MULTI}. Dossier de travail de l'étape~6 (mapping).
	\item \texttt{Synchro}. Aperçus des images et détections de la mire préparés pour l'étape~2 (\texttt{Thumbnails}, un cache recréé si besoin) ; fichiers de l'étape~3 : images de travail \texttt{*\_RawTherapee.jpg} (sous-dossiers \texttt{VIS} et \texttt{NIR}), images de contrôle de la détection des mires (\texttt{Check\_ARUCO}) et graphiques \texttt{check\_time\_*.png}.
\end{enumerate}

\section{Encha\^inement des \'etapes}
%\todo{Ordre conseill\'e, pr\'erequis de chaque \'etape.}
Vous devez avoir sur un support numérique (carte SD, disque dur) les images des deux caméras prises lors du vol de la mission. Il est fortement conseillé de n'avoir qu'une mission par carte SD. Les images du DJI sont au format dng et stockées dans un dossier \texttt{DCIM/HYPERLAPSE}. Les images de la caméra SJCam M20 sont au format RAW, doublées d'une image jpeg. Comme les images RAW ne possèdent pas de vignette, on peut les choisir en regardant l'image jpeg associée. Les images sont stockées dans un dossier \texttt{DCIM}.

Pour une nouvelle mission il faut suivre dans l'ordre les étapes 1 à 6. 
Cependant les étapes 2 à 6 peuvent être lancées individuellement, à condition que le dossier de la mission comporte au minimum les éléments produits par les étapes qui la précèdent. C'est particulièrement utile pour les étapes 4 à 6.

Les étapes se lancent depuis la fenêtre principale, décrite au chapitre suivant.

%==========================================================
\chapter{La fenêtre principale}\label{ch:home}
La fenêtre principale a deux pages : la liste des missions, puis la mission ouverte. En bas, une barre rappelle le dossier de travail et ouvre ce manuel.

\section{Liste des missions}\label{btn:home_missions}
Sans mission ouverte, la fenêtre propose (figure~\ref{fig:home_welcome}) :
\begin{description}
	\item[\bouton{New mission}] crée une mission à partir des cartes SD du drone et de la caméra infrarouge (étape~1, chapitre~\ref{ch:step1}).
	\item[\bouton{Open a mission}] ouvre une mission déjà créée : choisir son dossier \texttt{FLY-YYYYMMDD-hhmm-location}.
	\item[Missions récentes] les missions du dossier de travail, les plus récentes en premier, avec leur avancement : un point vert par étape faite. Un clic ouvre la mission.
\end{description}
Le bouton \bouton{All missions} de la page de la mission revient à cette liste.

\figureIRD[fig:home_welcome]{home_welcome.png}{La liste des missions : nouvelle mission, ouvrir une mission, missions récentes et leur avancement.}{0.85}

\section{La mission ouverte}\label{btn:home_steps}
La page de la mission (figure~\ref{fig:home_mission}) montre :
\begin{itemize}
	\item en haut, le nom de la mission (un clic ouvre son dossier dans l'explorateur) et le lieu, la date et le pilote ;
	\item les sept étapes, de gauche à droite : coche verte et date pour une étape faite, cadre bleu pour l'étape suivante, gris pour les autres. \textbf{Un clic sur une étape la lance}, même déjà faite (pour la refaire). Deux exceptions :
	\begin{itemize}
		\item l'étape~1 faite ouvre seulement la fenêtre \og Mission details \fg~(§\ref{btn:step1_details}) : les cartes SD ne sont pas relues ;
		\item l'étape~3 faite montre le résultat enregistré (décalage des horloges, courbes alignées) sans rien recalculer (§\ref{sec:step3_chained}).
	\end{itemize}
	La touche \textsf{F1} sur une étape ouvre son chapitre de ce manuel ;
	\item l'étape suivante, expliquée en une phrase, avec son bouton (§\ref{btn:home_continue}) ;
	\item quelques chiffres : nombre d'images, de paires, d'images alignées, de points signalés mal alignés ;
	\item à droite, la carte du vol ou une image alignée (§\ref{btn:home_map}).
\end{itemize}

\figureIRD[fig:home_mission]{home_mission.png}{La mission ouverte : les étapes (un clic lance une étape), étape suivante, chiffres de la mission et carte du vol sur l'orthophoto et les photos aériennes IGN.}{0.9}

\section{Bouton \bouton{Continue with Step N}}\label{btn:home_continue}
Lance l'étape suivante. Les étapes 1 à 4 s'enchaînent d'elles-mêmes :
\begin{enumerate}
	\item à la fin de l'étape~1, \bouton{Flight timeline} ouvre l'étape~2 ;
	\item \bouton{Continue with synchronisation} de l'étape~2 (ou \bouton{Continue with Step 3}) ouvre l'étape~3 ; \bouton{Synchronise the cameras} la lance : détection des mires, décalage des horloges, puis une question, \og Do the two curves overlap? \fg~; \bouton{Yes} apparie les images ;
	\item \bouton{Continue with image processing} ouvre l'étape~4. L'alignement, long, ne démarre jamais tout seul : il se lance avec \bouton{Run alignment} et s'arrête à tout moment avec \bouton{Stop} ;
	\item à la fin, \bouton{Continue with the review} ouvre l'étape~5, la revue des images alignées (elle s'ouvre d'elle-même quand l'alignement a été enchaîné depuis l'étape~3) ;
	\item \bouton{Continue with mapping} de la revue ouvre l'étape~6, où \bouton{Build the map\ldots} lance la carte ;
	\item une fois la carte construite, \bouton{Continue with archiving} ouvre l'étape~7, qui indique l'espace libéré par l'archivage ; rien n'est supprimé avant votre confirmation.
\end{enumerate}
Il reste donc à vérifier les phases du vol, la superposition des courbes de synchronisation, puis le résultat. Une fois toutes les étapes faites, le bouton devient \bouton{Archive the mission} (étape~7).

\section{Carte du vol}\label{btn:home_map}
La carte est disponible dès l'étape~1 : les positions sont lues dans les EXIF des images visibles. Elle montre :
\begin{itemize}
	\item la trajectoire du drone et un point par image, coloré selon les phases confirmées à l'étape~2 (violet : synchronisation, vert : relevé) ;
	\item les images de la carte, et elles seules : la trajectoire les relie, le vol complet restant en pointillés discrets, et tant qu'aucune orthophoto n'est affichée leur emprise au sol est dessinée en vert : on voit le recouvrement de la carte. Quelles images :
	\begin{itemize}
		\item avant l'étape~4, les images retenues par le \emph{best mapping} (étapes~2 et~4, section~\ref{btn:step2_best_mapping}), s'il est choisi ;
		\item une fois l'étape~4 faite, les points exactement comme dans la revue des images alignées (chapitre~\ref{ch:review}) : chaque point de prise de vue dans la couleur de son statut (vert aligné, orange à vérifier, rouge signalé\ldots), les points que l'étape~4 n'a pas alignés en petits points gris ; une image alignée ou réalignée depuis la revue apparaît aussitôt ;
		\item l'orthophoto de l'étape~6, une fois construite, est dessinée sous ces points (les mêmes points que la revue, pas seulement les images qu'elle a utilisées) ;
	\end{itemize}
	les images signalées (\bouton{Flag}, revue des images alignées, chapitre~\ref{ch:review}) sont de petits points rouges ; en survoler un montre son emprise en pointillés rouges : on voit le trou qu'il laisse dans le recouvrement ; signalées après l'étape~6, elles sont encore dans l'orthophoto affichée et en seront retirées quand l'étape~6 sera relancée (\texttt{left out when Step 6 runs again}). Le nombre d'images et d'images signalées est écrit en bas à droite. \bouton{Layers} > \bouton{Only the images of the map (and the flagged ones)} (\bouton{Only the images of the best mapping} avant l'étape~4) remontre toutes les images ;
	\item l'orthophoto de l'étape~6, une fois le mapping fait : dès que l'étape~6 l'a construite, la fenêtre principale l'affiche sur la carte (même si \bouton{Layers} > \bouton{Hidden} était choisi avant) ;
	\item en fond, des photos aériennes : IGN pour la France, Esri World Imagery ailleurs, ou OpenStreetMap.
\end{itemize}
Survoler un point montre la vignette de son image ; un clic sur un point montre l'emprise au sol de l'image (le trait épais est le haut de l'image). Molette : zoom ; glisser : déplacement ; double-clic ou bouton \bouton{⤢} : tout le vol.

Le menu \bouton{Layers} choisit le fond (IGN, Esri, OpenStreetMap ou aucun) et l'affichage de l'orthophoto, en couleurs naturelles ou en fausses couleurs infrarouges. Le fond vient d'internet : hors ligne, la carte reste utilisable (vol, emprises, orthophoto), seul le fond manque, et rien n'est bloqué. Les tuiles déjà vues sont gardées dans un cache (\texttt{\textasciitilde/.cache/irdrone} sous Linux, \texttt{\%LOCALAPPDATA\%\textbackslash IRDrone} sous Windows) et s'affichent ensuite hors ligne.

L'emprise est calculée avec la hauteur au-dessus du sol (modèle de terrain de l'étape~6, ou altitudes IGN) et le cap de l'image ; une fois le mapping fait, le cap vient d'OpenDroneMap, plus juste que la boussole du drone (souvent fausse de 5 à 10\textdegree).


\section{Barre du bas : dossier de travail et manuel}\label{btn:home_workdir}
Sur la page de la liste des missions, la barre du bas rappelle le dossier de travail (le dossier qui contient les missions, par défaut \texttt{C:/Air-Mission}) et la place libre sur son disque ; un clic sur le dossier l'ouvre dans l'explorateur. La roue dentée \bouton{⚙} permet de changer de dossier de travail : la liste des missions est alors relue. Elle donne aussi une option avancée, décochée par défaut : saisir à la main le décalage d'horloge des caméras à l'étape~3, pour les vols sans cible ArUco (§\ref{btn:step3_clock_offset}). Le dossier est choisi au premier lancement ; la variable d'environnement \texttt{IRDRONE\_WORKING\_DIR}, si elle est définie, reste prioritaire. Sur la page d'une mission, la barre du bas indique à droite la taille totale du dossier de la mission (\texttt{Mission size}), mise à jour à chaque étape terminée et toutes les minutes.

\bouton{?}, à gauche, ouvre ce manuel (§\ref{sec:help}).

\section{Rapport de diagnostic}\label{btn:home_diagnostic}
Quand quelque chose ne marche pas sur un ordinateur, le rapport de diagnostic rassemble dans un seul fichier zip tout ce qu'il faut pour comprendre le problème ; il suffit d'envoyer ce zip aux développeurs. Il se crée de deux façons :
\begin{itemize}
  \item dans les options (roue dentée \bouton{⚙} en bas de la fenêtre principale, §\ref{btn:home_workdir}), bouton \bouton{Create a diagnostic report…} : on choisit une mission, ou l'installation seule ;
  \item sans ouvrir le logiciel, avec \texttt{irdrone\_diagnostic.bat} (à la racine du dossier IRDrone) : double-clic pour l'installation seule, ou glisser-déposer le dossier d'une mission (\texttt{FLY-...}) sur ce fichier.
\end{itemize}
Le zip (\texttt{irdrone\_diagnostic\_AAAAMMJJ\_hhmmss.zip}) est écrit dans le dossier de travail. Il contient :
\begin{itemize}
  \item \texttt{report.txt} : un résumé de ce qui semble anormal, puis l'installation (version d'IRDrone, Python utilisé, système, page de code Windows, versions des modules Python dont OpenCV, chemins et versions d'ExifTool, RawTherapee et du convertisseur SJCam, dossier de travail, place libre) et deux auto-tests (détection ArUco d'une cible parfaite ; lecture et écriture d'images et d'EXIF dans un dossier accentué) ;
  \item pour une mission : les étapes faites, les fenêtres de l'étape~2, les lignes de temps, le décalage et le nombre de paires de l'étape~3 ; chaque image de synchronisation est développée à nouveau avec RawTherapee dans un dossier temporaire et les marqueurs ArUco y sont cherchés, ainsi que sur l'image en cache (\texttt{Synchro}) utilisée par l'étape~3. Quelques images développées sont jointes. Cela prend environ une minute ; la mission n'est pas modifiée ;
  \item \texttt{report.json} (tous les détails) et les journaux des dernières sessions : chaque lancement du logiciel recopie la console dans \texttt{logs\textbackslash irdrone\_AAAAMMJJ\_hhmmss.log}, à côté du fichier des réglages (les dix derniers sont gardés).
\end{itemize}
Rien n'est envoyé automatiquement.

%==========================================================
\chapter{Step 1 -- D\'efinir la mission}\label{ch:step1}
\section{La fenêtre \og New mission \fg}\label{sec:step1_new_mission}
Le Step~1 cr\'ee le dossier de la mission et y range les images du vol. Il s'ouvre avec \bouton{New mission} dans la fenêtre principale et se fait dans une seule fenêtre, en trois pages : \textbf{Images}, \textbf{Mission}, \textbf{Copy}. L'avancement est rappelé en haut de la fenêtre. Les cartes SD sont seulement lues : rien n'y est écrit ni effacé.

\subsection{Type de mission : multispectrale ou visible seul}\label{btn:step1_mission_type}
Le haut de la page \og Images \fg~choisit le type de la mission (le dernier choisi est proposé) :
\begin{itemize}
	\item \bouton{Multispectral (visible + infrared)} : le drone (visible, DNG) et la caméra infrarouge (SJCam, RAW). Tout ce que décrit ce manuel : synchronisation, alignement, carte en quatre bandes (rouge, vert, bleu, proche infrarouge), VIR et NDVI ;
	\item \bouton{Visible only (drone DNG)} : les images du drone seules, sans la caméra infrarouge. La carte infrarouge disparaît, la copie ne prend que les images visibles et les étapes suivantes s'adaptent (section~\ref{sec:visible_only}) : phase de relevé seule à l'étape~2, pas d'étape~3, les images sont \emph{développées} à l'étape~4 au lieu d'être alignées, la carte de l'étape~6 est en couleurs (rouge, vert, bleu).
\end{itemize}
Le type ne se change plus une fois la mission créée.

\subsection{Page \og Images \fg}\label{btn:step1_images}
Les cartes SD insérées sont parcourues (tout volume qui contient un dossier \texttt{DCIM}) ; on y trouve les images visibles (DNG du drone) et infrarouges (RAW de la SJCam), figure~\ref{fig:step1_new_mission_images}. Chaque caméra a sa carte :
\begin{itemize}
	\item le dossier trouvé, le nombre d'images, leur taille, la date et l'heure de la première et de la dernière ;
	\item deux images : \textbf{au décollage} et \textbf{en vol} (aux deux tiers des images, en général pendant le relevé). Pour l'infrarouge, c'est le JPEG que la SJCam enregistre à côté de chaque RAW. Ces images permettent de vérifier d'un coup d'\oe il que la carte est la bonne ;
	\item le bouton \bouton{Choose folder…} (\bouton{Change…} une fois des images trouvées) choisit un autre dossier, par exemple une copie des cartes sur le disque. Un dossier choisi ainsi n'est jamais remplacé par une carte trouvée ensuite.
\end{itemize}
Tant qu'il manque les images d'une caméra, la fenêtre cherche les cartes SD toutes les 3~secondes : une carte insérée après l'ouverture de la fenêtre est trouvée d'elle-même, il suffit de l'insérer.

L'image de décollage est la première image visible ; le lien \og change… \fg~de la carte visible en choisit une autre. Elle donne la date, l'heure, la position et l'appareil photo de la mission.
Un bandeau signale un écart de plus de deux minutes entre les horloges des deux caméras : il est normal, l'étape~3 le mesure précisément. Un écart de plusieurs jours fait douter que les deux cartes viennent du même vol.

\figureIRD[fig:step1_new_mission_images]{step1_new_mission_images.png}{Page \og Images \fg : images visibles et infrarouges trouvées sur les cartes SD, au décollage et en vol.}{0.8}

\subsection{Page \og Mission \fg}\label{btn:step1_mission}
La date, l'heure et la position viennent de l'image de décollage ; l'altitude du sol et l'adresse sont demandées à l'IGN\footnote{\label{fn:ign}IGN : Institut national de l'information g\'eographique et foresti\`ere.} et à OpenStreetMap pendant que l'on remplit la description et le pilote (figure~\ref{fig:step1_new_mission_details}). Le nom du dossier de la mission (\texttt{FLY-YYYYMMDD-hhmm-location}) se met à jour avec le champ \og Location \fg.

\bouton{Cameras and filter} déplie les réglages des caméras : intervalle de prise de vue et décalage d'horloge de chaque caméra, fabricant et modèle de la caméra infrarouge, fabricant et bande du filtre. Le pilote, son identifiant, les caméras et le filtre sont mémorisés pour la mission suivante.

\figureIRD[fig:step1_new_mission_details]{step1_new_mission_details.png}{Page \og Mission \fg : lieu trouvé à partir de l'image de décollage, nom du dossier, description, pilote, caméras.}{0.8}

\subsection{Page \og Copy \fg}\label{btn:step1_copy}
Le dossier est créé, les images visibles copiées, les images infrarouges converties de RAW en DNG, les EXIF lus et l'altitude du sol calculée (détails au §\ref{btn:step1_validate}). Pendant la conversion des images infrarouges, l'étape~2 est préparée en arrière-plan (vignettes et détection de la mire, visible puis infrarouge) : la ligne de temps s'ouvre ensuite sans attendre. La fenêtre ne peut pas être fermée pendant la copie. À la fin, \bouton{Flight timeline} ouvre l'étape~2 ; si la préparation n'est pas finie, l'étape~2 la reprend là où elle en était.

Le fichier \texttt{mission\_parameters.json} contient les mêmes informations qu'avec le formulaire classique.

\subsection{Boutons du bas de la fenêtre}\label{btn:step1_wizard}
\begin{description}
	\item[\bouton{Next}] (page Images) passe à la page Mission ; il faut les images des deux caméras. \bouton{Create the mission and copy the images} (page Mission) lance la copie ; \bouton{Flight timeline} (page Copy) ouvre l'étape~2.
	\item[\bouton{Back}] revient à la page Images ; rien n'est perdu.
	\item[\bouton{Cancel}] ferme la fenêtre sans créer la mission.
	\item[\bouton{Classic form…}] (page Images) ouvre l'ancien Step~1 (§\ref{sec:step1_classic}).
	\item[\bouton{?}] ouvre ce chapitre.
\end{description}

\section{Modifier une mission : la fenêtre \og Mission details \fg}\label{btn:step1_details}
Une fois l'étape~1 faite, un clic sur l'étape~1 de la fenêtre principale n'ouvre plus la création de mission, mais la fenêtre \og Mission details \fg~(figure~\ref{fig:step1_mission_details}). Elle modifie les informations de la mission sans relire les cartes SD : lieu, description, pilote et son identifiant, intervalle et décalage d'horloge des caméras, caméra infrarouge, filtre. La date, l'heure, la position et la caméra visible viennent de l'image de décollage et ne changent pas.

Dès qu'une information change, \bouton{Close} devient \bouton{Save and close} : il écrit les changements dans \texttt{FlightAnalytics/mission\_parameters.json} (fermer la fenêtre fait de même) ; l'ancien fichier est gardé en \texttt{mission\_parameters.json.bak}. Le dossier de la mission garde son nom, même si le lieu change. \bouton{↺}, à côté de \bouton{?}, apparaît alors : il remet les informations telles qu'à l'ouverture de la fenêtre ; la fenêtre reste ouverte.

\figureIRD[fig:step1_mission_details]{step1_mission_details.png}{La fenêtre \og Mission details \fg : les images sont déjà dans la mission, seules ces informations changent.}{0.55}

\section{Formulaire classique}\label{sec:step1_classic}
\legacybegin{Afficher le formulaire classique (ancienne étape 1)}
Le bouton \bouton{Classic form…} de la page \og Images \fg~ouvre l'ancien Step~1, décrit ci-dessous. Il se d\'eroule en deux fen\^etres successives :
\begin{enumerate}
  \item \textbf{\og Choose the image taken by the drone during takeoff \fg} : on choisit la premi\`ere image prise par le drone
        au d\'ecollage. Elle fournit la date, l'heure, la position GPS et l'appareil photo utilis\'e.
  \item \textbf{\og Create a mission \fg} : on v\'erifie et on compl\`ete les informations de la mission, puis on lance la cr\'eation
        du dossier et le transfert des images VIS et NIR.
\end{enumerate}
\`A la fin, la mission devient la mission chargée de la fen\^etre principale, qui propose l'étape~2.

\paragraph{Pr\'erequis.}
\begin{itemize}
  \item Une connexion Internet : l'altitude du sol est demand\'ee \`a l'IGN\textsuperscript{\ref{fn:ign}}
  (France) ou \`a OpenTopoData (hors de France),
        et l'adresse du lieu \`a OpenStreetMap.
  \item Les images du drone (visible) au format DNG, et les images de la cam\'era proche infrarouge (SJCam M20, fichiers RAW)
        copi\'ees sur le disque de l'ordinateur.
  \item Un dossier de travail (par d\'efaut \texttt{C:/Air-Mission}), choisi au premier lancement et modifiable
        avec le bouton \og engrenage \fg{} en bas de la fen\^etre principale. La variable d'environnement
        \texttt{IRDRONE\_WORKING\_DIR}, si elle est d\'efinie, reste prioritaire.
\end{itemize}

% \figureIRD{step1_load_takeoff_1.png}{Premi\`ere fen\^etre du Step 1 avant chargement de l'image de d\'ecollage.}{0.8}

%----------------------------------------------------------
\section{Bouton \bouton{Load take-off image}}\label{btn:step1_load_takeoff}
Ouvre un explorateur de fichiers pour choisir l'image prise par le drone au d\'ecollage.
Seuls les fichiers \texttt{.dng} de la cam\'era visible sont accept\'es.

Une fois l'image choisie, le logiciel :
\figureIRD{step1_load_takeoff_2.png}{Premi\`ere fen\^etre du Step 1 apr\`es chargement de l'image de d\'ecollage.}{0.5}
\begin{enumerate}
  \item affiche l'image (le d\'ecodage du fichier DNG est assez long : la barre de progression indique l'avancement) ;
  \item lit les donn\'ees EXIF : date et heure de prise de vue, position GPS, fabricant, mod\`ele et num\'ero de s\'erie de l'appareil ;
  \item interroge l'IGN (ou OpenTopoData hors de France) pour obtenir l'altitude du sol \emph{au-dessus du niveau de la mer} ;
  \item interroge OpenStreetMap pour obtenir le lieu-dit, la commune, le code postal, le d\'epartement, la r\'egion et le pays.
\end{enumerate}
Ces informations s'affichent dans le cadre gris sous l'image. Lorsqu'elles existent, le lieu-dit, le code postal, la région et le pays sont aussi enregistrés dans le fichier
\texttt{FlightAnalytics\textbackslash mission\_parameters.json} sous les clés \texttt{Place name}, \texttt{Postal code}, \texttt{Region} et \texttt{Country}
(la commune est enregistrée sous la clé \texttt{Location}). Une information introuvable n'est pas enregistrée : la clé correspondante n'existe pas.
\begin{center}
\fbox{\parbox{0.9\linewidth}{\small\textbf{Attention.} L'altitude enregistr\'ee dans l'image par le drone est mesur\'ee \emph{par rapport
au point de d\'ecollage}, et non par rapport \`a la mer. C'est pourquoi l'altitude du sol est demand\'ee \`a un service externe tel que l'IGN\textsuperscript{\ref{fn:ign}}.}}
\end{center}


\paragraph{Si l'image ne se charge pas.} V\'erifiez la connexion Internet et que le fichier est bien un \texttt{.dng} du drone.
Le bouton \bouton{Next ▶} reste inactif tant qu'aucune image n'a \'et\'e charg\'ee correctement.

%----------------------------------------------------------
\section{Boutons \bouton{Cancel}, \bouton{Next ▶} et \bouton{◀ Back}}\label{btn:step1_navigation}
\begin{description}
  \item[\bouton{Next ▶}] n'est disponible qu'apr\`es le chargement de l'image de d\'ecollage. Il ouvre la fen\^etre \og Create a mission \fg.
  \item[\bouton{Cancel}] ferme la fen\^etre de l'image de d\'ecollage sans rien enregistrer.
  \item[\bouton{◀ Back}] (fen\^etre \og Create a mission \fg) revient \`a la fen\^etre pr\'ec\'edente ;
        l'image de d\'ecollage d\'ej\`a charg\'ee est conserv\'ee.
\end{description}

%----------------------------------------------------------
\section{La fen\^etre \og Create a mission \fg}\label{btn:step1_form}
Les champs sont remplis automatiquement \`a partir de l'image de d\'ecollage. La touche \textsf{Entr\'ee} passe d'un champ au suivant.
Les champs gris sont remplis automatiquement et ne sont pas modifiables.

\begin{center}
\begin{tabular}{@{}p{3.6cm}p{2.4cm}p{7.5cm}@{}}
\hline
\textbf{Champ} & \textbf{Modifiable} & \textbf{Contenu} \\ \hline
Date, Hour & non & Date et heure du d\'ecollage (EXIF). \\
Location & oui & Commune trouv\'ee automatiquement. Fait partie du nom du dossier de la mission. \\
GPS & non & Latitude et longitude du point de d\'ecollage. \\
Short description & oui & Br\`eve description libre de la mission. \\
Pilot & oui & Nom du t\'el\'epilote et son num\'ero d'identification. \\
Camera VIS & en partie & Fabricant et mod\`ele (lus dans l'image, non modifiables) ; intervalle de prise de vue (timelapse, en s) et d\'ecalage d'horloge (delta time, en s). \\
Image VIS format & non & Format des images (DNG). \\
Camera NIR & oui & Par d\'efaut SJCam M20, format RAW, intervalle 3\,s, d\'ecalage d'horloge $-30$\,s. \\
NIR Filter & oui & Fabricant du filtre et bande, en nm (par d\'efaut 810\,nm). \\ \hline
\end{tabular}
\end{center}

\figureIRD{step1_param_mission.png}{Fen\^etre \og Create a mission \fg.}{0.5}

Remarque: les paramètres de décalage d'horloge VIS et NIR sont réservés à un utilisateur expérimenté. Ce sont des paramètres qui permettent de \og récupérer\fg ~une mission dont la phase de synchronisation  sur le terrain se révèle inexploitable (impossibilité de reconnaître les ArUco, oubli de déclenchement de la mire, etc.).

%----------------------------------------------------------
\section{Bouton \bouton{Create the mission}}\label{btn:step1_validate}
Lance la cr\'eation de la mission. Le bouton se d\'esactive pendant le traitement ; une \'etiquette et une barre de progression
indiquent la phase en cours. Le logiciel encha\^ine, sans intervention de l'utilisateur sauf \`a l'\'etape~3 :
\begin{enumerate}
  \item \textbf{Cr\'eation du dossier} de mission et de ses sous-dossiers, \'ecriture de \texttt{mission\_parameters.json},
        copie de l'image de d\'ecollage.
  \item \textbf{Transfert des images VIS} : les fichiers \texttt{HYPERLAPSE\_XXXX.dng} du dossier de l'image de d\'ecollage sont copi\'es
        dans \texttt{AerialPhotography/VIS} sous le nom \texttt{VIS\_XXXX.dng}.
  \item \textbf{Choix du dossier NIR} : un message demande de choisir la premi\`ere image proche infrarouge (fichier RAW ou JPG
  \footnote{Chaque image RAW a une image jpeg associée. La SJCam nomme les images par leur date de prise de vue. Par exemple 2022\_0125\_130413\_001.RAW, 2022\_0125\_130415\_002.JPG, 2022\_0125\_130416\_003.RAW, 2022\_0125\_130418\_004.JPG, etc. L'image jpeg est toujours enregistrée après l'image RAW environ 2s après. À noter que dans certaines circonstances (carte SD trop lente, longue séquence de vol) le délai d'enregistrement du fichier RAW dépasse la fréquence de prise de vue (3s dans notre cas). Dans ce cas l'image suivante est prise 6s après. IRdrone gère automatiquement ce type de saut.} de la SJCam).
        Le dossier qui la contient est utilis\'e pour toutes les images NIR.
  \item \textbf{Conversion des images NIR} : les fichiers RAW sont convertis en DNG dans \texttt{AerialPhotography/NIR}
        (calcul en parall\`ele, \'eventuellement long).
  \item \textbf{Fichiers \texttt{.exif}} des images VIS puis NIR : un fichier de m\'etadonn\'ees par image,
        et un fichier \texttt{time\_line.json} par bande donnant la chronologie des prises de vue.
  \item \textbf{Altitudes et trajectoire} : pour chaque position GPS du vol, l'altitude du sol est demand\'ee \`a l'IGN (par paquets ;
        trois essais en cas d'\'echec, puis bascule automatique sur OpenTopoData). Le logiciel en d\'eduit la hauteur du drone au-dessus du sol,
        les distances et les caps entre prises de vue, et enregistre ces valeurs dans les fichiers \texttt{.exif}.
  \item \textbf{Profil vertical} : trac\'e de l'altitude du drone et du sol le long du trajet, enregistr\'e dans
        \texttt{FlightAnalytics/fly\_vertical\_profil.png}.
  \item \textbf{Message final} indiquant le dossier cr\'e\'e, puis fermeture des fen\^etres du Step~1.
\end{enumerate}

\figureIRD[fig:takeoff_1]{step1_nir_prompt.png}{Message demandant de choisir la premi\`ere image NIR.}{0.2}
\figureIRD[fig:fly_vertical_profil]{step1_fly_vertical_profil.png}{Profil vertical du vol (altitude du drone et du sol), produit \`a la fin du Step~1.}{0.8}

\subsection{Dossiers cr\'e\'es}
À l'issue de l'étape 1 (Step 1) IRDrone aura cr\'e\'e la structure de dossiers et fichiers de la mission. À ce stade vous n'aurez plus besoin d'accéder aux images originales stockées sur les cartes SD des caméras. La structure du dossier de la mission est la suivante :


\begin{verbatim}
FLY-YYYYMMDD-hhmm-<location>/
  AerialPhotography/
     VIS/   VIS_XXXX.dng, VIS_XXXX.exif, time_line.json,
            transfer_info_VIS_dng.json
     NIR/   NIR_XXXX.dng, NIR_XXXX.exif, time_line.json,
            transfer_info_NIR_dng.json
  FlightAnalytics/
     mission_parameters.json, fly_vertical_profil.png
  Synchro/        (Step 3)
  ImgIRdrone/     (Step 4)
  mapping_MULTI/  (Step 6)
  cameras/
\end{verbatim}
Les sous-dossiers ont une ic\^one sp\'ecifique dans l'explorateur Windows.
Le dossier de mission n'est jamais effac\'e : si le dossier existe d\'ej\`a, seuls les \'el\'ements manquants sont cr\'e\'es.
Pour les images VIS d\'ej\`a transf\'er\'ees (fichier \texttt{.exif} pr\'esent), le transfert est ignor\'e : relancer le Step~1 sur
une mission existante est donc rapide.

\figureIRD{Step1_mission_folder.png}{Dossier de mission vu dans l'explorateur Windows.}{0.7}

La durée des transferts dépend beaucoup des performances des supports (carte SD, disque dur) et des liaisons (USB B ou C, lecteur). De plus et contrairement aux images de la caméra DJI qui sont au format dng, les images de la caméra SJCam M20 sont dans un format \og RAW \fg~qui n'est pas normalisé. Chaque image doit être convertie au format dng d'ADOBE à l'aide d'un logiciel tiers \footnote{Il s'agit du code \texttt{sjcam\_raw2dng} disponible uniquement comme un fichier exe. \texttt{sjcam\_raw2dng}
	Convert SJCAM raw files to Adobe DNG and TIFF while also getting metadata
	such as date and ISO from JPEGs that correlate to the RAW files
} ce qui allonge le temps de transfert effectif. 

\subsection{En cas de probl\`eme}
\begin{itemize}
  \item \textbf{Aucune image NIR trouv\'ee} : le message \og No NIR RAW files found \fg\ signifie que le dossier choisi ne contient pas de
        fichiers \texttt{.RAW} ; recommencez en choisissant le bon dossier.
  \item \textbf{Erreur pendant le transfert VIS ou la lecture EXIF} : un message d'erreur s'affiche et la fen\^etre se ferme.
        Le bouton reste utilisable apr\`es r\'eouverture de la fen\^etre.
  \item \textbf{Altitudes IGN indisponibles} : le logiciel bascule seul sur OpenTopoData, moins pr\'ecis en France.
\end{itemize}
%\todo{Compl\'eter avec les erreurs rencontr\'ees sur le terrain (carte SD absente, dossier prot\'eg\'e en \'ecriture, \dots).}

%==========================================================
\legacyend

\chapter{Step 2 -- Ligne de temps du vol}\label{ch:step2}
\section{Principe}
L'étape 2 découpe le vol en phases : \emph{synchronisation} (vol stationnaire au-dessus de la mire ArUco tournante) et \emph{survey} (vol au-dessus de la zone étudiée). Ces phases sont maintenant \textbf{détectées automatiquement} ; l'utilisateur vérifie et ajuste si besoin (figure~\ref{fig:step2_flight_timeline}).

\figureIRD[fig:step2_flight_timeline]{step2_flight_timeline.png}{La fenêtre \og Flight timeline \fg. Au centre l'aperçu des images visible et infrarouge sous la souris, à droite la carte du vol avec les empreintes au sol des images gardées pour la carte (§\ref{btn:step2_best_mapping}), puis les quatre points clés avec leurs vraies images ; en bas la ligne de temps : profil d'altitude (points verts : images gardées pour la carte), phases (violet : synchronisation, vert : survey), détections de la mire et bandes d'images visible et infrarouge.}{0.95}

\subsection*{Ce que fait le logiciel}
\begin{enumerate}
	\item Chaque fichier DNG (DJI comme SJCam) contient un aperçu JPEG d'environ 1000 pixels. Il se lit en quelques millisecondes : toute la mission est parcourue en quelques secondes, sans RawTherapee. Les aperçus sont gardés dans \texttt{Synchro/Thumbnails} (cache, recréé si besoin).
	\item \emph{Visible} : l'altitude enregistrée par le drone donne le vol stationnaire du début (synchronisation) et le palier le plus long, le plus éloigné de l'altitude du décollage (survey). Un décollage depuis un sommet, suivi d'une descente vers la zone étudiée, est traité de la même façon qu'une montée.
	\item \emph{Synchronisation} : l'angle de la mire mobile (identifiant 0) est mesuré sur les aperçus, seulement pendant le stationnaire pour le visible, et pour l'infrarouge depuis la première image jusqu'à ce que la mire disparaisse (le drone s'est éloigné). Les images où la mire tourne régulièrement dans un sens forment la fenêtre de synchronisation ; les images où le plateau est encore arrêté, où il change de sens, et le décollage sont écartés.
	\item \emph{Infrarouge} : l'horloge de la caméra NIR est souvent fausse (fuseau horaire, dérive). La même rotation de la mire étant vue par les deux caméras, la comparaison des deux courbes d'angle donne un décalage approché (de l'ordre de la seconde) qui place les images infrarouges sur la ligne de temps du visible. La fenêtre de synchronisation infrarouge correspond ainsi au même tour de mire que la fenêtre visible.
	\item La fenêtre survey infrarouge \textbf{suit} la fenêtre survey visible (zone hachurée) : il n'est pas nécessaire de regarder les images infrarouges. Après l'étape 3, le décalage précis met à jour cette fenêtre.
\end{enumerate}

\subsection*{Vérifier et ajuster}\label{btn:step2_adjust}
\begin{itemize}
	\item Survoler la ligne de temps affiche en grand l'image visible et l'image infrarouge prises au même instant, avec l'altitude et l'état de la mire (\og target turning \fg, \og target still \fg, \og no target \fg). Laisser la souris sur cet aperçu l'affiche sur toute la fenêtre pour bien voir les détails.
	\item Les points de la ligne \og target \fg~indiquent les images analysées : violet, la mire tourne ; gris, la mire est vue mais immobile ; cercle vide, mire non vue.
	\item Faire glisser un bord de phase pour le déplacer (il se cale sur une image), ou cliquer sur un point clé (colonne de droite) puis utiliser \bouton{◀} \bouton{▶} ou les flèches du clavier.
	\item Molette : zoom ; glisser le fond : déplacement ; double-clic : tout le vol.
	\item \bouton{↺}, à côté de \bouton{?}, apparaît dès qu'une phase est déplacée : il annule les changements faits à la main et revient à la proposition automatique ; la fenêtre reste ouverte.
	\item L'info-bulle du bouton \bouton{?} rappelle ces commandes.
\end{itemize}
Si la mire n'est pas trouvée (mire trop petite, image trop sombre), un bandeau orange l'indique : la fenêtre de synchronisation visible est alors le vol stationnaire, et la fenêtre infrarouge doit être placée à la main.

\section{Best mapping : les images utiles à la carte}\label{btn:step2_best_mapping}
Un survol photographié toutes les deux secondes donne beaucoup plus de recouvrement que la carte n'en demande. Aligner toutes ces images à l'étape~4 prend des heures, et donner à ODM des centaines d'images presque identiques ralentit le calcul de la carte sans l'améliorer, souvent au contraire. La case \bouton{Best mapping: keep only the images the map needs}, cochée par défaut, ne garde que les images du survol dont les \emph{empreintes au sol} se recouvrent de la valeur choisie à droite (\bouton{75~\% overlap} par défaut, la valeur recommandée pour ODM ; plus de recouvrement garde plus d'images).

Le choix est visible tout de suite :
\begin{itemize}
	\item sur la ligne de temps, des points verts sur le profil d'altitude et des repères verts en haut de la bande d'images visibles marquent les images gardées ; la bulle du curseur et la légende de l'aperçu indiquent \og kept for the map \fg ;
	\item à droite des aperçus, la \textbf{carte du vol} montre le trajet (violet : synchronisation, vert : survey), les empreintes au sol des images gardées (rectangles verts : leur recouvrement se voit) et, en orange, l'empreinte de l'image sous la souris avec l'image elle-même. Un clic sur un point de la carte affiche cette image ; molette : zoom, glisser : déplacement, \bouton{Layers} : fond de carte (§\ref{btn:home_map}) ;
	\item sous la ligne de temps, le nombre d'images gardées, par exemple \og 16 of 120 survey images kept for the map · 87~\% less to align \fg.
\end{itemize}
La sélection suit les bords de la phase survey quand on les déplace.

\textbf{Comment les images sont choisies.} La position, la hauteur au-dessus du sol et l'orientation de chaque image sont lues dans les fichiers \texttt{.exif} écrits à l'étape~1 ; le champ de vue de la caméra donne la taille de l'empreinte au sol. La première image du survol est gardée. Pour chaque image suivante, IRDrone estime le recouvrement de son empreinte avec celle de la dernière image gardée, d'après le déplacement du drone (dans l'axe de l'image et en travers). Tant qu'il reste au-dessus de la valeur choisie, l'image n'est qu'une candidate ; dès qu'il passe en dessous, la candidate précédente est gardée : deux images gardées successives se recouvrent toujours d'au moins la valeur choisie. La dernière image du survol est gardée aussi. Pour un déplacement dans l'axe de l'image, c'est le critère \texttt{--selection best-mapping} de l'ancienne ligne de commande \texttt{run.py}.

\textbf{Et la synchronisation ?} À l'étape~2 le décalage précis entre les deux caméras n'est pas encore connu : l'aperçu porte sur toutes les images du survol. À l'étape~4, une fois les caméras synchronisées, la même sélection est refaite parmi les seules paires bien synchronisées (§\ref{btn:step4_best_mapping}) ; le nombre final peut donc différer un peu de l'aperçu. Le choix (case et recouvrement) est enregistré par \bouton{Continue with synchronisation} et \bouton{Close} dans \texttt{FlightAnalytics\textbackslash best\_mapping.json}.

Si les images n'ont pas de position GPS, la carte n'est pas affichée et toutes les images du survol sont gardées.

\section{Boutons \bouton{Continue with synchronisation} et \bouton{Close}}\label{btn:step2_confirm}
Les deux boutons enregistrent les phases dans les fichiers \texttt{transfer\_info\_VIS\_dng.json} et \texttt{transfer\_info\_NIR\_dng.json} (lus par les étapes 3, 4 et le mapping) et dans \texttt{FlightAnalytics/flight\_timeline.json}. Rouvrir l'étape 2 retrouve les phases confirmées.
\bouton{Continue with synchronisation} ouvre ensuite l'étape~3 ; rien ne démarre tout seul : la synchronisation se lance depuis sa fenêtre (chapitre~\ref{ch:step3}). \bouton{Close} ferme simplement la fenêtre : l'étape~3 se lancera plus tard depuis la fenêtre principale. Il devient \bouton{Save and close} quand les phases ou le choix best mapping ont changé, ou n'ont jamais été enregistrés : fermer la fenêtre les enregistre alors aussi. S'il manque une phase, fermer demande s'il faut fermer sans enregistrer.

\section{Sélection classique}\label{btn:step2_classic}
\legacybegin{Afficher la sélection classique (ancienne fenêtre de l'étape 2)}
Le bouton \bouton{Classic selection…} ouvre l'ancienne sélection manuelle des cinq images de référence de chaque caméra, décrite ci-dessous.

\figureIRD{step2_dialog_select_phase_images_1.png}{Fen\^etre de s\'election des images de r\'ef\'erence. Pour la bande spectrale visible la première image est par défaut celle du point de décollage qui a servi à définir la mission à l'étape 1. Mais elle peut être changée ici.}{0.95}
La fenêtre de dialogue permet la définition :
\begin{enumerate}
	\item du point de décollage,
	\item des points de début et fin de synchronisation,
	\item des points de début et fin de vol au-dessus de la zone étudiée par la mission.
\end{enumerate}
Les images sont sélectionnées pour les spectres visible et proche infrarouge  respectivement dans les dossiers
\begin{itemize} 
	\item \texttt{C:/Air-Mission/FLY-YYYYMMDD-hhmm-location/AerialPhotography/VIS}, 
   	\item  \texttt{C:/Air-Mission/FLY-YYYYMMDD-hhmm-location/AerialPhotography/NIR}.
\end{itemize} 
 Lors de cette étape de choix il est toujours possible après une première sélection de choisir une autre image. 
 L'ordre de sélection des cinq images n'a pas d'importance mais bien entendu first précède toujours last.
 Il est possible de sélectionner la même image pour \texttt{first mission image} et \texttt{first sync image}.  
  Pour les images DJI du spectre visible, si vous sélectionnez la même image pour \texttt{first sync image} et \texttt{last sync image} alors une seule image sera retenue (c'est un cas limite pour la synchronisation). 
  En revanche  si vous sélectionnez la même image pour \texttt{first sync image} et \texttt{last sync image} pour le spectre proche infrarouge la synchronisation automatique échouera à l'étape 3 et il vous sera demandé d'entrer manuellement l'offset de synchronisation des horloges. 
  Cela peut constituer une \og astuce \fg~ pour forcer l' offset de synchronisation des horloges à une valeur particulière en cas de synchronisation terrain problématique.
%\todo{\bouton{Load first mission image.}, \bouton{Load first sync image.}, \bouton{Load last sync image.},
%----------------------------------------------------------
\section{Bouton \bouton{Load first mission image.}}\label{btn:step2_Load_first_mission_image}
Le choix de cette image définit l'origine de la ligne de temps pour la synchronisation des horloges des caméras. En règle générale on choisira la première prise de vue VIS de la mission.
\section{Boutons \bouton{Load first sync image.}, \bouton{Load last sync image.}}\label{btn:step2_Load_synchro_image}
L'utilisateur doit choisir l'image de début du cycle de synchronisation. Pour les images visibles, si la procédure a été respectée sur le terrain, il s'agit de la première ou deuxième image enregistrée. Pour les images infrarouges il faut repérer la première image stable où le système de mire est visible en totalité et où la mire mobile est effectivement en rotation. Ensuite il faut choisir la dernière image de la phase de synchronisation. En général on la repère facilement car lorsque la phase de synchronisation est terminée le pilote prend rapidement de l'altitude et le changement d'aspect est bien visible. Les choix faits ici conditionnent en partie la réussite de l'étape 3 de synchronisation des horloges. Mais après l'étape 3 il est possible de revenir à l'étape 2 et de redéfinir ces images.
\figureIRD[fig:step2_dialog_select_VIS]{step2_dialog_select_phase_image_VIS.png}{Fen\^etre de s\'election des images de r\'ef\'erence. Bande spectrale visible (VIS).}{0.95}
\section{Boutons \bouton{Load first fly image.}, \bouton{Load last fly image.}}\label{btn:step2_Load_mission_image}
L'utilisateur doit choisir la première image prise au-dessus de la zone à étudier (vol mission ou hyperlapse survey). On élimine ainsi les images de l'approche entre le décollage/synchronisation à très basse altitude et le vol de la mission à \og haute altitude \fg~ (hyperlapse survey). À noter qu'en principe le vol mission devrait s'opérer à altitude constante du drone et au-dessus d'une zone dont les altitudes varient peu par rapport à l'altitude du drone par rapport au sol. L'observation de la figure \ref{fig:fly_vertical_profil} permet de juger de la qualité du plan de vol réalisé sur le terrain.		
\section{Bouton \bouton{Validate mission images}}\label{btn:step2_validate}
\bouton{Cancel}, à sa gauche, ferme la fenêtre (après confirmation si les images ne sont pas validées). Ce bouton déclenche la construction des listes des images des trois phases de vol pertinentes pour IRDrone : takeoff, synchro et fly. Aucune image n'est déplacée. L'opération est très rapide.
\figureIRD[fig:step2_dialog_select_NIR]{step2_dialog_select_phase_images_NIR.png}{Fen\^etre de s\'election des images de r\'ef\'erence. Bande spectrale proche infrarouge (NIR).}{0.95}

%==========================================================
\legacyend

\chapter{Step 3 -- Synchronisation des caméras}\label{ch:step3}
\section{Principe}
Pour combiner différentes bandes spectrales, les images doivent être prises au même instant. Malheureusement ce n'est techniquement pas possible avec les matériels utilisés dans le projet IRdrone. C'est un choix de départ lié au caractère financièrement abordable (\emph{inexpensive and low cost}) du projet. Chaque caméra a sa propre horloge, qui définit la chronologie des images. On peut positionner les images sur une ligne représentant le temps. C'est la \og time line \fg~de la caméra. Dès lors il faut synchroniser les deux time lines VIS et NIR. Mettre les horloges à l'heure est insuffisant : d'une part elles ne donnent accès qu'à des points toutes les secondes, d'autre part on ne peut pas contrôler l'instant exact où l'opérateur sur le terrain a déclenché les caméras. On doit donc détecter des repères sur les images. Dans l'industrie cinématographique on synchronise la pellicule photo et la bande son à l'aide du \og clap \fg. Cette solution avait été envisagée dans les premières versions d'IRdrone, qui utilisaient la vidéo au lieu des images. Mais la vidéo posait de nombreux problèmes de qualité (résolution, \og jelly \fg, ...) et a été abandonnée. Par contre le principe d'un objet mobile sur les images au début du vol a été conservé. L'idée est de repérer un objet (automatiquement) et de suivre sa trajectoire sur les deux caméras. En faisant coïncider les deux trajectoires on synchronise les horloges. On peut atteindre une précision de l'ordre de 0.05s, ce qui est amplement suffisant dans notre cas.
IRdrone utilise comme objet mobile une mire ArUco 
\footnote{Ce type de mire est conçu pour être détecté automatiquement par logiciel, par exemple avec OpenCV. IRdrone utilise le dictionnaire \texttt{DICT\_4X4\_50} : la mire mobile porte l'identifiant 0 et les quatre mires fixes les identifiants 5, 22, 16 et 13.} 
fixée à un plateau tournant
\footnote{Le plateau est entraîné par un moteur pas à pas commandé par un Arduino. On programme le système pour obtenir une vitesse de rotation parfaitement constante. Elle est fixée à la valeur de 1 rpm (un tour par minute) de telle sorte que lors de la phase de vol de synchronisation les caméras enregistrent au moins une dizaine de positions angulaires différentes.} comme le montre la figure \ref{fig:step3_syncro_aruco_1}. 

\figureIRD[fig:step3_syncro_aruco_1]{step3_syncro_aruco_1.pdf}{Dispositif de synchronisation des horloges des caméras VIS et NIR. Le principe repose sur la détection de la position angulaire de la mire ArUco 0 qui tourne à vitesse constante. Les petites mires fixes permettent de définir un repère absolu pour s'affranchir des mouvements parasites du drone, particulièrement les rotations.}{0.95}

On mesure la position angulaire de la mire. Mais pour cela il faut un repère fixe. Les bords de l'image ont été utilisés initialement. Cependant ils ne sont fiables que si le drone ne tourne pas. Il faut donc une zone très calme et que le pilote ne touche pas à la manette de rotation. Pour libérer cette contrainte sur le terrain, on a ajouté au dispositif 4 mires ArUco fixes
\footnote{Théoriquement 2 mires suffisent pour définir un axe de référence. L'utilisation de 4 mires permet une certaine redondance. À noter que le logiciel tente de synchroniser dans le repère absolu, mais s'il échoue il synchronise dans le repère relatif lié aux bords des images (moins précis). Si là encore il y a échec, en dernier recours le logiciel demande à l'utilisateur de donner la valeur de l'offset de synchronisation.}. Elles permettent de définir un repère absolu qui améliore la précision du dispositif. Il est évident que cette phase de synchronisation doit être réalisée sur le terrain avec le plus grand soin. C'est la condition pour concilier le coût modéré du matériel et une bonne qualité des images multispectrales.

\subsection{Les trois phases de l'étape 3}
Le Step~3 se déroule en trois phases successives, enchaînées automatiquement (§\ref{sec:step3_chained}). En mode manuel, chacune est lancée par un bouton qui n'est activé qu'une fois la phase précédente terminée (figure~\ref{fig:step3_window}) :
\begin{enumerate}
  \item \emph{Détection des mires} (bouton \bouton{Run ArUco detection (VIS + NIR)}) : sur chaque image de synchronisation, on mesure l'angle de la mire mobile (§\ref{btn:step3_aruco}) ;
  \item \emph{Calcul du décalage temporel} (bouton \bouton{Compute time alignment}) : on cherche le décalage $\Delta t$ qui fait coïncider les deux séries d'angles (§\ref{btn:step3_time_shift}) ;
  \item \emph{Appariement} (bouton \bouton{Pair VIS / NIR images}) : on associe à chaque image VIS les images NIR les plus proches dans le temps (§\ref{btn:step3_pairing}).
\end{enumerate}

\subsection{Enchaînement automatique}\label{sec:step3_chained}
L'étape~3 ne démarre jamais seule : qu'elle soit ouverte depuis l'étape~2 (\bouton{Continue with synchronisation}) ou depuis la fenêtre principale, le bouton \bouton{Synchronise the cameras} (bleu foncé, au-dessus de la barre de progression) enchaîne les trois phases (§\ref{btn:step3_run}). Une ligne d'avancement montre la phase en cours et son résultat : nombre de mires vues, décalage, nombre de paires (figure~\ref{fig:step3_chained}). Une seule question est posée, après l'alignement des horloges (§\ref{btn:step3_overlap}).

Une fois les images appariées, rien ne démarre seul : \bouton{Continue with image processing} ouvre l'étape~4 quand vous êtes prêt (§\ref{btn:step3_continue}).
Si le décalage ne peut pas être mesuré, la saisie manuelle est proposée et l'enchaînement s'arrête.

\figureIRD[fig:step3_chained]{step3_chained.png}{Étape 3 enchaînée : mires détectées, horloges alignées ; la question \og Do the two curves overlap? \fg~attend la réponse avant l'appariement.}{0.7}

Rouverte une fois l'étape faite (clic sur l'étape~3 de la fenêtre principale), la fenêtre ne recalcule rien : elle montre le résultat enregistré, c'est-à-dire le décalage des horloges, le nombre de paires et les courbes alignées (§\ref{btn:step3_before}). \bouton{Synchronise again} la refait ; \bouton{Manual mode…} montre les boutons des trois phases.

\section{La question \og Do the two curves overlap? \fg}\label{btn:step3_overlap}
Après l'alignement des horloges, la fenêtre montre l'angle de la mire en fonction du temps : courbe verte pour le visible, rouge pour l'infrarouge, la courbe infrarouge étant décalée du décalage calculé.
\begin{description}
	\item[\bouton{Yes} (touche Y)] les deux courbes se superposent : le décalage est bon, les images sont appariées (§\ref{btn:step3_pairing}).
	\item[\bouton{No} (touche N)] ouvre le réglage à la main (§\ref{btn:step3_adjust}) pour glisser la courbe infrarouge sur la courbe visible ; la question est ensuite reposée avec le nouveau décalage.
\end{description}
Des courbes bien superposées ont des points verts sur la ligne rouge, dents de scie comprises. Les exemples du §\ref{sec:step3_examples} montrent un bon et un mauvais cas.

\section{Bouton \bouton{Synchronise the cameras}}\label{btn:step3_run}
L'action principale de l'étape~3 : elle lance la détection des mires, l'alignement des horloges puis, une fois répondu à la question sur les courbes, l'appariement. Elle devient \bouton{Synchronise again} quand l'étape a déjà été faite.

\section{Bouton \bouton{Manual mode…}}\label{btn:step3_manual}
Arrête l'enchaînement et fait apparaître les boutons des trois phases (§\ref{btn:step3_aruco} à \ref{btn:step3_pairing}), à lancer un par un, par exemple pour corriger le décalage avec \bouton{Adjust alignment...} avant l'appariement.

\section{Bouton \bouton{Enter the clock offset…}}\label{btn:step3_clock_offset}
Option avancée, pour un vol sans la cible ArUco (ancienne synchronisation par les dates). Le bouton n'apparaît que si la case \emph{Advanced – Step 3: let me enter the camera clock offset by hand} est cochée dans les options (roue dentée \bouton{⚙} en bas de la fenêtre principale, §\ref{btn:home_workdir}). Avec cette option, l'étape~3 ne démarre pas d'elle-même en venant de l'étape~2 : on choisit entre \bouton{Synchronise the cameras} (vol avec cible) et ce bouton.

Une petite fenêtre demande de combien de secondes l'horloge de la caméra infrarouge avance sur celle du drone (négatif si elle retarde, 0 si les deux horloges ont été réglées à la même heure). C'est le délai de l'ancienne synchronisation. Elle affiche le décalage $\Delta t$ de la ligne de temps infrarouge qui en découle, calculé à partir des dates des premières images. \bouton{Pair the images} saute la détection des cibles et lance directement l'appariement (§\ref{btn:step3_pairing}). Aucune courbe ne peut être vérifiée : la qualité de l'appariement dépend de la justesse du délai saisi.

\section{Bouton \bouton{Continue with image processing}}\label{btn:step3_continue}
Apparaît en bas à droite, à côté de \bouton{Close}, une fois les images appariées. Il ferme la fenêtre et ouvre l'étape~4. L'alignement, long, n'y démarre pas tout seul : il se lance avec \bouton{Run alignment}.

\section{Bouton \bouton{Before alignment}}\label{btn:step3_before}
Quand l'étape~3 est rouverte une fois faite, la fenêtre montre les courbes après alignement (\texttt{Synchro/check\_time\_alignment.png}). \bouton{Before alignment} montre les courbes avant correction (\texttt{check\_time\_before\_alignment.png}) ; un second clic revient aux courbes alignées.

\subsection{Principe du calcul}
L'angle de la mire mobile est mesuré sur chaque image, puis déplié (voir §\ref{btn:step3_time_shift}). Comme la vitesse de rotation $\omega$ est constante, l'angle est une fonction affine du temps. La time line NIR, qui compte au moins deux points, fournit la droite de référence $\theta = \omega t + b_{NIR}$. Chaque point VIS fournit une estimation de l'ordonnée à l'origine $b_{VIS} = \theta_{VIS} - \omega t_{VIS}$ ; on en prend la médiane. Le décalage s'écrit
\begin{equation}
  \Delta t = \frac{b_{NIR} - b_{VIS}}{\omega},
\end{equation}
et la time line NIR corrigée est $t_{NIR} + \Delta t$. C'est pourquoi une seule image VIS suffit en théorie (cas~2 ci-dessous), à condition que la vitesse de rotation soit constante et que l'on travaille dans le repère absolu.

\subsection{Exemples de synchronisation}\label{sec:step3_examples}
Les images de la phase de synchronisation viennent de l'étape~2. On présente quelques cas ici.

\subsubsection{Cas 1 : bonne synchronisation}
Les figures~\ref{fig:synchro_cas1} (\subref{fig:synchro_cas1_avant} et \subref{fig:synchro_cas1_apres}) représentent une bonne synchronisation. Le pilote a déclenché manuellement l'hyperlapse de la caméra NIR puis il a fait décoller le drone, ce qui a déclenché automatiquement la rotation de la mire. Lorsque le pilote a placé le drone au-dessus de la mire, la caméra NIR enregistrait la mire en rotation. Les images entre l'instant de déclenchement de l'hyperlapse NIR et l'instant où le drone a été stabilisé au-dessus de la mire ont été ignorées au step~2. Ensuite le pilote a réglé les paramètres de l'hyperlapse de la caméra VIS puis l'a déclenché. Cela prend en général une trentaine de secondes pour un pilote entraîné. Dès que l'hyperlapse VIS commence, il enregistre la mire en mouvement. À partir de cet instant et avec une fréquence d'enregistrement de 2s, il suffit au pilote de maintenir le drone moins d'une dizaine de secondes au-dessus de la mire pour obtenir 4 à 5 points, avant de prendre rapidement de l'altitude pour rejoindre la zone à étudier (hyperlapse survey). L'utilisateur du logiciel doit veiller lors de l'étape 2 à éliminer les images NIR prises avant la stabilisation du drone au-dessus du dispositif de synchronisation et le cas échéant celles qui après stabilisation correspondraient à une mire non encore en rotation. 

\figurePaireIRD{fig:synchro_cas1}
  {Step3_time_alignment_failed_cas1_0.png}{Trajectoire de la mire avant synchronisation.}
  {Step3_time_alignment_failed_cas1_1.png}{Synchronisation réussie.}
  {Cas de bonne synchronisation.}

\subsubsection{Cas 2 : cas limite}
Les figures~\ref{fig:synchro_cas2} (\subref{fig:synchro_cas2_avant} et \subref{fig:synchro_cas2_apres}) sont un cas limite pour la synchronisation. Le pilote est resté très peu de temps au-dessus de la mire et une seule image de cette mire est disponible sur les images VIS. Ici la synchronisation fonctionne tout de même car la vitesse de rotation est maintenue constante et le repère absolu a pu être utilisé.

\figurePaireIRD{fig:synchro_cas2}
  {Step3_time_alignment_failed_cas2_0.png}{Trajectoire de la mire avant synchronisation.}
  {Step3_time_alignment_failed_cas2_1.png}{Cas limite de synchronisation.}
  {Cas limite de synchronisation.}

\subsubsection{Cas 3 : échec de synchronisation}
Les figures~\ref{fig:synchro_cas3} (\subref{fig:synchro_cas3_avant} et \subref{fig:synchro_cas3_apres}) montrent la trajectoire de la mire détectée sur les deux caméras. Ici la synchronisation échoue car les premières images retenues pour la caméra NIR correspondent à des instants où la mire n'est pas en rotation.

\figurePaireIRD{fig:synchro_cas3}
  {Step3_time_alignment_failed_cas3_0.png}{Trajectoire de la mire avant synchronisation.}
  {Step3_time_alignment_failed_cas3_1.png}{Échec de synchronisation.}
  {Cas d'échec de synchronisation.}
   
\figureIRD[fig:step3_window]{step3_window.png}{Fen\^etre de synchronisation.}{0.8}

\section{Bouton \bouton{Run ArUco detection (VIS + NIR)}}\label{btn:step3_aruco}
Ce bouton lance la détection des mires ArUco sur toutes les images de la phase de synchronisation choisies à l'étape~2, d'abord pour la bande VIS puis, automatiquement, pour la bande NIR. Une barre de progression indique l'avancement. Les calculs sont répartis sur plusieurs c\oe urs du processeur pour accélérer le traitement.

\subsection*{Ce que fait le logiciel pour chaque image}
\begin{enumerate}
  \item Les images DNG sont converties en TIF par RawTherapee avec un profil de développement dédié à la détection (fichiers \texttt{*\_RawTherapee.tif} dans les dossiers \texttt{Synchro\textbackslash VIS} et \texttt{Synchro\textbackslash NIR} de la mission). Cette conversion n'est faite qu'une fois : si le fichier existe déjà, il est réutilisé.
  \item Les mires ArUco visibles sont détectées avec OpenCV. Pour chacune, on retient son identifiant, la position de son centre et l'angle de son axe principal dans l'image.
  \item La mire mobile (identifiant~0) donne l'\emph{angle relatif} (repère de l'image, \texttt{angle\_img}).
  \item Si les quatre mires fixes (identifiants 5, 22, 16 et 13) sont aussi détectées, elles définissent un repère lié au sol : deux axes obtenus en moyennant les côtés opposés du quadrilatère formé par leurs centres. L'angle de la mire mobile est alors mesuré dans ce repère. C'est l'\emph{angle absolu} (\texttt{angle\_abs}), insensible aux rotations du drone autour de son axe vertical.
  \item Si la mire mobile ou l'une des mires fixes manque, l'image est signalée dans le journal. L'angle absolu n'est pas calculé (mode dégradé : seul l'angle relatif est conservé si la mire mobile est visible) ; si la mire mobile n'est pas détectée, l'image est ignorée.
  \item À chaque image est associée sa date relative (en secondes) sur la time line de sa caméra.
\end{enumerate}
Les résultats de détection sont conservés dans les métadonn\'ees de chaque image.

\subsection*{Images de contrôle}
Pendant la détection, chaque image traitée est enregistrée avec une surcouche graphique dans le dossier \texttt{Synchro\textbackslash Check\_ARUCO}. La fenêtre en affiche une régulièrement (environ toutes les 5 secondes) pour que l'utilisateur puisse juger de la qualité de la détection. La surcouche montre :
\begin{itemize}
  \item les centres des mires fixes (disques rouges) ;
  \item le centre de la mire mobile ;
  \item l'orientation de la mire mobile dans l'image (trait vert) ;
  \item l'orientation projetée dans le repère absolu (trait orange en pointillés) et les deux axes du repère absolu (flèches noires) lorsque les quatre mires fixes sont détectées ;
  \item une légende : angle relatif, angle absolu et différence entre les deux, ou le message \og absolute reference: unavailable \fg~si le repère absolu n'a pas pu être construit.
\end{itemize}
\textbf{Conseil :} parcourir ce dossier dans l'explorateur Windows est le meilleur moyen de repérer des détections douteuses avant de passer à la phase suivante.

% Figure à créer par l'utilisateur (PowerPoint) : 2 lignes de 4 images (VIS en haut, NIR en bas) extraites de Synchro\Check_ARUCO.
\figureIRD[fig:check_aruco]{Step3_check_aruco_VIS_NIR.pdf}{Exemples d'images de contrôle (dossier \texttt{Synchro\textbackslash Check\_ARUCO}). Ligne du haut : bande VIS ; ligne du bas : bande NIR. La mire mobile est repérée par son axe vert ; les centres des mires fixes par des disques rouges.}{0.95}

\section{Bouton \bouton{Compute time alignment}}\label{btn:step3_time_shift}
Ce bouton, actif seulement lorsque la détection est terminée, calcule le décalage temporel entre les deux caméras.

\subsection*{Choix des angles}
Par défaut le logiciel utilise les angles absolus. Il faut au moins 2 images NIR et 1 image VIS où l'angle absolu est disponible. À défaut, il utilise les angles relatifs (repère de l'image), avec les mêmes minima, ce qui est moins précis si le drone a tourné pendant la synchronisation. Si aucune des deux solutions n'est possible, le message \og Insufficient data to automatically align VIS and NIR camera timelines \fg~s'affiche et le logiciel demande à l'utilisateur de saisir lui-même le décalage en secondes (valeur par défaut : 0). Dans ce cas aucun graphique n'est produit.

\subsection*{Calcul}
\begin{enumerate}
  \item Les angles de chaque bande sont \og dépliés \fg : on supprime les sauts de $\pm 360^\circ$ pour obtenir une courbe continue qui croît (ou décroît) régulièrement avec le temps. Le dépliage est fait de la même façon pour VIS et NIR, en partant de la dernière image, c'est-à-dire de l'image où la mire est en principe en rotation.
  \item On ajuste une droite sur les points NIR : sa pente est la vitesse de rotation $\omega$.
  \item On en déduit, avec les points VIS, le décalage $\Delta t$ (voir le principe en début de chapitre). L'estimation utilise la médiane des points VIS pour être robuste aux valeurs aberrantes.
\end{enumerate}
Le résultat est affiché dans le journal de la fenêtre :
\begin{itemize}
  \item le décalage estimé $\Delta t$ (en secondes) ;
  \item la vitesse de rotation mesurée sur NIR et sur VIS (en rpm). Elle doit être proche de la valeur nominale de 1~rpm et les deux valeurs doivent être voisines ; l'écart relatif est affiché ;
  \item le nombre de points VIS et NIR utilisés ;
  \item la qualité de l'ajustement NIR ($R^2$, proche de 1) ;
  \item la dispersion des décalages individuels (écart interquartile et écart-type, en secondes).
\end{itemize}

\subsection*{Graphiques de contrôle}
Deux graphiques sont enregistrés dans le dossier \texttt{Synchro} de la mission. Pendant la question \og Do the two curves overlap? \fg, la fenêtre montre les courbes après alignement ; en mode manuel, elle montre les deux graphiques en alternance (changement toutes les 1,5 s) :
\begin{itemize}
  \item \texttt{check\_time\_before\_alignment.png} : angle de la mire en fonction du temps pour VIS (vert) et NIR (rouge) avant correction ;
  \item \texttt{check\_time\_alignment.png} : mêmes courbes après avoir décalé la time line NIR de $\Delta t$.
\end{itemize}
Si la synchronisation est bonne, les points verts se placent sur la droite des points rouges. Ce sont ces deux graphiques qui sont présentés dans les trois cas de ce chapitre.

\textbf{Que faire en cas d'échec ?} Si les points VIS et NIR ne se superposent pas après alignement, ou si les vitesses de rotation VIS et NIR sont très différentes, revenir au Step~2 et redéfinir les images de synchronisation (en particulier la première image NIR, qui doit montrer la mire effectivement en rotation), puis relancer le Step~3.

\section{Bouton \bouton{Adjust alignment...}}\label{btn:step3_adjust}
Ouvre le réglage du décalage à la main, aussi ouvert par la réponse \bouton{No} à la question \og Do the two curves overlap? \fg. La fenêtre trace la courbe visible (verte), la courbe infrarouge au décalage calculé (rouge, en pointillés) et la courbe infrarouge au décalage réglé (rouge). Le curseur et la case numérique déplacent la courbe infrarouge : flèches du clavier par pas de 1\,ms, touches page précédente / suivante par pas de 0,1\,s.
\begin{description}
	\item[\bouton{Save and close}] (\bouton{Close} tant que le décalage n'a pas changé ; fermer la fenêtre fait de même) garde ce décalage : les graphiques de contrôle sont refaits. Si les images étaient déjà appariées, il faut refaire l'appariement.
	\item[\bouton{↺}] (en bas à gauche, affiché quand le décalage diffère du décalage calculé) revient au décalage calculé à partir des mires.
\end{description}
Si le décalage n'a pas pu être calculé, une petite fenêtre demande directement sa valeur en secondes (ajoutée aux dates infrarouges).

\section{Bouton \bouton{Pair VIS / NIR images}}\label{btn:step3_pairing}
Ce bouton est actif une fois le décalage temporel connu (calculé ou saisi à la main). Il construit la liste des couples d'images qui seront alignés au Step~4.

\begin{enumerate}
  \item La time line VIS sert de référence. Toutes les dates NIR sont corrigées : $t'_{NIR} = t_{NIR} + \Delta t$.
  \item Pour chaque image VIS, on retient les \emph{deux} images NIR les plus proches dans le temps. On mémorise pour chacune l'écart signé $dt = t'_{NIR} - t_{VIS}$ : $dt<0$ si la NIR précède la VIS, $dt>0$ si elle la suit.
  \item Les métadonnées de vol (position GPS, coordonnées UTM, altitudes, angles de lacet, tangage et roulis du drone et de la nacelle, distances, cap, etc.), connues pour les images VIS, sont interpolées aux instants corrigés des images NIR. Les angles (lacet, cap) sont interpolés par le chemin angulaire le plus court. Hors de la plage de temps VIS, la valeur de l'extrémité est conservée.
  \item Le résultat est enregistré dans le fichier \texttt{AerialPhotography\textbackslash shoot\_points.json}, utilisé par les étapes suivantes. Le journal indique le nombre de couples créés.
\end{enumerate}
Un même fichier NIR peut être utilisé pour plusieurs images VIS (par exemple si la fréquence de prise de vue NIR est plus élevée) et aucun écart temporel maximal n'est imposé : la qualité de l'appariement dépend donc de celle de la synchronisation.

\section{Bouton \bouton{Close}}\label{btn:step3_close}
Ferme la fenêtre et rend la main au dialogue principal. Les fichiers produits (graphiques, images de contrôle, \texttt{shoot\_points.json}) restent dans le dossier de la mission. Le bouton \bouton{?} ouvre cette aide à la page de ce chapitre.

%==========================================================
\chapter{Step 4 -- Traitement des images}\label{ch:step4}
\section{Principe}
Le traitement des images consiste pour une paire VIS-NIR à aligner ces deux images pour que les bandes spectrales Red Green Blue et Nir se superposent. L'opération est complexe pour plusieurs raisons :
\begin{itemize}
	\item Même si les time lines sont synchronisées, les instants de prise de vue sont proches mais distincts.
	\item Les deux caméras ne pointent pas vers le même point au sol. La caméra VIS du drone DJI est fixée sur un gimbal très performant qui maintient la visée vers le nadir. À noter que le vol hyperlapse est fait en mode \og gimbal lock \fg~autrement dit l'angle de roll du gimbal reste fixe. Seuls ses angles de pitch et yaw varient pour maintenir la visée vers le nadir (à $90^{\circ}$ vers le bas). En revanche la caméra NIR SJCam M20 étant fixée rigidement sous le corps du drone sa visée est totalement tributaire de l'attitude du drone (angles de pitch et yaw en particulier). 
	\item Les distances focales des deux caméras sont différentes. Pour la caméra VIS on a un angle (FOV) de $84^{\circ}$ alors que pour la caméra NIR la focale est ultra-courte avec un angle (FOV) supérieur à $150^{\circ}$ ce qui correspond quasiment à un \og fish eye \fg.
	\item Les optiques présentent de la distorsion et du vignetage. Pour la distorsion les deux caméras ont été préalablement calibrées par la méthode de Zhang.
\end{itemize} 



En pratique on commence par convertir les images du format raw-dng au format tif (via le logiciel tiers RawTherapee
\footnote{Via le logiciel tiers open source RawTherapee. L'opération est coûteuse la première fois mais l'image est mise en cache pour des runs ultérieurs.}).  On élimine la distorsion des deux images. On détermine les déplacements théoriques à appliquer pour ramener le centre de visée NIR sur le centre de visée VIS. Ces déplacements sont traduits en termes d'angles pitch, yaw et roll. Ils servent à initialiser la procédure d'alignement. L'alignement se déroule très schématiquement en deux temps. On détermine en premier les paramètres d'une homographie à appliquer sur l'image NIR pour qu'elle \og match \fg~avec l'image VIS. C'est un processus de minimisation complexe spécialement développé pour les images multispectrales 
	\footnote{Les techniques habituelles pour aligner deux images, qui sont extrêmement performantes dans le cas d'une bande spectrale unique ou en RGB, comme par exemple les \og sift \fg, échouent.}. 
Ce premier alignement (coarse process) est suivi d'un alignement fin (fine process) qui applique des homographies locales pour diminuer encore le critère d'alignement. L'algorithme est décrit en annexe (§\ref{sec:ann_algorithme}). L'utilisateur doit être conscient que la méthode n'est pas infaillible et surtout qu'elle est tributaire des conditions de vol sur le terrain. 

\figurePaireIRD{fig:step4_pair_VIS_NIR}
{step3_img_VIS.png}{Spectre visible (VIS).}
{step3_img_NIR.png}{Spectre proche infrarouge (NIR).}
{Couple d'images d'un shoot point. Les deux bandes spectrales enregistrent des informations différentes. Par exemple la surface de la rivière est totalement noire en proche infrarouge. L'eau ne réfléchit pas les longueurs d'onde au-delà de $750\,nm$.}  

\section{Dialogue de l'étape 4}
La fenêtre de l'étape~4 (figure~\ref{fig:step4_window}) rappelle le nom de la mission (le chemin complet est dans l'info-bulle) et le nombre de paires d'images. Elle comporte, de haut en bas : les options de sortie, le choix des images alignées avec la carte du vol, le bouton \bouton{Run alignment} (\bouton{Stop} pendant le calcul), l'état et la barre de progression du calcul, le journal, puis, en bas, \bouton{?} et \bouton{⚙} (réglages) à gauche, \bouton{Close} et, une fois des images alignées, \bouton{Continue with the review} à droite. La fenêtre s'adapte à la taille de l'écran.

\figureIRD[fig:step4_window]{step4_window.png}{Fen\^etre de traitement des \og shoot points \fg~(paires d'images VIS-NIR), avant le calcul.}{0.7}


\section{Options de sortie}\label{btn:step4_options}


\subsection{Case \bouton{first VIS image number} et \bouton{last VIS image number}}
Ces deux cases permettent de sélectionner des shoot points à traiter. Par défaut IRDrone propose la plage qui a été définie à l'étape~2. Il est possible de ne traiter qu'une image en entrant le même numéro pour first et last images, sous réserve que le \texttt{Maximum VIS/NIR time difference} soit compatible. En cas de doute fixez le \texttt{Maximum VIS/NIR time difference} à la valeur $3 \, s$.

Par défaut la procédure d'alignement produit  pour 

$\texttt{first VIS image number} \leq XXXX \leq \texttt{last VIS image number}$:
\begin{itemize}
	\item les quadruplets d'images multispectrales red (1), green (2), blue (3), nir (4) (\texttt{IRD\_XXXX\_n.tif}) placés dans le dossier \texttt{ImgIRdrone\textbackslash imgMultispectral}.
	\item les images \texttt{VIS\_XXXX.jpg}, \texttt{VIR\_XXXX\_local.jpg}.
	\item les fichiers \texttt{VIS\_XXXX\_REDO.bat} et deux fichiers techniques \texttt{VIS\_XXXX\_motion\_model.json}, \texttt{VIS\_XXXX\_motion\_model.npy}.
	\item le fichier \texttt{ImgIRdrone\textbackslash shoot\_points.json} qui contient un dictionnaire des shoot points traités avec en particulier la clé  \texttt{"failed"} dont la valeur est fixée à \texttt{false} par défaut. L'utilisateur vérifie cette hypothèse dans la revue des images alignées (chapitre~\ref{ch:review}). 
\end{itemize}


\subsection{Case \bouton{NDVI}, case \bouton{NIR}}
Par défaut la procédure d'alignement produit les images VIS et VIR au format jpeg. C'est utile pour vérifier visuellement la réussite des alignements. En cochant les cases \texttt{NDVI} et/ou \texttt{NIR} on peut générer les images NDVI et NIR au format jpeg. Ne cocher ces cases que si vous envisagez d'utiliser ces représentations directement dans un rapport. Notez qu'à l'étape~6 vous avez la possibilité de générer de telles représentations, mais cette fois géolocalisées et orthorectifiées, et visualisables dans le logiciel QGIS.
\subsection{Cases \bouton{Local alignment} et \bouton{Global alignment}}
Comme expliqué plus haut l'algorithme d'alignement procède en deux passes (coarse process et fine process). On peut voir le résultat de l'une ou l'autre (ou les deux) en cochant les cases \texttt{Global alignment} et \texttt{Local alignment}. La case \texttt{Global alignment} montre les résultats du coarse process alors que la case \texttt{Local alignment} montre le résultat du raffinement.  Une de ces deux cases doit impérativement être cochée. Nous suggérons de cocher systématiquement la case \texttt{Local alignment}. Il est à noter que dans certaines circonstances l'alignement global réussit mais l'alignement local échoue.

\subsection{Case \bouton{Maximum VIS/NIR time difference}}
Cette case permet de restreindre la liste des shoot points à aligner à ceux dont l'écart temporel de rang~1 vérifie $|dt|\leq \texttt{Maximum VIS/NIR time difference}$. La valeur par défaut est $0.5\,s$. Cela correspond aux images bien synchronisées. En pratique avec les fréquences de prise de vue de $2 \, s$ pour les images VIS et $3 \, s$ pour les images NIR cela revient à sélectionner un shoot point sur trois.


\section{Best mapping : n'aligner que les images utiles à la carte}\label{btn:step4_best_mapping}
Le cadre \bouton{Images aligned} propose deux choix, qui reprennent celui de l'étape~2 (§\ref{btn:step2_best_mapping}) : \bouton{Best mapping: only the images the map needs} (par défaut, avec un recouvrement de $75\,\%$) ou \bouton{All the synchronised images}. Parmi les shoot points bien synchronisés (case \texttt{Maximum VIS/NIR time difference}), IRDrone ne garde que ceux dont les empreintes au sol se recouvrent de la valeur choisie (case \bouton{Overlap between kept images}) : la première image du survol, puis, à chaque fois, la dernière image qui recouvre encore assez la précédente gardée, et la dernière image du survol. Toutes les images du survol servent à suivre le déplacement du drone ; seules les images bien synchronisées peuvent être gardées. Quand l'image qui serait gardée est mal synchronisée, IRDrone garde la plus proche bien synchronisée : la dernière avant elle ou, si son recouvrement est plus proche de la cible, la première après elle (recouvrement au plus 5 points sous la valeur choisie, 70\,\% pour 75\,\%). Le nombre d'images alignées reste ainsi proche de l'aperçu de l'étape~2 ; sans cela, chaque image mal synchronisée raccourcirait le pas et beaucoup plus d'images seraient gardées. La ligne sous les choix indique le nombre de shoot points qui seront alignés, par exemple \og 18 of 39 synchronised shoot points will be aligned \fg.

À droite, la carte du vol montre les shoot points : en vert ceux qui sont bien synchronisés et dans la plage d'images, en gris les autres ; les images gardées par le best mapping sont plus grosses, avec leur empreinte au sol en vert, ce qui montre le recouvrement. Elle se met à jour dès qu'un réglage change (plage, écart de temps, recouvrement, choix). Molette pour zoomer, glisser pour se déplacer, double-clic pour revoir tout le vol. Pendant l'alignement, la carte suit le calcul le long du vol : les empreintes des images alignées se remplissent de vert, celles en cours d'alignement (une par processus de calcul) sont entourées d'orange, celles qui restent à faire sont en pointillés ; le compte \og 7 of 18 done \fg{} s'affiche en bas de la carte. Le résultat reste affiché à la fin, jusqu'au prochain changement de réglage.

C'est en général le meilleur choix : beaucoup moins de calcul (l'alignement est l'étape la plus longue), moins de données, et une carte ODM souvent meilleure (trop d'images presque identiques ralentissent ODM et dégradent son résultat). Choisissez \bouton{All the synchronised images} pour aligner toutes les images synchronisées, par exemple pour les regarder une à une dans la revue. Le choix est enregistré dans \texttt{FlightAnalytics\textbackslash best\_mapping.json} au lancement de l'alignement ; l'étape~6 l'utilise pour la sélection des images d'ODM (§\ref{btn:step6_selection}).

La même sélection existe en ligne de commande : \texttt{run\_after\_preprocessing.py --best-mapping-overlap 75}.

\section{Bouton \bouton{Run alignment}}\label{btn:step4_run_alignment}
Lance la procédure d'alignement sur les shoot points vérifiant simultanément les deux contraintes 
\begin{itemize}
	\item $\texttt{first VIS image number} \leq XXXX \leq \texttt{last VIS image number}$,
	\item $|dt|\leq \texttt{Maximum VIS/NIR time difference}$,
\end{itemize} 
et, si \bouton{Best mapping: only the images the map needs} est choisi, seulement les images utiles à la carte (§\ref{btn:step4_best_mapping}).
Attention l'alignement des images est une procédure qui prend beaucoup de temps. Cela dépend bien entendu de la machine. Les  calculs sont parallélisés sur plusieurs processeurs si la machine est \og multi c\oe urs \fg. La barre de progression indique l'avancement des calculs ainsi qu'une estimation du temps de traitement restant. À titre d'exemple, sur un ordinateur portable avec un processeur Intel Core~i7 à quatre c\oe urs et 16\,Go de mémoire (3 workers d'alignement), l'alignement des images VIS et NIR d'un shoot point prend environ $15\,s$ : il faut donc environ une demi-heure pour traiter 120 images. Pour une grande mission, lancer ces calculs à l'heure du déjeuner reste une bonne idée.   

L'alignement ne démarre jamais tout seul, même quand l'étape~4 est ouverte par l'étape~3. Pendant le calcul, le bouton devient \bouton{Stop} (§\ref{btn:step4_stop}).

\section{Bouton \bouton{Stop}}\label{btn:step4_stop}
Pendant le calcul, \bouton{Run alignment} devient \bouton{Stop}. Après confirmation, le calcul et tous ses processus de travail s'arrêtent ; la fenêtre indique combien d'images ont été alignées. Les images déjà alignées sont gardées ; \bouton{Run alignment again} relance le calcul, et les points déjà faits réutilisent leur modèle de mouvement (plus rapide). Pour ne traiter que la suite, augmentez \texttt{First VIS image number}.

\section{Bouton \bouton{⚙} : réglages}\label{btn:step4_settings}
La roue dentée ouvre les réglages du calcul, gardés pour les sessions suivantes :
\begin{itemize}
	\item \texttt{Alignment worker processes} : nombre de shoot points alignés en parallèle. \texttt{Auto} utilise les c\oe urs physiques, dans la limite de la mémoire disponible au lancement (environ 3\,Go par processus). \bouton{Safe mode (1)} n'en garde qu'un : plus lent, mais le moins de mémoire. À utiliser si l'ordinateur se fige ou si le calcul est interrompu.
	\item \texttt{Log detail} : quantité de messages du calcul dans le journal de la fenêtre : avertissements seulement, étapes du calcul, ou détail des temps. Avec le détail des temps (\emph{Debug}), un réalignement à la main depuis la revue (section~\ref{btn:review_redo}) écrit aussi toutes ses images de diagnostic : plusieurs minutes par point au lieu d'environ une demi-minute.
\end{itemize}
\bouton{Save and close} (\bouton{Close} tant que rien n'a changé) garde les réglages. \bouton{↺}, affiché quand les réglages diffèrent des réglages par défaut, revient à \texttt{Auto} et aux avertissements seulement.

\section{Bouton \bouton{Close}}\label{btn:step4_close}
Ferme la fenêtre et rend la main au dialogue principal. Le bouton \bouton{?} ouvre cette aide à la page de ce chapitre. Il est grisé pendant l'alignement (arrêtez-le d'abord avec \bouton{Stop}).

\section{Bouton \bouton{Continue with the review}}\label{btn:step4_continue}
Apparaît en bas à droite, après \bouton{Close}, une fois des images alignées, et est mis en avant à la fin de l'alignement. Il ferme la fenêtre et ouvre l'étape~5, la revue des images alignées (chapitre~\ref{ch:review}). Quand l'alignement a été lancé depuis l'étape~3, la revue s'ouvre d'elle-même à la fin.

%==========================================================
\chapter{Step 5 -- Revue des images alignées}\label{ch:review}
La revue présente tous les points du relevé sur une ligne de temps pour vérifier l'alignement de chaque image (figure~\ref{fig:step4_review}). Elle s'ouvre avec \bouton{Continue with the review} de l'étape~4, ou avec l'étape~5 de la fenêtre principale (une fois des images alignées par l'étape~4). Elle remplace l'examen des images une par une dans l'explorateur. Rien n'est calculé à l'ouverture : l'étape compte comme faite une fois la revue fermée, les signalements et confirmations étant enregistrés au fur et à mesure.

\figureIRD[fig:step4_review]{step4_review.png}{Revue des images alignées : rideau entre visible et infrarouge aligné, fiche du point (décalage mesuré, image infrarouge associée), mini-carte, et en bas la ligne de temps des points (décalage résiduel, vignettes, état).}{0.95}

\section{Mesure automatique}\label{sec:review_measure}
Pour chaque point, le décalage résiduel entre les contours de l'image visible et ceux de l'infrarouge aligné est mesuré (en pixels de l'image complète, médiane sur 16 zones). Un point bien aligné mesure environ 1~px. Vert : moins de 3~px ; orange clair : entre 3 et 6~px, à regarder ; orange : plus de 6~px, probablement mal aligné. Une zone sans texture (eau, champ uniforme) ne se mesure pas. La fenêtre s'ouvre sur le point le plus douteux.

Les mesures et les points vérifiés sont gardés dans \texttt{ImgIRdrone/review.json} : rouvrir la revue ne refait pas les mesures, sauf pour les images refaites.

\section{Comparer le visible et l'infrarouge}\label{btn:review_compare}
Quatre façons de comparer, avec les boutons au-dessus de l'image ou les touches 1 à 4 ; la revue s'ouvre en \bouton{Overlay}. Un cadre autour de l'image est dans la couleur de l'état du point (vert aligné ou confirmé, rouge signalé, violet réaligné\ldots), les étiquettes sur l'image restent neutres (les rouges et cyans sont ceux des images) :
\begin{description}
	\item[\bouton{Swipe} (1)] un rideau sépare le visible (à gauche) de l'infrarouge aligné en fausses couleurs (à droite) ; routes, haies et toits doivent se prolonger à travers le rideau. Glisser le rideau pour le déplacer.
	\item[\bouton{Blink} (2)] visible et infrarouge alternent : une image mal alignée \og saute \fg.
	\item[\bouton{Overlay} (3)] visible en cyan, infrarouge en rouge (la végétation est naturellement rouge) ; un mauvais alignement double chaque contour.
	\item[\bouton{Side by side} (4)] les deux images côte à côte.
\end{description}
Molette : zoom (un décalage de quelques pixels ne se voit qu'en zoomant) ; glisser : déplacement ; double-clic : zoom $\times 4$ ou image entière. \bouton{Open full size} ouvre l'image alignée dans la visionneuse d'images de l'ordinateur.

\section{Situer l'image sur la carte}\label{btn:review_map}
La mini-carte (dans la fiche du point) montre le vol, le point courant et l'emprise au sol de l'image (le trait épais est le haut de l'image), dans la couleur de son état (vert aligné ou confirmé, rouge signalé\ldots) ; les shoot points que l'étape~4 n'a pas alignés sont de petits points gris, en taille normale quand on zoome fort. Un clic l'agrandit (\bouton{Map}, touche 5). Sous la mini-carte, la légende ne donne que les états présents, avec leur nombre. La grande carte affiche l'image alignée dans son emprise, par-dessus l'orthophoto de l'étape~6 quand elle existe et les photos aériennes en fond : l'infrarouge doit prolonger les routes et les haies du fond. Survoler un point montre sa vignette, un clic le sélectionne ; un clic hors des points ne laisse aucun point sélectionné sur la carte : la trajectoire seule, sans emprise ni image (le point reste celui de la fiche ; un clic sur un point, ou un déplacement dans le bandeau, en sélectionne un à nouveau) ; molette et glisser pour zoomer et se déplacer, double-clic pour voir tout le vol. Le menu \bouton{Layers}, le fond de carte hors ligne et le calcul de l'emprise sont décrits au §\ref{btn:home_map}. Sans position GPS dans les shoot points, la carte n'est pas disponible.

\paragraph{Aligner un point laissé de côté par l'étape~4.} Un clic droit sur un point gris de la carte (mini-carte ou grande carte) affiche son écart de temps avec l'infrarouge et propose \bouton{Align this point} : il est aligné comme le fait l'étape~4. Quand une image est irrécupérable, même à la main, un voisin écarté par le best mapping ou le filtre temporel peut être bon : alignez-le ainsi. Le point est aligné en arrière-plan (orange sur la carte, avancement dans le bandeau, \bouton{Stop aligning} pour arrêter) pendant que vous continuez la revue ; il apparaît ensuite dans la frise, à revoir comme les autres, et l'étape~6 le cartographie toujours (sauf si vous le signalez), quels que soient sa tolérance temporelle et son choix d'images. Le filtre temporel est le plus large de l'étape~4 ($|dt|<6$~s) ; un point sans image infrarouge assez proche dans le temps ne peut pas être aligné.

\figureIRD[fig:step4_review_map]{step4_review_map.png}{Mode \og Map \fg~de la revue : vol, emprise et image alignée du point courant, orthophoto de l'étape~6 sur les photos aériennes IGN.}{0.9}

\section{Se déplacer}\label{btn:review_navigation}
\begin{itemize}
	\item \bouton{◀} et \bouton{▶} (flèches du clavier) : point précédent, point suivant.
	\item \bouton{Next to review} (touche N) : prochain point qui demande un coup d'\oe il (réaligné à confirmer, mesuré douteux ou non mesuré). Un point mesuré aligné (vert) ne demande rien : il n'y a rien à marquer comme bon.
	\item La ligne de temps du bas : survoler montre l'image, un clic sélectionne le point ; molette pour zoomer, glisser pour se déplacer, double-clic pour tout le relevé.
	\item \bouton{Close} ferme la revue ; ce qui a été signalé et confirmé est déjà enregistré. \bouton{Continue with mapping} la ferme aussi, puis ouvre l'étape~6 (§\ref{btn:review_continue}).
	\item L'info-bulle du bouton \bouton{?} rappelle toutes les touches.
\end{itemize}

\section{Décider : signaler un point}\label{btn:review_decide}
\begin{description}
	\item[\bouton{Flag: leave out of the map} (touche F)] signale le point, \textbf{ce qui l'écarte de la carte} : mal aligné, mais aussi flou, ou avec une zone noire au bord de l'image infrarouge (drone incliné par un coup de vent : la caméra infrarouge ne voit pas tout le champ du visible). Un second appui retire le signalement. Le signalement est écrit tout de suite dans \texttt{ImgIRdrone\textbackslash shoot\_points.json}\footnote{Ce fichier ne contient pas les mêmes informations que le fichier \texttt{AerialPhotography\textbackslash shoot\_points.json}.} et l'étape~6 ignore ce point. Un point signalé peut être réaligné à la main (§\ref{btn:review_redo}) ; si cela ne le rattrape pas, le laisser signalé.
	\item[\bouton{Flag the N probably misaligned}] signale d'un coup tous les points dont le défaut d'alignement mesuré dépasse 6~px (points rouges \emph{probably misaligned}). Il n'apparaît que s'il y en a : la mesure ne voit ni un flou ni une zone noire, que vous signalez vous-même.
\end{description}

L'exemple ci-dessous montre le début du fichier \texttt{ImgIRdrone\textbackslash shoot\_points.json} qui stocke ces informations. Toutes les images ont été traitées en choisissant le critère peu discriminant $|dt|<4\,s$ et la méthode "au plus près". Cela correspond à 501 images VIS et 350 images NIR. Ici l'utilisateur a estimé visuellement que le shoot point 120 était mal aligné et l'a signalé, ce qui a fait basculer la valeur de la clé \texttt{failed} de \texttt{false} à \texttt{true}.
\begin{verbatim}
	{
		"mission_path": "C:\\Air-Mission\\FLY-20220524-1715-Chalivoy-Milon",
		"time_shift": -48.116,
		"n_VIS": 501,
		"n_NIR": 350,
		"pairing_method": "nearest_time",
		"created_at": "2026-10-29T20:44:13.203434",
		"schema_version": "1.0",
		"shoot_points": [
		{
			"FileName": "VIS_0119.dng",
			"dt": -0.5390370000000075,
			"failed": false
		},
		{
			"FileName": "VIS_0120.dng",
			"dt": 0.6902460000000019,
			"failed": true
		},
		{
			"FileName": "VIS_0121.dng",
			"dt": -1.3297540000000083,
			"failed": false
		},
		...
		]
	}	 
\end{verbatim}

Que faire ensuite des shoot points mal alignés ? Deux options :
\begin{itemize}
	\item \og aider \fg~l'algorithme d'IRDrone à les aligner : \bouton{Realign N flagged by hand} (§\ref{btn:review_redo}). Selon le résultat, le signalement est gardé ou retiré.
	\item les écarter des traitements ultérieurs : c'est l'effet par défaut du signalement, l'étape~6 ignore les points \texttt{"failed": true}.
\end{itemize}

\subsection{Exemple de succès d'alignement}
\figurePaireIRD{fig:goodalignment}
{step4_goodalignment_VIS.jpg}{Image dans le spectre visible.}
{step4_goodalignment_NIR.jpg}{Image la plus proche $|dt|=0.165~s$ dans le spectre infrarouge.}
{Cas de succès d'alignement : l'écart temporel est faible et l'attitude du drone est restée dans les spécifications des paramètres de vol.}

\figureIRD[fig:step4_goodalignment_VIR]{step4_goodalignment_VIR.jpg}{Alignement réussi. Les deux images ont été prises presque au même moment et en l'absence de vent. }{0.5}

\subsection{Exemple d'échec d'alignement}
Si le drone sort de sa zone des paramètres de vol, par exemple en cas de fortes rafales de vent qui l'obligent à accentuer ses angles de pitch et yaw, alors les centres de visée des deux caméras peuvent être tellement éloignés qu'il est impossible de trouver un alignement. L'alignement peut aussi être défaillant si l'écart temporel entre les prises de vue VIS et NIR est important (il peut être parfois supérieur à 2\,s) car le déplacement horizontal du drone est alors mal corrigé. L'alignement peut aussi être imparfait sans être un échec s'il y a un écart temporel et que les végétaux bougent sous l'action du vent. L'image VIR présente alors des \og fantômes \fg. Les figures \ref{fig:misalignment} et \ref{fig:step4_misalignment_VIR} montrent un cas d'échec de l'alignement lié à une forte rafale de vent qui déstabilise le drone. 

\figurePaireIRD{fig:misalignment}
{step4_misalignment_VIS.jpg}{Image dans le spectre visible.}
{Step4_misalignment_NIR.jpg}{Image la plus proche $|dt|=1.256~s$ dans le spectre infrarouge.}
{Cas d'échec d'alignement lié à une attitude du drone hors des spécifications des paramètres de vol et liée à une rafale de vent au moment de la prise de vue.}

\figureIRD[fig:step4_misalignment_VIR]{step4_misalignment_VIR.jpg}{Échec de l'alignement. L'effet de \og trou \fg est particulièrement caractéristique. Il peut aisément être repéré par l'utilisateur sur les images VIR. }{0.5}

\section{Réalignement à la main}\label{btn:review_redo}
Le bouton \bouton{Realign N flagged by hand} rattrape les shoot points signalés mal alignés. Il n'y a pas de miracle : si une rafale de vent a écarté le drone de sa zone de vol, les deux caméras ne voient pas la même scène et aucun alignement n'est possible. Mais parfois l'information existe et il suffit de \og pousser \fg~les paramètres de l'alignement grossier. En conditions normales une dizaine de shoot points au plus sont à reprendre.

\subsection*{Déroulement}
\begin{enumerate}
	\item Signaler les points mal alignés, puis appuyer sur \bouton{Realign N flagged by hand}. La fenêtre d'alignement s'ouvre avec tous les points signalés (figure~\ref{fig:step4_align_by_hand}).
	\item Pour chaque point, déplacer l'image infrarouge sur l'image visible jusqu'à faire coïncider les mêmes éléments (voir ci-dessous), puis \bouton{Placed, next} (touche Entrée). \textbf{L'alignement de ce point démarre aussitôt, en arrière-plan}, pendant que vous placez le suivant : pas d'attente entre deux points (le point passe à \emph{⚙ Computing} dans la liste et ne peut plus être déplacé). Les images des points suivants sont chargées en arrière-plan. \bouton{Skip this point} (touche S) laisse un point de côté : il reste signalé, donc écarté de la carte.
	\item \bouton{Compute N alignments} envoie au calcul les points placés sans \bouton{Placed, next} (quand tout est déjà envoyé, \bouton{Cancel} devient \bouton{Close}) et ferme la fenêtre. Les alignements sont calculés en arrière-plan, plusieurs à la fois (autant que de workers d'alignement de l'étape~4, bouton engrenage). Compter environ une demi-minute par point. Le bandeau du haut de la revue suit l'avancement : points terminés, points en cours avec leur étape et le temps écoulé, points en attente. \textbf{La revue reste utilisable} : on peut continuer à vérifier les autres points, ou faire autre chose. Les points en cours de calcul ne peuvent pas être signalés ni dé-signalés. \bouton{Stop} ne lance plus de nouvel alignement : ceux en cours se terminent, les autres restent signalés. La revue ne peut pas être fermée tant que des alignements sont en cours.
	\item Chaque point réaligné passe en violet, \emph{Realigned, to confirm} : il reste signalé (l'étape~6 l'ignore). Quand tout est calculé, la revue affiche le premier. Le vérifier avec le rideau, le clignotement ou la superposition et répondre \bouton{Yes, aligned} (touche Y) ou \bouton{Still misaligned} : le point suivant à confirmer s'affiche alors. \bouton{All N aligned} les confirme tous d'un coup. La touche N (point suivant à revoir) y va en premier. Les points à confirmer sont mémorisés si la revue est fermée.
	\item Avec \bouton{Yes, aligned}, la clé \texttt{"failed"} du point passe à \texttt{false} dans \texttt{shoot\_points.json} et le fichier \texttt{shoot\_points\_failed\_summary.json} est mis à jour (une copie \texttt{shoot\_points.json.bak} est conservée). Avec \bouton{Still misaligned}, le point reste signalé et pourra être réaligné de nouveau ; s'il ne peut pas être rattrapé (image infrarouge trop décalée, flou), il suffit de le laisser signalé : \textbf{il est écarté de la carte} (étape~6).
\end{enumerate}

Les sorties des alignements précédents sont archivées dans le dossier \texttt{VIS\_XXXX\_align\_traces\_manual\textbackslash original}. Le point retrouve les images qu'il avait (celles cochées à l'étape~4, section~\ref{btn:step4_options}) ; les images de diagnostic de l'alignement (dossier \texttt{Tech}, plusieurs minutes de calcul par point) ne sont écrites que si \texttt{Log detail} est réglé sur \emph{Debug} dans les réglages de l'étape~4 (bouton engrenage, section~\ref{btn:step4_settings}), ou depuis la ligne de commande (ci-dessous).

\subsection*{La fenêtre d'alignement}
\figureIRD[fig:step4_align_by_hand]{step4_align_by_hand.png}{Alignement à la main : les points signalés à gauche (coche : placé, rond : à placer), les deux images superposées au centre (ici le visible en cyan et l'infrarouge en rouge), les angles à droite. Ici l'infrarouge est encore environ 140~px trop haut : les contours sont doublés.}{0.9}

\begin{itemize}
	\item À gauche, la liste des points : un rond pour les points à placer, une coche pour les points placés, un tiret pour les points laissés de côté. Un clic va sur un point.
	\item Au centre, les deux images, affichées comme dans la revue : \bouton{Overlay} (les deux images en une : des contours doublés signalent un mauvais alignement ; affichage par défaut). \bouton{Colours} choisit ses couleurs, gardées pour la fois suivante : visible en couleur et infrarouge en gris (par défaut ; avec \bouton{Edges only}, les contours du visible sont verts), visible en vert et infrarouge en gris, visible en bleu et infrarouge en rouge, ou visible en cyan et infrarouge en rouge (comme dans la revue, proche de l'image VIR). \bouton{Swipe}, \bouton{Blink}, \bouton{Side by side} (touches 1 à 4). \bouton{Edges only} (touche G) montre les contours des deux images, ce qui aide sur les champs et les forêts où l'infrarouge et le visible se ressemblent peu. Molette pour zoomer, glisser pour se déplacer, double-clic pour voir toute l'image.
	\item Pour déplacer l'infrarouge, \textbf{vite puis finement} : d'abord les \textbf{curseurs} de \texttt{Yaw}, \texttt{Pitch} et \texttt{Roll} (panneau de droite) pour un grand déplacement, puis les boutons à leurs extrémités pour affiner, dans le sens qu'ils montrent : \bouton{◀}~\bouton{▶} à gauche et à droite (lacet), \bouton{▲}~\bouton{▼} en haut et en bas (tangage), \bouton{↺}~\bouton{↻} rotation (roulis) ; maintenus enfoncés, ils se répètent. Les \textbf{flèches du clavier} font la même chose que ces boutons (0,05\textdegree{} par appui, environ 2 à 3 pixels de l'image complète ; avec Maj, 10 fois plus), où que soit le curseur dans la fenêtre, même dans un champ de valeur ; Q et E le font tourner. On peut aussi \textbf{faire glisser l'infrarouge avec le bouton droit de la souris} (ou Maj + glisser) dans l'image, ou taper une valeur dans un champ.
	\item \bouton{Fit} (touche F) mesure le décalage entre les contours des deux images (tuile par tuile) et déplace l'infrarouge d'autant (lacet et tangage). Il fonctionne dès que les images se recouvrent à peu près ; s'il ne trouve pas de meilleure position, il le dit et ne change rien. \bouton{Reset} revient aux angles de l'alignement précédent.
	\item Sous ces boutons, le décalage mesuré entre les deux images, en pixels de l'image complète : sous 10~px le point est \emph{well placed}.
	\item À peu près suffit : ces angles ne sont qu'un point de départ, l'alignement grossier et l'affinage pyramidal repartent de là (§\ref{sec:ann_algorithme}).
	\item Un point dont les images ne peuvent pas être lues est marqué d'un signe d'avertissement et laissé de côté.
\end{itemize}

\subsection*{En ligne de commande : la fenêtre à curseurs}
Chaque point aligné a un script \texttt{VIS\_XXXX\_REDO.bat} (Windows) ou \texttt{VIS\_XXXX\_REDO.sh} dans \texttt{ImgIRdrone}, qui relance son alignement avec une fenêtre à curseurs (matplotlib) : c'est l'ancienne procédure, toujours utilisable en dehors de l'assistant. Après un réalignement dans l'assistant, le script repart des angles placés à la main.


\figureIRD[fig:redo_info]{step4_redo_info.png}{Fenêtre d'information de la fénêtre d'aide manuelle à l'alignement. }{0.5}

La fenêtre affiche l'image VIS et l'image NIR, avec cinq curseurs : \texttt{YAW}, \texttt{PITCH}, \texttt{ROLL}, \texttt{Absgrad visualize} et \texttt{Alpha}.
\begin{itemize}
	\item Le curseur \texttt{Alpha} choisit ce que l'on voit. Entre $-1$ (tout à gauche) et $-0{,}5$ (petit trait rouge médian), un rectangle de l'image NIR, légèrement transparent, est superposé à l'image VIS : plus \texttt{Alpha} est proche de $-1$, plus ce rectangle est grand. La valeur initiale est $-0{,}75$.
	\item Les curseurs \texttt{YAW}, \texttt{PITCH} et \texttt{ROLL} déplacent l'image NIR. La méthode consiste à modifier un angle, puis à faire varier \texttt{Alpha} pour observer le résultat, et à recommencer avec les trois angles. La clé est de repérer des éléments identiques dans les deux bandes spectrales (arbre, bord de chemin, bâtiment) et de les faire coïncider, d'abord au centre puis sur les bords du rectangle.
	\item Lorsque \texttt{Alpha} est déplacé vers la droite, la fenêtre n'affiche plus un rectangle NIR : entre $-0{,}5$ et $0$ elle montre la différence entre les deux images, et entre $0$ et $1$ un fondu enchaîné VIS/NIR. Cet usage est plus rapide mais réservé à un opérateur expérimenté.
	\item \texttt{Absgrad visualize} (valeur initiale $0{,}3$) affiche les gradients des deux images en deux couleurs lorsqu'il dépasse $0{,}5$ : les contours coïncident quand l'alignement est bon.\\
\end{itemize}

Lorsque l'utilisateur considère que l'alignement manuel est satisfaisant il tape la touche Q et la procédure d'alignement semi automatique démarre. \\


Comme pour tout outil de ce type, c'est en l'expérimentant que l'on se l'approprie ; une fois la méthode comprise, le réalignement est rapide et efficace.
Voici quelques étapes en image de la procédure
\footnote{L'outil est riche et offre plusieurs possibilités de visualiser les couches. Donc il n'y a pas une méthode unique pour arriver au résultat souhaité. Le slider \texttt{Absgrad visualize} permet de changer les représentations. Par exemple en le déplaçant vers la droite on peut obtenir un effet de contours. Ce n'est pas toujours évident à interpréter avec des forêts par exemple et encore moins dans des zones sans relief et nues.}. 
Sur l'image il y a une route qui traverse en diagonale. 
Il y a aussi un gros arbre en haut à droite, un petit arbuste en bas à gauche le long de la route  ainsi qu'un fossé en haut à gauche (\og trait noir \fg). Ce sont ces éléments caractéristiques que l'on va essayer d'aligner. 
À noter qu'à l'initialisation la couche NIR est légèrement transparente
\footnote{En déplaçant le slider \texttt{Alpha} vers la droite entre 0 et 1 la couche NIR passe de totalement transparente à opaque.}. 
Il n'y a pas vraiment d'ordre pour arriver à un résultat satisfaisant. 
Ne cherchez pas la perfection  car on n'applique que des transformations globales linéaires. Le code affinera globalement puis localement.
\figurePaireIRD{fig:redo_A}
{step4_redo_1.png}{Situation de départ après un passage dans la procédure associée au bouton \texttt{run alignement}.}
{step4_redo_2.png}{On a joué sur la rotation, qui passe de $0$ à $-2.46$, pour aligner la route NIR parallèlement à la route VIS. Attention utiliser ce slider délicatement. Lorsqu'il n'y a pas de ligne évidente comme ici la route, on utilise plutôt la rotation en dernier pour mettre la \og touche finale \fg.}
{Première modification en jouant sur le slider de rotation (\texttt{ROLL}).}
\figurePaireIRD{fig:redo_B}
{step4_redo_3.png}{Translation verticale vers le bas avec le slider \texttt{PITCH}. La valeur passe de  $-1.14$ à $2.04$. C'est une amplitude relativement importante et surtout qui change de signe.}
{step4_redo_4.png}{Translation horizontale avec le slider \texttt{YAW}. La valeur passe de $0.96$ à $1.8$.}
{On affine en jouant sur le curseur \texttt{Alpha} (en diminuant puis en agrandissant) et en observant l'alignement des éléments caractéristiques. La sensation de léger flou à la fin est liée à la transparence et aux petits défauts d'alignement locaux qui seront corrigés par la procédure d'alignement fin.}

\section{Bouton \bouton{Continue with mapping}}\label{btn:review_continue}
En bas à droite, après \bouton{Close}. Il ferme la revue et ouvre l'étape~6. La carte, longue, n'y démarre pas toute seule : elle se lance avec \bouton{Build the map\ldots}. Les points signalés, et les points réalignés pas encore confirmés, sont écartés de la carte. Tant que des alignements sont en calcul, la revue ne peut pas être fermée : attendez-les (ou \bouton{Stop}).



%==========================================================
\chapter{Step 6 -- Cartographie et géoréférencement}\label{ch:step6}
\section{Principe}
Les images multispectrales produites à l'étape~4 sont des images \og à plat \fg~: elles ne sont pas positionnées sur le terrain.

\figurePaireIRD{fig:step6_ortho}
{step5_VIS_plan.png}{Image corrigée des distorsions de l'optique (DJI Mavic Air 2 et calibrage de Zhang).}
{step5_VIS_ortho.png}{Image Orthorectifiée avec un modèle de terrain au pas de $10\,m$.}
{L'effet de l'orthorectification est bien visible sur cette image où le terrain présente une forte déclivité (voir figure \ref{fig:step6_topo}).}

 L'étape~6 les rend exploitables dans un système d'information géographique (SIG) comme QGIS. Elle propose trois traitements indépendants, un par bouton :
\begin{itemize}
	\item \bouton{Build the map\ldots} (le bouton principal) : assemblage de l'ensemble des images retenues en une seule \emph{orthophoto} à 4 bandes (rouge, vert, bleu, proche infrarouge) avec le logiciel libre OpenDroneMap (ODM), puis le modèle de terrain et les courbes de niveau ;
	\item \bouton{Terrain model and contours\ldots} : construction à nouveau du MNT et des \emph{courbes de niveau}, par exemple avec un autre pas de grille, à partir d'altitudes fournies par un service en ligne (IGN) ;
	\item \bouton{Georeference images\ldots} : \emph{géoréférencement} et \emph{orthorectification} d'images isolées. Chaque image est projetée sur un modèle numérique de terrain (MNT) puis enregistrée au format GeoTIFF en coordonnées UTM.
\end{itemize}
Ces traitements s'appuient sur les métadonnées de vol enregistrées dans les fichiers \texttt{.exif} \og compagnons \fg~(position GPS, altitude, angles de lacet, tangage et roulis, voir étape~1).

\figureIRD{step5_window.png}{Fen\^etre de l'étape~6.}{0.6}

La fenêtre est organisée comme celle de l'étape~3. De haut en bas : le nom de l'étape et de la mission ; les étapes de la carte (\texttt{Image selection}, \texttt{Orthophoto (OpenDroneMap)}, \texttt{Terrain model and contours} : $\circ$ à faire, $\bullet$ en cours, \checkmark{} fait, $\times$ échec, avec un détail sous chacune) ; une zone de messages (journal) ; les trois boutons, \bouton{Build the map\ldots} en premier ; une barre de progression précédée d'un intitulé qui indique l'opération en cours ; une zone d'aperçu d'image ; et les boutons \bouton{?} et \bouton{Close}, puis \bouton{Continue with archiving} une fois la carte construite.

Si la carte a déjà été construite, les étapes sont déjà cochées et l'orthophoto s'affiche dès l'ouverture de la fenêtre. \bouton{Terrain model and contours\ldots} reste grisé tant que la carte n'existe pas.

\textbf{Prérequis}
\begin{itemize}
	\item l'étape~4 doit avoir été exécutée pour les images à traiter (dossier \texttt{ImgIRdrone\textbackslash imgMultispectral}) ;
	\item pour \bouton{Build the map\ldots} : le logiciel \emph{Docker Desktop} doit être installé et démarré ;
	\item pour \bouton{Terrain model and contours\ldots} (et pour l'altitude du sol lors du géoréférencement) : une connexion Internet est nécessaire pour interroger le service d'altitudes.
\end{itemize}

\section{Bouton \bouton{Georeference images\ldots}}\label{btn:step6_georef}
Ce bouton géoréférence et orthorectifie des images choisies par l'utilisateur.

Un message rappelle les types d'images acceptés, puis une fenêtre de sélection de fichiers s'ouvre. Tous les fichiers choisis doivent être du même type :
\begin{itemize}
	\item des images DNG de la caméra DJI (spectre visible) ;
	\item des images TIFF multispectrales par \textbf{groupes de quatre} : \texttt{IRD\_XXXX\_1.tif} à \texttt{IRD\_XXXX\_4.tif} (rouge, vert, bleu, proche infrarouge). Il faut sélectionner les quatre fichiers ; un groupe incomplet est rejeté avec un avertissement.
\end{itemize}

Dans une mission visible seul, le message et la fenêtre de sélection ne proposent que les images DNG DJI.

\subsection*{Décalages de géoréférencement}
Une fois les images choisies, une fenêtre demande trois corrections, toutes égales à \textbf{0} par défaut :
\begin{itemize}
	\item le \emph{décalage de lacet} (en degrés), qui corrige une erreur systématique du cap du drone ;
	\item le \emph{décalage UTM est} (en mètres) ;
	\item le \emph{décalage UTM nord} (en mètres).
\end{itemize}

Ces valeurs s'appliquent à toutes les images traitées lors de ce lancement. Les décalages sont utiles car les données GPS
\footnote{La précision de positionnement pourrait être grandement améliorée à l'aide d'une balise RTK mais ce type de matériel est onéreux.} et d'attitude du drone DJI ont une précision limitée, tout particulièrement si l'étalonnage du GPS et/ou le calibrage de l'orientation de la caméra n'ont pas été fait avant le décollage. Ces valeurs \og d'offset \fg ne peuvent être déterminer qu'expérimentalement, par exemple en utilisant QGIS pour affiner le positionnement. 

 Le bouton \emph{Cancel} annule le géoréférencement.

\subsection*{Ce que fait le logiciel pour chaque image}
\begin{enumerate}
	\item Les images DNG sont converties en TIFF avec RawTherapee. Pour les images multispectrales, les quatre bandes sont normalisées entre 0 et 1 puis empilées en une image à 4 canaux.
	\item Les métadonnées de vol sont lues dans le fichier \texttt{.exif} compagnon (créé si nécessaire). Elles donnent la position du drone, son altitude au-dessus du sol et ses angles d'attitude. Si l'altitude du point de décollage est inconnue, elle peut être demandée à l'utilisateur.
	\item L'emprise de l'image au sol est estimée à partir de la géométrie de la caméra, puis un maillage du terrain est construit. Les altitudes des points du maillage sont demandées à l'IGN.
	\item Chaque point du maillage est projeté dans l'image. L'image est rééchantillonnée sur une grille régulière au sol (orthorectification).
	\item Le résultat est enregistré au format GeoTIFF (coordonnées UTM).
\end{enumerate}

\subsection*{Fichiers produits}
Les fichiers sont créés dans le sous-dossier \texttt{geo\_ref} du dossier contenant les images sélectionnées. Un fichier existant n'est jamais écrasé : un suffixe \texttt{\_1}, \texttt{\_2}, \dots~est ajouté au nom (utile par exemple si le fichier précédent est ouvert dans QGIS).
\begin{itemize}
	\item \texttt{*\_georef.tif} : l'image orthorectifiée, à ouvrir dans QGIS. Pour un quadruplet multispectral le nom est \texttt{IRD\_XXXX\_0\_georef.tif} et l'image comporte 4 bandes (Red, Green, Blue, NIR) ;
	\item \texttt{*\_MNT.png} et \texttt{*\_Grid.png} : graphiques de contrôle du terrain et du maillage ;
	\item \texttt{*\_dtm\_map.png} : carte 2D ombrée du MNT sous l'image, avec courbes de niveau, quadrillage UTM et points d'altitude IGN 	(même présentation que pour le mapping).
	\end{itemize}
	Pendant le calcul la fenêtre ne répond plus : c'est normal. L'image en cours de traitement s'affiche dans la zone d'aperçu et le journal détaille l'avancement.

	La précision du placement dépend de la qualité des métadonnées de vol (GPS, lacet) et de celle du MNT. Un éventuel décalage systématique peut être corrigé grâce à la fenêtre des décalages (voir ci-dessus).
	
\figureIRD{step5_run_Geo_ref.png}{Fenêtre pendant le géoréférencement d'une image. C'est l'image de départ non orthorectifiée qui est affichée.}{0.6}


\section{Bouton \bouton{Build the map\ldots}}\label{btn:step6_mapping}
Ce bouton assemble les images alignées de l'étape~4 en une orthophoto unique de toute la zone survolée, grâce au logiciel OpenDroneMap (ODM). ODM s'exécute dans un conteneur \emph{Docker} : l'utilisateur n'a rien à configurer, mais Docker Desktop doit être démarré (voir plus bas).

\subsection{Choix des images}\label{btn:step6_selection}
En cliquant sur le bouton, une fenêtre de choix des critères s'ouvre. Une note en bas rappelle que le calcul peut être long et utiliser plusieurs gigaoctets de mémoire et d'espace disque. \bouton{Build the map} le lance aussitôt ; \bouton{Cancel} ferme la fenêtre sans rien lancer. Une fois une carte construite, la fenêtre propose les réglages de la dernière exécution (sinon les valeurs par défaut ci-dessous ; le best mapping choisi à nouveau aux étapes~2 ou~4 après cette exécution l'emporte). Les images peuvent venir de tout le vol, pas seulement de la phase de survol de l'étape~2.

\begin{itemize}
	\item \texttt{Maximum |dt| (NIR rank 1)} : écart temporel maximal admis entre l'image VIS et l'image NIR la plus proche. C'est le même critère que le \texttt{Maximum VIS/NIR time difference} de l'étape~4. Valeur par défaut : $0.5\,s$ (plage $0.05$ à $6\,s$).
	\item \texttt{Best mapping: select by image overlap (recommended)} : choix par défaut. Ne garde que les images dont les empreintes au sol se recouvrent de la valeur indiquée (\texttt{Overlap between successive images}, par défaut la valeur choisie à l'étape~2 ou~4, $75\,\%$ sinon). C'est la même sélection qu'aux étapes~2 et~4 (§\ref{btn:step2_best_mapping}). Si l'étape~4 n'a aligné que les images utiles à la carte, elles sont toutes gardées ; si elle a tout aligné, seules les images utiles à la carte sont données à ODM. Une image déclarée mal alignée dans la revue est écartée et la suivante prend sa place.
	\item \texttt{Select one image every n eligible images} : permet de ne retenir qu'une image sur $n$ ($n$ de 1 à 4, valeur par défaut 3) afin de limiter le temps de calcul et le recouvrement. C'est le choix par défaut quand le best mapping a été décoché à l'étape~2 ou~4.
	\item \texttt{All the aligned images (every point of the review not flagged)} : ignore les critères ci-dessus : toutes les images alignées par l'étape~4 (ou depuis la revue) et non signalées sont cartographiées, quels que soient leur écart de temps et le recouvrement. Le plus d'images, la cartographie la plus longue. Ce choix est proposé à nouveau la fois suivante s'il a été utilisé.
\end{itemize}

\figureIRD{step5_odm_selection.png}{Fenêtre des critères de sélection des images pour le mapping.}{0.4}

À droite, comme à l'étape~4, la carte du vol montre les images retenues pour la carte, en vert avec leurs empreintes au sol (leur recouvrement se voit), les images éligibles écartées par la sélection en gris, et les autres en petits points. Elle est mise à jour dès qu'un réglage change ($|dt|$, recouvrement, $n$, choix) ; le compte sous la carte donne le nombre d'images retenues parmi les éligibles. Molette pour zoomer, glisser pour se déplacer, double-clic pour revoir tout le vol.

\textbf{OpenDroneMap : version et GPU.} Sous les choix, le cadre \texttt{OpenDroneMap} indique, après une vérification de quelques secondes (Docker, Internet), la version d'ODM installée sur l'ordinateur et si une version plus récente a été publiée sur Docker Hub :
\begin{itemize}
	\item \texttt{Download the newer version before the map} : affichée seulement quand une version plus récente existe, décochée par défaut. Cochée, la dernière image d'ODM est téléchargée (\texttt{docker pull}, de quelques centaines de mégaoctets à quelques gigaoctets) avant le calcul. Une nouvelle version peut donner une carte un peu différente : elle n'est donc jamais installée sans qu'on le demande. Sans Internet, le cadre indique que la dernière version n'a pas pu être vérifiée et la version installée est utilisée.
	\item \texttt{Use the NVIDIA GPU (faster image matching)} : disponible, et cochée, seulement quand l'ordinateur a une carte graphique NVIDIA \emph{et} que Docker peut la donner à un conteneur (le logiciel lance un conteneur de test avec \texttt{--gpus all}). ODM utilise alors sa version GPU (\texttt{opendronemap/odm:gpu}, environ 3\,Go téléchargés la première fois), qui détecte et apparie les points caractéristiques des images sur la carte graphique ; le reste du calcul est le même. Sinon la case est grisée et la raison est donnée (pas de carte NVIDIA ; sous Linux, le \emph{NVIDIA Container Toolkit} manque ; sous Windows, Docker Desktop doit utiliser WSL~2 avec un pilote NVIDIA récent) ; son info-bulle donne la réponse de Docker. Si le GPU n'est pas utilisable au lancement du calcul, ODM tourne sur le processeur et le journal le dit.
\end{itemize}
Le journal donne la version d'ODM utilisée, et si elle a tourné sur le GPU ou sur le processeur.


Les images retenues (\emph{eligible images}) sont les shoot points qui vérifient simultanément :
\begin{enumerate}
	\item ils n'ont pas été signalés mal alignés dans la revue, étape~5 (clé \texttt{"failed": false}, voir §\ref{btn:review_decide}) ;
	\item l'écart $|dt|$ avec l'image NIR la plus proche est inférieur ou égal à la tolérance ;
	\item l'image VIS correspondante appartient à la phase de vol de la mission choisie à l'étape~2 (les images de décollage et de synchronisation sont donc exclues).
\end{enumerate}
Parmi elles, le best mapping garde les images utiles à la carte (ou, avec l'autre critère, une image éligible sur $n$).

\subsection{Lancement du calcul}
\bouton{Build the map} lance le calcul sans autre question (la sélection est écrite dans le journal) :
\begin{enumerate}
	\item le logiciel vérifie que Docker Desktop est disponible. Sinon un message invite à le démarrer puis à choisir \texttt{Retry}. ODM est ensuite téléchargé s'il n'est pas encore sur l'ordinateur, ou mis à jour si on l'a demandé (voir plus haut) ;
	\item la liste des images retenues est enregistrée dans \texttt{mapping\_MULTI\textbackslash img\_best\_mappping.json} ;
	\item la résolution de l'orthophoto est estimée automatiquement : c'est la taille au sol d'un pixel (GSD, en cm) obtenue à partir de l'altitude de vol au-dessus du sol et de la focale de la caméra, médiane sur les images retenues, arrondie au dixième de cm supérieur ;
	\item les images TIFF retenues sont rendues accessibles à ODM dans \texttt{mapping\_MULTI\textbackslash images} par des \emph{liens physiques} (hard links) : aucune image n'est copiée et l'espace disque n'est donc pas doublé. Si le disque ne gère pas les liens physiques (disque externe en exFAT ou FAT32, certains partages réseau), les images sont copiées, ce qui utilise davantage d'espace disque ; le journal l'indique ;
	\item ODM est lancé. Le journal de la fenêtre affiche les messages principaux ; la totalité de la sortie est enregistrée dans un fichier \texttt{mapping\_MULTI\textbackslash odm\_run\_<date>.log} ;
	\item à la fin, l'orthophoto est contrôlée (4 bandes, géoréférencée), les résultats utiles sont archivés dans \texttt{mapping\_result} et les dossiers de travail d'ODM, très volumineux, sont supprimés.
\end{enumerate}
Pendant le calcul la barre de progression est indéterminée, les boutons et la fermeture de la fenêtre sont inactifs. À l'issue du calcul la durée est indiquée dans le journal.

\textbf{Résultats}\label{sec:step6_mapping_result}

Les résultats sont dans \texttt{mapping\_MULTI\textbackslash mapping\_result} :
\begin{itemize}
	\item \texttt{odm\_orthophoto.tif} : l'orthophoto à 4 bandes (rouge, vert, bleu, proche infrarouge), en UTM ;
	\item \texttt{report.pdf} : rapport de qualité produit par ODM ;
	\item \texttt{odm\_25dmesh.ply} : maillage 2,5D du terrain ;
	\item \texttt{shots.geojson} (et \texttt{shots.qml}) : position estimée de chaque prise de vue, avec son numéro de shoot point, à ouvrir dans QGIS ;  \texttt{shots.qml} est un fichier de mise en forme.
	\item des fichiers techniques (\texttt{ortho.png}, informations de géoréférencement).
\end{itemize}
Un aperçu en couleurs (bandes rouge, vert, bleu) de l'orthophoto s'affiche ensuite dans la fenêtre, puis le calcul du MNT et des courbes de niveau démarre automatiquement (c'est équivalent au bouton \bouton{Terrain model and contours\ldots}, voir \ref{btn:step6_dtm}).

\figureIRD{step5_orthophoto_preview.png}{Aperçu de l'orthophoto dans la fenêtre du Step~6.}{0.6}

\subsection{Précautions}
\begin{itemize}
	\item Les résultats existants ne sont \textbf{jamais écrasés}. Quand une carte est déjà construite, elle est conservée : son dossier est renommé \texttt{mapping\_MULTI\_<date de son exécution>} et la nouvelle carte est construite dans un nouveau \texttt{mapping\_MULTI} (avec la même configuration de caméra) ; la fenêtre de choix des images l'indique. Si la nouvelle exécution échoue, la carte précédente est remise dans \texttt{mapping\_MULTI} et l'exécution échouée, avec son journal, est conservée dans \texttt{mapping\_MULTI\_failed\_<date>}. L'étape~7 conserve ces dossiers.
	\item Docker doit utiliser le moteur \og desktop-linux \fg~de Docker Desktop. Le premier lancement télécharge l'image ODM (plusieurs gigaoctets) : prévoir une connexion Internet.
	\item Si certaines images retenues sont incomplètes (quadruplet TIFF manquant ou de tailles différentes), le calcul est refusé et le journal indique les fichiers en cause.
\end{itemize}

\section{Bouton \bouton{Terrain model and contours\ldots}}\label{btn:step6_dtm}
Ce bouton construit un modèle numérique de terrain (MNT) et des courbes de niveau sur l'emprise de l'orthophoto. Il est lancé \emph{automatiquement} à la fin d'un mapping réussi (avec un pas de grille de 10\,m) ; ce bouton permet de le refaire, par exemple avec un autre pas. Il nécessite l'orthophoto \texttt{mapping\_MULTI\textbackslash mapping\_result\textbackslash odm\_orthophoto.tif}.

Une fenêtre demande la valeur \texttt{Altitude grid step}, c'est-à-dire l'espacement (en mètres) des points de la grille régulière dont les altitudes sont demandées à l'IGN (service d'altitudes de l'Institut national de l'information géographique et forestière) ; en cas d'indisponibilité, le service OpenTopoData est utilisé avec une moindre précision en France. Plus le pas est petit, plus le nombre de requêtes est grand. Le pas par défaut est de 10\,m (plage 1 à 100\,m). Le nombre de points est limité à 5000 : au-delà, le logiciel augmente automatiquement le pas. \bouton{Build} lance le calcul, \bouton{Cancel} ne lance rien.

\subsection*{Étapes du calcul}
\begin{enumerate}
	\item lecture de l'emprise de l'orthophoto (UTM) ;
	\item construction d'une grille régulière qui couvre l'emprise ;
	\item requêtes d'altitude (aucun point manquant n'est toléré : sinon le calcul échoue et un message l'indique) ;
	\item interpolation par une surface lisse (spline quadratique) puis échantillonnage sur un raster à 1\,m ;
	\item export du MNT, des courbes de niveau (équidistance de 1\,m, courbes maîtresses tous les 10\,m), d'une carte 2D ombrée et d'une vue 3D.
\end{enumerate}
En cas d'échec, l'orthophoto n'est pas affectée.

\subsection{Fichiers produits}
Les fichiers sont créés à côté de l'orthophoto, dans \texttt{mapping\_result} (jamais écrasés, un suffixe numérique est ajouté si besoin) :
\begin{itemize}
	\item \texttt{ortho\_dtm.tif} : MNT (altitude en mètres, UTM) ;
	\item \texttt{ortho\_contours.geojson} et \texttt{ortho\_contours.qml} : courbes de niveau et leur style pour QGIS ;
	\item \texttt{ortho\_dtm\_map.png} : carte 2D ombrée avec courbes de niveau, position des points de la grille (point noir) et quadrillage UTM (bleu clair);
	\item \texttt{ortho\_dtm\_3d.png} : surface 3D ;
	\item \texttt{ortho\_dtm\_grid.json} : points de la grille avec leur altitude et la source des altitudes.
\end{itemize}

\figurePaireIRD{fig:step6_topo}
{step5_dtm_map.png}{Carte 2D du MNT avec courbes de niveau et quadrillage UTM.}
{step5_dtm_3d.png}{Surface 3D du MNT.}
{Représentation topographique du terrain à partir du modèle numérique de terrain déduit des altimétries IGN.}

\subsection{Ouverture dans QGIS}
Dans la fenêtre de la carte QGIS, faire glisser l'orthophoto \texttt{odm\_orthophoto.tif}, les courbes de niveau \texttt{ortho\_contours.geojson} et les points des prises de vue \texttt{shots.geojson}    (les styles \texttt{.qml} homonymes sont appliqués automatiquement). \\

Quelques actions simples avec une orthophotographie à 4 bandes spectrales red,green,blue,nir

\begin{itemize}
	\item Image couleur dans le spectre visible : clic droit sur la couche> \underline{P}ropriétés>Symbologie puis associer les bandes 1, 2, 3 au rouge, vert, bleu pour une vue en couleurs naturelles. C'est l'option d'affichage par défaut d'une image multispectrale.
	\item Image fausse couleur VIR : clic droit>\underline{P}ropriétés>Symbologie puis associer les bandes 4~,1~,2 au rouge, vert, bleu. 
	\item Carte de l'indice NDVI\footnote{L'indice NDVI pour \textit{Normalized Difference Vegetation Index} mesure la vigueur et la densité de la végétation. On détermine cet indice par la formule simple : $NDVI = \frac{NIR - RED}{NIR + RED}$} : \underline{R}aster>Calculatrice raster> entrez la formule 
	\begin{equation}
		\frac{(VIS@4 - VIS@1)}{(VIS@4 + VIS@1)}
	\end{equation}
	puis double clic sur la couche et choisir dans symbologie type de rendu : Pseudo-couleur à bande unique> Palette de couleur: NDVI. Le résultat est montré sur la figure \ref{fig:step6_mapping_NDVI}
	\item Couche NIR. Sélectionner dans symbologie comme première bande NIR,  puis dans symbologie type de rendu :  Niveau de gris à bande unique.
\end{itemize}

\figurePaireIRD{fig:step6_mapping_1}
{step5_topo_QGIS.png}{Représentation fine de la topographie. Équidistance des isolignes : $10\,m$.}
{step5_mapping_VIS.png}{Mapping en spectre visible.}
{Représentation des données IRdrone avec le logiciel QGIS sur fond OpenStreetMap.}

\figurePaireIRD{fig:step6_mapping_NDVI}
{step5_mapping_VIR.png}{Image multispectrale VIR.  $\qquad\qquad\qquad\quad$ Décalage spectral $[Nir|R|G] \longrightarrow [R|G|B]$.}
{step5_mapping_NDVI.png}{Carte de l'indice NDVI.}
{Représentation des données IRdrone avec le logiciel QGIS sur fond Google Satellite.}


\section{Bouton \bouton{Close}}\label{btn:step6_close}
Ferme la fenêtre et rend la main au dialogue principal. La fermeture est refusée tant qu'un calcul ODM ou MNT est en cours. Le bouton \bouton{?} ouvre cette aide à la page de ce chapitre.

\section{Bouton \bouton{Continue with archiving}}\label{btn:step6_continue}
Apparaît en bas à droite, après \bouton{Close}, une fois la carte construite. Il ferme la fenêtre et ouvre l'étape~7, qui indique l'espace libéré par l'archivage. Rien n'est supprimé avant votre confirmation.

%==========================================================
\chapter{Step 7 -- Archivage de la mission}\label{ch:step7}
L'étape~7 est l'archivage de la mission. Elle se fait \emph{une fois l'exploitation des images terminée}. Elle nettoie le dossier de la mission \emph{sur place} :
les nombreux fichiers techniques (images intermédiaires, dossiers de travail, copies d'images de la synchronisation, etc.) sont supprimés,
et seuls les éléments essentiels pour pouvoir, si nécessaire, traiter de nouveau la mission sont conservés. Sur les missions testées, le volume est réduit
d'un facteur 2 à 5 (par exemple de 131\,Go à 25\,Go pour une mission de 600 images VIS DJI).

\begin{center}
\fbox{\parbox{0.9\linewidth}{\small\textbf{Attention : opération irréversible.} Les fichiers supprimés ne vont pas à la corbeille de Windows ; ils sont définitivement perdus.
Si vous souhaitez conserver une copie de la mission, d'un dossier \texttt{mapping\_MULTI} particulier ou de tout autre fichier, \emph{copiez-les ailleurs avant de lancer l'archivage}.
Tout élément du dossier de mission qui ne figure pas dans la liste ci-dessous est supprimé, y compris les copies de dossiers (par exemple \texttt{mapping\_MULTI - Copie}).}}
\end{center}

Comme pour les autres étapes, la mission archivée est la mission ouverte dans la fenêtre principale.

\section{Ce qui est conservé}\label{btn:step7_rules}
La comparaison des noms de fichiers ne tient pas compte des majuscules et minuscules (\texttt{.dng} et \texttt{.DNG} sont équivalents).
\begin{description}
  \item[Racine de la mission] les dossiers \texttt{AerialPhotography}, \texttt{cameras}, \texttt{FlightAnalytics}, \texttt{ImgIRdrone}, \texttt{mapping\_MULTI} (et les cartes précédentes \texttt{mapping\_MULTI\_<date>} mises de côté par l'étape~6), \texttt{Synchro}
        et le fichier \texttt{desktop.ini}. Tout le reste est supprimé.
  \item[\texttt{AerialPhotography}] les dossiers \texttt{VIS} et \texttt{NIR} et le fichier \texttt{shoot\_points.json}. Dans \texttt{VIS} et \texttt{NIR} : les images originales
        \texttt{.dng} avec leurs fichiers compagnons \texttt{.exif}, ainsi que \texttt{time\_line.json} et \texttt{transfer\_info\_VIS\_dng.json} (respectivement \texttt{transfer\_info\_NIR\_dng.json}).
        Les autres fichiers (notamment les \texttt{*\_RawTherapee.tif}) sont supprimés.
  \item[\texttt{cameras} et \texttt{FlightAnalytics}] laissés tels quels.
  \item[\texttt{ImgIRdrone}] les images \texttt{VIS\_XXXX.jpg} et \texttt{VIR\_XXXX\_local.jpg}, les fichiers \texttt{VIS\_XXXX\_motion\_model.npy}, les images
        \texttt{NDVI\_XXXX\_local.jpg} et \texttt{NIR\_XXXX\_local.jpg} (si elles existent), \texttt{shoot\_points.json} et \texttt{shoot\_points\_failed\_summary.json}.
        Le dossier volumineux \texttt{imgMultispectral} et les dossiers de \og redo \fg{} individuels sont supprimés.
  \item[\texttt{mapping\_MULTI}] laissé tel quel : il contient les éléments directement exploitables dans QGIS et a déjà été allégé à l'étape~6.
  \item[\texttt{Synchro}] seuls \texttt{check\_time\_alignment.png}, \texttt{check\_time\_before\_alignment.png} et \texttt{desktop.ini} sont conservés ; les dossiers \texttt{VIS}, \texttt{NIR} et \texttt{Check\_ARUCO} sont supprimés.
\end{description}

\paragraph{Pourquoi ce choix ?} Le décalage temporel trouvé à l'étape~3 et la liste des paires d'images (\texttt{AerialPhotography\textbackslash shoot\_points.json}) résument la synchronisation :
les images de \texttt{Synchro} ne sont plus utiles. Les fichiers \texttt{VIS\_XXXX\_motion\_model.npy} mémorisent l'alignement VIS--NIR de chaque \emph{shoot point}, y compris les alignements
rattrapés par la procédure de \og redo \fg{} (étape~4). Pour reprendre la mission, il suffit donc de refaire l'étape~4, sans perdre ces corrections, à partir des images originales conservées.
Les images \texttt{.jpg} conservées permettent aussi de retrouver rapidement une image pour illustrer un rapport.

\figureIRD[fig:step7]{step6_window.png}{Fenêtre de l'archivage : l'espace libéré par l'archivage s'affiche dès l'ouverture de la fenêtre.}{0.6}

\section{Ce que l'archivage libère (simulation)}\label{btn:step7_simulate}
Dès l'ouverture de la fenêtre, le dossier est analysé \emph{sans rien supprimer} (une simulation). Le cadre sous l'avertissement donne la réponse :
\begin{itemize}
  \item \texttt{12.3 GB can be freed}, avec, en grand, la taille du dossier de mission avant et après l'archivage (par exemple \texttt{15.1 GB → 2.8 GB}) et le nombre de fichiers qui seront supprimés. Le détail, dossier par dossier (fichiers à supprimer, fichiers conservés, dossiers supprimés), est dans la zone en dessous ;
  \item \texttt{Archiving impossible}, avec les fichiers indispensables qui manquent (voir plus bas) : rien ne peut être supprimé ;
  \item \texttt{Mission already archived}, avec la date de l'archivage et, en grand, la taille du dossier ;
  \item \texttt{Nothing to delete}, quand le dossier ne contient que ce que l'archivage conserve.
\end{itemize}
Vérifiez que la liste correspond à ce que vous attendez avant d'archiver.

\section{Bouton \bouton{Archive and free \ldots}}\label{btn:step7_archive}
Le bouton indique l'espace qui sera libéré (par exemple \bouton{Archive and free 12.3 GB\ldots}). Il est grisé pendant l'analyse du dossier, et quand rien ne peut être supprimé. Le dossier est analysé à nouveau au clic, puis un message de confirmation est demandé à l'utilisateur  comme le montre la figure \ref{fig:step7_confirm_archving}. Le message indique aussi l'espace disque qui sera libéré et précise que la suppression sera définitive. Deux boutons \texttt{OK} et \texttt{Cancel} permettent de confirmer ou d'abandonner. Par sécurité et pour éviter un clic trop rapide sur \texttt{Ok} l'utilisateur  doit recopier exactement le même nom du dossier de mission que celui indiqué (en respectant la casse). 

\figureIRD[fig:step7_confirm_archving]{step6_confirm_archiving.png}{Message de confirmation de la demande d'archivage.}{0.3}

 Si clic OK  alors IRDrone lance l'archivage réel :

\begin{enumerate}
  \item analyse le dossier et vérifie que les fichiers indispensables existent (\texttt{mission\_parameters.json}, \texttt{shoot\_points.json}, \texttt{time\_line.json}, \texttt{transfer\_info\_*\_dng.json}, etc.).
        S'il en manque, l'archivage est \textbf{refusé} et la liste des fichiers manquants est affichée ; rien n'est supprimé. Les dossiers manquants ne sont jamais recréés ;
  \item affiche le nombre de fichiers et le volume qui seront supprimés, puis demande confirmation : il faut \emph{saisir le nom du dossier de la mission} ;
  \item supprime les fichiers. La barre de progression indique l'avancement ;
  \item si un fichier est verrouillé (ouvert dans un autre logiciel, par exemple QGIS), l'erreur est notée et la suppression continue. Fermez alors les logiciels concernés et relancez l'archivage ;
  \item écrit à la racine de la mission les fichiers \texttt{archive\_manifest.json} et \texttt{archive\_log.txt}.
\end{enumerate}

Une mission déjà archivée est détectée et n'est pas archivée une seconde fois.


 Une icône particulière attribuée au dossiers archivés et permet de voir d'un coup d'\oe il qu'une mission est archivée.(figure \ref{fig:Air_Mission_folder}).


\paragraph{Archivage et avancement du traitement.} A la fin de chaque étape le dossier de mission reflète l'état de l'avancement du traitement de la mission. 

L'archivage ne doit être lancé qu'à la toute fin de l'étude. Notez que les images originales de la mission (converties au format dng pour la SJCam) sont systématiquement archivée. Vous conservez ainsi les données du terrain. Il n'est donc pas nécessaire d'archiver les images originales des cartes SD du DJI et de la SJCam (mais vous pouvez aussi les archiver par sécurité).

\section{Bouton \bouton{Close}}\label{btn:step7_close}
Ferme la fenêtre et rend la main au dialogue principal. Le bouton \bouton{?} ouvre cette aide à la page de ce chapitre.

%==========================================================
\chapter{Missions visible seul}\label{ch:visible_only}
\section{Sans la caméra infrarouge}\label{sec:visible_only}
Une mission peut être volée avec la seule caméra du drone, sans la SJCam : le type \bouton{Visible only (drone DNG)} se choisit en haut de la première page de l'étape~1 (section~\ref{btn:step1_mission_type}). Il est enregistré dans la mission (\texttt{"mission type": "visible"} dans \texttt{FlightAnalytics/mission\_parameters.json} ; les missions sans cette clé sont multispectrales). Les étapes sont les mêmes, moins ce qui concerne l'infrarouge :
\begin{itemize}
	\item \textbf{Étape 1} : seule la carte SD du drone est nécessaire ; la copie a trois étapes (copie, EXIF, altitude du sol). Le dossier de la mission n'a pas de dossier \texttt{AerialPhotography/NIR} (\texttt{Synchro} ne contient que les vignettes de la ligne de temps de l'étape~2).
	\item \textbf{Étape 2} : la ligne de temps ne montre que les images visibles ; il n'y a pas de cible tournante à chercher ni de phase de synchronisation : seule la phase de relevé est vérifiée. Le best mapping fonctionne comme pour une mission multispectrale. \bouton{Continue with development} enregistre le relevé et les points de prise de vue visibles (\texttt{AerialPhotography/shoot\_points.json}, un point par image visible, sans infrarouge) et ouvre l'étape~4.
	\item \textbf{Étape 3} : inutile ; sa tuile indique \og Not needed \fg~et \bouton{Continue} passe directement de l'étape~2 à l'étape~4.
	\item \textbf{Étape 4} (\og Develop \fg) développe les images DNG visibles du relevé (ou de la sélection du best mapping) : couleurs, correction de la distorsion de l'objectif et même recadrage que les images multispectrales, pour que le modèle de caméra d'OpenDroneMap soit le même. Chaque image donne \texttt{ImgIRdrone/VIS\_XXXX.jpg} (revue) et les trois bandes \texttt{ImgIRdrone/imgMultispectral/IRD\_XXXX\_1..3.tif} (rouge, vert, bleu) pour la carte.
	\item \textbf{Étape 5} : la revue montre les images visibles développées ; il n'y a rien à réaligner. \bouton{Good: keep for the map} (touche \texttt{Y}) marque l'image vérifiée et montre la suivante à vérifier ; une image floue ou mal exposée est signalée (touche \texttt{F}) : elle est écartée de la carte. Les points que l'étape~4 n'a pas développés (gris sur la carte) se développent d'un clic droit sur la carte.
	\item \textbf{Étape 6} : OpenDroneMap construit une orthophoto à trois bandes (rouge, vert, bleu) au lieu de quatre. Il n'y a ni couche en fausses couleurs infrarouges, ni VIR, ni NDVI ; le modèle de terrain est calculé comme d'habitude.
	\item \textbf{Étape 7} : l'archive garde les mêmes fichiers, sans ceux de l'infrarouge.
\end{itemize}

%==========================================================
\appendix
\chapter{Alignement des images VIS et NIR (Step 4)}\label{ann:alignement}\label{sec:ann_algorithme}
Cette annexe explique en détail comment IRdrone superpose l'image proche infrarouge à l'image visible. Elle s'adresse à un lecteur de formation scientifique (algèbre linéaire, un peu de traitement du signal) qui n'est pas spécialiste de vision par ordinateur. Les premières sections donnent les bases de géométrie et de traitement du signal, les suivantes décrivent chaque étape de l'algorithme avec les traces de diagnostic d'alignements réels, puis viennent les méthodes essayées puis abandonnées, l'historique du développement et les limites connues. La dernière section (§\ref{sec:ann_theorique}) traite de l'origine physique du désalignement des deux caméras (angles de \og pitch \fg{} et de \og yaw \fg{} théoriques).

Le code se trouve dans \texttt{automatic\_registration.py} (fonctions \texttt{align\_raw} et \texttt{coarse\_alignment}), dans le dossier \texttt{registration} (\texttt{rigid.py} : pyramide, \texttt{cost.py} : représentation et coût, \texttt{warp\_flow.py} : rééchantillonnage) et dans \texttt{irdrone/semi\_auto\_registration.py} (rotation de la caméra, \texttt{manual\_warp}). Les figures de cette annexe sont produites par \texttt{help/make\_alignment\_figures.py} à partir de deux shoot points :
\begin{itemize}
	\item \textbf{point A} : l'image \texttt{VIS\_0201} du vol du 25 janvier 2022 (Blassac, 90\,m au-dessus du sol, écart de temps entre les deux prises de vue 0,0\,s), un cas typique où l'alignement réussit ;
	\item \textbf{point B} : la paire d'exemple du dépôt (\texttt{samples/\_visible.DNG}, \texttt{samples/\_nir.RAW}, mai 2022), un cas difficile où l'alignement automatique échoue (§\ref{sec:ann_limits}).
\end{itemize}
Les deux ont été alignés en ligne de commande avec l'option \texttt{--debug} (§\ref{sec:ann_traces}), c'est-à-dire avec l'élargissement de champ 1,6 de la ligne de commande au lieu de 1,4 au Step~4.

\section{Le problème}\label{sec:ann_problem}
\subsection{Ce que signifie \og aligner \fg}
Le drone emporte deux caméras (figure~\ref{fig:ann_pair}) : la caméra du DJI Mavic Air~2, qui enregistre la lumière visible (VIS, rouge, vert et bleu), et une SJCam M20 dont le filtre anti-infrarouge a été retiré et remplacé par un filtre qui ne laisse passer que le proche infrarouge (NIR, au-delà d'environ 720\,nm). Les indices de végétation comme le NDVI combinent, \emph{pixel par pixel}, le rouge de l'image visible et le proche infrarouge :
\begin{equation}
	\text{NDVI} = \frac{\text{NIR} - \text{R}}{\text{NIR} + \text{R}}.
\end{equation}
Cette formule n'a de sens que si les deux valeurs d'un pixel proviennent du même point du sol. Aligner (on dit aussi \emph{recaler}) les deux images, c'est trouver, pour chaque pixel de l'image visible, la position du même point du sol dans l'image infrarouge, puis rééchantillonner l'image infrarouge sur la grille de l'image visible. Une erreur de quelques pixels suffit à créer de fausses valeurs d'indice le long de chaque bord contrasté (une haie, une route, l'ombre d'un arbre).

\figureIRD[fig:ann_pair]{align_pair_raw.jpg}{Le point B tel que l'enregistrent les deux caméras. Outre des champs de vision différents et la forte distorsion de la caméra infrarouge, les contrastes diffèrent : la route est sombre dans les deux bandes mais ses bas-côtés ne sont clairs que dans le visible, la végétation est verte et plus sombre que l'herbe sèche dans le visible, claire dans le proche infrarouge.}{1}

\subsection{Pourquoi c'est difficile}
Trois difficultés s'additionnent :
\begin{enumerate}
	\item \textbf{Les deux caméras ne regardent pas dans la même direction.} La caméra visible est maintenue verticale par sa nacelle (gimbal), la caméra infrarouge est fixée au corps du drone et suit ses inclinaisons (§\ref{sec:ann_theorique}). L'écart atteint plusieurs degrés, c'est-à-dire plusieurs centaines de pixels.
	\item \textbf{Les deux optiques sont très différentes.} La caméra infrarouge a un champ très large et une forte distorsion ; les lignes droites apparaissent courbes (§\ref{sec:ann_distortion}).
	\item \textbf{Les deux bandes ne voient pas la même chose.} La végétation, l'eau, le sol et les routes n'ont pas la même réflectance dans le visible et dans le proche infrarouge. Un contraste peut s'inverser d'une bande à l'autre (figure~\ref{fig:ann_pair}, et figure~\ref{fig:step4_pair_VIS_NIR} où la rivière est noire dans l'infrarouge). Les méthodes classiques, qui comparent des intensités ou apparient des points caractéristiques (SIFT, ORB), sont donc peu fiables.
\end{enumerate}
Les deux premières difficultés sont résolues par la \emph{géométrie} : un modèle de chaque caméra, obtenu par calibration, et le fait que les deux images diffèrent essentiellement d'une rotation. La troisième est résolue en \emph{comparant des bords plutôt que des intensités}.

\section{Géométrie d'une caméra}\label{sec:ann_camera}
Le livre de référence est \emph{Multiple View Geometry in Computer Vision} de Hartley et Zisserman~[\hyperref[ann_hz]{1}] (chapitres 2 et 6 pour le modèle de caméra, 7.4 pour la distorsion, 8.4 pour les caméras en rotation). Le livre de Szeliski~[\hyperref[ann_szeliski]{2}] est une introduction plus accessible (chapitre 2).

\subsection{Le modèle sténopé et les coordonnées homogènes}
En première approximation, une caméra est un \emph{sténopé} (pinhole) : tout rayon lumineux passe par le centre optique $C$ avant d'atteindre le capteur. Soient $(X,Y,Z)$ les coordonnées d'un point de la scène dans le repère de la caméra ($Z$ selon l'axe optique). Son image est au pixel
\begin{equation}
	u = f\,\frac{X}{Z} + c_x, \qquad v = f\,\frac{Y}{Z} + c_y,
\end{equation}
où $f$ est la focale \emph{exprimée en pixels} (focale physique divisée par la taille d'un pixel) et $(c_x, c_y)$ le \emph{point principal}, intersection de l'axe optique avec le capteur, proche du centre de l'image. La division par $Z$ est la perspective : un objet deux fois plus loin apparaît deux fois plus petit.

Cette division n'est pas linéaire ; elle le devient en écrivant le pixel en \emph{coordonnées homogènes}, un vecteur de trois nombres défini à un facteur près : $(u, v, 1) \sim (\lambda u, \lambda v, \lambda)$. Alors
\begin{equation}\label{eq:ann_pinhole}
	\lambda \begin{pmatrix} u \\ v \\ 1 \end{pmatrix} = K \begin{pmatrix} X \\ Y \\ Z \end{pmatrix},
	\qquad
	K = \begin{pmatrix} f & 0 & c_x \\ 0 & f & c_y \\ 0 & 0 & 1 \end{pmatrix}.
\end{equation}
La matrice $K$ des \emph{paramètres intrinsèques} (ou matrice de calibration) décrit la caméra seule. Inversement, $K^{-1}(u, v, 1)^T$ est la \emph{direction} du rayon qui arrive sur le pixel $(u,v)$ : une image est un ensemble de directions, la profondeur $Z$ est perdue.

\begin{table}[htbp]\centering
\begin{tabular}{@{}lll@{}}
\hline
 & \textbf{DJI Mavic Air~2 (VIS)} & \textbf{SJCam M20 (NIR)} \\ \hline
Taille de l'image (pixels) & $3992 \times 2992$ & $4600 \times 3448$ \\
Focale $f$ (pixels) & 2898 & 2311 \\
Point principal $(c_x, c_y)$ & (2032, 1511) & (2348, 1725) \\
Champ horizontal & $68^\circ$ & plus de $110^\circ$ \\
Distorsion $k_1, k_2, k_3$ & $0{,}185,\ -0{,}455,\ 0{,}245$ & $-0{,}337,\ 0{,}131,\ -0{,}025$ \\ \hline
\end{tabular}
\caption{Calibrations des deux caméras (\texttt{calibration/DJI\_RAW} et \texttt{calibration/M20\_RAW}, images RAW). À la focale de la caméra visible, une rotation de $1^\circ$ déplace l'image de $f \tan 1^\circ \approx 51$ pixels.}
\label{tab:ann_calib}
\end{table}

\subsection{La distorsion}\label{sec:ann_distortion}
Les objectifs réels s'écartent du modèle sténopé : l'image d'une droite est courbe, de plus en plus vers les bords. IRdrone utilise le modèle polynomial de Brown et Conrady~[\hyperref[ann_brown]{3}], celui d'OpenCV. Avec les coordonnées normalisées $x = X/Z$, $y = Y/Z$ et $r^2 = x^2 + y^2$, le point idéal $(x, y)$ est observé en
\begin{equation}\label{eq:ann_distortion}
	x_d = x\,(1 + k_1 r^2 + k_2 r^4 + k_3 r^6) + 2p_1 xy + p_2 (r^2 + 2x^2),
\end{equation}
et de même pour $y_d$ ; le pixel est alors $K\,(x_d, y_d, 1)^T$. Les coefficients $k_i$ décrivent la distorsion radiale (en barillet si $k_1 < 0$, en coussinet si $k_1 > 0$), $p_1, p_2$ un léger décentrement des lentilles.

La figure~\ref{fig:ann_distortion} donne les ordres de grandeur. Pour la caméra visible la correction ne dépasse pas 40 pixels. Pour la caméra infrarouge elle atteint 210 pixels aux bords gauche et droit et 375 pixels aux coins du champ de la caméra visible, et 800 pixels à $50^\circ$ de l'axe. Son $k_1$ fortement négatif est une distorsion en barillet proche de celle d'un \og fish-eye \fg. Le polynôme n'est valable que dans la zone où des mires de calibration ont été vues : au-delà d'environ $57^\circ$ de l'axe il redécroît, le modèle n'est donc plus inversible, et les bords de l'image infrarouge (plus de $110^\circ$ de champ horizontal) sont hors de son domaine de validité. Ces zones ne sont jamais utilisées, puisque le champ visible n'atteint que $41^\circ$ dans les coins.

\figureIRD[fig:ann_distortion]{align_distortion.png}{Distorsion des deux caméras. À gauche et au centre : un quadrillage de droites au sol, tel qu'il est enregistré sur le capteur (cadre noir) de chaque caméra ; en rouge, le champ de la caméra visible vu par la caméra infrarouge. À droite : décalage radial d'un point en fonction de son angle avec l'axe optique.}{1}

\subsection{La calibration}
Les paramètres intrinsèques $f, c_x, c_y$ et les coefficients de distorsion sont mesurés une fois pour chaque caméra par la méthode de Zhang~[\hyperref[ann_zhang]{4}] : une trentaine de photographies d'un damier plan dans diverses orientations ; les coins des cases sont détectés, et tous les paramètres (plus la position du damier sur chaque photo) sont estimés en minimisant la distance entre les coins observés et les coins prédits par les équations \eqref{eq:ann_pinhole} et \eqref{eq:ann_distortion} (\texttt{calibration/cameracalibration.py}, OpenCV). Cette calibration est fixe : elle ne dépend pas du vol.

\section{Deux caméras qui tournent : l'homographie}\label{sec:ann_homography}
\subsection{Une rotation de la caméra est une homographie de l'image}
Supposons que les deux caméras aient le \emph{même centre optique} et ne diffèrent que d'une rotation $R$ (une matrice orthogonale $3 \times 3$). Le rayon qui arrive sur le pixel $p_N$ de la caméra infrarouge a la direction $K_{NIR}^{-1} p_N$ dans le repère de cette caméra, $R\, K_{NIR}^{-1} p_N$ dans celui de la caméra visible, et arrive donc sur le pixel
\begin{equation}\label{eq:ann_rotation_homography}
	p_V \sim K_{VIS}\, R\, K_{NIR}^{-1}\, p_N = H\, p_N
\end{equation}
de la caméra visible (après correction des distorsions). La matrice $3 \times 3$ $H$, définie à un facteur près, est une \emph{homographie} : la transformation la plus générale qui envoie les droites sur des droites (Hartley et Zisserman~[\hyperref[ann_hz]{1}], §2.3 et §8.4). Elle a 8 degrés de liberté. Point remarquable : l'équation \eqref{eq:ann_rotation_homography} ne fait pas intervenir la profondeur $Z$ : quel que soit le relief, une rotation pure de la caméra est une homographie de l'image. C'est le principe des panoramas réalisés en tournant sur soi-même.

La rotation est paramétrée par trois angles : le \emph{lacet} (yaw, rotation autour de l'axe vertical de l'image, qui la déplace horizontalement), le \emph{tangage} (pitch, qui la déplace verticalement) et le \emph{roulis} (roll, rotation autour de l'axe optique, qui fait tourner l'image autour de son centre). Dans le code, $R$ est construite à partir du vecteur de rotation $(-\text{pitch}, \text{yaw}, \text{roll})$ (formule de Rodrigues, \texttt{manual\_warp}). Aux petits angles, un lacet $\theta$ déplace le centre de l'image de $f \tan\theta$ pixels, soit environ 51 pixels par degré ; un roulis $\varphi$ déplace un coin, à environ 2500 pixels du centre, de $2500\,\varphi$ (44 pixels par degré).

\subsection{Pourquoi la rotation suffit (presque) : la parallaxe}
Les deux centres optiques ne sont pas exactement au même endroit : ils sont à 46\,mm l'un de l'autre sur le drone (\texttt{CNIRCVIS\_0} dans \texttt{config.py}), et le drone se déplace entre les deux prises de vue. Une distance $B$ entre les centres (la \emph{base}) crée une \emph{parallaxe} : un point à la distance $Z$ est décalé de
\begin{equation}
	d = f\,\frac{B}{Z} \quad \text{pixels}.
\end{equation}
Seules les \emph{variations} de ce décalage comptent, sa partie constante étant absorbée par la rotation. Avec $B = 46$\,mm, à 90\,m : $d = 2898 \times 0{,}046 / 90 \approx 1{,}5$ pixel, et un relief de 10\,m le modifie de 0,2 pixel : négligeable. Si les deux prises de vue sont séparées de $\Delta t = 0{,}5$\,s et que le drone vole à 3\,m/s, $B = 1{,}5$\,m : $d \approx 48$ pixels, et un relief de 10\,m (une haie d'arbres, une pente) le modifie maintenant de $f B \Delta Z / Z^2 \approx 5$ pixels. C'est la raison pour laquelle le Step~4 n'aligne que les shoot points dont l'écart de temps est faible, et pour laquelle une correction \emph{locale} s'ajoute au modèle global (§\ref{sec:ann_local}). À très basse altitude (décollage, atterrissage), la parallaxe domine et l'alignement n'a pas de sens.

\subsection{La caméra virtuelle}
L'image visible est d'abord corrigée de sa distorsion : elle devient l'image d'une caméra sténopé idéale, la \emph{référence}. L'image infrarouge est ensuite rééchantillonnée directement dans cette caméra de référence, par l'équation \eqref{eq:ann_rotation_homography} et son propre modèle de distorsion, mais sur un canevas plus large que l'image visible (1,4 fois par défaut, 1,6 dans le script REDO) : c'est la \emph{caméra virtuelle} de la figure~\ref{fig:ann_virtual}. Si la rotation était l'identité, l'image visible correspondrait exactement au rectangle central ; chercher la rotation revient à trouver où elle se trouve réellement.

\figureIRD[fig:ann_virtual]{align_virtual_camera.jpg}{Point B. À gauche : l'image infrarouge corrigée de sa distorsion et rééchantillonnée dans la caméra virtuelle (focale de la caméra visible, champ 1,4 fois plus large). En pointillés : position de l'image visible si les deux caméras étaient parallèles ; en rouge : position trouvée par la recherche grossière (lacet $8{,}2^\circ$, tangage $1{,}9^\circ$). À droite : la référence visible. En dehors du champ visible la correction de la distorsion est moins précise (la route s'incurve à gauche).}{1}

\section{Comparer deux bandes spectrales}\label{sec:ann_similarity}
Pour trouver la rotation, il faut un \emph{critère de similarité} : un nombre qui dit si deux images sont bien superposées, et qui est optimal pour la bonne superposition.

\subsection{Pourquoi on ne peut pas comparer les intensités}
Pour deux images d'une même bande, le critère le plus simple est la somme des carrés des différences, $\text{SSD} = \sum_p (I_1(p) - I_2(p))^2$, ou la corrélation. Tous deux supposent qu'un même objet a la même intensité (ou une intensité proportionnelle) dans les deux images. Entre le visible et le proche infrarouge c'est faux : la végétation est sombre dans le visible et claire dans l'infrarouge, les bas-côtés de la route ne sont clairs que dans le visible (figures~\ref{fig:ann_pair} et \ref{fig:ann_features}). Une réponse classique à ces problèmes \og multimodaux \fg{} est l'information mutuelle~[\hyperref[ann_viola]{5}], qui mesure une dépendance statistique au lieu d'une égalité. Une autre approche combine des points caractéristiques et une similarité de gradients~[\hyperref[ann_vayssade]{18}]. IRdrone utilise une idée plus simple et plus locale : les \emph{bords} sont communs même quand les contrastes ne le sont pas.

\subsection{Les cartes de caractéristiques : énergies directionnelles du gradient}\label{sec:ann_msr}
Chaque image est convertie en luminance (niveaux de gris), lissée par un filtre gaussien d'écart-type $\sigma$ (5 pixels pour le visible, 3 pour l'infrarouge, moins net ; 3 et 1 pour la recherche grossière), puis remplacée par quatre cartes (\texttt{multispectral\_representation}, mode \texttt{LAPLACIAN\_ENERGIES}) :
\begin{equation}
	E_x = (\partial_x I)^2, \quad E_y = (\partial_y I)^2, \quad E_{\nearrow} = (\partial_{\nearrow} I)^2, \quad E_{\searrow} = (\partial_{\searrow} I)^2.
\end{equation}
Les dérivées horizontale et verticale sont calculées avec des filtres de Scharr~[\hyperref[ann_scharr]{6}] (une dérivée $3 \times 3$ avec un lissage optimisé, plus isotrope que Sobel), les diagonales avec la croix de Roberts $2 \times 2$~[\hyperref[ann_roberts]{7}]. Ces cartes sont appelées \emph{représentation multispectrale} (MSR) dans le code. L'idée vient d'Irani et Anandan~[\hyperref[ann_irani]{9}], qui alignaient des images de capteurs différents (visible et infrarouge entre autres) au moyen de telles énergies de dérivées directionnelles.

Le carré est essentiel : un bord du sombre vers le clair et un bord du clair vers le sombre donnent la même énergie. Une inversion de contraste entre les bandes devient donc invisible, tandis que la \emph{position} du bord est conservée. Les quatre directions gardent l'orientation des bords, ce qui aide à distinguer, par exemple, une route du bord perpendiculaire d'un champ. Le lissage supprime les textures fines (herbe, feuilles), différentes dans les deux bandes, et garde les grandes structures. La figure~\ref{fig:ann_features} montre ces cartes sur un détail du point B : les luminances sont très différentes, les énergies des bords de la route sont au même endroit.

\figureIRD[fig:ann_features]{align_feature_maps.jpg}{Énergies directionnelles du gradient sur un détail du point B ($600 \times 600$ pixels). Les luminances des deux bandes ne sont pas comparables (bas-côtés clairs dans le visible, sombres dans l'infrarouge), mais les bords de la route produisent les mêmes lignes dans les cartes d'énergie. Certains bords n'existent que dans une bande (les quatre lignes du visible contre deux dans l'infrarouge) : le critère doit les tolérer.}{1}

\subsection{Le critère de similarité : gradient total de la différence}\label{sec:ann_ntg}
Le critère dérive du \emph{Normalized Total Gradient} (NTG) de Chen et al.~[\hyperref[ann_ntg]{8}]. Pour deux images $f$ et $g$ :
\begin{equation}
	\text{NTG}(f, g) = \frac{\sum_p \|\nabla (f - g)(p)\|_1}{\sum_p \|\nabla f(p)\|_1 + \sum_p \|\nabla g(p)\|_1}.
\end{equation}
L'idée : quand deux images sont bien superposées, leurs bords communs s'annulent dans la différence $f-g$, qui devient lisse ; son gradient total (somme des valeurs absolues des dérivées) est alors minimal. Quand elles sont décalées, chaque bord apparaît deux fois dans la différence, avec des signes opposés, et le gradient total augmente. La norme $\|\cdot\|_1$ (valeurs absolues plutôt que carrés) rend le critère robuste aux bords qui n'existent que dans une image : ils ajoutent un coût à peu près constant, qui dépend peu du décalage, au lieu de le dominer. Le dénominateur rend le critère indépendant de la quantité de texture.

IRdrone l'applique aux cartes d'énergie : pour un décalage $u$ de la référence,
\begin{equation}\label{eq:ann_cost}
	C(u) = \sum_{c=1}^{4} \sum_p \Big( \big|S_x * (E^{NIR}_c - E^{VIS}_c(\cdot + u))\big|(p) + \big|S_y * (E^{NIR}_c - E^{VIS}_c(\cdot + u))\big|(p) \Big),
\end{equation}
où $S_x, S_y$ sont les filtres de dérivation de Sobel (\texttt{cost\_surface\_NTG} dans \texttt{cost.py}). Seul le numérateur du NTG est calculé. Le filtre de Sobel étant linéaire, $S * (a - b) = S*a - S*b$ : les gradients des deux images sont calculés une fois, et chaque décalage n'est plus qu'une somme de différences absolues (15 fois plus rapide, novembre 2026). Une somme des carrés des différences sur les mêmes cartes est aussi programmée (\texttt{SSD}) mais n'est pas utilisée.

La figure~\ref{fig:ann_profiles} compare les critères sur le point A, en fonction d'un décalage horizontal autour du bon alignement. La SSD des luminances a une vallée large : elle situe l'alignement, mais avec peu de précision, et elle suppose des intensités comparables, ce qui est le cas sur cette scène d'hiver mais pas en été. Le NTG a un puits étroit et profond : il est très précis, mais ne \og voit \fg{} la bonne position qu'à quelques pixels près, d'où le besoin d'un bon point de départ et d'une stratégie du grossier au fin.

\figureIRD[fig:ann_profiles]{align_cost_profiles.png}{Coût en fonction d'un décalage horizontal (point A, après alignement, calculé à la résolution 1/4), normalisé entre 0 et 1. (a) L'image visible contre elle-même. (b) Visible contre infrarouge sur un bloc avec des arbres et une haie. (c) Sur un bloc de prairie seule, avec peu de structures : le NTG a encore son minimum au bon endroit mais en dehors du puits étroit la courbe n'est plus informative.}{1}

\section{Recherche grossière}\label{sec:ann_coarse}
La recherche grossière trouve l'essentiel de l'écart d'orientation (fonction \texttt{coarse\_alignment}) :
\begin{enumerate}
	\item l'image infrarouge est rééchantillonnée dans la caméra virtuelle avec les angles initiaux : nuls au Step~4, ceux réglés à la main lors d'un réalignement (§\ref{btn:review_redo}) ;
	\item les cartes d'énergie des deux images sont réduites d'un facteur 32 ($125 \times 94$ pixels pour le visible), la carte visible est placée au centre d'un canevas de la taille de la carte infrarouge ;
	\item une \emph{recherche exhaustive} calcule le coût \eqref{eq:ann_cost} pour chaque translation de la carte visible à l'intérieur de la carte infrarouge : avec l'élargissement de champ 1,4 du Step~4, $\pm 25$ pixels à cette échelle, soit environ $\pm 800$ pixels en pleine résolution ou $\pm 15^\circ$ ; avec l'élargissement 1,6 de la ligne de commande, $\pm 37$ pixels, $\pm 1200$ pixels ou $\pm 22^\circ$ (figure~\ref{fig:ann_coarse}) ;
	\item la meilleure position est localisée avec une précision inférieure au pixel (ci-dessous) ;
	\item aux petits angles une translation horizontale $t_x$ est un lacet et une translation verticale $t_y$ un tangage : $\text{yaw} = \arctan(t_x/f)$, $\text{pitch} = \arctan(t_y/f)$ ;
	\item l'image infrarouge est rééchantillonnée une nouvelle fois, en pleine résolution, avec les angles corrigés.
\end{enumerate}
Le roulis n'est pas recherché : il garde sa valeur initiale (§\ref{sec:ann_fourier}).

\paragraph{Choix du minimum.} Le minimum absolu de la surface de coût n'est pas toujours le bon : une zone sans texture ou le bord noir du canevas peuvent produire une vallée large et basse. La bonne superposition produit un puits étroit et profond (figure~\ref{fig:ann_coarse}). Le code cherche donc la position où la surface est la plus fortement courbée : en chaque position, une surface quadratique est ajustée sur le voisinage $3\times3$, et la position du plus grand déterminant de sa matrice hessienne (produit des deux courbures principales) est retenue (\texttt{minimum\_cost\_max\_hessian}).

\paragraph{Raffinement sous-pixel.} Autour de la position retenue, le coût est approché par une forme quadratique~[\hyperref[ann_hasinoff]{10}]
\begin{equation}
	C(u) \approx \tfrac{1}{2}\, u^T A\, u + b^T u + c,
\end{equation}
dont les coefficients s'obtiennent par de petites convolutions du voisinage $3\times3$. Son minimum est en $u^* = -A^{-1} b$ (un pas de Newton), retenu s'il est à moins d'un pixel. À la résolution 1/32, un pixel vaut 32 pixels en pleine résolution : ce raffinement rend la recherche grossière précise à quelques pixels près. Pour le point A, la translation trouvée est de $(21{,}7 ; 163{,}9)$ pixels, soit un lacet de $0{,}43^\circ$ et un tangage de $3{,}24^\circ$.

\figureIRD[fig:ann_coarse]{align_coarse_search.png}{Recherche grossière du point A à la résolution 1/32. À gauche : la surface de coût sur toute la plage de recherche ($\pm 1200$ pixels) : un puits net dans un paysage lisse. Au centre : zoom sur le centre du coût. À droite : déterminant de la hessienne, maximal au fond du puits.}{1}

\section{Raffinement pyramidal}\label{sec:ann_pyramid}
Après la recherche grossière, le désalignement résiduel va de quelques pixels à quelques dizaines de pixels, et il n'est pas uniforme : roulis résiduel, erreurs des calibrations, parallaxe. Le raffinement l'estime par blocs, sur une pyramide d'images (\texttt{pyramidal\_search} dans \texttt{rigid.py}).

\subsection{Du grossier au fin}
Une \emph{pyramide gaussienne}~[\hyperref[ann_burt]{11}] est la suite d'une image lissée et réduite par 2, 4, 8, 16\ldots{} Un déplacement de 100 pixels en pleine résolution fait 6 pixels à 1/16 : on le trouve avec une petite fenêtre de recherche, vite et avec peu de risque de confusion. Une fois trouvé, il est appliqué, et le niveau suivant, plus fin, n'a plus qu'à corriger le reste. C'est la stratégie classique des méthodes d'alignement de Lucas--Kanade~[\hyperref[ann_lucas]{12}, \hyperref[ann_baker]{13}]. IRdrone utilise le schéma suivant (\texttt{iterative\_scheme} dans \texttt{align\_raw}) :

\begin{table}[htbp]\centering
\begin{tabular}{@{}cccccc@{}}
\hline
\textbf{Échelle} & \textbf{Itérations} & \textbf{Blocs} & \textbf{Taille d'un bloc (pleine rés.)} & \textbf{Fenêtre du coût} & \textbf{Plus grand vecteur (pleine rés.)} \\ \hline
1/16 & 2 & $4 \times 4$ & $998 \times 748$ & $\pm 8$ px & $\pm 80$ px \\
1/16 & 2 & $5 \times 5$ & $798 \times 598$ & $\pm 6$ px & $\pm 80$ px \\
1/4 & 3 & $5 \times 5$ & $798 \times 598$ & $\pm 6$ px & $\pm 20$ px \\ \hline
\end{tabular}
\caption{Schéma du raffinement pyramidal. Le coût est calculé sur la fenêtre indiquée, à l'échelle du niveau, mais le minimum est cherché à $\pm 5$ pixels du centre (\texttt{brute\_force\_vector\_field\_search}) : un vecteur de bloc ne peut pas dépasser 80 pixels à 1/16, 20 pixels à 1/4. Les itérations sont numérotées 1, 2, 4, 5, 7, 8, 9 dans les noms des fichiers de traces (le compteur est incrémenté une fois de plus à chaque changement de niveau).}
\label{tab:ann_pyramid}
\end{table}

\subsection{Une itération}
\begin{enumerate}
	\item Les cartes d'énergie de l'image infrarouge sont rééchantillonnées avec le modèle courant (l'identité au départ).
	\item L'image est découpée en une grille de blocs. Pour chaque bloc, le coût \eqref{eq:ann_cost} est calculé pour chaque décalage de la fenêtre de recherche : on obtient une \emph{surface de coût} par bloc (figures~\ref{fig:ann_blocks_first} et \ref{fig:ann_blocks_last}).
	\item Le minimum de chaque surface (à $\pm 5$ pixels du centre), affiné au sous-pixel comme ci-dessus, donne un vecteur de déplacement par bloc : un \emph{champ de vecteurs} grossier, ou flot optique.
	\item Une homographie est ajustée sur ces vecteurs : les 16 ou 25 couples (centre du bloc, centre + vecteur) forment un système surdéterminé pour ses 8 inconnues. Il est résolu par la transformation linéaire directe (DLT) dans une boucle RANSAC~[\hyperref[ann_ransac]{14}], qui ignore les blocs dont le vecteur s'écarte de plus de 3 pixels de la majorité (\texttt{geometric\_rigid\_transform\_estimation}, \texttt{findHomography} d'OpenCV ; Hartley et Zisserman~[\hyperref[ann_hz]{1}], chapitre 4).
	\item Le modèle est mis à jour par composition, $H_k = \Delta H_k\, H_{k-1}$, puis l'itération suivante commence.
\end{enumerate}
Une homographie plutôt que trois angles : elle représente la rotation d'une caméra \emph{et} les erreurs résiduelles des calibrations (pour le point A, l'homographie finale contient un facteur d'échelle de 1,9\,\% horizontalement et 1,1\,\% verticalement, par exemple une petite erreur sur les focales).

\figureIRD[fig:ann_blocks_first]{align_blocks_first.jpg}{Première itération du raffinement du point A (résolution 1/16, blocs $4\times4$, recherche $\pm 8$ pixels). À gauche : superposition après la recherche grossière (visible en vert, infrarouge en magenta ; gris là où ils concordent), avec les vecteurs mesurés (bleu, agrandis 20 fois) et les vecteurs de l'homographie ajustée (rouge). À droite : la surface de coût de chaque bloc ; le point noir est le décalage nul, la croix le minimum. Les vecteurs atteignent 50 pixels et varient régulièrement : un facteur d'échelle et un tangage résiduels, que l'homographie représente bien.}{1}

\figureIRD[fig:ann_blocks_last]{align_blocks_last.jpg}{Dernière itération (résolution 1/4, blocs $5\times5$, recherche $\pm 6$ pixels), sur l'image alignée globalement. Tous les minima sont au centre : les vecteurs résiduels (agrandis 200 fois) sont inférieurs à 2 pixels, l'homographie ne change plus. Ces vecteurs résiduels forment la correction locale (§\ref{sec:ann_local}).}{1}

\figureIRD[fig:ann_convergence]{align_convergence.png}{Longueur des vecteurs des blocs à chaque itération, pour les deux points. Le point A (bleu) converge : de 23 pixels en moyenne à 0,6 pixel. Les marches aux changements de niveau sont dues au changement de taille des blocs et de résolution (un pixel à 1/16 vaut 16 pixels). Le point B (rouge) ne converge pas : les vecteurs restent à leur limite, 80 pixels à 1/16 et 20 pixels à 1/4 (§\ref{sec:ann_limits}).}{0.7}

\subsection{Alignement local}\label{sec:ann_local}
Une homographie est un modèle global : elle ne peut pas représenter la parallaxe du relief ni une imperfection locale du modèle de distorsion. Le champ de vecteurs de la dernière itération ($5 \times 5$ vecteurs) est donc conservé comme \emph{correction locale} : il est interpolé bilinéairement entre les centres des blocs (constant au-delà) en un déplacement pour chaque pixel, composé avec le modèle global (\texttt{warp\_from\_sparse\_vector\_field}). Deux résultats sont produits :
\begin{itemize}
	\item alignement \emph{global} : rotation et homographie finale ;
	\item alignement \emph{local} : rotation, homographie de l'itération précédente et champ de déplacement local. C'est en général le plus précis, et celui des images multispectrales \texttt{IRD\_XXXX\_n.tif}.
\end{itemize}

\subsection{Un seul rééchantillonnage}
Chaque rééchantillonnage d'une image interpole entre ses pixels, et la rend donc un peu floue. Au lieu de rééchantillonner l'image infrarouge à chaque étape, la transformation finale est composée en une seule carte : pour chaque pixel de l'image visible, la position à lire dans l'image infrarouge \emph{brute} s'obtient en enchaînant le déplacement local, l'homographie, la rotation de l'équation \eqref{eq:ann_rotation_homography} et la distorsion de la caméra infrarouge ; l'image n'est ensuite interpolée qu'une fois (interpolation bicubique, \texttt{cv2.remap}). Dans le code ce principe s'appelle \og You Only Warp Once \fg{} (\texttt{FULLRES\_REGISTERED\_REFINED\_WARPED\_ONCE\_ONLY.jpg}). Le modèle de mouvement (angles, homographies, champ de vecteurs) est enregistré dans \texttt{VIS\_XXXX\_motion\_model.npy} et sa version lisible \texttt{.json} ; s'il existe, relancer le Step~4 le relit au lieu de chercher à nouveau.

La figure~\ref{fig:ann_stages} montre la progression de l'alignement sur un détail du point A.

\figureIRD[fig:ann_stages]{align_stages.jpg}{Les quatre étapes sur un détail du point A (visible en vert, infrarouge en magenta, gris là où ils concordent). (a) Calibrations seules : les images sont décalées d'environ 160 pixels. (b) Après la recherche grossière : il reste une dizaine de pixels (franges colorées le long des ombres). (c) Avec l'homographie. (d) Avec la correction locale : les bords coïncident.}{1}

\section{Pistes explorées : corrélation de phase, Fourier--Mellin et appariement polaire}\label{sec:ann_fourier}
\subsection{Corrélation de phase}
Les premières versions (août--septembre 2021) estimaient la translation par \emph{corrélation de phase}~[\hyperref[ann_kuglin]{15}]. Si $g(x) = f(x - t)$, leurs transformées de Fourier vérifient $\hat g(\omega) = \hat f(\omega)\, e^{-i\omega \cdot t}$ (théorème du retard). L'interspectre normalisé
\begin{equation}
	\frac{\hat f\, \overline{\hat g}}{|\hat f\, \overline{\hat g}|} = e^{i\omega \cdot t}
\end{equation}
ne garde que la phase, et sa transformée de Fourier inverse est un pic de Dirac à la position $t$. La méthode est rapide (deux FFT), insensible à un changement global de luminosité, et trouve de grandes translations. Mais elle suppose que les deux images sont le \emph{même} signal décalé : une rotation, un changement d'échelle, une distorsion ou une inversion de contraste étalent le pic. La figure~\ref{fig:ann_phase} le montre : sur le point A, après la recherche grossière, le pic est net ; sur le point B il est noyé dans le bruit, sur les luminances comme sur les cartes d'énergie.

\figureIRD[fig:ann_phase]{align_phase_correlation.png}{Surfaces de corrélation de phase entre l'image visible et l'image infrarouge après la recherche grossière (résolution 1/8, zoom sur $\pm 320$ pixels). En haut, point A : un pic net, environ 2,5 fois plus haut que les valeurs suivantes. En bas, point B (rotation résiduelle, contrastes inversés, grandes zones texturées) : le pic émerge à peine.}{0.65}

\subsection{Fourier--Mellin : rotation et échelle}
La méthode de Fourier--Mellin~[\hyperref[ann_reddy]{16}] étend la corrélation de phase aux rotations et aux changements d'échelle. Le module de la transformée de Fourier ne dépend pas de la translation ; une rotation de l'image le fait tourner, un changement d'échelle le dilate. Rééchantillonné en coordonnées \emph{log-polaires} $(\log r, \theta)$, le module transforme une rotation en translation selon $\theta$ et un changement d'échelle en translation selon $\log r$, qu'une seconde corrélation de phase retrouve. Rabatel et Labbé~[\hyperref[ann_rabatel]{17}] l'ont appliquée à des images de drone visibles et proche infrarouge. Dans IRdrone, l'échelle est connue par les calibrations, et seul le roulis était cherché de cette façon : l'alignement semi-automatique de 2021 rééchantillonnait les deux images en coordonnées polaires et estimait le roulis par corrélation de phase selon l'angle (\texttt{mellin\_transform} et \texttt{estimate\_rotation} dans \texttt{semi\_auto\_registration.py}, toujours présentes pour une démonstration mais non utilisées par le Step~4). Combinée à la translation par corrélation de phase, ou à des points SIFT, elle n'a pas donné de résultats fiables sur les images multispectrales, pour les raisons ci-dessus. À partir de septembre 2021 elle a été remplacée par la minimisation du coût \eqref{eq:ann_cost} sur les cartes d'énergie.

\subsection{Appariement polaire des cartes d'énergie (essai de 2022)}
Le roulis est le point faible de l'algorithme actuel : la recherche grossière ne trouve que le lacet et le tangage, et le raffinement pyramidal ne corrige qu'un petit roulis. Lors d'un vol à Thaumiers (novembre 2022) la caméra infrarouge était montée avec un roulis constant d'environ $-2{,}4^\circ$, soit des déplacements de 100 pixels dans les coins, trop pour le raffinement (ticket GitHub \#33). Un essai a combiné les deux approches : les cartes d'énergie sont rééchantillonnées en coordonnées polaires autour du centre de l'image, où un roulis devient une translation verticale, et le coût \eqref{eq:ann_cost} est calculé pour chaque décalage de $1^\circ$ ; une parabole ajustée sur les trois valeurs les plus basses donne le roulis à mieux qu'un degré (\texttt{coarse\_roll\_estimation}, branche \texttt{local\_alignment\_improvements}, commit \texttt{86353cb}, non intégré). La figure~\ref{fig:ann_polar} le reproduit : sur le point A, tourné d'un roulis connu de $2{,}4^\circ$, l'estimation est de $2{,}26^\circ$. Sur le point B, la route est à environ $34^\circ$ dans l'image visible et à environ $28^\circ$ dans l'image infrarouge après la recherche grossière, et l'appariement polaire trouve un roulis de $-6{,}3^\circ$, mais avec un second minimum presque aussi profond à $+5^\circ$ (figure~\ref{fig:ann_polar_b}) : le profil est beaucoup moins tranché que sur le point A.

\figureIRD[fig:ann_polar]{align_polar_roll.jpg}{Appariement polaire sur le point A, dont l'image infrarouge a été tournée volontairement de $2{,}4^\circ$. À gauche : cartes d'énergie des deux images en coordonnées polaires (angle verticalement, distance au centre horizontalement) ; un roulis est un décalage vertical. À droite : coût en fonction du décalage angulaire, et parabole ajustée autour du minimum ($2{,}26^\circ$).}{1}

\figureIRD[fig:ann_polar_b]{align_polar_roll_b.jpg}{La même chose sur le point B, sans roulis artificiel, après la recherche grossière : minimum à $-6{,}3^\circ$, mais un second minimum près de $+5^\circ$ est presque aussi profond.}{1}

\section{Historique du développement}\label{sec:ann_history}
L'historique git du projet donne les jalons suivants.
\begin{itemize}
	\item \textbf{Août 2021} : alignement semi-automatique. L'utilisateur tourne l'image infrarouge avec trois curseurs (lacet, tangage, roulis, équation~\eqref{eq:ann_rotation_homography}) ; le résultat est affiné par corrélation de phase, par corrélation de phase polaire pour le roulis, ou par des points SIFT ou une corrélation de phase par blocs avec un ajustement affine. Comparaison sur la valeur absolue du gradient plutôt que sur les intensités.
	\item \textbf{Septembre 2021} : fonctions de coût (SSD, NTG) sur plusieurs représentations, recherche de Newton et ajustement quadratique du minimum (d'après le matériel supplémentaire de HDR+~[\hyperref[ann_hasinoff]{10}]), puis recherche itérative pyramidale ajustant une homographie sur les vecteurs des blocs, accélérée avec Numba.
	\item \textbf{Octobre 2021} : déformation locale suivant le champ de vecteurs, caméra virtuelle au champ élargi pour le fish-eye, chaîne complète recherche grossière + pyramide, choix du minimum par la hessienne, rééchantillonnage unique (\og You Only Warp Once \fg), sorties locale et globale.
	\item \textbf{Mars 2022} : ticket GitHub \#4, \og fix local alignment issues \fg{} (toujours ouvert) : utiliser la courbure du coût pour décider si un vecteur local est fiable, régulariser le champ avec ses voisins.
	\item \textbf{Novembre--décembre 2022} : ticket \#33, roulis de la caméra infrarouge ; appariement polaire des cartes d'énergie, raffinement parabolique du minimum (branche expérimentale).
	\item \textbf{Novembre 2026} : accélération sans changement des résultats (coût par linéarité du filtre de Sobel, pyramides en float32 avec OpenCV, réduction par 32 calculée sur les seules lignes et colonnes utiles) : environ 15\,s par shoot point ; réalignement à la main depuis la revue dans une fenêtre dédiée (§\ref{btn:review_redo}).
\end{itemize}

\section{Traces de diagnostic}\label{sec:ann_traces}
Les figures de cette annexe proviennent des traces de diagnostic de l'alignement. Elles ne sont pas écrites par défaut au Step~4 (environ 100\,Mo de figures par point). On les obtient :
\begin{itemize}
	\item en lançant le script \texttt{VIS\_XXXX\_REDO.bat} (ou \texttt{.sh}) écrit à côté de chaque image alignée : un \og full recovery run \fg{} avec l'option \texttt{--debug}, qui ouvre la fenêtre à curseurs puis écrit tout ;
	\item depuis la revue, quand le niveau de détail du journal des réglages du Step~4 (\bouton{⚙}) est \og Debug \fg{} : un réalignement à la main fait alors le calcul complet ;
	\item en ligne de commande : \texttt{python automatic\_registration.py --images VIS\_XXXX.dng NIR\_YYYY.dng --outdir OUT --debug}.
\end{itemize}
Les traces sont écrites dans le dossier \texttt{VIS\_XXXX\_align\_traces} :
\begin{itemize}
	\item \texttt{Full Search\_\ldots} : la recherche grossière (surface de coût, tableau \texttt{.npy} des coûts) ;
	\item \texttt{\_PADDED\_LOWRES\_*}, \texttt{\_LOWRES\_*} : cartes d'énergie à 1/32 avant et après la recherche grossière ;
	\item \texttt{ds16\_\ldots\_cost\_surface\_Y\_X\_\_itNN.png} : les entrées et les quatre surfaces de coût du bloc $(Y, X)$ à l'itération NN ; \texttt{\ldots overview\_cost\_surfaces\ldots} : tous les blocs d'une itération ; \texttt{\ldots\_cost\_search\_\_itNN.npy} : les coûts, les centres des blocs et l'échelle, lisibles par \texttt{numpy.load(\ldots, allow\_pickle=True)} ;
	\item \texttt{\_msr\_itNN\_dsMM\_\ldots}, \texttt{\_image\_itNN\_\ldots} : cartes d'énergie et images à chaque itération ;
	\item \texttt{FLOW\_itNN\_\ldots} : champ de vecteurs et image déformée par celui-ci ;
	\item \texttt{FULLRES\_REF.jpg}, \texttt{FULLRES\_REGISTERED\_COARSE.jpg}, \texttt{FULLRES\_REGISTERED\_REFINED\_WARPED\_ONCE\_ONLY.jpg} : la référence et l'image infrarouge après chaque étape, à comparer en basculant de l'une à l'autre ;
	\item \texttt{motion\_model.npy} : le modèle de mouvement.
\end{itemize}

\section{Limites et problèmes ouverts}\label{sec:ann_limits}
L'alignement multispectral est un problème mal posé, et l'algorithme échoue dans des situations identifiables :
\begin{itemize}
	\item les deux caméras ne voient pas la même scène : rafale de vent, grand écart de temps, basse altitude (parallaxe) ;
	\item la zone a peu de structures : eau, champ uniforme, canopée forestière ; le coût est alors plat ou trompeur (figure~\ref{fig:ann_profiles}c) ;
	\item le désalignement dépasse ce que les étapes savent corriger : un roulis de plusieurs degrés, que la recherche grossière ne modélise pas ;
	\item la recherche grossière tombe dans un minimum trompeur.
\end{itemize}
Le point B cumule plusieurs de ces difficultés : grand lacet ($8{,}2^\circ$), roulis probable de plusieurs degrés, grandes zones de végétation dont les textures diffèrent entre les bandes. Les vecteurs des blocs restent à leur limite de $\pm 80$ puis $\pm 20$ pixels (figures~\ref{fig:ann_blocks_hard} et \ref{fig:ann_convergence}) ; l'ajustement RANSAC retient alors une homographie aux forts termes de perspective, qui déplace les coins de l'image de plusieurs centaines de pixels. Le résultat est acceptable près du centre et faux vers les coins (textures étirées en bas à gauche). Dans nos essais, partir du roulis estimé par appariement polaire ($\pm 6{,}3^\circ$) ne l'a pas fait converger non plus : ce point demande un réalignement à la main. De tels échecs sont détectés à la revue (chapitre~\ref{ch:review}), qui mesure le décalage résiduel des bords.

\figureIRD[fig:ann_blocks_hard]{align_blocks_hard.jpg}{Première itération du raffinement du point B : les surfaces de coût des blocs de végétation n'ont pas de minimum net, les minima sont à des positions aléatoires, souvent au bord de la fenêtre de recherche (vecteurs agrandis 10 fois).}{1}

Le réalignement à la main (§\ref{btn:review_redo}) remplace les angles initiaux nuls par des angles choisis à l'œil, roulis compris. Ces angles ne sont qu'un point de départ : la recherche grossière et le raffinement pyramidal repartent de là et trouvent en général un alignement précis.

Améliorations possibles, en partie listées dans les tickets GitHub \#4 et \#33 :
\begin{itemize}
	\item utiliser la courbure de chaque surface de coût (matrice $A$ de l'ajustement quadratique) comme indice de confiance : un vecteur mesuré sur une surface plate ne devrait compter ni dans l'homographie ni dans la correction locale ;
	\item régulariser le champ de vecteurs local (médiane des voisins, contrainte de régularité) au lieu de l'utiliser brut ;
	\item ajouter le roulis à la recherche grossière (appariement polaire, ou recherche sur quelques valeurs de roulis) ;
	\item normaliser le coût (dénominateur du NTG), pour ne pas favoriser un bloc moins texturé ;
	\item calibrer les couleurs des deux caméras (calibration radiométrique) pour calculer des indices physiquement significatifs.
\end{itemize}

\section{Références}\label{sec:ann_references}
{\renewcommand{\labelenumi}{[\arabic{enumi}]}
\begin{enumerate}
	\item\label{ann_hz} R. Hartley, A. Zisserman, \emph{Multiple View Geometry in Computer Vision}, 2\textsuperscript{e} édition, Cambridge University Press, 2004. Chapitres 2 (géométrie projective du plan, homographies), 4 (estimation d'une homographie, DLT, RANSAC), 6 (modèles de caméra), 7.4 (distorsion radiale), 8.4 (caméra en rotation, rotations synthétiques, mosaïques).
	\item\label{ann_szeliski} R. Szeliski, \emph{Computer Vision: Algorithms and Applications}, 2\textsuperscript{e} édition, Springer, 2022 (chapitre 2 : formation de l'image ; chapitres 8 et 9 : alignement et estimation du mouvement). Gratuit en ligne sur \texttt{szeliski.org/Book}.
	\item\label{ann_brown} D. C. Brown, \emph{Close-range camera calibration}, Photogrammetric Engineering 37(8), 1971.
	\item\label{ann_zhang} Z. Zhang, \emph{A flexible new technique for camera calibration}, IEEE Transactions on Pattern Analysis and Machine Intelligence 22(11), 2000.
	\item\label{ann_viola} P. Viola, W. M. Wells III, \emph{Alignment by maximization of mutual information}, International Journal of Computer Vision 24(2), 1997.
	\item\label{ann_scharr} H. Scharr, \emph{Optimale Operatoren in der digitalen Bildverarbeitung}, thèse de doctorat, université de Heidelberg, 2000.
	\item\label{ann_roberts} L. G. Roberts, \emph{Machine perception of three-dimensional solids}, thèse de doctorat, MIT, 1963.
	\item\label{ann_ntg} S.-J. Chen, H.-L. Shen, C. Li, J. H. Xin, \emph{Normalized Total Gradient: a new measure for multispectral image registration}, IEEE Transactions on Image Processing 27(3), 2018 (\texttt{arXiv:1702.04562}).
	\item\label{ann_irani} M. Irani, P. Anandan, \emph{Robust multi-sensor image alignment}, ICCV 1998.
	\item\label{ann_hasinoff} S. W. Hasinoff et al., \emph{Burst photography for high dynamic range and low-light imaging on mobile cameras}, SIGGRAPH Asia 2016 (matériel supplémentaire : alignement pyramidal par blocs, raffinement sous-pixel quadratique).
	\item\label{ann_burt} P. J. Burt, E. H. Adelson, \emph{The Laplacian pyramid as a compact image code}, IEEE Transactions on Communications 31(4), 1983.
	\item\label{ann_lucas} B. D. Lucas, T. Kanade, \emph{An iterative image registration technique with an application to stereo vision}, IJCAI 1981.
	\item\label{ann_baker} S. Baker, I. Matthews, \emph{Lucas-Kanade 20 years on: a unifying framework}, International Journal of Computer Vision 56(3), 2004.
	\item\label{ann_ransac} M. A. Fischler, R. C. Bolles, \emph{Random sample consensus: a paradigm for model fitting with applications to image analysis and automated cartography}, Communications of the ACM 24(6), 1981.
	\item\label{ann_kuglin} C. D. Kuglin, D. C. Hines, \emph{The phase correlation image alignment method}, Proceedings of the IEEE International Conference on Cybernetics and Society, 1975.
	\item\label{ann_reddy} B. S. Reddy, B. N. Chatterji, \emph{An FFT-based technique for translation, rotation, and scale-invariant image registration}, IEEE Transactions on Image Processing 5(8), 1996.
	\item\label{ann_rabatel} G. Rabatel, S. Labbé, \emph{Registration of visible and near infrared unmanned aerial vehicle images based on Fourier-Mellin transform}, Precision Agriculture 17, 2016.
	\item\label{ann_vayssade} J.-A. Vayssade, G. Jones, J.-N. Paoli, C. Gée, \emph{Two-step multi-spectral registration via key-point detector and gradient similarity}, 2020.
\end{enumerate}}

\section{Détermination des \og Pitch \fg{ } et \og Yaw \fg{ } théoriques.}\label{sec:ann_theorique}

\subsection{Description du positionnement des caméras.}

La première étape du recalage de l'image infrarouge sur l'image visible consiste à appliquer une translation à l'image NIR de façon à aligner son centre avec celui de l'image VIS. En effet à la prise de vue les axes de visée des deux caméras ne sont pas parallèles. De plus le déclenchement des caméras n'est pas simultané. La caméra du drone est montée sur un \og gimbal \fg{} qui maintient la direction de visée quels que soient les mouvements du drone
\footnote{Dans la limite des spécifications indiquées par le constructeur DJI et reproduites sur la figure \ref{AnglesDroneGimbal}}.

\figureIRD[AnglesDroneGimbal]{AnglesDroneGimbal.pdf}{Description schématique du \og gimbal \fg{} et des amplitudes angulaires autorisées.}{0.6}

A l'inverse la caméra NIR est fixée rigidement au drone (figure \ref{Definition_Reperes_Cameras}). Ainsi son angle de visée va suivre les mouvements de ce dernier. Le déclenchement se fait en continue avec une période  $\Delta T^{Nir} = 3s$. C'est la fréquence la plus élevée possible pour la SJCam M20.\\


\figureIRD[Definition_Reperes_Cameras]{Definition_Reperes_Cameras.pdf}{Définition des angles de lacet, tangage et roulis (yaw, pitch, roll) du drone, du gimbal et des deux caméras VIS et NIR. Les caméras pointent vers le sol alors que le drone est horizontal ce qui engendre la permutation des axes yaw et roll entre le drone et le gimbal d'une part et les caméras d'autre part.}{0.8}

Nous allons prendre l'exemple du pitch en admettant un découplage entre les angles. Ceci correspond à faire l'hypothèse implicite que les angles yaw et roll sont nuls (ou très faibles).  Nous définissons dans le plan vertical passant par l'axe  du drone les différents angles permettant de repérer la position du drone et des deux caméras à un instant donné (figure \ref{Definition_des_angles_de_Pitch}): 
\figureIRD[Definition_des_angles_de_Pitch]{Definition_des_angles_de_Pitch.pdf}{Définition des angles dans le plan vertical.}{0.8}

\begin{itemize}
	\item $\alpha\qquad$  Flight Pitch. Angle de tangage du drone.
	\item $\delta\qquad$  NIR misalignement Pitch. Cet angle correspond à un défaut de positionnement du support de la caméra NIR sur le drone.
	\item $\varphi$ Angular wedge. Correspond à une éventuelle cale de correction du défaut $\delta$.
	\item $\zeta\qquad$ Residual deviation Pitch. $\quad \zeta = \delta + \phi$.
	\item $\theta_{Nir}\:\:\;$ Camera NIR Pitch.  $\quad \theta_{Nir} = \alpha + \zeta$.
	\item $\beta\qquad$ Gimbal Pitch. Angle de tangage de gimbal (et de la caméra VIS).
	\item $\delta_{Vis}\;\;\;$ Offset VIS. Par définition $\delta_{Vis} =90^\circ$.
	\item $\theta_{Vis}\;\;\;$ Camera VIS Pitch.  $\quad \theta_{Vis} = \beta + \delta_{Vis}$. 
	\item $\psi\qquad$  Theoritical NIR Pitch. $\quad \psi =\delta_{Vis}+\beta-\alpha$.
	\item $\psi_{r}\quad\;\;$ NIR Pitch. Angle entre les axes de visées des deux caméras au même instant $t$. $\qquad \qquad \zeta = \psi_{r} -\psi$. \\
\end{itemize}

Les données Exif des images de la caméra VIS indiquent les coordonnées GPS mais aussi les angles de yaw, pitch et roll du drone et du gimbal. Autrement dit $\alpha$ et $\beta$ sont connus à chaque acquisition d'une image de la caméra VIS. L'inclinaison de l'axe de visée engendre une déformation de l'image en trapèze comme le montre la figure \ref{Effet_du_pitch_sur_camera_pinhole}.

\figureIRD[Effet_du_pitch_sur_camera_pinhole]{Effet_du_pitch_sur_camera_pinhole.pdf}{Effet du tangage sur la déformation de l'image d'un plan horizontal formée sur le capteur d'une caméra sténopé (\og pinhole\fg{}). On notera que l'image sur le capteur est inversée.}{0.6}


Pour simplifier nous admettrons dans la suite que le  montage de la caméra NIR est sans défaut, autrement dit $\delta=0$. 
Nous supposons que les horloges des deux caméras VIS et NIR sont synchronisées. 
On dispose donc des instants exacts où les images de ces deux caméras ont été prises. 
Les paires d'images(VIS, NIR) sont formées de telle manière que l'intervalle de temps séparant les prises de vues de ces deux images soit minimum.

\subsection{Calcul du pitch théorique dans le cas synchronisé.}

\figureIRD[Pitch_cas_synchro]{Pitch_cas_synchro.pdf}{Position du drone et des deux caméras dans le cas d'images prises au même instant. Le pitch a été exagéré pour les besoins de la figure.}{0.5}

Si les images VIS et NIR sont prises au même instant (cas d'une prise de vue synchronisée
\footnote{On ne confondra pas la synchronisation des horloges et la synchronisation des prises de vue.}) il est possible de déterminer à chaque prise de vue l'angle $\psi $ entre les axes de visée des deux caméras. 

Notons que la distance entre les centres des caméras $C^{Nir}$ et $C^{Vis}$ (\og base line \fg{}) est négligeable devant l'altitude $H$ du drone par rapport au sol. 
Le rapport est inférieur à $4\;10^{-4}$. 
On peut donc  prendre $\phi \cong \theta_{Vis}$.

Enfin puisque nous supposons la caméra NIR fixée sans défaut $\theta_{Nir} =\alpha_k$. \\

En observant la figure \ref{Pitch_cas_synchro} et en utilisant des formules de géométrie élémentaires nous en déduisons la formule du pitch théorique pour l'image infrarouge dans le cas d'un déclenchement synchronisé à l'instant $t_k$
\begin{equation}
	\label{Theoritical_Pitch_syncho_2}
	\psi_k  =\phi-\theta_{Nir} =\beta_k + \delta_{Vis} - \alpha_k  \quad .
\end{equation}

La distance entre les centres des images, calculée au niveau du sol est
\begin{equation}
	L = H \tan \psi_k \quad ,
\end{equation}
et  la translation à appliquer à l'image NIR pour la centrer sur l'image VIS s'écrit, en notant $\mathrm{f}^{Nir}$ la distance focale de la caméra infrarouge 
\begin{equation}
	\mathbf{t}_y = \mathrm{f}^{Nir} \tan \psi_k \quad .
\end{equation} 

Bien entendu cette situation est un cas idéal
\footnote{C'est la configuration retenue pour le matériel professionnel.}.

\subsection{Estimation du pitch théorique dans le cas désynchronisé.}
En pratique le déclenchement des caméras VIS et NIR n'est pas simultané. 
Avec une acquisition en \og hyperlapse \fg{}
\footnote{Le terme \og time-lapse\fg{}  ou chronophotographie  correspond à la prise de clichés avec un intervalle de temps régulier et pour une position fixe de la caméra. En \og hyperlapse \fg{} au contraire la caméra est mobile avec un mouvement régulier.}
l'intervalle pour la caméra VIS est de $2s$
\footnote{En pratique l'intervalle est plutôt égal à $2,019s$. Autrement dit environ toutes les 50 images VIS la date d'enregistrement présente un saut de $1s$. Ce phénomène est corrigé par notre algorithme afin d'obtenir une \og ligne temporelle\fg{} précise au dixième de seconde.} alors qu'il est de $3s$ pour la caméra NIR. 
La figure \ref{Periodicite-du-pitch-NIR} montre un exemple d'une série d'images hyperlapse  replacées sur un même axe  temporel une fois les horloges synchronisées.\\

\figureIRD[Periodicite-du-pitch-NIR]{Periodicite-du-pitch-NIR.pdf}{Images hyperlapse des caméras VIS et NIR. Elles sont désynchronisées d'un délai variable positif ou négatif avec une périodicité de trois images visibles, soit $6s$. Pour vérifier la désynchronisation des prises de vues nous avons utilisé la présence d'un véhicule mobile. Vol FLY-20211109-1m/s.}{1}


Cette fois la base line $\mathbf{Bl}(t)$ à l'instant de la prise de vue de l'image NIR ne peut plus être négligée. Les formules de calcul du pitch deviennent
\footnote{ $(\mathbf{e}_1, \mathbf{e}_2, \mathbf{e}_3)$ repère orthonormé direct. $\mathbf{e}_2$ est dans l'axe du drone et $\mathbf{e}_3$ est sur l'axe vertical vers le haut. }

\begin{equation}
	\label{Theoritical_Pitch_non_syncho_0}
	\begin{aligned}
		\psi_k&=\phi-\alpha(t) \quad,\\
		\phi&= \arctan \left(\frac{\mathbf{Bl}(t) \cdot \mathbf{e}_2}{H} +\tan \theta_{Vis} \right)\quad.\\	
	\end{aligned} 
\end{equation}\\

Malheureusement le pitch du drone $\alpha(t)$ n'est pas connu à l'instant $t$ où l'image NIR a été prise. Il en est de même de la position exacte du drone à l'instant $t$, et donc de la \og base line \fg{} $\mathbf{Bl}(t)$. \\

Nous supposerons qu'entre deux images VIS espacées de $\Delta T^{Vis} =2s$ le drone se déplace en ligne droite et à une altitude constante.
On propose de déterminer $\alpha(t)$ et $\mathbf{Bl}(t)$ à l'aide d'une simple interpolation linéaire.\\

Le délai entre la prise de vue des images VIS et NIR est connu et noté $\delta t= t_k -\,t$ . 
Il peut être positif ou négatif selon que l'image NIR est prise avant ou après l'image VIS. 
Les positions GPS du drone par rapport au sol sont connues.  
Dans le repère géographique 
\footnote{$(\mathbf{e}_{\mathrm{W}}, \mathbf{e}_{\mathrm{N}}, \mathbf{e}_3)$ repère orthonormé direct.}
nous utiliserons les coordonnées UTM $(X, Y, Z)$. Bien entendu il faut déterminer les déplacements du drone dans son propre repère. L'orientation du drone par rapport au nord géographique est donnée par l'angle de lacet $\textit{Yaw}_k$. Soit\\ 
\begin{equation}
	\label{Theoritical_Pitch_non_syncho_11}
	\begin{aligned}
		{x}^{EW}_k&=X(t_{k+1})-X(t_{k}) \quad,\\
		{y}^{SN}_k&=Y(t_{k+1})-Y(t_{k}) \quad,\\
		{x_1}_k&=-{x}^{EW}_k\; \cos \textit{Yaw}_k + {y}^{SN}_k\; \sin \textit{Yaw}_k \quad,\\
		{x_2}_k&=\;\;{x}^{EW}_k\:\: \sin \textit{Yaw}_k + {y}^{SN}_k\; \cos \textit{Yaw}_k  \quad.\\
	\end{aligned} \\
\end{equation}


\begin{equation}
	\label{Theoritical_Pitch_non_syncho_1}
	\begin{aligned}
		si \;\delta t <0 \qquad& \qquad\alpha(t)&=&\;\alpha_k \left(\frac{\delta t}{ \Delta T^{Vis} }+1\right) - \alpha_{k+1}\,\frac{\delta t}{ \Delta T^{Vis}} \quad,\\
		&\mathbf{Bl}(t) \cdot \mathbf{e}_2&=&\; {x_2}_k\,\frac{\delta t}{ \Delta T^{Vis}} \quad,&\\
		&&&\\
		si \;\delta t >0 \qquad &\qquad \alpha(t)&=&\;\alpha_{k-1} \,\frac{\delta t}{ \Delta T^{Vis}}- \alpha_k\, \left(\frac{\delta t}{ \Delta T^{Vis} }-1\right)  \quad,\\
		&\mathbf{Bl}(t) \cdot \mathbf{e}_2&=&\; {x_2}_{k-1}\,\frac{\delta t}{ \Delta T^{Vis}} \quad.&\\
	\end{aligned} \\
\end{equation}\\

Un exemple des positions du drone et des caméras est représentée sur la figure \ref{Pitch_cas_non_synchro}.\\ 

\figureIRD[Pitch_cas_non_synchro]{Pitch_cas_non_synchro.pdf}{Cas d'une désynchronisation \og arrière \fg. L'image NIR a été prise après l'image visible. La \og base-line \fg{} a été exagérée. En pratique le point vert est situé avant le point rouge pour une désynchronisation arrière.}{0.9}


L'écart de désynchronisation $\delta t$ est périodique et nous noterons sa période  $T$. 
Dans notre exemple $T= 2 \Delta T^{Nir} = 3 \Delta T^{Vis}$. 
Le pitch de la caméra VIS, contrôlé par le gimbal, est quasi constant. 
En l'absence de vent et avec une vitesse constante le pitch du drone est lui aussi constant. 
Dans ces conditions idéales le signal $\psi(t)$ sera lui aussi de période T. 
Dans le cas d'un vol avec du vent le pitch du drone va varier pour compenser les turbulences. 
Ce phénomène peut être traduit par un pitch aléatoire 
\footnote{On ne rentrera pas dans des considérations très complexes sur la modélisation de la turbulence. Par exemple l'ergodicité,  la turbulence homogène  etc. }. 
Le pitch aura donc une composante périodique et une composante aléatoire. Le spectre de Fourier d'un tel signal présentera  un pic à la fréquence $0,166\, Hz$ ce qui correspond à la composante périodique de période $T=6\,s$.\\ 

\figureIRD[Image ViR]{_VIR__local_HYPERLAPSE_0057.jpg}{Image multi-spectrale VIR produite par IRDrone. FLY-20211109-1m/s.}{0.65}

\subsection{Analyse sur des données réelles.}

On utilise ici les données de la mission FLY-20220125 effectué le 25 janvier 2022 en fin de matinée (10 h 59, heure solaire). Le ciel était clair, le vent faible (< 3 m/s) du nord-ouest. Après la phase de synchronisation le drone s'est élevé à une altitude de 100m puis a débuté son vol d'acquisition des images \og utiles \fg{}. Pendant cette phase l'orientation du drone est restée constante ($-82^\circ$) de même que sa vitesse $1,5\,m/s$. Pendant les $232 s$ de la phase d'acquisition  les caméras VIS et NIR ont pris respectivement $116$ et $85$ images. La trajectoire du  drone se décompose en trois segment droits:
\begin{itemize}
	\item Vol vers l'avant en ligne droite. Direction $\sim 281^\circ$. Durée $166s$.
	\item Vol vers la droite. Direction  $\sim 8^\circ$ (normale à la direction précédente). Durée $26s$.
	\item Vol vers l'arrière. Direction  $\sim 95^\circ$ (parallèle à la direction initiale). Durée $40s$.\\
\end{itemize}
\figureIRD[Fly_Blassac_trajectoire]{Fly_Blassac_trajectoire.png}{Schéma de la mission FLY-20220125. Les shoot point qui ont servi au mapping sont repérés par les marques rouges. Ils ont un $|dt| \leq 0.5\,s$}{0.8}

Après synchronisation des horloges l'algorithme a formé 116 paires d'images (VIS, NIR). Notons qu'une image NIR sur deux sera associée à deux images VIS successives. 

\figureIRD[Pitch_Yaw_reel_vs_theorique]{Process-and-Theoretical_Offset_Pitch_vol_reel_01.pdf}{Évolution temporelle des angles yaw et pitch pendant la phase d'acquisition. Les angles sont déterminés par les formules théoriques et par le  \og coarse process \fg{}. Les angles théoriques ont été corrigés pour tenir compte des défauts de positionnement de la caméra NIR sur le drone.  FLY-20220125.}{0.99}

La caméra VIS du drone est placée en mode \og gimbal lock \fg{}
\footnote{Dans ce mode particulier il n'est pas certain que  l'information sur le lacet du gimbal soit fiable. La différence entre les angles de lacet du drone et du gimbal est constante et égale à $2,7^\circ$. Pour certain vols cette différence peut atteindre $7^\circ$. Ici l'angle de lacet correspond au roulis des caméras. Cependant en comparant les paires d'images VIS et NIR on ne constate pas de rotation entre les deux. C'est un problème de calibrage de la caméra au moment de passer en gimbal lock.}.
Les informations récupérées dans les données Exif des images du drone indiquent que les angles yaw, pitch et roll du gimbal sont constants ($-78,9^\circ$, $-90^\circ$  et $0^\circ$ respectivement). \\


Le yaw de la caméra NIR est pris égal au  roll du drone. 
Le calcul du yaw théorique à appliquer à l'image NIR pour la recentrer sur l'image visible utilise le même type de formules que celle mises en  évidence pour le pitch théorique (courbe \og Theoretical yaw \fg{} sur la figure \ref{Pitch_Yaw_reel_vs_theorique}). \\ 
L'évolution des valeurs théoriques des angles yaw et pitch suit globalement celle obtenue par le \og coarse process\fg{}.\\ 


Notons qu'il y  a un décalage systématique ($\sim +0,9^\circ$ ) sur le yaw. Ce décalage a été retrouvé sur l'ensemble des vols effectuées jusqu'à présent. Il peut être attribué à un défaut de positionnement de la caméra NIR sur le drone. Un décalage systématique existe aussi pour le pitch \og{}théorique \fg{} par rapport au pitch \og{} process \fg{}. Nous l'attribuons lui aussi à un défaut de positionnement de la caméra NIR sur le drone qui engendre une erreur systématique d'alignement de l'axe de visée de la caméra. Ce défaut résiduel est $\zeta = -1,3^\circ$.
\footnote{Ce que confirme une mesure réalisée à l'aide d'un inclinomètre numérique mesurant les angles avec une précision de $1/10\, ^\circ$. La mesure a été prise entre le support carbone et la surface du filtre NIR.} .\\

L'évolution des angles présente sur certaines parties un aspect oscillant ce que confirme l'analyse spectrale (DFFT) . 
Le spectre de Fourier du pitch représenté sur la figure \ref{Fourier_Pitch_process} montre  un pic net pour la fréquence $0,164\,Hz$ pour les valeurs théoriques ou obtenues par le  \og coarse process \fg{}.  
Cela correspond à une composante périodique de période $T=6\,s$. 
Nous avons montré que ce phénomène est liée à la variation périodique de la \og base line \fg{} (voir la figure \ref{Pitch_cas_non_synchro} et les formules  \eqref{Theoritical_Pitch_non_syncho_0}, \eqref{Theoritical_Pitch_non_syncho_1} ).
Le reste du spectre correspond à l'effet du vent.
Lorsque le drone vol en ligne droite vers l'avant ou l'arrière la projection de sa \og bas-line \fg{} sur son axe est périodique. 
Le pitch  a donc une composante périodique lors de cette phase de vol. 
La projection de sa \og bas-line \fg{} sur l'axe normal est nulle. 
Dans cette phase de vol le yaw a donc seulement la composante aléatoire associée au vent. 
Lorsque le drone se déplace perpendiculairement à son axe (en gardant toujours la même orientation) c'est le yaw qui possède une composante périodique alors que le pitch n'a qu'une composante aléatoire liée  au vent. 
Ce phénomène est clairement visible sur la figure \ref{Pitch_Yaw_reel_vs_theorique}. Les changements d'allure des courbes permettent même d'identifier les phases de vol.\\


\figureIRD[Fourier_Pitch_process]{Spectre-Pitch-process-vol-reel_test.png}{Transformée de Fourier du Pitch  calculé  par le \og coarse process\fg{} à partir des images numérique ou par les formules théoriques.  FLY-20222501.}{0.9}

Notons que le pitch théorique calculé avec la formule synchronisée \eqref{Theoritical_Pitch_syncho_2} ne fait intervenir que les angles du drone et du gimbal. Il ne présente pas de caractère périodique comme le confirme sa DFFT. Ceci est cohérent puisque la variation de ces angles est associée aux effets \og{}aléatoires \fg{} du vent sur le drone (turbulence).


\subsection{Conclusion.}


La méthode proposée pour déterminer a priori le lacet et le tangage (angles yaw et pitch) à appliquer à l'image NIR pour la recentrée sur l'image VIS repose sur plusieurs éléments :\\

\begin{itemize}
	\item Les images VIS et NIR doivent être prises en mode time-lapse. 
	La fréquence d'enregistrement de la caméra VIS (celle qui a le pas le plus petit, ici $2\,s$) doit être régulière. 
	Ceci nécessite l'utilisation de carte SD rapides. Toutefois pour le DJI les dates sont enregistrées à la seconde près. On remarque périodiquement des sauts d'une seconde. Il semble que le pas ne soit pas exactement un entier. Pour notre vol FLY-20211109-1m/s on trouve $2,01865 \, s$. Ce phénomène est retrouvé sur tout les vols effectués jusqu'à présent avec une valeur de $2,02 \, s$. Le logiciel IRDrone détermine automatiquement pour chaque vol la période exacte pour obtenir une fréquence constante.  
	On utilise le format de date de Python en  renseignant son attribut \texttt{micosecond}. 
	Notons que les données Exif des dates prévoient (pour certaines) une valeur décimale des secondes  
	La situation pour la caméra NIR est différente. Il existe aussi des sauts mais de la valeur du timelapse ($3\,s$). Les dates sont donc toujours des multiples du timelapse. Ici c'est clairement un défaut d'enregistrement d'une image
	\footnote{Comme on a pu le vérifier lors de la synchronisation manuelle des images, la séquence retenue  comportant justement un tel défaut d'enregistrement.}. Les dates n'ont donc pas été rectifiées. Lors d'un saut de ce type la \og base-line\fg{} est plus grande. Cela reste acceptable si la vitesse du drone est faible. Notons toutefois qu'il faut adapter les formules de calcul pour en tenir compte. A nouveau le code gère automatiquement ce phénomène.\\
	
	\item Les données d'attitude  (yaw, pitch, roll) doivent être disponibles pour les deux caméras. C'est bien le cas avec notre montage puisque les données Exif des images VIS (format dng) du DJI fournissent ces angles pour le drone (donc pour la caméra NIR qui lui est fixée rigidement) et pour le gimbal (donc pour la caméra VIS qui est fixée dessus).\\  
	\item Le calcul de la \og base-line \fg{} nécessite de connaître les coordonnées UTM. Actuellement l'altitude $H$ du drone par rapport au sol est déterminée en connaissant l'altitude du point de décollage, l'altitude du drone par rapport au point de décollage et l'altitude à l'aplomb de la prise d'image (via une API de l'IGN pour la France). .\\
	\item La synchronisation des horloges des deux caméras est un point critique. On utilise une séquence de synchronisation à effectuer à chaque décollage. Elle consiste à faire tourner
	\footnote{La cible est placée sur un plateau mis en rotation par un moteur pas à pas. Ce moteur est commandé par un micro contrôleur ARDUINO.} un QR code (ARUCO marker) sous les deux caméras. Un algorithme détecte ensuite ce QR code sur les images et en déduit l'orientation. Une méthode de minimisation permet de rapprocher  les courbes angle/date. La  précision de cette synchronisation est supérieure au $1/10^{\grave{e}me}s$.  Soulignons qu'il est impératif de ne pas interrompre la prise de vue en mode hyperlapse pendant toute la durée de la mission (phase de synchronisation et phase d'acquisition). En effet la valeur de synchronisation n'est valable (avec cette précision) que pour une séquence hyperlapse donnée. Entre deux séquences hyperlapse la valeur de synchronisation peut varier de plus d'une seconde (allant même jusqu'à 3s dans le cas le plus défavorable).
	Un tel défaut de synchronisation des horloges risque de fausser l'appariement des images (et le calcul de la  \og base-line \fg{}).
	Pour l'exemple présenté ici nous avons réalisé aussi cette synchronisation manuellement à partir d'un objet mobile présent sur 25 images consécutives (soit 50s). Le mobile avait été déplacé spécialement pour cet utilisation lors de notre vol d'essai. Cette technique est précise (précision de l'ordre du $1/4 s$) mais fastidieuse. Nous avons retrouvé le même appariement des images que celui obtenu avec la synchronisation par cible ARUCO. \\
	\item L'angle de pitch du drone ainsi que la \og base-line \fg{} à l'instant où l'image NIR est capturée ont été  interpolés linéairement. Cette méthode donne d'excellents résultats.\\
\end{itemize}

Nous avons vérifié la pertinence de l'ensemble de ces éléments sur plus d'une dizaine de vol avec des conditions météorologiques variées, notamment le vent.  
Par exemple les figures \ref{Pitch_Yaw_reel_vs_theorique_FLY2021109} et \ref{Fourier_Pitch_process_FLY2021109} correspondent à un vol programmé en hyperlapse (\og mode line\fg{}) avec une vitesse de $1\,m/s$ et environ $600$ images.

\figureIRD[Pitch_Yaw_reel_vs_theorique_FLY2021109]{Process-and-Theoretical_Offset_Pitch_vol_reel_02.png}{Évolution temporelle des angles yaw et pitch pendant la phase d'acquisition. Ici  le drone se déplace à cap constant ($179^\circ$) parallèlement à son axe. FLY-20211109.}{1}

\figureIRD[Fourier_Pitch_process_FLY2021109]{Spectre-Pitch-process-vol-reel_02.png}{Transformée de Fourier du Pitch. FLY-20211109.}{.75}




\chapter{Géoréférencement et orthorectification (Step 6)}\label{ann:georef}
\todo{Texte technique : géométrie caméra/drone, repères, projection sur le MNT, précision, destiné aux utilisateurs avertis.}

\end{document}
"
% (voir build_help.bat). Les images viennent alors de images_light/ au lieu de images/.
\ifdefined\IRDlight
  \newcommand{\IRDimgdir}{images_light}
  \graphicspath{{images_light/}}
\else
  \newcommand{\IRDimgdir}{images}
  \graphicspath{{images/}}
\fi

% Symboles des boutons du logiciel (pdflatex ne les connait pas)
\DeclareUnicodeCharacter{25C0}{\ensuremath{\blacktriangleleft}}
\DeclareUnicodeCharacter{25B6}{\ensuremath{\blacktriangleright}}
\DeclareUnicodeCharacter{2699}{\ensuremath{\circledast}}
\DeclareUnicodeCharacter{2922}{\ensuremath{\nearrow}}
\DeclareUnicodeCharacter{21BA}{\ensuremath{\circlearrowleft}}
\DeclareUnicodeCharacter{21BB}{\ensuremath{\circlearrowright}}
\DeclareUnicodeCharacter{25B2}{\ensuremath{\blacktriangle}}
\DeclareUnicodeCharacter{25BC}{\ensuremath{\blacktriangledown}}
% Bouton du logiciel
\DeclareRobustCommand{\bouton}[1]{\texorpdfstring{\fbox{\small\textsf{#1}}}{#1}}
% A completer
\newcommand{\todo}[1]{\textcolor{red}{[\`A r\'ediger : #1]}}
% Sections historiques : repliees dans le manuel web, texte normal dans le PDF
\newcommand{\legacybegin}[1]{}
\newcommand{\legacyend}{}
% Figure tolerante : affiche un cadre si l'image n'existe pas encore.
% Usage : \figureIRD[label]{fichier}{legende}{largeur}
%   - label optionnel (ex. fig:img_VIR_1), utilisable avec \ref{...} / \autoref{...}
%   - fichier avec ou sans extension ; sans extension, essaie .png, .jpg, .jpeg puis .pdf
\makeatletter
\newcommand{\ird@file}{}
\newcommand{\ird@find}[1]{%
  \def\ird@file{}%
  \IfFileExists{\IRDimgdir/#1}{\def\ird@file{#1}}{%
  \IfFileExists{\IRDimgdir/#1.png}{\def\ird@file{#1.png}}{%
  \IfFileExists{\IRDimgdir/#1.jpg}{\def\ird@file{#1.jpg}}{%
  \IfFileExists{\IRDimgdir/#1.jpeg}{\def\ird@file{#1.jpeg}}{%
  \IfFileExists{\IRDimgdir/#1.pdf}{\def\ird@file{#1.pdf}}{}}}}}}
\newcommand{\figureIRD}[4][]{%
  \ird@find{#2}%
  \begin{figure}[htbp]\centering
    \ifx\ird@file\@empty
      \fbox{\parbox{0.6\linewidth}{\centering\small\textcolor{gray}{Image manquante : \detokenize{#2}}}}%
    \else
      \includegraphics[width=#4\linewidth]{\ird@file}%
    \fi
    \caption{#3}%
    \if\relax\detokenize{#1}\relax\else\label{#1}\fi
  \end{figure}}
% Paire de figures cote a cote, placee a l'endroit exact ou elle est ecrite ([H]).
% Usage : \figurePaireIRD{label}{fichier gauche}{legende gauche}{fichier droit}{legende droite}{legende globale}
%   - le label designe la paire ; les sous-figures ont pour labels <label>_avant et <label>_apres
\newcommand{\ird@incsub}[1]{%
  \ird@find{#1}%
  \ifx\ird@file\@empty
    \fbox{\parbox{0.8\linewidth}{\centering\small\textcolor{gray}{Image manquante : \detokenize{#1}}}}%
  \else
    \includegraphics[width=\linewidth]{\ird@file}%
  \fi}
\newcommand{\figurePaireIRD}[6]{%
  \begin{figure}[H]\centering
    \begin{subfigure}[t]{0.48\linewidth}\centering
      \ird@incsub{#2}%
      \caption{#3}\label{#1_avant}
    \end{subfigure}\hfill
    \begin{subfigure}[t]{0.48\linewidth}\centering
      \ird@incsub{#4}%
      \caption{#5}\label{#1_apres}
    \end{subfigure}
    \caption{#6}\label{#1}
  \end{figure}}
\makeatother

\title{IRdrone\\Manuel de l'utilisateur}
\author{}
\date{\today}


\begin{document}
	
	% --------------------------------------------------
	% Cover page
	% --------------------------------------------------
	\begin{titlepage}
		\centering
						
		\vspace*{2cm}		
		%{\Huge\bfseries IRdrone\par}		
		%\vspace{0.5cm}		
		{\LARGE Manuel de l'utilisateur\par}
		\vspace{6cm}		
		\includegraphics[width=0.95\textwidth]{Logo-IRdrone.png}		
		
		\vfill		
		{\large Version du \today\par}		
	\end{titlepage}
	
	% --------------------------------------------------
	% Table of contents
	% --------------------------------------------------
	\tableofcontents



%==========================================================
\chapter{Le projet IRDrone}\label{ch:overview}


\section{Observer le couvert végétal.}

\subsection{Contexte général.}
La cartographie de la végétation est aujourd'hui classiquement employée dans de nombreux types de diagnostics environnementaux (site Natura 2000, ZNIEFF, études d'impact environnemental etc.). 
Cette cartographie est basée sur une segmentation (détourage d'ensembles homogènes au plan de la végétation) du couvert végétal en vue plane aérienne. 
Cette approche est généralement complétée par des prospections au sol (relevés botaniques, phytosociologiques etc.). 
La segmentation d'ensembles homogènes repose sur l'analyse des photographies aériennes en couleur naturelle, notamment sur leur couleur, leur texture, leur structure apparente, leur opacité etc.

\figurePaireIRD{fig:etude_de_cas}
{step0_carto_aerienne.jpg}{Délimitation des habitats sur fond image aérienne.}
{step0_habitats.jpg}{Carte restituée des habitats sur fond IGN topo 25.}
{Exemple d'exploitation de données terrain à l'aide de cartes et photos aériennes.}

Il est toutefois délicat de distinguer les essences (arbres) et les espèces végétales (espèces herbacées) sur des photographies aériennes classiques. Les dynamiques et les éléments fins du paysage sont également difficiles à cerner. 
Les photographies aériennes IR disponibles sont souvent anciennes, donc difficilement utilisables dans le cadre de diagnostics actuels. 
D'autre part la couverture nationale est lacunaire, avec des régions mal couvertes. 
Ajoutons que toutes les photos existantes sont réalisées en altitude (3-4000 m) et manquent donc de précision pour des diagnostics réalisés couramment à l'échelle du 1/5000, voire moins. \\




L'utilisation d'un drone pour acquérir des photographies présente les avantages de centrer les clichés sur les secteurs d'études, d'obtenir des images d'une précision largement supérieure (altitude de vol de 120 m), d'obtenir autant d'images récentes qu'on le souhaite, de pouvoir les réaliser à différentes saisons. 


\subsection{Intérêt des images multispectrales.}

Les photographies aériennes dans le proche infrarouge
\footnote{On ne confondra pas ce rayonnement infrarouge proche qui en journée provient du rayonnement solaire et est réfléchi par les objets avec le rayonnement infrarouge lointain qui lui est émis par tous les objets.} NIR sont plus riches en informations écologiques que les photographies en couleur naturelle. 
L'objectif principal de la photographie aérienne NIR est donc de faciliter l'interprétation des photographies aériennes en se basant sur les différences de facteur de réflexion du rayonnement infrarouge entre les types de végétations (en raison de différences structurelles entre les tissus foliaires des espèces végétales).

\figureIRD[fig:img_NDVI_1]{img_NDVI_1.png}{Carte de l'indice NDVI de développement des végétaux.}{0.7}
\figureIRD[fig:img_VIR_1]{img_VIR_1.png}{Exemple d'image composite de type VIR. Les couches R/G/B sont remplac\'ees par les couches NIR/R/G. Ce type de représentation est très utile, entre autres, pour distinguer entre résineux et feuillus.}{0.7} 


\subsection{Que cherchons-nous à observer ?}

L'emploi des photographies aériennes se justifie pour plusieurs raisons, qui ont trait essentiellement à la précision et à la rapidité d'analyse. 
D'une manière générale, sur les images aériennes NIR les contrastes entre les différents types d'habitats sont exacerbés, ce qui permet une lecture plus fine du paysage. 
Dans les peuplements forestiers denses, il est difficile d'estimer les proportions de recouvrement résineux/feuillus en se basant sur une prospection de terrain en raison des grandes surfaces à parcourir et de la hauteur des arbres. 
La différenciation entre résineux/feuillus sur les photographies en couleur naturelle est souvent ambiguë. 
Au contraire, la photographie aérienne NIR permet de discerner précisément et rapidement les contributions de chacune des essences dans les peuplements mélangés couvrant des superficies considérables, et ainsi d'avoir accès à une information concernant la gestion spécifique des peuplements (types de traitements sylvicoles, classes d'âge etc.).  D'autre part, certains peuplements arborescents mêlant feuillus et résineux en proportions variables peuvent, selon l'importance prise par les feuillus, relever ou non d'une protection européenne (directive \og Habitats \fg). 
L'évaluation des proportions des essences devient objective et précisément quantifiable à l'aide des photographies aériennes NIR. 
L'analyse des photographies aériennes NIR permet en outre une meilleure appréhension des phénomènes de dynamique de la végétation.
 
La détection des très jeunes ligneux devient possible en raison de leur vive coloration et de leur position au sein de milieux ouverts (fronts de colonisation ligneux au sein d'espaces ouverts) alors que leur présence est difficilement détectable en couleur naturelle en raison de la petite taille de leur houppier et de leur couleur comparable à celle des habitats ouverts colonisés. 
Une autre application importante concerne la localisation instantanée et exhaustive des masses d'eau (qui peuvent être discrètes, de petite taille et donc passer inaperçues en couleur naturelle) qui apparaissent comme des taches sombres  en NIR. 
Cette application est particulièrement importante en raison du statut réglementaire des zones humides, qui constituent généralement un des enjeux principaux des diagnostics environnementaux. 


\subsection{Principe de l'observation.}

Les végétaux produisent de la chlorophylle en absorbant la lumière visible du soleil (VIS). Seule une petite partie de ce rayonnement est utilisée dans ce mécanisme. Or le rayonnement solaire arrivant au sol comporte aussi des ultra-violets (UV) et des infrarouges proches (NIR) \footnote{Les rayonnements UV et NIR ne sont pas visibles par l'\oe il humain.} 

\figureIRD[Facteur_reflexion_chloro-carot]{Spectres-absorption-pigments-photosynthetiques.pdf}{Facteur d'absorption des pigments photosynthétiques dans la bande spectrale visible.}{1}

C'est essentiellement le rayonnement visible dans les bandes \og bleu \fg{} ($[400\,nm ,\, 450\,nm]$) et \og rouge \fg{} ($[630\,nm ,\, 680\,nm]$) qui contribue à la production chlorophyllienne.  
La partie non absorbée du rayonnement par la plante est transmise (la plante est \og transparente \fg{}) ou bien réfléchie (renvoyée vers l'environnement). 
Le rendement chlorophyllien étant nul dans la bande spectrale \og jaune-vert \fg ($[500\,nm ,\, 600\,nm]$) celle-ci est réfléchie par la plante (\ref{Facteur_reflexion_chloro-carot}). 
C'est de là que provient la couleur verte de ces végétaux. De même les rayonnements UV et IR ne participent pas à la production de chlorophylle. 
En revanche ils produiraient son échauffement s'ils étaient absorbés. 
Ainsi pour se protéger les plantes n'absorbent pas les rayonnements UV et IR émis par le soleil, qui sont alors transmis ou réfléchis. 
Les plantes sont totalement transparentes aux rayonnements ultraviolets.
En revanche elles réfléchissent une partie du rayonnement infrarouge.
C'est cette propriété qui va être exploitée dans le projet. 
Les résineux et les feuillus n'ont pas le même facteur de réflexion du rayonnement infrarouge. 
La figure \ref{Facteur_reflexion_vegetaux} montre les courbes typiques du facteur de réflexion spectral des résineux et des feuillus. 
Notons que les feuillus produisent aussi de la bêta-carotène et de la fucoxanthine qui absorbent le bleu et une partie du vert mais réfléchissent le rouge. 
C'est pour cette raison qu'à l'automne, lorsque les feuillus ne produisent plus de chlorophylle mais que la bêta-carotène est encore présente, ces espèces arborent des teintes allant du marron orangé à l'ocre rouge.

\figureIRD[Facteur_reflexion_vegetaux]{Facteur_reflexion_vegetaux.pdf}{Facteur de réflexion des résineux et des feuillus. Dans la bande spectrale visible $[400\,nm ,\,700\,nm]$ le comportement des deux espèces est très similaire. En revanche dans la bande spectrale infrarouge IR-A $[700\,nm ,\, 1.4\,\mu m]$ bien présente dans le spectre solaire les deux espèces se distinguent nettement.}{1}


On remarque que les feuilles des arbres réfléchissent le rayonnement autour de 550 nm qui correspond à la couleur verte perçue par l'\oe il humain. Mais les courbes étant identiques dans la partie visible du spectre (400 nm à 700 nm) on ne peut pas distinguer entre les différentes essences. En revanche les facteurs de réflexion dans le domaine infrarouge sont bien distincts entre résineux et feuillus. L'idée est d'exploiter cette propriété en produisant une image infrarouge  du couvert végétal
\footnote{Éventuellement combinée à une image dans le spectre visible.}. Sur cette image les feuillus devraient apparaître  \og plus clairs\fg{}  que les résineux. On pourra ensuite tracer une carte du couvert végétal. Signalons toutefois que selon l'époque de l'année et l'âge des sujets ces courbes peuvent évoluer rendant plus délicate l'interprétation des images.

\figureIRD[Facteur_absorption_vegetaux]{Par_action_spectrum.jpg}{Facteur d'absorption des végétaux et rendement de photosynthèse.}{0.5}


\subsection{Le matériel pour l'observation.}
Pour tracer une carte il faut des images \og vues du ciel \fg{}. Les professionnels utilisent des photos satellites (alt. 300km) ou prises d'avion (alt. 2000-3000 m). Mais comme expliqué plus haut, la résolution spatiale et temporelle n'est pas adaptée à des études fines de terrain. Dans le cadre de ce projet on envisage donc d'utiliser un petit drone qui permet d'obtenir avec une très grande souplesse des vues aériennes des zones à étudier (alt. 120 m) avec une très bonne résolution spatiale et une maîtrise totale des périodes de prise de vue.

La plupart des \og petits \fg{} drones existants ne permettent pas l'obtention simultanée de clichés en couleur naturelle et en IR. Or il est indispensable de disposer des deux types de clichés géoréférencés (IR pour interprétation de la végétation et couleur naturelle pour la réalisation de cartes standardisées). \\ 

Il existe des drones et des caméras multi-spectrales pour drones
\footnote{Par exemple DJI Phantom 4 multispectral, Parrot Sequoia, Micasense RedEdge.} . 
Il s'agit de caméras multi canaux  qui permettent d'acquérir, outre l'image visible classique, des bandes spectrales étroites dans le rouge, le vert, le bleu et dans une ou deux bandes dans la proche infrarouge. Il existe même des versions équipées de caméra thermiques \footnote{Par exemple la caméra Micasense Altum  possède un capteur Flir IR. Son poids est de 400g. Cependant il s'agit ici de rayonnement infrarouge thermique (LWIR) dont les longueurs d'onde sont dix fois plus grandes que celles du  proche infrarouge solaire.  Soulignons que ce matériel est très onéreux.}.
D'autres drones
sont équipés de deux caméras l'une dans le visible et l'autre dans l'infrarouge. 
Enfin la caméra de certains drones
\footnote{Par exemple le DJI Mavic Pro.2 équipé d'une très performante caméra Hasselblad. Son filtre \og hot mirror\fg{} est fixé directement sur le capteur. Des drones \og lourds\fg{} munis de nacelles peuvent embarquer des reflex numériques qui eux aussi peuvent être \og défiltrés \fg{}. }  
peut être rendue sensible dans le proche infrarouge. Cela nécessite de démonter
\footnote{Cette opération de \og défiltrage \fg est délicate. Elle peut être réalisée par des sociétés spécialisées.} le filtre interne du capteur numérique qui coupe le rayonnement infrarouge et de munir le système optique d'un filtre adapté
\footnote{Avec un jeu de filtres externes on peut choisir le type de bande spectrale IR. Avec un filtre passe haut IR $760\,\mathrm{nm}$ on obtiendra une bande large $[760, 1100]$. Il existe aussi des filtres à bande étroite, par exemple $850\,\mathrm{nm}\pm 25\mathrm{nm}$). On peut même remettre un filtre \og hot mirror\fg{} qui rétablira l'imagerie dans le spectre visible.}.  
De plus le drone n'enregistre alors qu'un seul type d'image selon le filtre dont il est équipé. \\  


L'idée du projet \textsf{IRDrone} est d'équiper un petit  drone de coût abordable 
\footnote{Ici un DJI Mavic Air.2  dont le poids n'excède pas 600 g.} d'une seconde  caméra. La caméra équipant le drone fournira des images en couleur dans le spectre visible (RGB). La deuxième caméra embarquée, que nous désignerons par caméra IR, devra produire des images dans la bande spectrale du proche infrarouge \og solaire \fg 
\footnote{Concrètement avec les capteurs CMOS disponibles la bande spectrale accessible s'étend de $720 \,\mathrm{nm}$ à $1100\,\mathrm{nm}$. Elle englobe les bandes étroites RedEdge et NIR (pour Near Infra Red)}.\\

La caméra NIR devra être légère pour ne pas perturber le vol du drone et ne pas diminuer son autonomie. Elle devra être d'un coût réduit ce qui conduit naturellement à se tourner soit vers des caméras de surveillance (\og vision de nuit\fg{}) soit vers des \og action cam \fg{} grand public. Les premières ne conviennent pas car elles offrent des performances médiocres en résolution. Les secondes sont conçues pour donner de belles images dans le domaine visible. Il faudra donc adapter l'optique par ablation du filtre \og hot mirror\fg{} et mise en place d'un filtre IR.  \\


Après chaque vol il sera nécessaire de faire un traitement informatique des images collectées (corrections géométriques, colorimétriques, calage et fusion d'images visibles et IR ...). Ce sera le rôle dévolu au code Python  IRDrone.py associé au projet \textsf{IRDrone} dont les fonctionnalités sont  exposées ci-après.


\section{Organisation d'une mission IRDrone}

\subsection{Préparation au bureau}

Avant le début de la partie active de la mission le pilote doit préparer soigneusement son vol. 
D'abord au bureau, en repérant les chemins d'accès, la topographie du terrain, les obstacles et la dimension de la zone à étudier. 
Cette dimension conditionne l'altitude de vol. On peut imaginer ne faire qu'un vol unique à haute altitude pour couvrir en une fois la zone, ou bien plusieurs vols successifs. Il faut garder à l'esprit que le drone DJI chargé de sa caméra infrarouge a une durée de vol limitée à 25 minutes. 
En pratique la phase de synchronisation prend 2 minutes. Avec une vitesse de $3\,\mathrm{m/s}$ le drone peut parcourir environ 4\,km, ce qui n'est pas énorme compte tenu de la nécessité de \og balayer \fg~la zone de survol afin que les clichés se superposent (si on souhaite faire un assemblage des images). 
On suppose aussi que l'opérateur a sa qualification de pilote de drone et les autorisations de survol.

\subsection{Préparatifs du vol sur le terrain}
Une fois sur le terrain le pilote installera la zone de décollage. Après balisage il positionne le stand de décollage et le système de synchronisation par mire ArUco avec son alimentation
\footnote{Que faire si on a oublié le dispositif de synchronisation ? On peut \og sauver \fg~la mission en déplaçant lentement un objet sous le drone. Il faudra ensuite estimer manuellement l'offset de synchronisation des deux caméras. Ça peut marcher\dots{} des fois !}. 
Un siège pliant est utile pour piloter en regardant le retour caméra afin de suivre au mieux le plan de vol. 
De même une petite manche à air pour apprécier la direction et la vitesse du vent est utile. 
Enfin, mesurer précisément la position du point de décollage avec un GPS en position fixe est appréciable\footnote{Voir les possibilités d'un GPS Garmin 750, qui a ce type de fonction.}. 
Avant chaque mission sur le terrain \textbf{le GPS du drone doit être calibré et sa caméra étalonnée} (direction de visée). 
Ces deux actions sont indispensables pour la réussite de l'exploitation des images par le logiciel IRDrone. 
Le pilote vérifie que les batteries sont chargées et que les cartes SD (vides) sont bien présentes sur le drone DJI et la caméra SJCam M20. 
Il est judicieux de prérégler les paramètres de la caméra du DJI (ceux qui peuvent l'être) avant le décollage. 
À ce stade le pilote peut déclencher l'hyperlapse de la caméra infrarouge (SJCam M20), positionner le drone sur son stand de décollage et enfin mettre sous tension le dispositif de synchronisation (dans la dernière version du dispositif, la mire ne se met à tourner que lorsque le drone a quitté le stand de décollage).

\subsection{Le vol de la mission}
Une fois ces préparatifs faits, la mission comporte 4 phases successives :
\begin{enumerate}
	\item \emph{Takeoff Point -- Synchronization} : synchronisation à proximité immédiate du point de décollage, à une altitude d'environ 2\,m. 
	Au début de cette phase la caméra NIR est déjà déclenchée en mode hyperlapse (fréquence de prise de vue 3\,s). 
	Pendant cette phase le pilote stabilise le drone au-dessus du dispositif de synchronisation, règle les paramètres de la caméra VIS et déclenche la prise de vue VIS en mode hyperlapse. Il reste au point fixe une dizaine de secondes.
	\item \emph{Approach} : déplacement vers la zone à survoler pour l'étude. 
	Le pilote prend rapidement de l'altitude. 
	Il se déplace à la vitesse maximum permise par le mode hyperlapse du DJI. 
	Une valeur de $3\,\mathrm{m/s}$ est un bon choix. Lors de la préparation de la mission il est fortement conseillé de choisir un point de décollage au plus près de la zone de survol.
	\item \emph{Hyperlapse Survey} : acquisition des images multispectrales au-dessus de la zone de survol (\emph{survey area}). 
	Le drone se déplace en gardant une orientation constante. 
	Il avance puis se déplace latéralement puis recule, etc. (voir figure~\ref{fig:4_step_of_the_mission}). 
	Lors du décollage il est conseillé d'orienter le drone dans le sens de déplacement principal au-dessus de la zone de survol. 
	À la sortie de la zone de survol le pilote peut interrompre l'enregistrement hyperlapse de la caméra VIS. 
	Ceci a deux avantages : cela évite de saturer la carte SD, les images étant en haute résolution, mais surtout cela libère la contrainte de vitesse du drone imposée par le mode hyperlapse.
	\item \emph{Return to Home} (RTH) : c'est la phase de retour au point de décollage. Ce retour peut se faire à vitesse maximum. Le mode RTH peut être enclenché pour un vol automatique.
\end{enumerate}

\figureIRD[fig:4_step_of_the_mission]{4_step_of_the_mission.pdf}{Les phases de vol d'une mission IRDrone sur le terrain.}{0.8} 

\subsection{Exploitation des images}

Une fois l'acquisition terminée, la phase suivante consiste à traiter numériquement les images des deux caméras à l'aide du logiciel IRDrone. 


\section{Objectif du logiciel IRDrone}
IRDrone.py est une application interactive et conviviale qui permet de générer des images multispectrales pour des études d'impact environnemental ou de botanique. 
Elle utilise pour cela les images prises lors des campagnes sur le terrain par deux caméras, l'une dans le spectre visible (\og VIS \fg) et l'autre dans le spectre du rayonnement dans le proche infrarouge (\og NIR \fg). 
Les deux caméras sont fixées sur un drone. 
Actuellement le système supporte un drone DJI Mavic Air 2 muni de sa caméra VIS montée sur un gimbal et une caméra SJCam M20 modifiée pour capturer des images dans le spectre du proche infrarouge (typiquement au-delà de 810 nm). 
L'application synchronise les horloges des deux caméras puis aligne les images VIS et NIR.
Le processus principal produit des images multispectrales alignées (4 bandes red, green, blue, nir au format tif). 
Un post-traitement permet d'orthorectifier et géoréférencer les images, soit unitairement, soit sous forme d'assemblage (mapping avec le code ODM ; figure~\ref{fig:img_mapping_1}).
\figureIRD[fig:img_mapping_1]{img_mapping_sample.png}{Assemblage à l'aide d'Open Drone Map (ODM) d'images dans le spectre VIS. La visualisation est faite avec le logiciel QGIS. Les courbes de niveau ont une équidistance de 1\,m. Les points rouges indiquent l'emplacement des clich\'es et le numéro de la prise de vue.}{0.8}
L'exploitation peut être réalisée avec le logiciel QGIS afin d'obtenir des images spécialisées (indice NDVI, figure~\ref{fig:img_NDVI_1} ; image VIR nir/red/green,
figure~\ref{fig:img_VIR_1}). 


\section{Fonctionnalités principales du logiciel IRDrone}
Au lancement, une petite fenêtre \og Loading, please wait \fg{} s'affiche pendant le chargement (quelques dizaines de secondes), puis la fenêtre principale s'ouvre (chapitre~\ref{ch:home}). Elle liste les missions du dossier de travail et, une fois une mission ouverte, ses sept étapes avec leur avancement. Un clic sur une étape la lance ; les étapes 1 à 4 s'enchaînent d'elles-mêmes.

L'exploitation d'une mission comporte 7 étapes :
\begin{enumerate}
	\item \textbf{Mission} : les images des deux caméras sont trouvées sur les cartes SD, copiées dans un nouveau dossier de mission, et les paramètres de la mission sont enregistrés.
	Le nom du dossier de mission est généré automatiquement en fonction de la localisation et de la date
	\footnote{\label{fn:MissionFolderName} Par défaut le nom du dossier est FLY-YYYYMMDD-hhmm-location. Il peut être modifié une fois pour toutes à la création de la mission dans le champ \og Location \fg. Par exemple si le champ \og Location \fg~a la valeur Blassac on peut préciser \og Blassac Bourg \fg.}.
	\item \textbf{Ligne de temps du vol} : les images de synchronisation (mire ArUco au décollage) et celles du relevé au-dessus de la zone d'étude sont détectées automatiquement ; l'utilisateur vérifie.
	\item \textbf{Synchronisation} : le décalage entre les horloges des deux caméras est mesuré et les images VIS et NIR sont appariées.
	\item \textbf{Alignement}, c\oe ur d'IRDrone : chaque paire d'images est alignée pour obtenir 4 couches spectrales superposées.
	\item \textbf{Revue} des images alignées (chapitre~\ref{ch:review}) : chaque image est vérifiée sur une ligne de temps, les mal alignées sont signalées (écartées de la carte) ou réalignées à la main.
	\item \textbf{Cartographie et géoréférencement} : géoréférencement, orthorectification, assemblage en orthophoto (mapping), modèle de terrain et courbes de niveau.
	\item \textbf{Archivage} de la mission.
\end{enumerate}

\figureIRD[fig:Air_Mission_folder]{Air_Mission_folder.png}{Le dossier principal \texttt{Air-Mission} avec deux missions en cours de traitement (ici Chalivoy-Milon et Labout) et une mission archivée (ici Blassac). Sous Windows les icônes sont attribués automatiquement.}{0.6}

\subsection*{Aide en ligne}\label{sec:help}
\begin{itemize}
	\item \textbf{Info-bulles.} Laisser le pointeur de la souris quelques instants sur un bouton (ou une case à cocher) affiche une info-bulle qui résume son action.
	\item \textbf{Touche \textsf{F1}.} La touche \textsf{F1}, le pointeur de la souris sur un bouton, ouvre ce manuel à la section qui décrit ce bouton. Sur une étape de la page d'accueil, elle ouvre le chapitre de l'étape. Ailleurs dans une fenêtre, elle ouvre le chapitre de la fenêtre, comme le bouton \bouton{?}.
	\item \textbf{Bouton \bouton{?}.} Chaque fenêtre a un bouton \bouton{?} qui ouvre ce manuel au chapitre de la fenêtre. Dans les fenêtres qui se pilotent au clavier (ligne de temps, revue), son info-bulle rappelle les touches. Il est en bas à gauche ; le bouton \bouton{Close} des fenêtres d'étape et de la revue est en bas à droite ; aux étapes 2 à 5, le bouton qui passe à l'étape suivante (\bouton{Continue with \ldots ▶}) est à sa droite et ne lance jamais un calcul tout seul. Les actions d'une étape (\bouton{Run alignment}, \bouton{Build the map\ldots}, \ldots) sont dans la fenêtre, au-dessus de la barre de progression ; la principale est bleu foncé, les autres blanches, et celle qui est irréversible (l'archivage) rouge. Les étapes~3 à~6 affichent en haut le nom de l'étape et, dessous, la mission. Dans la fenêtre principale aussi, en bas à gauche, il ouvre le manuel entier.
	\item \textbf{Boutons \bouton{↺} et \bouton{Close} des fenêtres qui modifient quelque chose} (détails de la mission, ligne de temps, réglage du décalage, réglages de l'étape~4). Fermer garde ce qui a été changé : dès qu'une valeur change, \bouton{Close} devient \bouton{Save and close} ; fermer la fenêtre par sa barre de titre ou avec \textsf{Échap} fait de même. Le bouton \bouton{↺}, à côté de \bouton{?}, n'apparaît que s'il y a quelque chose à remettre : il remet les valeurs et la fenêtre reste ouverte.
\end{itemize}
Le manuel s'ouvre dans le lecteur PDF par défaut (navigateur, Acrobat), à la bonne page quand le lecteur le permet.

\section{Organisation numérique d'une mission IRDrone}

%\todo{Dossier \texttt{FLY-YYYYMMDD-HHmm-localit\'e}, sous-dossiers \texttt{ImgIRdrone}, \texttt{mapping\_MULTI}, etc.}
On suppose ici que votre ordinateur sous Windows comporte un dossier \texttt{C:/Air-Mission}.
Ce dossier de travail est choisi au premier lancement de l'application ; il est rappel\'e en bas de la fen\^etre
principale et peut \^etre chang\'e avec le bouton \og engrenage \fg{} plac\'e \`a c\^ot\'e.

Chaque mission est placée dans un dossier \texttt{C:/Air-Mission/FLY-YYYYMMDD-hhmm-location}. 

Ce dossier comporte les sous-dossiers :
\begin{enumerate}
	\item \texttt{AerialPhotography}
	      \begin{itemize}
	      	\item \texttt{VIS}
	      	  \begin{itemize}
	      	  	\item \texttt{VIS\_XXXX.dng}. Images de la caméra DJI transférées,
	      	  	\item \texttt{VIS\_XXXX.exif}. Fichiers \og compagnons \fg~des images contenant des données exif,
	      	  	\item \texttt{time\_line.json}. Fichier qui donne la chronologie des images par rapport à l'horloge de la caméra VIS,
	      	  	\item \texttt{transfer\_info\_VIS\_dng.json}. Fichier qui résume les choix de l'étape~2,
	      	  	\item \texttt{VIS\_XXXX\_RawTherapee.tif}. Images\footnote{\label{fn:rawtherapee_1}Ces fichiers très volumineux sont indispensables au traitement numérique des images. Ils sont générés une seule fois à la première utilisation de l'étape~4. Lorsque la mission est entièrement traitée et doit être archivée, ils peuvent être supprimés.} créées à l'étape~4.
	      	  \end{itemize}
	      		
	      	\item \texttt{NIR}
	      	\begin{itemize}
	      		\item \texttt{NIR\_XXXX.dng}. Images de la caméra infrarouge transférées et converties (RAW propriétaire vers dng), 
	      		\item \texttt{NIR\_XXXX.exif}. Fichiers \og compagnons \fg~des images contenant des données exif,   
	      		\item \texttt{time\_line.json}. Fichier qui donne la chronologie des images par rapport à l'horloge de la caméra NIR,
	      		\item \texttt{transfer\_info\_NIR\_dng.json}. Fichier qui résume les choix de l'étape~2,
	      		\item \texttt{NIR\_XXXX\_RawTherapee.tif}. Images\textsuperscript{\ref{fn:rawtherapee_1}} créées à l'étape~4.
	      		\end{itemize}
      		\item \texttt{shoot\_points.json}. Dictionnaire des \emph{shoot points} (prise de vue avec comme référence l'image VIS et l'image NIR la plus proche de l'instant de la prise de vue VIS). Il est créé à l'étape~3.
	    \end{itemize}
	\item \texttt{cameras}
	         \begin{itemize}
	         	\item Non utilisé actuellement.
	         \end{itemize}
	\item \texttt{FlightAnalytics}. C'est un dossier technique.
			\begin{itemize}
				\item \texttt{mission\_parameters.json}. Fichier contenant les paramètres de la mission (modifiable avec la fenêtre \og Mission details \fg, §\ref{btn:step1_details}).
				\item \texttt{flight\_timeline.json}. Phases du vol confirmées à l'étape~2.
				\item \texttt{fly\_vertical\_profil.png}. Graphe qui permet de visualiser les altitudes de prise de vue par rapport au sol.
				\item \texttt{Angles analysis.png},  \texttt{Angles analysis dispersion.png}, \texttt{Angles analysis spectrum.png}. Graphes d'analyse de la qualité du vol. Le premier de ces trois graphes permet par exemple de juger si l'attitude du drone pendant la mission est restée dans les spécifications de vol.
			\end{itemize}
	\item \texttt{ImgIRdrone}
		\begin{itemize}
		\item Pour chaque shoot point de numéro XXXX on trouve : 
			\begin{itemize}
				\item \texttt{VIS\_XXXX\_REDO.bat},
				\item \texttt{VIS\_XXXX.jpg},
				\item \texttt{VIR\_XXXX\_local.jpg},
				\item \texttt{VIS\_XXXX\_motion\_model.json},
				\item \texttt{VIS\_XXXX\_motion\_model.npy}.
			\end{itemize}
		\item \texttt{review.json}. Mesures et points vérifiés de la revue des images alignées (chapitre~\ref{ch:review}).
		\item \texttt{shoot\_points.json}. La liste des shoot points traités avec un indicateur de réussite. Attention : cet indicateur suppose par défaut que l'alignement est correct. L'utilisateur dispose d'une procédure pour examiner les alignements et les déclarer, le cas échéant, défaillants.
		\item \texttt{imgMultispectral}. Dossier contenant les quadruplets d'images multispectrales au format tif. Ici ces images ne sont pas géoréférencées. Pour chaque shoot point XXXX :
			\begin{itemize}
				\item \texttt{IRD\_XXXX\_1.tif}. Image de la couche rouge. 
				\item \texttt{IRD\_XXXX\_2.tif}. Image de la couche verte.
				\item \texttt{IRD\_XXXX\_3.tif}. Image de la couche bleue.
				\item \texttt{IRD\_XXXX\_4.tif}. Image de la couche infrarouge proche.
				\item \texttt{IRD\_XXXX\_1.exif}. Fichier \og compagnon \fg~contenant des données exif\footnote{Les données exif sont aussi présentes dans les images \texttt{IRD\_XXXX\_Y.tif}. Toutefois l'extraction de certaines de ces données exif nécessite l'utilisation d'un logiciel tiers spécifique comme ExifTool. C'est une opération coûteuse qu'on ne réalise qu'une fois. On peut donc considérer les fichiers exif compagnons comme des données qui ont été mises en cache.}.
			\end{itemize}
		\item \texttt{VIS\_XXXX\_align\_traces\_manual}. Dossier contenant toutes les informations et images suite à une opération d'alignement manuel des images VIS et NIR du shoot point XXXX.
		\end{itemize}
	\item \texttt{mapping\_MULTI}. Dossier de travail de l'étape~6 (mapping).
	\item \texttt{Synchro}. Aperçus des images et détections de la mire préparés pour l'étape~2 (\texttt{Thumbnails}, un cache recréé si besoin) ; fichiers de l'étape~3 : images de travail \texttt{*\_RawTherapee.jpg} (sous-dossiers \texttt{VIS} et \texttt{NIR}), images de contrôle de la détection des mires (\texttt{Check\_ARUCO}) et graphiques \texttt{check\_time\_*.png}.
\end{enumerate}

\section{Encha\^inement des \'etapes}
%\todo{Ordre conseill\'e, pr\'erequis de chaque \'etape.}
Vous devez avoir sur un support numérique (carte SD, disque dur) les images des deux caméras prises lors du vol de la mission. Il est fortement conseillé de n'avoir qu'une mission par carte SD. Les images du DJI sont au format dng et stockées dans un dossier \texttt{DCIM/HYPERLAPSE}. Les images de la caméra SJCam M20 sont au format RAW, doublées d'une image jpeg. Comme les images RAW ne possèdent pas de vignette, on peut les choisir en regardant l'image jpeg associée. Les images sont stockées dans un dossier \texttt{DCIM}.

Pour une nouvelle mission il faut suivre dans l'ordre les étapes 1 à 6. 
Cependant les étapes 2 à 6 peuvent être lancées individuellement, à condition que le dossier de la mission comporte au minimum les éléments produits par les étapes qui la précèdent. C'est particulièrement utile pour les étapes 4 à 6.

Les étapes se lancent depuis la fenêtre principale, décrite au chapitre suivant.

%==========================================================
\chapter{La fenêtre principale}\label{ch:home}
La fenêtre principale a deux pages : la liste des missions, puis la mission ouverte. En bas, une barre rappelle le dossier de travail et ouvre ce manuel.

\section{Liste des missions}\label{btn:home_missions}
Sans mission ouverte, la fenêtre propose (figure~\ref{fig:home_welcome}) :
\begin{description}
	\item[\bouton{New mission}] crée une mission à partir des cartes SD du drone et de la caméra infrarouge (étape~1, chapitre~\ref{ch:step1}).
	\item[\bouton{Open a mission}] ouvre une mission déjà créée : choisir son dossier \texttt{FLY-YYYYMMDD-hhmm-location}.
	\item[Missions récentes] les missions du dossier de travail, les plus récentes en premier, avec leur avancement : un point vert par étape faite. Un clic ouvre la mission.
\end{description}
Le bouton \bouton{All missions} de la page de la mission revient à cette liste.

\figureIRD[fig:home_welcome]{home_welcome.png}{La liste des missions : nouvelle mission, ouvrir une mission, missions récentes et leur avancement.}{0.85}

\section{La mission ouverte}\label{btn:home_steps}
La page de la mission (figure~\ref{fig:home_mission}) montre :
\begin{itemize}
	\item en haut, le nom de la mission (un clic ouvre son dossier dans l'explorateur) et le lieu, la date et le pilote ;
	\item les sept étapes, de gauche à droite : coche verte et date pour une étape faite, cadre bleu pour l'étape suivante, gris pour les autres. \textbf{Un clic sur une étape la lance}, même déjà faite (pour la refaire). Deux exceptions :
	\begin{itemize}
		\item l'étape~1 faite ouvre seulement la fenêtre \og Mission details \fg~(§\ref{btn:step1_details}) : les cartes SD ne sont pas relues ;
		\item l'étape~3 faite montre le résultat enregistré (décalage des horloges, courbes alignées) sans rien recalculer (§\ref{sec:step3_chained}).
	\end{itemize}
	La touche \textsf{F1} sur une étape ouvre son chapitre de ce manuel ;
	\item l'étape suivante, expliquée en une phrase, avec son bouton (§\ref{btn:home_continue}) ;
	\item quelques chiffres : nombre d'images, de paires, d'images alignées, de points signalés mal alignés ;
	\item à droite, la carte du vol ou une image alignée (§\ref{btn:home_map}).
\end{itemize}

\figureIRD[fig:home_mission]{home_mission.png}{La mission ouverte : les étapes (un clic lance une étape), étape suivante, chiffres de la mission et carte du vol sur l'orthophoto et les photos aériennes IGN.}{0.9}

\section{Bouton \bouton{Continue with Step N}}\label{btn:home_continue}
Lance l'étape suivante. Les étapes 1 à 4 s'enchaînent d'elles-mêmes :
\begin{enumerate}
	\item à la fin de l'étape~1, \bouton{Flight timeline} ouvre l'étape~2 ;
	\item \bouton{Continue with synchronisation} de l'étape~2 (ou \bouton{Continue with Step 3}) ouvre l'étape~3 ; \bouton{Synchronise the cameras} la lance : détection des mires, décalage des horloges, puis une question, \og Do the two curves overlap? \fg~; \bouton{Yes} apparie les images ;
	\item \bouton{Continue with image processing} ouvre l'étape~4. L'alignement, long, ne démarre jamais tout seul : il se lance avec \bouton{Run alignment} et s'arrête à tout moment avec \bouton{Stop} ;
	\item à la fin, \bouton{Continue with the review} ouvre l'étape~5, la revue des images alignées (elle s'ouvre d'elle-même quand l'alignement a été enchaîné depuis l'étape~3) ;
	\item \bouton{Continue with mapping} de la revue ouvre l'étape~6, où \bouton{Build the map\ldots} lance la carte ;
	\item une fois la carte construite, \bouton{Continue with archiving} ouvre l'étape~7, qui indique l'espace libéré par l'archivage ; rien n'est supprimé avant votre confirmation.
\end{enumerate}
Il reste donc à vérifier les phases du vol, la superposition des courbes de synchronisation, puis le résultat. Une fois toutes les étapes faites, le bouton devient \bouton{Archive the mission} (étape~7).

\section{Carte du vol}\label{btn:home_map}
La carte est disponible dès l'étape~1 : les positions sont lues dans les EXIF des images visibles. Elle montre :
\begin{itemize}
	\item la trajectoire du drone et un point par image, coloré selon les phases confirmées à l'étape~2 (violet : synchronisation, vert : relevé) ;
	\item les images de la carte, et elles seules : la trajectoire les relie, le vol complet restant en pointillés discrets, et tant qu'aucune orthophoto n'est affichée leur emprise au sol est dessinée en vert : on voit le recouvrement de la carte. Quelles images :
	\begin{itemize}
		\item avant l'étape~4, les images retenues par le \emph{best mapping} (étapes~2 et~4, section~\ref{btn:step2_best_mapping}), s'il est choisi ;
		\item une fois l'étape~4 faite, les points exactement comme dans la revue des images alignées (chapitre~\ref{ch:review}) : chaque point de prise de vue dans la couleur de son statut (vert aligné, orange à vérifier, rouge signalé\ldots), les points que l'étape~4 n'a pas alignés en petits points gris ; une image alignée ou réalignée depuis la revue apparaît aussitôt ;
		\item l'orthophoto de l'étape~6, une fois construite, est dessinée sous ces points (les mêmes points que la revue, pas seulement les images qu'elle a utilisées) ;
	\end{itemize}
	les images signalées (\bouton{Flag}, revue des images alignées, chapitre~\ref{ch:review}) sont de petits points rouges ; en survoler un montre son emprise en pointillés rouges : on voit le trou qu'il laisse dans le recouvrement ; signalées après l'étape~6, elles sont encore dans l'orthophoto affichée et en seront retirées quand l'étape~6 sera relancée (\texttt{left out when Step 6 runs again}). Le nombre d'images et d'images signalées est écrit en bas à droite. \bouton{Layers} > \bouton{Only the images of the map (and the flagged ones)} (\bouton{Only the images of the best mapping} avant l'étape~4) remontre toutes les images ;
	\item l'orthophoto de l'étape~6, une fois le mapping fait : dès que l'étape~6 l'a construite, la fenêtre principale l'affiche sur la carte (même si \bouton{Layers} > \bouton{Hidden} était choisi avant) ;
	\item en fond, des photos aériennes : IGN pour la France, Esri World Imagery ailleurs, ou OpenStreetMap.
\end{itemize}
Survoler un point montre la vignette de son image ; un clic sur un point montre l'emprise au sol de l'image (le trait épais est le haut de l'image). Molette : zoom ; glisser : déplacement ; double-clic ou bouton \bouton{⤢} : tout le vol.

Le menu \bouton{Layers} choisit le fond (IGN, Esri, OpenStreetMap ou aucun) et l'affichage de l'orthophoto, en couleurs naturelles ou en fausses couleurs infrarouges. Le fond vient d'internet : hors ligne, la carte reste utilisable (vol, emprises, orthophoto), seul le fond manque, et rien n'est bloqué. Les tuiles déjà vues sont gardées dans un cache (\texttt{\textasciitilde/.cache/irdrone} sous Linux, \texttt{\%LOCALAPPDATA\%\textbackslash IRDrone} sous Windows) et s'affichent ensuite hors ligne.

L'emprise est calculée avec la hauteur au-dessus du sol (modèle de terrain de l'étape~6, ou altitudes IGN) et le cap de l'image ; une fois le mapping fait, le cap vient d'OpenDroneMap, plus juste que la boussole du drone (souvent fausse de 5 à 10\textdegree).


\section{Barre du bas : dossier de travail et manuel}\label{btn:home_workdir}
Sur la page de la liste des missions, la barre du bas rappelle le dossier de travail (le dossier qui contient les missions, par défaut \texttt{C:/Air-Mission}) et la place libre sur son disque ; un clic sur le dossier l'ouvre dans l'explorateur. La roue dentée \bouton{⚙} permet de changer de dossier de travail : la liste des missions est alors relue. Elle donne aussi une option avancée, décochée par défaut : saisir à la main le décalage d'horloge des caméras à l'étape~3, pour les vols sans cible ArUco (§\ref{btn:step3_clock_offset}). Le dossier est choisi au premier lancement ; la variable d'environnement \texttt{IRDRONE\_WORKING\_DIR}, si elle est définie, reste prioritaire. Sur la page d'une mission, la barre du bas indique à droite la taille totale du dossier de la mission (\texttt{Mission size}), mise à jour à chaque étape terminée et toutes les minutes.

\bouton{?}, à gauche, ouvre ce manuel (§\ref{sec:help}).

\section{Rapport de diagnostic}\label{btn:home_diagnostic}
Quand quelque chose ne marche pas sur un ordinateur, le rapport de diagnostic rassemble dans un seul fichier zip tout ce qu'il faut pour comprendre le problème ; il suffit d'envoyer ce zip aux développeurs. Il se crée de deux façons :
\begin{itemize}
  \item dans les options (roue dentée \bouton{⚙} en bas de la fenêtre principale, §\ref{btn:home_workdir}), bouton \bouton{Create a diagnostic report…} : on choisit une mission, ou l'installation seule ;
  \item sans ouvrir le logiciel, avec \texttt{irdrone\_diagnostic.bat} (à la racine du dossier IRDrone) : double-clic pour l'installation seule, ou glisser-déposer le dossier d'une mission (\texttt{FLY-...}) sur ce fichier.
\end{itemize}
Le zip (\texttt{irdrone\_diagnostic\_AAAAMMJJ\_hhmmss.zip}) est écrit dans le dossier de travail. Il contient :
\begin{itemize}
  \item \texttt{report.txt} : un résumé de ce qui semble anormal, puis l'installation (version d'IRDrone, Python utilisé, système, page de code Windows, versions des modules Python dont OpenCV, chemins et versions d'ExifTool, RawTherapee et du convertisseur SJCam, dossier de travail, place libre) et deux auto-tests (détection ArUco d'une cible parfaite ; lecture et écriture d'images et d'EXIF dans un dossier accentué) ;
  \item pour une mission : les étapes faites, les fenêtres de l'étape~2, les lignes de temps, le décalage et le nombre de paires de l'étape~3 ; chaque image de synchronisation est développée à nouveau avec RawTherapee dans un dossier temporaire et les marqueurs ArUco y sont cherchés, ainsi que sur l'image en cache (\texttt{Synchro}) utilisée par l'étape~3. Quelques images développées sont jointes. Cela prend environ une minute ; la mission n'est pas modifiée ;
  \item \texttt{report.json} (tous les détails) et les journaux des dernières sessions : chaque lancement du logiciel recopie la console dans \texttt{logs\textbackslash irdrone\_AAAAMMJJ\_hhmmss.log}, à côté du fichier des réglages (les dix derniers sont gardés).
\end{itemize}
Rien n'est envoyé automatiquement.

%==========================================================
\chapter{Step 1 -- D\'efinir la mission}\label{ch:step1}
\section{La fenêtre \og New mission \fg}\label{sec:step1_new_mission}
Le Step~1 cr\'ee le dossier de la mission et y range les images du vol. Il s'ouvre avec \bouton{New mission} dans la fenêtre principale et se fait dans une seule fenêtre, en trois pages : \textbf{Images}, \textbf{Mission}, \textbf{Copy}. L'avancement est rappelé en haut de la fenêtre. Les cartes SD sont seulement lues : rien n'y est écrit ni effacé.

\subsection{Type de mission : multispectrale ou visible seul}\label{btn:step1_mission_type}
Le haut de la page \og Images \fg~choisit le type de la mission (le dernier choisi est proposé) :
\begin{itemize}
	\item \bouton{Multispectral (visible + infrared)} : le drone (visible, DNG) et la caméra infrarouge (SJCam, RAW). Tout ce que décrit ce manuel : synchronisation, alignement, carte en quatre bandes (rouge, vert, bleu, proche infrarouge), VIR et NDVI ;
	\item \bouton{Visible only (drone DNG)} : les images du drone seules, sans la caméra infrarouge. La carte infrarouge disparaît, la copie ne prend que les images visibles et les étapes suivantes s'adaptent (section~\ref{sec:visible_only}) : phase de relevé seule à l'étape~2, pas d'étape~3, les images sont \emph{développées} à l'étape~4 au lieu d'être alignées, la carte de l'étape~6 est en couleurs (rouge, vert, bleu).
\end{itemize}
Le type ne se change plus une fois la mission créée.

\subsection{Page \og Images \fg}\label{btn:step1_images}
Les cartes SD insérées sont parcourues (tout volume qui contient un dossier \texttt{DCIM}) ; on y trouve les images visibles (DNG du drone) et infrarouges (RAW de la SJCam), figure~\ref{fig:step1_new_mission_images}. Chaque caméra a sa carte :
\begin{itemize}
	\item le dossier trouvé, le nombre d'images, leur taille, la date et l'heure de la première et de la dernière ;
	\item deux images : \textbf{au décollage} et \textbf{en vol} (aux deux tiers des images, en général pendant le relevé). Pour l'infrarouge, c'est le JPEG que la SJCam enregistre à côté de chaque RAW. Ces images permettent de vérifier d'un coup d'\oe il que la carte est la bonne ;
	\item le bouton \bouton{Choose folder…} (\bouton{Change…} une fois des images trouvées) choisit un autre dossier, par exemple une copie des cartes sur le disque. Un dossier choisi ainsi n'est jamais remplacé par une carte trouvée ensuite.
\end{itemize}
Tant qu'il manque les images d'une caméra, la fenêtre cherche les cartes SD toutes les 3~secondes : une carte insérée après l'ouverture de la fenêtre est trouvée d'elle-même, il suffit de l'insérer.

L'image de décollage est la première image visible ; le lien \og change… \fg~de la carte visible en choisit une autre. Elle donne la date, l'heure, la position et l'appareil photo de la mission.
Un bandeau signale un écart de plus de deux minutes entre les horloges des deux caméras : il est normal, l'étape~3 le mesure précisément. Un écart de plusieurs jours fait douter que les deux cartes viennent du même vol.

\figureIRD[fig:step1_new_mission_images]{step1_new_mission_images.png}{Page \og Images \fg : images visibles et infrarouges trouvées sur les cartes SD, au décollage et en vol.}{0.8}

\subsection{Page \og Mission \fg}\label{btn:step1_mission}
La date, l'heure et la position viennent de l'image de décollage ; l'altitude du sol et l'adresse sont demandées à l'IGN\footnote{\label{fn:ign}IGN : Institut national de l'information g\'eographique et foresti\`ere.} et à OpenStreetMap pendant que l'on remplit la description et le pilote (figure~\ref{fig:step1_new_mission_details}). Le nom du dossier de la mission (\texttt{FLY-YYYYMMDD-hhmm-location}) se met à jour avec le champ \og Location \fg.

\bouton{Cameras and filter} déplie les réglages des caméras : intervalle de prise de vue et décalage d'horloge de chaque caméra, fabricant et modèle de la caméra infrarouge, fabricant et bande du filtre. Le pilote, son identifiant, les caméras et le filtre sont mémorisés pour la mission suivante.

\figureIRD[fig:step1_new_mission_details]{step1_new_mission_details.png}{Page \og Mission \fg : lieu trouvé à partir de l'image de décollage, nom du dossier, description, pilote, caméras.}{0.8}

\subsection{Page \og Copy \fg}\label{btn:step1_copy}
Le dossier est créé, les images visibles copiées, les images infrarouges converties de RAW en DNG, les EXIF lus et l'altitude du sol calculée (détails au §\ref{btn:step1_validate}). Pendant la conversion des images infrarouges, l'étape~2 est préparée en arrière-plan (vignettes et détection de la mire, visible puis infrarouge) : la ligne de temps s'ouvre ensuite sans attendre. La fenêtre ne peut pas être fermée pendant la copie. À la fin, \bouton{Flight timeline} ouvre l'étape~2 ; si la préparation n'est pas finie, l'étape~2 la reprend là où elle en était.

Le fichier \texttt{mission\_parameters.json} contient les mêmes informations qu'avec le formulaire classique.

\subsection{Boutons du bas de la fenêtre}\label{btn:step1_wizard}
\begin{description}
	\item[\bouton{Next}] (page Images) passe à la page Mission ; il faut les images des deux caméras. \bouton{Create the mission and copy the images} (page Mission) lance la copie ; \bouton{Flight timeline} (page Copy) ouvre l'étape~2.
	\item[\bouton{Back}] revient à la page Images ; rien n'est perdu.
	\item[\bouton{Cancel}] ferme la fenêtre sans créer la mission.
	\item[\bouton{Classic form…}] (page Images) ouvre l'ancien Step~1 (§\ref{sec:step1_classic}).
	\item[\bouton{?}] ouvre ce chapitre.
\end{description}

\section{Modifier une mission : la fenêtre \og Mission details \fg}\label{btn:step1_details}
Une fois l'étape~1 faite, un clic sur l'étape~1 de la fenêtre principale n'ouvre plus la création de mission, mais la fenêtre \og Mission details \fg~(figure~\ref{fig:step1_mission_details}). Elle modifie les informations de la mission sans relire les cartes SD : lieu, description, pilote et son identifiant, intervalle et décalage d'horloge des caméras, caméra infrarouge, filtre. La date, l'heure, la position et la caméra visible viennent de l'image de décollage et ne changent pas.

Dès qu'une information change, \bouton{Close} devient \bouton{Save and close} : il écrit les changements dans \texttt{FlightAnalytics/mission\_parameters.json} (fermer la fenêtre fait de même) ; l'ancien fichier est gardé en \texttt{mission\_parameters.json.bak}. Le dossier de la mission garde son nom, même si le lieu change. \bouton{↺}, à côté de \bouton{?}, apparaît alors : il remet les informations telles qu'à l'ouverture de la fenêtre ; la fenêtre reste ouverte.

\figureIRD[fig:step1_mission_details]{step1_mission_details.png}{La fenêtre \og Mission details \fg : les images sont déjà dans la mission, seules ces informations changent.}{0.55}

\section{Formulaire classique}\label{sec:step1_classic}
\legacybegin{Afficher le formulaire classique (ancienne étape 1)}
Le bouton \bouton{Classic form…} de la page \og Images \fg~ouvre l'ancien Step~1, décrit ci-dessous. Il se d\'eroule en deux fen\^etres successives :
\begin{enumerate}
  \item \textbf{\og Choose the image taken by the drone during takeoff \fg} : on choisit la premi\`ere image prise par le drone
        au d\'ecollage. Elle fournit la date, l'heure, la position GPS et l'appareil photo utilis\'e.
  \item \textbf{\og Create a mission \fg} : on v\'erifie et on compl\`ete les informations de la mission, puis on lance la cr\'eation
        du dossier et le transfert des images VIS et NIR.
\end{enumerate}
\`A la fin, la mission devient la mission chargée de la fen\^etre principale, qui propose l'étape~2.

\paragraph{Pr\'erequis.}
\begin{itemize}
  \item Une connexion Internet : l'altitude du sol est demand\'ee \`a l'IGN\textsuperscript{\ref{fn:ign}}
  (France) ou \`a OpenTopoData (hors de France),
        et l'adresse du lieu \`a OpenStreetMap.
  \item Les images du drone (visible) au format DNG, et les images de la cam\'era proche infrarouge (SJCam M20, fichiers RAW)
        copi\'ees sur le disque de l'ordinateur.
  \item Un dossier de travail (par d\'efaut \texttt{C:/Air-Mission}), choisi au premier lancement et modifiable
        avec le bouton \og engrenage \fg{} en bas de la fen\^etre principale. La variable d'environnement
        \texttt{IRDRONE\_WORKING\_DIR}, si elle est d\'efinie, reste prioritaire.
\end{itemize}

% \figureIRD{step1_load_takeoff_1.png}{Premi\`ere fen\^etre du Step 1 avant chargement de l'image de d\'ecollage.}{0.8}

%----------------------------------------------------------
\section{Bouton \bouton{Load take-off image}}\label{btn:step1_load_takeoff}
Ouvre un explorateur de fichiers pour choisir l'image prise par le drone au d\'ecollage.
Seuls les fichiers \texttt{.dng} de la cam\'era visible sont accept\'es.

Une fois l'image choisie, le logiciel :
\figureIRD{step1_load_takeoff_2.png}{Premi\`ere fen\^etre du Step 1 apr\`es chargement de l'image de d\'ecollage.}{0.5}
\begin{enumerate}
  \item affiche l'image (le d\'ecodage du fichier DNG est assez long : la barre de progression indique l'avancement) ;
  \item lit les donn\'ees EXIF : date et heure de prise de vue, position GPS, fabricant, mod\`ele et num\'ero de s\'erie de l'appareil ;
  \item interroge l'IGN (ou OpenTopoData hors de France) pour obtenir l'altitude du sol \emph{au-dessus du niveau de la mer} ;
  \item interroge OpenStreetMap pour obtenir le lieu-dit, la commune, le code postal, le d\'epartement, la r\'egion et le pays.
\end{enumerate}
Ces informations s'affichent dans le cadre gris sous l'image. Lorsqu'elles existent, le lieu-dit, le code postal, la région et le pays sont aussi enregistrés dans le fichier
\texttt{FlightAnalytics\textbackslash mission\_parameters.json} sous les clés \texttt{Place name}, \texttt{Postal code}, \texttt{Region} et \texttt{Country}
(la commune est enregistrée sous la clé \texttt{Location}). Une information introuvable n'est pas enregistrée : la clé correspondante n'existe pas.
\begin{center}
\fbox{\parbox{0.9\linewidth}{\small\textbf{Attention.} L'altitude enregistr\'ee dans l'image par le drone est mesur\'ee \emph{par rapport
au point de d\'ecollage}, et non par rapport \`a la mer. C'est pourquoi l'altitude du sol est demand\'ee \`a un service externe tel que l'IGN\textsuperscript{\ref{fn:ign}}.}}
\end{center}


\paragraph{Si l'image ne se charge pas.} V\'erifiez la connexion Internet et que le fichier est bien un \texttt{.dng} du drone.
Le bouton \bouton{Next ▶} reste inactif tant qu'aucune image n'a \'et\'e charg\'ee correctement.

%----------------------------------------------------------
\section{Boutons \bouton{Cancel}, \bouton{Next ▶} et \bouton{◀ Back}}\label{btn:step1_navigation}
\begin{description}
  \item[\bouton{Next ▶}] n'est disponible qu'apr\`es le chargement de l'image de d\'ecollage. Il ouvre la fen\^etre \og Create a mission \fg.
  \item[\bouton{Cancel}] ferme la fen\^etre de l'image de d\'ecollage sans rien enregistrer.
  \item[\bouton{◀ Back}] (fen\^etre \og Create a mission \fg) revient \`a la fen\^etre pr\'ec\'edente ;
        l'image de d\'ecollage d\'ej\`a charg\'ee est conserv\'ee.
\end{description}

%----------------------------------------------------------
\section{La fen\^etre \og Create a mission \fg}\label{btn:step1_form}
Les champs sont remplis automatiquement \`a partir de l'image de d\'ecollage. La touche \textsf{Entr\'ee} passe d'un champ au suivant.
Les champs gris sont remplis automatiquement et ne sont pas modifiables.

\begin{center}
\begin{tabular}{@{}p{3.6cm}p{2.4cm}p{7.5cm}@{}}
\hline
\textbf{Champ} & \textbf{Modifiable} & \textbf{Contenu} \\ \hline
Date, Hour & non & Date et heure du d\'ecollage (EXIF). \\
Location & oui & Commune trouv\'ee automatiquement. Fait partie du nom du dossier de la mission. \\
GPS & non & Latitude et longitude du point de d\'ecollage. \\
Short description & oui & Br\`eve description libre de la mission. \\
Pilot & oui & Nom du t\'el\'epilote et son num\'ero d'identification. \\
Camera VIS & en partie & Fabricant et mod\`ele (lus dans l'image, non modifiables) ; intervalle de prise de vue (timelapse, en s) et d\'ecalage d'horloge (delta time, en s). \\
Image VIS format & non & Format des images (DNG). \\
Camera NIR & oui & Par d\'efaut SJCam M20, format RAW, intervalle 3\,s, d\'ecalage d'horloge $-30$\,s. \\
NIR Filter & oui & Fabricant du filtre et bande, en nm (par d\'efaut 810\,nm). \\ \hline
\end{tabular}
\end{center}

\figureIRD{step1_param_mission.png}{Fen\^etre \og Create a mission \fg.}{0.5}

Remarque: les paramètres de décalage d'horloge VIS et NIR sont réservés à un utilisateur expérimenté. Ce sont des paramètres qui permettent de \og récupérer\fg ~une mission dont la phase de synchronisation  sur le terrain se révèle inexploitable (impossibilité de reconnaître les ArUco, oubli de déclenchement de la mire, etc.).

%----------------------------------------------------------
\section{Bouton \bouton{Create the mission}}\label{btn:step1_validate}
Lance la cr\'eation de la mission. Le bouton se d\'esactive pendant le traitement ; une \'etiquette et une barre de progression
indiquent la phase en cours. Le logiciel encha\^ine, sans intervention de l'utilisateur sauf \`a l'\'etape~3 :
\begin{enumerate}
  \item \textbf{Cr\'eation du dossier} de mission et de ses sous-dossiers, \'ecriture de \texttt{mission\_parameters.json},
        copie de l'image de d\'ecollage.
  \item \textbf{Transfert des images VIS} : les fichiers \texttt{HYPERLAPSE\_XXXX.dng} du dossier de l'image de d\'ecollage sont copi\'es
        dans \texttt{AerialPhotography/VIS} sous le nom \texttt{VIS\_XXXX.dng}.
  \item \textbf{Choix du dossier NIR} : un message demande de choisir la premi\`ere image proche infrarouge (fichier RAW ou JPG
  \footnote{Chaque image RAW a une image jpeg associée. La SJCam nomme les images par leur date de prise de vue. Par exemple 2022\_0125\_130413\_001.RAW, 2022\_0125\_130415\_002.JPG, 2022\_0125\_130416\_003.RAW, 2022\_0125\_130418\_004.JPG, etc. L'image jpeg est toujours enregistrée après l'image RAW environ 2s après. À noter que dans certaines circonstances (carte SD trop lente, longue séquence de vol) le délai d'enregistrement du fichier RAW dépasse la fréquence de prise de vue (3s dans notre cas). Dans ce cas l'image suivante est prise 6s après. IRdrone gère automatiquement ce type de saut.} de la SJCam).
        Le dossier qui la contient est utilis\'e pour toutes les images NIR.
  \item \textbf{Conversion des images NIR} : les fichiers RAW sont convertis en DNG dans \texttt{AerialPhotography/NIR}
        (calcul en parall\`ele, \'eventuellement long).
  \item \textbf{Fichiers \texttt{.exif}} des images VIS puis NIR : un fichier de m\'etadonn\'ees par image,
        et un fichier \texttt{time\_line.json} par bande donnant la chronologie des prises de vue.
  \item \textbf{Altitudes et trajectoire} : pour chaque position GPS du vol, l'altitude du sol est demand\'ee \`a l'IGN (par paquets ;
        trois essais en cas d'\'echec, puis bascule automatique sur OpenTopoData). Le logiciel en d\'eduit la hauteur du drone au-dessus du sol,
        les distances et les caps entre prises de vue, et enregistre ces valeurs dans les fichiers \texttt{.exif}.
  \item \textbf{Profil vertical} : trac\'e de l'altitude du drone et du sol le long du trajet, enregistr\'e dans
        \texttt{FlightAnalytics/fly\_vertical\_profil.png}.
  \item \textbf{Message final} indiquant le dossier cr\'e\'e, puis fermeture des fen\^etres du Step~1.
\end{enumerate}

\figureIRD[fig:takeoff_1]{step1_nir_prompt.png}{Message demandant de choisir la premi\`ere image NIR.}{0.2}
\figureIRD[fig:fly_vertical_profil]{step1_fly_vertical_profil.png}{Profil vertical du vol (altitude du drone et du sol), produit \`a la fin du Step~1.}{0.8}

\subsection{Dossiers cr\'e\'es}
À l'issue de l'étape 1 (Step 1) IRDrone aura cr\'e\'e la structure de dossiers et fichiers de la mission. À ce stade vous n'aurez plus besoin d'accéder aux images originales stockées sur les cartes SD des caméras. La structure du dossier de la mission est la suivante :


\begin{verbatim}
FLY-YYYYMMDD-hhmm-<location>/
  AerialPhotography/
     VIS/   VIS_XXXX.dng, VIS_XXXX.exif, time_line.json,
            transfer_info_VIS_dng.json
     NIR/   NIR_XXXX.dng, NIR_XXXX.exif, time_line.json,
            transfer_info_NIR_dng.json
  FlightAnalytics/
     mission_parameters.json, fly_vertical_profil.png
  Synchro/        (Step 3)
  ImgIRdrone/     (Step 4)
  mapping_MULTI/  (Step 6)
  cameras/
\end{verbatim}
Les sous-dossiers ont une ic\^one sp\'ecifique dans l'explorateur Windows.
Le dossier de mission n'est jamais effac\'e : si le dossier existe d\'ej\`a, seuls les \'el\'ements manquants sont cr\'e\'es.
Pour les images VIS d\'ej\`a transf\'er\'ees (fichier \texttt{.exif} pr\'esent), le transfert est ignor\'e : relancer le Step~1 sur
une mission existante est donc rapide.

\figureIRD{Step1_mission_folder.png}{Dossier de mission vu dans l'explorateur Windows.}{0.7}

La durée des transferts dépend beaucoup des performances des supports (carte SD, disque dur) et des liaisons (USB B ou C, lecteur). De plus et contrairement aux images de la caméra DJI qui sont au format dng, les images de la caméra SJCam M20 sont dans un format \og RAW \fg~qui n'est pas normalisé. Chaque image doit être convertie au format dng d'ADOBE à l'aide d'un logiciel tiers \footnote{Il s'agit du code \texttt{sjcam\_raw2dng} disponible uniquement comme un fichier exe. \texttt{sjcam\_raw2dng}
	Convert SJCAM raw files to Adobe DNG and TIFF while also getting metadata
	such as date and ISO from JPEGs that correlate to the RAW files
} ce qui allonge le temps de transfert effectif. 

\subsection{En cas de probl\`eme}
\begin{itemize}
  \item \textbf{Aucune image NIR trouv\'ee} : le message \og No NIR RAW files found \fg\ signifie que le dossier choisi ne contient pas de
        fichiers \texttt{.RAW} ; recommencez en choisissant le bon dossier.
  \item \textbf{Erreur pendant le transfert VIS ou la lecture EXIF} : un message d'erreur s'affiche et la fen\^etre se ferme.
        Le bouton reste utilisable apr\`es r\'eouverture de la fen\^etre.
  \item \textbf{Altitudes IGN indisponibles} : le logiciel bascule seul sur OpenTopoData, moins pr\'ecis en France.
\end{itemize}
%\todo{Compl\'eter avec les erreurs rencontr\'ees sur le terrain (carte SD absente, dossier prot\'eg\'e en \'ecriture, \dots).}

%==========================================================
\legacyend

\chapter{Step 2 -- Ligne de temps du vol}\label{ch:step2}
\section{Principe}
L'étape 2 découpe le vol en phases : \emph{synchronisation} (vol stationnaire au-dessus de la mire ArUco tournante) et \emph{survey} (vol au-dessus de la zone étudiée). Ces phases sont maintenant \textbf{détectées automatiquement} ; l'utilisateur vérifie et ajuste si besoin (figure~\ref{fig:step2_flight_timeline}).

\figureIRD[fig:step2_flight_timeline]{step2_flight_timeline.png}{La fenêtre \og Flight timeline \fg. Au centre l'aperçu des images visible et infrarouge sous la souris, à droite la carte du vol avec les empreintes au sol des images gardées pour la carte (§\ref{btn:step2_best_mapping}), puis les quatre points clés avec leurs vraies images ; en bas la ligne de temps : profil d'altitude (points verts : images gardées pour la carte), phases (violet : synchronisation, vert : survey), détections de la mire et bandes d'images visible et infrarouge.}{0.95}

\subsection*{Ce que fait le logiciel}
\begin{enumerate}
	\item Chaque fichier DNG (DJI comme SJCam) contient un aperçu JPEG d'environ 1000 pixels. Il se lit en quelques millisecondes : toute la mission est parcourue en quelques secondes, sans RawTherapee. Les aperçus sont gardés dans \texttt{Synchro/Thumbnails} (cache, recréé si besoin).
	\item \emph{Visible} : l'altitude enregistrée par le drone donne le vol stationnaire du début (synchronisation) et le palier le plus long, le plus éloigné de l'altitude du décollage (survey). Un décollage depuis un sommet, suivi d'une descente vers la zone étudiée, est traité de la même façon qu'une montée.
	\item \emph{Synchronisation} : l'angle de la mire mobile (identifiant 0) est mesuré sur les aperçus, seulement pendant le stationnaire pour le visible, et pour l'infrarouge depuis la première image jusqu'à ce que la mire disparaisse (le drone s'est éloigné). Les images où la mire tourne régulièrement dans un sens forment la fenêtre de synchronisation ; les images où le plateau est encore arrêté, où il change de sens, et le décollage sont écartés.
	\item \emph{Infrarouge} : l'horloge de la caméra NIR est souvent fausse (fuseau horaire, dérive). La même rotation de la mire étant vue par les deux caméras, la comparaison des deux courbes d'angle donne un décalage approché (de l'ordre de la seconde) qui place les images infrarouges sur la ligne de temps du visible. La fenêtre de synchronisation infrarouge correspond ainsi au même tour de mire que la fenêtre visible.
	\item La fenêtre survey infrarouge \textbf{suit} la fenêtre survey visible (zone hachurée) : il n'est pas nécessaire de regarder les images infrarouges. Après l'étape 3, le décalage précis met à jour cette fenêtre.
\end{enumerate}

\subsection*{Vérifier et ajuster}\label{btn:step2_adjust}
\begin{itemize}
	\item Survoler la ligne de temps affiche en grand l'image visible et l'image infrarouge prises au même instant, avec l'altitude et l'état de la mire (\og target turning \fg, \og target still \fg, \og no target \fg). Laisser la souris sur cet aperçu l'affiche sur toute la fenêtre pour bien voir les détails.
	\item Les points de la ligne \og target \fg~indiquent les images analysées : violet, la mire tourne ; gris, la mire est vue mais immobile ; cercle vide, mire non vue.
	\item Faire glisser un bord de phase pour le déplacer (il se cale sur une image), ou cliquer sur un point clé (colonne de droite) puis utiliser \bouton{◀} \bouton{▶} ou les flèches du clavier.
	\item Molette : zoom ; glisser le fond : déplacement ; double-clic : tout le vol.
	\item \bouton{↺}, à côté de \bouton{?}, apparaît dès qu'une phase est déplacée : il annule les changements faits à la main et revient à la proposition automatique ; la fenêtre reste ouverte.
	\item L'info-bulle du bouton \bouton{?} rappelle ces commandes.
\end{itemize}
Si la mire n'est pas trouvée (mire trop petite, image trop sombre), un bandeau orange l'indique : la fenêtre de synchronisation visible est alors le vol stationnaire, et la fenêtre infrarouge doit être placée à la main.

\section{Best mapping : les images utiles à la carte}\label{btn:step2_best_mapping}
Un survol photographié toutes les deux secondes donne beaucoup plus de recouvrement que la carte n'en demande. Aligner toutes ces images à l'étape~4 prend des heures, et donner à ODM des centaines d'images presque identiques ralentit le calcul de la carte sans l'améliorer, souvent au contraire. La case \bouton{Best mapping: keep only the images the map needs}, cochée par défaut, ne garde que les images du survol dont les \emph{empreintes au sol} se recouvrent de la valeur choisie à droite (\bouton{75~\% overlap} par défaut, la valeur recommandée pour ODM ; plus de recouvrement garde plus d'images).

Le choix est visible tout de suite :
\begin{itemize}
	\item sur la ligne de temps, des points verts sur le profil d'altitude et des repères verts en haut de la bande d'images visibles marquent les images gardées ; la bulle du curseur et la légende de l'aperçu indiquent \og kept for the map \fg ;
	\item à droite des aperçus, la \textbf{carte du vol} montre le trajet (violet : synchronisation, vert : survey), les empreintes au sol des images gardées (rectangles verts : leur recouvrement se voit) et, en orange, l'empreinte de l'image sous la souris avec l'image elle-même. Un clic sur un point de la carte affiche cette image ; molette : zoom, glisser : déplacement, \bouton{Layers} : fond de carte (§\ref{btn:home_map}) ;
	\item sous la ligne de temps, le nombre d'images gardées, par exemple \og 16 of 120 survey images kept for the map · 87~\% less to align \fg.
\end{itemize}
La sélection suit les bords de la phase survey quand on les déplace.

\textbf{Comment les images sont choisies.} La position, la hauteur au-dessus du sol et l'orientation de chaque image sont lues dans les fichiers \texttt{.exif} écrits à l'étape~1 ; le champ de vue de la caméra donne la taille de l'empreinte au sol. La première image du survol est gardée. Pour chaque image suivante, IRDrone estime le recouvrement de son empreinte avec celle de la dernière image gardée, d'après le déplacement du drone (dans l'axe de l'image et en travers). Tant qu'il reste au-dessus de la valeur choisie, l'image n'est qu'une candidate ; dès qu'il passe en dessous, la candidate précédente est gardée : deux images gardées successives se recouvrent toujours d'au moins la valeur choisie. La dernière image du survol est gardée aussi. Pour un déplacement dans l'axe de l'image, c'est le critère \texttt{--selection best-mapping} de l'ancienne ligne de commande \texttt{run.py}.

\textbf{Et la synchronisation ?} À l'étape~2 le décalage précis entre les deux caméras n'est pas encore connu : l'aperçu porte sur toutes les images du survol. À l'étape~4, une fois les caméras synchronisées, la même sélection est refaite parmi les seules paires bien synchronisées (§\ref{btn:step4_best_mapping}) ; le nombre final peut donc différer un peu de l'aperçu. Le choix (case et recouvrement) est enregistré par \bouton{Continue with synchronisation} et \bouton{Close} dans \texttt{FlightAnalytics\textbackslash best\_mapping.json}.

Si les images n'ont pas de position GPS, la carte n'est pas affichée et toutes les images du survol sont gardées.

\section{Boutons \bouton{Continue with synchronisation} et \bouton{Close}}\label{btn:step2_confirm}
Les deux boutons enregistrent les phases dans les fichiers \texttt{transfer\_info\_VIS\_dng.json} et \texttt{transfer\_info\_NIR\_dng.json} (lus par les étapes 3, 4 et le mapping) et dans \texttt{FlightAnalytics/flight\_timeline.json}. Rouvrir l'étape 2 retrouve les phases confirmées.
\bouton{Continue with synchronisation} ouvre ensuite l'étape~3 ; rien ne démarre tout seul : la synchronisation se lance depuis sa fenêtre (chapitre~\ref{ch:step3}). \bouton{Close} ferme simplement la fenêtre : l'étape~3 se lancera plus tard depuis la fenêtre principale. Il devient \bouton{Save and close} quand les phases ou le choix best mapping ont changé, ou n'ont jamais été enregistrés : fermer la fenêtre les enregistre alors aussi. S'il manque une phase, fermer demande s'il faut fermer sans enregistrer.

\section{Sélection classique}\label{btn:step2_classic}
\legacybegin{Afficher la sélection classique (ancienne fenêtre de l'étape 2)}
Le bouton \bouton{Classic selection…} ouvre l'ancienne sélection manuelle des cinq images de référence de chaque caméra, décrite ci-dessous.

\figureIRD{step2_dialog_select_phase_images_1.png}{Fen\^etre de s\'election des images de r\'ef\'erence. Pour la bande spectrale visible la première image est par défaut celle du point de décollage qui a servi à définir la mission à l'étape 1. Mais elle peut être changée ici.}{0.95}
La fenêtre de dialogue permet la définition :
\begin{enumerate}
	\item du point de décollage,
	\item des points de début et fin de synchronisation,
	\item des points de début et fin de vol au-dessus de la zone étudiée par la mission.
\end{enumerate}
Les images sont sélectionnées pour les spectres visible et proche infrarouge  respectivement dans les dossiers
\begin{itemize} 
	\item \texttt{C:/Air-Mission/FLY-YYYYMMDD-hhmm-location/AerialPhotography/VIS}, 
   	\item  \texttt{C:/Air-Mission/FLY-YYYYMMDD-hhmm-location/AerialPhotography/NIR}.
\end{itemize} 
 Lors de cette étape de choix il est toujours possible après une première sélection de choisir une autre image. 
 L'ordre de sélection des cinq images n'a pas d'importance mais bien entendu first précède toujours last.
 Il est possible de sélectionner la même image pour \texttt{first mission image} et \texttt{first sync image}.  
  Pour les images DJI du spectre visible, si vous sélectionnez la même image pour \texttt{first sync image} et \texttt{last sync image} alors une seule image sera retenue (c'est un cas limite pour la synchronisation). 
  En revanche  si vous sélectionnez la même image pour \texttt{first sync image} et \texttt{last sync image} pour le spectre proche infrarouge la synchronisation automatique échouera à l'étape 3 et il vous sera demandé d'entrer manuellement l'offset de synchronisation des horloges. 
  Cela peut constituer une \og astuce \fg~ pour forcer l' offset de synchronisation des horloges à une valeur particulière en cas de synchronisation terrain problématique.
%\todo{\bouton{Load first mission image.}, \bouton{Load first sync image.}, \bouton{Load last sync image.},
%----------------------------------------------------------
\section{Bouton \bouton{Load first mission image.}}\label{btn:step2_Load_first_mission_image}
Le choix de cette image définit l'origine de la ligne de temps pour la synchronisation des horloges des caméras. En règle générale on choisira la première prise de vue VIS de la mission.
\section{Boutons \bouton{Load first sync image.}, \bouton{Load last sync image.}}\label{btn:step2_Load_synchro_image}
L'utilisateur doit choisir l'image de début du cycle de synchronisation. Pour les images visibles, si la procédure a été respectée sur le terrain, il s'agit de la première ou deuxième image enregistrée. Pour les images infrarouges il faut repérer la première image stable où le système de mire est visible en totalité et où la mire mobile est effectivement en rotation. Ensuite il faut choisir la dernière image de la phase de synchronisation. En général on la repère facilement car lorsque la phase de synchronisation est terminée le pilote prend rapidement de l'altitude et le changement d'aspect est bien visible. Les choix faits ici conditionnent en partie la réussite de l'étape 3 de synchronisation des horloges. Mais après l'étape 3 il est possible de revenir à l'étape 2 et de redéfinir ces images.
\figureIRD[fig:step2_dialog_select_VIS]{step2_dialog_select_phase_image_VIS.png}{Fen\^etre de s\'election des images de r\'ef\'erence. Bande spectrale visible (VIS).}{0.95}
\section{Boutons \bouton{Load first fly image.}, \bouton{Load last fly image.}}\label{btn:step2_Load_mission_image}
L'utilisateur doit choisir la première image prise au-dessus de la zone à étudier (vol mission ou hyperlapse survey). On élimine ainsi les images de l'approche entre le décollage/synchronisation à très basse altitude et le vol de la mission à \og haute altitude \fg~ (hyperlapse survey). À noter qu'en principe le vol mission devrait s'opérer à altitude constante du drone et au-dessus d'une zone dont les altitudes varient peu par rapport à l'altitude du drone par rapport au sol. L'observation de la figure \ref{fig:fly_vertical_profil} permet de juger de la qualité du plan de vol réalisé sur le terrain.		
\section{Bouton \bouton{Validate mission images}}\label{btn:step2_validate}
\bouton{Cancel}, à sa gauche, ferme la fenêtre (après confirmation si les images ne sont pas validées). Ce bouton déclenche la construction des listes des images des trois phases de vol pertinentes pour IRDrone : takeoff, synchro et fly. Aucune image n'est déplacée. L'opération est très rapide.
\figureIRD[fig:step2_dialog_select_NIR]{step2_dialog_select_phase_images_NIR.png}{Fen\^etre de s\'election des images de r\'ef\'erence. Bande spectrale proche infrarouge (NIR).}{0.95}

%==========================================================
\legacyend

\chapter{Step 3 -- Synchronisation des caméras}\label{ch:step3}
\section{Principe}
Pour combiner différentes bandes spectrales, les images doivent être prises au même instant. Malheureusement ce n'est techniquement pas possible avec les matériels utilisés dans le projet IRdrone. C'est un choix de départ lié au caractère financièrement abordable (\emph{inexpensive and low cost}) du projet. Chaque caméra a sa propre horloge, qui définit la chronologie des images. On peut positionner les images sur une ligne représentant le temps. C'est la \og time line \fg~de la caméra. Dès lors il faut synchroniser les deux time lines VIS et NIR. Mettre les horloges à l'heure est insuffisant : d'une part elles ne donnent accès qu'à des points toutes les secondes, d'autre part on ne peut pas contrôler l'instant exact où l'opérateur sur le terrain a déclenché les caméras. On doit donc détecter des repères sur les images. Dans l'industrie cinématographique on synchronise la pellicule photo et la bande son à l'aide du \og clap \fg. Cette solution avait été envisagée dans les premières versions d'IRdrone, qui utilisaient la vidéo au lieu des images. Mais la vidéo posait de nombreux problèmes de qualité (résolution, \og jelly \fg, ...) et a été abandonnée. Par contre le principe d'un objet mobile sur les images au début du vol a été conservé. L'idée est de repérer un objet (automatiquement) et de suivre sa trajectoire sur les deux caméras. En faisant coïncider les deux trajectoires on synchronise les horloges. On peut atteindre une précision de l'ordre de 0.05s, ce qui est amplement suffisant dans notre cas.
IRdrone utilise comme objet mobile une mire ArUco 
\footnote{Ce type de mire est conçu pour être détecté automatiquement par logiciel, par exemple avec OpenCV. IRdrone utilise le dictionnaire \texttt{DICT\_4X4\_50} : la mire mobile porte l'identifiant 0 et les quatre mires fixes les identifiants 5, 22, 16 et 13.} 
fixée à un plateau tournant
\footnote{Le plateau est entraîné par un moteur pas à pas commandé par un Arduino. On programme le système pour obtenir une vitesse de rotation parfaitement constante. Elle est fixée à la valeur de 1 rpm (un tour par minute) de telle sorte que lors de la phase de vol de synchronisation les caméras enregistrent au moins une dizaine de positions angulaires différentes.} comme le montre la figure \ref{fig:step3_syncro_aruco_1}. 

\figureIRD[fig:step3_syncro_aruco_1]{step3_syncro_aruco_1.pdf}{Dispositif de synchronisation des horloges des caméras VIS et NIR. Le principe repose sur la détection de la position angulaire de la mire ArUco 0 qui tourne à vitesse constante. Les petites mires fixes permettent de définir un repère absolu pour s'affranchir des mouvements parasites du drone, particulièrement les rotations.}{0.95}

On mesure la position angulaire de la mire. Mais pour cela il faut un repère fixe. Les bords de l'image ont été utilisés initialement. Cependant ils ne sont fiables que si le drone ne tourne pas. Il faut donc une zone très calme et que le pilote ne touche pas à la manette de rotation. Pour libérer cette contrainte sur le terrain, on a ajouté au dispositif 4 mires ArUco fixes
\footnote{Théoriquement 2 mires suffisent pour définir un axe de référence. L'utilisation de 4 mires permet une certaine redondance. À noter que le logiciel tente de synchroniser dans le repère absolu, mais s'il échoue il synchronise dans le repère relatif lié aux bords des images (moins précis). Si là encore il y a échec, en dernier recours le logiciel demande à l'utilisateur de donner la valeur de l'offset de synchronisation.}. Elles permettent de définir un repère absolu qui améliore la précision du dispositif. Il est évident que cette phase de synchronisation doit être réalisée sur le terrain avec le plus grand soin. C'est la condition pour concilier le coût modéré du matériel et une bonne qualité des images multispectrales.

\subsection{Les trois phases de l'étape 3}
Le Step~3 se déroule en trois phases successives, enchaînées automatiquement (§\ref{sec:step3_chained}). En mode manuel, chacune est lancée par un bouton qui n'est activé qu'une fois la phase précédente terminée (figure~\ref{fig:step3_window}) :
\begin{enumerate}
  \item \emph{Détection des mires} (bouton \bouton{Run ArUco detection (VIS + NIR)}) : sur chaque image de synchronisation, on mesure l'angle de la mire mobile (§\ref{btn:step3_aruco}) ;
  \item \emph{Calcul du décalage temporel} (bouton \bouton{Compute time alignment}) : on cherche le décalage $\Delta t$ qui fait coïncider les deux séries d'angles (§\ref{btn:step3_time_shift}) ;
  \item \emph{Appariement} (bouton \bouton{Pair VIS / NIR images}) : on associe à chaque image VIS les images NIR les plus proches dans le temps (§\ref{btn:step3_pairing}).
\end{enumerate}

\subsection{Enchaînement automatique}\label{sec:step3_chained}
L'étape~3 ne démarre jamais seule : qu'elle soit ouverte depuis l'étape~2 (\bouton{Continue with synchronisation}) ou depuis la fenêtre principale, le bouton \bouton{Synchronise the cameras} (bleu foncé, au-dessus de la barre de progression) enchaîne les trois phases (§\ref{btn:step3_run}). Une ligne d'avancement montre la phase en cours et son résultat : nombre de mires vues, décalage, nombre de paires (figure~\ref{fig:step3_chained}). Une seule question est posée, après l'alignement des horloges (§\ref{btn:step3_overlap}).

Une fois les images appariées, rien ne démarre seul : \bouton{Continue with image processing} ouvre l'étape~4 quand vous êtes prêt (§\ref{btn:step3_continue}).
Si le décalage ne peut pas être mesuré, la saisie manuelle est proposée et l'enchaînement s'arrête.

\figureIRD[fig:step3_chained]{step3_chained.png}{Étape 3 enchaînée : mires détectées, horloges alignées ; la question \og Do the two curves overlap? \fg~attend la réponse avant l'appariement.}{0.7}

Rouverte une fois l'étape faite (clic sur l'étape~3 de la fenêtre principale), la fenêtre ne recalcule rien : elle montre le résultat enregistré, c'est-à-dire le décalage des horloges, le nombre de paires et les courbes alignées (§\ref{btn:step3_before}). \bouton{Synchronise again} la refait ; \bouton{Manual mode…} montre les boutons des trois phases.

\section{La question \og Do the two curves overlap? \fg}\label{btn:step3_overlap}
Après l'alignement des horloges, la fenêtre montre l'angle de la mire en fonction du temps : courbe verte pour le visible, rouge pour l'infrarouge, la courbe infrarouge étant décalée du décalage calculé.
\begin{description}
	\item[\bouton{Yes} (touche Y)] les deux courbes se superposent : le décalage est bon, les images sont appariées (§\ref{btn:step3_pairing}).
	\item[\bouton{No} (touche N)] ouvre le réglage à la main (§\ref{btn:step3_adjust}) pour glisser la courbe infrarouge sur la courbe visible ; la question est ensuite reposée avec le nouveau décalage.
\end{description}
Des courbes bien superposées ont des points verts sur la ligne rouge, dents de scie comprises. Les exemples du §\ref{sec:step3_examples} montrent un bon et un mauvais cas.

\section{Bouton \bouton{Synchronise the cameras}}\label{btn:step3_run}
L'action principale de l'étape~3 : elle lance la détection des mires, l'alignement des horloges puis, une fois répondu à la question sur les courbes, l'appariement. Elle devient \bouton{Synchronise again} quand l'étape a déjà été faite.

\section{Bouton \bouton{Manual mode…}}\label{btn:step3_manual}
Arrête l'enchaînement et fait apparaître les boutons des trois phases (§\ref{btn:step3_aruco} à \ref{btn:step3_pairing}), à lancer un par un, par exemple pour corriger le décalage avec \bouton{Adjust alignment...} avant l'appariement.

\section{Bouton \bouton{Enter the clock offset…}}\label{btn:step3_clock_offset}
Option avancée, pour un vol sans la cible ArUco (ancienne synchronisation par les dates). Le bouton n'apparaît que si la case \emph{Advanced – Step 3: let me enter the camera clock offset by hand} est cochée dans les options (roue dentée \bouton{⚙} en bas de la fenêtre principale, §\ref{btn:home_workdir}). Avec cette option, l'étape~3 ne démarre pas d'elle-même en venant de l'étape~2 : on choisit entre \bouton{Synchronise the cameras} (vol avec cible) et ce bouton.

Une petite fenêtre demande de combien de secondes l'horloge de la caméra infrarouge avance sur celle du drone (négatif si elle retarde, 0 si les deux horloges ont été réglées à la même heure). C'est le délai de l'ancienne synchronisation. Elle affiche le décalage $\Delta t$ de la ligne de temps infrarouge qui en découle, calculé à partir des dates des premières images. \bouton{Pair the images} saute la détection des cibles et lance directement l'appariement (§\ref{btn:step3_pairing}). Aucune courbe ne peut être vérifiée : la qualité de l'appariement dépend de la justesse du délai saisi.

\section{Bouton \bouton{Continue with image processing}}\label{btn:step3_continue}
Apparaît en bas à droite, à côté de \bouton{Close}, une fois les images appariées. Il ferme la fenêtre et ouvre l'étape~4. L'alignement, long, n'y démarre pas tout seul : il se lance avec \bouton{Run alignment}.

\section{Bouton \bouton{Before alignment}}\label{btn:step3_before}
Quand l'étape~3 est rouverte une fois faite, la fenêtre montre les courbes après alignement (\texttt{Synchro/check\_time\_alignment.png}). \bouton{Before alignment} montre les courbes avant correction (\texttt{check\_time\_before\_alignment.png}) ; un second clic revient aux courbes alignées.

\subsection{Principe du calcul}
L'angle de la mire mobile est mesuré sur chaque image, puis déplié (voir §\ref{btn:step3_time_shift}). Comme la vitesse de rotation $\omega$ est constante, l'angle est une fonction affine du temps. La time line NIR, qui compte au moins deux points, fournit la droite de référence $\theta = \omega t + b_{NIR}$. Chaque point VIS fournit une estimation de l'ordonnée à l'origine $b_{VIS} = \theta_{VIS} - \omega t_{VIS}$ ; on en prend la médiane. Le décalage s'écrit
\begin{equation}
  \Delta t = \frac{b_{NIR} - b_{VIS}}{\omega},
\end{equation}
et la time line NIR corrigée est $t_{NIR} + \Delta t$. C'est pourquoi une seule image VIS suffit en théorie (cas~2 ci-dessous), à condition que la vitesse de rotation soit constante et que l'on travaille dans le repère absolu.

\subsection{Exemples de synchronisation}\label{sec:step3_examples}
Les images de la phase de synchronisation viennent de l'étape~2. On présente quelques cas ici.

\subsubsection{Cas 1 : bonne synchronisation}
Les figures~\ref{fig:synchro_cas1} (\subref{fig:synchro_cas1_avant} et \subref{fig:synchro_cas1_apres}) représentent une bonne synchronisation. Le pilote a déclenché manuellement l'hyperlapse de la caméra NIR puis il a fait décoller le drone, ce qui a déclenché automatiquement la rotation de la mire. Lorsque le pilote a placé le drone au-dessus de la mire, la caméra NIR enregistrait la mire en rotation. Les images entre l'instant de déclenchement de l'hyperlapse NIR et l'instant où le drone a été stabilisé au-dessus de la mire ont été ignorées au step~2. Ensuite le pilote a réglé les paramètres de l'hyperlapse de la caméra VIS puis l'a déclenché. Cela prend en général une trentaine de secondes pour un pilote entraîné. Dès que l'hyperlapse VIS commence, il enregistre la mire en mouvement. À partir de cet instant et avec une fréquence d'enregistrement de 2s, il suffit au pilote de maintenir le drone moins d'une dizaine de secondes au-dessus de la mire pour obtenir 4 à 5 points, avant de prendre rapidement de l'altitude pour rejoindre la zone à étudier (hyperlapse survey). L'utilisateur du logiciel doit veiller lors de l'étape 2 à éliminer les images NIR prises avant la stabilisation du drone au-dessus du dispositif de synchronisation et le cas échéant celles qui après stabilisation correspondraient à une mire non encore en rotation. 

\figurePaireIRD{fig:synchro_cas1}
  {Step3_time_alignment_failed_cas1_0.png}{Trajectoire de la mire avant synchronisation.}
  {Step3_time_alignment_failed_cas1_1.png}{Synchronisation réussie.}
  {Cas de bonne synchronisation.}

\subsubsection{Cas 2 : cas limite}
Les figures~\ref{fig:synchro_cas2} (\subref{fig:synchro_cas2_avant} et \subref{fig:synchro_cas2_apres}) sont un cas limite pour la synchronisation. Le pilote est resté très peu de temps au-dessus de la mire et une seule image de cette mire est disponible sur les images VIS. Ici la synchronisation fonctionne tout de même car la vitesse de rotation est maintenue constante et le repère absolu a pu être utilisé.

\figurePaireIRD{fig:synchro_cas2}
  {Step3_time_alignment_failed_cas2_0.png}{Trajectoire de la mire avant synchronisation.}
  {Step3_time_alignment_failed_cas2_1.png}{Cas limite de synchronisation.}
  {Cas limite de synchronisation.}

\subsubsection{Cas 3 : échec de synchronisation}
Les figures~\ref{fig:synchro_cas3} (\subref{fig:synchro_cas3_avant} et \subref{fig:synchro_cas3_apres}) montrent la trajectoire de la mire détectée sur les deux caméras. Ici la synchronisation échoue car les premières images retenues pour la caméra NIR correspondent à des instants où la mire n'est pas en rotation.

\figurePaireIRD{fig:synchro_cas3}
  {Step3_time_alignment_failed_cas3_0.png}{Trajectoire de la mire avant synchronisation.}
  {Step3_time_alignment_failed_cas3_1.png}{Échec de synchronisation.}
  {Cas d'échec de synchronisation.}
   
\figureIRD[fig:step3_window]{step3_window.png}{Fen\^etre de synchronisation.}{0.8}

\section{Bouton \bouton{Run ArUco detection (VIS + NIR)}}\label{btn:step3_aruco}
Ce bouton lance la détection des mires ArUco sur toutes les images de la phase de synchronisation choisies à l'étape~2, d'abord pour la bande VIS puis, automatiquement, pour la bande NIR. Une barre de progression indique l'avancement. Les calculs sont répartis sur plusieurs c\oe urs du processeur pour accélérer le traitement.

\subsection*{Ce que fait le logiciel pour chaque image}
\begin{enumerate}
  \item Les images DNG sont converties en TIF par RawTherapee avec un profil de développement dédié à la détection (fichiers \texttt{*\_RawTherapee.tif} dans les dossiers \texttt{Synchro\textbackslash VIS} et \texttt{Synchro\textbackslash NIR} de la mission). Cette conversion n'est faite qu'une fois : si le fichier existe déjà, il est réutilisé.
  \item Les mires ArUco visibles sont détectées avec OpenCV. Pour chacune, on retient son identifiant, la position de son centre et l'angle de son axe principal dans l'image.
  \item La mire mobile (identifiant~0) donne l'\emph{angle relatif} (repère de l'image, \texttt{angle\_img}).
  \item Si les quatre mires fixes (identifiants 5, 22, 16 et 13) sont aussi détectées, elles définissent un repère lié au sol : deux axes obtenus en moyennant les côtés opposés du quadrilatère formé par leurs centres. L'angle de la mire mobile est alors mesuré dans ce repère. C'est l'\emph{angle absolu} (\texttt{angle\_abs}), insensible aux rotations du drone autour de son axe vertical.
  \item Si la mire mobile ou l'une des mires fixes manque, l'image est signalée dans le journal. L'angle absolu n'est pas calculé (mode dégradé : seul l'angle relatif est conservé si la mire mobile est visible) ; si la mire mobile n'est pas détectée, l'image est ignorée.
  \item À chaque image est associée sa date relative (en secondes) sur la time line de sa caméra.
\end{enumerate}
Les résultats de détection sont conservés dans les métadonn\'ees de chaque image.

\subsection*{Images de contrôle}
Pendant la détection, chaque image traitée est enregistrée avec une surcouche graphique dans le dossier \texttt{Synchro\textbackslash Check\_ARUCO}. La fenêtre en affiche une régulièrement (environ toutes les 5 secondes) pour que l'utilisateur puisse juger de la qualité de la détection. La surcouche montre :
\begin{itemize}
  \item les centres des mires fixes (disques rouges) ;
  \item le centre de la mire mobile ;
  \item l'orientation de la mire mobile dans l'image (trait vert) ;
  \item l'orientation projetée dans le repère absolu (trait orange en pointillés) et les deux axes du repère absolu (flèches noires) lorsque les quatre mires fixes sont détectées ;
  \item une légende : angle relatif, angle absolu et différence entre les deux, ou le message \og absolute reference: unavailable \fg~si le repère absolu n'a pas pu être construit.
\end{itemize}
\textbf{Conseil :} parcourir ce dossier dans l'explorateur Windows est le meilleur moyen de repérer des détections douteuses avant de passer à la phase suivante.

% Figure à créer par l'utilisateur (PowerPoint) : 2 lignes de 4 images (VIS en haut, NIR en bas) extraites de Synchro\Check_ARUCO.
\figureIRD[fig:check_aruco]{Step3_check_aruco_VIS_NIR.pdf}{Exemples d'images de contrôle (dossier \texttt{Synchro\textbackslash Check\_ARUCO}). Ligne du haut : bande VIS ; ligne du bas : bande NIR. La mire mobile est repérée par son axe vert ; les centres des mires fixes par des disques rouges.}{0.95}

\section{Bouton \bouton{Compute time alignment}}\label{btn:step3_time_shift}
Ce bouton, actif seulement lorsque la détection est terminée, calcule le décalage temporel entre les deux caméras.

\subsection*{Choix des angles}
Par défaut le logiciel utilise les angles absolus. Il faut au moins 2 images NIR et 1 image VIS où l'angle absolu est disponible. À défaut, il utilise les angles relatifs (repère de l'image), avec les mêmes minima, ce qui est moins précis si le drone a tourné pendant la synchronisation. Si aucune des deux solutions n'est possible, le message \og Insufficient data to automatically align VIS and NIR camera timelines \fg~s'affiche et le logiciel demande à l'utilisateur de saisir lui-même le décalage en secondes (valeur par défaut : 0). Dans ce cas aucun graphique n'est produit.

\subsection*{Calcul}
\begin{enumerate}
  \item Les angles de chaque bande sont \og dépliés \fg : on supprime les sauts de $\pm 360^\circ$ pour obtenir une courbe continue qui croît (ou décroît) régulièrement avec le temps. Le dépliage est fait de la même façon pour VIS et NIR, en partant de la dernière image, c'est-à-dire de l'image où la mire est en principe en rotation.
  \item On ajuste une droite sur les points NIR : sa pente est la vitesse de rotation $\omega$.
  \item On en déduit, avec les points VIS, le décalage $\Delta t$ (voir le principe en début de chapitre). L'estimation utilise la médiane des points VIS pour être robuste aux valeurs aberrantes.
\end{enumerate}
Le résultat est affiché dans le journal de la fenêtre :
\begin{itemize}
  \item le décalage estimé $\Delta t$ (en secondes) ;
  \item la vitesse de rotation mesurée sur NIR et sur VIS (en rpm). Elle doit être proche de la valeur nominale de 1~rpm et les deux valeurs doivent être voisines ; l'écart relatif est affiché ;
  \item le nombre de points VIS et NIR utilisés ;
  \item la qualité de l'ajustement NIR ($R^2$, proche de 1) ;
  \item la dispersion des décalages individuels (écart interquartile et écart-type, en secondes).
\end{itemize}

\subsection*{Graphiques de contrôle}
Deux graphiques sont enregistrés dans le dossier \texttt{Synchro} de la mission. Pendant la question \og Do the two curves overlap? \fg, la fenêtre montre les courbes après alignement ; en mode manuel, elle montre les deux graphiques en alternance (changement toutes les 1,5 s) :
\begin{itemize}
  \item \texttt{check\_time\_before\_alignment.png} : angle de la mire en fonction du temps pour VIS (vert) et NIR (rouge) avant correction ;
  \item \texttt{check\_time\_alignment.png} : mêmes courbes après avoir décalé la time line NIR de $\Delta t$.
\end{itemize}
Si la synchronisation est bonne, les points verts se placent sur la droite des points rouges. Ce sont ces deux graphiques qui sont présentés dans les trois cas de ce chapitre.

\textbf{Que faire en cas d'échec ?} Si les points VIS et NIR ne se superposent pas après alignement, ou si les vitesses de rotation VIS et NIR sont très différentes, revenir au Step~2 et redéfinir les images de synchronisation (en particulier la première image NIR, qui doit montrer la mire effectivement en rotation), puis relancer le Step~3.

\section{Bouton \bouton{Adjust alignment...}}\label{btn:step3_adjust}
Ouvre le réglage du décalage à la main, aussi ouvert par la réponse \bouton{No} à la question \og Do the two curves overlap? \fg. La fenêtre trace la courbe visible (verte), la courbe infrarouge au décalage calculé (rouge, en pointillés) et la courbe infrarouge au décalage réglé (rouge). Le curseur et la case numérique déplacent la courbe infrarouge : flèches du clavier par pas de 1\,ms, touches page précédente / suivante par pas de 0,1\,s.
\begin{description}
	\item[\bouton{Save and close}] (\bouton{Close} tant que le décalage n'a pas changé ; fermer la fenêtre fait de même) garde ce décalage : les graphiques de contrôle sont refaits. Si les images étaient déjà appariées, il faut refaire l'appariement.
	\item[\bouton{↺}] (en bas à gauche, affiché quand le décalage diffère du décalage calculé) revient au décalage calculé à partir des mires.
\end{description}
Si le décalage n'a pas pu être calculé, une petite fenêtre demande directement sa valeur en secondes (ajoutée aux dates infrarouges).

\section{Bouton \bouton{Pair VIS / NIR images}}\label{btn:step3_pairing}
Ce bouton est actif une fois le décalage temporel connu (calculé ou saisi à la main). Il construit la liste des couples d'images qui seront alignés au Step~4.

\begin{enumerate}
  \item La time line VIS sert de référence. Toutes les dates NIR sont corrigées : $t'_{NIR} = t_{NIR} + \Delta t$.
  \item Pour chaque image VIS, on retient les \emph{deux} images NIR les plus proches dans le temps. On mémorise pour chacune l'écart signé $dt = t'_{NIR} - t_{VIS}$ : $dt<0$ si la NIR précède la VIS, $dt>0$ si elle la suit.
  \item Les métadonnées de vol (position GPS, coordonnées UTM, altitudes, angles de lacet, tangage et roulis du drone et de la nacelle, distances, cap, etc.), connues pour les images VIS, sont interpolées aux instants corrigés des images NIR. Les angles (lacet, cap) sont interpolés par le chemin angulaire le plus court. Hors de la plage de temps VIS, la valeur de l'extrémité est conservée.
  \item Le résultat est enregistré dans le fichier \texttt{AerialPhotography\textbackslash shoot\_points.json}, utilisé par les étapes suivantes. Le journal indique le nombre de couples créés.
\end{enumerate}
Un même fichier NIR peut être utilisé pour plusieurs images VIS (par exemple si la fréquence de prise de vue NIR est plus élevée) et aucun écart temporel maximal n'est imposé : la qualité de l'appariement dépend donc de celle de la synchronisation.

\section{Bouton \bouton{Close}}\label{btn:step3_close}
Ferme la fenêtre et rend la main au dialogue principal. Les fichiers produits (graphiques, images de contrôle, \texttt{shoot\_points.json}) restent dans le dossier de la mission. Le bouton \bouton{?} ouvre cette aide à la page de ce chapitre.

%==========================================================
\chapter{Step 4 -- Traitement des images}\label{ch:step4}
\section{Principe}
Le traitement des images consiste pour une paire VIS-NIR à aligner ces deux images pour que les bandes spectrales Red Green Blue et Nir se superposent. L'opération est complexe pour plusieurs raisons :
\begin{itemize}
	\item Même si les time lines sont synchronisées, les instants de prise de vue sont proches mais distincts.
	\item Les deux caméras ne pointent pas vers le même point au sol. La caméra VIS du drone DJI est fixée sur un gimbal très performant qui maintient la visée vers le nadir. À noter que le vol hyperlapse est fait en mode \og gimbal lock \fg~autrement dit l'angle de roll du gimbal reste fixe. Seuls ses angles de pitch et yaw varient pour maintenir la visée vers le nadir (à $90^{\circ}$ vers le bas). En revanche la caméra NIR SJCam M20 étant fixée rigidement sous le corps du drone sa visée est totalement tributaire de l'attitude du drone (angles de pitch et yaw en particulier). 
	\item Les distances focales des deux caméras sont différentes. Pour la caméra VIS on a un angle (FOV) de $84^{\circ}$ alors que pour la caméra NIR la focale est ultra-courte avec un angle (FOV) supérieur à $150^{\circ}$ ce qui correspond quasiment à un \og fish eye \fg.
	\item Les optiques présentent de la distorsion et du vignetage. Pour la distorsion les deux caméras ont été préalablement calibrées par la méthode de Zhang.
\end{itemize} 



En pratique on commence par convertir les images du format raw-dng au format tif (via le logiciel tiers RawTherapee
\footnote{Via le logiciel tiers open source RawTherapee. L'opération est coûteuse la première fois mais l'image est mise en cache pour des runs ultérieurs.}).  On élimine la distorsion des deux images. On détermine les déplacements théoriques à appliquer pour ramener le centre de visée NIR sur le centre de visée VIS. Ces déplacements sont traduits en termes d'angles pitch, yaw et roll. Ils servent à initialiser la procédure d'alignement. L'alignement se déroule très schématiquement en deux temps. On détermine en premier les paramètres d'une homographie à appliquer sur l'image NIR pour qu'elle \og match \fg~avec l'image VIS. C'est un processus de minimisation complexe spécialement développé pour les images multispectrales 
	\footnote{Les techniques habituelles pour aligner deux images, qui sont extrêmement performantes dans le cas d'une bande spectrale unique ou en RGB, comme par exemple les \og sift \fg, échouent.}. 
Ce premier alignement (coarse process) est suivi d'un alignement fin (fine process) qui applique des homographies locales pour diminuer encore le critère d'alignement. L'algorithme est décrit en annexe (§\ref{sec:ann_algorithme}). L'utilisateur doit être conscient que la méthode n'est pas infaillible et surtout qu'elle est tributaire des conditions de vol sur le terrain. 

\figurePaireIRD{fig:step4_pair_VIS_NIR}
{step3_img_VIS.png}{Spectre visible (VIS).}
{step3_img_NIR.png}{Spectre proche infrarouge (NIR).}
{Couple d'images d'un shoot point. Les deux bandes spectrales enregistrent des informations différentes. Par exemple la surface de la rivière est totalement noire en proche infrarouge. L'eau ne réfléchit pas les longueurs d'onde au-delà de $750\,nm$.}  

\section{Dialogue de l'étape 4}
La fenêtre de l'étape~4 (figure~\ref{fig:step4_window}) rappelle le nom de la mission (le chemin complet est dans l'info-bulle) et le nombre de paires d'images. Elle comporte, de haut en bas : les options de sortie, le choix des images alignées avec la carte du vol, le bouton \bouton{Run alignment} (\bouton{Stop} pendant le calcul), l'état et la barre de progression du calcul, le journal, puis, en bas, \bouton{?} et \bouton{⚙} (réglages) à gauche, \bouton{Close} et, une fois des images alignées, \bouton{Continue with the review} à droite. La fenêtre s'adapte à la taille de l'écran.

\figureIRD[fig:step4_window]{step4_window.png}{Fen\^etre de traitement des \og shoot points \fg~(paires d'images VIS-NIR), avant le calcul.}{0.7}


\section{Options de sortie}\label{btn:step4_options}


\subsection{Case \bouton{first VIS image number} et \bouton{last VIS image number}}
Ces deux cases permettent de sélectionner des shoot points à traiter. Par défaut IRDrone propose la plage qui a été définie à l'étape~2. Il est possible de ne traiter qu'une image en entrant le même numéro pour first et last images, sous réserve que le \texttt{Maximum VIS/NIR time difference} soit compatible. En cas de doute fixez le \texttt{Maximum VIS/NIR time difference} à la valeur $3 \, s$.

Par défaut la procédure d'alignement produit  pour 

$\texttt{first VIS image number} \leq XXXX \leq \texttt{last VIS image number}$:
\begin{itemize}
	\item les quadruplets d'images multispectrales red (1), green (2), blue (3), nir (4) (\texttt{IRD\_XXXX\_n.tif}) placés dans le dossier \texttt{ImgIRdrone\textbackslash imgMultispectral}.
	\item les images \texttt{VIS\_XXXX.jpg}, \texttt{VIR\_XXXX\_local.jpg}.
	\item les fichiers \texttt{VIS\_XXXX\_REDO.bat} et deux fichiers techniques \texttt{VIS\_XXXX\_motion\_model.json}, \texttt{VIS\_XXXX\_motion\_model.npy}.
	\item le fichier \texttt{ImgIRdrone\textbackslash shoot\_points.json} qui contient un dictionnaire des shoot points traités avec en particulier la clé  \texttt{"failed"} dont la valeur est fixée à \texttt{false} par défaut. L'utilisateur vérifie cette hypothèse dans la revue des images alignées (chapitre~\ref{ch:review}). 
\end{itemize}


\subsection{Case \bouton{NDVI}, case \bouton{NIR}}
Par défaut la procédure d'alignement produit les images VIS et VIR au format jpeg. C'est utile pour vérifier visuellement la réussite des alignements. En cochant les cases \texttt{NDVI} et/ou \texttt{NIR} on peut générer les images NDVI et NIR au format jpeg. Ne cocher ces cases que si vous envisagez d'utiliser ces représentations directement dans un rapport. Notez qu'à l'étape~6 vous avez la possibilité de générer de telles représentations, mais cette fois géolocalisées et orthorectifiées, et visualisables dans le logiciel QGIS.
\subsection{Cases \bouton{Local alignment} et \bouton{Global alignment}}
Comme expliqué plus haut l'algorithme d'alignement procède en deux passes (coarse process et fine process). On peut voir le résultat de l'une ou l'autre (ou les deux) en cochant les cases \texttt{Global alignment} et \texttt{Local alignment}. La case \texttt{Global alignment} montre les résultats du coarse process alors que la case \texttt{Local alignment} montre le résultat du raffinement.  Une de ces deux cases doit impérativement être cochée. Nous suggérons de cocher systématiquement la case \texttt{Local alignment}. Il est à noter que dans certaines circonstances l'alignement global réussit mais l'alignement local échoue.

\subsection{Case \bouton{Maximum VIS/NIR time difference}}
Cette case permet de restreindre la liste des shoot points à aligner à ceux dont l'écart temporel de rang~1 vérifie $|dt|\leq \texttt{Maximum VIS/NIR time difference}$. La valeur par défaut est $0.5\,s$. Cela correspond aux images bien synchronisées. En pratique avec les fréquences de prise de vue de $2 \, s$ pour les images VIS et $3 \, s$ pour les images NIR cela revient à sélectionner un shoot point sur trois.


\section{Best mapping : n'aligner que les images utiles à la carte}\label{btn:step4_best_mapping}
Le cadre \bouton{Images aligned} propose deux choix, qui reprennent celui de l'étape~2 (§\ref{btn:step2_best_mapping}) : \bouton{Best mapping: only the images the map needs} (par défaut, avec un recouvrement de $75\,\%$) ou \bouton{All the synchronised images}. Parmi les shoot points bien synchronisés (case \texttt{Maximum VIS/NIR time difference}), IRDrone ne garde que ceux dont les empreintes au sol se recouvrent de la valeur choisie (case \bouton{Overlap between kept images}) : la première image du survol, puis, à chaque fois, la dernière image qui recouvre encore assez la précédente gardée, et la dernière image du survol. Toutes les images du survol servent à suivre le déplacement du drone ; seules les images bien synchronisées peuvent être gardées. Quand l'image qui serait gardée est mal synchronisée, IRDrone garde la plus proche bien synchronisée : la dernière avant elle ou, si son recouvrement est plus proche de la cible, la première après elle (recouvrement au plus 5 points sous la valeur choisie, 70\,\% pour 75\,\%). Le nombre d'images alignées reste ainsi proche de l'aperçu de l'étape~2 ; sans cela, chaque image mal synchronisée raccourcirait le pas et beaucoup plus d'images seraient gardées. La ligne sous les choix indique le nombre de shoot points qui seront alignés, par exemple \og 18 of 39 synchronised shoot points will be aligned \fg.

À droite, la carte du vol montre les shoot points : en vert ceux qui sont bien synchronisés et dans la plage d'images, en gris les autres ; les images gardées par le best mapping sont plus grosses, avec leur empreinte au sol en vert, ce qui montre le recouvrement. Elle se met à jour dès qu'un réglage change (plage, écart de temps, recouvrement, choix). Molette pour zoomer, glisser pour se déplacer, double-clic pour revoir tout le vol. Pendant l'alignement, la carte suit le calcul le long du vol : les empreintes des images alignées se remplissent de vert, celles en cours d'alignement (une par processus de calcul) sont entourées d'orange, celles qui restent à faire sont en pointillés ; le compte \og 7 of 18 done \fg{} s'affiche en bas de la carte. Le résultat reste affiché à la fin, jusqu'au prochain changement de réglage.

C'est en général le meilleur choix : beaucoup moins de calcul (l'alignement est l'étape la plus longue), moins de données, et une carte ODM souvent meilleure (trop d'images presque identiques ralentissent ODM et dégradent son résultat). Choisissez \bouton{All the synchronised images} pour aligner toutes les images synchronisées, par exemple pour les regarder une à une dans la revue. Le choix est enregistré dans \texttt{FlightAnalytics\textbackslash best\_mapping.json} au lancement de l'alignement ; l'étape~6 l'utilise pour la sélection des images d'ODM (§\ref{btn:step6_selection}).

La même sélection existe en ligne de commande : \texttt{run\_after\_preprocessing.py --best-mapping-overlap 75}.

\section{Bouton \bouton{Run alignment}}\label{btn:step4_run_alignment}
Lance la procédure d'alignement sur les shoot points vérifiant simultanément les deux contraintes 
\begin{itemize}
	\item $\texttt{first VIS image number} \leq XXXX \leq \texttt{last VIS image number}$,
	\item $|dt|\leq \texttt{Maximum VIS/NIR time difference}$,
\end{itemize} 
et, si \bouton{Best mapping: only the images the map needs} est choisi, seulement les images utiles à la carte (§\ref{btn:step4_best_mapping}).
Attention l'alignement des images est une procédure qui prend beaucoup de temps. Cela dépend bien entendu de la machine. Les  calculs sont parallélisés sur plusieurs processeurs si la machine est \og multi c\oe urs \fg. La barre de progression indique l'avancement des calculs ainsi qu'une estimation du temps de traitement restant. À titre d'exemple, sur un ordinateur portable avec un processeur Intel Core~i7 à quatre c\oe urs et 16\,Go de mémoire (3 workers d'alignement), l'alignement des images VIS et NIR d'un shoot point prend environ $15\,s$ : il faut donc environ une demi-heure pour traiter 120 images. Pour une grande mission, lancer ces calculs à l'heure du déjeuner reste une bonne idée.   

L'alignement ne démarre jamais tout seul, même quand l'étape~4 est ouverte par l'étape~3. Pendant le calcul, le bouton devient \bouton{Stop} (§\ref{btn:step4_stop}).

\section{Bouton \bouton{Stop}}\label{btn:step4_stop}
Pendant le calcul, \bouton{Run alignment} devient \bouton{Stop}. Après confirmation, le calcul et tous ses processus de travail s'arrêtent ; la fenêtre indique combien d'images ont été alignées. Les images déjà alignées sont gardées ; \bouton{Run alignment again} relance le calcul, et les points déjà faits réutilisent leur modèle de mouvement (plus rapide). Pour ne traiter que la suite, augmentez \texttt{First VIS image number}.

\section{Bouton \bouton{⚙} : réglages}\label{btn:step4_settings}
La roue dentée ouvre les réglages du calcul, gardés pour les sessions suivantes :
\begin{itemize}
	\item \texttt{Alignment worker processes} : nombre de shoot points alignés en parallèle. \texttt{Auto} utilise les c\oe urs physiques, dans la limite de la mémoire disponible au lancement (environ 3\,Go par processus). \bouton{Safe mode (1)} n'en garde qu'un : plus lent, mais le moins de mémoire. À utiliser si l'ordinateur se fige ou si le calcul est interrompu.
	\item \texttt{Log detail} : quantité de messages du calcul dans le journal de la fenêtre : avertissements seulement, étapes du calcul, ou détail des temps. Avec le détail des temps (\emph{Debug}), un réalignement à la main depuis la revue (section~\ref{btn:review_redo}) écrit aussi toutes ses images de diagnostic : plusieurs minutes par point au lieu d'environ une demi-minute.
\end{itemize}
\bouton{Save and close} (\bouton{Close} tant que rien n'a changé) garde les réglages. \bouton{↺}, affiché quand les réglages diffèrent des réglages par défaut, revient à \texttt{Auto} et aux avertissements seulement.

\section{Bouton \bouton{Close}}\label{btn:step4_close}
Ferme la fenêtre et rend la main au dialogue principal. Le bouton \bouton{?} ouvre cette aide à la page de ce chapitre. Il est grisé pendant l'alignement (arrêtez-le d'abord avec \bouton{Stop}).

\section{Bouton \bouton{Continue with the review}}\label{btn:step4_continue}
Apparaît en bas à droite, après \bouton{Close}, une fois des images alignées, et est mis en avant à la fin de l'alignement. Il ferme la fenêtre et ouvre l'étape~5, la revue des images alignées (chapitre~\ref{ch:review}). Quand l'alignement a été lancé depuis l'étape~3, la revue s'ouvre d'elle-même à la fin.

%==========================================================
\chapter{Step 5 -- Revue des images alignées}\label{ch:review}
La revue présente tous les points du relevé sur une ligne de temps pour vérifier l'alignement de chaque image (figure~\ref{fig:step4_review}). Elle s'ouvre avec \bouton{Continue with the review} de l'étape~4, ou avec l'étape~5 de la fenêtre principale (une fois des images alignées par l'étape~4). Elle remplace l'examen des images une par une dans l'explorateur. Rien n'est calculé à l'ouverture : l'étape compte comme faite une fois la revue fermée, les signalements et confirmations étant enregistrés au fur et à mesure.

\figureIRD[fig:step4_review]{step4_review.png}{Revue des images alignées : rideau entre visible et infrarouge aligné, fiche du point (décalage mesuré, image infrarouge associée), mini-carte, et en bas la ligne de temps des points (décalage résiduel, vignettes, état).}{0.95}

\section{Mesure automatique}\label{sec:review_measure}
Pour chaque point, le décalage résiduel entre les contours de l'image visible et ceux de l'infrarouge aligné est mesuré (en pixels de l'image complète, médiane sur 16 zones). Un point bien aligné mesure environ 1~px. Vert : moins de 3~px ; orange clair : entre 3 et 6~px, à regarder ; orange : plus de 6~px, probablement mal aligné. Une zone sans texture (eau, champ uniforme) ne se mesure pas. La fenêtre s'ouvre sur le point le plus douteux.

Les mesures et les points vérifiés sont gardés dans \texttt{ImgIRdrone/review.json} : rouvrir la revue ne refait pas les mesures, sauf pour les images refaites.

\section{Comparer le visible et l'infrarouge}\label{btn:review_compare}
Quatre façons de comparer, avec les boutons au-dessus de l'image ou les touches 1 à 4 ; la revue s'ouvre en \bouton{Overlay}. Un cadre autour de l'image est dans la couleur de l'état du point (vert aligné ou confirmé, rouge signalé, violet réaligné\ldots), les étiquettes sur l'image restent neutres (les rouges et cyans sont ceux des images) :
\begin{description}
	\item[\bouton{Swipe} (1)] un rideau sépare le visible (à gauche) de l'infrarouge aligné en fausses couleurs (à droite) ; routes, haies et toits doivent se prolonger à travers le rideau. Glisser le rideau pour le déplacer.
	\item[\bouton{Blink} (2)] visible et infrarouge alternent : une image mal alignée \og saute \fg.
	\item[\bouton{Overlay} (3)] visible en cyan, infrarouge en rouge (la végétation est naturellement rouge) ; un mauvais alignement double chaque contour.
	\item[\bouton{Side by side} (4)] les deux images côte à côte.
\end{description}
Molette : zoom (un décalage de quelques pixels ne se voit qu'en zoomant) ; glisser : déplacement ; double-clic : zoom $\times 4$ ou image entière. \bouton{Open full size} ouvre l'image alignée dans la visionneuse d'images de l'ordinateur.

\section{Situer l'image sur la carte}\label{btn:review_map}
La mini-carte (dans la fiche du point) montre le vol, le point courant et l'emprise au sol de l'image (le trait épais est le haut de l'image), dans la couleur de son état (vert aligné ou confirmé, rouge signalé\ldots) ; les shoot points que l'étape~4 n'a pas alignés sont de petits points gris, en taille normale quand on zoome fort. Un clic l'agrandit (\bouton{Map}, touche 5). Sous la mini-carte, la légende ne donne que les états présents, avec leur nombre. La grande carte affiche l'image alignée dans son emprise, par-dessus l'orthophoto de l'étape~6 quand elle existe et les photos aériennes en fond : l'infrarouge doit prolonger les routes et les haies du fond. Survoler un point montre sa vignette, un clic le sélectionne ; un clic hors des points ne laisse aucun point sélectionné sur la carte : la trajectoire seule, sans emprise ni image (le point reste celui de la fiche ; un clic sur un point, ou un déplacement dans le bandeau, en sélectionne un à nouveau) ; molette et glisser pour zoomer et se déplacer, double-clic pour voir tout le vol. Le menu \bouton{Layers}, le fond de carte hors ligne et le calcul de l'emprise sont décrits au §\ref{btn:home_map}. Sans position GPS dans les shoot points, la carte n'est pas disponible.

\paragraph{Aligner un point laissé de côté par l'étape~4.} Un clic droit sur un point gris de la carte (mini-carte ou grande carte) affiche son écart de temps avec l'infrarouge et propose \bouton{Align this point} : il est aligné comme le fait l'étape~4. Quand une image est irrécupérable, même à la main, un voisin écarté par le best mapping ou le filtre temporel peut être bon : alignez-le ainsi. Le point est aligné en arrière-plan (orange sur la carte, avancement dans le bandeau, \bouton{Stop aligning} pour arrêter) pendant que vous continuez la revue ; il apparaît ensuite dans la frise, à revoir comme les autres, et l'étape~6 le cartographie toujours (sauf si vous le signalez), quels que soient sa tolérance temporelle et son choix d'images. Le filtre temporel est le plus large de l'étape~4 ($|dt|<6$~s) ; un point sans image infrarouge assez proche dans le temps ne peut pas être aligné.

\figureIRD[fig:step4_review_map]{step4_review_map.png}{Mode \og Map \fg~de la revue : vol, emprise et image alignée du point courant, orthophoto de l'étape~6 sur les photos aériennes IGN.}{0.9}

\section{Se déplacer}\label{btn:review_navigation}
\begin{itemize}
	\item \bouton{◀} et \bouton{▶} (flèches du clavier) : point précédent, point suivant.
	\item \bouton{Next to review} (touche N) : prochain point qui demande un coup d'\oe il (réaligné à confirmer, mesuré douteux ou non mesuré). Un point mesuré aligné (vert) ne demande rien : il n'y a rien à marquer comme bon.
	\item La ligne de temps du bas : survoler montre l'image, un clic sélectionne le point ; molette pour zoomer, glisser pour se déplacer, double-clic pour tout le relevé.
	\item \bouton{Close} ferme la revue ; ce qui a été signalé et confirmé est déjà enregistré. \bouton{Continue with mapping} la ferme aussi, puis ouvre l'étape~6 (§\ref{btn:review_continue}).
	\item L'info-bulle du bouton \bouton{?} rappelle toutes les touches.
\end{itemize}

\section{Décider : signaler un point}\label{btn:review_decide}
\begin{description}
	\item[\bouton{Flag: leave out of the map} (touche F)] signale le point, \textbf{ce qui l'écarte de la carte} : mal aligné, mais aussi flou, ou avec une zone noire au bord de l'image infrarouge (drone incliné par un coup de vent : la caméra infrarouge ne voit pas tout le champ du visible). Un second appui retire le signalement. Le signalement est écrit tout de suite dans \texttt{ImgIRdrone\textbackslash shoot\_points.json}\footnote{Ce fichier ne contient pas les mêmes informations que le fichier \texttt{AerialPhotography\textbackslash shoot\_points.json}.} et l'étape~6 ignore ce point. Un point signalé peut être réaligné à la main (§\ref{btn:review_redo}) ; si cela ne le rattrape pas, le laisser signalé.
	\item[\bouton{Flag the N probably misaligned}] signale d'un coup tous les points dont le défaut d'alignement mesuré dépasse 6~px (points rouges \emph{probably misaligned}). Il n'apparaît que s'il y en a : la mesure ne voit ni un flou ni une zone noire, que vous signalez vous-même.
\end{description}

L'exemple ci-dessous montre le début du fichier \texttt{ImgIRdrone\textbackslash shoot\_points.json} qui stocke ces informations. Toutes les images ont été traitées en choisissant le critère peu discriminant $|dt|<4\,s$ et la méthode "au plus près". Cela correspond à 501 images VIS et 350 images NIR. Ici l'utilisateur a estimé visuellement que le shoot point 120 était mal aligné et l'a signalé, ce qui a fait basculer la valeur de la clé \texttt{failed} de \texttt{false} à \texttt{true}.
\begin{verbatim}
	{
		"mission_path": "C:\\Air-Mission\\FLY-20220524-1715-Chalivoy-Milon",
		"time_shift": -48.116,
		"n_VIS": 501,
		"n_NIR": 350,
		"pairing_method": "nearest_time",
		"created_at": "2026-10-29T20:44:13.203434",
		"schema_version": "1.0",
		"shoot_points": [
		{
			"FileName": "VIS_0119.dng",
			"dt": -0.5390370000000075,
			"failed": false
		},
		{
			"FileName": "VIS_0120.dng",
			"dt": 0.6902460000000019,
			"failed": true
		},
		{
			"FileName": "VIS_0121.dng",
			"dt": -1.3297540000000083,
			"failed": false
		},
		...
		]
	}	 
\end{verbatim}

Que faire ensuite des shoot points mal alignés ? Deux options :
\begin{itemize}
	\item \og aider \fg~l'algorithme d'IRDrone à les aligner : \bouton{Realign N flagged by hand} (§\ref{btn:review_redo}). Selon le résultat, le signalement est gardé ou retiré.
	\item les écarter des traitements ultérieurs : c'est l'effet par défaut du signalement, l'étape~6 ignore les points \texttt{"failed": true}.
\end{itemize}

\subsection{Exemple de succès d'alignement}
\figurePaireIRD{fig:goodalignment}
{step4_goodalignment_VIS.jpg}{Image dans le spectre visible.}
{step4_goodalignment_NIR.jpg}{Image la plus proche $|dt|=0.165~s$ dans le spectre infrarouge.}
{Cas de succès d'alignement : l'écart temporel est faible et l'attitude du drone est restée dans les spécifications des paramètres de vol.}

\figureIRD[fig:step4_goodalignment_VIR]{step4_goodalignment_VIR.jpg}{Alignement réussi. Les deux images ont été prises presque au même moment et en l'absence de vent. }{0.5}

\subsection{Exemple d'échec d'alignement}
Si le drone sort de sa zone des paramètres de vol, par exemple en cas de fortes rafales de vent qui l'obligent à accentuer ses angles de pitch et yaw, alors les centres de visée des deux caméras peuvent être tellement éloignés qu'il est impossible de trouver un alignement. L'alignement peut aussi être défaillant si l'écart temporel entre les prises de vue VIS et NIR est important (il peut être parfois supérieur à 2\,s) car le déplacement horizontal du drone est alors mal corrigé. L'alignement peut aussi être imparfait sans être un échec s'il y a un écart temporel et que les végétaux bougent sous l'action du vent. L'image VIR présente alors des \og fantômes \fg. Les figures \ref{fig:misalignment} et \ref{fig:step4_misalignment_VIR} montrent un cas d'échec de l'alignement lié à une forte rafale de vent qui déstabilise le drone. 

\figurePaireIRD{fig:misalignment}
{step4_misalignment_VIS.jpg}{Image dans le spectre visible.}
{Step4_misalignment_NIR.jpg}{Image la plus proche $|dt|=1.256~s$ dans le spectre infrarouge.}
{Cas d'échec d'alignement lié à une attitude du drone hors des spécifications des paramètres de vol et liée à une rafale de vent au moment de la prise de vue.}

\figureIRD[fig:step4_misalignment_VIR]{step4_misalignment_VIR.jpg}{Échec de l'alignement. L'effet de \og trou \fg est particulièrement caractéristique. Il peut aisément être repéré par l'utilisateur sur les images VIR. }{0.5}

\section{Réalignement à la main}\label{btn:review_redo}
Le bouton \bouton{Realign N flagged by hand} rattrape les shoot points signalés mal alignés. Il n'y a pas de miracle : si une rafale de vent a écarté le drone de sa zone de vol, les deux caméras ne voient pas la même scène et aucun alignement n'est possible. Mais parfois l'information existe et il suffit de \og pousser \fg~les paramètres de l'alignement grossier. En conditions normales une dizaine de shoot points au plus sont à reprendre.

\subsection*{Déroulement}
\begin{enumerate}
	\item Signaler les points mal alignés, puis appuyer sur \bouton{Realign N flagged by hand}. La fenêtre d'alignement s'ouvre avec tous les points signalés (figure~\ref{fig:step4_align_by_hand}).
	\item Pour chaque point, déplacer l'image infrarouge sur l'image visible jusqu'à faire coïncider les mêmes éléments (voir ci-dessous), puis \bouton{Placed, next} (touche Entrée). \textbf{L'alignement de ce point démarre aussitôt, en arrière-plan}, pendant que vous placez le suivant : pas d'attente entre deux points (le point passe à \emph{⚙ Computing} dans la liste et ne peut plus être déplacé). Les images des points suivants sont chargées en arrière-plan. \bouton{Skip this point} (touche S) laisse un point de côté : il reste signalé, donc écarté de la carte.
	\item \bouton{Compute N alignments} envoie au calcul les points placés sans \bouton{Placed, next} (quand tout est déjà envoyé, \bouton{Cancel} devient \bouton{Close}) et ferme la fenêtre. Les alignements sont calculés en arrière-plan, plusieurs à la fois (autant que de workers d'alignement de l'étape~4, bouton engrenage). Compter environ une demi-minute par point. Le bandeau du haut de la revue suit l'avancement : points terminés, points en cours avec leur étape et le temps écoulé, points en attente. \textbf{La revue reste utilisable} : on peut continuer à vérifier les autres points, ou faire autre chose. Les points en cours de calcul ne peuvent pas être signalés ni dé-signalés. \bouton{Stop} ne lance plus de nouvel alignement : ceux en cours se terminent, les autres restent signalés. La revue ne peut pas être fermée tant que des alignements sont en cours.
	\item Chaque point réaligné passe en violet, \emph{Realigned, to confirm} : il reste signalé (l'étape~6 l'ignore). Quand tout est calculé, la revue affiche le premier. Le vérifier avec le rideau, le clignotement ou la superposition et répondre \bouton{Yes, aligned} (touche Y) ou \bouton{Still misaligned} : le point suivant à confirmer s'affiche alors. \bouton{All N aligned} les confirme tous d'un coup. La touche N (point suivant à revoir) y va en premier. Les points à confirmer sont mémorisés si la revue est fermée.
	\item Avec \bouton{Yes, aligned}, la clé \texttt{"failed"} du point passe à \texttt{false} dans \texttt{shoot\_points.json} et le fichier \texttt{shoot\_points\_failed\_summary.json} est mis à jour (une copie \texttt{shoot\_points.json.bak} est conservée). Avec \bouton{Still misaligned}, le point reste signalé et pourra être réaligné de nouveau ; s'il ne peut pas être rattrapé (image infrarouge trop décalée, flou), il suffit de le laisser signalé : \textbf{il est écarté de la carte} (étape~6).
\end{enumerate}

Les sorties des alignements précédents sont archivées dans le dossier \texttt{VIS\_XXXX\_align\_traces\_manual\textbackslash original}. Le point retrouve les images qu'il avait (celles cochées à l'étape~4, section~\ref{btn:step4_options}) ; les images de diagnostic de l'alignement (dossier \texttt{Tech}, plusieurs minutes de calcul par point) ne sont écrites que si \texttt{Log detail} est réglé sur \emph{Debug} dans les réglages de l'étape~4 (bouton engrenage, section~\ref{btn:step4_settings}), ou depuis la ligne de commande (ci-dessous).

\subsection*{La fenêtre d'alignement}
\figureIRD[fig:step4_align_by_hand]{step4_align_by_hand.png}{Alignement à la main : les points signalés à gauche (coche : placé, rond : à placer), les deux images superposées au centre (ici le visible en cyan et l'infrarouge en rouge), les angles à droite. Ici l'infrarouge est encore environ 140~px trop haut : les contours sont doublés.}{0.9}

\begin{itemize}
	\item À gauche, la liste des points : un rond pour les points à placer, une coche pour les points placés, un tiret pour les points laissés de côté. Un clic va sur un point.
	\item Au centre, les deux images, affichées comme dans la revue : \bouton{Overlay} (les deux images en une : des contours doublés signalent un mauvais alignement ; affichage par défaut). \bouton{Colours} choisit ses couleurs, gardées pour la fois suivante : visible en couleur et infrarouge en gris (par défaut ; avec \bouton{Edges only}, les contours du visible sont verts), visible en vert et infrarouge en gris, visible en bleu et infrarouge en rouge, ou visible en cyan et infrarouge en rouge (comme dans la revue, proche de l'image VIR). \bouton{Swipe}, \bouton{Blink}, \bouton{Side by side} (touches 1 à 4). \bouton{Edges only} (touche G) montre les contours des deux images, ce qui aide sur les champs et les forêts où l'infrarouge et le visible se ressemblent peu. Molette pour zoomer, glisser pour se déplacer, double-clic pour voir toute l'image.
	\item Pour déplacer l'infrarouge, \textbf{vite puis finement} : d'abord les \textbf{curseurs} de \texttt{Yaw}, \texttt{Pitch} et \texttt{Roll} (panneau de droite) pour un grand déplacement, puis les boutons à leurs extrémités pour affiner, dans le sens qu'ils montrent : \bouton{◀}~\bouton{▶} à gauche et à droite (lacet), \bouton{▲}~\bouton{▼} en haut et en bas (tangage), \bouton{↺}~\bouton{↻} rotation (roulis) ; maintenus enfoncés, ils se répètent. Les \textbf{flèches du clavier} font la même chose que ces boutons (0,05\textdegree{} par appui, environ 2 à 3 pixels de l'image complète ; avec Maj, 10 fois plus), où que soit le curseur dans la fenêtre, même dans un champ de valeur ; Q et E le font tourner. On peut aussi \textbf{faire glisser l'infrarouge avec le bouton droit de la souris} (ou Maj + glisser) dans l'image, ou taper une valeur dans un champ.
	\item \bouton{Fit} (touche F) mesure le décalage entre les contours des deux images (tuile par tuile) et déplace l'infrarouge d'autant (lacet et tangage). Il fonctionne dès que les images se recouvrent à peu près ; s'il ne trouve pas de meilleure position, il le dit et ne change rien. \bouton{Reset} revient aux angles de l'alignement précédent.
	\item Sous ces boutons, le décalage mesuré entre les deux images, en pixels de l'image complète : sous 10~px le point est \emph{well placed}.
	\item À peu près suffit : ces angles ne sont qu'un point de départ, l'alignement grossier et l'affinage pyramidal repartent de là (§\ref{sec:ann_algorithme}).
	\item Un point dont les images ne peuvent pas être lues est marqué d'un signe d'avertissement et laissé de côté.
\end{itemize}

\subsection*{En ligne de commande : la fenêtre à curseurs}
Chaque point aligné a un script \texttt{VIS\_XXXX\_REDO.bat} (Windows) ou \texttt{VIS\_XXXX\_REDO.sh} dans \texttt{ImgIRdrone}, qui relance son alignement avec une fenêtre à curseurs (matplotlib) : c'est l'ancienne procédure, toujours utilisable en dehors de l'assistant. Après un réalignement dans l'assistant, le script repart des angles placés à la main.


\figureIRD[fig:redo_info]{step4_redo_info.png}{Fenêtre d'information de la fénêtre d'aide manuelle à l'alignement. }{0.5}

La fenêtre affiche l'image VIS et l'image NIR, avec cinq curseurs : \texttt{YAW}, \texttt{PITCH}, \texttt{ROLL}, \texttt{Absgrad visualize} et \texttt{Alpha}.
\begin{itemize}
	\item Le curseur \texttt{Alpha} choisit ce que l'on voit. Entre $-1$ (tout à gauche) et $-0{,}5$ (petit trait rouge médian), un rectangle de l'image NIR, légèrement transparent, est superposé à l'image VIS : plus \texttt{Alpha} est proche de $-1$, plus ce rectangle est grand. La valeur initiale est $-0{,}75$.
	\item Les curseurs \texttt{YAW}, \texttt{PITCH} et \texttt{ROLL} déplacent l'image NIR. La méthode consiste à modifier un angle, puis à faire varier \texttt{Alpha} pour observer le résultat, et à recommencer avec les trois angles. La clé est de repérer des éléments identiques dans les deux bandes spectrales (arbre, bord de chemin, bâtiment) et de les faire coïncider, d'abord au centre puis sur les bords du rectangle.
	\item Lorsque \texttt{Alpha} est déplacé vers la droite, la fenêtre n'affiche plus un rectangle NIR : entre $-0{,}5$ et $0$ elle montre la différence entre les deux images, et entre $0$ et $1$ un fondu enchaîné VIS/NIR. Cet usage est plus rapide mais réservé à un opérateur expérimenté.
	\item \texttt{Absgrad visualize} (valeur initiale $0{,}3$) affiche les gradients des deux images en deux couleurs lorsqu'il dépasse $0{,}5$ : les contours coïncident quand l'alignement est bon.\\
\end{itemize}

Lorsque l'utilisateur considère que l'alignement manuel est satisfaisant il tape la touche Q et la procédure d'alignement semi automatique démarre. \\


Comme pour tout outil de ce type, c'est en l'expérimentant que l'on se l'approprie ; une fois la méthode comprise, le réalignement est rapide et efficace.
Voici quelques étapes en image de la procédure
\footnote{L'outil est riche et offre plusieurs possibilités de visualiser les couches. Donc il n'y a pas une méthode unique pour arriver au résultat souhaité. Le slider \texttt{Absgrad visualize} permet de changer les représentations. Par exemple en le déplaçant vers la droite on peut obtenir un effet de contours. Ce n'est pas toujours évident à interpréter avec des forêts par exemple et encore moins dans des zones sans relief et nues.}. 
Sur l'image il y a une route qui traverse en diagonale. 
Il y a aussi un gros arbre en haut à droite, un petit arbuste en bas à gauche le long de la route  ainsi qu'un fossé en haut à gauche (\og trait noir \fg). Ce sont ces éléments caractéristiques que l'on va essayer d'aligner. 
À noter qu'à l'initialisation la couche NIR est légèrement transparente
\footnote{En déplaçant le slider \texttt{Alpha} vers la droite entre 0 et 1 la couche NIR passe de totalement transparente à opaque.}. 
Il n'y a pas vraiment d'ordre pour arriver à un résultat satisfaisant. 
Ne cherchez pas la perfection  car on n'applique que des transformations globales linéaires. Le code affinera globalement puis localement.
\figurePaireIRD{fig:redo_A}
{step4_redo_1.png}{Situation de départ après un passage dans la procédure associée au bouton \texttt{run alignement}.}
{step4_redo_2.png}{On a joué sur la rotation, qui passe de $0$ à $-2.46$, pour aligner la route NIR parallèlement à la route VIS. Attention utiliser ce slider délicatement. Lorsqu'il n'y a pas de ligne évidente comme ici la route, on utilise plutôt la rotation en dernier pour mettre la \og touche finale \fg.}
{Première modification en jouant sur le slider de rotation (\texttt{ROLL}).}
\figurePaireIRD{fig:redo_B}
{step4_redo_3.png}{Translation verticale vers le bas avec le slider \texttt{PITCH}. La valeur passe de  $-1.14$ à $2.04$. C'est une amplitude relativement importante et surtout qui change de signe.}
{step4_redo_4.png}{Translation horizontale avec le slider \texttt{YAW}. La valeur passe de $0.96$ à $1.8$.}
{On affine en jouant sur le curseur \texttt{Alpha} (en diminuant puis en agrandissant) et en observant l'alignement des éléments caractéristiques. La sensation de léger flou à la fin est liée à la transparence et aux petits défauts d'alignement locaux qui seront corrigés par la procédure d'alignement fin.}

\section{Bouton \bouton{Continue with mapping}}\label{btn:review_continue}
En bas à droite, après \bouton{Close}. Il ferme la revue et ouvre l'étape~6. La carte, longue, n'y démarre pas toute seule : elle se lance avec \bouton{Build the map\ldots}. Les points signalés, et les points réalignés pas encore confirmés, sont écartés de la carte. Tant que des alignements sont en calcul, la revue ne peut pas être fermée : attendez-les (ou \bouton{Stop}).



%==========================================================
\chapter{Step 6 -- Cartographie et géoréférencement}\label{ch:step6}
\section{Principe}
Les images multispectrales produites à l'étape~4 sont des images \og à plat \fg~: elles ne sont pas positionnées sur le terrain.

\figurePaireIRD{fig:step6_ortho}
{step5_VIS_plan.png}{Image corrigée des distorsions de l'optique (DJI Mavic Air 2 et calibrage de Zhang).}
{step5_VIS_ortho.png}{Image Orthorectifiée avec un modèle de terrain au pas de $10\,m$.}
{L'effet de l'orthorectification est bien visible sur cette image où le terrain présente une forte déclivité (voir figure \ref{fig:step6_topo}).}

 L'étape~6 les rend exploitables dans un système d'information géographique (SIG) comme QGIS. Elle propose trois traitements indépendants, un par bouton :
\begin{itemize}
	\item \bouton{Build the map\ldots} (le bouton principal) : assemblage de l'ensemble des images retenues en une seule \emph{orthophoto} à 4 bandes (rouge, vert, bleu, proche infrarouge) avec le logiciel libre OpenDroneMap (ODM), puis le modèle de terrain et les courbes de niveau ;
	\item \bouton{Terrain model and contours\ldots} : construction à nouveau du MNT et des \emph{courbes de niveau}, par exemple avec un autre pas de grille, à partir d'altitudes fournies par un service en ligne (IGN) ;
	\item \bouton{Georeference images\ldots} : \emph{géoréférencement} et \emph{orthorectification} d'images isolées. Chaque image est projetée sur un modèle numérique de terrain (MNT) puis enregistrée au format GeoTIFF en coordonnées UTM.
\end{itemize}
Ces traitements s'appuient sur les métadonnées de vol enregistrées dans les fichiers \texttt{.exif} \og compagnons \fg~(position GPS, altitude, angles de lacet, tangage et roulis, voir étape~1).

\figureIRD{step5_window.png}{Fen\^etre de l'étape~6.}{0.6}

La fenêtre est organisée comme celle de l'étape~3. De haut en bas : le nom de l'étape et de la mission ; les étapes de la carte (\texttt{Image selection}, \texttt{Orthophoto (OpenDroneMap)}, \texttt{Terrain model and contours} : $\circ$ à faire, $\bullet$ en cours, \checkmark{} fait, $\times$ échec, avec un détail sous chacune) ; une zone de messages (journal) ; les trois boutons, \bouton{Build the map\ldots} en premier ; une barre de progression précédée d'un intitulé qui indique l'opération en cours ; une zone d'aperçu d'image ; et les boutons \bouton{?} et \bouton{Close}, puis \bouton{Continue with archiving} une fois la carte construite.

Si la carte a déjà été construite, les étapes sont déjà cochées et l'orthophoto s'affiche dès l'ouverture de la fenêtre. \bouton{Terrain model and contours\ldots} reste grisé tant que la carte n'existe pas.

\textbf{Prérequis}
\begin{itemize}
	\item l'étape~4 doit avoir été exécutée pour les images à traiter (dossier \texttt{ImgIRdrone\textbackslash imgMultispectral}) ;
	\item pour \bouton{Build the map\ldots} : le logiciel \emph{Docker Desktop} doit être installé et démarré ;
	\item pour \bouton{Terrain model and contours\ldots} (et pour l'altitude du sol lors du géoréférencement) : une connexion Internet est nécessaire pour interroger le service d'altitudes.
\end{itemize}

\section{Bouton \bouton{Georeference images\ldots}}\label{btn:step6_georef}
Ce bouton géoréférence et orthorectifie des images choisies par l'utilisateur.

Un message rappelle les types d'images acceptés, puis une fenêtre de sélection de fichiers s'ouvre. Tous les fichiers choisis doivent être du même type :
\begin{itemize}
	\item des images DNG de la caméra DJI (spectre visible) ;
	\item des images TIFF multispectrales par \textbf{groupes de quatre} : \texttt{IRD\_XXXX\_1.tif} à \texttt{IRD\_XXXX\_4.tif} (rouge, vert, bleu, proche infrarouge). Il faut sélectionner les quatre fichiers ; un groupe incomplet est rejeté avec un avertissement.
\end{itemize}

Dans une mission visible seul, le message et la fenêtre de sélection ne proposent que les images DNG DJI.

\subsection*{Décalages de géoréférencement}
Une fois les images choisies, une fenêtre demande trois corrections, toutes égales à \textbf{0} par défaut :
\begin{itemize}
	\item le \emph{décalage de lacet} (en degrés), qui corrige une erreur systématique du cap du drone ;
	\item le \emph{décalage UTM est} (en mètres) ;
	\item le \emph{décalage UTM nord} (en mètres).
\end{itemize}

Ces valeurs s'appliquent à toutes les images traitées lors de ce lancement. Les décalages sont utiles car les données GPS
\footnote{La précision de positionnement pourrait être grandement améliorée à l'aide d'une balise RTK mais ce type de matériel est onéreux.} et d'attitude du drone DJI ont une précision limitée, tout particulièrement si l'étalonnage du GPS et/ou le calibrage de l'orientation de la caméra n'ont pas été fait avant le décollage. Ces valeurs \og d'offset \fg ne peuvent être déterminer qu'expérimentalement, par exemple en utilisant QGIS pour affiner le positionnement. 

 Le bouton \emph{Cancel} annule le géoréférencement.

\subsection*{Ce que fait le logiciel pour chaque image}
\begin{enumerate}
	\item Les images DNG sont converties en TIFF avec RawTherapee. Pour les images multispectrales, les quatre bandes sont normalisées entre 0 et 1 puis empilées en une image à 4 canaux.
	\item Les métadonnées de vol sont lues dans le fichier \texttt{.exif} compagnon (créé si nécessaire). Elles donnent la position du drone, son altitude au-dessus du sol et ses angles d'attitude. Si l'altitude du point de décollage est inconnue, elle peut être demandée à l'utilisateur.
	\item L'emprise de l'image au sol est estimée à partir de la géométrie de la caméra, puis un maillage du terrain est construit. Les altitudes des points du maillage sont demandées à l'IGN.
	\item Chaque point du maillage est projeté dans l'image. L'image est rééchantillonnée sur une grille régulière au sol (orthorectification).
	\item Le résultat est enregistré au format GeoTIFF (coordonnées UTM).
\end{enumerate}

\subsection*{Fichiers produits}
Les fichiers sont créés dans le sous-dossier \texttt{geo\_ref} du dossier contenant les images sélectionnées. Un fichier existant n'est jamais écrasé : un suffixe \texttt{\_1}, \texttt{\_2}, \dots~est ajouté au nom (utile par exemple si le fichier précédent est ouvert dans QGIS).
\begin{itemize}
	\item \texttt{*\_georef.tif} : l'image orthorectifiée, à ouvrir dans QGIS. Pour un quadruplet multispectral le nom est \texttt{IRD\_XXXX\_0\_georef.tif} et l'image comporte 4 bandes (Red, Green, Blue, NIR) ;
	\item \texttt{*\_MNT.png} et \texttt{*\_Grid.png} : graphiques de contrôle du terrain et du maillage ;
	\item \texttt{*\_dtm\_map.png} : carte 2D ombrée du MNT sous l'image, avec courbes de niveau, quadrillage UTM et points d'altitude IGN 	(même présentation que pour le mapping).
	\end{itemize}
	Pendant le calcul la fenêtre ne répond plus : c'est normal. L'image en cours de traitement s'affiche dans la zone d'aperçu et le journal détaille l'avancement.

	La précision du placement dépend de la qualité des métadonnées de vol (GPS, lacet) et de celle du MNT. Un éventuel décalage systématique peut être corrigé grâce à la fenêtre des décalages (voir ci-dessus).
	
\figureIRD{step5_run_Geo_ref.png}{Fenêtre pendant le géoréférencement d'une image. C'est l'image de départ non orthorectifiée qui est affichée.}{0.6}


\section{Bouton \bouton{Build the map\ldots}}\label{btn:step6_mapping}
Ce bouton assemble les images alignées de l'étape~4 en une orthophoto unique de toute la zone survolée, grâce au logiciel OpenDroneMap (ODM). ODM s'exécute dans un conteneur \emph{Docker} : l'utilisateur n'a rien à configurer, mais Docker Desktop doit être démarré (voir plus bas).

\subsection{Choix des images}\label{btn:step6_selection}
En cliquant sur le bouton, une fenêtre de choix des critères s'ouvre. Une note en bas rappelle que le calcul peut être long et utiliser plusieurs gigaoctets de mémoire et d'espace disque. \bouton{Build the map} le lance aussitôt ; \bouton{Cancel} ferme la fenêtre sans rien lancer. Une fois une carte construite, la fenêtre propose les réglages de la dernière exécution (sinon les valeurs par défaut ci-dessous ; le best mapping choisi à nouveau aux étapes~2 ou~4 après cette exécution l'emporte). Les images peuvent venir de tout le vol, pas seulement de la phase de survol de l'étape~2.

\begin{itemize}
	\item \texttt{Maximum |dt| (NIR rank 1)} : écart temporel maximal admis entre l'image VIS et l'image NIR la plus proche. C'est le même critère que le \texttt{Maximum VIS/NIR time difference} de l'étape~4. Valeur par défaut : $0.5\,s$ (plage $0.05$ à $6\,s$).
	\item \texttt{Best mapping: select by image overlap (recommended)} : choix par défaut. Ne garde que les images dont les empreintes au sol se recouvrent de la valeur indiquée (\texttt{Overlap between successive images}, par défaut la valeur choisie à l'étape~2 ou~4, $75\,\%$ sinon). C'est la même sélection qu'aux étapes~2 et~4 (§\ref{btn:step2_best_mapping}). Si l'étape~4 n'a aligné que les images utiles à la carte, elles sont toutes gardées ; si elle a tout aligné, seules les images utiles à la carte sont données à ODM. Une image déclarée mal alignée dans la revue est écartée et la suivante prend sa place.
	\item \texttt{Select one image every n eligible images} : permet de ne retenir qu'une image sur $n$ ($n$ de 1 à 4, valeur par défaut 3) afin de limiter le temps de calcul et le recouvrement. C'est le choix par défaut quand le best mapping a été décoché à l'étape~2 ou~4.
	\item \texttt{All the aligned images (every point of the review not flagged)} : ignore les critères ci-dessus : toutes les images alignées par l'étape~4 (ou depuis la revue) et non signalées sont cartographiées, quels que soient leur écart de temps et le recouvrement. Le plus d'images, la cartographie la plus longue. Ce choix est proposé à nouveau la fois suivante s'il a été utilisé.
\end{itemize}

\figureIRD{step5_odm_selection.png}{Fenêtre des critères de sélection des images pour le mapping.}{0.4}

À droite, comme à l'étape~4, la carte du vol montre les images retenues pour la carte, en vert avec leurs empreintes au sol (leur recouvrement se voit), les images éligibles écartées par la sélection en gris, et les autres en petits points. Elle est mise à jour dès qu'un réglage change ($|dt|$, recouvrement, $n$, choix) ; le compte sous la carte donne le nombre d'images retenues parmi les éligibles. Molette pour zoomer, glisser pour se déplacer, double-clic pour revoir tout le vol.

\textbf{OpenDroneMap : version et GPU.} Sous les choix, le cadre \texttt{OpenDroneMap} indique, après une vérification de quelques secondes (Docker, Internet), la version d'ODM installée sur l'ordinateur et si une version plus récente a été publiée sur Docker Hub :
\begin{itemize}
	\item \texttt{Download the newer version before the map} : affichée seulement quand une version plus récente existe, décochée par défaut. Cochée, la dernière image d'ODM est téléchargée (\texttt{docker pull}, de quelques centaines de mégaoctets à quelques gigaoctets) avant le calcul. Une nouvelle version peut donner une carte un peu différente : elle n'est donc jamais installée sans qu'on le demande. Sans Internet, le cadre indique que la dernière version n'a pas pu être vérifiée et la version installée est utilisée.
	\item \texttt{Use the NVIDIA GPU (faster image matching)} : disponible, et cochée, seulement quand l'ordinateur a une carte graphique NVIDIA \emph{et} que Docker peut la donner à un conteneur (le logiciel lance un conteneur de test avec \texttt{--gpus all}). ODM utilise alors sa version GPU (\texttt{opendronemap/odm:gpu}, environ 3\,Go téléchargés la première fois), qui détecte et apparie les points caractéristiques des images sur la carte graphique ; le reste du calcul est le même. Sinon la case est grisée et la raison est donnée (pas de carte NVIDIA ; sous Linux, le \emph{NVIDIA Container Toolkit} manque ; sous Windows, Docker Desktop doit utiliser WSL~2 avec un pilote NVIDIA récent) ; son info-bulle donne la réponse de Docker. Si le GPU n'est pas utilisable au lancement du calcul, ODM tourne sur le processeur et le journal le dit.
\end{itemize}
Le journal donne la version d'ODM utilisée, et si elle a tourné sur le GPU ou sur le processeur.


Les images retenues (\emph{eligible images}) sont les shoot points qui vérifient simultanément :
\begin{enumerate}
	\item ils n'ont pas été signalés mal alignés dans la revue, étape~5 (clé \texttt{"failed": false}, voir §\ref{btn:review_decide}) ;
	\item l'écart $|dt|$ avec l'image NIR la plus proche est inférieur ou égal à la tolérance ;
	\item l'image VIS correspondante appartient à la phase de vol de la mission choisie à l'étape~2 (les images de décollage et de synchronisation sont donc exclues).
\end{enumerate}
Parmi elles, le best mapping garde les images utiles à la carte (ou, avec l'autre critère, une image éligible sur $n$).

\subsection{Lancement du calcul}
\bouton{Build the map} lance le calcul sans autre question (la sélection est écrite dans le journal) :
\begin{enumerate}
	\item le logiciel vérifie que Docker Desktop est disponible. Sinon un message invite à le démarrer puis à choisir \texttt{Retry}. ODM est ensuite téléchargé s'il n'est pas encore sur l'ordinateur, ou mis à jour si on l'a demandé (voir plus haut) ;
	\item la liste des images retenues est enregistrée dans \texttt{mapping\_MULTI\textbackslash img\_best\_mappping.json} ;
	\item la résolution de l'orthophoto est estimée automatiquement : c'est la taille au sol d'un pixel (GSD, en cm) obtenue à partir de l'altitude de vol au-dessus du sol et de la focale de la caméra, médiane sur les images retenues, arrondie au dixième de cm supérieur ;
	\item les images TIFF retenues sont rendues accessibles à ODM dans \texttt{mapping\_MULTI\textbackslash images} par des \emph{liens physiques} (hard links) : aucune image n'est copiée et l'espace disque n'est donc pas doublé. Si le disque ne gère pas les liens physiques (disque externe en exFAT ou FAT32, certains partages réseau), les images sont copiées, ce qui utilise davantage d'espace disque ; le journal l'indique ;
	\item ODM est lancé. Le journal de la fenêtre affiche les messages principaux ; la totalité de la sortie est enregistrée dans un fichier \texttt{mapping\_MULTI\textbackslash odm\_run\_<date>.log} ;
	\item à la fin, l'orthophoto est contrôlée (4 bandes, géoréférencée), les résultats utiles sont archivés dans \texttt{mapping\_result} et les dossiers de travail d'ODM, très volumineux, sont supprimés.
\end{enumerate}
Pendant le calcul la barre de progression est indéterminée, les boutons et la fermeture de la fenêtre sont inactifs. À l'issue du calcul la durée est indiquée dans le journal.

\textbf{Résultats}\label{sec:step6_mapping_result}

Les résultats sont dans \texttt{mapping\_MULTI\textbackslash mapping\_result} :
\begin{itemize}
	\item \texttt{odm\_orthophoto.tif} : l'orthophoto à 4 bandes (rouge, vert, bleu, proche infrarouge), en UTM ;
	\item \texttt{report.pdf} : rapport de qualité produit par ODM ;
	\item \texttt{odm\_25dmesh.ply} : maillage 2,5D du terrain ;
	\item \texttt{shots.geojson} (et \texttt{shots.qml}) : position estimée de chaque prise de vue, avec son numéro de shoot point, à ouvrir dans QGIS ;  \texttt{shots.qml} est un fichier de mise en forme.
	\item des fichiers techniques (\texttt{ortho.png}, informations de géoréférencement).
\end{itemize}
Un aperçu en couleurs (bandes rouge, vert, bleu) de l'orthophoto s'affiche ensuite dans la fenêtre, puis le calcul du MNT et des courbes de niveau démarre automatiquement (c'est équivalent au bouton \bouton{Terrain model and contours\ldots}, voir \ref{btn:step6_dtm}).

\figureIRD{step5_orthophoto_preview.png}{Aperçu de l'orthophoto dans la fenêtre du Step~6.}{0.6}

\subsection{Précautions}
\begin{itemize}
	\item Les résultats existants ne sont \textbf{jamais écrasés}. Quand une carte est déjà construite, elle est conservée : son dossier est renommé \texttt{mapping\_MULTI\_<date de son exécution>} et la nouvelle carte est construite dans un nouveau \texttt{mapping\_MULTI} (avec la même configuration de caméra) ; la fenêtre de choix des images l'indique. Si la nouvelle exécution échoue, la carte précédente est remise dans \texttt{mapping\_MULTI} et l'exécution échouée, avec son journal, est conservée dans \texttt{mapping\_MULTI\_failed\_<date>}. L'étape~7 conserve ces dossiers.
	\item Docker doit utiliser le moteur \og desktop-linux \fg~de Docker Desktop. Le premier lancement télécharge l'image ODM (plusieurs gigaoctets) : prévoir une connexion Internet.
	\item Si certaines images retenues sont incomplètes (quadruplet TIFF manquant ou de tailles différentes), le calcul est refusé et le journal indique les fichiers en cause.
\end{itemize}

\section{Bouton \bouton{Terrain model and contours\ldots}}\label{btn:step6_dtm}
Ce bouton construit un modèle numérique de terrain (MNT) et des courbes de niveau sur l'emprise de l'orthophoto. Il est lancé \emph{automatiquement} à la fin d'un mapping réussi (avec un pas de grille de 10\,m) ; ce bouton permet de le refaire, par exemple avec un autre pas. Il nécessite l'orthophoto \texttt{mapping\_MULTI\textbackslash mapping\_result\textbackslash odm\_orthophoto.tif}.

Une fenêtre demande la valeur \texttt{Altitude grid step}, c'est-à-dire l'espacement (en mètres) des points de la grille régulière dont les altitudes sont demandées à l'IGN (service d'altitudes de l'Institut national de l'information géographique et forestière) ; en cas d'indisponibilité, le service OpenTopoData est utilisé avec une moindre précision en France. Plus le pas est petit, plus le nombre de requêtes est grand. Le pas par défaut est de 10\,m (plage 1 à 100\,m). Le nombre de points est limité à 5000 : au-delà, le logiciel augmente automatiquement le pas. \bouton{Build} lance le calcul, \bouton{Cancel} ne lance rien.

\subsection*{Étapes du calcul}
\begin{enumerate}
	\item lecture de l'emprise de l'orthophoto (UTM) ;
	\item construction d'une grille régulière qui couvre l'emprise ;
	\item requêtes d'altitude (aucun point manquant n'est toléré : sinon le calcul échoue et un message l'indique) ;
	\item interpolation par une surface lisse (spline quadratique) puis échantillonnage sur un raster à 1\,m ;
	\item export du MNT, des courbes de niveau (équidistance de 1\,m, courbes maîtresses tous les 10\,m), d'une carte 2D ombrée et d'une vue 3D.
\end{enumerate}
En cas d'échec, l'orthophoto n'est pas affectée.

\subsection{Fichiers produits}
Les fichiers sont créés à côté de l'orthophoto, dans \texttt{mapping\_result} (jamais écrasés, un suffixe numérique est ajouté si besoin) :
\begin{itemize}
	\item \texttt{ortho\_dtm.tif} : MNT (altitude en mètres, UTM) ;
	\item \texttt{ortho\_contours.geojson} et \texttt{ortho\_contours.qml} : courbes de niveau et leur style pour QGIS ;
	\item \texttt{ortho\_dtm\_map.png} : carte 2D ombrée avec courbes de niveau, position des points de la grille (point noir) et quadrillage UTM (bleu clair);
	\item \texttt{ortho\_dtm\_3d.png} : surface 3D ;
	\item \texttt{ortho\_dtm\_grid.json} : points de la grille avec leur altitude et la source des altitudes.
\end{itemize}

\figurePaireIRD{fig:step6_topo}
{step5_dtm_map.png}{Carte 2D du MNT avec courbes de niveau et quadrillage UTM.}
{step5_dtm_3d.png}{Surface 3D du MNT.}
{Représentation topographique du terrain à partir du modèle numérique de terrain déduit des altimétries IGN.}

\subsection{Ouverture dans QGIS}
Dans la fenêtre de la carte QGIS, faire glisser l'orthophoto \texttt{odm\_orthophoto.tif}, les courbes de niveau \texttt{ortho\_contours.geojson} et les points des prises de vue \texttt{shots.geojson}    (les styles \texttt{.qml} homonymes sont appliqués automatiquement). \\

Quelques actions simples avec une orthophotographie à 4 bandes spectrales red,green,blue,nir

\begin{itemize}
	\item Image couleur dans le spectre visible : clic droit sur la couche> \underline{P}ropriétés>Symbologie puis associer les bandes 1, 2, 3 au rouge, vert, bleu pour une vue en couleurs naturelles. C'est l'option d'affichage par défaut d'une image multispectrale.
	\item Image fausse couleur VIR : clic droit>\underline{P}ropriétés>Symbologie puis associer les bandes 4~,1~,2 au rouge, vert, bleu. 
	\item Carte de l'indice NDVI\footnote{L'indice NDVI pour \textit{Normalized Difference Vegetation Index} mesure la vigueur et la densité de la végétation. On détermine cet indice par la formule simple : $NDVI = \frac{NIR - RED}{NIR + RED}$} : \underline{R}aster>Calculatrice raster> entrez la formule 
	\begin{equation}
		\frac{(VIS@4 - VIS@1)}{(VIS@4 + VIS@1)}
	\end{equation}
	puis double clic sur la couche et choisir dans symbologie type de rendu : Pseudo-couleur à bande unique> Palette de couleur: NDVI. Le résultat est montré sur la figure \ref{fig:step6_mapping_NDVI}
	\item Couche NIR. Sélectionner dans symbologie comme première bande NIR,  puis dans symbologie type de rendu :  Niveau de gris à bande unique.
\end{itemize}

\figurePaireIRD{fig:step6_mapping_1}
{step5_topo_QGIS.png}{Représentation fine de la topographie. Équidistance des isolignes : $10\,m$.}
{step5_mapping_VIS.png}{Mapping en spectre visible.}
{Représentation des données IRdrone avec le logiciel QGIS sur fond OpenStreetMap.}

\figurePaireIRD{fig:step6_mapping_NDVI}
{step5_mapping_VIR.png}{Image multispectrale VIR.  $\qquad\qquad\qquad\quad$ Décalage spectral $[Nir|R|G] \longrightarrow [R|G|B]$.}
{step5_mapping_NDVI.png}{Carte de l'indice NDVI.}
{Représentation des données IRdrone avec le logiciel QGIS sur fond Google Satellite.}


\section{Bouton \bouton{Close}}\label{btn:step6_close}
Ferme la fenêtre et rend la main au dialogue principal. La fermeture est refusée tant qu'un calcul ODM ou MNT est en cours. Le bouton \bouton{?} ouvre cette aide à la page de ce chapitre.

\section{Bouton \bouton{Continue with archiving}}\label{btn:step6_continue}
Apparaît en bas à droite, après \bouton{Close}, une fois la carte construite. Il ferme la fenêtre et ouvre l'étape~7, qui indique l'espace libéré par l'archivage. Rien n'est supprimé avant votre confirmation.

%==========================================================
\chapter{Step 7 -- Archivage de la mission}\label{ch:step7}
L'étape~7 est l'archivage de la mission. Elle se fait \emph{une fois l'exploitation des images terminée}. Elle nettoie le dossier de la mission \emph{sur place} :
les nombreux fichiers techniques (images intermédiaires, dossiers de travail, copies d'images de la synchronisation, etc.) sont supprimés,
et seuls les éléments essentiels pour pouvoir, si nécessaire, traiter de nouveau la mission sont conservés. Sur les missions testées, le volume est réduit
d'un facteur 2 à 5 (par exemple de 131\,Go à 25\,Go pour une mission de 600 images VIS DJI).

\begin{center}
\fbox{\parbox{0.9\linewidth}{\small\textbf{Attention : opération irréversible.} Les fichiers supprimés ne vont pas à la corbeille de Windows ; ils sont définitivement perdus.
Si vous souhaitez conserver une copie de la mission, d'un dossier \texttt{mapping\_MULTI} particulier ou de tout autre fichier, \emph{copiez-les ailleurs avant de lancer l'archivage}.
Tout élément du dossier de mission qui ne figure pas dans la liste ci-dessous est supprimé, y compris les copies de dossiers (par exemple \texttt{mapping\_MULTI - Copie}).}}
\end{center}

Comme pour les autres étapes, la mission archivée est la mission ouverte dans la fenêtre principale.

\section{Ce qui est conservé}\label{btn:step7_rules}
La comparaison des noms de fichiers ne tient pas compte des majuscules et minuscules (\texttt{.dng} et \texttt{.DNG} sont équivalents).
\begin{description}
  \item[Racine de la mission] les dossiers \texttt{AerialPhotography}, \texttt{cameras}, \texttt{FlightAnalytics}, \texttt{ImgIRdrone}, \texttt{mapping\_MULTI} (et les cartes précédentes \texttt{mapping\_MULTI\_<date>} mises de côté par l'étape~6), \texttt{Synchro}
        et le fichier \texttt{desktop.ini}. Tout le reste est supprimé.
  \item[\texttt{AerialPhotography}] les dossiers \texttt{VIS} et \texttt{NIR} et le fichier \texttt{shoot\_points.json}. Dans \texttt{VIS} et \texttt{NIR} : les images originales
        \texttt{.dng} avec leurs fichiers compagnons \texttt{.exif}, ainsi que \texttt{time\_line.json} et \texttt{transfer\_info\_VIS\_dng.json} (respectivement \texttt{transfer\_info\_NIR\_dng.json}).
        Les autres fichiers (notamment les \texttt{*\_RawTherapee.tif}) sont supprimés.
  \item[\texttt{cameras} et \texttt{FlightAnalytics}] laissés tels quels.
  \item[\texttt{ImgIRdrone}] les images \texttt{VIS\_XXXX.jpg} et \texttt{VIR\_XXXX\_local.jpg}, les fichiers \texttt{VIS\_XXXX\_motion\_model.npy}, les images
        \texttt{NDVI\_XXXX\_local.jpg} et \texttt{NIR\_XXXX\_local.jpg} (si elles existent), \texttt{shoot\_points.json} et \texttt{shoot\_points\_failed\_summary.json}.
        Le dossier volumineux \texttt{imgMultispectral} et les dossiers de \og redo \fg{} individuels sont supprimés.
  \item[\texttt{mapping\_MULTI}] laissé tel quel : il contient les éléments directement exploitables dans QGIS et a déjà été allégé à l'étape~6.
  \item[\texttt{Synchro}] seuls \texttt{check\_time\_alignment.png}, \texttt{check\_time\_before\_alignment.png} et \texttt{desktop.ini} sont conservés ; les dossiers \texttt{VIS}, \texttt{NIR} et \texttt{Check\_ARUCO} sont supprimés.
\end{description}

\paragraph{Pourquoi ce choix ?} Le décalage temporel trouvé à l'étape~3 et la liste des paires d'images (\texttt{AerialPhotography\textbackslash shoot\_points.json}) résument la synchronisation :
les images de \texttt{Synchro} ne sont plus utiles. Les fichiers \texttt{VIS\_XXXX\_motion\_model.npy} mémorisent l'alignement VIS--NIR de chaque \emph{shoot point}, y compris les alignements
rattrapés par la procédure de \og redo \fg{} (étape~4). Pour reprendre la mission, il suffit donc de refaire l'étape~4, sans perdre ces corrections, à partir des images originales conservées.
Les images \texttt{.jpg} conservées permettent aussi de retrouver rapidement une image pour illustrer un rapport.

\figureIRD[fig:step7]{step6_window.png}{Fenêtre de l'archivage : l'espace libéré par l'archivage s'affiche dès l'ouverture de la fenêtre.}{0.6}

\section{Ce que l'archivage libère (simulation)}\label{btn:step7_simulate}
Dès l'ouverture de la fenêtre, le dossier est analysé \emph{sans rien supprimer} (une simulation). Le cadre sous l'avertissement donne la réponse :
\begin{itemize}
  \item \texttt{12.3 GB can be freed}, avec, en grand, la taille du dossier de mission avant et après l'archivage (par exemple \texttt{15.1 GB → 2.8 GB}) et le nombre de fichiers qui seront supprimés. Le détail, dossier par dossier (fichiers à supprimer, fichiers conservés, dossiers supprimés), est dans la zone en dessous ;
  \item \texttt{Archiving impossible}, avec les fichiers indispensables qui manquent (voir plus bas) : rien ne peut être supprimé ;
  \item \texttt{Mission already archived}, avec la date de l'archivage et, en grand, la taille du dossier ;
  \item \texttt{Nothing to delete}, quand le dossier ne contient que ce que l'archivage conserve.
\end{itemize}
Vérifiez que la liste correspond à ce que vous attendez avant d'archiver.

\section{Bouton \bouton{Archive and free \ldots}}\label{btn:step7_archive}
Le bouton indique l'espace qui sera libéré (par exemple \bouton{Archive and free 12.3 GB\ldots}). Il est grisé pendant l'analyse du dossier, et quand rien ne peut être supprimé. Le dossier est analysé à nouveau au clic, puis un message de confirmation est demandé à l'utilisateur  comme le montre la figure \ref{fig:step7_confirm_archving}. Le message indique aussi l'espace disque qui sera libéré et précise que la suppression sera définitive. Deux boutons \texttt{OK} et \texttt{Cancel} permettent de confirmer ou d'abandonner. Par sécurité et pour éviter un clic trop rapide sur \texttt{Ok} l'utilisateur  doit recopier exactement le même nom du dossier de mission que celui indiqué (en respectant la casse). 

\figureIRD[fig:step7_confirm_archving]{step6_confirm_archiving.png}{Message de confirmation de la demande d'archivage.}{0.3}

 Si clic OK  alors IRDrone lance l'archivage réel :

\begin{enumerate}
  \item analyse le dossier et vérifie que les fichiers indispensables existent (\texttt{mission\_parameters.json}, \texttt{shoot\_points.json}, \texttt{time\_line.json}, \texttt{transfer\_info\_*\_dng.json}, etc.).
        S'il en manque, l'archivage est \textbf{refusé} et la liste des fichiers manquants est affichée ; rien n'est supprimé. Les dossiers manquants ne sont jamais recréés ;
  \item affiche le nombre de fichiers et le volume qui seront supprimés, puis demande confirmation : il faut \emph{saisir le nom du dossier de la mission} ;
  \item supprime les fichiers. La barre de progression indique l'avancement ;
  \item si un fichier est verrouillé (ouvert dans un autre logiciel, par exemple QGIS), l'erreur est notée et la suppression continue. Fermez alors les logiciels concernés et relancez l'archivage ;
  \item écrit à la racine de la mission les fichiers \texttt{archive\_manifest.json} et \texttt{archive\_log.txt}.
\end{enumerate}

Une mission déjà archivée est détectée et n'est pas archivée une seconde fois.


 Une icône particulière attribuée au dossiers archivés et permet de voir d'un coup d'\oe il qu'une mission est archivée.(figure \ref{fig:Air_Mission_folder}).


\paragraph{Archivage et avancement du traitement.} A la fin de chaque étape le dossier de mission reflète l'état de l'avancement du traitement de la mission. 

L'archivage ne doit être lancé qu'à la toute fin de l'étude. Notez que les images originales de la mission (converties au format dng pour la SJCam) sont systématiquement archivée. Vous conservez ainsi les données du terrain. Il n'est donc pas nécessaire d'archiver les images originales des cartes SD du DJI et de la SJCam (mais vous pouvez aussi les archiver par sécurité).

\section{Bouton \bouton{Close}}\label{btn:step7_close}
Ferme la fenêtre et rend la main au dialogue principal. Le bouton \bouton{?} ouvre cette aide à la page de ce chapitre.

%==========================================================
\chapter{Missions visible seul}\label{ch:visible_only}
\section{Sans la caméra infrarouge}\label{sec:visible_only}
Une mission peut être volée avec la seule caméra du drone, sans la SJCam : le type \bouton{Visible only (drone DNG)} se choisit en haut de la première page de l'étape~1 (section~\ref{btn:step1_mission_type}). Il est enregistré dans la mission (\texttt{"mission type": "visible"} dans \texttt{FlightAnalytics/mission\_parameters.json} ; les missions sans cette clé sont multispectrales). Les étapes sont les mêmes, moins ce qui concerne l'infrarouge :
\begin{itemize}
	\item \textbf{Étape 1} : seule la carte SD du drone est nécessaire ; la copie a trois étapes (copie, EXIF, altitude du sol). Le dossier de la mission n'a pas de dossier \texttt{AerialPhotography/NIR} (\texttt{Synchro} ne contient que les vignettes de la ligne de temps de l'étape~2).
	\item \textbf{Étape 2} : la ligne de temps ne montre que les images visibles ; il n'y a pas de cible tournante à chercher ni de phase de synchronisation : seule la phase de relevé est vérifiée. Le best mapping fonctionne comme pour une mission multispectrale. \bouton{Continue with development} enregistre le relevé et les points de prise de vue visibles (\texttt{AerialPhotography/shoot\_points.json}, un point par image visible, sans infrarouge) et ouvre l'étape~4.
	\item \textbf{Étape 3} : inutile ; sa tuile indique \og Not needed \fg~et \bouton{Continue} passe directement de l'étape~2 à l'étape~4.
	\item \textbf{Étape 4} (\og Develop \fg) développe les images DNG visibles du relevé (ou de la sélection du best mapping) : couleurs, correction de la distorsion de l'objectif et même recadrage que les images multispectrales, pour que le modèle de caméra d'OpenDroneMap soit le même. Chaque image donne \texttt{ImgIRdrone/VIS\_XXXX.jpg} (revue) et les trois bandes \texttt{ImgIRdrone/imgMultispectral/IRD\_XXXX\_1..3.tif} (rouge, vert, bleu) pour la carte.
	\item \textbf{Étape 5} : la revue montre les images visibles développées ; il n'y a rien à réaligner. \bouton{Good: keep for the map} (touche \texttt{Y}) marque l'image vérifiée et montre la suivante à vérifier ; une image floue ou mal exposée est signalée (touche \texttt{F}) : elle est écartée de la carte. Les points que l'étape~4 n'a pas développés (gris sur la carte) se développent d'un clic droit sur la carte.
	\item \textbf{Étape 6} : OpenDroneMap construit une orthophoto à trois bandes (rouge, vert, bleu) au lieu de quatre. Il n'y a ni couche en fausses couleurs infrarouges, ni VIR, ni NDVI ; le modèle de terrain est calculé comme d'habitude.
	\item \textbf{Étape 7} : l'archive garde les mêmes fichiers, sans ceux de l'infrarouge.
\end{itemize}

%==========================================================
\appendix
\chapter{Alignement des images VIS et NIR (Step 4)}\label{ann:alignement}\label{sec:ann_algorithme}
Cette annexe explique en détail comment IRdrone superpose l'image proche infrarouge à l'image visible. Elle s'adresse à un lecteur de formation scientifique (algèbre linéaire, un peu de traitement du signal) qui n'est pas spécialiste de vision par ordinateur. Les premières sections donnent les bases de géométrie et de traitement du signal, les suivantes décrivent chaque étape de l'algorithme avec les traces de diagnostic d'alignements réels, puis viennent les méthodes essayées puis abandonnées, l'historique du développement et les limites connues. La dernière section (§\ref{sec:ann_theorique}) traite de l'origine physique du désalignement des deux caméras (angles de \og pitch \fg{} et de \og yaw \fg{} théoriques).

Le code se trouve dans \texttt{automatic\_registration.py} (fonctions \texttt{align\_raw} et \texttt{coarse\_alignment}), dans le dossier \texttt{registration} (\texttt{rigid.py} : pyramide, \texttt{cost.py} : représentation et coût, \texttt{warp\_flow.py} : rééchantillonnage) et dans \texttt{irdrone/semi\_auto\_registration.py} (rotation de la caméra, \texttt{manual\_warp}). Les figures de cette annexe sont produites par \texttt{help/make\_alignment\_figures.py} à partir de deux shoot points :
\begin{itemize}
	\item \textbf{point A} : l'image \texttt{VIS\_0201} du vol du 25 janvier 2022 (Blassac, 90\,m au-dessus du sol, écart de temps entre les deux prises de vue 0,0\,s), un cas typique où l'alignement réussit ;
	\item \textbf{point B} : la paire d'exemple du dépôt (\texttt{samples/\_visible.DNG}, \texttt{samples/\_nir.RAW}, mai 2022), un cas difficile où l'alignement automatique échoue (§\ref{sec:ann_limits}).
\end{itemize}
Les deux ont été alignés en ligne de commande avec l'option \texttt{--debug} (§\ref{sec:ann_traces}), c'est-à-dire avec l'élargissement de champ 1,6 de la ligne de commande au lieu de 1,4 au Step~4.

\section{Le problème}\label{sec:ann_problem}
\subsection{Ce que signifie \og aligner \fg}
Le drone emporte deux caméras (figure~\ref{fig:ann_pair}) : la caméra du DJI Mavic Air~2, qui enregistre la lumière visible (VIS, rouge, vert et bleu), et une SJCam M20 dont le filtre anti-infrarouge a été retiré et remplacé par un filtre qui ne laisse passer que le proche infrarouge (NIR, au-delà d'environ 720\,nm). Les indices de végétation comme le NDVI combinent, \emph{pixel par pixel}, le rouge de l'image visible et le proche infrarouge :
\begin{equation}
	\text{NDVI} = \frac{\text{NIR} - \text{R}}{\text{NIR} + \text{R}}.
\end{equation}
Cette formule n'a de sens que si les deux valeurs d'un pixel proviennent du même point du sol. Aligner (on dit aussi \emph{recaler}) les deux images, c'est trouver, pour chaque pixel de l'image visible, la position du même point du sol dans l'image infrarouge, puis rééchantillonner l'image infrarouge sur la grille de l'image visible. Une erreur de quelques pixels suffit à créer de fausses valeurs d'indice le long de chaque bord contrasté (une haie, une route, l'ombre d'un arbre).

\figureIRD[fig:ann_pair]{align_pair_raw.jpg}{Le point B tel que l'enregistrent les deux caméras. Outre des champs de vision différents et la forte distorsion de la caméra infrarouge, les contrastes diffèrent : la route est sombre dans les deux bandes mais ses bas-côtés ne sont clairs que dans le visible, la végétation est verte et plus sombre que l'herbe sèche dans le visible, claire dans le proche infrarouge.}{1}

\subsection{Pourquoi c'est difficile}
Trois difficultés s'additionnent :
\begin{enumerate}
	\item \textbf{Les deux caméras ne regardent pas dans la même direction.} La caméra visible est maintenue verticale par sa nacelle (gimbal), la caméra infrarouge est fixée au corps du drone et suit ses inclinaisons (§\ref{sec:ann_theorique}). L'écart atteint plusieurs degrés, c'est-à-dire plusieurs centaines de pixels.
	\item \textbf{Les deux optiques sont très différentes.} La caméra infrarouge a un champ très large et une forte distorsion ; les lignes droites apparaissent courbes (§\ref{sec:ann_distortion}).
	\item \textbf{Les deux bandes ne voient pas la même chose.} La végétation, l'eau, le sol et les routes n'ont pas la même réflectance dans le visible et dans le proche infrarouge. Un contraste peut s'inverser d'une bande à l'autre (figure~\ref{fig:ann_pair}, et figure~\ref{fig:step4_pair_VIS_NIR} où la rivière est noire dans l'infrarouge). Les méthodes classiques, qui comparent des intensités ou apparient des points caractéristiques (SIFT, ORB), sont donc peu fiables.
\end{enumerate}
Les deux premières difficultés sont résolues par la \emph{géométrie} : un modèle de chaque caméra, obtenu par calibration, et le fait que les deux images diffèrent essentiellement d'une rotation. La troisième est résolue en \emph{comparant des bords plutôt que des intensités}.

\section{Géométrie d'une caméra}\label{sec:ann_camera}
Le livre de référence est \emph{Multiple View Geometry in Computer Vision} de Hartley et Zisserman~[\hyperref[ann_hz]{1}] (chapitres 2 et 6 pour le modèle de caméra, 7.4 pour la distorsion, 8.4 pour les caméras en rotation). Le livre de Szeliski~[\hyperref[ann_szeliski]{2}] est une introduction plus accessible (chapitre 2).

\subsection{Le modèle sténopé et les coordonnées homogènes}
En première approximation, une caméra est un \emph{sténopé} (pinhole) : tout rayon lumineux passe par le centre optique $C$ avant d'atteindre le capteur. Soient $(X,Y,Z)$ les coordonnées d'un point de la scène dans le repère de la caméra ($Z$ selon l'axe optique). Son image est au pixel
\begin{equation}
	u = f\,\frac{X}{Z} + c_x, \qquad v = f\,\frac{Y}{Z} + c_y,
\end{equation}
où $f$ est la focale \emph{exprimée en pixels} (focale physique divisée par la taille d'un pixel) et $(c_x, c_y)$ le \emph{point principal}, intersection de l'axe optique avec le capteur, proche du centre de l'image. La division par $Z$ est la perspective : un objet deux fois plus loin apparaît deux fois plus petit.

Cette division n'est pas linéaire ; elle le devient en écrivant le pixel en \emph{coordonnées homogènes}, un vecteur de trois nombres défini à un facteur près : $(u, v, 1) \sim (\lambda u, \lambda v, \lambda)$. Alors
\begin{equation}\label{eq:ann_pinhole}
	\lambda \begin{pmatrix} u \\ v \\ 1 \end{pmatrix} = K \begin{pmatrix} X \\ Y \\ Z \end{pmatrix},
	\qquad
	K = \begin{pmatrix} f & 0 & c_x \\ 0 & f & c_y \\ 0 & 0 & 1 \end{pmatrix}.
\end{equation}
La matrice $K$ des \emph{paramètres intrinsèques} (ou matrice de calibration) décrit la caméra seule. Inversement, $K^{-1}(u, v, 1)^T$ est la \emph{direction} du rayon qui arrive sur le pixel $(u,v)$ : une image est un ensemble de directions, la profondeur $Z$ est perdue.

\begin{table}[htbp]\centering
\begin{tabular}{@{}lll@{}}
\hline
 & \textbf{DJI Mavic Air~2 (VIS)} & \textbf{SJCam M20 (NIR)} \\ \hline
Taille de l'image (pixels) & $3992 \times 2992$ & $4600 \times 3448$ \\
Focale $f$ (pixels) & 2898 & 2311 \\
Point principal $(c_x, c_y)$ & (2032, 1511) & (2348, 1725) \\
Champ horizontal & $68^\circ$ & plus de $110^\circ$ \\
Distorsion $k_1, k_2, k_3$ & $0{,}185,\ -0{,}455,\ 0{,}245$ & $-0{,}337,\ 0{,}131,\ -0{,}025$ \\ \hline
\end{tabular}
\caption{Calibrations des deux caméras (\texttt{calibration/DJI\_RAW} et \texttt{calibration/M20\_RAW}, images RAW). À la focale de la caméra visible, une rotation de $1^\circ$ déplace l'image de $f \tan 1^\circ \approx 51$ pixels.}
\label{tab:ann_calib}
\end{table}

\subsection{La distorsion}\label{sec:ann_distortion}
Les objectifs réels s'écartent du modèle sténopé : l'image d'une droite est courbe, de plus en plus vers les bords. IRdrone utilise le modèle polynomial de Brown et Conrady~[\hyperref[ann_brown]{3}], celui d'OpenCV. Avec les coordonnées normalisées $x = X/Z$, $y = Y/Z$ et $r^2 = x^2 + y^2$, le point idéal $(x, y)$ est observé en
\begin{equation}\label{eq:ann_distortion}
	x_d = x\,(1 + k_1 r^2 + k_2 r^4 + k_3 r^6) + 2p_1 xy + p_2 (r^2 + 2x^2),
\end{equation}
et de même pour $y_d$ ; le pixel est alors $K\,(x_d, y_d, 1)^T$. Les coefficients $k_i$ décrivent la distorsion radiale (en barillet si $k_1 < 0$, en coussinet si $k_1 > 0$), $p_1, p_2$ un léger décentrement des lentilles.

La figure~\ref{fig:ann_distortion} donne les ordres de grandeur. Pour la caméra visible la correction ne dépasse pas 40 pixels. Pour la caméra infrarouge elle atteint 210 pixels aux bords gauche et droit et 375 pixels aux coins du champ de la caméra visible, et 800 pixels à $50^\circ$ de l'axe. Son $k_1$ fortement négatif est une distorsion en barillet proche de celle d'un \og fish-eye \fg. Le polynôme n'est valable que dans la zone où des mires de calibration ont été vues : au-delà d'environ $57^\circ$ de l'axe il redécroît, le modèle n'est donc plus inversible, et les bords de l'image infrarouge (plus de $110^\circ$ de champ horizontal) sont hors de son domaine de validité. Ces zones ne sont jamais utilisées, puisque le champ visible n'atteint que $41^\circ$ dans les coins.

\figureIRD[fig:ann_distortion]{align_distortion.png}{Distorsion des deux caméras. À gauche et au centre : un quadrillage de droites au sol, tel qu'il est enregistré sur le capteur (cadre noir) de chaque caméra ; en rouge, le champ de la caméra visible vu par la caméra infrarouge. À droite : décalage radial d'un point en fonction de son angle avec l'axe optique.}{1}

\subsection{La calibration}
Les paramètres intrinsèques $f, c_x, c_y$ et les coefficients de distorsion sont mesurés une fois pour chaque caméra par la méthode de Zhang~[\hyperref[ann_zhang]{4}] : une trentaine de photographies d'un damier plan dans diverses orientations ; les coins des cases sont détectés, et tous les paramètres (plus la position du damier sur chaque photo) sont estimés en minimisant la distance entre les coins observés et les coins prédits par les équations \eqref{eq:ann_pinhole} et \eqref{eq:ann_distortion} (\texttt{calibration/cameracalibration.py}, OpenCV). Cette calibration est fixe : elle ne dépend pas du vol.

\section{Deux caméras qui tournent : l'homographie}\label{sec:ann_homography}
\subsection{Une rotation de la caméra est une homographie de l'image}
Supposons que les deux caméras aient le \emph{même centre optique} et ne diffèrent que d'une rotation $R$ (une matrice orthogonale $3 \times 3$). Le rayon qui arrive sur le pixel $p_N$ de la caméra infrarouge a la direction $K_{NIR}^{-1} p_N$ dans le repère de cette caméra, $R\, K_{NIR}^{-1} p_N$ dans celui de la caméra visible, et arrive donc sur le pixel
\begin{equation}\label{eq:ann_rotation_homography}
	p_V \sim K_{VIS}\, R\, K_{NIR}^{-1}\, p_N = H\, p_N
\end{equation}
de la caméra visible (après correction des distorsions). La matrice $3 \times 3$ $H$, définie à un facteur près, est une \emph{homographie} : la transformation la plus générale qui envoie les droites sur des droites (Hartley et Zisserman~[\hyperref[ann_hz]{1}], §2.3 et §8.4). Elle a 8 degrés de liberté. Point remarquable : l'équation \eqref{eq:ann_rotation_homography} ne fait pas intervenir la profondeur $Z$ : quel que soit le relief, une rotation pure de la caméra est une homographie de l'image. C'est le principe des panoramas réalisés en tournant sur soi-même.

La rotation est paramétrée par trois angles : le \emph{lacet} (yaw, rotation autour de l'axe vertical de l'image, qui la déplace horizontalement), le \emph{tangage} (pitch, qui la déplace verticalement) et le \emph{roulis} (roll, rotation autour de l'axe optique, qui fait tourner l'image autour de son centre). Dans le code, $R$ est construite à partir du vecteur de rotation $(-\text{pitch}, \text{yaw}, \text{roll})$ (formule de Rodrigues, \texttt{manual\_warp}). Aux petits angles, un lacet $\theta$ déplace le centre de l'image de $f \tan\theta$ pixels, soit environ 51 pixels par degré ; un roulis $\varphi$ déplace un coin, à environ 2500 pixels du centre, de $2500\,\varphi$ (44 pixels par degré).

\subsection{Pourquoi la rotation suffit (presque) : la parallaxe}
Les deux centres optiques ne sont pas exactement au même endroit : ils sont à 46\,mm l'un de l'autre sur le drone (\texttt{CNIRCVIS\_0} dans \texttt{config.py}), et le drone se déplace entre les deux prises de vue. Une distance $B$ entre les centres (la \emph{base}) crée une \emph{parallaxe} : un point à la distance $Z$ est décalé de
\begin{equation}
	d = f\,\frac{B}{Z} \quad \text{pixels}.
\end{equation}
Seules les \emph{variations} de ce décalage comptent, sa partie constante étant absorbée par la rotation. Avec $B = 46$\,mm, à 90\,m : $d = 2898 \times 0{,}046 / 90 \approx 1{,}5$ pixel, et un relief de 10\,m le modifie de 0,2 pixel : négligeable. Si les deux prises de vue sont séparées de $\Delta t = 0{,}5$\,s et que le drone vole à 3\,m/s, $B = 1{,}5$\,m : $d \approx 48$ pixels, et un relief de 10\,m (une haie d'arbres, une pente) le modifie maintenant de $f B \Delta Z / Z^2 \approx 5$ pixels. C'est la raison pour laquelle le Step~4 n'aligne que les shoot points dont l'écart de temps est faible, et pour laquelle une correction \emph{locale} s'ajoute au modèle global (§\ref{sec:ann_local}). À très basse altitude (décollage, atterrissage), la parallaxe domine et l'alignement n'a pas de sens.

\subsection{La caméra virtuelle}
L'image visible est d'abord corrigée de sa distorsion : elle devient l'image d'une caméra sténopé idéale, la \emph{référence}. L'image infrarouge est ensuite rééchantillonnée directement dans cette caméra de référence, par l'équation \eqref{eq:ann_rotation_homography} et son propre modèle de distorsion, mais sur un canevas plus large que l'image visible (1,4 fois par défaut, 1,6 dans le script REDO) : c'est la \emph{caméra virtuelle} de la figure~\ref{fig:ann_virtual}. Si la rotation était l'identité, l'image visible correspondrait exactement au rectangle central ; chercher la rotation revient à trouver où elle se trouve réellement.

\figureIRD[fig:ann_virtual]{align_virtual_camera.jpg}{Point B. À gauche : l'image infrarouge corrigée de sa distorsion et rééchantillonnée dans la caméra virtuelle (focale de la caméra visible, champ 1,4 fois plus large). En pointillés : position de l'image visible si les deux caméras étaient parallèles ; en rouge : position trouvée par la recherche grossière (lacet $8{,}2^\circ$, tangage $1{,}9^\circ$). À droite : la référence visible. En dehors du champ visible la correction de la distorsion est moins précise (la route s'incurve à gauche).}{1}

\section{Comparer deux bandes spectrales}\label{sec:ann_similarity}
Pour trouver la rotation, il faut un \emph{critère de similarité} : un nombre qui dit si deux images sont bien superposées, et qui est optimal pour la bonne superposition.

\subsection{Pourquoi on ne peut pas comparer les intensités}
Pour deux images d'une même bande, le critère le plus simple est la somme des carrés des différences, $\text{SSD} = \sum_p (I_1(p) - I_2(p))^2$, ou la corrélation. Tous deux supposent qu'un même objet a la même intensité (ou une intensité proportionnelle) dans les deux images. Entre le visible et le proche infrarouge c'est faux : la végétation est sombre dans le visible et claire dans l'infrarouge, les bas-côtés de la route ne sont clairs que dans le visible (figures~\ref{fig:ann_pair} et \ref{fig:ann_features}). Une réponse classique à ces problèmes \og multimodaux \fg{} est l'information mutuelle~[\hyperref[ann_viola]{5}], qui mesure une dépendance statistique au lieu d'une égalité. Une autre approche combine des points caractéristiques et une similarité de gradients~[\hyperref[ann_vayssade]{18}]. IRdrone utilise une idée plus simple et plus locale : les \emph{bords} sont communs même quand les contrastes ne le sont pas.

\subsection{Les cartes de caractéristiques : énergies directionnelles du gradient}\label{sec:ann_msr}
Chaque image est convertie en luminance (niveaux de gris), lissée par un filtre gaussien d'écart-type $\sigma$ (5 pixels pour le visible, 3 pour l'infrarouge, moins net ; 3 et 1 pour la recherche grossière), puis remplacée par quatre cartes (\texttt{multispectral\_representation}, mode \texttt{LAPLACIAN\_ENERGIES}) :
\begin{equation}
	E_x = (\partial_x I)^2, \quad E_y = (\partial_y I)^2, \quad E_{\nearrow} = (\partial_{\nearrow} I)^2, \quad E_{\searrow} = (\partial_{\searrow} I)^2.
\end{equation}
Les dérivées horizontale et verticale sont calculées avec des filtres de Scharr~[\hyperref[ann_scharr]{6}] (une dérivée $3 \times 3$ avec un lissage optimisé, plus isotrope que Sobel), les diagonales avec la croix de Roberts $2 \times 2$~[\hyperref[ann_roberts]{7}]. Ces cartes sont appelées \emph{représentation multispectrale} (MSR) dans le code. L'idée vient d'Irani et Anandan~[\hyperref[ann_irani]{9}], qui alignaient des images de capteurs différents (visible et infrarouge entre autres) au moyen de telles énergies de dérivées directionnelles.

Le carré est essentiel : un bord du sombre vers le clair et un bord du clair vers le sombre donnent la même énergie. Une inversion de contraste entre les bandes devient donc invisible, tandis que la \emph{position} du bord est conservée. Les quatre directions gardent l'orientation des bords, ce qui aide à distinguer, par exemple, une route du bord perpendiculaire d'un champ. Le lissage supprime les textures fines (herbe, feuilles), différentes dans les deux bandes, et garde les grandes structures. La figure~\ref{fig:ann_features} montre ces cartes sur un détail du point B : les luminances sont très différentes, les énergies des bords de la route sont au même endroit.

\figureIRD[fig:ann_features]{align_feature_maps.jpg}{Énergies directionnelles du gradient sur un détail du point B ($600 \times 600$ pixels). Les luminances des deux bandes ne sont pas comparables (bas-côtés clairs dans le visible, sombres dans l'infrarouge), mais les bords de la route produisent les mêmes lignes dans les cartes d'énergie. Certains bords n'existent que dans une bande (les quatre lignes du visible contre deux dans l'infrarouge) : le critère doit les tolérer.}{1}

\subsection{Le critère de similarité : gradient total de la différence}\label{sec:ann_ntg}
Le critère dérive du \emph{Normalized Total Gradient} (NTG) de Chen et al.~[\hyperref[ann_ntg]{8}]. Pour deux images $f$ et $g$ :
\begin{equation}
	\text{NTG}(f, g) = \frac{\sum_p \|\nabla (f - g)(p)\|_1}{\sum_p \|\nabla f(p)\|_1 + \sum_p \|\nabla g(p)\|_1}.
\end{equation}
L'idée : quand deux images sont bien superposées, leurs bords communs s'annulent dans la différence $f-g$, qui devient lisse ; son gradient total (somme des valeurs absolues des dérivées) est alors minimal. Quand elles sont décalées, chaque bord apparaît deux fois dans la différence, avec des signes opposés, et le gradient total augmente. La norme $\|\cdot\|_1$ (valeurs absolues plutôt que carrés) rend le critère robuste aux bords qui n'existent que dans une image : ils ajoutent un coût à peu près constant, qui dépend peu du décalage, au lieu de le dominer. Le dénominateur rend le critère indépendant de la quantité de texture.

IRdrone l'applique aux cartes d'énergie : pour un décalage $u$ de la référence,
\begin{equation}\label{eq:ann_cost}
	C(u) = \sum_{c=1}^{4} \sum_p \Big( \big|S_x * (E^{NIR}_c - E^{VIS}_c(\cdot + u))\big|(p) + \big|S_y * (E^{NIR}_c - E^{VIS}_c(\cdot + u))\big|(p) \Big),
\end{equation}
où $S_x, S_y$ sont les filtres de dérivation de Sobel (\texttt{cost\_surface\_NTG} dans \texttt{cost.py}). Seul le numérateur du NTG est calculé. Le filtre de Sobel étant linéaire, $S * (a - b) = S*a - S*b$ : les gradients des deux images sont calculés une fois, et chaque décalage n'est plus qu'une somme de différences absolues (15 fois plus rapide, novembre 2026). Une somme des carrés des différences sur les mêmes cartes est aussi programmée (\texttt{SSD}) mais n'est pas utilisée.

La figure~\ref{fig:ann_profiles} compare les critères sur le point A, en fonction d'un décalage horizontal autour du bon alignement. La SSD des luminances a une vallée large : elle situe l'alignement, mais avec peu de précision, et elle suppose des intensités comparables, ce qui est le cas sur cette scène d'hiver mais pas en été. Le NTG a un puits étroit et profond : il est très précis, mais ne \og voit \fg{} la bonne position qu'à quelques pixels près, d'où le besoin d'un bon point de départ et d'une stratégie du grossier au fin.

\figureIRD[fig:ann_profiles]{align_cost_profiles.png}{Coût en fonction d'un décalage horizontal (point A, après alignement, calculé à la résolution 1/4), normalisé entre 0 et 1. (a) L'image visible contre elle-même. (b) Visible contre infrarouge sur un bloc avec des arbres et une haie. (c) Sur un bloc de prairie seule, avec peu de structures : le NTG a encore son minimum au bon endroit mais en dehors du puits étroit la courbe n'est plus informative.}{1}

\section{Recherche grossière}\label{sec:ann_coarse}
La recherche grossière trouve l'essentiel de l'écart d'orientation (fonction \texttt{coarse\_alignment}) :
\begin{enumerate}
	\item l'image infrarouge est rééchantillonnée dans la caméra virtuelle avec les angles initiaux : nuls au Step~4, ceux réglés à la main lors d'un réalignement (§\ref{btn:review_redo}) ;
	\item les cartes d'énergie des deux images sont réduites d'un facteur 32 ($125 \times 94$ pixels pour le visible), la carte visible est placée au centre d'un canevas de la taille de la carte infrarouge ;
	\item une \emph{recherche exhaustive} calcule le coût \eqref{eq:ann_cost} pour chaque translation de la carte visible à l'intérieur de la carte infrarouge : avec l'élargissement de champ 1,4 du Step~4, $\pm 25$ pixels à cette échelle, soit environ $\pm 800$ pixels en pleine résolution ou $\pm 15^\circ$ ; avec l'élargissement 1,6 de la ligne de commande, $\pm 37$ pixels, $\pm 1200$ pixels ou $\pm 22^\circ$ (figure~\ref{fig:ann_coarse}) ;
	\item la meilleure position est localisée avec une précision inférieure au pixel (ci-dessous) ;
	\item aux petits angles une translation horizontale $t_x$ est un lacet et une translation verticale $t_y$ un tangage : $\text{yaw} = \arctan(t_x/f)$, $\text{pitch} = \arctan(t_y/f)$ ;
	\item l'image infrarouge est rééchantillonnée une nouvelle fois, en pleine résolution, avec les angles corrigés.
\end{enumerate}
Le roulis n'est pas recherché : il garde sa valeur initiale (§\ref{sec:ann_fourier}).

\paragraph{Choix du minimum.} Le minimum absolu de la surface de coût n'est pas toujours le bon : une zone sans texture ou le bord noir du canevas peuvent produire une vallée large et basse. La bonne superposition produit un puits étroit et profond (figure~\ref{fig:ann_coarse}). Le code cherche donc la position où la surface est la plus fortement courbée : en chaque position, une surface quadratique est ajustée sur le voisinage $3\times3$, et la position du plus grand déterminant de sa matrice hessienne (produit des deux courbures principales) est retenue (\texttt{minimum\_cost\_max\_hessian}).

\paragraph{Raffinement sous-pixel.} Autour de la position retenue, le coût est approché par une forme quadratique~[\hyperref[ann_hasinoff]{10}]
\begin{equation}
	C(u) \approx \tfrac{1}{2}\, u^T A\, u + b^T u + c,
\end{equation}
dont les coefficients s'obtiennent par de petites convolutions du voisinage $3\times3$. Son minimum est en $u^* = -A^{-1} b$ (un pas de Newton), retenu s'il est à moins d'un pixel. À la résolution 1/32, un pixel vaut 32 pixels en pleine résolution : ce raffinement rend la recherche grossière précise à quelques pixels près. Pour le point A, la translation trouvée est de $(21{,}7 ; 163{,}9)$ pixels, soit un lacet de $0{,}43^\circ$ et un tangage de $3{,}24^\circ$.

\figureIRD[fig:ann_coarse]{align_coarse_search.png}{Recherche grossière du point A à la résolution 1/32. À gauche : la surface de coût sur toute la plage de recherche ($\pm 1200$ pixels) : un puits net dans un paysage lisse. Au centre : zoom sur le centre du coût. À droite : déterminant de la hessienne, maximal au fond du puits.}{1}

\section{Raffinement pyramidal}\label{sec:ann_pyramid}
Après la recherche grossière, le désalignement résiduel va de quelques pixels à quelques dizaines de pixels, et il n'est pas uniforme : roulis résiduel, erreurs des calibrations, parallaxe. Le raffinement l'estime par blocs, sur une pyramide d'images (\texttt{pyramidal\_search} dans \texttt{rigid.py}).

\subsection{Du grossier au fin}
Une \emph{pyramide gaussienne}~[\hyperref[ann_burt]{11}] est la suite d'une image lissée et réduite par 2, 4, 8, 16\ldots{} Un déplacement de 100 pixels en pleine résolution fait 6 pixels à 1/16 : on le trouve avec une petite fenêtre de recherche, vite et avec peu de risque de confusion. Une fois trouvé, il est appliqué, et le niveau suivant, plus fin, n'a plus qu'à corriger le reste. C'est la stratégie classique des méthodes d'alignement de Lucas--Kanade~[\hyperref[ann_lucas]{12}, \hyperref[ann_baker]{13}]. IRdrone utilise le schéma suivant (\texttt{iterative\_scheme} dans \texttt{align\_raw}) :

\begin{table}[htbp]\centering
\begin{tabular}{@{}cccccc@{}}
\hline
\textbf{Échelle} & \textbf{Itérations} & \textbf{Blocs} & \textbf{Taille d'un bloc (pleine rés.)} & \textbf{Fenêtre du coût} & \textbf{Plus grand vecteur (pleine rés.)} \\ \hline
1/16 & 2 & $4 \times 4$ & $998 \times 748$ & $\pm 8$ px & $\pm 80$ px \\
1/16 & 2 & $5 \times 5$ & $798 \times 598$ & $\pm 6$ px & $\pm 80$ px \\
1/4 & 3 & $5 \times 5$ & $798 \times 598$ & $\pm 6$ px & $\pm 20$ px \\ \hline
\end{tabular}
\caption{Schéma du raffinement pyramidal. Le coût est calculé sur la fenêtre indiquée, à l'échelle du niveau, mais le minimum est cherché à $\pm 5$ pixels du centre (\texttt{brute\_force\_vector\_field\_search}) : un vecteur de bloc ne peut pas dépasser 80 pixels à 1/16, 20 pixels à 1/4. Les itérations sont numérotées 1, 2, 4, 5, 7, 8, 9 dans les noms des fichiers de traces (le compteur est incrémenté une fois de plus à chaque changement de niveau).}
\label{tab:ann_pyramid}
\end{table}

\subsection{Une itération}
\begin{enumerate}
	\item Les cartes d'énergie de l'image infrarouge sont rééchantillonnées avec le modèle courant (l'identité au départ).
	\item L'image est découpée en une grille de blocs. Pour chaque bloc, le coût \eqref{eq:ann_cost} est calculé pour chaque décalage de la fenêtre de recherche : on obtient une \emph{surface de coût} par bloc (figures~\ref{fig:ann_blocks_first} et \ref{fig:ann_blocks_last}).
	\item Le minimum de chaque surface (à $\pm 5$ pixels du centre), affiné au sous-pixel comme ci-dessus, donne un vecteur de déplacement par bloc : un \emph{champ de vecteurs} grossier, ou flot optique.
	\item Une homographie est ajustée sur ces vecteurs : les 16 ou 25 couples (centre du bloc, centre + vecteur) forment un système surdéterminé pour ses 8 inconnues. Il est résolu par la transformation linéaire directe (DLT) dans une boucle RANSAC~[\hyperref[ann_ransac]{14}], qui ignore les blocs dont le vecteur s'écarte de plus de 3 pixels de la majorité (\texttt{geometric\_rigid\_transform\_estimation}, \texttt{findHomography} d'OpenCV ; Hartley et Zisserman~[\hyperref[ann_hz]{1}], chapitre 4).
	\item Le modèle est mis à jour par composition, $H_k = \Delta H_k\, H_{k-1}$, puis l'itération suivante commence.
\end{enumerate}
Une homographie plutôt que trois angles : elle représente la rotation d'une caméra \emph{et} les erreurs résiduelles des calibrations (pour le point A, l'homographie finale contient un facteur d'échelle de 1,9\,\% horizontalement et 1,1\,\% verticalement, par exemple une petite erreur sur les focales).

\figureIRD[fig:ann_blocks_first]{align_blocks_first.jpg}{Première itération du raffinement du point A (résolution 1/16, blocs $4\times4$, recherche $\pm 8$ pixels). À gauche : superposition après la recherche grossière (visible en vert, infrarouge en magenta ; gris là où ils concordent), avec les vecteurs mesurés (bleu, agrandis 20 fois) et les vecteurs de l'homographie ajustée (rouge). À droite : la surface de coût de chaque bloc ; le point noir est le décalage nul, la croix le minimum. Les vecteurs atteignent 50 pixels et varient régulièrement : un facteur d'échelle et un tangage résiduels, que l'homographie représente bien.}{1}

\figureIRD[fig:ann_blocks_last]{align_blocks_last.jpg}{Dernière itération (résolution 1/4, blocs $5\times5$, recherche $\pm 6$ pixels), sur l'image alignée globalement. Tous les minima sont au centre : les vecteurs résiduels (agrandis 200 fois) sont inférieurs à 2 pixels, l'homographie ne change plus. Ces vecteurs résiduels forment la correction locale (§\ref{sec:ann_local}).}{1}

\figureIRD[fig:ann_convergence]{align_convergence.png}{Longueur des vecteurs des blocs à chaque itération, pour les deux points. Le point A (bleu) converge : de 23 pixels en moyenne à 0,6 pixel. Les marches aux changements de niveau sont dues au changement de taille des blocs et de résolution (un pixel à 1/16 vaut 16 pixels). Le point B (rouge) ne converge pas : les vecteurs restent à leur limite, 80 pixels à 1/16 et 20 pixels à 1/4 (§\ref{sec:ann_limits}).}{0.7}

\subsection{Alignement local}\label{sec:ann_local}
Une homographie est un modèle global : elle ne peut pas représenter la parallaxe du relief ni une imperfection locale du modèle de distorsion. Le champ de vecteurs de la dernière itération ($5 \times 5$ vecteurs) est donc conservé comme \emph{correction locale} : il est interpolé bilinéairement entre les centres des blocs (constant au-delà) en un déplacement pour chaque pixel, composé avec le modèle global (\texttt{warp\_from\_sparse\_vector\_field}). Deux résultats sont produits :
\begin{itemize}
	\item alignement \emph{global} : rotation et homographie finale ;
	\item alignement \emph{local} : rotation, homographie de l'itération précédente et champ de déplacement local. C'est en général le plus précis, et celui des images multispectrales \texttt{IRD\_XXXX\_n.tif}.
\end{itemize}

\subsection{Un seul rééchantillonnage}
Chaque rééchantillonnage d'une image interpole entre ses pixels, et la rend donc un peu floue. Au lieu de rééchantillonner l'image infrarouge à chaque étape, la transformation finale est composée en une seule carte : pour chaque pixel de l'image visible, la position à lire dans l'image infrarouge \emph{brute} s'obtient en enchaînant le déplacement local, l'homographie, la rotation de l'équation \eqref{eq:ann_rotation_homography} et la distorsion de la caméra infrarouge ; l'image n'est ensuite interpolée qu'une fois (interpolation bicubique, \texttt{cv2.remap}). Dans le code ce principe s'appelle \og You Only Warp Once \fg{} (\texttt{FULLRES\_REGISTERED\_REFINED\_WARPED\_ONCE\_ONLY.jpg}). Le modèle de mouvement (angles, homographies, champ de vecteurs) est enregistré dans \texttt{VIS\_XXXX\_motion\_model.npy} et sa version lisible \texttt{.json} ; s'il existe, relancer le Step~4 le relit au lieu de chercher à nouveau.

La figure~\ref{fig:ann_stages} montre la progression de l'alignement sur un détail du point A.

\figureIRD[fig:ann_stages]{align_stages.jpg}{Les quatre étapes sur un détail du point A (visible en vert, infrarouge en magenta, gris là où ils concordent). (a) Calibrations seules : les images sont décalées d'environ 160 pixels. (b) Après la recherche grossière : il reste une dizaine de pixels (franges colorées le long des ombres). (c) Avec l'homographie. (d) Avec la correction locale : les bords coïncident.}{1}

\section{Pistes explorées : corrélation de phase, Fourier--Mellin et appariement polaire}\label{sec:ann_fourier}
\subsection{Corrélation de phase}
Les premières versions (août--septembre 2021) estimaient la translation par \emph{corrélation de phase}~[\hyperref[ann_kuglin]{15}]. Si $g(x) = f(x - t)$, leurs transformées de Fourier vérifient $\hat g(\omega) = \hat f(\omega)\, e^{-i\omega \cdot t}$ (théorème du retard). L'interspectre normalisé
\begin{equation}
	\frac{\hat f\, \overline{\hat g}}{|\hat f\, \overline{\hat g}|} = e^{i\omega \cdot t}
\end{equation}
ne garde que la phase, et sa transformée de Fourier inverse est un pic de Dirac à la position $t$. La méthode est rapide (deux FFT), insensible à un changement global de luminosité, et trouve de grandes translations. Mais elle suppose que les deux images sont le \emph{même} signal décalé : une rotation, un changement d'échelle, une distorsion ou une inversion de contraste étalent le pic. La figure~\ref{fig:ann_phase} le montre : sur le point A, après la recherche grossière, le pic est net ; sur le point B il est noyé dans le bruit, sur les luminances comme sur les cartes d'énergie.

\figureIRD[fig:ann_phase]{align_phase_correlation.png}{Surfaces de corrélation de phase entre l'image visible et l'image infrarouge après la recherche grossière (résolution 1/8, zoom sur $\pm 320$ pixels). En haut, point A : un pic net, environ 2,5 fois plus haut que les valeurs suivantes. En bas, point B (rotation résiduelle, contrastes inversés, grandes zones texturées) : le pic émerge à peine.}{0.65}

\subsection{Fourier--Mellin : rotation et échelle}
La méthode de Fourier--Mellin~[\hyperref[ann_reddy]{16}] étend la corrélation de phase aux rotations et aux changements d'échelle. Le module de la transformée de Fourier ne dépend pas de la translation ; une rotation de l'image le fait tourner, un changement d'échelle le dilate. Rééchantillonné en coordonnées \emph{log-polaires} $(\log r, \theta)$, le module transforme une rotation en translation selon $\theta$ et un changement d'échelle en translation selon $\log r$, qu'une seconde corrélation de phase retrouve. Rabatel et Labbé~[\hyperref[ann_rabatel]{17}] l'ont appliquée à des images de drone visibles et proche infrarouge. Dans IRdrone, l'échelle est connue par les calibrations, et seul le roulis était cherché de cette façon : l'alignement semi-automatique de 2021 rééchantillonnait les deux images en coordonnées polaires et estimait le roulis par corrélation de phase selon l'angle (\texttt{mellin\_transform} et \texttt{estimate\_rotation} dans \texttt{semi\_auto\_registration.py}, toujours présentes pour une démonstration mais non utilisées par le Step~4). Combinée à la translation par corrélation de phase, ou à des points SIFT, elle n'a pas donné de résultats fiables sur les images multispectrales, pour les raisons ci-dessus. À partir de septembre 2021 elle a été remplacée par la minimisation du coût \eqref{eq:ann_cost} sur les cartes d'énergie.

\subsection{Appariement polaire des cartes d'énergie (essai de 2022)}
Le roulis est le point faible de l'algorithme actuel : la recherche grossière ne trouve que le lacet et le tangage, et le raffinement pyramidal ne corrige qu'un petit roulis. Lors d'un vol à Thaumiers (novembre 2022) la caméra infrarouge était montée avec un roulis constant d'environ $-2{,}4^\circ$, soit des déplacements de 100 pixels dans les coins, trop pour le raffinement (ticket GitHub \#33). Un essai a combiné les deux approches : les cartes d'énergie sont rééchantillonnées en coordonnées polaires autour du centre de l'image, où un roulis devient une translation verticale, et le coût \eqref{eq:ann_cost} est calculé pour chaque décalage de $1^\circ$ ; une parabole ajustée sur les trois valeurs les plus basses donne le roulis à mieux qu'un degré (\texttt{coarse\_roll\_estimation}, branche \texttt{local\_alignment\_improvements}, commit \texttt{86353cb}, non intégré). La figure~\ref{fig:ann_polar} le reproduit : sur le point A, tourné d'un roulis connu de $2{,}4^\circ$, l'estimation est de $2{,}26^\circ$. Sur le point B, la route est à environ $34^\circ$ dans l'image visible et à environ $28^\circ$ dans l'image infrarouge après la recherche grossière, et l'appariement polaire trouve un roulis de $-6{,}3^\circ$, mais avec un second minimum presque aussi profond à $+5^\circ$ (figure~\ref{fig:ann_polar_b}) : le profil est beaucoup moins tranché que sur le point A.

\figureIRD[fig:ann_polar]{align_polar_roll.jpg}{Appariement polaire sur le point A, dont l'image infrarouge a été tournée volontairement de $2{,}4^\circ$. À gauche : cartes d'énergie des deux images en coordonnées polaires (angle verticalement, distance au centre horizontalement) ; un roulis est un décalage vertical. À droite : coût en fonction du décalage angulaire, et parabole ajustée autour du minimum ($2{,}26^\circ$).}{1}

\figureIRD[fig:ann_polar_b]{align_polar_roll_b.jpg}{La même chose sur le point B, sans roulis artificiel, après la recherche grossière : minimum à $-6{,}3^\circ$, mais un second minimum près de $+5^\circ$ est presque aussi profond.}{1}

\section{Historique du développement}\label{sec:ann_history}
L'historique git du projet donne les jalons suivants.
\begin{itemize}
	\item \textbf{Août 2021} : alignement semi-automatique. L'utilisateur tourne l'image infrarouge avec trois curseurs (lacet, tangage, roulis, équation~\eqref{eq:ann_rotation_homography}) ; le résultat est affiné par corrélation de phase, par corrélation de phase polaire pour le roulis, ou par des points SIFT ou une corrélation de phase par blocs avec un ajustement affine. Comparaison sur la valeur absolue du gradient plutôt que sur les intensités.
	\item \textbf{Septembre 2021} : fonctions de coût (SSD, NTG) sur plusieurs représentations, recherche de Newton et ajustement quadratique du minimum (d'après le matériel supplémentaire de HDR+~[\hyperref[ann_hasinoff]{10}]), puis recherche itérative pyramidale ajustant une homographie sur les vecteurs des blocs, accélérée avec Numba.
	\item \textbf{Octobre 2021} : déformation locale suivant le champ de vecteurs, caméra virtuelle au champ élargi pour le fish-eye, chaîne complète recherche grossière + pyramide, choix du minimum par la hessienne, rééchantillonnage unique (\og You Only Warp Once \fg), sorties locale et globale.
	\item \textbf{Mars 2022} : ticket GitHub \#4, \og fix local alignment issues \fg{} (toujours ouvert) : utiliser la courbure du coût pour décider si un vecteur local est fiable, régulariser le champ avec ses voisins.
	\item \textbf{Novembre--décembre 2022} : ticket \#33, roulis de la caméra infrarouge ; appariement polaire des cartes d'énergie, raffinement parabolique du minimum (branche expérimentale).
	\item \textbf{Novembre 2026} : accélération sans changement des résultats (coût par linéarité du filtre de Sobel, pyramides en float32 avec OpenCV, réduction par 32 calculée sur les seules lignes et colonnes utiles) : environ 15\,s par shoot point ; réalignement à la main depuis la revue dans une fenêtre dédiée (§\ref{btn:review_redo}).
\end{itemize}

\section{Traces de diagnostic}\label{sec:ann_traces}
Les figures de cette annexe proviennent des traces de diagnostic de l'alignement. Elles ne sont pas écrites par défaut au Step~4 (environ 100\,Mo de figures par point). On les obtient :
\begin{itemize}
	\item en lançant le script \texttt{VIS\_XXXX\_REDO.bat} (ou \texttt{.sh}) écrit à côté de chaque image alignée : un \og full recovery run \fg{} avec l'option \texttt{--debug}, qui ouvre la fenêtre à curseurs puis écrit tout ;
	\item depuis la revue, quand le niveau de détail du journal des réglages du Step~4 (\bouton{⚙}) est \og Debug \fg{} : un réalignement à la main fait alors le calcul complet ;
	\item en ligne de commande : \texttt{python automatic\_registration.py --images VIS\_XXXX.dng NIR\_YYYY.dng --outdir OUT --debug}.
\end{itemize}
Les traces sont écrites dans le dossier \texttt{VIS\_XXXX\_align\_traces} :
\begin{itemize}
	\item \texttt{Full Search\_\ldots} : la recherche grossière (surface de coût, tableau \texttt{.npy} des coûts) ;
	\item \texttt{\_PADDED\_LOWRES\_*}, \texttt{\_LOWRES\_*} : cartes d'énergie à 1/32 avant et après la recherche grossière ;
	\item \texttt{ds16\_\ldots\_cost\_surface\_Y\_X\_\_itNN.png} : les entrées et les quatre surfaces de coût du bloc $(Y, X)$ à l'itération NN ; \texttt{\ldots overview\_cost\_surfaces\ldots} : tous les blocs d'une itération ; \texttt{\ldots\_cost\_search\_\_itNN.npy} : les coûts, les centres des blocs et l'échelle, lisibles par \texttt{numpy.load(\ldots, allow\_pickle=True)} ;
	\item \texttt{\_msr\_itNN\_dsMM\_\ldots}, \texttt{\_image\_itNN\_\ldots} : cartes d'énergie et images à chaque itération ;
	\item \texttt{FLOW\_itNN\_\ldots} : champ de vecteurs et image déformée par celui-ci ;
	\item \texttt{FULLRES\_REF.jpg}, \texttt{FULLRES\_REGISTERED\_COARSE.jpg}, \texttt{FULLRES\_REGISTERED\_REFINED\_WARPED\_ONCE\_ONLY.jpg} : la référence et l'image infrarouge après chaque étape, à comparer en basculant de l'une à l'autre ;
	\item \texttt{motion\_model.npy} : le modèle de mouvement.
\end{itemize}

\section{Limites et problèmes ouverts}\label{sec:ann_limits}
L'alignement multispectral est un problème mal posé, et l'algorithme échoue dans des situations identifiables :
\begin{itemize}
	\item les deux caméras ne voient pas la même scène : rafale de vent, grand écart de temps, basse altitude (parallaxe) ;
	\item la zone a peu de structures : eau, champ uniforme, canopée forestière ; le coût est alors plat ou trompeur (figure~\ref{fig:ann_profiles}c) ;
	\item le désalignement dépasse ce que les étapes savent corriger : un roulis de plusieurs degrés, que la recherche grossière ne modélise pas ;
	\item la recherche grossière tombe dans un minimum trompeur.
\end{itemize}
Le point B cumule plusieurs de ces difficultés : grand lacet ($8{,}2^\circ$), roulis probable de plusieurs degrés, grandes zones de végétation dont les textures diffèrent entre les bandes. Les vecteurs des blocs restent à leur limite de $\pm 80$ puis $\pm 20$ pixels (figures~\ref{fig:ann_blocks_hard} et \ref{fig:ann_convergence}) ; l'ajustement RANSAC retient alors une homographie aux forts termes de perspective, qui déplace les coins de l'image de plusieurs centaines de pixels. Le résultat est acceptable près du centre et faux vers les coins (textures étirées en bas à gauche). Dans nos essais, partir du roulis estimé par appariement polaire ($\pm 6{,}3^\circ$) ne l'a pas fait converger non plus : ce point demande un réalignement à la main. De tels échecs sont détectés à la revue (chapitre~\ref{ch:review}), qui mesure le décalage résiduel des bords.

\figureIRD[fig:ann_blocks_hard]{align_blocks_hard.jpg}{Première itération du raffinement du point B : les surfaces de coût des blocs de végétation n'ont pas de minimum net, les minima sont à des positions aléatoires, souvent au bord de la fenêtre de recherche (vecteurs agrandis 10 fois).}{1}

Le réalignement à la main (§\ref{btn:review_redo}) remplace les angles initiaux nuls par des angles choisis à l'œil, roulis compris. Ces angles ne sont qu'un point de départ : la recherche grossière et le raffinement pyramidal repartent de là et trouvent en général un alignement précis.

Améliorations possibles, en partie listées dans les tickets GitHub \#4 et \#33 :
\begin{itemize}
	\item utiliser la courbure de chaque surface de coût (matrice $A$ de l'ajustement quadratique) comme indice de confiance : un vecteur mesuré sur une surface plate ne devrait compter ni dans l'homographie ni dans la correction locale ;
	\item régulariser le champ de vecteurs local (médiane des voisins, contrainte de régularité) au lieu de l'utiliser brut ;
	\item ajouter le roulis à la recherche grossière (appariement polaire, ou recherche sur quelques valeurs de roulis) ;
	\item normaliser le coût (dénominateur du NTG), pour ne pas favoriser un bloc moins texturé ;
	\item calibrer les couleurs des deux caméras (calibration radiométrique) pour calculer des indices physiquement significatifs.
\end{itemize}

\section{Références}\label{sec:ann_references}
{\renewcommand{\labelenumi}{[\arabic{enumi}]}
\begin{enumerate}
	\item\label{ann_hz} R. Hartley, A. Zisserman, \emph{Multiple View Geometry in Computer Vision}, 2\textsuperscript{e} édition, Cambridge University Press, 2004. Chapitres 2 (géométrie projective du plan, homographies), 4 (estimation d'une homographie, DLT, RANSAC), 6 (modèles de caméra), 7.4 (distorsion radiale), 8.4 (caméra en rotation, rotations synthétiques, mosaïques).
	\item\label{ann_szeliski} R. Szeliski, \emph{Computer Vision: Algorithms and Applications}, 2\textsuperscript{e} édition, Springer, 2022 (chapitre 2 : formation de l'image ; chapitres 8 et 9 : alignement et estimation du mouvement). Gratuit en ligne sur \texttt{szeliski.org/Book}.
	\item\label{ann_brown} D. C. Brown, \emph{Close-range camera calibration}, Photogrammetric Engineering 37(8), 1971.
	\item\label{ann_zhang} Z. Zhang, \emph{A flexible new technique for camera calibration}, IEEE Transactions on Pattern Analysis and Machine Intelligence 22(11), 2000.
	\item\label{ann_viola} P. Viola, W. M. Wells III, \emph{Alignment by maximization of mutual information}, International Journal of Computer Vision 24(2), 1997.
	\item\label{ann_scharr} H. Scharr, \emph{Optimale Operatoren in der digitalen Bildverarbeitung}, thèse de doctorat, université de Heidelberg, 2000.
	\item\label{ann_roberts} L. G. Roberts, \emph{Machine perception of three-dimensional solids}, thèse de doctorat, MIT, 1963.
	\item\label{ann_ntg} S.-J. Chen, H.-L. Shen, C. Li, J. H. Xin, \emph{Normalized Total Gradient: a new measure for multispectral image registration}, IEEE Transactions on Image Processing 27(3), 2018 (\texttt{arXiv:1702.04562}).
	\item\label{ann_irani} M. Irani, P. Anandan, \emph{Robust multi-sensor image alignment}, ICCV 1998.
	\item\label{ann_hasinoff} S. W. Hasinoff et al., \emph{Burst photography for high dynamic range and low-light imaging on mobile cameras}, SIGGRAPH Asia 2016 (matériel supplémentaire : alignement pyramidal par blocs, raffinement sous-pixel quadratique).
	\item\label{ann_burt} P. J. Burt, E. H. Adelson, \emph{The Laplacian pyramid as a compact image code}, IEEE Transactions on Communications 31(4), 1983.
	\item\label{ann_lucas} B. D. Lucas, T. Kanade, \emph{An iterative image registration technique with an application to stereo vision}, IJCAI 1981.
	\item\label{ann_baker} S. Baker, I. Matthews, \emph{Lucas-Kanade 20 years on: a unifying framework}, International Journal of Computer Vision 56(3), 2004.
	\item\label{ann_ransac} M. A. Fischler, R. C. Bolles, \emph{Random sample consensus: a paradigm for model fitting with applications to image analysis and automated cartography}, Communications of the ACM 24(6), 1981.
	\item\label{ann_kuglin} C. D. Kuglin, D. C. Hines, \emph{The phase correlation image alignment method}, Proceedings of the IEEE International Conference on Cybernetics and Society, 1975.
	\item\label{ann_reddy} B. S. Reddy, B. N. Chatterji, \emph{An FFT-based technique for translation, rotation, and scale-invariant image registration}, IEEE Transactions on Image Processing 5(8), 1996.
	\item\label{ann_rabatel} G. Rabatel, S. Labbé, \emph{Registration of visible and near infrared unmanned aerial vehicle images based on Fourier-Mellin transform}, Precision Agriculture 17, 2016.
	\item\label{ann_vayssade} J.-A. Vayssade, G. Jones, J.-N. Paoli, C. Gée, \emph{Two-step multi-spectral registration via key-point detector and gradient similarity}, 2020.
\end{enumerate}}

\section{Détermination des \og Pitch \fg{ } et \og Yaw \fg{ } théoriques.}\label{sec:ann_theorique}

\subsection{Description du positionnement des caméras.}

La première étape du recalage de l'image infrarouge sur l'image visible consiste à appliquer une translation à l'image NIR de façon à aligner son centre avec celui de l'image VIS. En effet à la prise de vue les axes de visée des deux caméras ne sont pas parallèles. De plus le déclenchement des caméras n'est pas simultané. La caméra du drone est montée sur un \og gimbal \fg{} qui maintient la direction de visée quels que soient les mouvements du drone
\footnote{Dans la limite des spécifications indiquées par le constructeur DJI et reproduites sur la figure \ref{AnglesDroneGimbal}}.

\figureIRD[AnglesDroneGimbal]{AnglesDroneGimbal.pdf}{Description schématique du \og gimbal \fg{} et des amplitudes angulaires autorisées.}{0.6}

A l'inverse la caméra NIR est fixée rigidement au drone (figure \ref{Definition_Reperes_Cameras}). Ainsi son angle de visée va suivre les mouvements de ce dernier. Le déclenchement se fait en continue avec une période  $\Delta T^{Nir} = 3s$. C'est la fréquence la plus élevée possible pour la SJCam M20.\\


\figureIRD[Definition_Reperes_Cameras]{Definition_Reperes_Cameras.pdf}{Définition des angles de lacet, tangage et roulis (yaw, pitch, roll) du drone, du gimbal et des deux caméras VIS et NIR. Les caméras pointent vers le sol alors que le drone est horizontal ce qui engendre la permutation des axes yaw et roll entre le drone et le gimbal d'une part et les caméras d'autre part.}{0.8}

Nous allons prendre l'exemple du pitch en admettant un découplage entre les angles. Ceci correspond à faire l'hypothèse implicite que les angles yaw et roll sont nuls (ou très faibles).  Nous définissons dans le plan vertical passant par l'axe  du drone les différents angles permettant de repérer la position du drone et des deux caméras à un instant donné (figure \ref{Definition_des_angles_de_Pitch}): 
\figureIRD[Definition_des_angles_de_Pitch]{Definition_des_angles_de_Pitch.pdf}{Définition des angles dans le plan vertical.}{0.8}

\begin{itemize}
	\item $\alpha\qquad$  Flight Pitch. Angle de tangage du drone.
	\item $\delta\qquad$  NIR misalignement Pitch. Cet angle correspond à un défaut de positionnement du support de la caméra NIR sur le drone.
	\item $\varphi$ Angular wedge. Correspond à une éventuelle cale de correction du défaut $\delta$.
	\item $\zeta\qquad$ Residual deviation Pitch. $\quad \zeta = \delta + \phi$.
	\item $\theta_{Nir}\:\:\;$ Camera NIR Pitch.  $\quad \theta_{Nir} = \alpha + \zeta$.
	\item $\beta\qquad$ Gimbal Pitch. Angle de tangage de gimbal (et de la caméra VIS).
	\item $\delta_{Vis}\;\;\;$ Offset VIS. Par définition $\delta_{Vis} =90^\circ$.
	\item $\theta_{Vis}\;\;\;$ Camera VIS Pitch.  $\quad \theta_{Vis} = \beta + \delta_{Vis}$. 
	\item $\psi\qquad$  Theoritical NIR Pitch. $\quad \psi =\delta_{Vis}+\beta-\alpha$.
	\item $\psi_{r}\quad\;\;$ NIR Pitch. Angle entre les axes de visées des deux caméras au même instant $t$. $\qquad \qquad \zeta = \psi_{r} -\psi$. \\
\end{itemize}

Les données Exif des images de la caméra VIS indiquent les coordonnées GPS mais aussi les angles de yaw, pitch et roll du drone et du gimbal. Autrement dit $\alpha$ et $\beta$ sont connus à chaque acquisition d'une image de la caméra VIS. L'inclinaison de l'axe de visée engendre une déformation de l'image en trapèze comme le montre la figure \ref{Effet_du_pitch_sur_camera_pinhole}.

\figureIRD[Effet_du_pitch_sur_camera_pinhole]{Effet_du_pitch_sur_camera_pinhole.pdf}{Effet du tangage sur la déformation de l'image d'un plan horizontal formée sur le capteur d'une caméra sténopé (\og pinhole\fg{}). On notera que l'image sur le capteur est inversée.}{0.6}


Pour simplifier nous admettrons dans la suite que le  montage de la caméra NIR est sans défaut, autrement dit $\delta=0$. 
Nous supposons que les horloges des deux caméras VIS et NIR sont synchronisées. 
On dispose donc des instants exacts où les images de ces deux caméras ont été prises. 
Les paires d'images(VIS, NIR) sont formées de telle manière que l'intervalle de temps séparant les prises de vues de ces deux images soit minimum.

\subsection{Calcul du pitch théorique dans le cas synchronisé.}

\figureIRD[Pitch_cas_synchro]{Pitch_cas_synchro.pdf}{Position du drone et des deux caméras dans le cas d'images prises au même instant. Le pitch a été exagéré pour les besoins de la figure.}{0.5}

Si les images VIS et NIR sont prises au même instant (cas d'une prise de vue synchronisée
\footnote{On ne confondra pas la synchronisation des horloges et la synchronisation des prises de vue.}) il est possible de déterminer à chaque prise de vue l'angle $\psi $ entre les axes de visée des deux caméras. 

Notons que la distance entre les centres des caméras $C^{Nir}$ et $C^{Vis}$ (\og base line \fg{}) est négligeable devant l'altitude $H$ du drone par rapport au sol. 
Le rapport est inférieur à $4\;10^{-4}$. 
On peut donc  prendre $\phi \cong \theta_{Vis}$.

Enfin puisque nous supposons la caméra NIR fixée sans défaut $\theta_{Nir} =\alpha_k$. \\

En observant la figure \ref{Pitch_cas_synchro} et en utilisant des formules de géométrie élémentaires nous en déduisons la formule du pitch théorique pour l'image infrarouge dans le cas d'un déclenchement synchronisé à l'instant $t_k$
\begin{equation}
	\label{Theoritical_Pitch_syncho_2}
	\psi_k  =\phi-\theta_{Nir} =\beta_k + \delta_{Vis} - \alpha_k  \quad .
\end{equation}

La distance entre les centres des images, calculée au niveau du sol est
\begin{equation}
	L = H \tan \psi_k \quad ,
\end{equation}
et  la translation à appliquer à l'image NIR pour la centrer sur l'image VIS s'écrit, en notant $\mathrm{f}^{Nir}$ la distance focale de la caméra infrarouge 
\begin{equation}
	\mathbf{t}_y = \mathrm{f}^{Nir} \tan \psi_k \quad .
\end{equation} 

Bien entendu cette situation est un cas idéal
\footnote{C'est la configuration retenue pour le matériel professionnel.}.

\subsection{Estimation du pitch théorique dans le cas désynchronisé.}
En pratique le déclenchement des caméras VIS et NIR n'est pas simultané. 
Avec une acquisition en \og hyperlapse \fg{}
\footnote{Le terme \og time-lapse\fg{}  ou chronophotographie  correspond à la prise de clichés avec un intervalle de temps régulier et pour une position fixe de la caméra. En \og hyperlapse \fg{} au contraire la caméra est mobile avec un mouvement régulier.}
l'intervalle pour la caméra VIS est de $2s$
\footnote{En pratique l'intervalle est plutôt égal à $2,019s$. Autrement dit environ toutes les 50 images VIS la date d'enregistrement présente un saut de $1s$. Ce phénomène est corrigé par notre algorithme afin d'obtenir une \og ligne temporelle\fg{} précise au dixième de seconde.} alors qu'il est de $3s$ pour la caméra NIR. 
La figure \ref{Periodicite-du-pitch-NIR} montre un exemple d'une série d'images hyperlapse  replacées sur un même axe  temporel une fois les horloges synchronisées.\\

\figureIRD[Periodicite-du-pitch-NIR]{Periodicite-du-pitch-NIR.pdf}{Images hyperlapse des caméras VIS et NIR. Elles sont désynchronisées d'un délai variable positif ou négatif avec une périodicité de trois images visibles, soit $6s$. Pour vérifier la désynchronisation des prises de vues nous avons utilisé la présence d'un véhicule mobile. Vol FLY-20211109-1m/s.}{1}


Cette fois la base line $\mathbf{Bl}(t)$ à l'instant de la prise de vue de l'image NIR ne peut plus être négligée. Les formules de calcul du pitch deviennent
\footnote{ $(\mathbf{e}_1, \mathbf{e}_2, \mathbf{e}_3)$ repère orthonormé direct. $\mathbf{e}_2$ est dans l'axe du drone et $\mathbf{e}_3$ est sur l'axe vertical vers le haut. }

\begin{equation}
	\label{Theoritical_Pitch_non_syncho_0}
	\begin{aligned}
		\psi_k&=\phi-\alpha(t) \quad,\\
		\phi&= \arctan \left(\frac{\mathbf{Bl}(t) \cdot \mathbf{e}_2}{H} +\tan \theta_{Vis} \right)\quad.\\	
	\end{aligned} 
\end{equation}\\

Malheureusement le pitch du drone $\alpha(t)$ n'est pas connu à l'instant $t$ où l'image NIR a été prise. Il en est de même de la position exacte du drone à l'instant $t$, et donc de la \og base line \fg{} $\mathbf{Bl}(t)$. \\

Nous supposerons qu'entre deux images VIS espacées de $\Delta T^{Vis} =2s$ le drone se déplace en ligne droite et à une altitude constante.
On propose de déterminer $\alpha(t)$ et $\mathbf{Bl}(t)$ à l'aide d'une simple interpolation linéaire.\\

Le délai entre la prise de vue des images VIS et NIR est connu et noté $\delta t= t_k -\,t$ . 
Il peut être positif ou négatif selon que l'image NIR est prise avant ou après l'image VIS. 
Les positions GPS du drone par rapport au sol sont connues.  
Dans le repère géographique 
\footnote{$(\mathbf{e}_{\mathrm{W}}, \mathbf{e}_{\mathrm{N}}, \mathbf{e}_3)$ repère orthonormé direct.}
nous utiliserons les coordonnées UTM $(X, Y, Z)$. Bien entendu il faut déterminer les déplacements du drone dans son propre repère. L'orientation du drone par rapport au nord géographique est donnée par l'angle de lacet $\textit{Yaw}_k$. Soit\\ 
\begin{equation}
	\label{Theoritical_Pitch_non_syncho_11}
	\begin{aligned}
		{x}^{EW}_k&=X(t_{k+1})-X(t_{k}) \quad,\\
		{y}^{SN}_k&=Y(t_{k+1})-Y(t_{k}) \quad,\\
		{x_1}_k&=-{x}^{EW}_k\; \cos \textit{Yaw}_k + {y}^{SN}_k\; \sin \textit{Yaw}_k \quad,\\
		{x_2}_k&=\;\;{x}^{EW}_k\:\: \sin \textit{Yaw}_k + {y}^{SN}_k\; \cos \textit{Yaw}_k  \quad.\\
	\end{aligned} \\
\end{equation}


\begin{equation}
	\label{Theoritical_Pitch_non_syncho_1}
	\begin{aligned}
		si \;\delta t <0 \qquad& \qquad\alpha(t)&=&\;\alpha_k \left(\frac{\delta t}{ \Delta T^{Vis} }+1\right) - \alpha_{k+1}\,\frac{\delta t}{ \Delta T^{Vis}} \quad,\\
		&\mathbf{Bl}(t) \cdot \mathbf{e}_2&=&\; {x_2}_k\,\frac{\delta t}{ \Delta T^{Vis}} \quad,&\\
		&&&\\
		si \;\delta t >0 \qquad &\qquad \alpha(t)&=&\;\alpha_{k-1} \,\frac{\delta t}{ \Delta T^{Vis}}- \alpha_k\, \left(\frac{\delta t}{ \Delta T^{Vis} }-1\right)  \quad,\\
		&\mathbf{Bl}(t) \cdot \mathbf{e}_2&=&\; {x_2}_{k-1}\,\frac{\delta t}{ \Delta T^{Vis}} \quad.&\\
	\end{aligned} \\
\end{equation}\\

Un exemple des positions du drone et des caméras est représentée sur la figure \ref{Pitch_cas_non_synchro}.\\ 

\figureIRD[Pitch_cas_non_synchro]{Pitch_cas_non_synchro.pdf}{Cas d'une désynchronisation \og arrière \fg. L'image NIR a été prise après l'image visible. La \og base-line \fg{} a été exagérée. En pratique le point vert est situé avant le point rouge pour une désynchronisation arrière.}{0.9}


L'écart de désynchronisation $\delta t$ est périodique et nous noterons sa période  $T$. 
Dans notre exemple $T= 2 \Delta T^{Nir} = 3 \Delta T^{Vis}$. 
Le pitch de la caméra VIS, contrôlé par le gimbal, est quasi constant. 
En l'absence de vent et avec une vitesse constante le pitch du drone est lui aussi constant. 
Dans ces conditions idéales le signal $\psi(t)$ sera lui aussi de période T. 
Dans le cas d'un vol avec du vent le pitch du drone va varier pour compenser les turbulences. 
Ce phénomène peut être traduit par un pitch aléatoire 
\footnote{On ne rentrera pas dans des considérations très complexes sur la modélisation de la turbulence. Par exemple l'ergodicité,  la turbulence homogène  etc. }. 
Le pitch aura donc une composante périodique et une composante aléatoire. Le spectre de Fourier d'un tel signal présentera  un pic à la fréquence $0,166\, Hz$ ce qui correspond à la composante périodique de période $T=6\,s$.\\ 

\figureIRD[Image ViR]{_VIR__local_HYPERLAPSE_0057.jpg}{Image multi-spectrale VIR produite par IRDrone. FLY-20211109-1m/s.}{0.65}

\subsection{Analyse sur des données réelles.}

On utilise ici les données de la mission FLY-20220125 effectué le 25 janvier 2022 en fin de matinée (10 h 59, heure solaire). Le ciel était clair, le vent faible (< 3 m/s) du nord-ouest. Après la phase de synchronisation le drone s'est élevé à une altitude de 100m puis a débuté son vol d'acquisition des images \og utiles \fg{}. Pendant cette phase l'orientation du drone est restée constante ($-82^\circ$) de même que sa vitesse $1,5\,m/s$. Pendant les $232 s$ de la phase d'acquisition  les caméras VIS et NIR ont pris respectivement $116$ et $85$ images. La trajectoire du  drone se décompose en trois segment droits:
\begin{itemize}
	\item Vol vers l'avant en ligne droite. Direction $\sim 281^\circ$. Durée $166s$.
	\item Vol vers la droite. Direction  $\sim 8^\circ$ (normale à la direction précédente). Durée $26s$.
	\item Vol vers l'arrière. Direction  $\sim 95^\circ$ (parallèle à la direction initiale). Durée $40s$.\\
\end{itemize}
\figureIRD[Fly_Blassac_trajectoire]{Fly_Blassac_trajectoire.png}{Schéma de la mission FLY-20220125. Les shoot point qui ont servi au mapping sont repérés par les marques rouges. Ils ont un $|dt| \leq 0.5\,s$}{0.8}

Après synchronisation des horloges l'algorithme a formé 116 paires d'images (VIS, NIR). Notons qu'une image NIR sur deux sera associée à deux images VIS successives. 

\figureIRD[Pitch_Yaw_reel_vs_theorique]{Process-and-Theoretical_Offset_Pitch_vol_reel_01.pdf}{Évolution temporelle des angles yaw et pitch pendant la phase d'acquisition. Les angles sont déterminés par les formules théoriques et par le  \og coarse process \fg{}. Les angles théoriques ont été corrigés pour tenir compte des défauts de positionnement de la caméra NIR sur le drone.  FLY-20220125.}{0.99}

La caméra VIS du drone est placée en mode \og gimbal lock \fg{}
\footnote{Dans ce mode particulier il n'est pas certain que  l'information sur le lacet du gimbal soit fiable. La différence entre les angles de lacet du drone et du gimbal est constante et égale à $2,7^\circ$. Pour certain vols cette différence peut atteindre $7^\circ$. Ici l'angle de lacet correspond au roulis des caméras. Cependant en comparant les paires d'images VIS et NIR on ne constate pas de rotation entre les deux. C'est un problème de calibrage de la caméra au moment de passer en gimbal lock.}.
Les informations récupérées dans les données Exif des images du drone indiquent que les angles yaw, pitch et roll du gimbal sont constants ($-78,9^\circ$, $-90^\circ$  et $0^\circ$ respectivement). \\


Le yaw de la caméra NIR est pris égal au  roll du drone. 
Le calcul du yaw théorique à appliquer à l'image NIR pour la recentrer sur l'image visible utilise le même type de formules que celle mises en  évidence pour le pitch théorique (courbe \og Theoretical yaw \fg{} sur la figure \ref{Pitch_Yaw_reel_vs_theorique}). \\ 
L'évolution des valeurs théoriques des angles yaw et pitch suit globalement celle obtenue par le \og coarse process\fg{}.\\ 


Notons qu'il y  a un décalage systématique ($\sim +0,9^\circ$ ) sur le yaw. Ce décalage a été retrouvé sur l'ensemble des vols effectuées jusqu'à présent. Il peut être attribué à un défaut de positionnement de la caméra NIR sur le drone. Un décalage systématique existe aussi pour le pitch \og{}théorique \fg{} par rapport au pitch \og{} process \fg{}. Nous l'attribuons lui aussi à un défaut de positionnement de la caméra NIR sur le drone qui engendre une erreur systématique d'alignement de l'axe de visée de la caméra. Ce défaut résiduel est $\zeta = -1,3^\circ$.
\footnote{Ce que confirme une mesure réalisée à l'aide d'un inclinomètre numérique mesurant les angles avec une précision de $1/10\, ^\circ$. La mesure a été prise entre le support carbone et la surface du filtre NIR.} .\\

L'évolution des angles présente sur certaines parties un aspect oscillant ce que confirme l'analyse spectrale (DFFT) . 
Le spectre de Fourier du pitch représenté sur la figure \ref{Fourier_Pitch_process} montre  un pic net pour la fréquence $0,164\,Hz$ pour les valeurs théoriques ou obtenues par le  \og coarse process \fg{}.  
Cela correspond à une composante périodique de période $T=6\,s$. 
Nous avons montré que ce phénomène est liée à la variation périodique de la \og base line \fg{} (voir la figure \ref{Pitch_cas_non_synchro} et les formules  \eqref{Theoritical_Pitch_non_syncho_0}, \eqref{Theoritical_Pitch_non_syncho_1} ).
Le reste du spectre correspond à l'effet du vent.
Lorsque le drone vol en ligne droite vers l'avant ou l'arrière la projection de sa \og bas-line \fg{} sur son axe est périodique. 
Le pitch  a donc une composante périodique lors de cette phase de vol. 
La projection de sa \og bas-line \fg{} sur l'axe normal est nulle. 
Dans cette phase de vol le yaw a donc seulement la composante aléatoire associée au vent. 
Lorsque le drone se déplace perpendiculairement à son axe (en gardant toujours la même orientation) c'est le yaw qui possède une composante périodique alors que le pitch n'a qu'une composante aléatoire liée  au vent. 
Ce phénomène est clairement visible sur la figure \ref{Pitch_Yaw_reel_vs_theorique}. Les changements d'allure des courbes permettent même d'identifier les phases de vol.\\


\figureIRD[Fourier_Pitch_process]{Spectre-Pitch-process-vol-reel_test.png}{Transformée de Fourier du Pitch  calculé  par le \og coarse process\fg{} à partir des images numérique ou par les formules théoriques.  FLY-20222501.}{0.9}

Notons que le pitch théorique calculé avec la formule synchronisée \eqref{Theoritical_Pitch_syncho_2} ne fait intervenir que les angles du drone et du gimbal. Il ne présente pas de caractère périodique comme le confirme sa DFFT. Ceci est cohérent puisque la variation de ces angles est associée aux effets \og{}aléatoires \fg{} du vent sur le drone (turbulence).


\subsection{Conclusion.}


La méthode proposée pour déterminer a priori le lacet et le tangage (angles yaw et pitch) à appliquer à l'image NIR pour la recentrée sur l'image VIS repose sur plusieurs éléments :\\

\begin{itemize}
	\item Les images VIS et NIR doivent être prises en mode time-lapse. 
	La fréquence d'enregistrement de la caméra VIS (celle qui a le pas le plus petit, ici $2\,s$) doit être régulière. 
	Ceci nécessite l'utilisation de carte SD rapides. Toutefois pour le DJI les dates sont enregistrées à la seconde près. On remarque périodiquement des sauts d'une seconde. Il semble que le pas ne soit pas exactement un entier. Pour notre vol FLY-20211109-1m/s on trouve $2,01865 \, s$. Ce phénomène est retrouvé sur tout les vols effectués jusqu'à présent avec une valeur de $2,02 \, s$. Le logiciel IRDrone détermine automatiquement pour chaque vol la période exacte pour obtenir une fréquence constante.  
	On utilise le format de date de Python en  renseignant son attribut \texttt{micosecond}. 
	Notons que les données Exif des dates prévoient (pour certaines) une valeur décimale des secondes  
	La situation pour la caméra NIR est différente. Il existe aussi des sauts mais de la valeur du timelapse ($3\,s$). Les dates sont donc toujours des multiples du timelapse. Ici c'est clairement un défaut d'enregistrement d'une image
	\footnote{Comme on a pu le vérifier lors de la synchronisation manuelle des images, la séquence retenue  comportant justement un tel défaut d'enregistrement.}. Les dates n'ont donc pas été rectifiées. Lors d'un saut de ce type la \og base-line\fg{} est plus grande. Cela reste acceptable si la vitesse du drone est faible. Notons toutefois qu'il faut adapter les formules de calcul pour en tenir compte. A nouveau le code gère automatiquement ce phénomène.\\
	
	\item Les données d'attitude  (yaw, pitch, roll) doivent être disponibles pour les deux caméras. C'est bien le cas avec notre montage puisque les données Exif des images VIS (format dng) du DJI fournissent ces angles pour le drone (donc pour la caméra NIR qui lui est fixée rigidement) et pour le gimbal (donc pour la caméra VIS qui est fixée dessus).\\  
	\item Le calcul de la \og base-line \fg{} nécessite de connaître les coordonnées UTM. Actuellement l'altitude $H$ du drone par rapport au sol est déterminée en connaissant l'altitude du point de décollage, l'altitude du drone par rapport au point de décollage et l'altitude à l'aplomb de la prise d'image (via une API de l'IGN pour la France). .\\
	\item La synchronisation des horloges des deux caméras est un point critique. On utilise une séquence de synchronisation à effectuer à chaque décollage. Elle consiste à faire tourner
	\footnote{La cible est placée sur un plateau mis en rotation par un moteur pas à pas. Ce moteur est commandé par un micro contrôleur ARDUINO.} un QR code (ARUCO marker) sous les deux caméras. Un algorithme détecte ensuite ce QR code sur les images et en déduit l'orientation. Une méthode de minimisation permet de rapprocher  les courbes angle/date. La  précision de cette synchronisation est supérieure au $1/10^{\grave{e}me}s$.  Soulignons qu'il est impératif de ne pas interrompre la prise de vue en mode hyperlapse pendant toute la durée de la mission (phase de synchronisation et phase d'acquisition). En effet la valeur de synchronisation n'est valable (avec cette précision) que pour une séquence hyperlapse donnée. Entre deux séquences hyperlapse la valeur de synchronisation peut varier de plus d'une seconde (allant même jusqu'à 3s dans le cas le plus défavorable).
	Un tel défaut de synchronisation des horloges risque de fausser l'appariement des images (et le calcul de la  \og base-line \fg{}).
	Pour l'exemple présenté ici nous avons réalisé aussi cette synchronisation manuellement à partir d'un objet mobile présent sur 25 images consécutives (soit 50s). Le mobile avait été déplacé spécialement pour cet utilisation lors de notre vol d'essai. Cette technique est précise (précision de l'ordre du $1/4 s$) mais fastidieuse. Nous avons retrouvé le même appariement des images que celui obtenu avec la synchronisation par cible ARUCO. \\
	\item L'angle de pitch du drone ainsi que la \og base-line \fg{} à l'instant où l'image NIR est capturée ont été  interpolés linéairement. Cette méthode donne d'excellents résultats.\\
\end{itemize}

Nous avons vérifié la pertinence de l'ensemble de ces éléments sur plus d'une dizaine de vol avec des conditions météorologiques variées, notamment le vent.  
Par exemple les figures \ref{Pitch_Yaw_reel_vs_theorique_FLY2021109} et \ref{Fourier_Pitch_process_FLY2021109} correspondent à un vol programmé en hyperlapse (\og mode line\fg{}) avec une vitesse de $1\,m/s$ et environ $600$ images.

\figureIRD[Pitch_Yaw_reel_vs_theorique_FLY2021109]{Process-and-Theoretical_Offset_Pitch_vol_reel_02.png}{Évolution temporelle des angles yaw et pitch pendant la phase d'acquisition. Ici  le drone se déplace à cap constant ($179^\circ$) parallèlement à son axe. FLY-20211109.}{1}

\figureIRD[Fourier_Pitch_process_FLY2021109]{Spectre-Pitch-process-vol-reel_02.png}{Transformée de Fourier du Pitch. FLY-20211109.}{.75}




\chapter{Géoréférencement et orthorectification (Step 6)}\label{ann:georef}
\todo{Texte technique : géométrie caméra/drone, repères, projection sur le MNT, précision, destiné aux utilisateurs avertis.}

\end{document}
"
% (voir build_help.bat). Les images viennent alors de images_light/ au lieu de images/.
\ifdefined\IRDlight
  \newcommand{\IRDimgdir}{images_light}
  \graphicspath{{images_light/}}
\else
  \newcommand{\IRDimgdir}{images}
  \graphicspath{{images/}}
\fi

% Symboles des boutons du logiciel (pdflatex ne les connait pas)
\DeclareUnicodeCharacter{25C0}{\ensuremath{\blacktriangleleft}}
\DeclareUnicodeCharacter{25B6}{\ensuremath{\blacktriangleright}}
\DeclareUnicodeCharacter{2699}{\ensuremath{\circledast}}
\DeclareUnicodeCharacter{2922}{\ensuremath{\nearrow}}
\DeclareUnicodeCharacter{21BA}{\ensuremath{\circlearrowleft}}
\DeclareUnicodeCharacter{21BB}{\ensuremath{\circlearrowright}}
\DeclareUnicodeCharacter{25B2}{\ensuremath{\blacktriangle}}
\DeclareUnicodeCharacter{25BC}{\ensuremath{\blacktriangledown}}
% Bouton du logiciel
\DeclareRobustCommand{\bouton}[1]{\texorpdfstring{\fbox{\small\textsf{#1}}}{#1}}
% A completer
\newcommand{\todo}[1]{\textcolor{red}{[\`A r\'ediger : #1]}}
% Sections historiques : repliees dans le manuel web, texte normal dans le PDF
\newcommand{\legacybegin}[1]{}
\newcommand{\legacyend}{}
% Figure tolerante : affiche un cadre si l'image n'existe pas encore.
% Usage : \figureIRD[label]{fichier}{legende}{largeur}
%   - label optionnel (ex. fig:img_VIR_1), utilisable avec \ref{...} / \autoref{...}
%   - fichier avec ou sans extension ; sans extension, essaie .png, .jpg, .jpeg puis .pdf
\makeatletter
\newcommand{\ird@file}{}
\newcommand{\ird@find}[1]{%
  \def\ird@file{}%
  \IfFileExists{\IRDimgdir/#1}{\def\ird@file{#1}}{%
  \IfFileExists{\IRDimgdir/#1.png}{\def\ird@file{#1.png}}{%
  \IfFileExists{\IRDimgdir/#1.jpg}{\def\ird@file{#1.jpg}}{%
  \IfFileExists{\IRDimgdir/#1.jpeg}{\def\ird@file{#1.jpeg}}{%
  \IfFileExists{\IRDimgdir/#1.pdf}{\def\ird@file{#1.pdf}}{}}}}}}
\newcommand{\figureIRD}[4][]{%
  \ird@find{#2}%
  \begin{figure}[htbp]\centering
    \ifx\ird@file\@empty
      \fbox{\parbox{0.6\linewidth}{\centering\small\textcolor{gray}{Image manquante : \detokenize{#2}}}}%
    \else
      \includegraphics[width=#4\linewidth]{\ird@file}%
    \fi
    \caption{#3}%
    \if\relax\detokenize{#1}\relax\else\label{#1}\fi
  \end{figure}}
% Paire de figures cote a cote, placee a l'endroit exact ou elle est ecrite ([H]).
% Usage : \figurePaireIRD{label}{fichier gauche}{legende gauche}{fichier droit}{legende droite}{legende globale}
%   - le label designe la paire ; les sous-figures ont pour labels <label>_avant et <label>_apres
\newcommand{\ird@incsub}[1]{%
  \ird@find{#1}%
  \ifx\ird@file\@empty
    \fbox{\parbox{0.8\linewidth}{\centering\small\textcolor{gray}{Image manquante : \detokenize{#1}}}}%
  \else
    \includegraphics[width=\linewidth]{\ird@file}%
  \fi}
\newcommand{\figurePaireIRD}[6]{%
  \begin{figure}[H]\centering
    \begin{subfigure}[t]{0.48\linewidth}\centering
      \ird@incsub{#2}%
      \caption{#3}\label{#1_avant}
    \end{subfigure}\hfill
    \begin{subfigure}[t]{0.48\linewidth}\centering
      \ird@incsub{#4}%
      \caption{#5}\label{#1_apres}
    \end{subfigure}
    \caption{#6}\label{#1}
  \end{figure}}
\makeatother

\title{IRdrone\\Manuel de l'utilisateur}
\author{}
\date{\today}


\begin{document}
	
	% --------------------------------------------------
	% Cover page
	% --------------------------------------------------
	\begin{titlepage}
		\centering
						
		\vspace*{2cm}		
		%{\Huge\bfseries IRdrone\par}		
		%\vspace{0.5cm}		
		{\LARGE Manuel de l'utilisateur\par}
		\vspace{6cm}		
		\includegraphics[width=0.95\textwidth]{Logo-IRdrone.png}		
		
		\vfill		
		{\large Version du \today\par}		
	\end{titlepage}
	
	% --------------------------------------------------
	% Table of contents
	% --------------------------------------------------
	\tableofcontents



%==========================================================
\chapter{Le projet IRDrone}\label{ch:overview}


\section{Observer le couvert végétal.}

\subsection{Contexte général.}
La cartographie de la végétation est aujourd'hui classiquement employée dans de nombreux types de diagnostics environnementaux (site Natura 2000, ZNIEFF, études d'impact environnemental etc.). 
Cette cartographie est basée sur une segmentation (détourage d'ensembles homogènes au plan de la végétation) du couvert végétal en vue plane aérienne. 
Cette approche est généralement complétée par des prospections au sol (relevés botaniques, phytosociologiques etc.). 
La segmentation d'ensembles homogènes repose sur l'analyse des photographies aériennes en couleur naturelle, notamment sur leur couleur, leur texture, leur structure apparente, leur opacité etc.

\figurePaireIRD{fig:etude_de_cas}
{step0_carto_aerienne.jpg}{Délimitation des habitats sur fond image aérienne.}
{step0_habitats.jpg}{Carte restituée des habitats sur fond IGN topo 25.}
{Exemple d'exploitation de données terrain à l'aide de cartes et photos aériennes.}

Il est toutefois délicat de distinguer les essences (arbres) et les espèces végétales (espèces herbacées) sur des photographies aériennes classiques. Les dynamiques et les éléments fins du paysage sont également difficiles à cerner. 
Les photographies aériennes IR disponibles sont souvent anciennes, donc difficilement utilisables dans le cadre de diagnostics actuels. 
D'autre part la couverture nationale est lacunaire, avec des régions mal couvertes. 
Ajoutons que toutes les photos existantes sont réalisées en altitude (3-4000 m) et manquent donc de précision pour des diagnostics réalisés couramment à l'échelle du 1/5000, voire moins. \\




L'utilisation d'un drone pour acquérir des photographies présente les avantages de centrer les clichés sur les secteurs d'études, d'obtenir des images d'une précision largement supérieure (altitude de vol de 120 m), d'obtenir autant d'images récentes qu'on le souhaite, de pouvoir les réaliser à différentes saisons. 


\subsection{Intérêt des images multispectrales.}

Les photographies aériennes dans le proche infrarouge
\footnote{On ne confondra pas ce rayonnement infrarouge proche qui en journée provient du rayonnement solaire et est réfléchi par les objets avec le rayonnement infrarouge lointain qui lui est émis par tous les objets.} NIR sont plus riches en informations écologiques que les photographies en couleur naturelle. 
L'objectif principal de la photographie aérienne NIR est donc de faciliter l'interprétation des photographies aériennes en se basant sur les différences de facteur de réflexion du rayonnement infrarouge entre les types de végétations (en raison de différences structurelles entre les tissus foliaires des espèces végétales).

\figureIRD[fig:img_NDVI_1]{img_NDVI_1.png}{Carte de l'indice NDVI de développement des végétaux.}{0.7}
\figureIRD[fig:img_VIR_1]{img_VIR_1.png}{Exemple d'image composite de type VIR. Les couches R/G/B sont remplac\'ees par les couches NIR/R/G. Ce type de représentation est très utile, entre autres, pour distinguer entre résineux et feuillus.}{0.7} 


\subsection{Que cherchons-nous à observer ?}

L'emploi des photographies aériennes se justifie pour plusieurs raisons, qui ont trait essentiellement à la précision et à la rapidité d'analyse. 
D'une manière générale, sur les images aériennes NIR les contrastes entre les différents types d'habitats sont exacerbés, ce qui permet une lecture plus fine du paysage. 
Dans les peuplements forestiers denses, il est difficile d'estimer les proportions de recouvrement résineux/feuillus en se basant sur une prospection de terrain en raison des grandes surfaces à parcourir et de la hauteur des arbres. 
La différenciation entre résineux/feuillus sur les photographies en couleur naturelle est souvent ambiguë. 
Au contraire, la photographie aérienne NIR permet de discerner précisément et rapidement les contributions de chacune des essences dans les peuplements mélangés couvrant des superficies considérables, et ainsi d'avoir accès à une information concernant la gestion spécifique des peuplements (types de traitements sylvicoles, classes d'âge etc.).  D'autre part, certains peuplements arborescents mêlant feuillus et résineux en proportions variables peuvent, selon l'importance prise par les feuillus, relever ou non d'une protection européenne (directive \og Habitats \fg). 
L'évaluation des proportions des essences devient objective et précisément quantifiable à l'aide des photographies aériennes NIR. 
L'analyse des photographies aériennes NIR permet en outre une meilleure appréhension des phénomènes de dynamique de la végétation.
 
La détection des très jeunes ligneux devient possible en raison de leur vive coloration et de leur position au sein de milieux ouverts (fronts de colonisation ligneux au sein d'espaces ouverts) alors que leur présence est difficilement détectable en couleur naturelle en raison de la petite taille de leur houppier et de leur couleur comparable à celle des habitats ouverts colonisés. 
Une autre application importante concerne la localisation instantanée et exhaustive des masses d'eau (qui peuvent être discrètes, de petite taille et donc passer inaperçues en couleur naturelle) qui apparaissent comme des taches sombres  en NIR. 
Cette application est particulièrement importante en raison du statut réglementaire des zones humides, qui constituent généralement un des enjeux principaux des diagnostics environnementaux. 


\subsection{Principe de l'observation.}

Les végétaux produisent de la chlorophylle en absorbant la lumière visible du soleil (VIS). Seule une petite partie de ce rayonnement est utilisée dans ce mécanisme. Or le rayonnement solaire arrivant au sol comporte aussi des ultra-violets (UV) et des infrarouges proches (NIR) \footnote{Les rayonnements UV et NIR ne sont pas visibles par l'\oe il humain.} 

\figureIRD[Facteur_reflexion_chloro-carot]{Spectres-absorption-pigments-photosynthetiques.pdf}{Facteur d'absorption des pigments photosynthétiques dans la bande spectrale visible.}{1}

C'est essentiellement le rayonnement visible dans les bandes \og bleu \fg{} ($[400\,nm ,\, 450\,nm]$) et \og rouge \fg{} ($[630\,nm ,\, 680\,nm]$) qui contribue à la production chlorophyllienne.  
La partie non absorbée du rayonnement par la plante est transmise (la plante est \og transparente \fg{}) ou bien réfléchie (renvoyée vers l'environnement). 
Le rendement chlorophyllien étant nul dans la bande spectrale \og jaune-vert \fg ($[500\,nm ,\, 600\,nm]$) celle-ci est réfléchie par la plante (\ref{Facteur_reflexion_chloro-carot}). 
C'est de là que provient la couleur verte de ces végétaux. De même les rayonnements UV et IR ne participent pas à la production de chlorophylle. 
En revanche ils produiraient son échauffement s'ils étaient absorbés. 
Ainsi pour se protéger les plantes n'absorbent pas les rayonnements UV et IR émis par le soleil, qui sont alors transmis ou réfléchis. 
Les plantes sont totalement transparentes aux rayonnements ultraviolets.
En revanche elles réfléchissent une partie du rayonnement infrarouge.
C'est cette propriété qui va être exploitée dans le projet. 
Les résineux et les feuillus n'ont pas le même facteur de réflexion du rayonnement infrarouge. 
La figure \ref{Facteur_reflexion_vegetaux} montre les courbes typiques du facteur de réflexion spectral des résineux et des feuillus. 
Notons que les feuillus produisent aussi de la bêta-carotène et de la fucoxanthine qui absorbent le bleu et une partie du vert mais réfléchissent le rouge. 
C'est pour cette raison qu'à l'automne, lorsque les feuillus ne produisent plus de chlorophylle mais que la bêta-carotène est encore présente, ces espèces arborent des teintes allant du marron orangé à l'ocre rouge.

\figureIRD[Facteur_reflexion_vegetaux]{Facteur_reflexion_vegetaux.pdf}{Facteur de réflexion des résineux et des feuillus. Dans la bande spectrale visible $[400\,nm ,\,700\,nm]$ le comportement des deux espèces est très similaire. En revanche dans la bande spectrale infrarouge IR-A $[700\,nm ,\, 1.4\,\mu m]$ bien présente dans le spectre solaire les deux espèces se distinguent nettement.}{1}


On remarque que les feuilles des arbres réfléchissent le rayonnement autour de 550 nm qui correspond à la couleur verte perçue par l'\oe il humain. Mais les courbes étant identiques dans la partie visible du spectre (400 nm à 700 nm) on ne peut pas distinguer entre les différentes essences. En revanche les facteurs de réflexion dans le domaine infrarouge sont bien distincts entre résineux et feuillus. L'idée est d'exploiter cette propriété en produisant une image infrarouge  du couvert végétal
\footnote{Éventuellement combinée à une image dans le spectre visible.}. Sur cette image les feuillus devraient apparaître  \og plus clairs\fg{}  que les résineux. On pourra ensuite tracer une carte du couvert végétal. Signalons toutefois que selon l'époque de l'année et l'âge des sujets ces courbes peuvent évoluer rendant plus délicate l'interprétation des images.

\figureIRD[Facteur_absorption_vegetaux]{Par_action_spectrum.jpg}{Facteur d'absorption des végétaux et rendement de photosynthèse.}{0.5}


\subsection{Le matériel pour l'observation.}
Pour tracer une carte il faut des images \og vues du ciel \fg{}. Les professionnels utilisent des photos satellites (alt. 300km) ou prises d'avion (alt. 2000-3000 m). Mais comme expliqué plus haut, la résolution spatiale et temporelle n'est pas adaptée à des études fines de terrain. Dans le cadre de ce projet on envisage donc d'utiliser un petit drone qui permet d'obtenir avec une très grande souplesse des vues aériennes des zones à étudier (alt. 120 m) avec une très bonne résolution spatiale et une maîtrise totale des périodes de prise de vue.

La plupart des \og petits \fg{} drones existants ne permettent pas l'obtention simultanée de clichés en couleur naturelle et en IR. Or il est indispensable de disposer des deux types de clichés géoréférencés (IR pour interprétation de la végétation et couleur naturelle pour la réalisation de cartes standardisées). \\ 

Il existe des drones et des caméras multi-spectrales pour drones
\footnote{Par exemple DJI Phantom 4 multispectral, Parrot Sequoia, Micasense RedEdge.} . 
Il s'agit de caméras multi canaux  qui permettent d'acquérir, outre l'image visible classique, des bandes spectrales étroites dans le rouge, le vert, le bleu et dans une ou deux bandes dans la proche infrarouge. Il existe même des versions équipées de caméra thermiques \footnote{Par exemple la caméra Micasense Altum  possède un capteur Flir IR. Son poids est de 400g. Cependant il s'agit ici de rayonnement infrarouge thermique (LWIR) dont les longueurs d'onde sont dix fois plus grandes que celles du  proche infrarouge solaire.  Soulignons que ce matériel est très onéreux.}.
D'autres drones
sont équipés de deux caméras l'une dans le visible et l'autre dans l'infrarouge. 
Enfin la caméra de certains drones
\footnote{Par exemple le DJI Mavic Pro.2 équipé d'une très performante caméra Hasselblad. Son filtre \og hot mirror\fg{} est fixé directement sur le capteur. Des drones \og lourds\fg{} munis de nacelles peuvent embarquer des reflex numériques qui eux aussi peuvent être \og défiltrés \fg{}. }  
peut être rendue sensible dans le proche infrarouge. Cela nécessite de démonter
\footnote{Cette opération de \og défiltrage \fg est délicate. Elle peut être réalisée par des sociétés spécialisées.} le filtre interne du capteur numérique qui coupe le rayonnement infrarouge et de munir le système optique d'un filtre adapté
\footnote{Avec un jeu de filtres externes on peut choisir le type de bande spectrale IR. Avec un filtre passe haut IR $760\,\mathrm{nm}$ on obtiendra une bande large $[760, 1100]$. Il existe aussi des filtres à bande étroite, par exemple $850\,\mathrm{nm}\pm 25\mathrm{nm}$). On peut même remettre un filtre \og hot mirror\fg{} qui rétablira l'imagerie dans le spectre visible.}.  
De plus le drone n'enregistre alors qu'un seul type d'image selon le filtre dont il est équipé. \\  


L'idée du projet \textsf{IRDrone} est d'équiper un petit  drone de coût abordable 
\footnote{Ici un DJI Mavic Air.2  dont le poids n'excède pas 600 g.} d'une seconde  caméra. La caméra équipant le drone fournira des images en couleur dans le spectre visible (RGB). La deuxième caméra embarquée, que nous désignerons par caméra IR, devra produire des images dans la bande spectrale du proche infrarouge \og solaire \fg 
\footnote{Concrètement avec les capteurs CMOS disponibles la bande spectrale accessible s'étend de $720 \,\mathrm{nm}$ à $1100\,\mathrm{nm}$. Elle englobe les bandes étroites RedEdge et NIR (pour Near Infra Red)}.\\

La caméra NIR devra être légère pour ne pas perturber le vol du drone et ne pas diminuer son autonomie. Elle devra être d'un coût réduit ce qui conduit naturellement à se tourner soit vers des caméras de surveillance (\og vision de nuit\fg{}) soit vers des \og action cam \fg{} grand public. Les premières ne conviennent pas car elles offrent des performances médiocres en résolution. Les secondes sont conçues pour donner de belles images dans le domaine visible. Il faudra donc adapter l'optique par ablation du filtre \og hot mirror\fg{} et mise en place d'un filtre IR.  \\


Après chaque vol il sera nécessaire de faire un traitement informatique des images collectées (corrections géométriques, colorimétriques, calage et fusion d'images visibles et IR ...). Ce sera le rôle dévolu au code Python  IRDrone.py associé au projet \textsf{IRDrone} dont les fonctionnalités sont  exposées ci-après.


\section{Organisation d'une mission IRDrone}

\subsection{Préparation au bureau}

Avant le début de la partie active de la mission le pilote doit préparer soigneusement son vol. 
D'abord au bureau, en repérant les chemins d'accès, la topographie du terrain, les obstacles et la dimension de la zone à étudier. 
Cette dimension conditionne l'altitude de vol. On peut imaginer ne faire qu'un vol unique à haute altitude pour couvrir en une fois la zone, ou bien plusieurs vols successifs. Il faut garder à l'esprit que le drone DJI chargé de sa caméra infrarouge a une durée de vol limitée à 25 minutes. 
En pratique la phase de synchronisation prend 2 minutes. Avec une vitesse de $3\,\mathrm{m/s}$ le drone peut parcourir environ 4\,km, ce qui n'est pas énorme compte tenu de la nécessité de \og balayer \fg~la zone de survol afin que les clichés se superposent (si on souhaite faire un assemblage des images). 
On suppose aussi que l'opérateur a sa qualification de pilote de drone et les autorisations de survol.

\subsection{Préparatifs du vol sur le terrain}
Une fois sur le terrain le pilote installera la zone de décollage. Après balisage il positionne le stand de décollage et le système de synchronisation par mire ArUco avec son alimentation
\footnote{Que faire si on a oublié le dispositif de synchronisation ? On peut \og sauver \fg~la mission en déplaçant lentement un objet sous le drone. Il faudra ensuite estimer manuellement l'offset de synchronisation des deux caméras. Ça peut marcher\dots{} des fois !}. 
Un siège pliant est utile pour piloter en regardant le retour caméra afin de suivre au mieux le plan de vol. 
De même une petite manche à air pour apprécier la direction et la vitesse du vent est utile. 
Enfin, mesurer précisément la position du point de décollage avec un GPS en position fixe est appréciable\footnote{Voir les possibilités d'un GPS Garmin 750, qui a ce type de fonction.}. 
Avant chaque mission sur le terrain \textbf{le GPS du drone doit être calibré et sa caméra étalonnée} (direction de visée). 
Ces deux actions sont indispensables pour la réussite de l'exploitation des images par le logiciel IRDrone. 
Le pilote vérifie que les batteries sont chargées et que les cartes SD (vides) sont bien présentes sur le drone DJI et la caméra SJCam M20. 
Il est judicieux de prérégler les paramètres de la caméra du DJI (ceux qui peuvent l'être) avant le décollage. 
À ce stade le pilote peut déclencher l'hyperlapse de la caméra infrarouge (SJCam M20), positionner le drone sur son stand de décollage et enfin mettre sous tension le dispositif de synchronisation (dans la dernière version du dispositif, la mire ne se met à tourner que lorsque le drone a quitté le stand de décollage).

\subsection{Le vol de la mission}
Une fois ces préparatifs faits, la mission comporte 4 phases successives :
\begin{enumerate}
	\item \emph{Takeoff Point -- Synchronization} : synchronisation à proximité immédiate du point de décollage, à une altitude d'environ 2\,m. 
	Au début de cette phase la caméra NIR est déjà déclenchée en mode hyperlapse (fréquence de prise de vue 3\,s). 
	Pendant cette phase le pilote stabilise le drone au-dessus du dispositif de synchronisation, règle les paramètres de la caméra VIS et déclenche la prise de vue VIS en mode hyperlapse. Il reste au point fixe une dizaine de secondes.
	\item \emph{Approach} : déplacement vers la zone à survoler pour l'étude. 
	Le pilote prend rapidement de l'altitude. 
	Il se déplace à la vitesse maximum permise par le mode hyperlapse du DJI. 
	Une valeur de $3\,\mathrm{m/s}$ est un bon choix. Lors de la préparation de la mission il est fortement conseillé de choisir un point de décollage au plus près de la zone de survol.
	\item \emph{Hyperlapse Survey} : acquisition des images multispectrales au-dessus de la zone de survol (\emph{survey area}). 
	Le drone se déplace en gardant une orientation constante. 
	Il avance puis se déplace latéralement puis recule, etc. (voir figure~\ref{fig:4_step_of_the_mission}). 
	Lors du décollage il est conseillé d'orienter le drone dans le sens de déplacement principal au-dessus de la zone de survol. 
	À la sortie de la zone de survol le pilote peut interrompre l'enregistrement hyperlapse de la caméra VIS. 
	Ceci a deux avantages : cela évite de saturer la carte SD, les images étant en haute résolution, mais surtout cela libère la contrainte de vitesse du drone imposée par le mode hyperlapse.
	\item \emph{Return to Home} (RTH) : c'est la phase de retour au point de décollage. Ce retour peut se faire à vitesse maximum. Le mode RTH peut être enclenché pour un vol automatique.
\end{enumerate}

\figureIRD[fig:4_step_of_the_mission]{4_step_of_the_mission.pdf}{Les phases de vol d'une mission IRDrone sur le terrain.}{0.8} 

\subsection{Exploitation des images}

Une fois l'acquisition terminée, la phase suivante consiste à traiter numériquement les images des deux caméras à l'aide du logiciel IRDrone. 


\section{Objectif du logiciel IRDrone}
IRDrone.py est une application interactive et conviviale qui permet de générer des images multispectrales pour des études d'impact environnemental ou de botanique. 
Elle utilise pour cela les images prises lors des campagnes sur le terrain par deux caméras, l'une dans le spectre visible (\og VIS \fg) et l'autre dans le spectre du rayonnement dans le proche infrarouge (\og NIR \fg). 
Les deux caméras sont fixées sur un drone. 
Actuellement le système supporte un drone DJI Mavic Air 2 muni de sa caméra VIS montée sur un gimbal et une caméra SJCam M20 modifiée pour capturer des images dans le spectre du proche infrarouge (typiquement au-delà de 810 nm). 
L'application synchronise les horloges des deux caméras puis aligne les images VIS et NIR.
Le processus principal produit des images multispectrales alignées (4 bandes red, green, blue, nir au format tif). 
Un post-traitement permet d'orthorectifier et géoréférencer les images, soit unitairement, soit sous forme d'assemblage (mapping avec le code ODM ; figure~\ref{fig:img_mapping_1}).
\figureIRD[fig:img_mapping_1]{img_mapping_sample.png}{Assemblage à l'aide d'Open Drone Map (ODM) d'images dans le spectre VIS. La visualisation est faite avec le logiciel QGIS. Les courbes de niveau ont une équidistance de 1\,m. Les points rouges indiquent l'emplacement des clich\'es et le numéro de la prise de vue.}{0.8}
L'exploitation peut être réalisée avec le logiciel QGIS afin d'obtenir des images spécialisées (indice NDVI, figure~\ref{fig:img_NDVI_1} ; image VIR nir/red/green,
figure~\ref{fig:img_VIR_1}). 


\section{Fonctionnalités principales du logiciel IRDrone}
Au lancement, une petite fenêtre \og Loading, please wait \fg{} s'affiche pendant le chargement (quelques dizaines de secondes), puis la fenêtre principale s'ouvre (chapitre~\ref{ch:home}). Elle liste les missions du dossier de travail et, une fois une mission ouverte, ses sept étapes avec leur avancement. Un clic sur une étape la lance ; les étapes 1 à 4 s'enchaînent d'elles-mêmes.

L'exploitation d'une mission comporte 7 étapes :
\begin{enumerate}
	\item \textbf{Mission} : les images des deux caméras sont trouvées sur les cartes SD, copiées dans un nouveau dossier de mission, et les paramètres de la mission sont enregistrés.
	Le nom du dossier de mission est généré automatiquement en fonction de la localisation et de la date
	\footnote{\label{fn:MissionFolderName} Par défaut le nom du dossier est FLY-YYYYMMDD-hhmm-location. Il peut être modifié une fois pour toutes à la création de la mission dans le champ \og Location \fg. Par exemple si le champ \og Location \fg~a la valeur Blassac on peut préciser \og Blassac Bourg \fg.}.
	\item \textbf{Ligne de temps du vol} : les images de synchronisation (mire ArUco au décollage) et celles du relevé au-dessus de la zone d'étude sont détectées automatiquement ; l'utilisateur vérifie.
	\item \textbf{Synchronisation} : le décalage entre les horloges des deux caméras est mesuré et les images VIS et NIR sont appariées.
	\item \textbf{Alignement}, c\oe ur d'IRDrone : chaque paire d'images est alignée pour obtenir 4 couches spectrales superposées.
	\item \textbf{Revue} des images alignées (chapitre~\ref{ch:review}) : chaque image est vérifiée sur une ligne de temps, les mal alignées sont signalées (écartées de la carte) ou réalignées à la main.
	\item \textbf{Cartographie et géoréférencement} : géoréférencement, orthorectification, assemblage en orthophoto (mapping), modèle de terrain et courbes de niveau.
	\item \textbf{Archivage} de la mission.
\end{enumerate}

\figureIRD[fig:Air_Mission_folder]{Air_Mission_folder.png}{Le dossier principal \texttt{Air-Mission} avec deux missions en cours de traitement (ici Chalivoy-Milon et Labout) et une mission archivée (ici Blassac). Sous Windows les icônes sont attribués automatiquement.}{0.6}

\subsection*{Aide en ligne}\label{sec:help}
\begin{itemize}
	\item \textbf{Info-bulles.} Laisser le pointeur de la souris quelques instants sur un bouton (ou une case à cocher) affiche une info-bulle qui résume son action.
	\item \textbf{Touche \textsf{F1}.} La touche \textsf{F1}, le pointeur de la souris sur un bouton, ouvre ce manuel à la section qui décrit ce bouton. Sur une étape de la page d'accueil, elle ouvre le chapitre de l'étape. Ailleurs dans une fenêtre, elle ouvre le chapitre de la fenêtre, comme le bouton \bouton{?}.
	\item \textbf{Bouton \bouton{?}.} Chaque fenêtre a un bouton \bouton{?} qui ouvre ce manuel au chapitre de la fenêtre. Dans les fenêtres qui se pilotent au clavier (ligne de temps, revue), son info-bulle rappelle les touches. Il est en bas à gauche ; le bouton \bouton{Close} des fenêtres d'étape et de la revue est en bas à droite ; aux étapes 2 à 5, le bouton qui passe à l'étape suivante (\bouton{Continue with \ldots ▶}) est à sa droite et ne lance jamais un calcul tout seul. Les actions d'une étape (\bouton{Run alignment}, \bouton{Build the map\ldots}, \ldots) sont dans la fenêtre, au-dessus de la barre de progression ; la principale est bleu foncé, les autres blanches, et celle qui est irréversible (l'archivage) rouge. Les étapes~3 à~6 affichent en haut le nom de l'étape et, dessous, la mission. Dans la fenêtre principale aussi, en bas à gauche, il ouvre le manuel entier.
	\item \textbf{Boutons \bouton{↺} et \bouton{Close} des fenêtres qui modifient quelque chose} (détails de la mission, ligne de temps, réglage du décalage, réglages de l'étape~4). Fermer garde ce qui a été changé : dès qu'une valeur change, \bouton{Close} devient \bouton{Save and close} ; fermer la fenêtre par sa barre de titre ou avec \textsf{Échap} fait de même. Le bouton \bouton{↺}, à côté de \bouton{?}, n'apparaît que s'il y a quelque chose à remettre : il remet les valeurs et la fenêtre reste ouverte.
\end{itemize}
Le manuel s'ouvre dans le lecteur PDF par défaut (navigateur, Acrobat), à la bonne page quand le lecteur le permet.

\section{Organisation numérique d'une mission IRDrone}

%\todo{Dossier \texttt{FLY-YYYYMMDD-HHmm-localit\'e}, sous-dossiers \texttt{ImgIRdrone}, \texttt{mapping\_MULTI}, etc.}
On suppose ici que votre ordinateur sous Windows comporte un dossier \texttt{C:/Air-Mission}.
Ce dossier de travail est choisi au premier lancement de l'application ; il est rappel\'e en bas de la fen\^etre
principale et peut \^etre chang\'e avec le bouton \og engrenage \fg{} plac\'e \`a c\^ot\'e.

Chaque mission est placée dans un dossier \texttt{C:/Air-Mission/FLY-YYYYMMDD-hhmm-location}. 

Ce dossier comporte les sous-dossiers :
\begin{enumerate}
	\item \texttt{AerialPhotography}
	      \begin{itemize}
	      	\item \texttt{VIS}
	      	  \begin{itemize}
	      	  	\item \texttt{VIS\_XXXX.dng}. Images de la caméra DJI transférées,
	      	  	\item \texttt{VIS\_XXXX.exif}. Fichiers \og compagnons \fg~des images contenant des données exif,
	      	  	\item \texttt{time\_line.json}. Fichier qui donne la chronologie des images par rapport à l'horloge de la caméra VIS,
	      	  	\item \texttt{transfer\_info\_VIS\_dng.json}. Fichier qui résume les choix de l'étape~2,
	      	  	\item \texttt{VIS\_XXXX\_RawTherapee.tif}. Images\footnote{\label{fn:rawtherapee_1}Ces fichiers très volumineux sont indispensables au traitement numérique des images. Ils sont générés une seule fois à la première utilisation de l'étape~4. Lorsque la mission est entièrement traitée et doit être archivée, ils peuvent être supprimés.} créées à l'étape~4.
	      	  \end{itemize}
	      		
	      	\item \texttt{NIR}
	      	\begin{itemize}
	      		\item \texttt{NIR\_XXXX.dng}. Images de la caméra infrarouge transférées et converties (RAW propriétaire vers dng), 
	      		\item \texttt{NIR\_XXXX.exif}. Fichiers \og compagnons \fg~des images contenant des données exif,   
	      		\item \texttt{time\_line.json}. Fichier qui donne la chronologie des images par rapport à l'horloge de la caméra NIR,
	      		\item \texttt{transfer\_info\_NIR\_dng.json}. Fichier qui résume les choix de l'étape~2,
	      		\item \texttt{NIR\_XXXX\_RawTherapee.tif}. Images\textsuperscript{\ref{fn:rawtherapee_1}} créées à l'étape~4.
	      		\end{itemize}
      		\item \texttt{shoot\_points.json}. Dictionnaire des \emph{shoot points} (prise de vue avec comme référence l'image VIS et l'image NIR la plus proche de l'instant de la prise de vue VIS). Il est créé à l'étape~3.
	    \end{itemize}
	\item \texttt{cameras}
	         \begin{itemize}
	         	\item Non utilisé actuellement.
	         \end{itemize}
	\item \texttt{FlightAnalytics}. C'est un dossier technique.
			\begin{itemize}
				\item \texttt{mission\_parameters.json}. Fichier contenant les paramètres de la mission (modifiable avec la fenêtre \og Mission details \fg, §\ref{btn:step1_details}).
				\item \texttt{flight\_timeline.json}. Phases du vol confirmées à l'étape~2.
				\item \texttt{fly\_vertical\_profil.png}. Graphe qui permet de visualiser les altitudes de prise de vue par rapport au sol.
				\item \texttt{Angles analysis.png},  \texttt{Angles analysis dispersion.png}, \texttt{Angles analysis spectrum.png}. Graphes d'analyse de la qualité du vol. Le premier de ces trois graphes permet par exemple de juger si l'attitude du drone pendant la mission est restée dans les spécifications de vol.
			\end{itemize}
	\item \texttt{ImgIRdrone}
		\begin{itemize}
		\item Pour chaque shoot point de numéro XXXX on trouve : 
			\begin{itemize}
				\item \texttt{VIS\_XXXX\_REDO.bat},
				\item \texttt{VIS\_XXXX.jpg},
				\item \texttt{VIR\_XXXX\_local.jpg},
				\item \texttt{VIS\_XXXX\_motion\_model.json},
				\item \texttt{VIS\_XXXX\_motion\_model.npy}.
			\end{itemize}
		\item \texttt{review.json}. Mesures et points vérifiés de la revue des images alignées (chapitre~\ref{ch:review}).
		\item \texttt{shoot\_points.json}. La liste des shoot points traités avec un indicateur de réussite. Attention : cet indicateur suppose par défaut que l'alignement est correct. L'utilisateur dispose d'une procédure pour examiner les alignements et les déclarer, le cas échéant, défaillants.
		\item \texttt{imgMultispectral}. Dossier contenant les quadruplets d'images multispectrales au format tif. Ici ces images ne sont pas géoréférencées. Pour chaque shoot point XXXX :
			\begin{itemize}
				\item \texttt{IRD\_XXXX\_1.tif}. Image de la couche rouge. 
				\item \texttt{IRD\_XXXX\_2.tif}. Image de la couche verte.
				\item \texttt{IRD\_XXXX\_3.tif}. Image de la couche bleue.
				\item \texttt{IRD\_XXXX\_4.tif}. Image de la couche infrarouge proche.
				\item \texttt{IRD\_XXXX\_1.exif}. Fichier \og compagnon \fg~contenant des données exif\footnote{Les données exif sont aussi présentes dans les images \texttt{IRD\_XXXX\_Y.tif}. Toutefois l'extraction de certaines de ces données exif nécessite l'utilisation d'un logiciel tiers spécifique comme ExifTool. C'est une opération coûteuse qu'on ne réalise qu'une fois. On peut donc considérer les fichiers exif compagnons comme des données qui ont été mises en cache.}.
			\end{itemize}
		\item \texttt{VIS\_XXXX\_align\_traces\_manual}. Dossier contenant toutes les informations et images suite à une opération d'alignement manuel des images VIS et NIR du shoot point XXXX.
		\end{itemize}
	\item \texttt{mapping\_MULTI}. Dossier de travail de l'étape~6 (mapping).
	\item \texttt{Synchro}. Aperçus des images et détections de la mire préparés pour l'étape~2 (\texttt{Thumbnails}, un cache recréé si besoin) ; fichiers de l'étape~3 : images de travail \texttt{*\_RawTherapee.jpg} (sous-dossiers \texttt{VIS} et \texttt{NIR}), images de contrôle de la détection des mires (\texttt{Check\_ARUCO}) et graphiques \texttt{check\_time\_*.png}.
\end{enumerate}

\section{Encha\^inement des \'etapes}
%\todo{Ordre conseill\'e, pr\'erequis de chaque \'etape.}
Vous devez avoir sur un support numérique (carte SD, disque dur) les images des deux caméras prises lors du vol de la mission. Il est fortement conseillé de n'avoir qu'une mission par carte SD. Les images du DJI sont au format dng et stockées dans un dossier \texttt{DCIM/HYPERLAPSE}. Les images de la caméra SJCam M20 sont au format RAW, doublées d'une image jpeg. Comme les images RAW ne possèdent pas de vignette, on peut les choisir en regardant l'image jpeg associée. Les images sont stockées dans un dossier \texttt{DCIM}.

Pour une nouvelle mission il faut suivre dans l'ordre les étapes 1 à 6. 
Cependant les étapes 2 à 6 peuvent être lancées individuellement, à condition que le dossier de la mission comporte au minimum les éléments produits par les étapes qui la précèdent. C'est particulièrement utile pour les étapes 4 à 6.

Les étapes se lancent depuis la fenêtre principale, décrite au chapitre suivant.

%==========================================================
\chapter{La fenêtre principale}\label{ch:home}
La fenêtre principale a deux pages : la liste des missions, puis la mission ouverte. En bas, une barre rappelle le dossier de travail et ouvre ce manuel.

\section{Liste des missions}\label{btn:home_missions}
Sans mission ouverte, la fenêtre propose (figure~\ref{fig:home_welcome}) :
\begin{description}
	\item[\bouton{New mission}] crée une mission à partir des cartes SD du drone et de la caméra infrarouge (étape~1, chapitre~\ref{ch:step1}).
	\item[\bouton{Open a mission}] ouvre une mission déjà créée : choisir son dossier \texttt{FLY-YYYYMMDD-hhmm-location}.
	\item[Missions récentes] les missions du dossier de travail, les plus récentes en premier, avec leur avancement : un point vert par étape faite. Un clic ouvre la mission.
\end{description}
Le bouton \bouton{All missions} de la page de la mission revient à cette liste.

\figureIRD[fig:home_welcome]{home_welcome.png}{La liste des missions : nouvelle mission, ouvrir une mission, missions récentes et leur avancement.}{0.85}

\section{La mission ouverte}\label{btn:home_steps}
La page de la mission (figure~\ref{fig:home_mission}) montre :
\begin{itemize}
	\item en haut, le nom de la mission (un clic ouvre son dossier dans l'explorateur) et le lieu, la date et le pilote ;
	\item les sept étapes, de gauche à droite : coche verte et date pour une étape faite, cadre bleu pour l'étape suivante, gris pour les autres. \textbf{Un clic sur une étape la lance}, même déjà faite (pour la refaire). Deux exceptions :
	\begin{itemize}
		\item l'étape~1 faite ouvre seulement la fenêtre \og Mission details \fg~(§\ref{btn:step1_details}) : les cartes SD ne sont pas relues ;
		\item l'étape~3 faite montre le résultat enregistré (décalage des horloges, courbes alignées) sans rien recalculer (§\ref{sec:step3_chained}).
	\end{itemize}
	La touche \textsf{F1} sur une étape ouvre son chapitre de ce manuel ;
	\item l'étape suivante, expliquée en une phrase, avec son bouton (§\ref{btn:home_continue}) ;
	\item quelques chiffres : nombre d'images, de paires, d'images alignées, de points signalés mal alignés ;
	\item à droite, la carte du vol ou une image alignée (§\ref{btn:home_map}).
\end{itemize}

\figureIRD[fig:home_mission]{home_mission.png}{La mission ouverte : les étapes (un clic lance une étape), étape suivante, chiffres de la mission et carte du vol sur l'orthophoto et les photos aériennes IGN.}{0.9}

\section{Bouton \bouton{Continue with Step N}}\label{btn:home_continue}
Lance l'étape suivante. Les étapes 1 à 4 s'enchaînent d'elles-mêmes :
\begin{enumerate}
	\item à la fin de l'étape~1, \bouton{Flight timeline} ouvre l'étape~2 ;
	\item \bouton{Continue with synchronisation} de l'étape~2 (ou \bouton{Continue with Step 3}) ouvre l'étape~3 ; \bouton{Synchronise the cameras} la lance : détection des mires, décalage des horloges, puis une question, \og Do the two curves overlap? \fg~; \bouton{Yes} apparie les images ;
	\item \bouton{Continue with image processing} ouvre l'étape~4. L'alignement, long, ne démarre jamais tout seul : il se lance avec \bouton{Run alignment} et s'arrête à tout moment avec \bouton{Stop} ;
	\item à la fin, \bouton{Continue with the review} ouvre l'étape~5, la revue des images alignées (elle s'ouvre d'elle-même quand l'alignement a été enchaîné depuis l'étape~3) ;
	\item \bouton{Continue with mapping} de la revue ouvre l'étape~6, où \bouton{Build the map\ldots} lance la carte ;
	\item une fois la carte construite, \bouton{Continue with archiving} ouvre l'étape~7, qui indique l'espace libéré par l'archivage ; rien n'est supprimé avant votre confirmation.
\end{enumerate}
Il reste donc à vérifier les phases du vol, la superposition des courbes de synchronisation, puis le résultat. Une fois toutes les étapes faites, le bouton devient \bouton{Archive the mission} (étape~7).

\section{Carte du vol}\label{btn:home_map}
La carte est disponible dès l'étape~1 : les positions sont lues dans les EXIF des images visibles. Elle montre :
\begin{itemize}
	\item la trajectoire du drone et un point par image, coloré selon les phases confirmées à l'étape~2 (violet : synchronisation, vert : relevé) ;
	\item les images de la carte, et elles seules : la trajectoire les relie, le vol complet restant en pointillés discrets, et tant qu'aucune orthophoto n'est affichée leur emprise au sol est dessinée en vert : on voit le recouvrement de la carte. Quelles images :
	\begin{itemize}
		\item avant l'étape~4, les images retenues par le \emph{best mapping} (étapes~2 et~4, section~\ref{btn:step2_best_mapping}), s'il est choisi ;
		\item une fois l'étape~4 faite, les points exactement comme dans la revue des images alignées (chapitre~\ref{ch:review}) : chaque point de prise de vue dans la couleur de son statut (vert aligné, orange à vérifier, rouge signalé\ldots), les points que l'étape~4 n'a pas alignés en petits points gris ; une image alignée ou réalignée depuis la revue apparaît aussitôt ;
		\item l'orthophoto de l'étape~6, une fois construite, est dessinée sous ces points (les mêmes points que la revue, pas seulement les images qu'elle a utilisées) ;
	\end{itemize}
	les images signalées (\bouton{Flag}, revue des images alignées, chapitre~\ref{ch:review}) sont de petits points rouges ; en survoler un montre son emprise en pointillés rouges : on voit le trou qu'il laisse dans le recouvrement ; signalées après l'étape~6, elles sont encore dans l'orthophoto affichée et en seront retirées quand l'étape~6 sera relancée (\texttt{left out when Step 6 runs again}). Le nombre d'images et d'images signalées est écrit en bas à droite. \bouton{Layers} > \bouton{Only the images of the map (and the flagged ones)} (\bouton{Only the images of the best mapping} avant l'étape~4) remontre toutes les images ;
	\item l'orthophoto de l'étape~6, une fois le mapping fait : dès que l'étape~6 l'a construite, la fenêtre principale l'affiche sur la carte (même si \bouton{Layers} > \bouton{Hidden} était choisi avant) ;
	\item en fond, des photos aériennes : IGN pour la France, Esri World Imagery ailleurs, ou OpenStreetMap.
\end{itemize}
Survoler un point montre la vignette de son image ; un clic sur un point montre l'emprise au sol de l'image (le trait épais est le haut de l'image). Molette : zoom ; glisser : déplacement ; double-clic ou bouton \bouton{⤢} : tout le vol.

Le menu \bouton{Layers} choisit le fond (IGN, Esri, OpenStreetMap ou aucun) et l'affichage de l'orthophoto, en couleurs naturelles ou en fausses couleurs infrarouges. Le fond vient d'internet : hors ligne, la carte reste utilisable (vol, emprises, orthophoto), seul le fond manque, et rien n'est bloqué. Les tuiles déjà vues sont gardées dans un cache (\texttt{\textasciitilde/.cache/irdrone} sous Linux, \texttt{\%LOCALAPPDATA\%\textbackslash IRDrone} sous Windows) et s'affichent ensuite hors ligne.

L'emprise est calculée avec la hauteur au-dessus du sol (modèle de terrain de l'étape~6, ou altitudes IGN) et le cap de l'image ; une fois le mapping fait, le cap vient d'OpenDroneMap, plus juste que la boussole du drone (souvent fausse de 5 à 10\textdegree).


\section{Barre du bas : dossier de travail et manuel}\label{btn:home_workdir}
Sur la page de la liste des missions, la barre du bas rappelle le dossier de travail (le dossier qui contient les missions, par défaut \texttt{C:/Air-Mission}) et la place libre sur son disque ; un clic sur le dossier l'ouvre dans l'explorateur. La roue dentée \bouton{⚙} permet de changer de dossier de travail : la liste des missions est alors relue. Elle donne aussi une option avancée, décochée par défaut : saisir à la main le décalage d'horloge des caméras à l'étape~3, pour les vols sans cible ArUco (§\ref{btn:step3_clock_offset}). Le dossier est choisi au premier lancement ; la variable d'environnement \texttt{IRDRONE\_WORKING\_DIR}, si elle est définie, reste prioritaire. Sur la page d'une mission, la barre du bas indique à droite la taille totale du dossier de la mission (\texttt{Mission size}), mise à jour à chaque étape terminée et toutes les minutes.

\bouton{?}, à gauche, ouvre ce manuel (§\ref{sec:help}).

\section{Rapport de diagnostic}\label{btn:home_diagnostic}
Quand quelque chose ne marche pas sur un ordinateur, le rapport de diagnostic rassemble dans un seul fichier zip tout ce qu'il faut pour comprendre le problème ; il suffit d'envoyer ce zip aux développeurs. Il se crée de deux façons :
\begin{itemize}
  \item dans les options (roue dentée \bouton{⚙} en bas de la fenêtre principale, §\ref{btn:home_workdir}), bouton \bouton{Create a diagnostic report…} : on choisit une mission, ou l'installation seule ;
  \item sans ouvrir le logiciel, avec \texttt{irdrone\_diagnostic.bat} (à la racine du dossier IRDrone) : double-clic pour l'installation seule, ou glisser-déposer le dossier d'une mission (\texttt{FLY-...}) sur ce fichier.
\end{itemize}
Le zip (\texttt{irdrone\_diagnostic\_AAAAMMJJ\_hhmmss.zip}) est écrit dans le dossier de travail. Il contient :
\begin{itemize}
  \item \texttt{report.txt} : un résumé de ce qui semble anormal, puis l'installation (version d'IRDrone, Python utilisé, système, page de code Windows, versions des modules Python dont OpenCV, chemins et versions d'ExifTool, RawTherapee et du convertisseur SJCam, dossier de travail, place libre) et deux auto-tests (détection ArUco d'une cible parfaite ; lecture et écriture d'images et d'EXIF dans un dossier accentué) ;
  \item pour une mission : les étapes faites, les fenêtres de l'étape~2, les lignes de temps, le décalage et le nombre de paires de l'étape~3 ; chaque image de synchronisation est développée à nouveau avec RawTherapee dans un dossier temporaire et les marqueurs ArUco y sont cherchés, ainsi que sur l'image en cache (\texttt{Synchro}) utilisée par l'étape~3. Quelques images développées sont jointes. Cela prend environ une minute ; la mission n'est pas modifiée ;
  \item \texttt{report.json} (tous les détails) et les journaux des dernières sessions : chaque lancement du logiciel recopie la console dans \texttt{logs\textbackslash irdrone\_AAAAMMJJ\_hhmmss.log}, à côté du fichier des réglages (les dix derniers sont gardés).
\end{itemize}
Rien n'est envoyé automatiquement.

%==========================================================
\chapter{Step 1 -- D\'efinir la mission}\label{ch:step1}
\section{La fenêtre \og New mission \fg}\label{sec:step1_new_mission}
Le Step~1 cr\'ee le dossier de la mission et y range les images du vol. Il s'ouvre avec \bouton{New mission} dans la fenêtre principale et se fait dans une seule fenêtre, en trois pages : \textbf{Images}, \textbf{Mission}, \textbf{Copy}. L'avancement est rappelé en haut de la fenêtre. Les cartes SD sont seulement lues : rien n'y est écrit ni effacé.

\subsection{Type de mission : multispectrale ou visible seul}\label{btn:step1_mission_type}
Le haut de la page \og Images \fg~choisit le type de la mission (le dernier choisi est proposé) :
\begin{itemize}
	\item \bouton{Multispectral (visible + infrared)} : le drone (visible, DNG) et la caméra infrarouge (SJCam, RAW). Tout ce que décrit ce manuel : synchronisation, alignement, carte en quatre bandes (rouge, vert, bleu, proche infrarouge), VIR et NDVI ;
	\item \bouton{Visible only (drone DNG)} : les images du drone seules, sans la caméra infrarouge. La carte infrarouge disparaît, la copie ne prend que les images visibles et les étapes suivantes s'adaptent (section~\ref{sec:visible_only}) : phase de relevé seule à l'étape~2, pas d'étape~3, les images sont \emph{développées} à l'étape~4 au lieu d'être alignées, la carte de l'étape~6 est en couleurs (rouge, vert, bleu).
\end{itemize}
Le type ne se change plus une fois la mission créée.

\subsection{Page \og Images \fg}\label{btn:step1_images}
Les cartes SD insérées sont parcourues (tout volume qui contient un dossier \texttt{DCIM}) ; on y trouve les images visibles (DNG du drone) et infrarouges (RAW de la SJCam), figure~\ref{fig:step1_new_mission_images}. Chaque caméra a sa carte :
\begin{itemize}
	\item le dossier trouvé, le nombre d'images, leur taille, la date et l'heure de la première et de la dernière ;
	\item deux images : \textbf{au décollage} et \textbf{en vol} (aux deux tiers des images, en général pendant le relevé). Pour l'infrarouge, c'est le JPEG que la SJCam enregistre à côté de chaque RAW. Ces images permettent de vérifier d'un coup d'\oe il que la carte est la bonne ;
	\item le bouton \bouton{Choose folder…} (\bouton{Change…} une fois des images trouvées) choisit un autre dossier, par exemple une copie des cartes sur le disque. Un dossier choisi ainsi n'est jamais remplacé par une carte trouvée ensuite.
\end{itemize}
Tant qu'il manque les images d'une caméra, la fenêtre cherche les cartes SD toutes les 3~secondes : une carte insérée après l'ouverture de la fenêtre est trouvée d'elle-même, il suffit de l'insérer.

L'image de décollage est la première image visible ; le lien \og change… \fg~de la carte visible en choisit une autre. Elle donne la date, l'heure, la position et l'appareil photo de la mission.
Un bandeau signale un écart de plus de deux minutes entre les horloges des deux caméras : il est normal, l'étape~3 le mesure précisément. Un écart de plusieurs jours fait douter que les deux cartes viennent du même vol.

\figureIRD[fig:step1_new_mission_images]{step1_new_mission_images.png}{Page \og Images \fg : images visibles et infrarouges trouvées sur les cartes SD, au décollage et en vol.}{0.8}

\subsection{Page \og Mission \fg}\label{btn:step1_mission}
La date, l'heure et la position viennent de l'image de décollage ; l'altitude du sol et l'adresse sont demandées à l'IGN\footnote{\label{fn:ign}IGN : Institut national de l'information g\'eographique et foresti\`ere.} et à OpenStreetMap pendant que l'on remplit la description et le pilote (figure~\ref{fig:step1_new_mission_details}). Le nom du dossier de la mission (\texttt{FLY-YYYYMMDD-hhmm-location}) se met à jour avec le champ \og Location \fg.

\bouton{Cameras and filter} déplie les réglages des caméras : intervalle de prise de vue et décalage d'horloge de chaque caméra, fabricant et modèle de la caméra infrarouge, fabricant et bande du filtre. Le pilote, son identifiant, les caméras et le filtre sont mémorisés pour la mission suivante.

\figureIRD[fig:step1_new_mission_details]{step1_new_mission_details.png}{Page \og Mission \fg : lieu trouvé à partir de l'image de décollage, nom du dossier, description, pilote, caméras.}{0.8}

\subsection{Page \og Copy \fg}\label{btn:step1_copy}
Le dossier est créé, les images visibles copiées, les images infrarouges converties de RAW en DNG, les EXIF lus et l'altitude du sol calculée (détails au §\ref{btn:step1_validate}). Pendant la conversion des images infrarouges, l'étape~2 est préparée en arrière-plan (vignettes et détection de la mire, visible puis infrarouge) : la ligne de temps s'ouvre ensuite sans attendre. La fenêtre ne peut pas être fermée pendant la copie. À la fin, \bouton{Flight timeline} ouvre l'étape~2 ; si la préparation n'est pas finie, l'étape~2 la reprend là où elle en était.

Le fichier \texttt{mission\_parameters.json} contient les mêmes informations qu'avec le formulaire classique.

\subsection{Boutons du bas de la fenêtre}\label{btn:step1_wizard}
\begin{description}
	\item[\bouton{Next}] (page Images) passe à la page Mission ; il faut les images des deux caméras. \bouton{Create the mission and copy the images} (page Mission) lance la copie ; \bouton{Flight timeline} (page Copy) ouvre l'étape~2.
	\item[\bouton{Back}] revient à la page Images ; rien n'est perdu.
	\item[\bouton{Cancel}] ferme la fenêtre sans créer la mission.
	\item[\bouton{Classic form…}] (page Images) ouvre l'ancien Step~1 (§\ref{sec:step1_classic}).
	\item[\bouton{?}] ouvre ce chapitre.
\end{description}

\section{Modifier une mission : la fenêtre \og Mission details \fg}\label{btn:step1_details}
Une fois l'étape~1 faite, un clic sur l'étape~1 de la fenêtre principale n'ouvre plus la création de mission, mais la fenêtre \og Mission details \fg~(figure~\ref{fig:step1_mission_details}). Elle modifie les informations de la mission sans relire les cartes SD : lieu, description, pilote et son identifiant, intervalle et décalage d'horloge des caméras, caméra infrarouge, filtre. La date, l'heure, la position et la caméra visible viennent de l'image de décollage et ne changent pas.

Dès qu'une information change, \bouton{Close} devient \bouton{Save and close} : il écrit les changements dans \texttt{FlightAnalytics/mission\_parameters.json} (fermer la fenêtre fait de même) ; l'ancien fichier est gardé en \texttt{mission\_parameters.json.bak}. Le dossier de la mission garde son nom, même si le lieu change. \bouton{↺}, à côté de \bouton{?}, apparaît alors : il remet les informations telles qu'à l'ouverture de la fenêtre ; la fenêtre reste ouverte.

\figureIRD[fig:step1_mission_details]{step1_mission_details.png}{La fenêtre \og Mission details \fg : les images sont déjà dans la mission, seules ces informations changent.}{0.55}

\section{Formulaire classique}\label{sec:step1_classic}
\legacybegin{Afficher le formulaire classique (ancienne étape 1)}
Le bouton \bouton{Classic form…} de la page \og Images \fg~ouvre l'ancien Step~1, décrit ci-dessous. Il se d\'eroule en deux fen\^etres successives :
\begin{enumerate}
  \item \textbf{\og Choose the image taken by the drone during takeoff \fg} : on choisit la premi\`ere image prise par le drone
        au d\'ecollage. Elle fournit la date, l'heure, la position GPS et l'appareil photo utilis\'e.
  \item \textbf{\og Create a mission \fg} : on v\'erifie et on compl\`ete les informations de la mission, puis on lance la cr\'eation
        du dossier et le transfert des images VIS et NIR.
\end{enumerate}
\`A la fin, la mission devient la mission chargée de la fen\^etre principale, qui propose l'étape~2.

\paragraph{Pr\'erequis.}
\begin{itemize}
  \item Une connexion Internet : l'altitude du sol est demand\'ee \`a l'IGN\textsuperscript{\ref{fn:ign}}
  (France) ou \`a OpenTopoData (hors de France),
        et l'adresse du lieu \`a OpenStreetMap.
  \item Les images du drone (visible) au format DNG, et les images de la cam\'era proche infrarouge (SJCam M20, fichiers RAW)
        copi\'ees sur le disque de l'ordinateur.
  \item Un dossier de travail (par d\'efaut \texttt{C:/Air-Mission}), choisi au premier lancement et modifiable
        avec le bouton \og engrenage \fg{} en bas de la fen\^etre principale. La variable d'environnement
        \texttt{IRDRONE\_WORKING\_DIR}, si elle est d\'efinie, reste prioritaire.
\end{itemize}

% \figureIRD{step1_load_takeoff_1.png}{Premi\`ere fen\^etre du Step 1 avant chargement de l'image de d\'ecollage.}{0.8}

%----------------------------------------------------------
\section{Bouton \bouton{Load take-off image}}\label{btn:step1_load_takeoff}
Ouvre un explorateur de fichiers pour choisir l'image prise par le drone au d\'ecollage.
Seuls les fichiers \texttt{.dng} de la cam\'era visible sont accept\'es.

Une fois l'image choisie, le logiciel :
\figureIRD{step1_load_takeoff_2.png}{Premi\`ere fen\^etre du Step 1 apr\`es chargement de l'image de d\'ecollage.}{0.5}
\begin{enumerate}
  \item affiche l'image (le d\'ecodage du fichier DNG est assez long : la barre de progression indique l'avancement) ;
  \item lit les donn\'ees EXIF : date et heure de prise de vue, position GPS, fabricant, mod\`ele et num\'ero de s\'erie de l'appareil ;
  \item interroge l'IGN (ou OpenTopoData hors de France) pour obtenir l'altitude du sol \emph{au-dessus du niveau de la mer} ;
  \item interroge OpenStreetMap pour obtenir le lieu-dit, la commune, le code postal, le d\'epartement, la r\'egion et le pays.
\end{enumerate}
Ces informations s'affichent dans le cadre gris sous l'image. Lorsqu'elles existent, le lieu-dit, le code postal, la région et le pays sont aussi enregistrés dans le fichier
\texttt{FlightAnalytics\textbackslash mission\_parameters.json} sous les clés \texttt{Place name}, \texttt{Postal code}, \texttt{Region} et \texttt{Country}
(la commune est enregistrée sous la clé \texttt{Location}). Une information introuvable n'est pas enregistrée : la clé correspondante n'existe pas.
\begin{center}
\fbox{\parbox{0.9\linewidth}{\small\textbf{Attention.} L'altitude enregistr\'ee dans l'image par le drone est mesur\'ee \emph{par rapport
au point de d\'ecollage}, et non par rapport \`a la mer. C'est pourquoi l'altitude du sol est demand\'ee \`a un service externe tel que l'IGN\textsuperscript{\ref{fn:ign}}.}}
\end{center}


\paragraph{Si l'image ne se charge pas.} V\'erifiez la connexion Internet et que le fichier est bien un \texttt{.dng} du drone.
Le bouton \bouton{Next ▶} reste inactif tant qu'aucune image n'a \'et\'e charg\'ee correctement.

%----------------------------------------------------------
\section{Boutons \bouton{Cancel}, \bouton{Next ▶} et \bouton{◀ Back}}\label{btn:step1_navigation}
\begin{description}
  \item[\bouton{Next ▶}] n'est disponible qu'apr\`es le chargement de l'image de d\'ecollage. Il ouvre la fen\^etre \og Create a mission \fg.
  \item[\bouton{Cancel}] ferme la fen\^etre de l'image de d\'ecollage sans rien enregistrer.
  \item[\bouton{◀ Back}] (fen\^etre \og Create a mission \fg) revient \`a la fen\^etre pr\'ec\'edente ;
        l'image de d\'ecollage d\'ej\`a charg\'ee est conserv\'ee.
\end{description}

%----------------------------------------------------------
\section{La fen\^etre \og Create a mission \fg}\label{btn:step1_form}
Les champs sont remplis automatiquement \`a partir de l'image de d\'ecollage. La touche \textsf{Entr\'ee} passe d'un champ au suivant.
Les champs gris sont remplis automatiquement et ne sont pas modifiables.

\begin{center}
\begin{tabular}{@{}p{3.6cm}p{2.4cm}p{7.5cm}@{}}
\hline
\textbf{Champ} & \textbf{Modifiable} & \textbf{Contenu} \\ \hline
Date, Hour & non & Date et heure du d\'ecollage (EXIF). \\
Location & oui & Commune trouv\'ee automatiquement. Fait partie du nom du dossier de la mission. \\
GPS & non & Latitude et longitude du point de d\'ecollage. \\
Short description & oui & Br\`eve description libre de la mission. \\
Pilot & oui & Nom du t\'el\'epilote et son num\'ero d'identification. \\
Camera VIS & en partie & Fabricant et mod\`ele (lus dans l'image, non modifiables) ; intervalle de prise de vue (timelapse, en s) et d\'ecalage d'horloge (delta time, en s). \\
Image VIS format & non & Format des images (DNG). \\
Camera NIR & oui & Par d\'efaut SJCam M20, format RAW, intervalle 3\,s, d\'ecalage d'horloge $-30$\,s. \\
NIR Filter & oui & Fabricant du filtre et bande, en nm (par d\'efaut 810\,nm). \\ \hline
\end{tabular}
\end{center}

\figureIRD{step1_param_mission.png}{Fen\^etre \og Create a mission \fg.}{0.5}

Remarque: les paramètres de décalage d'horloge VIS et NIR sont réservés à un utilisateur expérimenté. Ce sont des paramètres qui permettent de \og récupérer\fg ~une mission dont la phase de synchronisation  sur le terrain se révèle inexploitable (impossibilité de reconnaître les ArUco, oubli de déclenchement de la mire, etc.).

%----------------------------------------------------------
\section{Bouton \bouton{Create the mission}}\label{btn:step1_validate}
Lance la cr\'eation de la mission. Le bouton se d\'esactive pendant le traitement ; une \'etiquette et une barre de progression
indiquent la phase en cours. Le logiciel encha\^ine, sans intervention de l'utilisateur sauf \`a l'\'etape~3 :
\begin{enumerate}
  \item \textbf{Cr\'eation du dossier} de mission et de ses sous-dossiers, \'ecriture de \texttt{mission\_parameters.json},
        copie de l'image de d\'ecollage.
  \item \textbf{Transfert des images VIS} : les fichiers \texttt{HYPERLAPSE\_XXXX.dng} du dossier de l'image de d\'ecollage sont copi\'es
        dans \texttt{AerialPhotography/VIS} sous le nom \texttt{VIS\_XXXX.dng}.
  \item \textbf{Choix du dossier NIR} : un message demande de choisir la premi\`ere image proche infrarouge (fichier RAW ou JPG
  \footnote{Chaque image RAW a une image jpeg associée. La SJCam nomme les images par leur date de prise de vue. Par exemple 2022\_0125\_130413\_001.RAW, 2022\_0125\_130415\_002.JPG, 2022\_0125\_130416\_003.RAW, 2022\_0125\_130418\_004.JPG, etc. L'image jpeg est toujours enregistrée après l'image RAW environ 2s après. À noter que dans certaines circonstances (carte SD trop lente, longue séquence de vol) le délai d'enregistrement du fichier RAW dépasse la fréquence de prise de vue (3s dans notre cas). Dans ce cas l'image suivante est prise 6s après. IRdrone gère automatiquement ce type de saut.} de la SJCam).
        Le dossier qui la contient est utilis\'e pour toutes les images NIR.
  \item \textbf{Conversion des images NIR} : les fichiers RAW sont convertis en DNG dans \texttt{AerialPhotography/NIR}
        (calcul en parall\`ele, \'eventuellement long).
  \item \textbf{Fichiers \texttt{.exif}} des images VIS puis NIR : un fichier de m\'etadonn\'ees par image,
        et un fichier \texttt{time\_line.json} par bande donnant la chronologie des prises de vue.
  \item \textbf{Altitudes et trajectoire} : pour chaque position GPS du vol, l'altitude du sol est demand\'ee \`a l'IGN (par paquets ;
        trois essais en cas d'\'echec, puis bascule automatique sur OpenTopoData). Le logiciel en d\'eduit la hauteur du drone au-dessus du sol,
        les distances et les caps entre prises de vue, et enregistre ces valeurs dans les fichiers \texttt{.exif}.
  \item \textbf{Profil vertical} : trac\'e de l'altitude du drone et du sol le long du trajet, enregistr\'e dans
        \texttt{FlightAnalytics/fly\_vertical\_profil.png}.
  \item \textbf{Message final} indiquant le dossier cr\'e\'e, puis fermeture des fen\^etres du Step~1.
\end{enumerate}

\figureIRD[fig:takeoff_1]{step1_nir_prompt.png}{Message demandant de choisir la premi\`ere image NIR.}{0.2}
\figureIRD[fig:fly_vertical_profil]{step1_fly_vertical_profil.png}{Profil vertical du vol (altitude du drone et du sol), produit \`a la fin du Step~1.}{0.8}

\subsection{Dossiers cr\'e\'es}
À l'issue de l'étape 1 (Step 1) IRDrone aura cr\'e\'e la structure de dossiers et fichiers de la mission. À ce stade vous n'aurez plus besoin d'accéder aux images originales stockées sur les cartes SD des caméras. La structure du dossier de la mission est la suivante :


\begin{verbatim}
FLY-YYYYMMDD-hhmm-<location>/
  AerialPhotography/
     VIS/   VIS_XXXX.dng, VIS_XXXX.exif, time_line.json,
            transfer_info_VIS_dng.json
     NIR/   NIR_XXXX.dng, NIR_XXXX.exif, time_line.json,
            transfer_info_NIR_dng.json
  FlightAnalytics/
     mission_parameters.json, fly_vertical_profil.png
  Synchro/        (Step 3)
  ImgIRdrone/     (Step 4)
  mapping_MULTI/  (Step 6)
  cameras/
\end{verbatim}
Les sous-dossiers ont une ic\^one sp\'ecifique dans l'explorateur Windows.
Le dossier de mission n'est jamais effac\'e : si le dossier existe d\'ej\`a, seuls les \'el\'ements manquants sont cr\'e\'es.
Pour les images VIS d\'ej\`a transf\'er\'ees (fichier \texttt{.exif} pr\'esent), le transfert est ignor\'e : relancer le Step~1 sur
une mission existante est donc rapide.

\figureIRD{Step1_mission_folder.png}{Dossier de mission vu dans l'explorateur Windows.}{0.7}

La durée des transferts dépend beaucoup des performances des supports (carte SD, disque dur) et des liaisons (USB B ou C, lecteur). De plus et contrairement aux images de la caméra DJI qui sont au format dng, les images de la caméra SJCam M20 sont dans un format \og RAW \fg~qui n'est pas normalisé. Chaque image doit être convertie au format dng d'ADOBE à l'aide d'un logiciel tiers \footnote{Il s'agit du code \texttt{sjcam\_raw2dng} disponible uniquement comme un fichier exe. \texttt{sjcam\_raw2dng}
	Convert SJCAM raw files to Adobe DNG and TIFF while also getting metadata
	such as date and ISO from JPEGs that correlate to the RAW files
} ce qui allonge le temps de transfert effectif. 

\subsection{En cas de probl\`eme}
\begin{itemize}
  \item \textbf{Aucune image NIR trouv\'ee} : le message \og No NIR RAW files found \fg\ signifie que le dossier choisi ne contient pas de
        fichiers \texttt{.RAW} ; recommencez en choisissant le bon dossier.
  \item \textbf{Erreur pendant le transfert VIS ou la lecture EXIF} : un message d'erreur s'affiche et la fen\^etre se ferme.
        Le bouton reste utilisable apr\`es r\'eouverture de la fen\^etre.
  \item \textbf{Altitudes IGN indisponibles} : le logiciel bascule seul sur OpenTopoData, moins pr\'ecis en France.
\end{itemize}
%\todo{Compl\'eter avec les erreurs rencontr\'ees sur le terrain (carte SD absente, dossier prot\'eg\'e en \'ecriture, \dots).}

%==========================================================
\legacyend

\chapter{Step 2 -- Ligne de temps du vol}\label{ch:step2}
\section{Principe}
L'étape 2 découpe le vol en phases : \emph{synchronisation} (vol stationnaire au-dessus de la mire ArUco tournante) et \emph{survey} (vol au-dessus de la zone étudiée). Ces phases sont maintenant \textbf{détectées automatiquement} ; l'utilisateur vérifie et ajuste si besoin (figure~\ref{fig:step2_flight_timeline}).

\figureIRD[fig:step2_flight_timeline]{step2_flight_timeline.png}{La fenêtre \og Flight timeline \fg. Au centre l'aperçu des images visible et infrarouge sous la souris, à droite la carte du vol avec les empreintes au sol des images gardées pour la carte (§\ref{btn:step2_best_mapping}), puis les quatre points clés avec leurs vraies images ; en bas la ligne de temps : profil d'altitude (points verts : images gardées pour la carte), phases (violet : synchronisation, vert : survey), détections de la mire et bandes d'images visible et infrarouge.}{0.95}

\subsection*{Ce que fait le logiciel}
\begin{enumerate}
	\item Chaque fichier DNG (DJI comme SJCam) contient un aperçu JPEG d'environ 1000 pixels. Il se lit en quelques millisecondes : toute la mission est parcourue en quelques secondes, sans RawTherapee. Les aperçus sont gardés dans \texttt{Synchro/Thumbnails} (cache, recréé si besoin).
	\item \emph{Visible} : l'altitude enregistrée par le drone donne le vol stationnaire du début (synchronisation) et le palier le plus long, le plus éloigné de l'altitude du décollage (survey). Un décollage depuis un sommet, suivi d'une descente vers la zone étudiée, est traité de la même façon qu'une montée.
	\item \emph{Synchronisation} : l'angle de la mire mobile (identifiant 0) est mesuré sur les aperçus, seulement pendant le stationnaire pour le visible, et pour l'infrarouge depuis la première image jusqu'à ce que la mire disparaisse (le drone s'est éloigné). Les images où la mire tourne régulièrement dans un sens forment la fenêtre de synchronisation ; les images où le plateau est encore arrêté, où il change de sens, et le décollage sont écartés.
	\item \emph{Infrarouge} : l'horloge de la caméra NIR est souvent fausse (fuseau horaire, dérive). La même rotation de la mire étant vue par les deux caméras, la comparaison des deux courbes d'angle donne un décalage approché (de l'ordre de la seconde) qui place les images infrarouges sur la ligne de temps du visible. La fenêtre de synchronisation infrarouge correspond ainsi au même tour de mire que la fenêtre visible.
	\item La fenêtre survey infrarouge \textbf{suit} la fenêtre survey visible (zone hachurée) : il n'est pas nécessaire de regarder les images infrarouges. Après l'étape 3, le décalage précis met à jour cette fenêtre.
\end{enumerate}

\subsection*{Vérifier et ajuster}\label{btn:step2_adjust}
\begin{itemize}
	\item Survoler la ligne de temps affiche en grand l'image visible et l'image infrarouge prises au même instant, avec l'altitude et l'état de la mire (\og target turning \fg, \og target still \fg, \og no target \fg). Laisser la souris sur cet aperçu l'affiche sur toute la fenêtre pour bien voir les détails.
	\item Les points de la ligne \og target \fg~indiquent les images analysées : violet, la mire tourne ; gris, la mire est vue mais immobile ; cercle vide, mire non vue.
	\item Faire glisser un bord de phase pour le déplacer (il se cale sur une image), ou cliquer sur un point clé (colonne de droite) puis utiliser \bouton{◀} \bouton{▶} ou les flèches du clavier.
	\item Molette : zoom ; glisser le fond : déplacement ; double-clic : tout le vol.
	\item \bouton{↺}, à côté de \bouton{?}, apparaît dès qu'une phase est déplacée : il annule les changements faits à la main et revient à la proposition automatique ; la fenêtre reste ouverte.
	\item L'info-bulle du bouton \bouton{?} rappelle ces commandes.
\end{itemize}
Si la mire n'est pas trouvée (mire trop petite, image trop sombre), un bandeau orange l'indique : la fenêtre de synchronisation visible est alors le vol stationnaire, et la fenêtre infrarouge doit être placée à la main.

\section{Best mapping : les images utiles à la carte}\label{btn:step2_best_mapping}
Un survol photographié toutes les deux secondes donne beaucoup plus de recouvrement que la carte n'en demande. Aligner toutes ces images à l'étape~4 prend des heures, et donner à ODM des centaines d'images presque identiques ralentit le calcul de la carte sans l'améliorer, souvent au contraire. La case \bouton{Best mapping: keep only the images the map needs}, cochée par défaut, ne garde que les images du survol dont les \emph{empreintes au sol} se recouvrent de la valeur choisie à droite (\bouton{75~\% overlap} par défaut, la valeur recommandée pour ODM ; plus de recouvrement garde plus d'images).

Le choix est visible tout de suite :
\begin{itemize}
	\item sur la ligne de temps, des points verts sur le profil d'altitude et des repères verts en haut de la bande d'images visibles marquent les images gardées ; la bulle du curseur et la légende de l'aperçu indiquent \og kept for the map \fg ;
	\item à droite des aperçus, la \textbf{carte du vol} montre le trajet (violet : synchronisation, vert : survey), les empreintes au sol des images gardées (rectangles verts : leur recouvrement se voit) et, en orange, l'empreinte de l'image sous la souris avec l'image elle-même. Un clic sur un point de la carte affiche cette image ; molette : zoom, glisser : déplacement, \bouton{Layers} : fond de carte (§\ref{btn:home_map}) ;
	\item sous la ligne de temps, le nombre d'images gardées, par exemple \og 16 of 120 survey images kept for the map · 87~\% less to align \fg.
\end{itemize}
La sélection suit les bords de la phase survey quand on les déplace.

\textbf{Comment les images sont choisies.} La position, la hauteur au-dessus du sol et l'orientation de chaque image sont lues dans les fichiers \texttt{.exif} écrits à l'étape~1 ; le champ de vue de la caméra donne la taille de l'empreinte au sol. La première image du survol est gardée. Pour chaque image suivante, IRDrone estime le recouvrement de son empreinte avec celle de la dernière image gardée, d'après le déplacement du drone (dans l'axe de l'image et en travers). Tant qu'il reste au-dessus de la valeur choisie, l'image n'est qu'une candidate ; dès qu'il passe en dessous, la candidate précédente est gardée : deux images gardées successives se recouvrent toujours d'au moins la valeur choisie. La dernière image du survol est gardée aussi. Pour un déplacement dans l'axe de l'image, c'est le critère \texttt{--selection best-mapping} de l'ancienne ligne de commande \texttt{run.py}.

\textbf{Et la synchronisation ?} À l'étape~2 le décalage précis entre les deux caméras n'est pas encore connu : l'aperçu porte sur toutes les images du survol. À l'étape~4, une fois les caméras synchronisées, la même sélection est refaite parmi les seules paires bien synchronisées (§\ref{btn:step4_best_mapping}) ; le nombre final peut donc différer un peu de l'aperçu. Le choix (case et recouvrement) est enregistré par \bouton{Continue with synchronisation} et \bouton{Close} dans \texttt{FlightAnalytics\textbackslash best\_mapping.json}.

Si les images n'ont pas de position GPS, la carte n'est pas affichée et toutes les images du survol sont gardées.

\section{Boutons \bouton{Continue with synchronisation} et \bouton{Close}}\label{btn:step2_confirm}
Les deux boutons enregistrent les phases dans les fichiers \texttt{transfer\_info\_VIS\_dng.json} et \texttt{transfer\_info\_NIR\_dng.json} (lus par les étapes 3, 4 et le mapping) et dans \texttt{FlightAnalytics/flight\_timeline.json}. Rouvrir l'étape 2 retrouve les phases confirmées.
\bouton{Continue with synchronisation} ouvre ensuite l'étape~3 ; rien ne démarre tout seul : la synchronisation se lance depuis sa fenêtre (chapitre~\ref{ch:step3}). \bouton{Close} ferme simplement la fenêtre : l'étape~3 se lancera plus tard depuis la fenêtre principale. Il devient \bouton{Save and close} quand les phases ou le choix best mapping ont changé, ou n'ont jamais été enregistrés : fermer la fenêtre les enregistre alors aussi. S'il manque une phase, fermer demande s'il faut fermer sans enregistrer.

\section{Sélection classique}\label{btn:step2_classic}
\legacybegin{Afficher la sélection classique (ancienne fenêtre de l'étape 2)}
Le bouton \bouton{Classic selection…} ouvre l'ancienne sélection manuelle des cinq images de référence de chaque caméra, décrite ci-dessous.

\figureIRD{step2_dialog_select_phase_images_1.png}{Fen\^etre de s\'election des images de r\'ef\'erence. Pour la bande spectrale visible la première image est par défaut celle du point de décollage qui a servi à définir la mission à l'étape 1. Mais elle peut être changée ici.}{0.95}
La fenêtre de dialogue permet la définition :
\begin{enumerate}
	\item du point de décollage,
	\item des points de début et fin de synchronisation,
	\item des points de début et fin de vol au-dessus de la zone étudiée par la mission.
\end{enumerate}
Les images sont sélectionnées pour les spectres visible et proche infrarouge  respectivement dans les dossiers
\begin{itemize} 
	\item \texttt{C:/Air-Mission/FLY-YYYYMMDD-hhmm-location/AerialPhotography/VIS}, 
   	\item  \texttt{C:/Air-Mission/FLY-YYYYMMDD-hhmm-location/AerialPhotography/NIR}.
\end{itemize} 
 Lors de cette étape de choix il est toujours possible après une première sélection de choisir une autre image. 
 L'ordre de sélection des cinq images n'a pas d'importance mais bien entendu first précède toujours last.
 Il est possible de sélectionner la même image pour \texttt{first mission image} et \texttt{first sync image}.  
  Pour les images DJI du spectre visible, si vous sélectionnez la même image pour \texttt{first sync image} et \texttt{last sync image} alors une seule image sera retenue (c'est un cas limite pour la synchronisation). 
  En revanche  si vous sélectionnez la même image pour \texttt{first sync image} et \texttt{last sync image} pour le spectre proche infrarouge la synchronisation automatique échouera à l'étape 3 et il vous sera demandé d'entrer manuellement l'offset de synchronisation des horloges. 
  Cela peut constituer une \og astuce \fg~ pour forcer l' offset de synchronisation des horloges à une valeur particulière en cas de synchronisation terrain problématique.
%\todo{\bouton{Load first mission image.}, \bouton{Load first sync image.}, \bouton{Load last sync image.},
%----------------------------------------------------------
\section{Bouton \bouton{Load first mission image.}}\label{btn:step2_Load_first_mission_image}
Le choix de cette image définit l'origine de la ligne de temps pour la synchronisation des horloges des caméras. En règle générale on choisira la première prise de vue VIS de la mission.
\section{Boutons \bouton{Load first sync image.}, \bouton{Load last sync image.}}\label{btn:step2_Load_synchro_image}
L'utilisateur doit choisir l'image de début du cycle de synchronisation. Pour les images visibles, si la procédure a été respectée sur le terrain, il s'agit de la première ou deuxième image enregistrée. Pour les images infrarouges il faut repérer la première image stable où le système de mire est visible en totalité et où la mire mobile est effectivement en rotation. Ensuite il faut choisir la dernière image de la phase de synchronisation. En général on la repère facilement car lorsque la phase de synchronisation est terminée le pilote prend rapidement de l'altitude et le changement d'aspect est bien visible. Les choix faits ici conditionnent en partie la réussite de l'étape 3 de synchronisation des horloges. Mais après l'étape 3 il est possible de revenir à l'étape 2 et de redéfinir ces images.
\figureIRD[fig:step2_dialog_select_VIS]{step2_dialog_select_phase_image_VIS.png}{Fen\^etre de s\'election des images de r\'ef\'erence. Bande spectrale visible (VIS).}{0.95}
\section{Boutons \bouton{Load first fly image.}, \bouton{Load last fly image.}}\label{btn:step2_Load_mission_image}
L'utilisateur doit choisir la première image prise au-dessus de la zone à étudier (vol mission ou hyperlapse survey). On élimine ainsi les images de l'approche entre le décollage/synchronisation à très basse altitude et le vol de la mission à \og haute altitude \fg~ (hyperlapse survey). À noter qu'en principe le vol mission devrait s'opérer à altitude constante du drone et au-dessus d'une zone dont les altitudes varient peu par rapport à l'altitude du drone par rapport au sol. L'observation de la figure \ref{fig:fly_vertical_profil} permet de juger de la qualité du plan de vol réalisé sur le terrain.		
\section{Bouton \bouton{Validate mission images}}\label{btn:step2_validate}
\bouton{Cancel}, à sa gauche, ferme la fenêtre (après confirmation si les images ne sont pas validées). Ce bouton déclenche la construction des listes des images des trois phases de vol pertinentes pour IRDrone : takeoff, synchro et fly. Aucune image n'est déplacée. L'opération est très rapide.
\figureIRD[fig:step2_dialog_select_NIR]{step2_dialog_select_phase_images_NIR.png}{Fen\^etre de s\'election des images de r\'ef\'erence. Bande spectrale proche infrarouge (NIR).}{0.95}

%==========================================================
\legacyend

\chapter{Step 3 -- Synchronisation des caméras}\label{ch:step3}
\section{Principe}
Pour combiner différentes bandes spectrales, les images doivent être prises au même instant. Malheureusement ce n'est techniquement pas possible avec les matériels utilisés dans le projet IRdrone. C'est un choix de départ lié au caractère financièrement abordable (\emph{inexpensive and low cost}) du projet. Chaque caméra a sa propre horloge, qui définit la chronologie des images. On peut positionner les images sur une ligne représentant le temps. C'est la \og time line \fg~de la caméra. Dès lors il faut synchroniser les deux time lines VIS et NIR. Mettre les horloges à l'heure est insuffisant : d'une part elles ne donnent accès qu'à des points toutes les secondes, d'autre part on ne peut pas contrôler l'instant exact où l'opérateur sur le terrain a déclenché les caméras. On doit donc détecter des repères sur les images. Dans l'industrie cinématographique on synchronise la pellicule photo et la bande son à l'aide du \og clap \fg. Cette solution avait été envisagée dans les premières versions d'IRdrone, qui utilisaient la vidéo au lieu des images. Mais la vidéo posait de nombreux problèmes de qualité (résolution, \og jelly \fg, ...) et a été abandonnée. Par contre le principe d'un objet mobile sur les images au début du vol a été conservé. L'idée est de repérer un objet (automatiquement) et de suivre sa trajectoire sur les deux caméras. En faisant coïncider les deux trajectoires on synchronise les horloges. On peut atteindre une précision de l'ordre de 0.05s, ce qui est amplement suffisant dans notre cas.
IRdrone utilise comme objet mobile une mire ArUco 
\footnote{Ce type de mire est conçu pour être détecté automatiquement par logiciel, par exemple avec OpenCV. IRdrone utilise le dictionnaire \texttt{DICT\_4X4\_50} : la mire mobile porte l'identifiant 0 et les quatre mires fixes les identifiants 5, 22, 16 et 13.} 
fixée à un plateau tournant
\footnote{Le plateau est entraîné par un moteur pas à pas commandé par un Arduino. On programme le système pour obtenir une vitesse de rotation parfaitement constante. Elle est fixée à la valeur de 1 rpm (un tour par minute) de telle sorte que lors de la phase de vol de synchronisation les caméras enregistrent au moins une dizaine de positions angulaires différentes.} comme le montre la figure \ref{fig:step3_syncro_aruco_1}. 

\figureIRD[fig:step3_syncro_aruco_1]{step3_syncro_aruco_1.pdf}{Dispositif de synchronisation des horloges des caméras VIS et NIR. Le principe repose sur la détection de la position angulaire de la mire ArUco 0 qui tourne à vitesse constante. Les petites mires fixes permettent de définir un repère absolu pour s'affranchir des mouvements parasites du drone, particulièrement les rotations.}{0.95}

On mesure la position angulaire de la mire. Mais pour cela il faut un repère fixe. Les bords de l'image ont été utilisés initialement. Cependant ils ne sont fiables que si le drone ne tourne pas. Il faut donc une zone très calme et que le pilote ne touche pas à la manette de rotation. Pour libérer cette contrainte sur le terrain, on a ajouté au dispositif 4 mires ArUco fixes
\footnote{Théoriquement 2 mires suffisent pour définir un axe de référence. L'utilisation de 4 mires permet une certaine redondance. À noter que le logiciel tente de synchroniser dans le repère absolu, mais s'il échoue il synchronise dans le repère relatif lié aux bords des images (moins précis). Si là encore il y a échec, en dernier recours le logiciel demande à l'utilisateur de donner la valeur de l'offset de synchronisation.}. Elles permettent de définir un repère absolu qui améliore la précision du dispositif. Il est évident que cette phase de synchronisation doit être réalisée sur le terrain avec le plus grand soin. C'est la condition pour concilier le coût modéré du matériel et une bonne qualité des images multispectrales.

\subsection{Les trois phases de l'étape 3}
Le Step~3 se déroule en trois phases successives, enchaînées automatiquement (§\ref{sec:step3_chained}). En mode manuel, chacune est lancée par un bouton qui n'est activé qu'une fois la phase précédente terminée (figure~\ref{fig:step3_window}) :
\begin{enumerate}
  \item \emph{Détection des mires} (bouton \bouton{Run ArUco detection (VIS + NIR)}) : sur chaque image de synchronisation, on mesure l'angle de la mire mobile (§\ref{btn:step3_aruco}) ;
  \item \emph{Calcul du décalage temporel} (bouton \bouton{Compute time alignment}) : on cherche le décalage $\Delta t$ qui fait coïncider les deux séries d'angles (§\ref{btn:step3_time_shift}) ;
  \item \emph{Appariement} (bouton \bouton{Pair VIS / NIR images}) : on associe à chaque image VIS les images NIR les plus proches dans le temps (§\ref{btn:step3_pairing}).
\end{enumerate}

\subsection{Enchaînement automatique}\label{sec:step3_chained}
L'étape~3 ne démarre jamais seule : qu'elle soit ouverte depuis l'étape~2 (\bouton{Continue with synchronisation}) ou depuis la fenêtre principale, le bouton \bouton{Synchronise the cameras} (bleu foncé, au-dessus de la barre de progression) enchaîne les trois phases (§\ref{btn:step3_run}). Une ligne d'avancement montre la phase en cours et son résultat : nombre de mires vues, décalage, nombre de paires (figure~\ref{fig:step3_chained}). Une seule question est posée, après l'alignement des horloges (§\ref{btn:step3_overlap}).

Une fois les images appariées, rien ne démarre seul : \bouton{Continue with image processing} ouvre l'étape~4 quand vous êtes prêt (§\ref{btn:step3_continue}).
Si le décalage ne peut pas être mesuré, la saisie manuelle est proposée et l'enchaînement s'arrête.

\figureIRD[fig:step3_chained]{step3_chained.png}{Étape 3 enchaînée : mires détectées, horloges alignées ; la question \og Do the two curves overlap? \fg~attend la réponse avant l'appariement.}{0.7}

Rouverte une fois l'étape faite (clic sur l'étape~3 de la fenêtre principale), la fenêtre ne recalcule rien : elle montre le résultat enregistré, c'est-à-dire le décalage des horloges, le nombre de paires et les courbes alignées (§\ref{btn:step3_before}). \bouton{Synchronise again} la refait ; \bouton{Manual mode…} montre les boutons des trois phases.

\section{La question \og Do the two curves overlap? \fg}\label{btn:step3_overlap}
Après l'alignement des horloges, la fenêtre montre l'angle de la mire en fonction du temps : courbe verte pour le visible, rouge pour l'infrarouge, la courbe infrarouge étant décalée du décalage calculé.
\begin{description}
	\item[\bouton{Yes} (touche Y)] les deux courbes se superposent : le décalage est bon, les images sont appariées (§\ref{btn:step3_pairing}).
	\item[\bouton{No} (touche N)] ouvre le réglage à la main (§\ref{btn:step3_adjust}) pour glisser la courbe infrarouge sur la courbe visible ; la question est ensuite reposée avec le nouveau décalage.
\end{description}
Des courbes bien superposées ont des points verts sur la ligne rouge, dents de scie comprises. Les exemples du §\ref{sec:step3_examples} montrent un bon et un mauvais cas.

\section{Bouton \bouton{Synchronise the cameras}}\label{btn:step3_run}
L'action principale de l'étape~3 : elle lance la détection des mires, l'alignement des horloges puis, une fois répondu à la question sur les courbes, l'appariement. Elle devient \bouton{Synchronise again} quand l'étape a déjà été faite.

\section{Bouton \bouton{Manual mode…}}\label{btn:step3_manual}
Arrête l'enchaînement et fait apparaître les boutons des trois phases (§\ref{btn:step3_aruco} à \ref{btn:step3_pairing}), à lancer un par un, par exemple pour corriger le décalage avec \bouton{Adjust alignment...} avant l'appariement.

\section{Bouton \bouton{Enter the clock offset…}}\label{btn:step3_clock_offset}
Option avancée, pour un vol sans la cible ArUco (ancienne synchronisation par les dates). Le bouton n'apparaît que si la case \emph{Advanced – Step 3: let me enter the camera clock offset by hand} est cochée dans les options (roue dentée \bouton{⚙} en bas de la fenêtre principale, §\ref{btn:home_workdir}). Avec cette option, l'étape~3 ne démarre pas d'elle-même en venant de l'étape~2 : on choisit entre \bouton{Synchronise the cameras} (vol avec cible) et ce bouton.

Une petite fenêtre demande de combien de secondes l'horloge de la caméra infrarouge avance sur celle du drone (négatif si elle retarde, 0 si les deux horloges ont été réglées à la même heure). C'est le délai de l'ancienne synchronisation. Elle affiche le décalage $\Delta t$ de la ligne de temps infrarouge qui en découle, calculé à partir des dates des premières images. \bouton{Pair the images} saute la détection des cibles et lance directement l'appariement (§\ref{btn:step3_pairing}). Aucune courbe ne peut être vérifiée : la qualité de l'appariement dépend de la justesse du délai saisi.

\section{Bouton \bouton{Continue with image processing}}\label{btn:step3_continue}
Apparaît en bas à droite, à côté de \bouton{Close}, une fois les images appariées. Il ferme la fenêtre et ouvre l'étape~4. L'alignement, long, n'y démarre pas tout seul : il se lance avec \bouton{Run alignment}.

\section{Bouton \bouton{Before alignment}}\label{btn:step3_before}
Quand l'étape~3 est rouverte une fois faite, la fenêtre montre les courbes après alignement (\texttt{Synchro/check\_time\_alignment.png}). \bouton{Before alignment} montre les courbes avant correction (\texttt{check\_time\_before\_alignment.png}) ; un second clic revient aux courbes alignées.

\subsection{Principe du calcul}
L'angle de la mire mobile est mesuré sur chaque image, puis déplié (voir §\ref{btn:step3_time_shift}). Comme la vitesse de rotation $\omega$ est constante, l'angle est une fonction affine du temps. La time line NIR, qui compte au moins deux points, fournit la droite de référence $\theta = \omega t + b_{NIR}$. Chaque point VIS fournit une estimation de l'ordonnée à l'origine $b_{VIS} = \theta_{VIS} - \omega t_{VIS}$ ; on en prend la médiane. Le décalage s'écrit
\begin{equation}
  \Delta t = \frac{b_{NIR} - b_{VIS}}{\omega},
\end{equation}
et la time line NIR corrigée est $t_{NIR} + \Delta t$. C'est pourquoi une seule image VIS suffit en théorie (cas~2 ci-dessous), à condition que la vitesse de rotation soit constante et que l'on travaille dans le repère absolu.

\subsection{Exemples de synchronisation}\label{sec:step3_examples}
Les images de la phase de synchronisation viennent de l'étape~2. On présente quelques cas ici.

\subsubsection{Cas 1 : bonne synchronisation}
Les figures~\ref{fig:synchro_cas1} (\subref{fig:synchro_cas1_avant} et \subref{fig:synchro_cas1_apres}) représentent une bonne synchronisation. Le pilote a déclenché manuellement l'hyperlapse de la caméra NIR puis il a fait décoller le drone, ce qui a déclenché automatiquement la rotation de la mire. Lorsque le pilote a placé le drone au-dessus de la mire, la caméra NIR enregistrait la mire en rotation. Les images entre l'instant de déclenchement de l'hyperlapse NIR et l'instant où le drone a été stabilisé au-dessus de la mire ont été ignorées au step~2. Ensuite le pilote a réglé les paramètres de l'hyperlapse de la caméra VIS puis l'a déclenché. Cela prend en général une trentaine de secondes pour un pilote entraîné. Dès que l'hyperlapse VIS commence, il enregistre la mire en mouvement. À partir de cet instant et avec une fréquence d'enregistrement de 2s, il suffit au pilote de maintenir le drone moins d'une dizaine de secondes au-dessus de la mire pour obtenir 4 à 5 points, avant de prendre rapidement de l'altitude pour rejoindre la zone à étudier (hyperlapse survey). L'utilisateur du logiciel doit veiller lors de l'étape 2 à éliminer les images NIR prises avant la stabilisation du drone au-dessus du dispositif de synchronisation et le cas échéant celles qui après stabilisation correspondraient à une mire non encore en rotation. 

\figurePaireIRD{fig:synchro_cas1}
  {Step3_time_alignment_failed_cas1_0.png}{Trajectoire de la mire avant synchronisation.}
  {Step3_time_alignment_failed_cas1_1.png}{Synchronisation réussie.}
  {Cas de bonne synchronisation.}

\subsubsection{Cas 2 : cas limite}
Les figures~\ref{fig:synchro_cas2} (\subref{fig:synchro_cas2_avant} et \subref{fig:synchro_cas2_apres}) sont un cas limite pour la synchronisation. Le pilote est resté très peu de temps au-dessus de la mire et une seule image de cette mire est disponible sur les images VIS. Ici la synchronisation fonctionne tout de même car la vitesse de rotation est maintenue constante et le repère absolu a pu être utilisé.

\figurePaireIRD{fig:synchro_cas2}
  {Step3_time_alignment_failed_cas2_0.png}{Trajectoire de la mire avant synchronisation.}
  {Step3_time_alignment_failed_cas2_1.png}{Cas limite de synchronisation.}
  {Cas limite de synchronisation.}

\subsubsection{Cas 3 : échec de synchronisation}
Les figures~\ref{fig:synchro_cas3} (\subref{fig:synchro_cas3_avant} et \subref{fig:synchro_cas3_apres}) montrent la trajectoire de la mire détectée sur les deux caméras. Ici la synchronisation échoue car les premières images retenues pour la caméra NIR correspondent à des instants où la mire n'est pas en rotation.

\figurePaireIRD{fig:synchro_cas3}
  {Step3_time_alignment_failed_cas3_0.png}{Trajectoire de la mire avant synchronisation.}
  {Step3_time_alignment_failed_cas3_1.png}{Échec de synchronisation.}
  {Cas d'échec de synchronisation.}
   
\figureIRD[fig:step3_window]{step3_window.png}{Fen\^etre de synchronisation.}{0.8}

\section{Bouton \bouton{Run ArUco detection (VIS + NIR)}}\label{btn:step3_aruco}
Ce bouton lance la détection des mires ArUco sur toutes les images de la phase de synchronisation choisies à l'étape~2, d'abord pour la bande VIS puis, automatiquement, pour la bande NIR. Une barre de progression indique l'avancement. Les calculs sont répartis sur plusieurs c\oe urs du processeur pour accélérer le traitement.

\subsection*{Ce que fait le logiciel pour chaque image}
\begin{enumerate}
  \item Les images DNG sont converties en TIF par RawTherapee avec un profil de développement dédié à la détection (fichiers \texttt{*\_RawTherapee.tif} dans les dossiers \texttt{Synchro\textbackslash VIS} et \texttt{Synchro\textbackslash NIR} de la mission). Cette conversion n'est faite qu'une fois : si le fichier existe déjà, il est réutilisé.
  \item Les mires ArUco visibles sont détectées avec OpenCV. Pour chacune, on retient son identifiant, la position de son centre et l'angle de son axe principal dans l'image.
  \item La mire mobile (identifiant~0) donne l'\emph{angle relatif} (repère de l'image, \texttt{angle\_img}).
  \item Si les quatre mires fixes (identifiants 5, 22, 16 et 13) sont aussi détectées, elles définissent un repère lié au sol : deux axes obtenus en moyennant les côtés opposés du quadrilatère formé par leurs centres. L'angle de la mire mobile est alors mesuré dans ce repère. C'est l'\emph{angle absolu} (\texttt{angle\_abs}), insensible aux rotations du drone autour de son axe vertical.
  \item Si la mire mobile ou l'une des mires fixes manque, l'image est signalée dans le journal. L'angle absolu n'est pas calculé (mode dégradé : seul l'angle relatif est conservé si la mire mobile est visible) ; si la mire mobile n'est pas détectée, l'image est ignorée.
  \item À chaque image est associée sa date relative (en secondes) sur la time line de sa caméra.
\end{enumerate}
Les résultats de détection sont conservés dans les métadonn\'ees de chaque image.

\subsection*{Images de contrôle}
Pendant la détection, chaque image traitée est enregistrée avec une surcouche graphique dans le dossier \texttt{Synchro\textbackslash Check\_ARUCO}. La fenêtre en affiche une régulièrement (environ toutes les 5 secondes) pour que l'utilisateur puisse juger de la qualité de la détection. La surcouche montre :
\begin{itemize}
  \item les centres des mires fixes (disques rouges) ;
  \item le centre de la mire mobile ;
  \item l'orientation de la mire mobile dans l'image (trait vert) ;
  \item l'orientation projetée dans le repère absolu (trait orange en pointillés) et les deux axes du repère absolu (flèches noires) lorsque les quatre mires fixes sont détectées ;
  \item une légende : angle relatif, angle absolu et différence entre les deux, ou le message \og absolute reference: unavailable \fg~si le repère absolu n'a pas pu être construit.
\end{itemize}
\textbf{Conseil :} parcourir ce dossier dans l'explorateur Windows est le meilleur moyen de repérer des détections douteuses avant de passer à la phase suivante.

% Figure à créer par l'utilisateur (PowerPoint) : 2 lignes de 4 images (VIS en haut, NIR en bas) extraites de Synchro\Check_ARUCO.
\figureIRD[fig:check_aruco]{Step3_check_aruco_VIS_NIR.pdf}{Exemples d'images de contrôle (dossier \texttt{Synchro\textbackslash Check\_ARUCO}). Ligne du haut : bande VIS ; ligne du bas : bande NIR. La mire mobile est repérée par son axe vert ; les centres des mires fixes par des disques rouges.}{0.95}

\section{Bouton \bouton{Compute time alignment}}\label{btn:step3_time_shift}
Ce bouton, actif seulement lorsque la détection est terminée, calcule le décalage temporel entre les deux caméras.

\subsection*{Choix des angles}
Par défaut le logiciel utilise les angles absolus. Il faut au moins 2 images NIR et 1 image VIS où l'angle absolu est disponible. À défaut, il utilise les angles relatifs (repère de l'image), avec les mêmes minima, ce qui est moins précis si le drone a tourné pendant la synchronisation. Si aucune des deux solutions n'est possible, le message \og Insufficient data to automatically align VIS and NIR camera timelines \fg~s'affiche et le logiciel demande à l'utilisateur de saisir lui-même le décalage en secondes (valeur par défaut : 0). Dans ce cas aucun graphique n'est produit.

\subsection*{Calcul}
\begin{enumerate}
  \item Les angles de chaque bande sont \og dépliés \fg : on supprime les sauts de $\pm 360^\circ$ pour obtenir une courbe continue qui croît (ou décroît) régulièrement avec le temps. Le dépliage est fait de la même façon pour VIS et NIR, en partant de la dernière image, c'est-à-dire de l'image où la mire est en principe en rotation.
  \item On ajuste une droite sur les points NIR : sa pente est la vitesse de rotation $\omega$.
  \item On en déduit, avec les points VIS, le décalage $\Delta t$ (voir le principe en début de chapitre). L'estimation utilise la médiane des points VIS pour être robuste aux valeurs aberrantes.
\end{enumerate}
Le résultat est affiché dans le journal de la fenêtre :
\begin{itemize}
  \item le décalage estimé $\Delta t$ (en secondes) ;
  \item la vitesse de rotation mesurée sur NIR et sur VIS (en rpm). Elle doit être proche de la valeur nominale de 1~rpm et les deux valeurs doivent être voisines ; l'écart relatif est affiché ;
  \item le nombre de points VIS et NIR utilisés ;
  \item la qualité de l'ajustement NIR ($R^2$, proche de 1) ;
  \item la dispersion des décalages individuels (écart interquartile et écart-type, en secondes).
\end{itemize}

\subsection*{Graphiques de contrôle}
Deux graphiques sont enregistrés dans le dossier \texttt{Synchro} de la mission. Pendant la question \og Do the two curves overlap? \fg, la fenêtre montre les courbes après alignement ; en mode manuel, elle montre les deux graphiques en alternance (changement toutes les 1,5 s) :
\begin{itemize}
  \item \texttt{check\_time\_before\_alignment.png} : angle de la mire en fonction du temps pour VIS (vert) et NIR (rouge) avant correction ;
  \item \texttt{check\_time\_alignment.png} : mêmes courbes après avoir décalé la time line NIR de $\Delta t$.
\end{itemize}
Si la synchronisation est bonne, les points verts se placent sur la droite des points rouges. Ce sont ces deux graphiques qui sont présentés dans les trois cas de ce chapitre.

\textbf{Que faire en cas d'échec ?} Si les points VIS et NIR ne se superposent pas après alignement, ou si les vitesses de rotation VIS et NIR sont très différentes, revenir au Step~2 et redéfinir les images de synchronisation (en particulier la première image NIR, qui doit montrer la mire effectivement en rotation), puis relancer le Step~3.

\section{Bouton \bouton{Adjust alignment...}}\label{btn:step3_adjust}
Ouvre le réglage du décalage à la main, aussi ouvert par la réponse \bouton{No} à la question \og Do the two curves overlap? \fg. La fenêtre trace la courbe visible (verte), la courbe infrarouge au décalage calculé (rouge, en pointillés) et la courbe infrarouge au décalage réglé (rouge). Le curseur et la case numérique déplacent la courbe infrarouge : flèches du clavier par pas de 1\,ms, touches page précédente / suivante par pas de 0,1\,s.
\begin{description}
	\item[\bouton{Save and close}] (\bouton{Close} tant que le décalage n'a pas changé ; fermer la fenêtre fait de même) garde ce décalage : les graphiques de contrôle sont refaits. Si les images étaient déjà appariées, il faut refaire l'appariement.
	\item[\bouton{↺}] (en bas à gauche, affiché quand le décalage diffère du décalage calculé) revient au décalage calculé à partir des mires.
\end{description}
Si le décalage n'a pas pu être calculé, une petite fenêtre demande directement sa valeur en secondes (ajoutée aux dates infrarouges).

\section{Bouton \bouton{Pair VIS / NIR images}}\label{btn:step3_pairing}
Ce bouton est actif une fois le décalage temporel connu (calculé ou saisi à la main). Il construit la liste des couples d'images qui seront alignés au Step~4.

\begin{enumerate}
  \item La time line VIS sert de référence. Toutes les dates NIR sont corrigées : $t'_{NIR} = t_{NIR} + \Delta t$.
  \item Pour chaque image VIS, on retient les \emph{deux} images NIR les plus proches dans le temps. On mémorise pour chacune l'écart signé $dt = t'_{NIR} - t_{VIS}$ : $dt<0$ si la NIR précède la VIS, $dt>0$ si elle la suit.
  \item Les métadonnées de vol (position GPS, coordonnées UTM, altitudes, angles de lacet, tangage et roulis du drone et de la nacelle, distances, cap, etc.), connues pour les images VIS, sont interpolées aux instants corrigés des images NIR. Les angles (lacet, cap) sont interpolés par le chemin angulaire le plus court. Hors de la plage de temps VIS, la valeur de l'extrémité est conservée.
  \item Le résultat est enregistré dans le fichier \texttt{AerialPhotography\textbackslash shoot\_points.json}, utilisé par les étapes suivantes. Le journal indique le nombre de couples créés.
\end{enumerate}
Un même fichier NIR peut être utilisé pour plusieurs images VIS (par exemple si la fréquence de prise de vue NIR est plus élevée) et aucun écart temporel maximal n'est imposé : la qualité de l'appariement dépend donc de celle de la synchronisation.

\section{Bouton \bouton{Close}}\label{btn:step3_close}
Ferme la fenêtre et rend la main au dialogue principal. Les fichiers produits (graphiques, images de contrôle, \texttt{shoot\_points.json}) restent dans le dossier de la mission. Le bouton \bouton{?} ouvre cette aide à la page de ce chapitre.

%==========================================================
\chapter{Step 4 -- Traitement des images}\label{ch:step4}
\section{Principe}
Le traitement des images consiste pour une paire VIS-NIR à aligner ces deux images pour que les bandes spectrales Red Green Blue et Nir se superposent. L'opération est complexe pour plusieurs raisons :
\begin{itemize}
	\item Même si les time lines sont synchronisées, les instants de prise de vue sont proches mais distincts.
	\item Les deux caméras ne pointent pas vers le même point au sol. La caméra VIS du drone DJI est fixée sur un gimbal très performant qui maintient la visée vers le nadir. À noter que le vol hyperlapse est fait en mode \og gimbal lock \fg~autrement dit l'angle de roll du gimbal reste fixe. Seuls ses angles de pitch et yaw varient pour maintenir la visée vers le nadir (à $90^{\circ}$ vers le bas). En revanche la caméra NIR SJCam M20 étant fixée rigidement sous le corps du drone sa visée est totalement tributaire de l'attitude du drone (angles de pitch et yaw en particulier). 
	\item Les distances focales des deux caméras sont différentes. Pour la caméra VIS on a un angle (FOV) de $84^{\circ}$ alors que pour la caméra NIR la focale est ultra-courte avec un angle (FOV) supérieur à $150^{\circ}$ ce qui correspond quasiment à un \og fish eye \fg.
	\item Les optiques présentent de la distorsion et du vignetage. Pour la distorsion les deux caméras ont été préalablement calibrées par la méthode de Zhang.
\end{itemize} 



En pratique on commence par convertir les images du format raw-dng au format tif (via le logiciel tiers RawTherapee
\footnote{Via le logiciel tiers open source RawTherapee. L'opération est coûteuse la première fois mais l'image est mise en cache pour des runs ultérieurs.}).  On élimine la distorsion des deux images. On détermine les déplacements théoriques à appliquer pour ramener le centre de visée NIR sur le centre de visée VIS. Ces déplacements sont traduits en termes d'angles pitch, yaw et roll. Ils servent à initialiser la procédure d'alignement. L'alignement se déroule très schématiquement en deux temps. On détermine en premier les paramètres d'une homographie à appliquer sur l'image NIR pour qu'elle \og match \fg~avec l'image VIS. C'est un processus de minimisation complexe spécialement développé pour les images multispectrales 
	\footnote{Les techniques habituelles pour aligner deux images, qui sont extrêmement performantes dans le cas d'une bande spectrale unique ou en RGB, comme par exemple les \og sift \fg, échouent.}. 
Ce premier alignement (coarse process) est suivi d'un alignement fin (fine process) qui applique des homographies locales pour diminuer encore le critère d'alignement. L'algorithme est décrit en annexe (§\ref{sec:ann_algorithme}). L'utilisateur doit être conscient que la méthode n'est pas infaillible et surtout qu'elle est tributaire des conditions de vol sur le terrain. 

\figurePaireIRD{fig:step4_pair_VIS_NIR}
{step3_img_VIS.png}{Spectre visible (VIS).}
{step3_img_NIR.png}{Spectre proche infrarouge (NIR).}
{Couple d'images d'un shoot point. Les deux bandes spectrales enregistrent des informations différentes. Par exemple la surface de la rivière est totalement noire en proche infrarouge. L'eau ne réfléchit pas les longueurs d'onde au-delà de $750\,nm$.}  

\section{Dialogue de l'étape 4}
La fenêtre de l'étape~4 (figure~\ref{fig:step4_window}) rappelle le nom de la mission (le chemin complet est dans l'info-bulle) et le nombre de paires d'images. Elle comporte, de haut en bas : les options de sortie, le choix des images alignées avec la carte du vol, le bouton \bouton{Run alignment} (\bouton{Stop} pendant le calcul), l'état et la barre de progression du calcul, le journal, puis, en bas, \bouton{?} et \bouton{⚙} (réglages) à gauche, \bouton{Close} et, une fois des images alignées, \bouton{Continue with the review} à droite. La fenêtre s'adapte à la taille de l'écran.

\figureIRD[fig:step4_window]{step4_window.png}{Fen\^etre de traitement des \og shoot points \fg~(paires d'images VIS-NIR), avant le calcul.}{0.7}


\section{Options de sortie}\label{btn:step4_options}


\subsection{Case \bouton{first VIS image number} et \bouton{last VIS image number}}
Ces deux cases permettent de sélectionner des shoot points à traiter. Par défaut IRDrone propose la plage qui a été définie à l'étape~2. Il est possible de ne traiter qu'une image en entrant le même numéro pour first et last images, sous réserve que le \texttt{Maximum VIS/NIR time difference} soit compatible. En cas de doute fixez le \texttt{Maximum VIS/NIR time difference} à la valeur $3 \, s$.

Par défaut la procédure d'alignement produit  pour 

$\texttt{first VIS image number} \leq XXXX \leq \texttt{last VIS image number}$:
\begin{itemize}
	\item les quadruplets d'images multispectrales red (1), green (2), blue (3), nir (4) (\texttt{IRD\_XXXX\_n.tif}) placés dans le dossier \texttt{ImgIRdrone\textbackslash imgMultispectral}.
	\item les images \texttt{VIS\_XXXX.jpg}, \texttt{VIR\_XXXX\_local.jpg}.
	\item les fichiers \texttt{VIS\_XXXX\_REDO.bat} et deux fichiers techniques \texttt{VIS\_XXXX\_motion\_model.json}, \texttt{VIS\_XXXX\_motion\_model.npy}.
	\item le fichier \texttt{ImgIRdrone\textbackslash shoot\_points.json} qui contient un dictionnaire des shoot points traités avec en particulier la clé  \texttt{"failed"} dont la valeur est fixée à \texttt{false} par défaut. L'utilisateur vérifie cette hypothèse dans la revue des images alignées (chapitre~\ref{ch:review}). 
\end{itemize}


\subsection{Case \bouton{NDVI}, case \bouton{NIR}}
Par défaut la procédure d'alignement produit les images VIS et VIR au format jpeg. C'est utile pour vérifier visuellement la réussite des alignements. En cochant les cases \texttt{NDVI} et/ou \texttt{NIR} on peut générer les images NDVI et NIR au format jpeg. Ne cocher ces cases que si vous envisagez d'utiliser ces représentations directement dans un rapport. Notez qu'à l'étape~6 vous avez la possibilité de générer de telles représentations, mais cette fois géolocalisées et orthorectifiées, et visualisables dans le logiciel QGIS.
\subsection{Cases \bouton{Local alignment} et \bouton{Global alignment}}
Comme expliqué plus haut l'algorithme d'alignement procède en deux passes (coarse process et fine process). On peut voir le résultat de l'une ou l'autre (ou les deux) en cochant les cases \texttt{Global alignment} et \texttt{Local alignment}. La case \texttt{Global alignment} montre les résultats du coarse process alors que la case \texttt{Local alignment} montre le résultat du raffinement.  Une de ces deux cases doit impérativement être cochée. Nous suggérons de cocher systématiquement la case \texttt{Local alignment}. Il est à noter que dans certaines circonstances l'alignement global réussit mais l'alignement local échoue.

\subsection{Case \bouton{Maximum VIS/NIR time difference}}
Cette case permet de restreindre la liste des shoot points à aligner à ceux dont l'écart temporel de rang~1 vérifie $|dt|\leq \texttt{Maximum VIS/NIR time difference}$. La valeur par défaut est $0.5\,s$. Cela correspond aux images bien synchronisées. En pratique avec les fréquences de prise de vue de $2 \, s$ pour les images VIS et $3 \, s$ pour les images NIR cela revient à sélectionner un shoot point sur trois.


\section{Best mapping : n'aligner que les images utiles à la carte}\label{btn:step4_best_mapping}
Le cadre \bouton{Images aligned} propose deux choix, qui reprennent celui de l'étape~2 (§\ref{btn:step2_best_mapping}) : \bouton{Best mapping: only the images the map needs} (par défaut, avec un recouvrement de $75\,\%$) ou \bouton{All the synchronised images}. Parmi les shoot points bien synchronisés (case \texttt{Maximum VIS/NIR time difference}), IRDrone ne garde que ceux dont les empreintes au sol se recouvrent de la valeur choisie (case \bouton{Overlap between kept images}) : la première image du survol, puis, à chaque fois, la dernière image qui recouvre encore assez la précédente gardée, et la dernière image du survol. Toutes les images du survol servent à suivre le déplacement du drone ; seules les images bien synchronisées peuvent être gardées. Quand l'image qui serait gardée est mal synchronisée, IRDrone garde la plus proche bien synchronisée : la dernière avant elle ou, si son recouvrement est plus proche de la cible, la première après elle (recouvrement au plus 5 points sous la valeur choisie, 70\,\% pour 75\,\%). Le nombre d'images alignées reste ainsi proche de l'aperçu de l'étape~2 ; sans cela, chaque image mal synchronisée raccourcirait le pas et beaucoup plus d'images seraient gardées. La ligne sous les choix indique le nombre de shoot points qui seront alignés, par exemple \og 18 of 39 synchronised shoot points will be aligned \fg.

À droite, la carte du vol montre les shoot points : en vert ceux qui sont bien synchronisés et dans la plage d'images, en gris les autres ; les images gardées par le best mapping sont plus grosses, avec leur empreinte au sol en vert, ce qui montre le recouvrement. Elle se met à jour dès qu'un réglage change (plage, écart de temps, recouvrement, choix). Molette pour zoomer, glisser pour se déplacer, double-clic pour revoir tout le vol. Pendant l'alignement, la carte suit le calcul le long du vol : les empreintes des images alignées se remplissent de vert, celles en cours d'alignement (une par processus de calcul) sont entourées d'orange, celles qui restent à faire sont en pointillés ; le compte \og 7 of 18 done \fg{} s'affiche en bas de la carte. Le résultat reste affiché à la fin, jusqu'au prochain changement de réglage.

C'est en général le meilleur choix : beaucoup moins de calcul (l'alignement est l'étape la plus longue), moins de données, et une carte ODM souvent meilleure (trop d'images presque identiques ralentissent ODM et dégradent son résultat). Choisissez \bouton{All the synchronised images} pour aligner toutes les images synchronisées, par exemple pour les regarder une à une dans la revue. Le choix est enregistré dans \texttt{FlightAnalytics\textbackslash best\_mapping.json} au lancement de l'alignement ; l'étape~6 l'utilise pour la sélection des images d'ODM (§\ref{btn:step6_selection}).

La même sélection existe en ligne de commande : \texttt{run\_after\_preprocessing.py --best-mapping-overlap 75}.

\section{Bouton \bouton{Run alignment}}\label{btn:step4_run_alignment}
Lance la procédure d'alignement sur les shoot points vérifiant simultanément les deux contraintes 
\begin{itemize}
	\item $\texttt{first VIS image number} \leq XXXX \leq \texttt{last VIS image number}$,
	\item $|dt|\leq \texttt{Maximum VIS/NIR time difference}$,
\end{itemize} 
et, si \bouton{Best mapping: only the images the map needs} est choisi, seulement les images utiles à la carte (§\ref{btn:step4_best_mapping}).
Attention l'alignement des images est une procédure qui prend beaucoup de temps. Cela dépend bien entendu de la machine. Les  calculs sont parallélisés sur plusieurs processeurs si la machine est \og multi c\oe urs \fg. La barre de progression indique l'avancement des calculs ainsi qu'une estimation du temps de traitement restant. À titre d'exemple, sur un ordinateur portable avec un processeur Intel Core~i7 à quatre c\oe urs et 16\,Go de mémoire (3 workers d'alignement), l'alignement des images VIS et NIR d'un shoot point prend environ $15\,s$ : il faut donc environ une demi-heure pour traiter 120 images. Pour une grande mission, lancer ces calculs à l'heure du déjeuner reste une bonne idée.   

L'alignement ne démarre jamais tout seul, même quand l'étape~4 est ouverte par l'étape~3. Pendant le calcul, le bouton devient \bouton{Stop} (§\ref{btn:step4_stop}).

\section{Bouton \bouton{Stop}}\label{btn:step4_stop}
Pendant le calcul, \bouton{Run alignment} devient \bouton{Stop}. Après confirmation, le calcul et tous ses processus de travail s'arrêtent ; la fenêtre indique combien d'images ont été alignées. Les images déjà alignées sont gardées ; \bouton{Run alignment again} relance le calcul, et les points déjà faits réutilisent leur modèle de mouvement (plus rapide). Pour ne traiter que la suite, augmentez \texttt{First VIS image number}.

\section{Bouton \bouton{⚙} : réglages}\label{btn:step4_settings}
La roue dentée ouvre les réglages du calcul, gardés pour les sessions suivantes :
\begin{itemize}
	\item \texttt{Alignment worker processes} : nombre de shoot points alignés en parallèle. \texttt{Auto} utilise les c\oe urs physiques, dans la limite de la mémoire disponible au lancement (environ 3\,Go par processus). \bouton{Safe mode (1)} n'en garde qu'un : plus lent, mais le moins de mémoire. À utiliser si l'ordinateur se fige ou si le calcul est interrompu.
	\item \texttt{Log detail} : quantité de messages du calcul dans le journal de la fenêtre : avertissements seulement, étapes du calcul, ou détail des temps. Avec le détail des temps (\emph{Debug}), un réalignement à la main depuis la revue (section~\ref{btn:review_redo}) écrit aussi toutes ses images de diagnostic : plusieurs minutes par point au lieu d'environ une demi-minute.
\end{itemize}
\bouton{Save and close} (\bouton{Close} tant que rien n'a changé) garde les réglages. \bouton{↺}, affiché quand les réglages diffèrent des réglages par défaut, revient à \texttt{Auto} et aux avertissements seulement.

\section{Bouton \bouton{Close}}\label{btn:step4_close}
Ferme la fenêtre et rend la main au dialogue principal. Le bouton \bouton{?} ouvre cette aide à la page de ce chapitre. Il est grisé pendant l'alignement (arrêtez-le d'abord avec \bouton{Stop}).

\section{Bouton \bouton{Continue with the review}}\label{btn:step4_continue}
Apparaît en bas à droite, après \bouton{Close}, une fois des images alignées, et est mis en avant à la fin de l'alignement. Il ferme la fenêtre et ouvre l'étape~5, la revue des images alignées (chapitre~\ref{ch:review}). Quand l'alignement a été lancé depuis l'étape~3, la revue s'ouvre d'elle-même à la fin.

%==========================================================
\chapter{Step 5 -- Revue des images alignées}\label{ch:review}
La revue présente tous les points du relevé sur une ligne de temps pour vérifier l'alignement de chaque image (figure~\ref{fig:step4_review}). Elle s'ouvre avec \bouton{Continue with the review} de l'étape~4, ou avec l'étape~5 de la fenêtre principale (une fois des images alignées par l'étape~4). Elle remplace l'examen des images une par une dans l'explorateur. Rien n'est calculé à l'ouverture : l'étape compte comme faite une fois la revue fermée, les signalements et confirmations étant enregistrés au fur et à mesure.

\figureIRD[fig:step4_review]{step4_review.png}{Revue des images alignées : rideau entre visible et infrarouge aligné, fiche du point (décalage mesuré, image infrarouge associée), mini-carte, et en bas la ligne de temps des points (décalage résiduel, vignettes, état).}{0.95}

\section{Mesure automatique}\label{sec:review_measure}
Pour chaque point, le décalage résiduel entre les contours de l'image visible et ceux de l'infrarouge aligné est mesuré (en pixels de l'image complète, médiane sur 16 zones). Un point bien aligné mesure environ 1~px. Vert : moins de 3~px ; orange clair : entre 3 et 6~px, à regarder ; orange : plus de 6~px, probablement mal aligné. Une zone sans texture (eau, champ uniforme) ne se mesure pas. La fenêtre s'ouvre sur le point le plus douteux.

Les mesures et les points vérifiés sont gardés dans \texttt{ImgIRdrone/review.json} : rouvrir la revue ne refait pas les mesures, sauf pour les images refaites.

\section{Comparer le visible et l'infrarouge}\label{btn:review_compare}
Quatre façons de comparer, avec les boutons au-dessus de l'image ou les touches 1 à 4 ; la revue s'ouvre en \bouton{Overlay}. Un cadre autour de l'image est dans la couleur de l'état du point (vert aligné ou confirmé, rouge signalé, violet réaligné\ldots), les étiquettes sur l'image restent neutres (les rouges et cyans sont ceux des images) :
\begin{description}
	\item[\bouton{Swipe} (1)] un rideau sépare le visible (à gauche) de l'infrarouge aligné en fausses couleurs (à droite) ; routes, haies et toits doivent se prolonger à travers le rideau. Glisser le rideau pour le déplacer.
	\item[\bouton{Blink} (2)] visible et infrarouge alternent : une image mal alignée \og saute \fg.
	\item[\bouton{Overlay} (3)] visible en cyan, infrarouge en rouge (la végétation est naturellement rouge) ; un mauvais alignement double chaque contour.
	\item[\bouton{Side by side} (4)] les deux images côte à côte.
\end{description}
Molette : zoom (un décalage de quelques pixels ne se voit qu'en zoomant) ; glisser : déplacement ; double-clic : zoom $\times 4$ ou image entière. \bouton{Open full size} ouvre l'image alignée dans la visionneuse d'images de l'ordinateur.

\section{Situer l'image sur la carte}\label{btn:review_map}
La mini-carte (dans la fiche du point) montre le vol, le point courant et l'emprise au sol de l'image (le trait épais est le haut de l'image), dans la couleur de son état (vert aligné ou confirmé, rouge signalé\ldots) ; les shoot points que l'étape~4 n'a pas alignés sont de petits points gris, en taille normale quand on zoome fort. Un clic l'agrandit (\bouton{Map}, touche 5). Sous la mini-carte, la légende ne donne que les états présents, avec leur nombre. La grande carte affiche l'image alignée dans son emprise, par-dessus l'orthophoto de l'étape~6 quand elle existe et les photos aériennes en fond : l'infrarouge doit prolonger les routes et les haies du fond. Survoler un point montre sa vignette, un clic le sélectionne ; un clic hors des points ne laisse aucun point sélectionné sur la carte : la trajectoire seule, sans emprise ni image (le point reste celui de la fiche ; un clic sur un point, ou un déplacement dans le bandeau, en sélectionne un à nouveau) ; molette et glisser pour zoomer et se déplacer, double-clic pour voir tout le vol. Le menu \bouton{Layers}, le fond de carte hors ligne et le calcul de l'emprise sont décrits au §\ref{btn:home_map}. Sans position GPS dans les shoot points, la carte n'est pas disponible.

\paragraph{Aligner un point laissé de côté par l'étape~4.} Un clic droit sur un point gris de la carte (mini-carte ou grande carte) affiche son écart de temps avec l'infrarouge et propose \bouton{Align this point} : il est aligné comme le fait l'étape~4. Quand une image est irrécupérable, même à la main, un voisin écarté par le best mapping ou le filtre temporel peut être bon : alignez-le ainsi. Le point est aligné en arrière-plan (orange sur la carte, avancement dans le bandeau, \bouton{Stop aligning} pour arrêter) pendant que vous continuez la revue ; il apparaît ensuite dans la frise, à revoir comme les autres, et l'étape~6 le cartographie toujours (sauf si vous le signalez), quels que soient sa tolérance temporelle et son choix d'images. Le filtre temporel est le plus large de l'étape~4 ($|dt|<6$~s) ; un point sans image infrarouge assez proche dans le temps ne peut pas être aligné.

\figureIRD[fig:step4_review_map]{step4_review_map.png}{Mode \og Map \fg~de la revue : vol, emprise et image alignée du point courant, orthophoto de l'étape~6 sur les photos aériennes IGN.}{0.9}

\section{Se déplacer}\label{btn:review_navigation}
\begin{itemize}
	\item \bouton{◀} et \bouton{▶} (flèches du clavier) : point précédent, point suivant.
	\item \bouton{Next to review} (touche N) : prochain point qui demande un coup d'\oe il (réaligné à confirmer, mesuré douteux ou non mesuré). Un point mesuré aligné (vert) ne demande rien : il n'y a rien à marquer comme bon.
	\item La ligne de temps du bas : survoler montre l'image, un clic sélectionne le point ; molette pour zoomer, glisser pour se déplacer, double-clic pour tout le relevé.
	\item \bouton{Close} ferme la revue ; ce qui a été signalé et confirmé est déjà enregistré. \bouton{Continue with mapping} la ferme aussi, puis ouvre l'étape~6 (§\ref{btn:review_continue}).
	\item L'info-bulle du bouton \bouton{?} rappelle toutes les touches.
\end{itemize}

\section{Décider : signaler un point}\label{btn:review_decide}
\begin{description}
	\item[\bouton{Flag: leave out of the map} (touche F)] signale le point, \textbf{ce qui l'écarte de la carte} : mal aligné, mais aussi flou, ou avec une zone noire au bord de l'image infrarouge (drone incliné par un coup de vent : la caméra infrarouge ne voit pas tout le champ du visible). Un second appui retire le signalement. Le signalement est écrit tout de suite dans \texttt{ImgIRdrone\textbackslash shoot\_points.json}\footnote{Ce fichier ne contient pas les mêmes informations que le fichier \texttt{AerialPhotography\textbackslash shoot\_points.json}.} et l'étape~6 ignore ce point. Un point signalé peut être réaligné à la main (§\ref{btn:review_redo}) ; si cela ne le rattrape pas, le laisser signalé.
	\item[\bouton{Flag the N probably misaligned}] signale d'un coup tous les points dont le défaut d'alignement mesuré dépasse 6~px (points rouges \emph{probably misaligned}). Il n'apparaît que s'il y en a : la mesure ne voit ni un flou ni une zone noire, que vous signalez vous-même.
\end{description}

L'exemple ci-dessous montre le début du fichier \texttt{ImgIRdrone\textbackslash shoot\_points.json} qui stocke ces informations. Toutes les images ont été traitées en choisissant le critère peu discriminant $|dt|<4\,s$ et la méthode "au plus près". Cela correspond à 501 images VIS et 350 images NIR. Ici l'utilisateur a estimé visuellement que le shoot point 120 était mal aligné et l'a signalé, ce qui a fait basculer la valeur de la clé \texttt{failed} de \texttt{false} à \texttt{true}.
\begin{verbatim}
	{
		"mission_path": "C:\\Air-Mission\\FLY-20220524-1715-Chalivoy-Milon",
		"time_shift": -48.116,
		"n_VIS": 501,
		"n_NIR": 350,
		"pairing_method": "nearest_time",
		"created_at": "2026-10-29T20:44:13.203434",
		"schema_version": "1.0",
		"shoot_points": [
		{
			"FileName": "VIS_0119.dng",
			"dt": -0.5390370000000075,
			"failed": false
		},
		{
			"FileName": "VIS_0120.dng",
			"dt": 0.6902460000000019,
			"failed": true
		},
		{
			"FileName": "VIS_0121.dng",
			"dt": -1.3297540000000083,
			"failed": false
		},
		...
		]
	}	 
\end{verbatim}

Que faire ensuite des shoot points mal alignés ? Deux options :
\begin{itemize}
	\item \og aider \fg~l'algorithme d'IRDrone à les aligner : \bouton{Realign N flagged by hand} (§\ref{btn:review_redo}). Selon le résultat, le signalement est gardé ou retiré.
	\item les écarter des traitements ultérieurs : c'est l'effet par défaut du signalement, l'étape~6 ignore les points \texttt{"failed": true}.
\end{itemize}

\subsection{Exemple de succès d'alignement}
\figurePaireIRD{fig:goodalignment}
{step4_goodalignment_VIS.jpg}{Image dans le spectre visible.}
{step4_goodalignment_NIR.jpg}{Image la plus proche $|dt|=0.165~s$ dans le spectre infrarouge.}
{Cas de succès d'alignement : l'écart temporel est faible et l'attitude du drone est restée dans les spécifications des paramètres de vol.}

\figureIRD[fig:step4_goodalignment_VIR]{step4_goodalignment_VIR.jpg}{Alignement réussi. Les deux images ont été prises presque au même moment et en l'absence de vent. }{0.5}

\subsection{Exemple d'échec d'alignement}
Si le drone sort de sa zone des paramètres de vol, par exemple en cas de fortes rafales de vent qui l'obligent à accentuer ses angles de pitch et yaw, alors les centres de visée des deux caméras peuvent être tellement éloignés qu'il est impossible de trouver un alignement. L'alignement peut aussi être défaillant si l'écart temporel entre les prises de vue VIS et NIR est important (il peut être parfois supérieur à 2\,s) car le déplacement horizontal du drone est alors mal corrigé. L'alignement peut aussi être imparfait sans être un échec s'il y a un écart temporel et que les végétaux bougent sous l'action du vent. L'image VIR présente alors des \og fantômes \fg. Les figures \ref{fig:misalignment} et \ref{fig:step4_misalignment_VIR} montrent un cas d'échec de l'alignement lié à une forte rafale de vent qui déstabilise le drone. 

\figurePaireIRD{fig:misalignment}
{step4_misalignment_VIS.jpg}{Image dans le spectre visible.}
{Step4_misalignment_NIR.jpg}{Image la plus proche $|dt|=1.256~s$ dans le spectre infrarouge.}
{Cas d'échec d'alignement lié à une attitude du drone hors des spécifications des paramètres de vol et liée à une rafale de vent au moment de la prise de vue.}

\figureIRD[fig:step4_misalignment_VIR]{step4_misalignment_VIR.jpg}{Échec de l'alignement. L'effet de \og trou \fg est particulièrement caractéristique. Il peut aisément être repéré par l'utilisateur sur les images VIR. }{0.5}

\section{Réalignement à la main}\label{btn:review_redo}
Le bouton \bouton{Realign N flagged by hand} rattrape les shoot points signalés mal alignés. Il n'y a pas de miracle : si une rafale de vent a écarté le drone de sa zone de vol, les deux caméras ne voient pas la même scène et aucun alignement n'est possible. Mais parfois l'information existe et il suffit de \og pousser \fg~les paramètres de l'alignement grossier. En conditions normales une dizaine de shoot points au plus sont à reprendre.

\subsection*{Déroulement}
\begin{enumerate}
	\item Signaler les points mal alignés, puis appuyer sur \bouton{Realign N flagged by hand}. La fenêtre d'alignement s'ouvre avec tous les points signalés (figure~\ref{fig:step4_align_by_hand}).
	\item Pour chaque point, déplacer l'image infrarouge sur l'image visible jusqu'à faire coïncider les mêmes éléments (voir ci-dessous), puis \bouton{Placed, next} (touche Entrée). \textbf{L'alignement de ce point démarre aussitôt, en arrière-plan}, pendant que vous placez le suivant : pas d'attente entre deux points (le point passe à \emph{⚙ Computing} dans la liste et ne peut plus être déplacé). Les images des points suivants sont chargées en arrière-plan. \bouton{Skip this point} (touche S) laisse un point de côté : il reste signalé, donc écarté de la carte.
	\item \bouton{Compute N alignments} envoie au calcul les points placés sans \bouton{Placed, next} (quand tout est déjà envoyé, \bouton{Cancel} devient \bouton{Close}) et ferme la fenêtre. Les alignements sont calculés en arrière-plan, plusieurs à la fois (autant que de workers d'alignement de l'étape~4, bouton engrenage). Compter environ une demi-minute par point. Le bandeau du haut de la revue suit l'avancement : points terminés, points en cours avec leur étape et le temps écoulé, points en attente. \textbf{La revue reste utilisable} : on peut continuer à vérifier les autres points, ou faire autre chose. Les points en cours de calcul ne peuvent pas être signalés ni dé-signalés. \bouton{Stop} ne lance plus de nouvel alignement : ceux en cours se terminent, les autres restent signalés. La revue ne peut pas être fermée tant que des alignements sont en cours.
	\item Chaque point réaligné passe en violet, \emph{Realigned, to confirm} : il reste signalé (l'étape~6 l'ignore). Quand tout est calculé, la revue affiche le premier. Le vérifier avec le rideau, le clignotement ou la superposition et répondre \bouton{Yes, aligned} (touche Y) ou \bouton{Still misaligned} : le point suivant à confirmer s'affiche alors. \bouton{All N aligned} les confirme tous d'un coup. La touche N (point suivant à revoir) y va en premier. Les points à confirmer sont mémorisés si la revue est fermée.
	\item Avec \bouton{Yes, aligned}, la clé \texttt{"failed"} du point passe à \texttt{false} dans \texttt{shoot\_points.json} et le fichier \texttt{shoot\_points\_failed\_summary.json} est mis à jour (une copie \texttt{shoot\_points.json.bak} est conservée). Avec \bouton{Still misaligned}, le point reste signalé et pourra être réaligné de nouveau ; s'il ne peut pas être rattrapé (image infrarouge trop décalée, flou), il suffit de le laisser signalé : \textbf{il est écarté de la carte} (étape~6).
\end{enumerate}

Les sorties des alignements précédents sont archivées dans le dossier \texttt{VIS\_XXXX\_align\_traces\_manual\textbackslash original}. Le point retrouve les images qu'il avait (celles cochées à l'étape~4, section~\ref{btn:step4_options}) ; les images de diagnostic de l'alignement (dossier \texttt{Tech}, plusieurs minutes de calcul par point) ne sont écrites que si \texttt{Log detail} est réglé sur \emph{Debug} dans les réglages de l'étape~4 (bouton engrenage, section~\ref{btn:step4_settings}), ou depuis la ligne de commande (ci-dessous).

\subsection*{La fenêtre d'alignement}
\figureIRD[fig:step4_align_by_hand]{step4_align_by_hand.png}{Alignement à la main : les points signalés à gauche (coche : placé, rond : à placer), les deux images superposées au centre (ici le visible en cyan et l'infrarouge en rouge), les angles à droite. Ici l'infrarouge est encore environ 140~px trop haut : les contours sont doublés.}{0.9}

\begin{itemize}
	\item À gauche, la liste des points : un rond pour les points à placer, une coche pour les points placés, un tiret pour les points laissés de côté. Un clic va sur un point.
	\item Au centre, les deux images, affichées comme dans la revue : \bouton{Overlay} (les deux images en une : des contours doublés signalent un mauvais alignement ; affichage par défaut). \bouton{Colours} choisit ses couleurs, gardées pour la fois suivante : visible en couleur et infrarouge en gris (par défaut ; avec \bouton{Edges only}, les contours du visible sont verts), visible en vert et infrarouge en gris, visible en bleu et infrarouge en rouge, ou visible en cyan et infrarouge en rouge (comme dans la revue, proche de l'image VIR). \bouton{Swipe}, \bouton{Blink}, \bouton{Side by side} (touches 1 à 4). \bouton{Edges only} (touche G) montre les contours des deux images, ce qui aide sur les champs et les forêts où l'infrarouge et le visible se ressemblent peu. Molette pour zoomer, glisser pour se déplacer, double-clic pour voir toute l'image.
	\item Pour déplacer l'infrarouge, \textbf{vite puis finement} : d'abord les \textbf{curseurs} de \texttt{Yaw}, \texttt{Pitch} et \texttt{Roll} (panneau de droite) pour un grand déplacement, puis les boutons à leurs extrémités pour affiner, dans le sens qu'ils montrent : \bouton{◀}~\bouton{▶} à gauche et à droite (lacet), \bouton{▲}~\bouton{▼} en haut et en bas (tangage), \bouton{↺}~\bouton{↻} rotation (roulis) ; maintenus enfoncés, ils se répètent. Les \textbf{flèches du clavier} font la même chose que ces boutons (0,05\textdegree{} par appui, environ 2 à 3 pixels de l'image complète ; avec Maj, 10 fois plus), où que soit le curseur dans la fenêtre, même dans un champ de valeur ; Q et E le font tourner. On peut aussi \textbf{faire glisser l'infrarouge avec le bouton droit de la souris} (ou Maj + glisser) dans l'image, ou taper une valeur dans un champ.
	\item \bouton{Fit} (touche F) mesure le décalage entre les contours des deux images (tuile par tuile) et déplace l'infrarouge d'autant (lacet et tangage). Il fonctionne dès que les images se recouvrent à peu près ; s'il ne trouve pas de meilleure position, il le dit et ne change rien. \bouton{Reset} revient aux angles de l'alignement précédent.
	\item Sous ces boutons, le décalage mesuré entre les deux images, en pixels de l'image complète : sous 10~px le point est \emph{well placed}.
	\item À peu près suffit : ces angles ne sont qu'un point de départ, l'alignement grossier et l'affinage pyramidal repartent de là (§\ref{sec:ann_algorithme}).
	\item Un point dont les images ne peuvent pas être lues est marqué d'un signe d'avertissement et laissé de côté.
\end{itemize}

\subsection*{En ligne de commande : la fenêtre à curseurs}
Chaque point aligné a un script \texttt{VIS\_XXXX\_REDO.bat} (Windows) ou \texttt{VIS\_XXXX\_REDO.sh} dans \texttt{ImgIRdrone}, qui relance son alignement avec une fenêtre à curseurs (matplotlib) : c'est l'ancienne procédure, toujours utilisable en dehors de l'assistant. Après un réalignement dans l'assistant, le script repart des angles placés à la main.


\figureIRD[fig:redo_info]{step4_redo_info.png}{Fenêtre d'information de la fénêtre d'aide manuelle à l'alignement. }{0.5}

La fenêtre affiche l'image VIS et l'image NIR, avec cinq curseurs : \texttt{YAW}, \texttt{PITCH}, \texttt{ROLL}, \texttt{Absgrad visualize} et \texttt{Alpha}.
\begin{itemize}
	\item Le curseur \texttt{Alpha} choisit ce que l'on voit. Entre $-1$ (tout à gauche) et $-0{,}5$ (petit trait rouge médian), un rectangle de l'image NIR, légèrement transparent, est superposé à l'image VIS : plus \texttt{Alpha} est proche de $-1$, plus ce rectangle est grand. La valeur initiale est $-0{,}75$.
	\item Les curseurs \texttt{YAW}, \texttt{PITCH} et \texttt{ROLL} déplacent l'image NIR. La méthode consiste à modifier un angle, puis à faire varier \texttt{Alpha} pour observer le résultat, et à recommencer avec les trois angles. La clé est de repérer des éléments identiques dans les deux bandes spectrales (arbre, bord de chemin, bâtiment) et de les faire coïncider, d'abord au centre puis sur les bords du rectangle.
	\item Lorsque \texttt{Alpha} est déplacé vers la droite, la fenêtre n'affiche plus un rectangle NIR : entre $-0{,}5$ et $0$ elle montre la différence entre les deux images, et entre $0$ et $1$ un fondu enchaîné VIS/NIR. Cet usage est plus rapide mais réservé à un opérateur expérimenté.
	\item \texttt{Absgrad visualize} (valeur initiale $0{,}3$) affiche les gradients des deux images en deux couleurs lorsqu'il dépasse $0{,}5$ : les contours coïncident quand l'alignement est bon.\\
\end{itemize}

Lorsque l'utilisateur considère que l'alignement manuel est satisfaisant il tape la touche Q et la procédure d'alignement semi automatique démarre. \\


Comme pour tout outil de ce type, c'est en l'expérimentant que l'on se l'approprie ; une fois la méthode comprise, le réalignement est rapide et efficace.
Voici quelques étapes en image de la procédure
\footnote{L'outil est riche et offre plusieurs possibilités de visualiser les couches. Donc il n'y a pas une méthode unique pour arriver au résultat souhaité. Le slider \texttt{Absgrad visualize} permet de changer les représentations. Par exemple en le déplaçant vers la droite on peut obtenir un effet de contours. Ce n'est pas toujours évident à interpréter avec des forêts par exemple et encore moins dans des zones sans relief et nues.}. 
Sur l'image il y a une route qui traverse en diagonale. 
Il y a aussi un gros arbre en haut à droite, un petit arbuste en bas à gauche le long de la route  ainsi qu'un fossé en haut à gauche (\og trait noir \fg). Ce sont ces éléments caractéristiques que l'on va essayer d'aligner. 
À noter qu'à l'initialisation la couche NIR est légèrement transparente
\footnote{En déplaçant le slider \texttt{Alpha} vers la droite entre 0 et 1 la couche NIR passe de totalement transparente à opaque.}. 
Il n'y a pas vraiment d'ordre pour arriver à un résultat satisfaisant. 
Ne cherchez pas la perfection  car on n'applique que des transformations globales linéaires. Le code affinera globalement puis localement.
\figurePaireIRD{fig:redo_A}
{step4_redo_1.png}{Situation de départ après un passage dans la procédure associée au bouton \texttt{run alignement}.}
{step4_redo_2.png}{On a joué sur la rotation, qui passe de $0$ à $-2.46$, pour aligner la route NIR parallèlement à la route VIS. Attention utiliser ce slider délicatement. Lorsqu'il n'y a pas de ligne évidente comme ici la route, on utilise plutôt la rotation en dernier pour mettre la \og touche finale \fg.}
{Première modification en jouant sur le slider de rotation (\texttt{ROLL}).}
\figurePaireIRD{fig:redo_B}
{step4_redo_3.png}{Translation verticale vers le bas avec le slider \texttt{PITCH}. La valeur passe de  $-1.14$ à $2.04$. C'est une amplitude relativement importante et surtout qui change de signe.}
{step4_redo_4.png}{Translation horizontale avec le slider \texttt{YAW}. La valeur passe de $0.96$ à $1.8$.}
{On affine en jouant sur le curseur \texttt{Alpha} (en diminuant puis en agrandissant) et en observant l'alignement des éléments caractéristiques. La sensation de léger flou à la fin est liée à la transparence et aux petits défauts d'alignement locaux qui seront corrigés par la procédure d'alignement fin.}

\section{Bouton \bouton{Continue with mapping}}\label{btn:review_continue}
En bas à droite, après \bouton{Close}. Il ferme la revue et ouvre l'étape~6. La carte, longue, n'y démarre pas toute seule : elle se lance avec \bouton{Build the map\ldots}. Les points signalés, et les points réalignés pas encore confirmés, sont écartés de la carte. Tant que des alignements sont en calcul, la revue ne peut pas être fermée : attendez-les (ou \bouton{Stop}).



%==========================================================
\chapter{Step 6 -- Cartographie et géoréférencement}\label{ch:step6}
\section{Principe}
Les images multispectrales produites à l'étape~4 sont des images \og à plat \fg~: elles ne sont pas positionnées sur le terrain.

\figurePaireIRD{fig:step6_ortho}
{step5_VIS_plan.png}{Image corrigée des distorsions de l'optique (DJI Mavic Air 2 et calibrage de Zhang).}
{step5_VIS_ortho.png}{Image Orthorectifiée avec un modèle de terrain au pas de $10\,m$.}
{L'effet de l'orthorectification est bien visible sur cette image où le terrain présente une forte déclivité (voir figure \ref{fig:step6_topo}).}

 L'étape~6 les rend exploitables dans un système d'information géographique (SIG) comme QGIS. Elle propose trois traitements indépendants, un par bouton :
\begin{itemize}
	\item \bouton{Build the map\ldots} (le bouton principal) : assemblage de l'ensemble des images retenues en une seule \emph{orthophoto} à 4 bandes (rouge, vert, bleu, proche infrarouge) avec le logiciel libre OpenDroneMap (ODM), puis le modèle de terrain et les courbes de niveau ;
	\item \bouton{Terrain model and contours\ldots} : construction à nouveau du MNT et des \emph{courbes de niveau}, par exemple avec un autre pas de grille, à partir d'altitudes fournies par un service en ligne (IGN) ;
	\item \bouton{Georeference images\ldots} : \emph{géoréférencement} et \emph{orthorectification} d'images isolées. Chaque image est projetée sur un modèle numérique de terrain (MNT) puis enregistrée au format GeoTIFF en coordonnées UTM.
\end{itemize}
Ces traitements s'appuient sur les métadonnées de vol enregistrées dans les fichiers \texttt{.exif} \og compagnons \fg~(position GPS, altitude, angles de lacet, tangage et roulis, voir étape~1).

\figureIRD{step5_window.png}{Fen\^etre de l'étape~6.}{0.6}

La fenêtre est organisée comme celle de l'étape~3. De haut en bas : le nom de l'étape et de la mission ; les étapes de la carte (\texttt{Image selection}, \texttt{Orthophoto (OpenDroneMap)}, \texttt{Terrain model and contours} : $\circ$ à faire, $\bullet$ en cours, \checkmark{} fait, $\times$ échec, avec un détail sous chacune) ; une zone de messages (journal) ; les trois boutons, \bouton{Build the map\ldots} en premier ; une barre de progression précédée d'un intitulé qui indique l'opération en cours ; une zone d'aperçu d'image ; et les boutons \bouton{?} et \bouton{Close}, puis \bouton{Continue with archiving} une fois la carte construite.

Si la carte a déjà été construite, les étapes sont déjà cochées et l'orthophoto s'affiche dès l'ouverture de la fenêtre. \bouton{Terrain model and contours\ldots} reste grisé tant que la carte n'existe pas.

\textbf{Prérequis}
\begin{itemize}
	\item l'étape~4 doit avoir été exécutée pour les images à traiter (dossier \texttt{ImgIRdrone\textbackslash imgMultispectral}) ;
	\item pour \bouton{Build the map\ldots} : le logiciel \emph{Docker Desktop} doit être installé et démarré ;
	\item pour \bouton{Terrain model and contours\ldots} (et pour l'altitude du sol lors du géoréférencement) : une connexion Internet est nécessaire pour interroger le service d'altitudes.
\end{itemize}

\section{Bouton \bouton{Georeference images\ldots}}\label{btn:step6_georef}
Ce bouton géoréférence et orthorectifie des images choisies par l'utilisateur.

Un message rappelle les types d'images acceptés, puis une fenêtre de sélection de fichiers s'ouvre. Tous les fichiers choisis doivent être du même type :
\begin{itemize}
	\item des images DNG de la caméra DJI (spectre visible) ;
	\item des images TIFF multispectrales par \textbf{groupes de quatre} : \texttt{IRD\_XXXX\_1.tif} à \texttt{IRD\_XXXX\_4.tif} (rouge, vert, bleu, proche infrarouge). Il faut sélectionner les quatre fichiers ; un groupe incomplet est rejeté avec un avertissement.
\end{itemize}

Dans une mission visible seul, le message et la fenêtre de sélection ne proposent que les images DNG DJI.

\subsection*{Décalages de géoréférencement}
Une fois les images choisies, une fenêtre demande trois corrections, toutes égales à \textbf{0} par défaut :
\begin{itemize}
	\item le \emph{décalage de lacet} (en degrés), qui corrige une erreur systématique du cap du drone ;
	\item le \emph{décalage UTM est} (en mètres) ;
	\item le \emph{décalage UTM nord} (en mètres).
\end{itemize}

Ces valeurs s'appliquent à toutes les images traitées lors de ce lancement. Les décalages sont utiles car les données GPS
\footnote{La précision de positionnement pourrait être grandement améliorée à l'aide d'une balise RTK mais ce type de matériel est onéreux.} et d'attitude du drone DJI ont une précision limitée, tout particulièrement si l'étalonnage du GPS et/ou le calibrage de l'orientation de la caméra n'ont pas été fait avant le décollage. Ces valeurs \og d'offset \fg ne peuvent être déterminer qu'expérimentalement, par exemple en utilisant QGIS pour affiner le positionnement. 

 Le bouton \emph{Cancel} annule le géoréférencement.

\subsection*{Ce que fait le logiciel pour chaque image}
\begin{enumerate}
	\item Les images DNG sont converties en TIFF avec RawTherapee. Pour les images multispectrales, les quatre bandes sont normalisées entre 0 et 1 puis empilées en une image à 4 canaux.
	\item Les métadonnées de vol sont lues dans le fichier \texttt{.exif} compagnon (créé si nécessaire). Elles donnent la position du drone, son altitude au-dessus du sol et ses angles d'attitude. Si l'altitude du point de décollage est inconnue, elle peut être demandée à l'utilisateur.
	\item L'emprise de l'image au sol est estimée à partir de la géométrie de la caméra, puis un maillage du terrain est construit. Les altitudes des points du maillage sont demandées à l'IGN.
	\item Chaque point du maillage est projeté dans l'image. L'image est rééchantillonnée sur une grille régulière au sol (orthorectification).
	\item Le résultat est enregistré au format GeoTIFF (coordonnées UTM).
\end{enumerate}

\subsection*{Fichiers produits}
Les fichiers sont créés dans le sous-dossier \texttt{geo\_ref} du dossier contenant les images sélectionnées. Un fichier existant n'est jamais écrasé : un suffixe \texttt{\_1}, \texttt{\_2}, \dots~est ajouté au nom (utile par exemple si le fichier précédent est ouvert dans QGIS).
\begin{itemize}
	\item \texttt{*\_georef.tif} : l'image orthorectifiée, à ouvrir dans QGIS. Pour un quadruplet multispectral le nom est \texttt{IRD\_XXXX\_0\_georef.tif} et l'image comporte 4 bandes (Red, Green, Blue, NIR) ;
	\item \texttt{*\_MNT.png} et \texttt{*\_Grid.png} : graphiques de contrôle du terrain et du maillage ;
	\item \texttt{*\_dtm\_map.png} : carte 2D ombrée du MNT sous l'image, avec courbes de niveau, quadrillage UTM et points d'altitude IGN 	(même présentation que pour le mapping).
	\end{itemize}
	Pendant le calcul la fenêtre ne répond plus : c'est normal. L'image en cours de traitement s'affiche dans la zone d'aperçu et le journal détaille l'avancement.

	La précision du placement dépend de la qualité des métadonnées de vol (GPS, lacet) et de celle du MNT. Un éventuel décalage systématique peut être corrigé grâce à la fenêtre des décalages (voir ci-dessus).
	
\figureIRD{step5_run_Geo_ref.png}{Fenêtre pendant le géoréférencement d'une image. C'est l'image de départ non orthorectifiée qui est affichée.}{0.6}


\section{Bouton \bouton{Build the map\ldots}}\label{btn:step6_mapping}
Ce bouton assemble les images alignées de l'étape~4 en une orthophoto unique de toute la zone survolée, grâce au logiciel OpenDroneMap (ODM). ODM s'exécute dans un conteneur \emph{Docker} : l'utilisateur n'a rien à configurer, mais Docker Desktop doit être démarré (voir plus bas).

\subsection{Choix des images}\label{btn:step6_selection}
En cliquant sur le bouton, une fenêtre de choix des critères s'ouvre. Une note en bas rappelle que le calcul peut être long et utiliser plusieurs gigaoctets de mémoire et d'espace disque. \bouton{Build the map} le lance aussitôt ; \bouton{Cancel} ferme la fenêtre sans rien lancer. Une fois une carte construite, la fenêtre propose les réglages de la dernière exécution (sinon les valeurs par défaut ci-dessous ; le best mapping choisi à nouveau aux étapes~2 ou~4 après cette exécution l'emporte). Les images peuvent venir de tout le vol, pas seulement de la phase de survol de l'étape~2.

\begin{itemize}
	\item \texttt{Maximum |dt| (NIR rank 1)} : écart temporel maximal admis entre l'image VIS et l'image NIR la plus proche. C'est le même critère que le \texttt{Maximum VIS/NIR time difference} de l'étape~4. Valeur par défaut : $0.5\,s$ (plage $0.05$ à $6\,s$).
	\item \texttt{Best mapping: select by image overlap (recommended)} : choix par défaut. Ne garde que les images dont les empreintes au sol se recouvrent de la valeur indiquée (\texttt{Overlap between successive images}, par défaut la valeur choisie à l'étape~2 ou~4, $75\,\%$ sinon). C'est la même sélection qu'aux étapes~2 et~4 (§\ref{btn:step2_best_mapping}). Si l'étape~4 n'a aligné que les images utiles à la carte, elles sont toutes gardées ; si elle a tout aligné, seules les images utiles à la carte sont données à ODM. Une image déclarée mal alignée dans la revue est écartée et la suivante prend sa place.
	\item \texttt{Select one image every n eligible images} : permet de ne retenir qu'une image sur $n$ ($n$ de 1 à 4, valeur par défaut 3) afin de limiter le temps de calcul et le recouvrement. C'est le choix par défaut quand le best mapping a été décoché à l'étape~2 ou~4.
	\item \texttt{All the aligned images (every point of the review not flagged)} : ignore les critères ci-dessus : toutes les images alignées par l'étape~4 (ou depuis la revue) et non signalées sont cartographiées, quels que soient leur écart de temps et le recouvrement. Le plus d'images, la cartographie la plus longue. Ce choix est proposé à nouveau la fois suivante s'il a été utilisé.
\end{itemize}

\figureIRD{step5_odm_selection.png}{Fenêtre des critères de sélection des images pour le mapping.}{0.4}

À droite, comme à l'étape~4, la carte du vol montre les images retenues pour la carte, en vert avec leurs empreintes au sol (leur recouvrement se voit), les images éligibles écartées par la sélection en gris, et les autres en petits points. Elle est mise à jour dès qu'un réglage change ($|dt|$, recouvrement, $n$, choix) ; le compte sous la carte donne le nombre d'images retenues parmi les éligibles. Molette pour zoomer, glisser pour se déplacer, double-clic pour revoir tout le vol.

\textbf{OpenDroneMap : version et GPU.} Sous les choix, le cadre \texttt{OpenDroneMap} indique, après une vérification de quelques secondes (Docker, Internet), la version d'ODM installée sur l'ordinateur et si une version plus récente a été publiée sur Docker Hub :
\begin{itemize}
	\item \texttt{Download the newer version before the map} : affichée seulement quand une version plus récente existe, décochée par défaut. Cochée, la dernière image d'ODM est téléchargée (\texttt{docker pull}, de quelques centaines de mégaoctets à quelques gigaoctets) avant le calcul. Une nouvelle version peut donner une carte un peu différente : elle n'est donc jamais installée sans qu'on le demande. Sans Internet, le cadre indique que la dernière version n'a pas pu être vérifiée et la version installée est utilisée.
	\item \texttt{Use the NVIDIA GPU (faster image matching)} : disponible, et cochée, seulement quand l'ordinateur a une carte graphique NVIDIA \emph{et} que Docker peut la donner à un conteneur (le logiciel lance un conteneur de test avec \texttt{--gpus all}). ODM utilise alors sa version GPU (\texttt{opendronemap/odm:gpu}, environ 3\,Go téléchargés la première fois), qui détecte et apparie les points caractéristiques des images sur la carte graphique ; le reste du calcul est le même. Sinon la case est grisée et la raison est donnée (pas de carte NVIDIA ; sous Linux, le \emph{NVIDIA Container Toolkit} manque ; sous Windows, Docker Desktop doit utiliser WSL~2 avec un pilote NVIDIA récent) ; son info-bulle donne la réponse de Docker. Si le GPU n'est pas utilisable au lancement du calcul, ODM tourne sur le processeur et le journal le dit.
\end{itemize}
Le journal donne la version d'ODM utilisée, et si elle a tourné sur le GPU ou sur le processeur.


Les images retenues (\emph{eligible images}) sont les shoot points qui vérifient simultanément :
\begin{enumerate}
	\item ils n'ont pas été signalés mal alignés dans la revue, étape~5 (clé \texttt{"failed": false}, voir §\ref{btn:review_decide}) ;
	\item l'écart $|dt|$ avec l'image NIR la plus proche est inférieur ou égal à la tolérance ;
	\item l'image VIS correspondante appartient à la phase de vol de la mission choisie à l'étape~2 (les images de décollage et de synchronisation sont donc exclues).
\end{enumerate}
Parmi elles, le best mapping garde les images utiles à la carte (ou, avec l'autre critère, une image éligible sur $n$).

\subsection{Lancement du calcul}
\bouton{Build the map} lance le calcul sans autre question (la sélection est écrite dans le journal) :
\begin{enumerate}
	\item le logiciel vérifie que Docker Desktop est disponible. Sinon un message invite à le démarrer puis à choisir \texttt{Retry}. ODM est ensuite téléchargé s'il n'est pas encore sur l'ordinateur, ou mis à jour si on l'a demandé (voir plus haut) ;
	\item la liste des images retenues est enregistrée dans \texttt{mapping\_MULTI\textbackslash img\_best\_mappping.json} ;
	\item la résolution de l'orthophoto est estimée automatiquement : c'est la taille au sol d'un pixel (GSD, en cm) obtenue à partir de l'altitude de vol au-dessus du sol et de la focale de la caméra, médiane sur les images retenues, arrondie au dixième de cm supérieur ;
	\item les images TIFF retenues sont rendues accessibles à ODM dans \texttt{mapping\_MULTI\textbackslash images} par des \emph{liens physiques} (hard links) : aucune image n'est copiée et l'espace disque n'est donc pas doublé. Si le disque ne gère pas les liens physiques (disque externe en exFAT ou FAT32, certains partages réseau), les images sont copiées, ce qui utilise davantage d'espace disque ; le journal l'indique ;
	\item ODM est lancé. Le journal de la fenêtre affiche les messages principaux ; la totalité de la sortie est enregistrée dans un fichier \texttt{mapping\_MULTI\textbackslash odm\_run\_<date>.log} ;
	\item à la fin, l'orthophoto est contrôlée (4 bandes, géoréférencée), les résultats utiles sont archivés dans \texttt{mapping\_result} et les dossiers de travail d'ODM, très volumineux, sont supprimés.
\end{enumerate}
Pendant le calcul la barre de progression est indéterminée, les boutons et la fermeture de la fenêtre sont inactifs. À l'issue du calcul la durée est indiquée dans le journal.

\textbf{Résultats}\label{sec:step6_mapping_result}

Les résultats sont dans \texttt{mapping\_MULTI\textbackslash mapping\_result} :
\begin{itemize}
	\item \texttt{odm\_orthophoto.tif} : l'orthophoto à 4 bandes (rouge, vert, bleu, proche infrarouge), en UTM ;
	\item \texttt{report.pdf} : rapport de qualité produit par ODM ;
	\item \texttt{odm\_25dmesh.ply} : maillage 2,5D du terrain ;
	\item \texttt{shots.geojson} (et \texttt{shots.qml}) : position estimée de chaque prise de vue, avec son numéro de shoot point, à ouvrir dans QGIS ;  \texttt{shots.qml} est un fichier de mise en forme.
	\item des fichiers techniques (\texttt{ortho.png}, informations de géoréférencement).
\end{itemize}
Un aperçu en couleurs (bandes rouge, vert, bleu) de l'orthophoto s'affiche ensuite dans la fenêtre, puis le calcul du MNT et des courbes de niveau démarre automatiquement (c'est équivalent au bouton \bouton{Terrain model and contours\ldots}, voir \ref{btn:step6_dtm}).

\figureIRD{step5_orthophoto_preview.png}{Aperçu de l'orthophoto dans la fenêtre du Step~6.}{0.6}

\subsection{Précautions}
\begin{itemize}
	\item Les résultats existants ne sont \textbf{jamais écrasés}. Quand une carte est déjà construite, elle est conservée : son dossier est renommé \texttt{mapping\_MULTI\_<date de son exécution>} et la nouvelle carte est construite dans un nouveau \texttt{mapping\_MULTI} (avec la même configuration de caméra) ; la fenêtre de choix des images l'indique. Si la nouvelle exécution échoue, la carte précédente est remise dans \texttt{mapping\_MULTI} et l'exécution échouée, avec son journal, est conservée dans \texttt{mapping\_MULTI\_failed\_<date>}. L'étape~7 conserve ces dossiers.
	\item Docker doit utiliser le moteur \og desktop-linux \fg~de Docker Desktop. Le premier lancement télécharge l'image ODM (plusieurs gigaoctets) : prévoir une connexion Internet.
	\item Si certaines images retenues sont incomplètes (quadruplet TIFF manquant ou de tailles différentes), le calcul est refusé et le journal indique les fichiers en cause.
\end{itemize}

\section{Bouton \bouton{Terrain model and contours\ldots}}\label{btn:step6_dtm}
Ce bouton construit un modèle numérique de terrain (MNT) et des courbes de niveau sur l'emprise de l'orthophoto. Il est lancé \emph{automatiquement} à la fin d'un mapping réussi (avec un pas de grille de 10\,m) ; ce bouton permet de le refaire, par exemple avec un autre pas. Il nécessite l'orthophoto \texttt{mapping\_MULTI\textbackslash mapping\_result\textbackslash odm\_orthophoto.tif}.

Une fenêtre demande la valeur \texttt{Altitude grid step}, c'est-à-dire l'espacement (en mètres) des points de la grille régulière dont les altitudes sont demandées à l'IGN (service d'altitudes de l'Institut national de l'information géographique et forestière) ; en cas d'indisponibilité, le service OpenTopoData est utilisé avec une moindre précision en France. Plus le pas est petit, plus le nombre de requêtes est grand. Le pas par défaut est de 10\,m (plage 1 à 100\,m). Le nombre de points est limité à 5000 : au-delà, le logiciel augmente automatiquement le pas. \bouton{Build} lance le calcul, \bouton{Cancel} ne lance rien.

\subsection*{Étapes du calcul}
\begin{enumerate}
	\item lecture de l'emprise de l'orthophoto (UTM) ;
	\item construction d'une grille régulière qui couvre l'emprise ;
	\item requêtes d'altitude (aucun point manquant n'est toléré : sinon le calcul échoue et un message l'indique) ;
	\item interpolation par une surface lisse (spline quadratique) puis échantillonnage sur un raster à 1\,m ;
	\item export du MNT, des courbes de niveau (équidistance de 1\,m, courbes maîtresses tous les 10\,m), d'une carte 2D ombrée et d'une vue 3D.
\end{enumerate}
En cas d'échec, l'orthophoto n'est pas affectée.

\subsection{Fichiers produits}
Les fichiers sont créés à côté de l'orthophoto, dans \texttt{mapping\_result} (jamais écrasés, un suffixe numérique est ajouté si besoin) :
\begin{itemize}
	\item \texttt{ortho\_dtm.tif} : MNT (altitude en mètres, UTM) ;
	\item \texttt{ortho\_contours.geojson} et \texttt{ortho\_contours.qml} : courbes de niveau et leur style pour QGIS ;
	\item \texttt{ortho\_dtm\_map.png} : carte 2D ombrée avec courbes de niveau, position des points de la grille (point noir) et quadrillage UTM (bleu clair);
	\item \texttt{ortho\_dtm\_3d.png} : surface 3D ;
	\item \texttt{ortho\_dtm\_grid.json} : points de la grille avec leur altitude et la source des altitudes.
\end{itemize}

\figurePaireIRD{fig:step6_topo}
{step5_dtm_map.png}{Carte 2D du MNT avec courbes de niveau et quadrillage UTM.}
{step5_dtm_3d.png}{Surface 3D du MNT.}
{Représentation topographique du terrain à partir du modèle numérique de terrain déduit des altimétries IGN.}

\subsection{Ouverture dans QGIS}
Dans la fenêtre de la carte QGIS, faire glisser l'orthophoto \texttt{odm\_orthophoto.tif}, les courbes de niveau \texttt{ortho\_contours.geojson} et les points des prises de vue \texttt{shots.geojson}    (les styles \texttt{.qml} homonymes sont appliqués automatiquement). \\

Quelques actions simples avec une orthophotographie à 4 bandes spectrales red,green,blue,nir

\begin{itemize}
	\item Image couleur dans le spectre visible : clic droit sur la couche> \underline{P}ropriétés>Symbologie puis associer les bandes 1, 2, 3 au rouge, vert, bleu pour une vue en couleurs naturelles. C'est l'option d'affichage par défaut d'une image multispectrale.
	\item Image fausse couleur VIR : clic droit>\underline{P}ropriétés>Symbologie puis associer les bandes 4~,1~,2 au rouge, vert, bleu. 
	\item Carte de l'indice NDVI\footnote{L'indice NDVI pour \textit{Normalized Difference Vegetation Index} mesure la vigueur et la densité de la végétation. On détermine cet indice par la formule simple : $NDVI = \frac{NIR - RED}{NIR + RED}$} : \underline{R}aster>Calculatrice raster> entrez la formule 
	\begin{equation}
		\frac{(VIS@4 - VIS@1)}{(VIS@4 + VIS@1)}
	\end{equation}
	puis double clic sur la couche et choisir dans symbologie type de rendu : Pseudo-couleur à bande unique> Palette de couleur: NDVI. Le résultat est montré sur la figure \ref{fig:step6_mapping_NDVI}
	\item Couche NIR. Sélectionner dans symbologie comme première bande NIR,  puis dans symbologie type de rendu :  Niveau de gris à bande unique.
\end{itemize}

\figurePaireIRD{fig:step6_mapping_1}
{step5_topo_QGIS.png}{Représentation fine de la topographie. Équidistance des isolignes : $10\,m$.}
{step5_mapping_VIS.png}{Mapping en spectre visible.}
{Représentation des données IRdrone avec le logiciel QGIS sur fond OpenStreetMap.}

\figurePaireIRD{fig:step6_mapping_NDVI}
{step5_mapping_VIR.png}{Image multispectrale VIR.  $\qquad\qquad\qquad\quad$ Décalage spectral $[Nir|R|G] \longrightarrow [R|G|B]$.}
{step5_mapping_NDVI.png}{Carte de l'indice NDVI.}
{Représentation des données IRdrone avec le logiciel QGIS sur fond Google Satellite.}


\section{Bouton \bouton{Close}}\label{btn:step6_close}
Ferme la fenêtre et rend la main au dialogue principal. La fermeture est refusée tant qu'un calcul ODM ou MNT est en cours. Le bouton \bouton{?} ouvre cette aide à la page de ce chapitre.

\section{Bouton \bouton{Continue with archiving}}\label{btn:step6_continue}
Apparaît en bas à droite, après \bouton{Close}, une fois la carte construite. Il ferme la fenêtre et ouvre l'étape~7, qui indique l'espace libéré par l'archivage. Rien n'est supprimé avant votre confirmation.

%==========================================================
\chapter{Step 7 -- Archivage de la mission}\label{ch:step7}
L'étape~7 est l'archivage de la mission. Elle se fait \emph{une fois l'exploitation des images terminée}. Elle nettoie le dossier de la mission \emph{sur place} :
les nombreux fichiers techniques (images intermédiaires, dossiers de travail, copies d'images de la synchronisation, etc.) sont supprimés,
et seuls les éléments essentiels pour pouvoir, si nécessaire, traiter de nouveau la mission sont conservés. Sur les missions testées, le volume est réduit
d'un facteur 2 à 5 (par exemple de 131\,Go à 25\,Go pour une mission de 600 images VIS DJI).

\begin{center}
\fbox{\parbox{0.9\linewidth}{\small\textbf{Attention : opération irréversible.} Les fichiers supprimés ne vont pas à la corbeille de Windows ; ils sont définitivement perdus.
Si vous souhaitez conserver une copie de la mission, d'un dossier \texttt{mapping\_MULTI} particulier ou de tout autre fichier, \emph{copiez-les ailleurs avant de lancer l'archivage}.
Tout élément du dossier de mission qui ne figure pas dans la liste ci-dessous est supprimé, y compris les copies de dossiers (par exemple \texttt{mapping\_MULTI - Copie}).}}
\end{center}

Comme pour les autres étapes, la mission archivée est la mission ouverte dans la fenêtre principale.

\section{Ce qui est conservé}\label{btn:step7_rules}
La comparaison des noms de fichiers ne tient pas compte des majuscules et minuscules (\texttt{.dng} et \texttt{.DNG} sont équivalents).
\begin{description}
  \item[Racine de la mission] les dossiers \texttt{AerialPhotography}, \texttt{cameras}, \texttt{FlightAnalytics}, \texttt{ImgIRdrone}, \texttt{mapping\_MULTI} (et les cartes précédentes \texttt{mapping\_MULTI\_<date>} mises de côté par l'étape~6), \texttt{Synchro}
        et le fichier \texttt{desktop.ini}. Tout le reste est supprimé.
  \item[\texttt{AerialPhotography}] les dossiers \texttt{VIS} et \texttt{NIR} et le fichier \texttt{shoot\_points.json}. Dans \texttt{VIS} et \texttt{NIR} : les images originales
        \texttt{.dng} avec leurs fichiers compagnons \texttt{.exif}, ainsi que \texttt{time\_line.json} et \texttt{transfer\_info\_VIS\_dng.json} (respectivement \texttt{transfer\_info\_NIR\_dng.json}).
        Les autres fichiers (notamment les \texttt{*\_RawTherapee.tif}) sont supprimés.
  \item[\texttt{cameras} et \texttt{FlightAnalytics}] laissés tels quels.
  \item[\texttt{ImgIRdrone}] les images \texttt{VIS\_XXXX.jpg} et \texttt{VIR\_XXXX\_local.jpg}, les fichiers \texttt{VIS\_XXXX\_motion\_model.npy}, les images
        \texttt{NDVI\_XXXX\_local.jpg} et \texttt{NIR\_XXXX\_local.jpg} (si elles existent), \texttt{shoot\_points.json} et \texttt{shoot\_points\_failed\_summary.json}.
        Le dossier volumineux \texttt{imgMultispectral} et les dossiers de \og redo \fg{} individuels sont supprimés.
  \item[\texttt{mapping\_MULTI}] laissé tel quel : il contient les éléments directement exploitables dans QGIS et a déjà été allégé à l'étape~6.
  \item[\texttt{Synchro}] seuls \texttt{check\_time\_alignment.png}, \texttt{check\_time\_before\_alignment.png} et \texttt{desktop.ini} sont conservés ; les dossiers \texttt{VIS}, \texttt{NIR} et \texttt{Check\_ARUCO} sont supprimés.
\end{description}

\paragraph{Pourquoi ce choix ?} Le décalage temporel trouvé à l'étape~3 et la liste des paires d'images (\texttt{AerialPhotography\textbackslash shoot\_points.json}) résument la synchronisation :
les images de \texttt{Synchro} ne sont plus utiles. Les fichiers \texttt{VIS\_XXXX\_motion\_model.npy} mémorisent l'alignement VIS--NIR de chaque \emph{shoot point}, y compris les alignements
rattrapés par la procédure de \og redo \fg{} (étape~4). Pour reprendre la mission, il suffit donc de refaire l'étape~4, sans perdre ces corrections, à partir des images originales conservées.
Les images \texttt{.jpg} conservées permettent aussi de retrouver rapidement une image pour illustrer un rapport.

\figureIRD[fig:step7]{step6_window.png}{Fenêtre de l'archivage : l'espace libéré par l'archivage s'affiche dès l'ouverture de la fenêtre.}{0.6}

\section{Ce que l'archivage libère (simulation)}\label{btn:step7_simulate}
Dès l'ouverture de la fenêtre, le dossier est analysé \emph{sans rien supprimer} (une simulation). Le cadre sous l'avertissement donne la réponse :
\begin{itemize}
  \item \texttt{12.3 GB can be freed}, avec, en grand, la taille du dossier de mission avant et après l'archivage (par exemple \texttt{15.1 GB → 2.8 GB}) et le nombre de fichiers qui seront supprimés. Le détail, dossier par dossier (fichiers à supprimer, fichiers conservés, dossiers supprimés), est dans la zone en dessous ;
  \item \texttt{Archiving impossible}, avec les fichiers indispensables qui manquent (voir plus bas) : rien ne peut être supprimé ;
  \item \texttt{Mission already archived}, avec la date de l'archivage et, en grand, la taille du dossier ;
  \item \texttt{Nothing to delete}, quand le dossier ne contient que ce que l'archivage conserve.
\end{itemize}
Vérifiez que la liste correspond à ce que vous attendez avant d'archiver.

\section{Bouton \bouton{Archive and free \ldots}}\label{btn:step7_archive}
Le bouton indique l'espace qui sera libéré (par exemple \bouton{Archive and free 12.3 GB\ldots}). Il est grisé pendant l'analyse du dossier, et quand rien ne peut être supprimé. Le dossier est analysé à nouveau au clic, puis un message de confirmation est demandé à l'utilisateur  comme le montre la figure \ref{fig:step7_confirm_archving}. Le message indique aussi l'espace disque qui sera libéré et précise que la suppression sera définitive. Deux boutons \texttt{OK} et \texttt{Cancel} permettent de confirmer ou d'abandonner. Par sécurité et pour éviter un clic trop rapide sur \texttt{Ok} l'utilisateur  doit recopier exactement le même nom du dossier de mission que celui indiqué (en respectant la casse). 

\figureIRD[fig:step7_confirm_archving]{step6_confirm_archiving.png}{Message de confirmation de la demande d'archivage.}{0.3}

 Si clic OK  alors IRDrone lance l'archivage réel :

\begin{enumerate}
  \item analyse le dossier et vérifie que les fichiers indispensables existent (\texttt{mission\_parameters.json}, \texttt{shoot\_points.json}, \texttt{time\_line.json}, \texttt{transfer\_info\_*\_dng.json}, etc.).
        S'il en manque, l'archivage est \textbf{refusé} et la liste des fichiers manquants est affichée ; rien n'est supprimé. Les dossiers manquants ne sont jamais recréés ;
  \item affiche le nombre de fichiers et le volume qui seront supprimés, puis demande confirmation : il faut \emph{saisir le nom du dossier de la mission} ;
  \item supprime les fichiers. La barre de progression indique l'avancement ;
  \item si un fichier est verrouillé (ouvert dans un autre logiciel, par exemple QGIS), l'erreur est notée et la suppression continue. Fermez alors les logiciels concernés et relancez l'archivage ;
  \item écrit à la racine de la mission les fichiers \texttt{archive\_manifest.json} et \texttt{archive\_log.txt}.
\end{enumerate}

Une mission déjà archivée est détectée et n'est pas archivée une seconde fois.


 Une icône particulière attribuée au dossiers archivés et permet de voir d'un coup d'\oe il qu'une mission est archivée.(figure \ref{fig:Air_Mission_folder}).


\paragraph{Archivage et avancement du traitement.} A la fin de chaque étape le dossier de mission reflète l'état de l'avancement du traitement de la mission. 

L'archivage ne doit être lancé qu'à la toute fin de l'étude. Notez que les images originales de la mission (converties au format dng pour la SJCam) sont systématiquement archivée. Vous conservez ainsi les données du terrain. Il n'est donc pas nécessaire d'archiver les images originales des cartes SD du DJI et de la SJCam (mais vous pouvez aussi les archiver par sécurité).

\section{Bouton \bouton{Close}}\label{btn:step7_close}
Ferme la fenêtre et rend la main au dialogue principal. Le bouton \bouton{?} ouvre cette aide à la page de ce chapitre.

%==========================================================
\chapter{Missions visible seul}\label{ch:visible_only}
\section{Sans la caméra infrarouge}\label{sec:visible_only}
Une mission peut être volée avec la seule caméra du drone, sans la SJCam : le type \bouton{Visible only (drone DNG)} se choisit en haut de la première page de l'étape~1 (section~\ref{btn:step1_mission_type}). Il est enregistré dans la mission (\texttt{"mission type": "visible"} dans \texttt{FlightAnalytics/mission\_parameters.json} ; les missions sans cette clé sont multispectrales). Les étapes sont les mêmes, moins ce qui concerne l'infrarouge :
\begin{itemize}
	\item \textbf{Étape 1} : seule la carte SD du drone est nécessaire ; la copie a trois étapes (copie, EXIF, altitude du sol). Le dossier de la mission n'a pas de dossier \texttt{AerialPhotography/NIR} (\texttt{Synchro} ne contient que les vignettes de la ligne de temps de l'étape~2).
	\item \textbf{Étape 2} : la ligne de temps ne montre que les images visibles ; il n'y a pas de cible tournante à chercher ni de phase de synchronisation : seule la phase de relevé est vérifiée. Le best mapping fonctionne comme pour une mission multispectrale. \bouton{Continue with development} enregistre le relevé et les points de prise de vue visibles (\texttt{AerialPhotography/shoot\_points.json}, un point par image visible, sans infrarouge) et ouvre l'étape~4.
	\item \textbf{Étape 3} : inutile ; sa tuile indique \og Not needed \fg~et \bouton{Continue} passe directement de l'étape~2 à l'étape~4.
	\item \textbf{Étape 4} (\og Develop \fg) développe les images DNG visibles du relevé (ou de la sélection du best mapping) : couleurs, correction de la distorsion de l'objectif et même recadrage que les images multispectrales, pour que le modèle de caméra d'OpenDroneMap soit le même. Chaque image donne \texttt{ImgIRdrone/VIS\_XXXX.jpg} (revue) et les trois bandes \texttt{ImgIRdrone/imgMultispectral/IRD\_XXXX\_1..3.tif} (rouge, vert, bleu) pour la carte.
	\item \textbf{Étape 5} : la revue montre les images visibles développées ; il n'y a rien à réaligner. \bouton{Good: keep for the map} (touche \texttt{Y}) marque l'image vérifiée et montre la suivante à vérifier ; une image floue ou mal exposée est signalée (touche \texttt{F}) : elle est écartée de la carte. Les points que l'étape~4 n'a pas développés (gris sur la carte) se développent d'un clic droit sur la carte.
	\item \textbf{Étape 6} : OpenDroneMap construit une orthophoto à trois bandes (rouge, vert, bleu) au lieu de quatre. Il n'y a ni couche en fausses couleurs infrarouges, ni VIR, ni NDVI ; le modèle de terrain est calculé comme d'habitude.
	\item \textbf{Étape 7} : l'archive garde les mêmes fichiers, sans ceux de l'infrarouge.
\end{itemize}

%==========================================================
\appendix
\chapter{Alignement des images VIS et NIR (Step 4)}\label{ann:alignement}\label{sec:ann_algorithme}
Cette annexe explique en détail comment IRdrone superpose l'image proche infrarouge à l'image visible. Elle s'adresse à un lecteur de formation scientifique (algèbre linéaire, un peu de traitement du signal) qui n'est pas spécialiste de vision par ordinateur. Les premières sections donnent les bases de géométrie et de traitement du signal, les suivantes décrivent chaque étape de l'algorithme avec les traces de diagnostic d'alignements réels, puis viennent les méthodes essayées puis abandonnées, l'historique du développement et les limites connues. La dernière section (§\ref{sec:ann_theorique}) traite de l'origine physique du désalignement des deux caméras (angles de \og pitch \fg{} et de \og yaw \fg{} théoriques).

Le code se trouve dans \texttt{automatic\_registration.py} (fonctions \texttt{align\_raw} et \texttt{coarse\_alignment}), dans le dossier \texttt{registration} (\texttt{rigid.py} : pyramide, \texttt{cost.py} : représentation et coût, \texttt{warp\_flow.py} : rééchantillonnage) et dans \texttt{irdrone/semi\_auto\_registration.py} (rotation de la caméra, \texttt{manual\_warp}). Les figures de cette annexe sont produites par \texttt{help/make\_alignment\_figures.py} à partir de deux shoot points :
\begin{itemize}
	\item \textbf{point A} : l'image \texttt{VIS\_0201} du vol du 25 janvier 2022 (Blassac, 90\,m au-dessus du sol, écart de temps entre les deux prises de vue 0,0\,s), un cas typique où l'alignement réussit ;
	\item \textbf{point B} : la paire d'exemple du dépôt (\texttt{samples/\_visible.DNG}, \texttt{samples/\_nir.RAW}, mai 2022), un cas difficile où l'alignement automatique échoue (§\ref{sec:ann_limits}).
\end{itemize}
Les deux ont été alignés en ligne de commande avec l'option \texttt{--debug} (§\ref{sec:ann_traces}), c'est-à-dire avec l'élargissement de champ 1,6 de la ligne de commande au lieu de 1,4 au Step~4.

\section{Le problème}\label{sec:ann_problem}
\subsection{Ce que signifie \og aligner \fg}
Le drone emporte deux caméras (figure~\ref{fig:ann_pair}) : la caméra du DJI Mavic Air~2, qui enregistre la lumière visible (VIS, rouge, vert et bleu), et une SJCam M20 dont le filtre anti-infrarouge a été retiré et remplacé par un filtre qui ne laisse passer que le proche infrarouge (NIR, au-delà d'environ 720\,nm). Les indices de végétation comme le NDVI combinent, \emph{pixel par pixel}, le rouge de l'image visible et le proche infrarouge :
\begin{equation}
	\text{NDVI} = \frac{\text{NIR} - \text{R}}{\text{NIR} + \text{R}}.
\end{equation}
Cette formule n'a de sens que si les deux valeurs d'un pixel proviennent du même point du sol. Aligner (on dit aussi \emph{recaler}) les deux images, c'est trouver, pour chaque pixel de l'image visible, la position du même point du sol dans l'image infrarouge, puis rééchantillonner l'image infrarouge sur la grille de l'image visible. Une erreur de quelques pixels suffit à créer de fausses valeurs d'indice le long de chaque bord contrasté (une haie, une route, l'ombre d'un arbre).

\figureIRD[fig:ann_pair]{align_pair_raw.jpg}{Le point B tel que l'enregistrent les deux caméras. Outre des champs de vision différents et la forte distorsion de la caméra infrarouge, les contrastes diffèrent : la route est sombre dans les deux bandes mais ses bas-côtés ne sont clairs que dans le visible, la végétation est verte et plus sombre que l'herbe sèche dans le visible, claire dans le proche infrarouge.}{1}

\subsection{Pourquoi c'est difficile}
Trois difficultés s'additionnent :
\begin{enumerate}
	\item \textbf{Les deux caméras ne regardent pas dans la même direction.} La caméra visible est maintenue verticale par sa nacelle (gimbal), la caméra infrarouge est fixée au corps du drone et suit ses inclinaisons (§\ref{sec:ann_theorique}). L'écart atteint plusieurs degrés, c'est-à-dire plusieurs centaines de pixels.
	\item \textbf{Les deux optiques sont très différentes.} La caméra infrarouge a un champ très large et une forte distorsion ; les lignes droites apparaissent courbes (§\ref{sec:ann_distortion}).
	\item \textbf{Les deux bandes ne voient pas la même chose.} La végétation, l'eau, le sol et les routes n'ont pas la même réflectance dans le visible et dans le proche infrarouge. Un contraste peut s'inverser d'une bande à l'autre (figure~\ref{fig:ann_pair}, et figure~\ref{fig:step4_pair_VIS_NIR} où la rivière est noire dans l'infrarouge). Les méthodes classiques, qui comparent des intensités ou apparient des points caractéristiques (SIFT, ORB), sont donc peu fiables.
\end{enumerate}
Les deux premières difficultés sont résolues par la \emph{géométrie} : un modèle de chaque caméra, obtenu par calibration, et le fait que les deux images diffèrent essentiellement d'une rotation. La troisième est résolue en \emph{comparant des bords plutôt que des intensités}.

\section{Géométrie d'une caméra}\label{sec:ann_camera}
Le livre de référence est \emph{Multiple View Geometry in Computer Vision} de Hartley et Zisserman~[\hyperref[ann_hz]{1}] (chapitres 2 et 6 pour le modèle de caméra, 7.4 pour la distorsion, 8.4 pour les caméras en rotation). Le livre de Szeliski~[\hyperref[ann_szeliski]{2}] est une introduction plus accessible (chapitre 2).

\subsection{Le modèle sténopé et les coordonnées homogènes}
En première approximation, une caméra est un \emph{sténopé} (pinhole) : tout rayon lumineux passe par le centre optique $C$ avant d'atteindre le capteur. Soient $(X,Y,Z)$ les coordonnées d'un point de la scène dans le repère de la caméra ($Z$ selon l'axe optique). Son image est au pixel
\begin{equation}
	u = f\,\frac{X}{Z} + c_x, \qquad v = f\,\frac{Y}{Z} + c_y,
\end{equation}
où $f$ est la focale \emph{exprimée en pixels} (focale physique divisée par la taille d'un pixel) et $(c_x, c_y)$ le \emph{point principal}, intersection de l'axe optique avec le capteur, proche du centre de l'image. La division par $Z$ est la perspective : un objet deux fois plus loin apparaît deux fois plus petit.

Cette division n'est pas linéaire ; elle le devient en écrivant le pixel en \emph{coordonnées homogènes}, un vecteur de trois nombres défini à un facteur près : $(u, v, 1) \sim (\lambda u, \lambda v, \lambda)$. Alors
\begin{equation}\label{eq:ann_pinhole}
	\lambda \begin{pmatrix} u \\ v \\ 1 \end{pmatrix} = K \begin{pmatrix} X \\ Y \\ Z \end{pmatrix},
	\qquad
	K = \begin{pmatrix} f & 0 & c_x \\ 0 & f & c_y \\ 0 & 0 & 1 \end{pmatrix}.
\end{equation}
La matrice $K$ des \emph{paramètres intrinsèques} (ou matrice de calibration) décrit la caméra seule. Inversement, $K^{-1}(u, v, 1)^T$ est la \emph{direction} du rayon qui arrive sur le pixel $(u,v)$ : une image est un ensemble de directions, la profondeur $Z$ est perdue.

\begin{table}[htbp]\centering
\begin{tabular}{@{}lll@{}}
\hline
 & \textbf{DJI Mavic Air~2 (VIS)} & \textbf{SJCam M20 (NIR)} \\ \hline
Taille de l'image (pixels) & $3992 \times 2992$ & $4600 \times 3448$ \\
Focale $f$ (pixels) & 2898 & 2311 \\
Point principal $(c_x, c_y)$ & (2032, 1511) & (2348, 1725) \\
Champ horizontal & $68^\circ$ & plus de $110^\circ$ \\
Distorsion $k_1, k_2, k_3$ & $0{,}185,\ -0{,}455,\ 0{,}245$ & $-0{,}337,\ 0{,}131,\ -0{,}025$ \\ \hline
\end{tabular}
\caption{Calibrations des deux caméras (\texttt{calibration/DJI\_RAW} et \texttt{calibration/M20\_RAW}, images RAW). À la focale de la caméra visible, une rotation de $1^\circ$ déplace l'image de $f \tan 1^\circ \approx 51$ pixels.}
\label{tab:ann_calib}
\end{table}

\subsection{La distorsion}\label{sec:ann_distortion}
Les objectifs réels s'écartent du modèle sténopé : l'image d'une droite est courbe, de plus en plus vers les bords. IRdrone utilise le modèle polynomial de Brown et Conrady~[\hyperref[ann_brown]{3}], celui d'OpenCV. Avec les coordonnées normalisées $x = X/Z$, $y = Y/Z$ et $r^2 = x^2 + y^2$, le point idéal $(x, y)$ est observé en
\begin{equation}\label{eq:ann_distortion}
	x_d = x\,(1 + k_1 r^2 + k_2 r^4 + k_3 r^6) + 2p_1 xy + p_2 (r^2 + 2x^2),
\end{equation}
et de même pour $y_d$ ; le pixel est alors $K\,(x_d, y_d, 1)^T$. Les coefficients $k_i$ décrivent la distorsion radiale (en barillet si $k_1 < 0$, en coussinet si $k_1 > 0$), $p_1, p_2$ un léger décentrement des lentilles.

La figure~\ref{fig:ann_distortion} donne les ordres de grandeur. Pour la caméra visible la correction ne dépasse pas 40 pixels. Pour la caméra infrarouge elle atteint 210 pixels aux bords gauche et droit et 375 pixels aux coins du champ de la caméra visible, et 800 pixels à $50^\circ$ de l'axe. Son $k_1$ fortement négatif est une distorsion en barillet proche de celle d'un \og fish-eye \fg. Le polynôme n'est valable que dans la zone où des mires de calibration ont été vues : au-delà d'environ $57^\circ$ de l'axe il redécroît, le modèle n'est donc plus inversible, et les bords de l'image infrarouge (plus de $110^\circ$ de champ horizontal) sont hors de son domaine de validité. Ces zones ne sont jamais utilisées, puisque le champ visible n'atteint que $41^\circ$ dans les coins.

\figureIRD[fig:ann_distortion]{align_distortion.png}{Distorsion des deux caméras. À gauche et au centre : un quadrillage de droites au sol, tel qu'il est enregistré sur le capteur (cadre noir) de chaque caméra ; en rouge, le champ de la caméra visible vu par la caméra infrarouge. À droite : décalage radial d'un point en fonction de son angle avec l'axe optique.}{1}

\subsection{La calibration}
Les paramètres intrinsèques $f, c_x, c_y$ et les coefficients de distorsion sont mesurés une fois pour chaque caméra par la méthode de Zhang~[\hyperref[ann_zhang]{4}] : une trentaine de photographies d'un damier plan dans diverses orientations ; les coins des cases sont détectés, et tous les paramètres (plus la position du damier sur chaque photo) sont estimés en minimisant la distance entre les coins observés et les coins prédits par les équations \eqref{eq:ann_pinhole} et \eqref{eq:ann_distortion} (\texttt{calibration/cameracalibration.py}, OpenCV). Cette calibration est fixe : elle ne dépend pas du vol.

\section{Deux caméras qui tournent : l'homographie}\label{sec:ann_homography}
\subsection{Une rotation de la caméra est une homographie de l'image}
Supposons que les deux caméras aient le \emph{même centre optique} et ne diffèrent que d'une rotation $R$ (une matrice orthogonale $3 \times 3$). Le rayon qui arrive sur le pixel $p_N$ de la caméra infrarouge a la direction $K_{NIR}^{-1} p_N$ dans le repère de cette caméra, $R\, K_{NIR}^{-1} p_N$ dans celui de la caméra visible, et arrive donc sur le pixel
\begin{equation}\label{eq:ann_rotation_homography}
	p_V \sim K_{VIS}\, R\, K_{NIR}^{-1}\, p_N = H\, p_N
\end{equation}
de la caméra visible (après correction des distorsions). La matrice $3 \times 3$ $H$, définie à un facteur près, est une \emph{homographie} : la transformation la plus générale qui envoie les droites sur des droites (Hartley et Zisserman~[\hyperref[ann_hz]{1}], §2.3 et §8.4). Elle a 8 degrés de liberté. Point remarquable : l'équation \eqref{eq:ann_rotation_homography} ne fait pas intervenir la profondeur $Z$ : quel que soit le relief, une rotation pure de la caméra est une homographie de l'image. C'est le principe des panoramas réalisés en tournant sur soi-même.

La rotation est paramétrée par trois angles : le \emph{lacet} (yaw, rotation autour de l'axe vertical de l'image, qui la déplace horizontalement), le \emph{tangage} (pitch, qui la déplace verticalement) et le \emph{roulis} (roll, rotation autour de l'axe optique, qui fait tourner l'image autour de son centre). Dans le code, $R$ est construite à partir du vecteur de rotation $(-\text{pitch}, \text{yaw}, \text{roll})$ (formule de Rodrigues, \texttt{manual\_warp}). Aux petits angles, un lacet $\theta$ déplace le centre de l'image de $f \tan\theta$ pixels, soit environ 51 pixels par degré ; un roulis $\varphi$ déplace un coin, à environ 2500 pixels du centre, de $2500\,\varphi$ (44 pixels par degré).

\subsection{Pourquoi la rotation suffit (presque) : la parallaxe}
Les deux centres optiques ne sont pas exactement au même endroit : ils sont à 46\,mm l'un de l'autre sur le drone (\texttt{CNIRCVIS\_0} dans \texttt{config.py}), et le drone se déplace entre les deux prises de vue. Une distance $B$ entre les centres (la \emph{base}) crée une \emph{parallaxe} : un point à la distance $Z$ est décalé de
\begin{equation}
	d = f\,\frac{B}{Z} \quad \text{pixels}.
\end{equation}
Seules les \emph{variations} de ce décalage comptent, sa partie constante étant absorbée par la rotation. Avec $B = 46$\,mm, à 90\,m : $d = 2898 \times 0{,}046 / 90 \approx 1{,}5$ pixel, et un relief de 10\,m le modifie de 0,2 pixel : négligeable. Si les deux prises de vue sont séparées de $\Delta t = 0{,}5$\,s et que le drone vole à 3\,m/s, $B = 1{,}5$\,m : $d \approx 48$ pixels, et un relief de 10\,m (une haie d'arbres, une pente) le modifie maintenant de $f B \Delta Z / Z^2 \approx 5$ pixels. C'est la raison pour laquelle le Step~4 n'aligne que les shoot points dont l'écart de temps est faible, et pour laquelle une correction \emph{locale} s'ajoute au modèle global (§\ref{sec:ann_local}). À très basse altitude (décollage, atterrissage), la parallaxe domine et l'alignement n'a pas de sens.

\subsection{La caméra virtuelle}
L'image visible est d'abord corrigée de sa distorsion : elle devient l'image d'une caméra sténopé idéale, la \emph{référence}. L'image infrarouge est ensuite rééchantillonnée directement dans cette caméra de référence, par l'équation \eqref{eq:ann_rotation_homography} et son propre modèle de distorsion, mais sur un canevas plus large que l'image visible (1,4 fois par défaut, 1,6 dans le script REDO) : c'est la \emph{caméra virtuelle} de la figure~\ref{fig:ann_virtual}. Si la rotation était l'identité, l'image visible correspondrait exactement au rectangle central ; chercher la rotation revient à trouver où elle se trouve réellement.

\figureIRD[fig:ann_virtual]{align_virtual_camera.jpg}{Point B. À gauche : l'image infrarouge corrigée de sa distorsion et rééchantillonnée dans la caméra virtuelle (focale de la caméra visible, champ 1,4 fois plus large). En pointillés : position de l'image visible si les deux caméras étaient parallèles ; en rouge : position trouvée par la recherche grossière (lacet $8{,}2^\circ$, tangage $1{,}9^\circ$). À droite : la référence visible. En dehors du champ visible la correction de la distorsion est moins précise (la route s'incurve à gauche).}{1}

\section{Comparer deux bandes spectrales}\label{sec:ann_similarity}
Pour trouver la rotation, il faut un \emph{critère de similarité} : un nombre qui dit si deux images sont bien superposées, et qui est optimal pour la bonne superposition.

\subsection{Pourquoi on ne peut pas comparer les intensités}
Pour deux images d'une même bande, le critère le plus simple est la somme des carrés des différences, $\text{SSD} = \sum_p (I_1(p) - I_2(p))^2$, ou la corrélation. Tous deux supposent qu'un même objet a la même intensité (ou une intensité proportionnelle) dans les deux images. Entre le visible et le proche infrarouge c'est faux : la végétation est sombre dans le visible et claire dans l'infrarouge, les bas-côtés de la route ne sont clairs que dans le visible (figures~\ref{fig:ann_pair} et \ref{fig:ann_features}). Une réponse classique à ces problèmes \og multimodaux \fg{} est l'information mutuelle~[\hyperref[ann_viola]{5}], qui mesure une dépendance statistique au lieu d'une égalité. Une autre approche combine des points caractéristiques et une similarité de gradients~[\hyperref[ann_vayssade]{18}]. IRdrone utilise une idée plus simple et plus locale : les \emph{bords} sont communs même quand les contrastes ne le sont pas.

\subsection{Les cartes de caractéristiques : énergies directionnelles du gradient}\label{sec:ann_msr}
Chaque image est convertie en luminance (niveaux de gris), lissée par un filtre gaussien d'écart-type $\sigma$ (5 pixels pour le visible, 3 pour l'infrarouge, moins net ; 3 et 1 pour la recherche grossière), puis remplacée par quatre cartes (\texttt{multispectral\_representation}, mode \texttt{LAPLACIAN\_ENERGIES}) :
\begin{equation}
	E_x = (\partial_x I)^2, \quad E_y = (\partial_y I)^2, \quad E_{\nearrow} = (\partial_{\nearrow} I)^2, \quad E_{\searrow} = (\partial_{\searrow} I)^2.
\end{equation}
Les dérivées horizontale et verticale sont calculées avec des filtres de Scharr~[\hyperref[ann_scharr]{6}] (une dérivée $3 \times 3$ avec un lissage optimisé, plus isotrope que Sobel), les diagonales avec la croix de Roberts $2 \times 2$~[\hyperref[ann_roberts]{7}]. Ces cartes sont appelées \emph{représentation multispectrale} (MSR) dans le code. L'idée vient d'Irani et Anandan~[\hyperref[ann_irani]{9}], qui alignaient des images de capteurs différents (visible et infrarouge entre autres) au moyen de telles énergies de dérivées directionnelles.

Le carré est essentiel : un bord du sombre vers le clair et un bord du clair vers le sombre donnent la même énergie. Une inversion de contraste entre les bandes devient donc invisible, tandis que la \emph{position} du bord est conservée. Les quatre directions gardent l'orientation des bords, ce qui aide à distinguer, par exemple, une route du bord perpendiculaire d'un champ. Le lissage supprime les textures fines (herbe, feuilles), différentes dans les deux bandes, et garde les grandes structures. La figure~\ref{fig:ann_features} montre ces cartes sur un détail du point B : les luminances sont très différentes, les énergies des bords de la route sont au même endroit.

\figureIRD[fig:ann_features]{align_feature_maps.jpg}{Énergies directionnelles du gradient sur un détail du point B ($600 \times 600$ pixels). Les luminances des deux bandes ne sont pas comparables (bas-côtés clairs dans le visible, sombres dans l'infrarouge), mais les bords de la route produisent les mêmes lignes dans les cartes d'énergie. Certains bords n'existent que dans une bande (les quatre lignes du visible contre deux dans l'infrarouge) : le critère doit les tolérer.}{1}

\subsection{Le critère de similarité : gradient total de la différence}\label{sec:ann_ntg}
Le critère dérive du \emph{Normalized Total Gradient} (NTG) de Chen et al.~[\hyperref[ann_ntg]{8}]. Pour deux images $f$ et $g$ :
\begin{equation}
	\text{NTG}(f, g) = \frac{\sum_p \|\nabla (f - g)(p)\|_1}{\sum_p \|\nabla f(p)\|_1 + \sum_p \|\nabla g(p)\|_1}.
\end{equation}
L'idée : quand deux images sont bien superposées, leurs bords communs s'annulent dans la différence $f-g$, qui devient lisse ; son gradient total (somme des valeurs absolues des dérivées) est alors minimal. Quand elles sont décalées, chaque bord apparaît deux fois dans la différence, avec des signes opposés, et le gradient total augmente. La norme $\|\cdot\|_1$ (valeurs absolues plutôt que carrés) rend le critère robuste aux bords qui n'existent que dans une image : ils ajoutent un coût à peu près constant, qui dépend peu du décalage, au lieu de le dominer. Le dénominateur rend le critère indépendant de la quantité de texture.

IRdrone l'applique aux cartes d'énergie : pour un décalage $u$ de la référence,
\begin{equation}\label{eq:ann_cost}
	C(u) = \sum_{c=1}^{4} \sum_p \Big( \big|S_x * (E^{NIR}_c - E^{VIS}_c(\cdot + u))\big|(p) + \big|S_y * (E^{NIR}_c - E^{VIS}_c(\cdot + u))\big|(p) \Big),
\end{equation}
où $S_x, S_y$ sont les filtres de dérivation de Sobel (\texttt{cost\_surface\_NTG} dans \texttt{cost.py}). Seul le numérateur du NTG est calculé. Le filtre de Sobel étant linéaire, $S * (a - b) = S*a - S*b$ : les gradients des deux images sont calculés une fois, et chaque décalage n'est plus qu'une somme de différences absolues (15 fois plus rapide, novembre 2026). Une somme des carrés des différences sur les mêmes cartes est aussi programmée (\texttt{SSD}) mais n'est pas utilisée.

La figure~\ref{fig:ann_profiles} compare les critères sur le point A, en fonction d'un décalage horizontal autour du bon alignement. La SSD des luminances a une vallée large : elle situe l'alignement, mais avec peu de précision, et elle suppose des intensités comparables, ce qui est le cas sur cette scène d'hiver mais pas en été. Le NTG a un puits étroit et profond : il est très précis, mais ne \og voit \fg{} la bonne position qu'à quelques pixels près, d'où le besoin d'un bon point de départ et d'une stratégie du grossier au fin.

\figureIRD[fig:ann_profiles]{align_cost_profiles.png}{Coût en fonction d'un décalage horizontal (point A, après alignement, calculé à la résolution 1/4), normalisé entre 0 et 1. (a) L'image visible contre elle-même. (b) Visible contre infrarouge sur un bloc avec des arbres et une haie. (c) Sur un bloc de prairie seule, avec peu de structures : le NTG a encore son minimum au bon endroit mais en dehors du puits étroit la courbe n'est plus informative.}{1}

\section{Recherche grossière}\label{sec:ann_coarse}
La recherche grossière trouve l'essentiel de l'écart d'orientation (fonction \texttt{coarse\_alignment}) :
\begin{enumerate}
	\item l'image infrarouge est rééchantillonnée dans la caméra virtuelle avec les angles initiaux : nuls au Step~4, ceux réglés à la main lors d'un réalignement (§\ref{btn:review_redo}) ;
	\item les cartes d'énergie des deux images sont réduites d'un facteur 32 ($125 \times 94$ pixels pour le visible), la carte visible est placée au centre d'un canevas de la taille de la carte infrarouge ;
	\item une \emph{recherche exhaustive} calcule le coût \eqref{eq:ann_cost} pour chaque translation de la carte visible à l'intérieur de la carte infrarouge : avec l'élargissement de champ 1,4 du Step~4, $\pm 25$ pixels à cette échelle, soit environ $\pm 800$ pixels en pleine résolution ou $\pm 15^\circ$ ; avec l'élargissement 1,6 de la ligne de commande, $\pm 37$ pixels, $\pm 1200$ pixels ou $\pm 22^\circ$ (figure~\ref{fig:ann_coarse}) ;
	\item la meilleure position est localisée avec une précision inférieure au pixel (ci-dessous) ;
	\item aux petits angles une translation horizontale $t_x$ est un lacet et une translation verticale $t_y$ un tangage : $\text{yaw} = \arctan(t_x/f)$, $\text{pitch} = \arctan(t_y/f)$ ;
	\item l'image infrarouge est rééchantillonnée une nouvelle fois, en pleine résolution, avec les angles corrigés.
\end{enumerate}
Le roulis n'est pas recherché : il garde sa valeur initiale (§\ref{sec:ann_fourier}).

\paragraph{Choix du minimum.} Le minimum absolu de la surface de coût n'est pas toujours le bon : une zone sans texture ou le bord noir du canevas peuvent produire une vallée large et basse. La bonne superposition produit un puits étroit et profond (figure~\ref{fig:ann_coarse}). Le code cherche donc la position où la surface est la plus fortement courbée : en chaque position, une surface quadratique est ajustée sur le voisinage $3\times3$, et la position du plus grand déterminant de sa matrice hessienne (produit des deux courbures principales) est retenue (\texttt{minimum\_cost\_max\_hessian}).

\paragraph{Raffinement sous-pixel.} Autour de la position retenue, le coût est approché par une forme quadratique~[\hyperref[ann_hasinoff]{10}]
\begin{equation}
	C(u) \approx \tfrac{1}{2}\, u^T A\, u + b^T u + c,
\end{equation}
dont les coefficients s'obtiennent par de petites convolutions du voisinage $3\times3$. Son minimum est en $u^* = -A^{-1} b$ (un pas de Newton), retenu s'il est à moins d'un pixel. À la résolution 1/32, un pixel vaut 32 pixels en pleine résolution : ce raffinement rend la recherche grossière précise à quelques pixels près. Pour le point A, la translation trouvée est de $(21{,}7 ; 163{,}9)$ pixels, soit un lacet de $0{,}43^\circ$ et un tangage de $3{,}24^\circ$.

\figureIRD[fig:ann_coarse]{align_coarse_search.png}{Recherche grossière du point A à la résolution 1/32. À gauche : la surface de coût sur toute la plage de recherche ($\pm 1200$ pixels) : un puits net dans un paysage lisse. Au centre : zoom sur le centre du coût. À droite : déterminant de la hessienne, maximal au fond du puits.}{1}

\section{Raffinement pyramidal}\label{sec:ann_pyramid}
Après la recherche grossière, le désalignement résiduel va de quelques pixels à quelques dizaines de pixels, et il n'est pas uniforme : roulis résiduel, erreurs des calibrations, parallaxe. Le raffinement l'estime par blocs, sur une pyramide d'images (\texttt{pyramidal\_search} dans \texttt{rigid.py}).

\subsection{Du grossier au fin}
Une \emph{pyramide gaussienne}~[\hyperref[ann_burt]{11}] est la suite d'une image lissée et réduite par 2, 4, 8, 16\ldots{} Un déplacement de 100 pixels en pleine résolution fait 6 pixels à 1/16 : on le trouve avec une petite fenêtre de recherche, vite et avec peu de risque de confusion. Une fois trouvé, il est appliqué, et le niveau suivant, plus fin, n'a plus qu'à corriger le reste. C'est la stratégie classique des méthodes d'alignement de Lucas--Kanade~[\hyperref[ann_lucas]{12}, \hyperref[ann_baker]{13}]. IRdrone utilise le schéma suivant (\texttt{iterative\_scheme} dans \texttt{align\_raw}) :

\begin{table}[htbp]\centering
\begin{tabular}{@{}cccccc@{}}
\hline
\textbf{Échelle} & \textbf{Itérations} & \textbf{Blocs} & \textbf{Taille d'un bloc (pleine rés.)} & \textbf{Fenêtre du coût} & \textbf{Plus grand vecteur (pleine rés.)} \\ \hline
1/16 & 2 & $4 \times 4$ & $998 \times 748$ & $\pm 8$ px & $\pm 80$ px \\
1/16 & 2 & $5 \times 5$ & $798 \times 598$ & $\pm 6$ px & $\pm 80$ px \\
1/4 & 3 & $5 \times 5$ & $798 \times 598$ & $\pm 6$ px & $\pm 20$ px \\ \hline
\end{tabular}
\caption{Schéma du raffinement pyramidal. Le coût est calculé sur la fenêtre indiquée, à l'échelle du niveau, mais le minimum est cherché à $\pm 5$ pixels du centre (\texttt{brute\_force\_vector\_field\_search}) : un vecteur de bloc ne peut pas dépasser 80 pixels à 1/16, 20 pixels à 1/4. Les itérations sont numérotées 1, 2, 4, 5, 7, 8, 9 dans les noms des fichiers de traces (le compteur est incrémenté une fois de plus à chaque changement de niveau).}
\label{tab:ann_pyramid}
\end{table}

\subsection{Une itération}
\begin{enumerate}
	\item Les cartes d'énergie de l'image infrarouge sont rééchantillonnées avec le modèle courant (l'identité au départ).
	\item L'image est découpée en une grille de blocs. Pour chaque bloc, le coût \eqref{eq:ann_cost} est calculé pour chaque décalage de la fenêtre de recherche : on obtient une \emph{surface de coût} par bloc (figures~\ref{fig:ann_blocks_first} et \ref{fig:ann_blocks_last}).
	\item Le minimum de chaque surface (à $\pm 5$ pixels du centre), affiné au sous-pixel comme ci-dessus, donne un vecteur de déplacement par bloc : un \emph{champ de vecteurs} grossier, ou flot optique.
	\item Une homographie est ajustée sur ces vecteurs : les 16 ou 25 couples (centre du bloc, centre + vecteur) forment un système surdéterminé pour ses 8 inconnues. Il est résolu par la transformation linéaire directe (DLT) dans une boucle RANSAC~[\hyperref[ann_ransac]{14}], qui ignore les blocs dont le vecteur s'écarte de plus de 3 pixels de la majorité (\texttt{geometric\_rigid\_transform\_estimation}, \texttt{findHomography} d'OpenCV ; Hartley et Zisserman~[\hyperref[ann_hz]{1}], chapitre 4).
	\item Le modèle est mis à jour par composition, $H_k = \Delta H_k\, H_{k-1}$, puis l'itération suivante commence.
\end{enumerate}
Une homographie plutôt que trois angles : elle représente la rotation d'une caméra \emph{et} les erreurs résiduelles des calibrations (pour le point A, l'homographie finale contient un facteur d'échelle de 1,9\,\% horizontalement et 1,1\,\% verticalement, par exemple une petite erreur sur les focales).

\figureIRD[fig:ann_blocks_first]{align_blocks_first.jpg}{Première itération du raffinement du point A (résolution 1/16, blocs $4\times4$, recherche $\pm 8$ pixels). À gauche : superposition après la recherche grossière (visible en vert, infrarouge en magenta ; gris là où ils concordent), avec les vecteurs mesurés (bleu, agrandis 20 fois) et les vecteurs de l'homographie ajustée (rouge). À droite : la surface de coût de chaque bloc ; le point noir est le décalage nul, la croix le minimum. Les vecteurs atteignent 50 pixels et varient régulièrement : un facteur d'échelle et un tangage résiduels, que l'homographie représente bien.}{1}

\figureIRD[fig:ann_blocks_last]{align_blocks_last.jpg}{Dernière itération (résolution 1/4, blocs $5\times5$, recherche $\pm 6$ pixels), sur l'image alignée globalement. Tous les minima sont au centre : les vecteurs résiduels (agrandis 200 fois) sont inférieurs à 2 pixels, l'homographie ne change plus. Ces vecteurs résiduels forment la correction locale (§\ref{sec:ann_local}).}{1}

\figureIRD[fig:ann_convergence]{align_convergence.png}{Longueur des vecteurs des blocs à chaque itération, pour les deux points. Le point A (bleu) converge : de 23 pixels en moyenne à 0,6 pixel. Les marches aux changements de niveau sont dues au changement de taille des blocs et de résolution (un pixel à 1/16 vaut 16 pixels). Le point B (rouge) ne converge pas : les vecteurs restent à leur limite, 80 pixels à 1/16 et 20 pixels à 1/4 (§\ref{sec:ann_limits}).}{0.7}

\subsection{Alignement local}\label{sec:ann_local}
Une homographie est un modèle global : elle ne peut pas représenter la parallaxe du relief ni une imperfection locale du modèle de distorsion. Le champ de vecteurs de la dernière itération ($5 \times 5$ vecteurs) est donc conservé comme \emph{correction locale} : il est interpolé bilinéairement entre les centres des blocs (constant au-delà) en un déplacement pour chaque pixel, composé avec le modèle global (\texttt{warp\_from\_sparse\_vector\_field}). Deux résultats sont produits :
\begin{itemize}
	\item alignement \emph{global} : rotation et homographie finale ;
	\item alignement \emph{local} : rotation, homographie de l'itération précédente et champ de déplacement local. C'est en général le plus précis, et celui des images multispectrales \texttt{IRD\_XXXX\_n.tif}.
\end{itemize}

\subsection{Un seul rééchantillonnage}
Chaque rééchantillonnage d'une image interpole entre ses pixels, et la rend donc un peu floue. Au lieu de rééchantillonner l'image infrarouge à chaque étape, la transformation finale est composée en une seule carte : pour chaque pixel de l'image visible, la position à lire dans l'image infrarouge \emph{brute} s'obtient en enchaînant le déplacement local, l'homographie, la rotation de l'équation \eqref{eq:ann_rotation_homography} et la distorsion de la caméra infrarouge ; l'image n'est ensuite interpolée qu'une fois (interpolation bicubique, \texttt{cv2.remap}). Dans le code ce principe s'appelle \og You Only Warp Once \fg{} (\texttt{FULLRES\_REGISTERED\_REFINED\_WARPED\_ONCE\_ONLY.jpg}). Le modèle de mouvement (angles, homographies, champ de vecteurs) est enregistré dans \texttt{VIS\_XXXX\_motion\_model.npy} et sa version lisible \texttt{.json} ; s'il existe, relancer le Step~4 le relit au lieu de chercher à nouveau.

La figure~\ref{fig:ann_stages} montre la progression de l'alignement sur un détail du point A.

\figureIRD[fig:ann_stages]{align_stages.jpg}{Les quatre étapes sur un détail du point A (visible en vert, infrarouge en magenta, gris là où ils concordent). (a) Calibrations seules : les images sont décalées d'environ 160 pixels. (b) Après la recherche grossière : il reste une dizaine de pixels (franges colorées le long des ombres). (c) Avec l'homographie. (d) Avec la correction locale : les bords coïncident.}{1}

\section{Pistes explorées : corrélation de phase, Fourier--Mellin et appariement polaire}\label{sec:ann_fourier}
\subsection{Corrélation de phase}
Les premières versions (août--septembre 2021) estimaient la translation par \emph{corrélation de phase}~[\hyperref[ann_kuglin]{15}]. Si $g(x) = f(x - t)$, leurs transformées de Fourier vérifient $\hat g(\omega) = \hat f(\omega)\, e^{-i\omega \cdot t}$ (théorème du retard). L'interspectre normalisé
\begin{equation}
	\frac{\hat f\, \overline{\hat g}}{|\hat f\, \overline{\hat g}|} = e^{i\omega \cdot t}
\end{equation}
ne garde que la phase, et sa transformée de Fourier inverse est un pic de Dirac à la position $t$. La méthode est rapide (deux FFT), insensible à un changement global de luminosité, et trouve de grandes translations. Mais elle suppose que les deux images sont le \emph{même} signal décalé : une rotation, un changement d'échelle, une distorsion ou une inversion de contraste étalent le pic. La figure~\ref{fig:ann_phase} le montre : sur le point A, après la recherche grossière, le pic est net ; sur le point B il est noyé dans le bruit, sur les luminances comme sur les cartes d'énergie.

\figureIRD[fig:ann_phase]{align_phase_correlation.png}{Surfaces de corrélation de phase entre l'image visible et l'image infrarouge après la recherche grossière (résolution 1/8, zoom sur $\pm 320$ pixels). En haut, point A : un pic net, environ 2,5 fois plus haut que les valeurs suivantes. En bas, point B (rotation résiduelle, contrastes inversés, grandes zones texturées) : le pic émerge à peine.}{0.65}

\subsection{Fourier--Mellin : rotation et échelle}
La méthode de Fourier--Mellin~[\hyperref[ann_reddy]{16}] étend la corrélation de phase aux rotations et aux changements d'échelle. Le module de la transformée de Fourier ne dépend pas de la translation ; une rotation de l'image le fait tourner, un changement d'échelle le dilate. Rééchantillonné en coordonnées \emph{log-polaires} $(\log r, \theta)$, le module transforme une rotation en translation selon $\theta$ et un changement d'échelle en translation selon $\log r$, qu'une seconde corrélation de phase retrouve. Rabatel et Labbé~[\hyperref[ann_rabatel]{17}] l'ont appliquée à des images de drone visibles et proche infrarouge. Dans IRdrone, l'échelle est connue par les calibrations, et seul le roulis était cherché de cette façon : l'alignement semi-automatique de 2021 rééchantillonnait les deux images en coordonnées polaires et estimait le roulis par corrélation de phase selon l'angle (\texttt{mellin\_transform} et \texttt{estimate\_rotation} dans \texttt{semi\_auto\_registration.py}, toujours présentes pour une démonstration mais non utilisées par le Step~4). Combinée à la translation par corrélation de phase, ou à des points SIFT, elle n'a pas donné de résultats fiables sur les images multispectrales, pour les raisons ci-dessus. À partir de septembre 2021 elle a été remplacée par la minimisation du coût \eqref{eq:ann_cost} sur les cartes d'énergie.

\subsection{Appariement polaire des cartes d'énergie (essai de 2022)}
Le roulis est le point faible de l'algorithme actuel : la recherche grossière ne trouve que le lacet et le tangage, et le raffinement pyramidal ne corrige qu'un petit roulis. Lors d'un vol à Thaumiers (novembre 2022) la caméra infrarouge était montée avec un roulis constant d'environ $-2{,}4^\circ$, soit des déplacements de 100 pixels dans les coins, trop pour le raffinement (ticket GitHub \#33). Un essai a combiné les deux approches : les cartes d'énergie sont rééchantillonnées en coordonnées polaires autour du centre de l'image, où un roulis devient une translation verticale, et le coût \eqref{eq:ann_cost} est calculé pour chaque décalage de $1^\circ$ ; une parabole ajustée sur les trois valeurs les plus basses donne le roulis à mieux qu'un degré (\texttt{coarse\_roll\_estimation}, branche \texttt{local\_alignment\_improvements}, commit \texttt{86353cb}, non intégré). La figure~\ref{fig:ann_polar} le reproduit : sur le point A, tourné d'un roulis connu de $2{,}4^\circ$, l'estimation est de $2{,}26^\circ$. Sur le point B, la route est à environ $34^\circ$ dans l'image visible et à environ $28^\circ$ dans l'image infrarouge après la recherche grossière, et l'appariement polaire trouve un roulis de $-6{,}3^\circ$, mais avec un second minimum presque aussi profond à $+5^\circ$ (figure~\ref{fig:ann_polar_b}) : le profil est beaucoup moins tranché que sur le point A.

\figureIRD[fig:ann_polar]{align_polar_roll.jpg}{Appariement polaire sur le point A, dont l'image infrarouge a été tournée volontairement de $2{,}4^\circ$. À gauche : cartes d'énergie des deux images en coordonnées polaires (angle verticalement, distance au centre horizontalement) ; un roulis est un décalage vertical. À droite : coût en fonction du décalage angulaire, et parabole ajustée autour du minimum ($2{,}26^\circ$).}{1}

\figureIRD[fig:ann_polar_b]{align_polar_roll_b.jpg}{La même chose sur le point B, sans roulis artificiel, après la recherche grossière : minimum à $-6{,}3^\circ$, mais un second minimum près de $+5^\circ$ est presque aussi profond.}{1}

\section{Historique du développement}\label{sec:ann_history}
L'historique git du projet donne les jalons suivants.
\begin{itemize}
	\item \textbf{Août 2021} : alignement semi-automatique. L'utilisateur tourne l'image infrarouge avec trois curseurs (lacet, tangage, roulis, équation~\eqref{eq:ann_rotation_homography}) ; le résultat est affiné par corrélation de phase, par corrélation de phase polaire pour le roulis, ou par des points SIFT ou une corrélation de phase par blocs avec un ajustement affine. Comparaison sur la valeur absolue du gradient plutôt que sur les intensités.
	\item \textbf{Septembre 2021} : fonctions de coût (SSD, NTG) sur plusieurs représentations, recherche de Newton et ajustement quadratique du minimum (d'après le matériel supplémentaire de HDR+~[\hyperref[ann_hasinoff]{10}]), puis recherche itérative pyramidale ajustant une homographie sur les vecteurs des blocs, accélérée avec Numba.
	\item \textbf{Octobre 2021} : déformation locale suivant le champ de vecteurs, caméra virtuelle au champ élargi pour le fish-eye, chaîne complète recherche grossière + pyramide, choix du minimum par la hessienne, rééchantillonnage unique (\og You Only Warp Once \fg), sorties locale et globale.
	\item \textbf{Mars 2022} : ticket GitHub \#4, \og fix local alignment issues \fg{} (toujours ouvert) : utiliser la courbure du coût pour décider si un vecteur local est fiable, régulariser le champ avec ses voisins.
	\item \textbf{Novembre--décembre 2022} : ticket \#33, roulis de la caméra infrarouge ; appariement polaire des cartes d'énergie, raffinement parabolique du minimum (branche expérimentale).
	\item \textbf{Novembre 2026} : accélération sans changement des résultats (coût par linéarité du filtre de Sobel, pyramides en float32 avec OpenCV, réduction par 32 calculée sur les seules lignes et colonnes utiles) : environ 15\,s par shoot point ; réalignement à la main depuis la revue dans une fenêtre dédiée (§\ref{btn:review_redo}).
\end{itemize}

\section{Traces de diagnostic}\label{sec:ann_traces}
Les figures de cette annexe proviennent des traces de diagnostic de l'alignement. Elles ne sont pas écrites par défaut au Step~4 (environ 100\,Mo de figures par point). On les obtient :
\begin{itemize}
	\item en lançant le script \texttt{VIS\_XXXX\_REDO.bat} (ou \texttt{.sh}) écrit à côté de chaque image alignée : un \og full recovery run \fg{} avec l'option \texttt{--debug}, qui ouvre la fenêtre à curseurs puis écrit tout ;
	\item depuis la revue, quand le niveau de détail du journal des réglages du Step~4 (\bouton{⚙}) est \og Debug \fg{} : un réalignement à la main fait alors le calcul complet ;
	\item en ligne de commande : \texttt{python automatic\_registration.py --images VIS\_XXXX.dng NIR\_YYYY.dng --outdir OUT --debug}.
\end{itemize}
Les traces sont écrites dans le dossier \texttt{VIS\_XXXX\_align\_traces} :
\begin{itemize}
	\item \texttt{Full Search\_\ldots} : la recherche grossière (surface de coût, tableau \texttt{.npy} des coûts) ;
	\item \texttt{\_PADDED\_LOWRES\_*}, \texttt{\_LOWRES\_*} : cartes d'énergie à 1/32 avant et après la recherche grossière ;
	\item \texttt{ds16\_\ldots\_cost\_surface\_Y\_X\_\_itNN.png} : les entrées et les quatre surfaces de coût du bloc $(Y, X)$ à l'itération NN ; \texttt{\ldots overview\_cost\_surfaces\ldots} : tous les blocs d'une itération ; \texttt{\ldots\_cost\_search\_\_itNN.npy} : les coûts, les centres des blocs et l'échelle, lisibles par \texttt{numpy.load(\ldots, allow\_pickle=True)} ;
	\item \texttt{\_msr\_itNN\_dsMM\_\ldots}, \texttt{\_image\_itNN\_\ldots} : cartes d'énergie et images à chaque itération ;
	\item \texttt{FLOW\_itNN\_\ldots} : champ de vecteurs et image déformée par celui-ci ;
	\item \texttt{FULLRES\_REF.jpg}, \texttt{FULLRES\_REGISTERED\_COARSE.jpg}, \texttt{FULLRES\_REGISTERED\_REFINED\_WARPED\_ONCE\_ONLY.jpg} : la référence et l'image infrarouge après chaque étape, à comparer en basculant de l'une à l'autre ;
	\item \texttt{motion\_model.npy} : le modèle de mouvement.
\end{itemize}

\section{Limites et problèmes ouverts}\label{sec:ann_limits}
L'alignement multispectral est un problème mal posé, et l'algorithme échoue dans des situations identifiables :
\begin{itemize}
	\item les deux caméras ne voient pas la même scène : rafale de vent, grand écart de temps, basse altitude (parallaxe) ;
	\item la zone a peu de structures : eau, champ uniforme, canopée forestière ; le coût est alors plat ou trompeur (figure~\ref{fig:ann_profiles}c) ;
	\item le désalignement dépasse ce que les étapes savent corriger : un roulis de plusieurs degrés, que la recherche grossière ne modélise pas ;
	\item la recherche grossière tombe dans un minimum trompeur.
\end{itemize}
Le point B cumule plusieurs de ces difficultés : grand lacet ($8{,}2^\circ$), roulis probable de plusieurs degrés, grandes zones de végétation dont les textures diffèrent entre les bandes. Les vecteurs des blocs restent à leur limite de $\pm 80$ puis $\pm 20$ pixels (figures~\ref{fig:ann_blocks_hard} et \ref{fig:ann_convergence}) ; l'ajustement RANSAC retient alors une homographie aux forts termes de perspective, qui déplace les coins de l'image de plusieurs centaines de pixels. Le résultat est acceptable près du centre et faux vers les coins (textures étirées en bas à gauche). Dans nos essais, partir du roulis estimé par appariement polaire ($\pm 6{,}3^\circ$) ne l'a pas fait converger non plus : ce point demande un réalignement à la main. De tels échecs sont détectés à la revue (chapitre~\ref{ch:review}), qui mesure le décalage résiduel des bords.

\figureIRD[fig:ann_blocks_hard]{align_blocks_hard.jpg}{Première itération du raffinement du point B : les surfaces de coût des blocs de végétation n'ont pas de minimum net, les minima sont à des positions aléatoires, souvent au bord de la fenêtre de recherche (vecteurs agrandis 10 fois).}{1}

Le réalignement à la main (§\ref{btn:review_redo}) remplace les angles initiaux nuls par des angles choisis à l'œil, roulis compris. Ces angles ne sont qu'un point de départ : la recherche grossière et le raffinement pyramidal repartent de là et trouvent en général un alignement précis.

Améliorations possibles, en partie listées dans les tickets GitHub \#4 et \#33 :
\begin{itemize}
	\item utiliser la courbure de chaque surface de coût (matrice $A$ de l'ajustement quadratique) comme indice de confiance : un vecteur mesuré sur une surface plate ne devrait compter ni dans l'homographie ni dans la correction locale ;
	\item régulariser le champ de vecteurs local (médiane des voisins, contrainte de régularité) au lieu de l'utiliser brut ;
	\item ajouter le roulis à la recherche grossière (appariement polaire, ou recherche sur quelques valeurs de roulis) ;
	\item normaliser le coût (dénominateur du NTG), pour ne pas favoriser un bloc moins texturé ;
	\item calibrer les couleurs des deux caméras (calibration radiométrique) pour calculer des indices physiquement significatifs.
\end{itemize}

\section{Références}\label{sec:ann_references}
{\renewcommand{\labelenumi}{[\arabic{enumi}]}
\begin{enumerate}
	\item\label{ann_hz} R. Hartley, A. Zisserman, \emph{Multiple View Geometry in Computer Vision}, 2\textsuperscript{e} édition, Cambridge University Press, 2004. Chapitres 2 (géométrie projective du plan, homographies), 4 (estimation d'une homographie, DLT, RANSAC), 6 (modèles de caméra), 7.4 (distorsion radiale), 8.4 (caméra en rotation, rotations synthétiques, mosaïques).
	\item\label{ann_szeliski} R. Szeliski, \emph{Computer Vision: Algorithms and Applications}, 2\textsuperscript{e} édition, Springer, 2022 (chapitre 2 : formation de l'image ; chapitres 8 et 9 : alignement et estimation du mouvement). Gratuit en ligne sur \texttt{szeliski.org/Book}.
	\item\label{ann_brown} D. C. Brown, \emph{Close-range camera calibration}, Photogrammetric Engineering 37(8), 1971.
	\item\label{ann_zhang} Z. Zhang, \emph{A flexible new technique for camera calibration}, IEEE Transactions on Pattern Analysis and Machine Intelligence 22(11), 2000.
	\item\label{ann_viola} P. Viola, W. M. Wells III, \emph{Alignment by maximization of mutual information}, International Journal of Computer Vision 24(2), 1997.
	\item\label{ann_scharr} H. Scharr, \emph{Optimale Operatoren in der digitalen Bildverarbeitung}, thèse de doctorat, université de Heidelberg, 2000.
	\item\label{ann_roberts} L. G. Roberts, \emph{Machine perception of three-dimensional solids}, thèse de doctorat, MIT, 1963.
	\item\label{ann_ntg} S.-J. Chen, H.-L. Shen, C. Li, J. H. Xin, \emph{Normalized Total Gradient: a new measure for multispectral image registration}, IEEE Transactions on Image Processing 27(3), 2018 (\texttt{arXiv:1702.04562}).
	\item\label{ann_irani} M. Irani, P. Anandan, \emph{Robust multi-sensor image alignment}, ICCV 1998.
	\item\label{ann_hasinoff} S. W. Hasinoff et al., \emph{Burst photography for high dynamic range and low-light imaging on mobile cameras}, SIGGRAPH Asia 2016 (matériel supplémentaire : alignement pyramidal par blocs, raffinement sous-pixel quadratique).
	\item\label{ann_burt} P. J. Burt, E. H. Adelson, \emph{The Laplacian pyramid as a compact image code}, IEEE Transactions on Communications 31(4), 1983.
	\item\label{ann_lucas} B. D. Lucas, T. Kanade, \emph{An iterative image registration technique with an application to stereo vision}, IJCAI 1981.
	\item\label{ann_baker} S. Baker, I. Matthews, \emph{Lucas-Kanade 20 years on: a unifying framework}, International Journal of Computer Vision 56(3), 2004.
	\item\label{ann_ransac} M. A. Fischler, R. C. Bolles, \emph{Random sample consensus: a paradigm for model fitting with applications to image analysis and automated cartography}, Communications of the ACM 24(6), 1981.
	\item\label{ann_kuglin} C. D. Kuglin, D. C. Hines, \emph{The phase correlation image alignment method}, Proceedings of the IEEE International Conference on Cybernetics and Society, 1975.
	\item\label{ann_reddy} B. S. Reddy, B. N. Chatterji, \emph{An FFT-based technique for translation, rotation, and scale-invariant image registration}, IEEE Transactions on Image Processing 5(8), 1996.
	\item\label{ann_rabatel} G. Rabatel, S. Labbé, \emph{Registration of visible and near infrared unmanned aerial vehicle images based on Fourier-Mellin transform}, Precision Agriculture 17, 2016.
	\item\label{ann_vayssade} J.-A. Vayssade, G. Jones, J.-N. Paoli, C. Gée, \emph{Two-step multi-spectral registration via key-point detector and gradient similarity}, 2020.
\end{enumerate}}

\section{Détermination des \og Pitch \fg{ } et \og Yaw \fg{ } théoriques.}\label{sec:ann_theorique}

\subsection{Description du positionnement des caméras.}

La première étape du recalage de l'image infrarouge sur l'image visible consiste à appliquer une translation à l'image NIR de façon à aligner son centre avec celui de l'image VIS. En effet à la prise de vue les axes de visée des deux caméras ne sont pas parallèles. De plus le déclenchement des caméras n'est pas simultané. La caméra du drone est montée sur un \og gimbal \fg{} qui maintient la direction de visée quels que soient les mouvements du drone
\footnote{Dans la limite des spécifications indiquées par le constructeur DJI et reproduites sur la figure \ref{AnglesDroneGimbal}}.

\figureIRD[AnglesDroneGimbal]{AnglesDroneGimbal.pdf}{Description schématique du \og gimbal \fg{} et des amplitudes angulaires autorisées.}{0.6}

A l'inverse la caméra NIR est fixée rigidement au drone (figure \ref{Definition_Reperes_Cameras}). Ainsi son angle de visée va suivre les mouvements de ce dernier. Le déclenchement se fait en continue avec une période  $\Delta T^{Nir} = 3s$. C'est la fréquence la plus élevée possible pour la SJCam M20.\\


\figureIRD[Definition_Reperes_Cameras]{Definition_Reperes_Cameras.pdf}{Définition des angles de lacet, tangage et roulis (yaw, pitch, roll) du drone, du gimbal et des deux caméras VIS et NIR. Les caméras pointent vers le sol alors que le drone est horizontal ce qui engendre la permutation des axes yaw et roll entre le drone et le gimbal d'une part et les caméras d'autre part.}{0.8}

Nous allons prendre l'exemple du pitch en admettant un découplage entre les angles. Ceci correspond à faire l'hypothèse implicite que les angles yaw et roll sont nuls (ou très faibles).  Nous définissons dans le plan vertical passant par l'axe  du drone les différents angles permettant de repérer la position du drone et des deux caméras à un instant donné (figure \ref{Definition_des_angles_de_Pitch}): 
\figureIRD[Definition_des_angles_de_Pitch]{Definition_des_angles_de_Pitch.pdf}{Définition des angles dans le plan vertical.}{0.8}

\begin{itemize}
	\item $\alpha\qquad$  Flight Pitch. Angle de tangage du drone.
	\item $\delta\qquad$  NIR misalignement Pitch. Cet angle correspond à un défaut de positionnement du support de la caméra NIR sur le drone.
	\item $\varphi$ Angular wedge. Correspond à une éventuelle cale de correction du défaut $\delta$.
	\item $\zeta\qquad$ Residual deviation Pitch. $\quad \zeta = \delta + \phi$.
	\item $\theta_{Nir}\:\:\;$ Camera NIR Pitch.  $\quad \theta_{Nir} = \alpha + \zeta$.
	\item $\beta\qquad$ Gimbal Pitch. Angle de tangage de gimbal (et de la caméra VIS).
	\item $\delta_{Vis}\;\;\;$ Offset VIS. Par définition $\delta_{Vis} =90^\circ$.
	\item $\theta_{Vis}\;\;\;$ Camera VIS Pitch.  $\quad \theta_{Vis} = \beta + \delta_{Vis}$. 
	\item $\psi\qquad$  Theoritical NIR Pitch. $\quad \psi =\delta_{Vis}+\beta-\alpha$.
	\item $\psi_{r}\quad\;\;$ NIR Pitch. Angle entre les axes de visées des deux caméras au même instant $t$. $\qquad \qquad \zeta = \psi_{r} -\psi$. \\
\end{itemize}

Les données Exif des images de la caméra VIS indiquent les coordonnées GPS mais aussi les angles de yaw, pitch et roll du drone et du gimbal. Autrement dit $\alpha$ et $\beta$ sont connus à chaque acquisition d'une image de la caméra VIS. L'inclinaison de l'axe de visée engendre une déformation de l'image en trapèze comme le montre la figure \ref{Effet_du_pitch_sur_camera_pinhole}.

\figureIRD[Effet_du_pitch_sur_camera_pinhole]{Effet_du_pitch_sur_camera_pinhole.pdf}{Effet du tangage sur la déformation de l'image d'un plan horizontal formée sur le capteur d'une caméra sténopé (\og pinhole\fg{}). On notera que l'image sur le capteur est inversée.}{0.6}


Pour simplifier nous admettrons dans la suite que le  montage de la caméra NIR est sans défaut, autrement dit $\delta=0$. 
Nous supposons que les horloges des deux caméras VIS et NIR sont synchronisées. 
On dispose donc des instants exacts où les images de ces deux caméras ont été prises. 
Les paires d'images(VIS, NIR) sont formées de telle manière que l'intervalle de temps séparant les prises de vues de ces deux images soit minimum.

\subsection{Calcul du pitch théorique dans le cas synchronisé.}

\figureIRD[Pitch_cas_synchro]{Pitch_cas_synchro.pdf}{Position du drone et des deux caméras dans le cas d'images prises au même instant. Le pitch a été exagéré pour les besoins de la figure.}{0.5}

Si les images VIS et NIR sont prises au même instant (cas d'une prise de vue synchronisée
\footnote{On ne confondra pas la synchronisation des horloges et la synchronisation des prises de vue.}) il est possible de déterminer à chaque prise de vue l'angle $\psi $ entre les axes de visée des deux caméras. 

Notons que la distance entre les centres des caméras $C^{Nir}$ et $C^{Vis}$ (\og base line \fg{}) est négligeable devant l'altitude $H$ du drone par rapport au sol. 
Le rapport est inférieur à $4\;10^{-4}$. 
On peut donc  prendre $\phi \cong \theta_{Vis}$.

Enfin puisque nous supposons la caméra NIR fixée sans défaut $\theta_{Nir} =\alpha_k$. \\

En observant la figure \ref{Pitch_cas_synchro} et en utilisant des formules de géométrie élémentaires nous en déduisons la formule du pitch théorique pour l'image infrarouge dans le cas d'un déclenchement synchronisé à l'instant $t_k$
\begin{equation}
	\label{Theoritical_Pitch_syncho_2}
	\psi_k  =\phi-\theta_{Nir} =\beta_k + \delta_{Vis} - \alpha_k  \quad .
\end{equation}

La distance entre les centres des images, calculée au niveau du sol est
\begin{equation}
	L = H \tan \psi_k \quad ,
\end{equation}
et  la translation à appliquer à l'image NIR pour la centrer sur l'image VIS s'écrit, en notant $\mathrm{f}^{Nir}$ la distance focale de la caméra infrarouge 
\begin{equation}
	\mathbf{t}_y = \mathrm{f}^{Nir} \tan \psi_k \quad .
\end{equation} 

Bien entendu cette situation est un cas idéal
\footnote{C'est la configuration retenue pour le matériel professionnel.}.

\subsection{Estimation du pitch théorique dans le cas désynchronisé.}
En pratique le déclenchement des caméras VIS et NIR n'est pas simultané. 
Avec une acquisition en \og hyperlapse \fg{}
\footnote{Le terme \og time-lapse\fg{}  ou chronophotographie  correspond à la prise de clichés avec un intervalle de temps régulier et pour une position fixe de la caméra. En \og hyperlapse \fg{} au contraire la caméra est mobile avec un mouvement régulier.}
l'intervalle pour la caméra VIS est de $2s$
\footnote{En pratique l'intervalle est plutôt égal à $2,019s$. Autrement dit environ toutes les 50 images VIS la date d'enregistrement présente un saut de $1s$. Ce phénomène est corrigé par notre algorithme afin d'obtenir une \og ligne temporelle\fg{} précise au dixième de seconde.} alors qu'il est de $3s$ pour la caméra NIR. 
La figure \ref{Periodicite-du-pitch-NIR} montre un exemple d'une série d'images hyperlapse  replacées sur un même axe  temporel une fois les horloges synchronisées.\\

\figureIRD[Periodicite-du-pitch-NIR]{Periodicite-du-pitch-NIR.pdf}{Images hyperlapse des caméras VIS et NIR. Elles sont désynchronisées d'un délai variable positif ou négatif avec une périodicité de trois images visibles, soit $6s$. Pour vérifier la désynchronisation des prises de vues nous avons utilisé la présence d'un véhicule mobile. Vol FLY-20211109-1m/s.}{1}


Cette fois la base line $\mathbf{Bl}(t)$ à l'instant de la prise de vue de l'image NIR ne peut plus être négligée. Les formules de calcul du pitch deviennent
\footnote{ $(\mathbf{e}_1, \mathbf{e}_2, \mathbf{e}_3)$ repère orthonormé direct. $\mathbf{e}_2$ est dans l'axe du drone et $\mathbf{e}_3$ est sur l'axe vertical vers le haut. }

\begin{equation}
	\label{Theoritical_Pitch_non_syncho_0}
	\begin{aligned}
		\psi_k&=\phi-\alpha(t) \quad,\\
		\phi&= \arctan \left(\frac{\mathbf{Bl}(t) \cdot \mathbf{e}_2}{H} +\tan \theta_{Vis} \right)\quad.\\	
	\end{aligned} 
\end{equation}\\

Malheureusement le pitch du drone $\alpha(t)$ n'est pas connu à l'instant $t$ où l'image NIR a été prise. Il en est de même de la position exacte du drone à l'instant $t$, et donc de la \og base line \fg{} $\mathbf{Bl}(t)$. \\

Nous supposerons qu'entre deux images VIS espacées de $\Delta T^{Vis} =2s$ le drone se déplace en ligne droite et à une altitude constante.
On propose de déterminer $\alpha(t)$ et $\mathbf{Bl}(t)$ à l'aide d'une simple interpolation linéaire.\\

Le délai entre la prise de vue des images VIS et NIR est connu et noté $\delta t= t_k -\,t$ . 
Il peut être positif ou négatif selon que l'image NIR est prise avant ou après l'image VIS. 
Les positions GPS du drone par rapport au sol sont connues.  
Dans le repère géographique 
\footnote{$(\mathbf{e}_{\mathrm{W}}, \mathbf{e}_{\mathrm{N}}, \mathbf{e}_3)$ repère orthonormé direct.}
nous utiliserons les coordonnées UTM $(X, Y, Z)$. Bien entendu il faut déterminer les déplacements du drone dans son propre repère. L'orientation du drone par rapport au nord géographique est donnée par l'angle de lacet $\textit{Yaw}_k$. Soit\\ 
\begin{equation}
	\label{Theoritical_Pitch_non_syncho_11}
	\begin{aligned}
		{x}^{EW}_k&=X(t_{k+1})-X(t_{k}) \quad,\\
		{y}^{SN}_k&=Y(t_{k+1})-Y(t_{k}) \quad,\\
		{x_1}_k&=-{x}^{EW}_k\; \cos \textit{Yaw}_k + {y}^{SN}_k\; \sin \textit{Yaw}_k \quad,\\
		{x_2}_k&=\;\;{x}^{EW}_k\:\: \sin \textit{Yaw}_k + {y}^{SN}_k\; \cos \textit{Yaw}_k  \quad.\\
	\end{aligned} \\
\end{equation}


\begin{equation}
	\label{Theoritical_Pitch_non_syncho_1}
	\begin{aligned}
		si \;\delta t <0 \qquad& \qquad\alpha(t)&=&\;\alpha_k \left(\frac{\delta t}{ \Delta T^{Vis} }+1\right) - \alpha_{k+1}\,\frac{\delta t}{ \Delta T^{Vis}} \quad,\\
		&\mathbf{Bl}(t) \cdot \mathbf{e}_2&=&\; {x_2}_k\,\frac{\delta t}{ \Delta T^{Vis}} \quad,&\\
		&&&\\
		si \;\delta t >0 \qquad &\qquad \alpha(t)&=&\;\alpha_{k-1} \,\frac{\delta t}{ \Delta T^{Vis}}- \alpha_k\, \left(\frac{\delta t}{ \Delta T^{Vis} }-1\right)  \quad,\\
		&\mathbf{Bl}(t) \cdot \mathbf{e}_2&=&\; {x_2}_{k-1}\,\frac{\delta t}{ \Delta T^{Vis}} \quad.&\\
	\end{aligned} \\
\end{equation}\\

Un exemple des positions du drone et des caméras est représentée sur la figure \ref{Pitch_cas_non_synchro}.\\ 

\figureIRD[Pitch_cas_non_synchro]{Pitch_cas_non_synchro.pdf}{Cas d'une désynchronisation \og arrière \fg. L'image NIR a été prise après l'image visible. La \og base-line \fg{} a été exagérée. En pratique le point vert est situé avant le point rouge pour une désynchronisation arrière.}{0.9}


L'écart de désynchronisation $\delta t$ est périodique et nous noterons sa période  $T$. 
Dans notre exemple $T= 2 \Delta T^{Nir} = 3 \Delta T^{Vis}$. 
Le pitch de la caméra VIS, contrôlé par le gimbal, est quasi constant. 
En l'absence de vent et avec une vitesse constante le pitch du drone est lui aussi constant. 
Dans ces conditions idéales le signal $\psi(t)$ sera lui aussi de période T. 
Dans le cas d'un vol avec du vent le pitch du drone va varier pour compenser les turbulences. 
Ce phénomène peut être traduit par un pitch aléatoire 
\footnote{On ne rentrera pas dans des considérations très complexes sur la modélisation de la turbulence. Par exemple l'ergodicité,  la turbulence homogène  etc. }. 
Le pitch aura donc une composante périodique et une composante aléatoire. Le spectre de Fourier d'un tel signal présentera  un pic à la fréquence $0,166\, Hz$ ce qui correspond à la composante périodique de période $T=6\,s$.\\ 

\figureIRD[Image ViR]{_VIR__local_HYPERLAPSE_0057.jpg}{Image multi-spectrale VIR produite par IRDrone. FLY-20211109-1m/s.}{0.65}

\subsection{Analyse sur des données réelles.}

On utilise ici les données de la mission FLY-20220125 effectué le 25 janvier 2022 en fin de matinée (10 h 59, heure solaire). Le ciel était clair, le vent faible (< 3 m/s) du nord-ouest. Après la phase de synchronisation le drone s'est élevé à une altitude de 100m puis a débuté son vol d'acquisition des images \og utiles \fg{}. Pendant cette phase l'orientation du drone est restée constante ($-82^\circ$) de même que sa vitesse $1,5\,m/s$. Pendant les $232 s$ de la phase d'acquisition  les caméras VIS et NIR ont pris respectivement $116$ et $85$ images. La trajectoire du  drone se décompose en trois segment droits:
\begin{itemize}
	\item Vol vers l'avant en ligne droite. Direction $\sim 281^\circ$. Durée $166s$.
	\item Vol vers la droite. Direction  $\sim 8^\circ$ (normale à la direction précédente). Durée $26s$.
	\item Vol vers l'arrière. Direction  $\sim 95^\circ$ (parallèle à la direction initiale). Durée $40s$.\\
\end{itemize}
\figureIRD[Fly_Blassac_trajectoire]{Fly_Blassac_trajectoire.png}{Schéma de la mission FLY-20220125. Les shoot point qui ont servi au mapping sont repérés par les marques rouges. Ils ont un $|dt| \leq 0.5\,s$}{0.8}

Après synchronisation des horloges l'algorithme a formé 116 paires d'images (VIS, NIR). Notons qu'une image NIR sur deux sera associée à deux images VIS successives. 

\figureIRD[Pitch_Yaw_reel_vs_theorique]{Process-and-Theoretical_Offset_Pitch_vol_reel_01.pdf}{Évolution temporelle des angles yaw et pitch pendant la phase d'acquisition. Les angles sont déterminés par les formules théoriques et par le  \og coarse process \fg{}. Les angles théoriques ont été corrigés pour tenir compte des défauts de positionnement de la caméra NIR sur le drone.  FLY-20220125.}{0.99}

La caméra VIS du drone est placée en mode \og gimbal lock \fg{}
\footnote{Dans ce mode particulier il n'est pas certain que  l'information sur le lacet du gimbal soit fiable. La différence entre les angles de lacet du drone et du gimbal est constante et égale à $2,7^\circ$. Pour certain vols cette différence peut atteindre $7^\circ$. Ici l'angle de lacet correspond au roulis des caméras. Cependant en comparant les paires d'images VIS et NIR on ne constate pas de rotation entre les deux. C'est un problème de calibrage de la caméra au moment de passer en gimbal lock.}.
Les informations récupérées dans les données Exif des images du drone indiquent que les angles yaw, pitch et roll du gimbal sont constants ($-78,9^\circ$, $-90^\circ$  et $0^\circ$ respectivement). \\


Le yaw de la caméra NIR est pris égal au  roll du drone. 
Le calcul du yaw théorique à appliquer à l'image NIR pour la recentrer sur l'image visible utilise le même type de formules que celle mises en  évidence pour le pitch théorique (courbe \og Theoretical yaw \fg{} sur la figure \ref{Pitch_Yaw_reel_vs_theorique}). \\ 
L'évolution des valeurs théoriques des angles yaw et pitch suit globalement celle obtenue par le \og coarse process\fg{}.\\ 


Notons qu'il y  a un décalage systématique ($\sim +0,9^\circ$ ) sur le yaw. Ce décalage a été retrouvé sur l'ensemble des vols effectuées jusqu'à présent. Il peut être attribué à un défaut de positionnement de la caméra NIR sur le drone. Un décalage systématique existe aussi pour le pitch \og{}théorique \fg{} par rapport au pitch \og{} process \fg{}. Nous l'attribuons lui aussi à un défaut de positionnement de la caméra NIR sur le drone qui engendre une erreur systématique d'alignement de l'axe de visée de la caméra. Ce défaut résiduel est $\zeta = -1,3^\circ$.
\footnote{Ce que confirme une mesure réalisée à l'aide d'un inclinomètre numérique mesurant les angles avec une précision de $1/10\, ^\circ$. La mesure a été prise entre le support carbone et la surface du filtre NIR.} .\\

L'évolution des angles présente sur certaines parties un aspect oscillant ce que confirme l'analyse spectrale (DFFT) . 
Le spectre de Fourier du pitch représenté sur la figure \ref{Fourier_Pitch_process} montre  un pic net pour la fréquence $0,164\,Hz$ pour les valeurs théoriques ou obtenues par le  \og coarse process \fg{}.  
Cela correspond à une composante périodique de période $T=6\,s$. 
Nous avons montré que ce phénomène est liée à la variation périodique de la \og base line \fg{} (voir la figure \ref{Pitch_cas_non_synchro} et les formules  \eqref{Theoritical_Pitch_non_syncho_0}, \eqref{Theoritical_Pitch_non_syncho_1} ).
Le reste du spectre correspond à l'effet du vent.
Lorsque le drone vol en ligne droite vers l'avant ou l'arrière la projection de sa \og bas-line \fg{} sur son axe est périodique. 
Le pitch  a donc une composante périodique lors de cette phase de vol. 
La projection de sa \og bas-line \fg{} sur l'axe normal est nulle. 
Dans cette phase de vol le yaw a donc seulement la composante aléatoire associée au vent. 
Lorsque le drone se déplace perpendiculairement à son axe (en gardant toujours la même orientation) c'est le yaw qui possède une composante périodique alors que le pitch n'a qu'une composante aléatoire liée  au vent. 
Ce phénomène est clairement visible sur la figure \ref{Pitch_Yaw_reel_vs_theorique}. Les changements d'allure des courbes permettent même d'identifier les phases de vol.\\


\figureIRD[Fourier_Pitch_process]{Spectre-Pitch-process-vol-reel_test.png}{Transformée de Fourier du Pitch  calculé  par le \og coarse process\fg{} à partir des images numérique ou par les formules théoriques.  FLY-20222501.}{0.9}

Notons que le pitch théorique calculé avec la formule synchronisée \eqref{Theoritical_Pitch_syncho_2} ne fait intervenir que les angles du drone et du gimbal. Il ne présente pas de caractère périodique comme le confirme sa DFFT. Ceci est cohérent puisque la variation de ces angles est associée aux effets \og{}aléatoires \fg{} du vent sur le drone (turbulence).


\subsection{Conclusion.}


La méthode proposée pour déterminer a priori le lacet et le tangage (angles yaw et pitch) à appliquer à l'image NIR pour la recentrée sur l'image VIS repose sur plusieurs éléments :\\

\begin{itemize}
	\item Les images VIS et NIR doivent être prises en mode time-lapse. 
	La fréquence d'enregistrement de la caméra VIS (celle qui a le pas le plus petit, ici $2\,s$) doit être régulière. 
	Ceci nécessite l'utilisation de carte SD rapides. Toutefois pour le DJI les dates sont enregistrées à la seconde près. On remarque périodiquement des sauts d'une seconde. Il semble que le pas ne soit pas exactement un entier. Pour notre vol FLY-20211109-1m/s on trouve $2,01865 \, s$. Ce phénomène est retrouvé sur tout les vols effectués jusqu'à présent avec une valeur de $2,02 \, s$. Le logiciel IRDrone détermine automatiquement pour chaque vol la période exacte pour obtenir une fréquence constante.  
	On utilise le format de date de Python en  renseignant son attribut \texttt{micosecond}. 
	Notons que les données Exif des dates prévoient (pour certaines) une valeur décimale des secondes  
	La situation pour la caméra NIR est différente. Il existe aussi des sauts mais de la valeur du timelapse ($3\,s$). Les dates sont donc toujours des multiples du timelapse. Ici c'est clairement un défaut d'enregistrement d'une image
	\footnote{Comme on a pu le vérifier lors de la synchronisation manuelle des images, la séquence retenue  comportant justement un tel défaut d'enregistrement.}. Les dates n'ont donc pas été rectifiées. Lors d'un saut de ce type la \og base-line\fg{} est plus grande. Cela reste acceptable si la vitesse du drone est faible. Notons toutefois qu'il faut adapter les formules de calcul pour en tenir compte. A nouveau le code gère automatiquement ce phénomène.\\
	
	\item Les données d'attitude  (yaw, pitch, roll) doivent être disponibles pour les deux caméras. C'est bien le cas avec notre montage puisque les données Exif des images VIS (format dng) du DJI fournissent ces angles pour le drone (donc pour la caméra NIR qui lui est fixée rigidement) et pour le gimbal (donc pour la caméra VIS qui est fixée dessus).\\  
	\item Le calcul de la \og base-line \fg{} nécessite de connaître les coordonnées UTM. Actuellement l'altitude $H$ du drone par rapport au sol est déterminée en connaissant l'altitude du point de décollage, l'altitude du drone par rapport au point de décollage et l'altitude à l'aplomb de la prise d'image (via une API de l'IGN pour la France). .\\
	\item La synchronisation des horloges des deux caméras est un point critique. On utilise une séquence de synchronisation à effectuer à chaque décollage. Elle consiste à faire tourner
	\footnote{La cible est placée sur un plateau mis en rotation par un moteur pas à pas. Ce moteur est commandé par un micro contrôleur ARDUINO.} un QR code (ARUCO marker) sous les deux caméras. Un algorithme détecte ensuite ce QR code sur les images et en déduit l'orientation. Une méthode de minimisation permet de rapprocher  les courbes angle/date. La  précision de cette synchronisation est supérieure au $1/10^{\grave{e}me}s$.  Soulignons qu'il est impératif de ne pas interrompre la prise de vue en mode hyperlapse pendant toute la durée de la mission (phase de synchronisation et phase d'acquisition). En effet la valeur de synchronisation n'est valable (avec cette précision) que pour une séquence hyperlapse donnée. Entre deux séquences hyperlapse la valeur de synchronisation peut varier de plus d'une seconde (allant même jusqu'à 3s dans le cas le plus défavorable).
	Un tel défaut de synchronisation des horloges risque de fausser l'appariement des images (et le calcul de la  \og base-line \fg{}).
	Pour l'exemple présenté ici nous avons réalisé aussi cette synchronisation manuellement à partir d'un objet mobile présent sur 25 images consécutives (soit 50s). Le mobile avait été déplacé spécialement pour cet utilisation lors de notre vol d'essai. Cette technique est précise (précision de l'ordre du $1/4 s$) mais fastidieuse. Nous avons retrouvé le même appariement des images que celui obtenu avec la synchronisation par cible ARUCO. \\
	\item L'angle de pitch du drone ainsi que la \og base-line \fg{} à l'instant où l'image NIR est capturée ont été  interpolés linéairement. Cette méthode donne d'excellents résultats.\\
\end{itemize}

Nous avons vérifié la pertinence de l'ensemble de ces éléments sur plus d'une dizaine de vol avec des conditions météorologiques variées, notamment le vent.  
Par exemple les figures \ref{Pitch_Yaw_reel_vs_theorique_FLY2021109} et \ref{Fourier_Pitch_process_FLY2021109} correspondent à un vol programmé en hyperlapse (\og mode line\fg{}) avec une vitesse de $1\,m/s$ et environ $600$ images.

\figureIRD[Pitch_Yaw_reel_vs_theorique_FLY2021109]{Process-and-Theoretical_Offset_Pitch_vol_reel_02.png}{Évolution temporelle des angles yaw et pitch pendant la phase d'acquisition. Ici  le drone se déplace à cap constant ($179^\circ$) parallèlement à son axe. FLY-20211109.}{1}

\figureIRD[Fourier_Pitch_process_FLY2021109]{Spectre-Pitch-process-vol-reel_02.png}{Transformée de Fourier du Pitch. FLY-20211109.}{.75}




\chapter{Géoréférencement et orthorectification (Step 6)}\label{ann:georef}
\todo{Texte technique : géométrie caméra/drone, repères, projection sur le MNT, précision, destiné aux utilisateurs avertis.}

\end{document}
"
% (see build_help.bat). Images are then taken from images_light/ instead of images/.
\ifdefined\IRDlight
  \newcommand{\IRDimgdir}{images_light}
  \graphicspath{{images_light/}}
\else
  \newcommand{\IRDimgdir}{images}
  \graphicspath{{images/}}
\fi

% Symbols of the software buttons (pdflatex does not know them)
\DeclareUnicodeCharacter{25C0}{\ensuremath{\blacktriangleleft}}
\DeclareUnicodeCharacter{25B6}{\ensuremath{\blacktriangleright}}
\DeclareUnicodeCharacter{2699}{\ensuremath{\circledast}}
\DeclareUnicodeCharacter{2922}{\ensuremath{\nearrow}}
\DeclareUnicodeCharacter{21BA}{\ensuremath{\circlearrowleft}}
\DeclareUnicodeCharacter{21BB}{\ensuremath{\circlearrowright}}
\DeclareUnicodeCharacter{25B2}{\ensuremath{\blacktriangle}}
\DeclareUnicodeCharacter{25BC}{\ensuremath{\blacktriangledown}}
\DeclareUnicodeCharacter{2713}{\ensuremath{\checkmark}}
\DeclareUnicodeCharacter{2717}{\ensuremath{\times}}
% Software button
\DeclareRobustCommand{\bouton}[1]{\texorpdfstring{\fbox{\small\textsf{#1}}}{#1}}
% To be completed
\newcommand{\todo}[1]{\textcolor{red}{[To be written: #1]}}
% Legacy sections: collapsed on the web manual, plain text in the PDF
\newcommand{\legacybegin}[1]{}
\newcommand{\legacyend}{}
% Tolerant figure: shows a frame if the image does not exist yet.
% Usage: \figureIRD[label]{file}{caption}{width}
%   - optional label (e.g. fig:img_VIR_1), usable with \ref{...} / \autoref{...}
%   - file with or without extension; without extension, tries .png, .jpg, .jpeg then .pdf
\makeatletter
\newcommand{\ird@file}{}
\newcommand{\ird@find}[1]{%
  \def\ird@file{}%
  \IfFileExists{\IRDimgdir/#1}{\def\ird@file{#1}}{%
  \IfFileExists{\IRDimgdir/#1.png}{\def\ird@file{#1.png}}{%
  \IfFileExists{\IRDimgdir/#1.jpg}{\def\ird@file{#1.jpg}}{%
  \IfFileExists{\IRDimgdir/#1.jpeg}{\def\ird@file{#1.jpeg}}{%
  \IfFileExists{\IRDimgdir/#1.pdf}{\def\ird@file{#1.pdf}}{}}}}}}
\newcommand{\figureIRD}[4][]{%
  \ird@find{#2}%
  \begin{figure}[htbp]\centering
    \ifx\ird@file\@empty
      \fbox{\parbox{0.6\linewidth}{\centering\small\textcolor{gray}{Missing image: \detokenize{#2}}}}%
    \else
      \includegraphics[width=#4\linewidth]{\ird@file}%
    \fi
    \caption{#3}%
    \if\relax\detokenize{#1}\relax\else\label{#1}\fi
  \end{figure}}
% Pair of side-by-side figures, placed exactly where it is written ([H]).
% Usage: \figurePaireIRD{label}{left file}{left caption}{right file}{right caption}{global caption}
%   - the label refers to the pair; the sub-figures have the labels <label>_avant and <label>_apres
\newcommand{\ird@incsub}[1]{%
  \ird@find{#1}%
  \ifx\ird@file\@empty
    \fbox{\parbox{0.8\linewidth}{\centering\small\textcolor{gray}{Missing image: \detokenize{#1}}}}%
  \else
    \includegraphics[width=\linewidth]{\ird@file}%
  \fi}
\newcommand{\figurePaireIRD}[6]{%
  \begin{figure}[H]\centering
    \begin{subfigure}[t]{0.48\linewidth}\centering
      \ird@incsub{#2}%
      \caption{#3}\label{#1_avant}
    \end{subfigure}\hfill
    \begin{subfigure}[t]{0.48\linewidth}\centering
      \ird@incsub{#4}%
      \caption{#5}\label{#1_apres}
    \end{subfigure}
    \caption{#6}\label{#1}
  \end{figure}}
\makeatother

\title{IRdrone\\User manual}
\author{}
\date{\today}


\begin{document}

	% --------------------------------------------------
	% Cover page
	% --------------------------------------------------
	\begin{titlepage}
		\centering

		\vspace*{2cm}
		%{\Huge\bfseries IRdrone\par}
		%\vspace{0.5cm}
		{\LARGE User manual\par}
		\vspace{6cm}
		\includegraphics[width=0.95\textwidth]{Logo-IRdrone.png}

		\vfill
		{\large Version of \today\par}
	\end{titlepage}

	% --------------------------------------------------
	% Table of contents
	% --------------------------------------------------
	\tableofcontents



%==========================================================
\chapter{The IRDrone project}\label{ch:overview}


\section{Observing the vegetation cover.}

\subsection{General context.}
Vegetation mapping is nowadays routinely used in many kinds of environmental assessments (Natura 2000 sites, ZNIEFF, environmental impact studies, etc.).
This mapping is based on a segmentation (outlining of areas that are homogeneous in terms of vegetation) of the vegetation cover in aerial plan view.
This approach is usually complemented by field surveys (botanical, phytosociological surveys, etc.).
The segmentation into homogeneous areas relies on the analysis of natural-color aerial photographs, in particular their color, texture, apparent structure, opacity, etc.

\figurePaireIRD{fig:etude_de_cas}
{step0_carto_aerienne.jpg}{Delineation of habitats on an aerial image background.}
{step0_habitats.jpg}{Resulting habitat map on an IGN topo 25 background.}
{Example of use of field data with maps and aerial photographs.}

However, it is difficult to distinguish tree species and plant species (herbaceous species) on conventional aerial photographs. Landscape dynamics and fine landscape features are also hard to identify.
The available IR aerial photographs are often old, and therefore hard to use for current assessments.
Moreover, the national coverage is incomplete, with some regions poorly covered.
In addition, all existing photographs are taken at high altitude (3-4000 m) and therefore lack precision for assessments commonly carried out at a scale of 1/5000 or finer. \\




Using a drone to acquire photographs has several advantages: the shots are centered on the study areas, the images are far more precise (flight altitude of 120 m), as many recent images as needed can be obtained, and they can be taken in different seasons.


\subsection{Value of multispectral images.}

Aerial photographs in the near infrared
\footnote{This near infrared radiation, which in daytime comes from solar radiation and is reflected by objects, must not be confused with far infrared radiation, which is emitted by all objects.} NIR carry more ecological information than natural-color photographs.
The main purpose of NIR aerial photography is therefore to make aerial photographs easier to interpret, based on the differences in infrared reflectance between vegetation types (due to structural differences between the leaf tissues of plant species).

\figureIRD[fig:img_NDVI_1]{img_NDVI_1.png}{Map of the NDVI vegetation development index.}{0.7}
\figureIRD[fig:img_VIR_1]{img_VIR_1.png}{Example of a VIR composite image. The R/G/B layers are replaced by the NIR/R/G layers. This kind of representation is very useful, among other things, to distinguish between conifers and broadleaf trees.}{0.7}


\subsection{What are we trying to observe?}

The use of aerial photographs is justified for several reasons, mainly related to the precision and speed of the analysis.
In general, on NIR aerial images the contrasts between the different habitat types are enhanced, which allows a finer reading of the landscape.
In dense forest stands, it is difficult to estimate the conifer/broadleaf cover proportions from a field survey because of the large areas to cover and the height of the trees.
Distinguishing conifers from broadleaf trees on natural-color photographs is often ambiguous.
On the contrary, NIR aerial photography makes it possible to identify precisely and quickly the contribution of each tree species in mixed stands covering large areas, and thus to obtain information on the specific management of the stands (types of silvicultural treatment, age classes, etc.).  Moreover, some tree stands mixing broadleaf trees and conifers in variable proportions may or may not fall under European protection (``Habitats'' directive), depending on the share of broadleaf trees.
With NIR aerial photographs, the assessment of tree species proportions becomes objective and precisely quantifiable.
The analysis of NIR aerial photographs also gives a better understanding of vegetation dynamics.

The detection of very young woody plants becomes possible because of their bright coloring and their position within open habitats (woody colonization fronts within open areas), whereas their presence is hard to detect in natural color because of the small size of their crowns and their color, similar to that of the colonized open habitats.
Another important application is the immediate and exhaustive location of water bodies (which may be inconspicuous, small and therefore go unnoticed in natural color), which appear as dark spots  in NIR.
This application is particularly important because of the regulatory status of wetlands, which are generally one of the main issues of environmental assessments.


\subsection{Principle of the observation.}

Plants produce chlorophyll by absorbing visible light from the sun (VIS). Only a small part of this radiation is used in this mechanism. Yet the solar radiation reaching the ground also contains ultraviolet (UV) and near infrared (NIR) \footnote{UV and NIR radiation are not visible to the human eye.}

\figureIRD[Facteur_reflexion_chloro-carot]{Spectres-absorption-pigments-photosynthetiques.pdf}{Absorption factor of photosynthetic pigments in the visible spectral band.}{1}

It is mainly the visible radiation in the ``blue'' ($[400\,nm ,\, 450\,nm]$) and ``red'' ($[630\,nm ,\, 680\,nm]$) bands that contributes to chlorophyll production.
The part of the radiation not absorbed by the plant is transmitted (the plant is ``transparent'') or reflected (sent back to the environment).
Since the chlorophyll yield is zero in the ``yellow-green'' spectral band ($[500\,nm ,\, 600\,nm]$), this band is reflected by the plant (Figure~\ref{Facteur_reflexion_chloro-carot}).
This is where the green color of these plants comes from. Likewise, UV and IR radiation do not contribute to chlorophyll production.
On the other hand, they would heat the plant if they were absorbed.
So, to protect themselves, plants do not absorb the UV and IR radiation emitted by the sun, which is then transmitted or reflected.
Plants are totally transparent to ultraviolet radiation.
On the other hand, they reflect part of the infrared radiation.
It is this property that is exploited in the project.
Conifers and broadleaf trees do not have the same infrared reflectance.
Figure~\ref{Facteur_reflexion_vegetaux} shows typical spectral reflectance curves of conifers and broadleaf trees.
Note that broadleaf trees also produce beta-carotene and fucoxanthin, which absorb blue and part of green but reflect red.
This is why in autumn, when broadleaf trees no longer produce chlorophyll but beta-carotene is still present, these species display shades ranging from orange-brown to red ochre.

\figureIRD[Facteur_reflexion_vegetaux]{Facteur_reflexion_vegetaux.pdf}{Reflectance of conifers and broadleaf trees. In the visible spectral band $[400\,nm ,\,700\,nm]$ the behavior of the two species is very similar. On the other hand, in the IR-A infrared spectral band $[700\,nm ,\, 1.4\,\mu m]$, well present in the solar spectrum, the two species are clearly distinct.}{1}


Note that tree leaves reflect radiation around 550 nm, which corresponds to the green color perceived by the human eye. But since the curves are identical in the visible part of the spectrum (400 nm to 700 nm), the different tree species cannot be distinguished. On the other hand, the reflectances in the infrared range are clearly different between conifers and broadleaf trees. The idea is to exploit this property by producing an infrared image  of the vegetation cover
\footnote{Possibly combined with an image in the visible spectrum.}. On this image, broadleaf trees should appear  ``lighter''  than conifers. A map of the vegetation cover can then be drawn. Note however that, depending on the time of year and the age of the trees, these curves may change, making the interpretation of the images more difficult.

\figureIRD[Facteur_absorption_vegetaux]{Par_action_spectrum.jpg}{Absorption factor of plants and photosynthesis yield.}{0.5}


\subsection{Equipment for the observation.}
To draw a map, ``bird's-eye view'' images are needed. Professionals use satellite photographs (alt. 300km) or photographs taken from aircraft (alt. 2000-3000 m). But as explained above, their spatial and temporal resolution is not suited to detailed field studies. In this project we therefore plan to use a small drone, which gives aerial views of the study areas with great flexibility (alt. 120 m), with very good spatial resolution and full control over the shooting periods.

Most existing ``small'' drones cannot take natural-color and IR shots simultaneously. Yet it is essential to have both types of georeferenced shots (IR for interpreting the vegetation and natural color for producing standardized maps). \\

There are multispectral drones and multispectral cameras for drones
\footnote{For example DJI Phantom 4 multispectral, Parrot Sequoia, Micasense RedEdge.} .
These are multi-channel cameras that acquire, besides the conventional visible image, narrow spectral bands in the red, green and blue, and one or two bands in the near infrared. There are even versions fitted with thermal cameras \footnote{For example the Micasense Altum camera has a Flir IR sensor. It weighs 400g. However, this is thermal infrared radiation (LWIR), whose wavelengths are ten times longer than those of solar near infrared.  Note that this equipment is very expensive.}.
Other drones
are fitted with two cameras, one in the visible and the other in the infrared.
Finally, the camera of some drones
\footnote{For example the DJI Mavic Pro.2, fitted with a very high-performance Hasselblad camera. Its ``hot mirror'' filter is attached directly to the sensor. ``Heavy'' drones with gimbals can carry digital SLR cameras, which can also be ``unfiltered''. }
can be made sensitive in the near infrared. This requires removing
\footnote{This ``unfiltering'' operation is delicate. It can be done by specialized companies.} the internal filter of the digital sensor that blocks infrared radiation and fitting the optical system with a suitable filter
\footnote{With a set of external filters, the type of IR spectral band can be chosen. With a $760\,\mathrm{nm}$ IR high-pass filter, a wide band $[760, 1100]$ is obtained. There are also narrow-band filters, for example $850\,\mathrm{nm}\pm 25\mathrm{nm}$). A ``hot mirror'' filter can even be put back, which restores imaging in the visible spectrum.}.
Moreover, the drone then records only one type of image, depending on the filter it is fitted with. \\


The idea of the \textsf{IRDrone} project is to fit a small,  affordable drone
\footnote{Here a DJI Mavic Air.2, weighing no more than 600 g.} with a second  camera. The drone's own camera provides color images in the visible spectrum (RGB). The second on-board camera, which we will call the IR camera, must produce images in the ``solar'' near infrared spectral band
\footnote{In practice, with the available CMOS sensors, the accessible spectral band extends from $720 \,\mathrm{nm}$ to $1100\,\mathrm{nm}$. It includes the narrow RedEdge and NIR (for Near Infra Red) bands}.\\

The NIR camera must be light so as not to disturb the drone's flight or reduce its flight time. It must be low-cost, which naturally leads either to surveillance cameras (``night vision'') or to consumer ``action cams''. The former are not suitable because their resolution is poor. The latter are designed to give nice images in the visible range. The optics must therefore be adapted by removing the ``hot mirror'' filter and fitting an IR filter.  \\


After each flight, the collected images must be processed by computer (geometric and color corrections, registration and fusion of visible and IR images ...). This is the role of the Python code  IRDrone.py of the \textsf{IRDrone} project, whose features are  described below.


\section{Organization of an IRDrone mission}

\subsection{Preparation at the office}

Before the active part of the mission begins, the pilot must carefully prepare the flight.
First at the office, by identifying the access routes, the topography of the terrain, the obstacles and the size of the study area.
This size determines the flight altitude. One may plan a single high-altitude flight to cover the area at once, or several successive flights. Keep in mind that the DJI drone carrying its infrared camera has a flight time limited to 25 minutes.
In practice the synchronization phase takes 2 minutes. At a speed of $3\,\mathrm{m/s}$ the drone can travel about 4\,km, which is not much given the need to ``sweep''~the survey area so that the shots overlap (if the images are to be stitched together).
It is also assumed that the operator holds a drone pilot qualification and the flight authorizations.

\subsection{Flight preparation in the field}
Once in the field, the pilot sets up the take-off area. After marking it out, he places the take-off stand and the ArUco target synchronization system with its power supply
\footnote{What if the synchronization device has been forgotten? The mission can be ``saved''~by slowly moving an object under the drone. The synchronization offset of the two cameras will then have to be estimated manually. It can work\dots{} sometimes!}.
A folding chair is useful for piloting while watching the camera feed, so as to follow the flight plan as closely as possible.
Likewise, a small windsock to assess the wind direction and speed is useful.
Finally, precisely measuring the position of the take-off point with a GPS in fixed position is a plus\footnote{See the capabilities of a Garmin 750 GPS, which has this kind of function.}.
Before each field mission \textbf{the drone's GPS must be calibrated and its camera calibrated} (viewing direction).
These two actions are essential for the images to be successfully processed by the IRDrone software.
The pilot checks that the batteries are charged and that the (empty) SD cards are in the DJI drone and in the SJCam M20 camera.
It is wise to preset the DJI camera parameters (those that can be preset) before take-off.
At this point the pilot can start the hyperlapse of the infrared camera (SJCam M20), place the drone on its take-off stand and finally power on the synchronization device (in the latest version of the device, the target only starts rotating once the drone has left the take-off stand).

\subsection{The mission flight}
Once these preparations are done, the mission has 4 successive phases:
\begin{enumerate}
	\item \emph{Takeoff Point -- Synchronization}: synchronization right next to the take-off point, at an altitude of about 2\,m.
	At the start of this phase the NIR camera is already running in hyperlapse mode (shooting interval 3\,s).
	During this phase the pilot stabilizes the drone above the synchronization device, sets the VIS camera parameters and starts the VIS shooting in hyperlapse mode. He stays in a fixed position for about ten seconds.
	\item \emph{Approach}: moving to the area to be surveyed for the study.
	The pilot quickly gains altitude.
	He moves at the maximum speed allowed by the DJI hyperlapse mode.
	A value of $3\,\mathrm{m/s}$ is a good choice. When preparing the mission, it is strongly recommended to choose a take-off point as close as possible to the survey area.
	\item \emph{Hyperlapse Survey}: acquisition of the multispectral images above the survey area (\emph{survey area}).
	The drone moves while keeping a constant orientation.
	It moves forward, then sideways, then backward, etc. (see Figure~\ref{fig:4_step_of_the_mission}).
	At take-off it is recommended to orient the drone in the main direction of travel above the survey area.
	When leaving the survey area the pilot can stop the hyperlapse recording of the VIS camera.
	This has two advantages: it avoids filling up the SD card, since the images are high resolution, but above all it removes the speed constraint imposed on the drone by the hyperlapse mode.
	\item \emph{Return to Home} (RTH): this is the phase of return to the take-off point. This return can be done at maximum speed. The RTH mode can be engaged for an automatic flight.
\end{enumerate}

\figureIRD[fig:4_step_of_the_mission]{4_step_of_the_mission.pdf}{The flight phases of an IRDrone mission in the field.}{0.8}

\subsection{Processing the images}

Once the acquisition is over, the next phase is to process the images of the two cameras digitally with the IRDrone software.


\section{Purpose of the IRDrone software}
IRDrone.py is an interactive, user-friendly application that generates multispectral images for environmental impact or botanical studies.
To do so, it uses the images taken during field campaigns by two cameras, one in the visible spectrum (``VIS'') and the other in the near infrared spectrum (``NIR'').
Both cameras are mounted on a drone.
The system currently supports a DJI Mavic Air 2 drone with its VIS camera mounted on a gimbal, and an SJCam M20 camera modified to capture images in the near infrared spectrum (typically beyond 810 nm).
The application synchronizes the clocks of the two cameras, then aligns the VIS and NIR images.
The main process produces aligned multispectral images (4 bands red, green, blue, nir in tif format).
A post-processing step orthorectifies and georeferences the images, either individually or as a mosaic (mapping with the ODM code; Figure~\ref{fig:img_mapping_1}).
\figureIRD[fig:img_mapping_1]{img_mapping_sample.png}{Mosaic of images in the VIS spectrum made with Open Drone Map (ODM). The display is done with the QGIS software. The contour lines have an interval of 1\,m. The red dots show the location of the shots and the shot number.}{0.8}
The results can be used in the QGIS software to obtain specialized images (NDVI index, Figure~\ref{fig:img_NDVI_1}; VIR nir/red/green image,
Figure~\ref{fig:img_VIR_1}).


\section{Main features of the IRDrone software}
At startup, a small ``Loading, please wait'' window is shown while loading (a few tens of seconds), then the main window opens (Chapter~\ref{ch:home}). It lists the missions of the working folder and, once a mission is open, its seven steps with their progress. Clicking a step starts it; steps 1 to 4 follow on from one another automatically.

Processing a mission has 7 steps:
\begin{enumerate}
	\item \textbf{Mission}: the images of the two cameras are found on the SD cards, copied into a new mission folder, and the mission parameters are saved.
	The name of the mission folder is generated automatically from the location and the date
	\footnote{\label{fn:MissionFolderName} By default the folder name is FLY-YYYYMMDD-hhmm-location. It can be changed once and for all when the mission is created, in the ``Location'' field. For example, if the ``Location''~field has the value Blassac, one can specify ``Blassac Bourg''.}.
	\item \textbf{Flight timeline}: the synchronization images (ArUco target at take-off) and those of the survey above the study area are detected automatically; the user checks them.
	\item \textbf{Synchronization}: the offset between the clocks of the two cameras is measured and the VIS and NIR images are paired.
	\item \textbf{Alignment}, the core of IRDrone: each image pair is aligned to obtain 4 superimposed spectral layers.
	\item \textbf{Review} of the aligned images (Chapter~\ref{ch:review}): every image is checked on a timeline, the misaligned ones are flagged (left out of the map) or realigned by hand.
	\item \textbf{Mapping and georeferencing}: georeferencing, orthorectification, orthophoto mosaic (mapping), terrain model and contour lines.
	\item \textbf{Archiving} of the mission.
\end{enumerate}

\figureIRD[fig:Air_Mission_folder]{Air_Mission_folder.png}{The main \texttt{Air-Mission} folder with two missions being processed (here Chalivoy-Milon and Labout) and one archived mission (here Blassac). On Windows the icons are assigned automatically.}{0.6}

\subsection*{On-line help}\label{sec:help}
\begin{itemize}
	\item \textbf{Tool tips.} Leaving the mouse pointer for a moment on a button (or a check box) shows a tool tip that summarizes what it does.
	\item \textbf{\textsf{F1} key.} \textsf{F1}, with the mouse pointer over a button, opens this manual at the section that describes this button. On a step of the home page, it opens the chapter of the step. Elsewhere in a window, it opens the chapter of the window, like the \bouton{?} button.
	\item \textbf{\bouton{?} button.} Each window has a \bouton{?} button that opens this manual at the chapter of the window. In the windows driven from the keyboard (timeline, review), its tool tip recalls the keys. It sits at the bottom left; the \bouton{Close} button of the step windows and of the review is at the bottom right; in Steps 2 to 5, the button that continues with the next step (\bouton{Continue with \ldots ▶}) sits to its right and never starts a computation by itself. The actions of a step (\bouton{Run alignment}, \bouton{Build the map\ldots}, \ldots) are in the window, above the progress bar; the main one is dark blue, the others white, and the irreversible one (archiving) red. Steps~3 to~6 show at the top the name of the step and, under it, the mission. In the main window too, at the bottom left, it opens the whole manual.
	\item \textbf{\bouton{↺} and \bouton{Close} buttons of the windows that edit something} (mission details, timeline, offset adjustment, Step~4 settings). Closing keeps what was changed: as soon as something is changed, \bouton{Close} becomes \bouton{Save and close}; closing the window by its title bar or with \textsf{Esc} does the same. The \bouton{↺} button, next to \bouton{?}, appears only when there is something to put back: it puts the values back and the window stays open.
\end{itemize}
The manual opens in the default PDF reader (browser, Acrobat), at the right page when the reader allows it.

\section{Digital organization of an IRDrone mission}

%\todo{Folder \texttt{FLY-YYYYMMDD-HHmm-localit\'e}, sub-folders \texttt{ImgIRdrone}, \texttt{mapping\_MULTI}, etc.}
Here we assume that your Windows computer has a folder \texttt{C:/Air-Mission}.
This working folder is chosen the first time the application is started; it is shown at the bottom of the main
window and can be changed with the ``gear'' button next to it.

Each mission is placed in a folder \texttt{C:/Air-Mission/FLY-YYYYMMDD-hhmm-location}.

This folder contains the sub-folders:
\begin{enumerate}
	\item \texttt{AerialPhotography}
	      \begin{itemize}
	      	\item \texttt{VIS}
	      	  \begin{itemize}
	      	  	\item \texttt{VIS\_XXXX.dng}. Transferred images from the DJI camera,
	      	  	\item \texttt{VIS\_XXXX.exif}. ``Companion''~files of the images, containing exif data,
	      	  	\item \texttt{time\_line.json}. File giving the chronology of the images relative to the VIS camera clock,
	      	  	\item \texttt{transfer\_info\_VIS\_dng.json}. File summarizing the choices made in Step~2,
	      	  	\item \texttt{VIS\_XXXX\_RawTherapee.tif}. Images\footnote{\label{fn:rawtherapee_1}These very large files are essential for the digital processing of the images. They are generated only once, the first time Step~4 is used. When the mission is fully processed and is to be archived, they can be deleted.} created in Step~4.
	      	  \end{itemize}

	      	\item \texttt{NIR}
	      	\begin{itemize}
	      		\item \texttt{NIR\_XXXX.dng}. Transferred and converted images from the infrared camera (proprietary RAW to dng),
	      		\item \texttt{NIR\_XXXX.exif}. ``Companion''~files of the images, containing exif data,
	      		\item \texttt{time\_line.json}. File giving the chronology of the images relative to the NIR camera clock,
	      		\item \texttt{transfer\_info\_NIR\_dng.json}. File summarizing the choices made in Step~2,
	      		\item \texttt{NIR\_XXXX\_RawTherapee.tif}. Images\textsuperscript{\ref{fn:rawtherapee_1}} created in Step~4.
	      		\end{itemize}
      		\item \texttt{shoot\_points.json}. Dictionary of the \emph{shoot points} (shots, using as reference the VIS image and the NIR image closest in time to the VIS shot). It is created in Step~3.
	    \end{itemize}
	\item \texttt{cameras}
	         \begin{itemize}
	         	\item Currently not used.
	         \end{itemize}
	\item \texttt{FlightAnalytics}. This is a technical folder.
			\begin{itemize}
				\item \texttt{mission\_parameters.json}. File containing the mission parameters (editable with the ``Mission details'' window, section~\ref{btn:step1_details}).
				\item \texttt{flight\_timeline.json}. Flight phases confirmed in Step~2.
				\item \texttt{fly\_vertical\_profil.png}. Graph showing the shooting altitudes above the ground.
				\item \texttt{Angles analysis.png},  \texttt{Angles analysis dispersion.png}, \texttt{Angles analysis spectrum.png}. Graphs analyzing the flight quality. The first of these three graphs shows, for example, whether the attitude of the drone during the mission stayed within the flight specifications.
			\end{itemize}
	\item \texttt{ImgIRdrone}
		\begin{itemize}
		\item For each shoot point number XXXX there are:
			\begin{itemize}
				\item \texttt{VIS\_XXXX\_REDO.bat},
				\item \texttt{VIS\_XXXX.jpg},
				\item \texttt{VIR\_XXXX\_local.jpg},
				\item \texttt{VIS\_XXXX\_motion\_model.json},
				\item \texttt{VIS\_XXXX\_motion\_model.npy}.
			\end{itemize}
		\item \texttt{review.json}. Measurements and checked points from the review of the aligned images (Chapter~\ref{ch:review}).
		\item \texttt{shoot\_points.json}. The list of processed shoot points with a success flag. Caution: by default this flag assumes that the alignment is correct. The user has a procedure to examine the alignments and mark them as failed if needed.
		\item \texttt{imgMultispectral}. Folder containing the quadruplets of multispectral images in tif format. Here these images are not georeferenced. For each shoot point XXXX:
			\begin{itemize}
				\item \texttt{IRD\_XXXX\_1.tif}. Red layer image.
				\item \texttt{IRD\_XXXX\_2.tif}. Green layer image.
				\item \texttt{IRD\_XXXX\_3.tif}. Blue layer image.
				\item \texttt{IRD\_XXXX\_4.tif}. Near infrared layer image.
				\item \texttt{IRD\_XXXX\_1.exif}. ``Companion''~file containing exif data\footnote{The exif data are also present in the \texttt{IRD\_XXXX\_Y.tif} images. However, extracting some of these exif data requires a specific third-party tool such as ExifTool. This is a costly operation that is done only once. The companion exif files can therefore be seen as cached data.}.
			\end{itemize}
		\item \texttt{VIS\_XXXX\_align\_traces\_manual}. Folder containing all the information and images resulting from a manual alignment of the VIS and NIR images of shoot point XXXX.
		\end{itemize}
	\item \texttt{mapping\_MULTI}. Working folder of Step~6 (mapping).
	\item \texttt{Synchro}. Image previews and target detections prepared for Step~2 (\texttt{Thumbnails}, a cache rebuilt when needed); Step~3 files: working images \texttt{*\_RawTherapee.jpg} (sub-folders \texttt{VIS} and \texttt{NIR}), control images of the target detection (\texttt{Check\_ARUCO}) and \texttt{check\_time\_*.png} graphs.
\end{enumerate}

\section{Sequence of the steps}
%\todo{Recommended order, prerequisites of each step.}
You must have on a digital medium (SD card, hard disk) the images of the two cameras taken during the mission flight. It is strongly recommended to have only one mission per SD card. The DJI images are in dng format and stored in a \texttt{DCIM/HYPERLAPSE} folder. The SJCam M20 camera images are in RAW format, each paired with a jpeg image. Since the RAW images have no thumbnail, they can be selected by looking at the associated jpeg image. The images are stored in a \texttt{DCIM} folder.

For a new mission, steps 1 to 6 must be followed in order.
However, steps 2 to 6 can be run individually, provided that the mission folder contains at least the elements produced by the steps before it. This is particularly useful for steps 4 to 6.

The steps are started from the main window, described in the next chapter.

%==========================================================
\chapter{The main window}\label{ch:home}
The main window has two pages: the list of missions, then the open mission. At the bottom, a bar shows the working folder and opens this manual.

\section{List of missions}\label{btn:home_missions}
With no mission open, the window offers (Figure~\ref{fig:home_welcome}):
\begin{description}
	\item[\bouton{New mission}] creates a mission from the SD cards of the drone and of the infrared camera (Step~1, Chapter~\ref{ch:step1}).
	\item[\bouton{Open a mission}] opens a mission already created: choose its \texttt{FLY-YYYYMMDD-hhmm-location} folder.
	\item[Recent missions] the missions in the working folder, most recent first, with their progress: one green dot per completed step. A click opens the mission.
\end{description}
The \bouton{All missions} button on the mission page returns to this list.

\figureIRD[fig:home_welcome]{home_welcome.png}{The list of missions: new mission, open a mission, recent missions and their progress.}{0.85}

\section{The open mission}\label{btn:home_steps}
The mission page (Figure~\ref{fig:home_mission}) shows:
\begin{itemize}
	\item at the top, the name of the mission (a click opens its folder in the file explorer) and the place, the date and the pilot;
	\item the seven steps, from left to right: green tick and date for a completed step, blue frame for the next step, grey for the others. \textbf{A click on a step runs it}, even if already done (to redo it). Two exceptions:
	\begin{itemize}
		\item a completed Step~1 only opens the ``Mission details'' window (Section~\ref{btn:step1_details}): the SD cards are not read again;
		\item a completed Step~3 shows the saved result (clock offset, aligned curves) without recomputing anything (Section~\ref{sec:step3_chained}).
	\end{itemize}
	\textsf{F1} over a step opens its chapter in this manual;
	\item the next step, explained in one sentence, with its button (Section~\ref{btn:home_continue});
	\item a few figures: number of images, of pairs, of aligned images, of points flagged as misaligned;
	\item on the right, the flight map or an aligned image (Section~\ref{btn:home_map}).
\end{itemize}

\figureIRD[fig:home_mission]{home_mission.png}{The open mission: the steps (a click runs a step), next step, mission figures and flight map over the orthophoto and the IGN aerial photographs.}{0.9}

\section{\bouton{Continue with Step N} button}\label{btn:home_continue}
Runs the next step. Steps 1 to 4 follow one another by themselves:
\begin{enumerate}
	\item at the end of Step~1, \bouton{Flight timeline} opens Step~2;
	\item \bouton{Continue with synchronisation} in Step~2 (or \bouton{Continue with Step 3}) opens Step~3; \bouton{Synchronise the cameras} runs it: target detection, clock offset, then one question, ``Do the two curves overlap?''; \bouton{Yes} pairs the images;
	\item \bouton{Continue with image processing} opens Step~4. The alignment, which is long, never starts by itself: it is started with \bouton{Run alignment} and can be stopped at any time with \bouton{Stop};
	\item at the end, \bouton{Continue with the review} opens Step~5, the review of the aligned images (it opens by itself when the alignment was chained from Step~3);
	\item \bouton{Continue with mapping} in the review opens Step~6, where \bouton{Build the map\ldots} starts the map;
	\item once the map is built, \bouton{Continue with archiving} opens Step~7, which shows how much space archiving frees; nothing is deleted until you confirm.
\end{enumerate}
What remains for you is to check the flight phases, the overlap of the synchronization curves, and then the result. Once all the steps are done, the button becomes \bouton{Archive the mission} (Step~7).

\section{Flight map}\label{btn:home_map}
The map is available from Step~1 on: the positions are read from the EXIF data of the visible images. It shows:
\begin{itemize}
	\item the drone trajectory and one dot per image, colored according to the phases confirmed in Step~2 (purple: synchronization, green: survey);
	\item the images of the map, and only them: the trajectory joins them, the whole flight remaining as a faint dashed line, and while no orthophoto is shown their ground footprints are drawn in green, so you see the overlap of the map. Which images:
	\begin{itemize}
		\item before Step~4, the images kept by the best mapping (Steps~2 and~4, section~\ref{btn:step2_best_mapping}), when it is chosen;
		\item once Step~4 is done, the points exactly as in the review of the aligned images (chapter~\ref{ch:review}): every shoot point in the colour of its status (green aligned, orange to check, red flagged\ldots), the points Step~4 did not align as small grey dots; an image aligned or realigned from the review shows at once;
		\item the orthophoto of Step~6, once built, is drawn under these points (the same points as the review, not only the images it used);
	\end{itemize}
	the flagged images (\bouton{Flag}, review of the aligned images, chapter~\ref{ch:review}) are small red dots; hovering one shows its footprint dashed in red: you see the hole it leaves in the overlap; flagged after Step~6, they are still in the orthophoto shown and will be left out when Step~6 runs again (\texttt{left out when Step 6 runs again}). The number of images and of flagged images is written at the bottom right. \bouton{Layers} > \bouton{Only the images of the map (and the flagged ones)} (\bouton{Only the images of the best mapping} before Step~4) shows every image again;
	\item the orthophoto from Step~6, once the mapping is done: as soon as Step~6 has built it, the main window shows it on the map (even if \bouton{Layers} > \bouton{Hidden} was chosen before);
	\item in the background, aerial photographs: IGN for France, Esri World Imagery elsewhere, or OpenStreetMap.
\end{itemize}
Hovering over a dot shows the thumbnail of its image; a click on a dot shows the ground footprint of the image (the thick line is the top of the image). Mouse wheel: zoom; drag: pan; double-click or \bouton{⤢} button: the whole flight.

The \bouton{Layers} menu chooses the background (IGN, Esri, OpenStreetMap or none) and the display of the orthophoto, in natural colors or in infrared false colors. The background comes from the internet: offline, the map remains usable (flight, footprints, orthophoto), only the background is missing, and nothing is blocked. Tiles already seen are kept in a cache (\texttt{\textasciitilde/.cache/irdrone} on Linux, \texttt{\%LOCALAPPDATA\%\textbackslash IRDrone} on Windows) and are then displayed offline.

The footprint is computed from the height above ground (terrain model from Step~6, or IGN altitudes) and the heading of the image; once the mapping is done, the heading comes from OpenDroneMap, more accurate than the drone compass (often off by 5 to 10\textdegree).


\section{Bottom bar: working folder and manual}\label{btn:home_workdir}
On the page of the list of missions, the bottom bar shows the working folder (the folder that contains the missions, \texttt{C:/Air-Mission} by default) and the free space on its disk; a click on the folder opens it in the file explorer. The \bouton{⚙} gear wheel changes the working folder: the list of missions is then read again. It also gives an advanced option, unchecked by default: entering the camera clock offset by hand in Step~3, for flights without the ArUco target (Section~\ref{btn:step3_clock_offset}). The folder is chosen at first launch; the \texttt{IRDRONE\_WORKING\_DIR} environment variable, if set, takes precedence. On the page of a mission, the bottom bar shows on the right the total size of the mission folder (\texttt{Mission size}), updated when a step is done and every minute.

\bouton{?}, on the left, opens this manual (Section~\ref{sec:help}).

\section{Check the installation}\label{btn:home_check}
On the first page, the \bouton{Check the installation} card (also in the \bouton{⚙} options) tries each link of the processing chain on sample images shipped with IRDrone, the way the steps do, in a temporary folder whose name holds accents (about 20~seconds; missions are not touched):
\begin{itemize}
  \item Python (64~bits), Python packages (with the \texttt{pip} command to run when one is missing), OpenCV with ArUco and a single OpenCV package;
  \item ExifTool, RawTherapee and the SJCam converter: found \emph{and} run (version);
  \item working folder writable and free space; images and EXIF in an accented folder; 16~bits TIFF written and read back; worker processes of Step~4;
  \item synthetic ArUco target: the five markers found, the rotation of the mobile marker measured within 3°, the absolute angle independent of the drone yaw (Step~3);
  \item DJI DNG: date, gimbal and GPS read by ExifTool (Steps~1 and~2), developed by RawTherapee as in Step~3 (JPEG) and Step~4 (linear TIFF);
  \item SJCam RAW: converted to DNG with its date (Step~1), developed as in Steps~3 and~4;
  \item visible / near infrared alignment: a known shift found within 1.5~pixel (Step~4).
\end{itemize}
Each line turns green (\texttt{✓}), orange (warning, \texttt{!}) or red (failed, \texttt{✗}); a check needing a tool that failed is skipped (\texttt{–}). What is not green says \emph{what is wrong} and \emph{how to fix it}. \bouton{Show the details} shows everything the checks print; \bouton{Copy the results} copies the text; \bouton{Create a diagnostic report…} creates the zip below with these results. The text is also written in the console (and so in the session log). Without opening the software: \texttt{python selftests.py}.

\section{Diagnostic report}\label{btn:home_diagnostic}
When something goes wrong on a computer, the diagnostic report gathers in a single zip file everything needed to understand the problem; just send this zip to the developers. It is created in two ways:
\begin{itemize}
  \item in the options (\bouton{⚙} gear wheel at the bottom of the main window, Section~\ref{btn:home_workdir}), \bouton{Create a diagnostic report…} button: choose a mission, or the installation only;
  \item without opening the software, with \texttt{irdrone\_diagnostic.bat} (at the root of the IRDrone folder): double-click for the installation only, or drag and drop the folder of a mission (\texttt{FLY-...}) onto this file.
\end{itemize}
The zip (\texttt{irdrone\_diagnostic\_YYYYMMDD\_hhmmss.zip}) is written in the working folder. It contains:
\begin{itemize}
  \item \texttt{report.txt}: a summary of what looks abnormal, then the installation (IRDrone version, Python used, system, Windows code page, versions of the Python packages including OpenCV, paths and versions of ExifTool, RawTherapee and the SJCam converter, working folder, free space) and the installation checks (Section~\ref{btn:home_check}), each with how to fix it;
  \item for a mission: the steps done, the Step~2 windows, the time lines, the Step~3 offset and number of pairs; every synchronisation image is developed again with RawTherapee in a temporary folder and the ArUco markers are searched on it, and on the cached image (\texttt{Synchro}) used by Step~3. A few developed images are attached. This takes about a minute; the mission is not modified;
  \item \texttt{report.json} (all the details) and the logs of the last sessions: each launch of the software copies the console into \texttt{logs\textbackslash irdrone\_YYYYMMDD\_hhmmss.log}, next to the settings file (the last ten are kept).
\end{itemize}
Once the zip is written, the text of \texttt{report.txt} is shown in a window (to read, copy with \bouton{Copy the text} or screenshot) and printed in the console. Nothing is sent automatically.

%==========================================================
\chapter{Step 1 -- Define the mission}\label{ch:step1}
\section{The ``New mission'' window}\label{sec:step1_new_mission}
Step~1 creates the mission folder and stores the flight images in it. It is opened with \bouton{New mission} in the main window and takes place in a single window, in three pages: \textbf{Images}, \textbf{Mission}, \textbf{Copy}. The progress is shown at the top of the window. The SD cards are only read: nothing is written to or deleted from them.

\subsection{Mission type: multispectral or visible only}\label{btn:step1_mission_type}
The top of the ``Images'' page chooses the type of the mission (the last one chosen is proposed):
\begin{itemize}
	\item \bouton{Multispectral (visible + infrared)}: the drone (visible, DNG) and the infrared camera (SJCam, RAW). Everything described in this manual: synchronisation, alignment, map in four bands (red, green, blue, NIR), VIR and NDVI;
	\item \bouton{Visible only (drone DNG)}: the drone images only, without the infrared camera. The infrared card disappears, the copy only takes the visible images, and the next steps adapt (Section~\ref{sec:visible_only}): survey phase only in Step~2, no Step~3, the images are \emph{developed} in Step~4 instead of aligned, the map of Step~6 is in colour (red, green, blue).
\end{itemize}
The type cannot be changed once the mission is created.

\subsection{``Images'' page}\label{btn:step1_images}
The inserted SD cards are scanned (any volume that contains a \texttt{DCIM} folder); they hold the visible images (drone DNG) and the infrared images (SJCam RAW), Figure~\ref{fig:step1_new_mission_images}. Each camera has its card:
\begin{itemize}
	\item the folder found, the number of images, their size, the date and time of the first and of the last one;
	\item two images: \textbf{at take-off} and \textbf{in flight} (two thirds of the way through the images, usually during the survey). For the infrared, this is the JPEG that the SJCam saves next to each RAW. These images let you check at a glance that the card is the right one;
	\item the \bouton{Choose folder…} button (\bouton{Change…} once images are found) chooses another folder, for example a copy of the cards on the disk. A folder chosen this way is never replaced by a card found later.
\end{itemize}
As long as the images of one camera are missing, the window looks for SD cards every 3~seconds: a card inserted after the window was opened is found by itself, just insert it.

The take-off image is the first visible image; the ``change…'' link of the visible card chooses another one. It gives the date, time, position and camera of the mission.
A banner reports a difference of more than two minutes between the clocks of the two cameras: this is normal, Step~3 measures it precisely. A difference of several days casts doubt on whether the two cards come from the same flight.

\figureIRD[fig:step1_new_mission_images]{step1_new_mission_images.png}{``Images'' page: visible and infrared images found on the SD cards, at take-off and in flight.}{0.8}

\subsection{``Mission'' page}\label{btn:step1_mission}
The date, time and position come from the take-off image; the ground altitude and the address are requested from IGN\footnote{\label{fn:ign}IGN: Institut national de l'information g\'eographique et foresti\`ere (French national institute of geographic and forest information).} and OpenStreetMap while you fill in the description and the pilot (Figure~\ref{fig:step1_new_mission_details}). The name of the mission folder (\texttt{FLY-YYYYMMDD-hhmm-location}) is updated with the ``Location'' field.

\bouton{Cameras and filter} expands the camera settings: shooting interval and clock offset of each camera, manufacturer and model of the infrared camera, manufacturer and band of the filter. The pilot, the pilot's ID, the cameras and the filter are remembered for the next mission.

\figureIRD[fig:step1_new_mission_details]{step1_new_mission_details.png}{``Mission'' page: place found from the take-off image, folder name, description, pilot, cameras.}{0.8}

\subsection{``Copy'' page}\label{btn:step1_copy}
The folder is created, the visible images copied, the infrared images converted from RAW to DNG, the EXIF data read and the ground altitude computed (details in Section~\ref{btn:step1_validate}). While the infrared images are being converted, Step~2 is prepared in the background (thumbnails and target detection, visible then infrared): the timeline then opens without waiting. The window cannot be closed during the copy. At the end, \bouton{Flight timeline} opens Step~2; if the preparation is not finished, Step~2 resumes it where it left off.

The \texttt{mission\_parameters.json} file contains the same information as with the classic form.

\subsection{Buttons at the bottom of the window}\label{btn:step1_wizard}
\begin{description}
	\item[\bouton{Next}] (Images page) goes to the Mission page; the images of both cameras are required. \bouton{Create the mission and copy the images} (Mission page) starts the copy; \bouton{Flight timeline} (Copy page) opens Step~2.
	\item[\bouton{Back}] returns to the Images page; nothing is lost.
	\item[\bouton{Cancel}] closes the window without creating the mission.
	\item[\bouton{Classic form…}] (Images page) opens the old Step~1 (Section~\ref{sec:step1_classic}).
	\item[\bouton{?}] opens this chapter.
\end{description}

\section{Editing a mission: the ``Mission details'' window}\label{btn:step1_details}
Once Step~1 is done, a click on Step~1 in the main window no longer opens the mission creation, but the ``Mission details'' window (Figure~\ref{fig:step1_mission_details}). It edits the mission information without reading the SD cards again: location, description, pilot and pilot's ID, interval and clock offset of the cameras, infrared camera, filter. The date, time, position and visible camera come from the take-off image and do not change.

As soon as a piece of information is changed, \bouton{Close} becomes \bouton{Save and close}: it writes the changes to \texttt{FlightAnalytics/mission\_parameters.json} (closing the window does the same); the old file is kept as \texttt{mission\_parameters.json.bak}. The mission folder keeps its name, even if the location changes. \bouton{↺}, next to \bouton{?}, then appears: it puts back the information as it was when the window opened; the window stays open.

\figureIRD[fig:step1_mission_details]{step1_mission_details.png}{The ``Mission details'' window: the images are already in the mission, only this information changes.}{0.55}

\section{Classic form}\label{sec:step1_classic}
\legacybegin{Show the classic form (old Step 1)}
The \bouton{Classic form…} button of the ``Images'' page opens the old Step~1, described below. It takes place in two successive windows:
\begin{enumerate}
  \item \textbf{``Choose the image taken by the drone during takeoff''}: choose the first image taken by the drone
        at take-off. It provides the date, the time, the GPS position and the camera used.
  \item \textbf{``Create a mission''}: check and complete the mission information, then start the creation
        of the folder and the transfer of the VIS and NIR images.
\end{enumerate}
At the end, the mission becomes the mission loaded in the main window, which offers Step~2.

\paragraph{Prerequisites.}
\begin{itemize}
  \item An internet connection: the ground altitude is requested from IGN\textsuperscript{\ref{fn:ign}}
  (France) or from OpenTopoData (outside France),
        and the address of the place from OpenStreetMap.
  \item The drone images (visible) in DNG format, and the images of the near-infrared camera (SJCam M20, RAW files)
        copied to the computer's disk.
  \item A working folder (\texttt{C:/Air-Mission} by default), chosen at first launch and changeable
        with the ``gear'' button at the bottom of the main window. The \texttt{IRDRONE\_WORKING\_DIR}
        environment variable, if set, takes precedence.
\end{itemize}

% \figureIRD{step1_load_takeoff_1.png}{First window of Step 1 before loading the take-off image.}{0.8}

%----------------------------------------------------------
\section{\bouton{Load take-off image} button}\label{btn:step1_load_takeoff}
Opens a file browser to choose the image taken by the drone at take-off.
Only \texttt{.dng} files from the visible camera are accepted.

Once the image is chosen, the software:
\figureIRD{step1_load_takeoff_2.png}{First window of Step 1 after loading the take-off image.}{0.5}
\begin{enumerate}
  \item displays the image (decoding the DNG file takes a while: the progress bar shows the progress);
  \item reads the EXIF data: date and time of the shot, GPS position, manufacturer, model and serial number of the camera;
  \item queries IGN (or OpenTopoData outside France) to get the ground altitude \emph{above sea level};
  \item queries OpenStreetMap to get the place name, the municipality, the postal code, the department, the region and the country.
\end{enumerate}
This information is displayed in the grey frame below the image. When available, the place name, postal code, region and country are also saved in the file
\texttt{FlightAnalytics\textbackslash mission\_parameters.json} under the keys \texttt{Place name}, \texttt{Postal code}, \texttt{Region} and \texttt{Country}
(the municipality is saved under the key \texttt{Location}). Information that cannot be found is not saved: the corresponding key does not exist.
\begin{center}
\fbox{\parbox{0.9\linewidth}{\small\textbf{Warning.} The altitude recorded in the image by the drone is measured \emph{relative
to the take-off point}, not relative to sea level. This is why the ground altitude is requested from an external service such as IGN\textsuperscript{\ref{fn:ign}}.}}
\end{center}


\paragraph{If the image does not load.} Check the internet connection and that the file is indeed a \texttt{.dng} from the drone.
The \bouton{Next ▶} button stays disabled as long as no image has been loaded correctly.

%----------------------------------------------------------
\section{\bouton{Cancel}, \bouton{Next ▶} and \bouton{◀ Back} buttons}\label{btn:step1_navigation}
\begin{description}
  \item[\bouton{Next ▶}] is only available after the take-off image has been loaded. It opens the ``Create a mission'' window.
  \item[\bouton{Cancel}] closes the take-off image window without saving anything.
  \item[\bouton{◀ Back}] (``Create a mission'' window) returns to the previous window;
        the take-off image already loaded is kept.
\end{description}

%----------------------------------------------------------
\section{The ``Create a mission'' window}\label{btn:step1_form}
The fields are filled in automatically from the take-off image. The \textsf{Enter} key moves from one field to the next.
Grey fields are filled in automatically and cannot be edited.

\begin{center}
\begin{tabular}{@{}p{3.6cm}p{2.4cm}p{7.5cm}@{}}
\hline
\textbf{Field} & \textbf{Editable} & \textbf{Content} \\ \hline
Date, Hour & no & Date and time of take-off (EXIF). \\
Location & yes & Municipality found automatically. Part of the mission folder name. \\
GPS & no & Latitude and longitude of the take-off point. \\
Short description & yes & Short free-form description of the mission. \\
Pilot & yes & Name of the remote pilot and their ID number. \\
Camera VIS & partly & Manufacturer and model (read from the image, not editable); shooting interval (timelapse, in s) and clock offset (delta time, in s). \\
Image VIS format & no & Image format (DNG). \\
Camera NIR & yes & Default SJCam M20, RAW format, interval 3\,s, clock offset $-30$\,s. \\
NIR Filter & yes & Filter manufacturer and band, in nm (default 810\,nm). \\ \hline
\end{tabular}
\end{center}

\figureIRD{step1_param_mission.png}{``Create a mission'' window.}{0.5}

Note: the VIS and NIR clock offset parameters are reserved for experienced users. These parameters make it possible to ``recover'' a mission whose synchronization phase in the field turns out to be unusable (ArUco markers cannot be recognized, target rotation not started, etc.).

%----------------------------------------------------------
\section{\bouton{Create the mission} button}\label{btn:step1_validate}
Starts the creation of the mission. The button is disabled during processing; a label and a progress bar
show the current phase. The software runs the following, with no user intervention except at step~3:
\begin{enumerate}
  \item \textbf{Creation of the mission folder} and its subfolders, writing of \texttt{mission\_parameters.json},
        copy of the take-off image.
  \item \textbf{Transfer of the VIS images}: the \texttt{HYPERLAPSE\_XXXX.dng} files in the folder of the take-off image are copied
        to \texttt{AerialPhotography/VIS} under the name \texttt{VIS\_XXXX.dng}.
  \item \textbf{Choice of the NIR folder}: a message asks you to choose the first near-infrared image (RAW or JPG file
  \footnote{Each RAW image has an associated jpeg image. The SJCam names images by their shooting date. For example 2022\_0125\_130413\_001.RAW, 2022\_0125\_130415\_002.JPG, 2022\_0125\_130416\_003.RAW, 2022\_0125\_130418\_004.JPG, etc. The jpeg image is always saved after the RAW image, about 2s later. Note that in some circumstances (SD card too slow, long flight sequence) the time needed to save the RAW file exceeds the shooting interval (3s in our case). In that case the next image is taken 6s later. IRdrone handles this kind of gap automatically.} from the SJCam).
        The folder that contains it is used for all NIR images.
  \item \textbf{Conversion of the NIR images}: the RAW files are converted to DNG in \texttt{AerialPhotography/NIR}
        (parallel computation, possibly long).
  \item \textbf{\texttt{.exif} files} for the VIS then NIR images: one metadata file per image,
        and one \texttt{time\_line.json} file per band giving the chronology of the shots.
  \item \textbf{Altitudes and trajectory}: for each GPS position of the flight, the ground altitude is requested from IGN (in batches;
        three attempts on failure, then automatic switch to OpenTopoData). The software derives from it the height of the drone above ground,
        the distances and headings between shots, and saves these values in the \texttt{.exif} files.
  \item \textbf{Vertical profile}: plot of the drone and ground altitudes along the path, saved in
        \texttt{FlightAnalytics/fly\_vertical\_profil.png}.
  \item \textbf{Final message} giving the folder created, then the Step~1 windows close.
\end{enumerate}

\figureIRD[fig:takeoff_1]{step1_nir_prompt.png}{Message asking to choose the first NIR image.}{0.2}
\figureIRD[fig:fly_vertical_profil]{step1_fly_vertical_profil.png}{Vertical profile of the flight (drone and ground altitude), produced at the end of Step~1.}{0.8}

\subsection{Folders created}
At the end of Step~1, IRDrone has created the folder and file structure of the mission. From this point on you no longer need to access the original images stored on the cameras' SD cards. The structure of the mission folder is as follows:


\begin{verbatim}
FLY-YYYYMMDD-hhmm-<location>/
  AerialPhotography/
     VIS/   VIS_XXXX.dng, VIS_XXXX.exif, time_line.json,
            transfer_info_VIS_dng.json
     NIR/   NIR_XXXX.dng, NIR_XXXX.exif, time_line.json,
            transfer_info_NIR_dng.json
  FlightAnalytics/
     mission_parameters.json, fly_vertical_profil.png
  Synchro/        (Step 3)
  ImgIRdrone/     (Step 4)
  mapping_MULTI/  (Step 6)
  cameras/
\end{verbatim}
The subfolders have a specific icon in the Windows file explorer.
The mission folder is never deleted: if the folder already exists, only the missing items are created.
For VIS images already transferred (\texttt{.exif} file present), the transfer is skipped: running Step~1 again on
an existing mission is therefore fast.

\figureIRD{Step1_mission_folder.png}{Mission folder seen in the Windows file explorer.}{0.7}

The duration of the transfers depends a lot on the performance of the storage media (SD card, hard disk) and of the connections (USB B or C, card reader). Moreover, unlike the images of the DJI camera, which are in dng format, the images of the SJCam M20 camera are in a non-standard ``RAW'' format. Each image must be converted to Adobe's dng format with third-party software \footnote{This is the \texttt{sjcam\_raw2dng} code, available only as an exe file. \texttt{sjcam\_raw2dng}
	Convert SJCAM raw files to Adobe DNG and TIFF while also getting metadata
	such as date and ISO from JPEGs that correlate to the RAW files
} which lengthens the actual transfer time.

\subsection{Troubleshooting}
\begin{itemize}
  \item \textbf{No NIR image found}: the message ``No NIR RAW files found'' means that the chosen folder contains no
        \texttt{.RAW} files; start again and choose the right folder.
  \item \textbf{Error during the VIS transfer or the EXIF reading}: an error message is displayed and the window closes.
        The button can be used again after reopening the window.
  \item \textbf{IGN altitudes unavailable}: the software switches by itself to OpenTopoData, less accurate in France.
\end{itemize}
%\todo{Add the errors encountered in the field (missing SD card, write-protected folder, \dots).}

%==========================================================
\legacyend

\chapter{Step 2 -- Flight timeline}\label{ch:step2}
\section{Principle}
Step~2 splits the flight into phases: \emph{synchronization} (hovering above the rotating ArUco target) and \emph{survey} (flight over the study area). These phases are now \textbf{detected automatically}; the user checks them and adjusts them if needed (Figure~\ref{fig:step2_flight_timeline}).

\figureIRD[fig:step2_flight_timeline]{step2_flight_timeline.png}{The ``Flight timeline'' window. In the center, the preview of the visible and infrared images under the mouse; on the right, the flight map with the ground footprints of the images kept for the map (Section~\ref{btn:step2_best_mapping}), then the four key points with their actual images; at the bottom, the timeline: altitude profile (green dots: images kept for the map), phases (purple: synchronization, green: survey), target detections and strips of visible and infrared images.}{0.95}

\subsection*{What the software does}
\begin{enumerate}
	\item Each DNG file (DJI and SJCam alike) contains a JPEG preview of about 1000 pixels. It is read in a few milliseconds: the whole mission is scanned in a few seconds, without RawTherapee. The previews are kept in \texttt{Synchro/Thumbnails} (cache, recreated if needed).
	\item \emph{Visible}: the altitude recorded by the drone gives the initial hover (synchronization) and the longest plateau, the farthest from the take-off altitude (survey). A take-off from a summit, followed by a descent towards the study area, is handled in the same way as a climb.
	\item \emph{Synchronization}: the angle of the moving target (ID 0) is measured on the previews, only during the hover for the visible, and for the infrared from the first image until the target disappears (the drone has moved away). The images where the target turns steadily in one direction form the synchronization window; the images where the turntable is still stopped, where it changes direction, and the take-off are discarded.
	\item \emph{Infrared}: the clock of the NIR camera is often wrong (time zone, drift). Since the same target rotation is seen by both cameras, comparing the two angle curves gives an approximate offset (on the order of a second) that places the infrared images on the visible timeline. The infrared synchronization window thus corresponds to the same target revolution as the visible window.
	\item The infrared survey window \textbf{follows} the visible survey window (hatched area): there is no need to look at the infrared images. After Step~3, the precise offset updates this window.
\end{enumerate}

\subsection*{Check and adjust}\label{btn:step2_adjust}
\begin{itemize}
	\item Hovering over the timeline shows, enlarged, the visible image and the infrared image taken at the same moment, with the altitude and the state of the target (``target turning'', ``target still'', ``no target''). Leaving the mouse on this preview displays it over the whole window to see the details.
	\item The dots on the ``target'' line show the analyzed images: purple, the target is turning; grey, the target is seen but still; empty circle, target not seen.
	\item Drag a phase edge to move it (it snaps to an image), or click a key point (right-hand column) then use \bouton{◀} \bouton{▶} or the keyboard arrows.
	\item Mouse wheel: zoom; drag the background: pan; double-click: the whole flight.
	\item \bouton{↺}, next to \bouton{?}, appears once a phase is moved: it cancels the manual changes and returns to the automatic proposal; the window stays open.
	\item The tool tip of the \bouton{?} button recalls these controls.
\end{itemize}
If the target is not found (target too small, image too dark), an orange banner says so: the visible synchronization window is then the hover, and the infrared window must be placed by hand.

\section{Best mapping: the images the map needs}\label{btn:step2_best_mapping}
A survey photographed every two seconds gives far more overlap than the map needs. Aligning all these images in Step~4 takes hours, and feeding hundreds of nearly identical images to ODM slows down the computation of the map without improving it, often the opposite. The \bouton{Best mapping: keep only the images the map needs} check box, ticked by default, keeps only the survey images whose \emph{ground footprints} overlap by the value chosen on its right (\bouton{75~\% overlap} by default, the recommended value for ODM; more overlap keeps more images).

The choice shows at once:
\begin{itemize}
	\item on the timeline, green dots on the altitude profile and green marks at the top of the visible image strip show the images kept; the cursor bubble and the preview caption say ``kept for the map'';
	\item to the right of the previews, the \textbf{flight map} shows the path (purple: synchronization, green: survey), the ground footprints of the images kept (green rectangles: their overlap shows) and, in orange, the footprint of the image under the mouse with the image itself. Clicking a point of the map shows that image; wheel: zoom, drag: pan, \bouton{Layers}: background map (Section~\ref{btn:home_map});
	\item below the timeline, the number of images kept, for example ``16 of 120 survey images kept for the map · 87~\% less to align''.
\end{itemize}
The selection follows the edges of the survey phase when they are moved.

\textbf{How the images are chosen.} The position, height above the ground and heading of each image are read from the \texttt{.exif} files written in Step~1; the field of view of the camera gives the size of the ground footprint. The first survey image is kept. For each next image, IRDrone estimates the overlap of its footprint with the footprint of the last image kept, from the displacement of the drone (along the image axis and across it). As long as it stays above the chosen value, the image is only a candidate; as soon as it drops below, the previous candidate is kept: two successive images kept always overlap by at least the chosen value. The last survey image is kept as well. For a displacement along the image axis, this is the \texttt{--selection best-mapping} criterion of the former \texttt{run.py} command line.

\textbf{What about synchronization?} In Step~2 the precise offset between the two cameras is not known yet: the preview covers all the survey images. In Step~4, once the cameras are synchronized, the same selection is made again among the well synchronized pairs only (Section~\ref{btn:step4_best_mapping}); the final count can therefore differ a little from the preview. The choice (check box and overlap) is saved by \bouton{Continue with synchronisation} and \bouton{Close} in \texttt{FlightAnalytics\textbackslash best\_mapping.json}.

If the images have no GPS position, the map is not shown and all the survey images are kept.

\section{\bouton{Continue with synchronisation} and \bouton{Close} buttons}\label{btn:step2_confirm}
Both buttons save the phases in the files \texttt{transfer\_info\_VIS\_dng.json} and \texttt{transfer\_info\_NIR\_dng.json} (read by Steps 3, 4 and the mapping) and in \texttt{FlightAnalytics/flight\_timeline.json}. Reopening Step~2 brings back the confirmed phases.
\bouton{Continue with synchronisation} then opens Step~3; nothing starts by itself: the synchronization is started from its window (Chapter~\ref{ch:step3}). \bouton{Close} only closes the window: Step~3 will be run later from the main window. It reads \bouton{Save and close} when the phases or the best mapping choice are changed, or were never saved: closing the window then saves them too. If a phase is missing, closing asks whether to close without saving.

\section{Classic selection}\label{btn:step2_classic}
\legacybegin{Show the classic selection (old Step 2 window)}
The \bouton{Classic selection…} button opens the old manual selection of the five reference images of each camera, described below.

\figureIRD{step2_dialog_select_phase_images_1.png}{Reference image selection window. For the visible spectral band, the first image is by default the take-off image that was used to define the mission in Step~1. But it can be changed here.}{0.95}
The dialog window lets you define:
\begin{enumerate}
	\item the take-off point,
	\item the start and end points of the synchronization,
	\item the start and end points of the flight over the area studied by the mission.
\end{enumerate}
The images are selected for the visible and near-infrared spectra in the folders
\begin{itemize}
	\item \texttt{C:/Air-Mission/FLY-YYYYMMDD-hhmm-location/AerialPhotography/VIS},
   	\item  \texttt{C:/Air-Mission/FLY-YYYYMMDD-hhmm-location/AerialPhotography/NIR} respectively.
\end{itemize}
 During this selection, it is always possible to choose another image after a first selection.
 The order in which the five images are selected does not matter, but of course first always precedes last.
 The same image can be selected for \texttt{first mission image} and \texttt{first sync image}.
  For DJI images of the visible spectrum, if you select the same image for \texttt{first sync image} and \texttt{last sync image}, then a single image is kept (this is a borderline case for synchronization).
  However, if you select the same image for \texttt{first sync image} and \texttt{last sync image} for the near-infrared spectrum, the automatic synchronization will fail in Step~3 and you will be asked to enter the clock synchronization offset manually.
  This can be used as a ``trick'' to force the clock synchronization offset to a particular value when the synchronization in the field was problematic.
%\todo{\bouton{Load first mission image.}, \bouton{Load first sync image.}, \bouton{Load last sync image.},
%----------------------------------------------------------
\section{\bouton{Load first mission image.} button}\label{btn:step2_Load_first_mission_image}
The choice of this image defines the origin of the timeline for the synchronization of the camera clocks. As a rule, choose the first VIS shot of the mission.
\section{\bouton{Load first sync image.}, \bouton{Load last sync image.} buttons}\label{btn:step2_Load_synchro_image}
The user must choose the image at the start of the synchronization cycle. For the visible images, if the procedure was followed in the field, this is the first or second image recorded. For the infrared images, find the first stable image where the whole target system is visible and where the moving target is actually rotating. Then choose the last image of the synchronization phase. It is usually easy to spot, because when the synchronization phase is over the pilot climbs quickly and the change of appearance is clearly visible. The choices made here partly determine the success of the clock synchronization in Step~3. But after Step~3 it is possible to go back to Step~2 and redefine these images.
\figureIRD[fig:step2_dialog_select_VIS]{step2_dialog_select_phase_image_VIS.png}{Reference image selection window. Visible spectral band (VIS).}{0.95}
\section{\bouton{Load first fly image.}, \bouton{Load last fly image.} buttons}\label{btn:step2_Load_mission_image}
The user must choose the first image taken above the area to study (mission flight or hyperlapse survey). This discards the images of the approach between the take-off/synchronization at very low altitude and the mission flight at ``high altitude'' (hyperlapse survey). Note that in principle the mission flight should be flown at constant drone altitude and over an area whose altitudes vary little compared with the height of the drone above the ground. Looking at Figure \ref{fig:fly_vertical_profil} lets you judge the quality of the flight plan carried out in the field.
\section{\bouton{Validate mission images} button}\label{btn:step2_validate}
\bouton{Cancel}, on its left, closes the window (after a confirmation if the images are not validated). This button builds the lists of images of the three flight phases relevant to IRDrone: takeoff, synchro and fly. No image is moved. The operation is very fast.
\figureIRD[fig:step2_dialog_select_NIR]{step2_dialog_select_phase_images_NIR.png}{Reference image selection window. Near-infrared spectral band (NIR).}{0.95}

%==========================================================
\legacyend

\chapter{Step 3 -- Camera synchronization}\label{ch:step3}
\section{Principle}
To combine different spectral bands, the images must be taken at the same moment. Unfortunately this is not technically possible with the hardware used in the IRdrone project. This is an initial choice linked to the affordable nature (\emph{inexpensive and low cost}) of the project. Each camera has its own clock, which defines the chronology of the images. The images can be placed on a line representing time. This is the ``time line'' of the camera. The two time lines, VIS and NIR, must then be synchronized. Setting the clocks to the right time is not enough: on the one hand they only give points every second, on the other hand the exact moment when the operator in the field started the cameras cannot be controlled. Markers must therefore be detected in the images. In the film industry, the film and the soundtrack are synchronized with a ``clapperboard''. This solution was considered in the first versions of IRdrone, which used video instead of images. But video raised many quality problems (resolution, ``jelly'', ...) and was abandoned. However, the principle of a moving object in the images at the start of the flight was kept. The idea is to locate an object (automatically) and to track its trajectory on both cameras. Making the two trajectories coincide synchronizes the clocks. An accuracy of about 0.05s can be reached, which is more than enough in our case.
IRdrone uses as moving object an ArUco target
\footnote{This type of target is designed to be detected automatically by software, for example with OpenCV. IRdrone uses the \texttt{DICT\_4X4\_50} dictionary: the moving target has ID 0 and the four fixed targets have IDs 5, 22, 16 and 13.}
mounted on a turntable
\footnote{The turntable is driven by a stepper motor controlled by an Arduino. The system is programmed to obtain a perfectly constant rotation speed. It is set to 1 rpm (one revolution per minute) so that during the synchronization flight phase the cameras record at least about ten different angular positions.} as shown in Figure \ref{fig:step3_syncro_aruco_1}.

\figureIRD[fig:step3_syncro_aruco_1]{step3_syncro_aruco_1.pdf}{Device for synchronizing the clocks of the VIS and NIR cameras. The principle relies on detecting the angular position of ArUco target 0, which rotates at constant speed. The small fixed targets define an absolute reference frame that removes the effect of unwanted drone movements, especially rotations.}{0.95}

The angular position of the target is measured. But this requires a fixed reference frame. The image edges were used initially. However, they are only reliable if the drone does not rotate. This requires a very calm area and the pilot must not touch the rotation stick. To remove this constraint in the field, 4 fixed ArUco targets were added to the device
\footnote{In theory 2 targets are enough to define a reference axis. Using 4 targets provides some redundancy. Note that the software tries to synchronize in the absolute frame, but if it fails it synchronizes in the relative frame tied to the image edges (less accurate). If that also fails, as a last resort the software asks the user to enter the synchronization offset value.}. They define an absolute reference frame that improves the accuracy of the device. Obviously this synchronization phase must be carried out in the field with the greatest care. This is the condition for combining moderately priced hardware with good-quality multispectral images.

\subsection{The three phases of Step 3}
Step~3 runs in three successive phases, chained automatically (Section~\ref{sec:step3_chained}). In manual mode, each one is started by a button that is enabled only once the previous phase is finished (Figure~\ref{fig:step3_window}):
\begin{enumerate}
  \item \emph{Target detection} (\bouton{Run ArUco detection (VIS + NIR)} button): on each synchronization image, the angle of the moving target is measured (Section~\ref{btn:step3_aruco});
  \item \emph{Time offset computation} (\bouton{Compute time alignment} button): the offset $\Delta t$ that makes the two series of angles coincide is searched for (Section~\ref{btn:step3_time_shift});
  \item \emph{Pairing} (\bouton{Pair VIS / NIR images} button): each VIS image is associated with the NIR images closest in time (Section~\ref{btn:step3_pairing}).
\end{enumerate}

\subsection{Automatic chaining}\label{sec:step3_chained}
Step~3 never starts by itself: whether it is opened from Step~2 (\bouton{Continue with synchronisation}) or from the main window, the \bouton{Synchronise the cameras} button (dark blue, above the progress bar) runs the three phases in a row (Section~\ref{btn:step3_run}). A progress line shows the current phase and its result: number of targets seen, offset, number of pairs (Figure~\ref{fig:step3_chained}). A single question is asked, after the clock alignment (Section~\ref{btn:step3_overlap}).

Once the images are paired, nothing starts by itself: \bouton{Continue with image processing} opens Step~4 when you are ready (Section~\ref{btn:step3_continue}).
If the offset cannot be measured, manual entry is offered and the chaining stops.

\figureIRD[fig:step3_chained]{step3_chained.png}{Chained Step 3: targets detected, clocks aligned; the question ``Do the two curves overlap?'' waits for the answer before pairing.}{0.7}

When reopened once the step is done (click on Step~3 in the main window), the window recomputes nothing: it shows the saved result, that is the clock offset, the number of pairs and the aligned curves (Section~\ref{btn:step3_before}). \bouton{Synchronise again} redoes it; \bouton{Manual mode…} shows the buttons of the three phases.

\section{The question ``Do the two curves overlap?''}\label{btn:step3_overlap}
After the clock alignment, the window shows the target angle as a function of time: green curve for the visible, red for the infrared, the infrared curve being shifted by the computed offset.
\begin{description}
	\item[\bouton{Yes} (Y key)] the two curves overlap: the offset is correct, the images are paired (Section~\ref{btn:step3_pairing}).
	\item[\bouton{No} (N key)] opens the manual adjustment (Section~\ref{btn:step3_adjust}) to slide the infrared curve onto the visible curve; the question is then asked again with the new offset.
\end{description}
Well-overlapping curves have green dots on the red line, sawtooth included. The examples in Section~\ref{sec:step3_examples} show a good and a bad case.

\section{\bouton{Synchronise the cameras} button}\label{btn:step3_run}
The main action of Step~3: it runs the target detection, the clock alignment and, once you have answered the question on the curves, the pairing. It becomes \bouton{Synchronise again} when the step was done earlier.

\section{\bouton{Manual mode…} button}\label{btn:step3_manual}
Stops the chaining and shows the buttons of the three phases (Sections~\ref{btn:step3_aruco} to \ref{btn:step3_pairing}), to be run one by one, for example to correct the offset with \bouton{Adjust alignment...} before pairing.

\section{\bouton{Enter the clock offset…} button}\label{btn:step3_clock_offset}
Advanced option, for a flight without the ArUco target (former synchronisation by the dates). The button is shown only if the box \emph{Advanced – Step 3: let me enter the camera clock offset by hand} is checked in the options (\bouton{⚙} gear wheel at the bottom of the main window, Section~\ref{btn:home_workdir}). With this option, Step~3 does not start by itself when coming from Step~2: you choose between \bouton{Synchronise the cameras} (flight with the target) and this button.

A small window asks by how many seconds the infrared camera clock is ahead of the drone clock (negative if it is behind, 0 if both clocks were set to the same time). This is the delay of the former synchronisation. It shows the resulting offset $\Delta t$ of the infrared time line, computed from the dates of the first images. \bouton{Pair the images} skips the target detection and runs the pairing directly (Section~\ref{btn:step3_pairing}). No curve can be checked: the quality of the pairing depends on the accuracy of the delay entered.

\section{\bouton{Continue with image processing} button}\label{btn:step3_continue}
Appears at the bottom right, next to \bouton{Close}, once the images are paired. It closes the window and opens Step~4. The alignment, which is long, does not start there by itself: it is started with \bouton{Run alignment}.

\section{When Step~3 fails}\label{sec:step3_problems}
When something goes wrong, a window says \emph{what happened} and \emph{how to fix it} (numbered), with the technical message under \bouton{Details}; the same text is written in the console. Its buttons: \bouton{Check the installation} (Section~\ref{btn:home_check}), \bouton{Open the images of the target} (\texttt{Synchro/VIS}, \texttt{Synchro/NIR}: the images as developed for the detection; \texttt{Synchro/Check\_ARUCO}: the markers found, drawn), \bouton{Create a diagnostic report…} of this mission (Section~\ref{btn:home_diagnostic}). The usual causes:
\begin{itemize}
  \item the list of the images or the time line of a camera is missing: Step~1 was not completed, run it again;
  \item no synchronisation window, or an empty one: place it in Step~2 on the images where the target is filmed turning;
  \item RawTherapee is missing, or develops nothing: the window says so (instead of ``not enough points''); the installation check tells why;
  \item the offset could not be measured (target found on too few images): the window gives, for each camera, how many synchronisation images there are, on how many the target was found, and why the others were left out (missing, not developed, without any marker, not in the time line). If the target is not on the developed images, the synchronisation window of Step~2 is on the wrong images; if it was not filmed, \bouton{Enter the offset by hand…}.
\end{itemize}
An image that is not in the time line of its camera is left out (it used to be counted at $-9999$~s, which made the offset wrong).

\section{\bouton{Before alignment} button}\label{btn:step3_before}
When Step~3 is reopened once done, the window shows the curves after alignment (\texttt{Synchro/check\_time\_alignment.png}). \bouton{Before alignment} shows the curves before correction (\texttt{check\_time\_before\_alignment.png}); a second click returns to the aligned curves.

\subsection{Principle of the computation}
The angle of the moving target is measured on each image, then unwrapped (see Section~\ref{btn:step3_time_shift}). Since the rotation speed $\omega$ is constant, the angle is an affine function of time. The NIR time line, which has at least two points, provides the reference line $\theta = \omega t + b_{NIR}$. Each VIS point provides an estimate of the intercept $b_{VIS} = \theta_{VIS} - \omega t_{VIS}$; their median is taken. The offset is
\begin{equation}
  \Delta t = \frac{b_{NIR} - b_{VIS}}{\omega},
\end{equation}
and the corrected NIR time line is $t_{NIR} + \Delta t$. This is why a single VIS image is enough in theory (case~2 below), provided the rotation speed is constant and the absolute frame is used.

\subsection{Synchronization examples}\label{sec:step3_examples}
The images of the synchronization phase come from Step~2. A few cases are presented here.

\subsubsection{Case 1: good synchronization}
Figures~\ref{fig:synchro_cas1} (\subref{fig:synchro_cas1_avant} and \subref{fig:synchro_cas1_apres}) show a good synchronization. The pilot manually started the hyperlapse of the NIR camera, then took off with the drone, which automatically started the rotation of the target. When the pilot placed the drone above the target, the NIR camera was recording the rotating target. The images between the moment the NIR hyperlapse was started and the moment the drone was stabilized above the target were ignored in Step~2. Then the pilot set the hyperlapse parameters of the VIS camera and started it. This usually takes about thirty seconds for a trained pilot. As soon as the VIS hyperlapse begins, it records the moving target. From that moment, with a recording interval of 2s, the pilot only needs to hold the drone above the target for less than about ten seconds to get 4 to 5 points, before climbing quickly to reach the study area (hyperlapse survey). In Step~2, the software user must make sure to discard the NIR images taken before the drone was stabilized above the synchronization device and, where applicable, those that after stabilization correspond to a target not yet rotating.

\figurePaireIRD{fig:synchro_cas1}
  {Step3_time_alignment_failed_cas1_0.png}{Target trajectory before synchronization.}
  {Step3_time_alignment_failed_cas1_1.png}{Successful synchronization.}
  {Case of good synchronization.}

\subsubsection{Case 2: borderline case}
Figures~\ref{fig:synchro_cas2} (\subref{fig:synchro_cas2_avant} and \subref{fig:synchro_cas2_apres}) are a borderline case for synchronization. The pilot stayed above the target for a very short time and only one image of the target is available in the VIS images. Here the synchronization works anyway because the rotation speed is kept constant and the absolute frame could be used.

\figurePaireIRD{fig:synchro_cas2}
  {Step3_time_alignment_failed_cas2_0.png}{Target trajectory before synchronization.}
  {Step3_time_alignment_failed_cas2_1.png}{Borderline synchronization case.}
  {Borderline synchronization case.}

\subsubsection{Case 3: synchronization failure}
Figures~\ref{fig:synchro_cas3} (\subref{fig:synchro_cas3_avant} and \subref{fig:synchro_cas3_apres}) show the trajectory of the target detected on both cameras. Here the synchronization fails because the first images kept for the NIR camera correspond to moments when the target is not rotating.

\figurePaireIRD{fig:synchro_cas3}
  {Step3_time_alignment_failed_cas3_0.png}{Target trajectory before synchronization.}
  {Step3_time_alignment_failed_cas3_1.png}{Synchronization failure.}
  {Case of synchronization failure.}

\figureIRD[fig:step3_window]{step3_window.png}{Synchronization window.}{0.8}

\section{\bouton{Run ArUco detection (VIS + NIR)} button}\label{btn:step3_aruco}
This button starts the detection of the ArUco targets on all the images of the synchronization phase chosen in Step~2, first for the VIS band then, automatically, for the NIR band. A progress bar shows the progress. The computation is spread over several processor cores to speed up processing.

\subsection*{What the software does for each image}
\begin{enumerate}
  \item The DNG images are converted to TIF by RawTherapee with a development profile dedicated to detection (\texttt{*\_RawTherapee.tif} files in the \texttt{Synchro\textbackslash VIS} and \texttt{Synchro\textbackslash NIR} folders of the mission). This conversion is done only once: if the file already exists, it is reused.
  \item The visible ArUco targets are detected with OpenCV. For each one, its ID, the position of its center and the angle of its main axis in the image are kept.
  \item The moving target (ID~0) gives the \emph{relative angle} (image frame, \texttt{angle\_img}).
  \item If the four fixed targets (IDs 5, 22, 16 and 13) are also detected, they define a ground-fixed frame: two axes obtained by averaging the opposite sides of the quadrilateral formed by their centers. The angle of the moving target is then measured in this frame. This is the \emph{absolute angle} (\texttt{angle\_abs}), insensitive to rotations of the drone around its vertical axis.
  \item If the moving target or one of the fixed targets is missing, the image is reported in the log. The absolute angle is not computed (degraded mode: only the relative angle is kept if the moving target is visible); if the moving target is not detected, the image is ignored.
  \item Each image is associated with its relative date (in seconds) on the time line of its camera.
\end{enumerate}
The detection results are stored in the metadata of each image.

\subsection*{Check images}
During detection, each processed image is saved with a graphic overlay in the \texttt{Synchro\textbackslash Check\_ARUCO} folder. The window displays one regularly (about every 5 seconds) so that the user can judge the quality of the detection. The overlay shows:
\begin{itemize}
  \item the centers of the fixed targets (red disks);
  \item the center of the moving target;
  \item the orientation of the moving target in the image (green line);
  \item the orientation projected in the absolute frame (dashed orange line) and the two axes of the absolute frame (black arrows) when the four fixed targets are detected;
  \item a legend: relative angle, absolute angle and the difference between the two, or the message ``absolute reference: unavailable'' if the absolute frame could not be built.
\end{itemize}
\textbf{Tip:} browsing this folder in the Windows file explorer is the best way to spot doubtful detections before moving on to the next phase.

% Figure to be created by the user (PowerPoint): 2 rows of 4 images (VIS on top, NIR at the bottom) taken from Synchro\Check_ARUCO.
\figureIRD[fig:check_aruco]{Step3_check_aruco_VIS_NIR.pdf}{Examples of check images (\texttt{Synchro\textbackslash Check\_ARUCO} folder). Top row: VIS band; bottom row: NIR band. The moving target is marked by its green axis; the centers of the fixed targets by red disks.}{0.95}

\section{\bouton{Compute time alignment} button}\label{btn:step3_time_shift}
This button, enabled only when detection is finished, computes the time offset between the two cameras.

\subsection*{Choice of angles}
By default the software uses the absolute angles. At least 2 NIR images and 1 VIS image with an absolute angle are required. Otherwise, it uses the relative angles (image frame), with the same minimums, which is less accurate if the drone rotated during synchronization. If neither solution is possible, the message ``Insufficient data to automatically align VIS and NIR camera timelines'' is displayed and the software asks the user to enter the offset in seconds (default value: 0). In that case no plot is produced.

\subsection*{Computation}
\begin{enumerate}
  \item The angles of each band are ``unwrapped'': the $\pm 360^\circ$ jumps are removed to obtain a continuous curve that increases (or decreases) steadily with time. Unwrapping is done in the same way for VIS and NIR, starting from the last image, that is the image where the target is normally rotating.
  \item A line is fitted to the NIR points: its slope is the rotation speed $\omega$.
  \item From it and the VIS points, the offset $\Delta t$ is derived (see the principle at the beginning of the chapter). The estimate uses the median of the VIS points to be robust to outliers.
\end{enumerate}
The result is displayed in the window log:
\begin{itemize}
  \item the estimated offset $\Delta t$ (in seconds);
  \item the rotation speed measured on NIR and on VIS (in rpm). It should be close to the nominal value of 1~rpm and the two values should be close; the relative difference is displayed;
  \item the number of VIS and NIR points used;
  \item the quality of the NIR fit ($R^2$, close to 1);
  \item the spread of the individual offsets (interquartile range and standard deviation, in seconds).
\end{itemize}

\subsection*{Check plots}
Two plots are saved in the \texttt{Synchro} folder of the mission. During the question ``Do the two curves overlap?'', the window shows the curves after alignment; in manual mode, it shows the two plots alternately (switching every 1.5 s):
\begin{itemize}
  \item \texttt{check\_time\_before\_alignment.png}: target angle as a function of time for VIS (green) and NIR (red) before correction;
  \item \texttt{check\_time\_alignment.png}: the same curves after shifting the NIR time line by $\Delta t$.
\end{itemize}
If the synchronization is good, the green dots lie on the line of the red dots. These two plots are the ones shown in the three cases of this chapter.

\textbf{What to do in case of failure?} If the VIS and NIR points do not overlap after alignment, or if the VIS and NIR rotation speeds are very different, go back to Step~2 and redefine the synchronization images (in particular the first NIR image, which must show the target actually rotating), then run Step~3 again.

\section{\bouton{Adjust alignment...} button}\label{btn:step3_adjust}
Opens the manual offset adjustment, also opened by answering \bouton{No} to the question ``Do the two curves overlap?''. The window plots the visible curve (green), the infrared curve at the computed offset (red, dashed) and the infrared curve at the adjusted offset (red). The slider and the numeric box move the infrared curve: keyboard arrows in steps of 1\,ms, Page Up / Page Down keys in steps of 0.1\,s.
\begin{description}
	\item[\bouton{Save and close}] (\bouton{Close} while the offset is unchanged; closing the window does the same) keeps this offset: the check plots are redone. If the images were already paired, the pairing must be redone.
	\item[\bouton{↺}] (bottom left, shown when the offset differs from the computed one) returns to the offset computed from the targets.
\end{description}
If the offset could not be computed, a small window asks directly for its value in seconds (added to the infrared dates).

\section{\bouton{Pair VIS / NIR images} button}\label{btn:step3_pairing}
This button is enabled once the time offset is known (computed or entered by hand). It builds the list of image pairs that will be aligned in Step~4.

\begin{enumerate}
  \item The VIS time line is the reference. All NIR dates are corrected: $t'_{NIR} = t_{NIR} + \Delta t$.
  \item For each VIS image, the \emph{two} NIR images closest in time are kept. For each one, the signed difference $dt = t'_{NIR} - t_{VIS}$ is stored: $dt<0$ if the NIR precedes the VIS, $dt>0$ if it follows it.
  \item The flight metadata (GPS position, UTM coordinates, altitudes, yaw, pitch and roll angles of the drone and of the gimbal, distances, heading, etc.), known for the VIS images, are interpolated at the corrected times of the NIR images. Angles (yaw, heading) are interpolated along the shortest angular path. Outside the VIS time range, the value at the end point is kept.
  \item The result is saved in the \texttt{AerialPhotography\textbackslash shoot\_points.json} file, used by the following steps. The log gives the number of pairs created.
\end{enumerate}
The same NIR file can be used for several VIS images (for example if the NIR shooting rate is higher) and no maximum time difference is imposed: the quality of the pairing therefore depends on that of the synchronization.

\section{\bouton{Close} button}\label{btn:step3_close}
Closes the window and returns control to the main dialog. The files produced (plots, check images, \texttt{shoot\_points.json}) remain in the mission folder. The \bouton{?} button opens this help at the page of this chapter.

%==========================================================
\chapter{Step 4 -- Image processing}\label{ch:step4}
\section{Principle}
For a VIS-NIR pair, image processing consists of aligning the two images so that the Red, Green, Blue and NIR spectral bands overlap. This is difficult for several reasons:
\begin{itemize}
	\item Even when the time lines are synchronized, the shooting times are close but not identical.
	\item The two cameras do not point at the same spot on the ground. The VIS camera of the DJI drone is mounted on a very effective gimbal that keeps it pointing at the nadir. Note that the hyperlapse flight is done in ``gimbal lock'' mode, in other words the roll angle of the gimbal stays fixed. Only its pitch and yaw angles change to keep it pointing at the nadir ($90^{\circ}$ downwards). The NIR SJCam M20 camera, on the other hand, is rigidly fixed under the body of the drone, so where it points depends entirely on the attitude of the drone (pitch and yaw angles in particular).
	\item The focal lengths of the two cameras are different. The VIS camera has a field of view (FOV) of $84^{\circ}$ in the diagonal ($68^{\circ}$ horizontally), whereas the NIR camera has an ultra-short focal length with a horizontal field of view above $110^{\circ}$, which is almost a ``fish eye'' (\S\ref{sec:ann_camera}).
	\item The lenses show distortion and vignetting. For distortion, both cameras were calibrated beforehand with Zhang's method.
\end{itemize}



In practice, the images are first converted from raw-dng to tif (with the third-party software RawTherapee
\footnote{With the open source third-party software RawTherapee. The operation is costly the first time, but the image is cached for later runs.}).  The distortion is removed from both images. The theoretical shifts needed to bring the NIR aiming center onto the VIS aiming center are computed. These shifts are expressed as pitch, yaw and roll angles. They are used to initialize the alignment procedure. Very roughly, the alignment runs in two stages. First, the parameters of a homography to apply to the NIR image so that it ``matches'' the VIS image are determined. This is a complex minimization process developed specifically for multispectral images
	\footnote{The usual techniques for aligning two images, which work extremely well for a single spectral band or for RGB, such as ``sift'', fail.}.
This first alignment (coarse process) is followed by a fine alignment (fine process) that applies local homographies to lower the alignment criterion further. The algorithm is described in the appendix (Section~\ref{sec:ann_algorithme}). The user should be aware that the method is not foolproof and, above all, that it depends on the flight conditions in the field.

\figurePaireIRD{fig:step4_pair_VIS_NIR}
{step3_img_VIS.png}{Visible spectrum (VIS).}
{step3_img_NIR.png}{Near-infrared spectrum (NIR).}
{Pair of images of a shoot point. The two spectral bands record different information. For example, the surface of the river is completely black in near infrared. Water does not reflect wavelengths above $750\,nm$.}

\section{Step 4 dialog}
The Step~4 window (Figure~\ref{fig:step4_window}) shows the name of the mission (the full path is in the tool tip) and the number of image pairs. From top to bottom it contains: the output options, the choice of the images aligned with the flight map, the button \bouton{Run alignment} (\bouton{Stop} during the computation), the status and progress bar of the computation, the log, then, at the bottom, \bouton{?} and \bouton{⚙} (settings) on the left, \bouton{Close} and, once images are aligned, \bouton{Continue with the review} on the right. The window adapts to the screen size.

\figureIRD[fig:step4_window]{step4_window.png}{Window for processing the ``shoot points'' (VIS-NIR image pairs), before the computation.}{0.7}


\section{Output options}\label{btn:step4_options}


\subsection{\bouton{first VIS image number} and \bouton{last VIS image number} fields}
These two fields select the shoot points to process. By default IRDrone offers the range defined in Step~2. It is possible to process a single image by entering the same number for the first and last images, provided that the \texttt{Maximum VIS/NIR time difference} is compatible. If in doubt, set the \texttt{Maximum VIS/NIR time difference} to $3 \, s$.

By default the alignment procedure produces, for

$\texttt{first VIS image number} \leq XXXX \leq \texttt{last VIS image number}$:
\begin{itemize}
	\item the quadruplets of multispectral images red (1), green (2), blue (3), nir (4) (\texttt{IRD\_XXXX\_n.tif}) placed in the \texttt{ImgIRdrone\textbackslash imgMultispectral} folder.
	\item the images \texttt{VIS\_XXXX.jpg}, \texttt{VIR\_XXXX\_local.jpg}.
	\item the files \texttt{VIS\_XXXX\_REDO.bat} and two technical files \texttt{VIS\_XXXX\_motion\_model.json}, \texttt{VIS\_XXXX\_motion\_model.npy}.
	\item the file \texttt{ImgIRdrone\textbackslash shoot\_points.json}, which contains a dictionary of the processed shoot points with, in particular, the key \texttt{"failed"} whose value is set to \texttt{false} by default. The user checks this assumption in the review of the aligned images (Chapter~\ref{ch:review}).
\end{itemize}


\subsection{\bouton{NDVI} and \bouton{NIR} check boxes}
By default the alignment procedure produces the VIS and VIR images in jpeg format. This is useful for checking visually that the alignments succeeded. By ticking the \texttt{NDVI} and/or \texttt{NIR} boxes, the NDVI and NIR images can also be generated in jpeg format. Only tick these boxes if you plan to use these representations directly in a report. Note that in Step~6 you can generate such representations, this time georeferenced and orthorectified, and viewable in the QGIS software.
\subsection{\bouton{Local alignment} and \bouton{Global alignment} check boxes}
As explained above, the alignment algorithm works in two passes (coarse process and fine process). The result of either one (or both) can be shown by ticking the \texttt{Global alignment} and \texttt{Local alignment} boxes. The \texttt{Global alignment} box shows the results of the coarse process, while the \texttt{Local alignment} box shows the result of the refinement.  At least one of these two boxes must be ticked. We suggest always ticking the \texttt{Local alignment} box. Note that in some circumstances the global alignment succeeds but the local alignment fails.

\subsection{\bouton{Maximum VIS/NIR time difference} field}
This field restricts the list of shoot points to align to those whose rank~1 time difference satisfies $|dt|\leq \texttt{Maximum VIS/NIR time difference}$. The default value is $0.5\,s$. This corresponds to well-synchronized images. In practice, with shooting intervals of $2 \, s$ for the VIS images and $3 \, s$ for the NIR images, this amounts to selecting one shoot point out of three.


\section{Best mapping: align only the images the map needs}\label{btn:step4_best_mapping}
The \bouton{Images aligned} frame offers two choices, which follow the one made in Step~2 (Section~\ref{btn:step2_best_mapping}): \bouton{Best mapping: only the images the map needs} (the default, with a $75\,\%$ overlap) or \bouton{All the synchronised images}. Among the well synchronized shoot points (\texttt{Maximum VIS/NIR time difference} box), IRDrone keeps only those whose ground footprints overlap by the chosen value (\bouton{Overlap between kept images} box): the first survey image, then each time the last image that still overlaps the previous one kept enough, and the last survey image. All the survey images are used to follow the motion of the drone; only well synchronized images can be kept. When the image that would be kept is badly synchronized, IRDrone keeps the nearest well synchronized one: the last one before it or, if its overlap is closer to the target, the first one after it (overlap at most 5 points below the chosen value, 70\,\% for 75\,\%). The number of images aligned thus stays close to the preview of Step~2; without this, every badly synchronized image would shorten the step and many more images would be kept. The line below the choices gives the number of shoot points that will be aligned, for example ``18 of 39 synchronised shoot points will be aligned''.

On the right, the flight map shows the shoot points: in green those well synchronized and in the image range, in grey the others; the images kept by the best mapping are bigger, with their ground footprints in green, which shows the overlap. It is updated as soon as a setting changes (range, time difference, overlap, choice). Wheel to zoom, drag to pan, double-click to see the whole flight again. During the alignment, the map follows the computation along the flight: the footprints of the aligned images fill with green, those being aligned (one per worker process) are outlined in orange, those still to do are dashed; the count ``7 of 18 done'' is shown at the bottom of the map. The result stays shown at the end, until a setting changes.

This is usually the best choice: much less computation (alignment is the longest step), less data, and an ODM map that is often better (too many nearly identical images slow ODM down and degrade its result). Choose \bouton{All the synchronised images} to align all the synchronized images, for example to look at them one by one in the review. The choice is saved in \texttt{FlightAnalytics\textbackslash best\_mapping.json} when the alignment starts; Step~6 uses it for the selection of the ODM images (Section~\ref{btn:step6_selection}).

The same selection is available on the command line: \texttt{run\_after\_preprocessing.py --best-mapping-overlap 75}.

\section{\bouton{Run alignment} button}\label{btn:step4_run_alignment}
Runs the alignment procedure on the shoot points that satisfy both constraints
\begin{itemize}
	\item $\texttt{first VIS image number} \leq XXXX \leq \texttt{last VIS image number}$,
	\item $|dt|\leq \texttt{Maximum VIS/NIR time difference}$,
\end{itemize}
and, if \bouton{Best mapping: only the images the map needs} is chosen, only the images the map needs (Section~\ref{btn:step4_best_mapping}).
Note that image alignment is a very time-consuming procedure. This of course depends on the computer. The computations are spread over several processors if the computer is ``multi-core''. The progress bar shows how far the computation has progressed, along with an estimate of the remaining processing time. As an example, on a laptop with a quad-core Intel Core~i7 processor and 16\,GB of memory (3 alignment workers), aligning the VIS and NIR images of one shoot point takes about $15\,s$, so 120 images take about half an hour. For a large mission, running these computations at lunchtime is still a good idea.

The alignment never starts by itself, even when Step~4 is opened from Step~3. During the computation, the button becomes \bouton{Stop} (Section~\ref{btn:step4_stop}).

\section{\bouton{Stop} button}\label{btn:step4_stop}
During the computation, \bouton{Run alignment} becomes \bouton{Stop}. After confirmation, the computation and all its worker processes stop; the window shows how many images have been aligned. The images already aligned are kept; \bouton{Run alignment again} restarts the computation, and the points already done reuse their motion model (faster). To process only the remaining points, increase \texttt{First VIS image number}.

\section{\bouton{⚙} button: settings}\label{btn:step4_settings}
The gear wheel opens the computation settings, which are kept for later sessions:
\begin{itemize}
	\item \texttt{Alignment worker processes}: number of shoot points aligned in parallel. \texttt{Auto} uses the physical cores, within the limit of the memory available at launch (about 3\,GB per process). \bouton{Safe mode (1)} keeps only one: slower, but uses the least memory. Use it if the computer freezes or if the computation is interrupted.
	\item \texttt{Log detail}: amount of computation messages in the window log: warnings only, computation steps, or timing details. With the timing details (\emph{Debug}), a realignment by hand from the review (Section~\ref{btn:review_redo}) also writes all its diagnostic images: several minutes per point instead of about half a minute.
\end{itemize}
\bouton{Save and close} (\bouton{Close} while nothing is changed) keeps the settings. \bouton{↺}, shown when the settings differ from the default ones, goes back to \texttt{Auto} and warnings only.

\section{\bouton{Close} button}\label{btn:step4_close}
Closes the window and returns to the main dialog. The \bouton{?} button opens this help at the page of this chapter. It is greyed out while the alignment runs (stop it first with \bouton{Stop}).

\section{\bouton{Continue with the review} button}\label{btn:step4_continue}
Appears at the bottom right, after \bouton{Close}, once images are aligned, and is highlighted at the end of the alignment. It closes the window and opens Step~5, the review of the aligned images (Chapter~\ref{ch:review}). When the alignment was started from Step~3, the review opens by itself at the end.

%==========================================================
\chapter{Step 5 -- Review of the aligned images}\label{ch:review}
The review shows all the points of the survey on a timeline to check the alignment of each image (Figure~\ref{fig:step4_review}). It opens with \bouton{Continue with the review} in Step~4, or with Step~5 in the main window (once Step~4 has aligned images). It replaces looking at the images one by one in the file explorer. Nothing is computed when it opens: the step counts as done once the review has been closed, the flags and confirmations being saved as you go.

\figureIRD[fig:step4_review]{step4_review.png}{Review of the aligned images: curtain between visible and aligned infrared, point card (measured offset, associated infrared image), mini-map, and at the bottom the timeline of the points (residual offset, thumbnails, status).}{0.95}

\section{Automatic measurement}\label{sec:review_measure}
For each point, the residual offset between the edges of the visible image and those of the aligned infrared is measured (in pixels of the full image, median over 16 areas). A well-aligned point measures about 1~px. Green: less than 3~px; light orange: between 3 and 6~px, worth a look; orange: more than 6~px, probably misaligned. An area without texture (water, uniform field) cannot be measured. The window opens on the most doubtful point.

The measurements and the checked points are kept in \texttt{ImgIRdrone/review.json}: reopening the review does not redo the measurements, except for the images that were redone.

\section{Comparing visible and infrared}\label{btn:review_compare}
Four ways to compare, with the buttons above the image or the keys 1 to 4; the review opens in \bouton{Overlay}. A frame around the image is in the colour of the point's status (green aligned or confirmed, red flagged, purple realigned\ldots), the labels on the image stay neutral (the reds and cyans are those of the images):
\begin{description}
	\item[\bouton{Swipe} (1)] a curtain separates the visible (on the left) from the aligned infrared in false colors (on the right); roads, hedges and roofs should continue across the curtain. Drag the curtain to move it.
	\item[\bouton{Blink} (2)] visible and infrared alternate: a misaligned image ``jumps''.
	\item[\bouton{Overlay} (3)] visible in cyan, infrared in red (vegetation is naturally red); a misalignment doubles every edge.
	\item[\bouton{Side by side} (4)] the two images next to each other.
\end{description}
Mouse wheel: zoom (an offset of a few pixels is only visible when zoomed in); drag: pan; double-click: zoom $\times 4$ or whole image. \bouton{Open full size} opens the aligned image in the computer's image viewer.

\section{Locating the image on the map}\label{btn:review_map}
The mini-map (in the point card) shows the flight, the current point and the ground footprint of the image (the thick line is the top of the image), in the colour of its status (green aligned or confirmed, red flagged\ldots); the shoot points Step~4 did not align are small grey dots, full size when zoomed in closely. A click enlarges it (\bouton{Map}, key 5). Under the mini-map, the legend lists only the statuses present, with their number. The large map shows the aligned image in its footprint, over the Step~6 orthophoto when it exists and the aerial photos in the background: the infrared should continue the roads and hedges of the background. Hovering over a point shows its thumbnail, a click selects it; a click outside the points leaves no point selected on the map: the trajectory alone, no footprint and no image (the point card keeps its point; a click on a point, or a move along the strip, selects one again); mouse wheel and drag to zoom and pan, double-click to see the whole flight. The \bouton{Layers} menu, the offline base map and the footprint computation are described in Section~\ref{btn:home_map}. Without a GPS position in the shoot points, the map is not available.

\paragraph{Aligning a point left out by Step~4.} A right-click on a grey point of the map (mini-map or large map) shows its time difference with the infrared and offers \bouton{Align this point}: it is aligned like Step~4 does. When an image cannot be recovered, even by hand, a neighbour left out by the best mapping or the time filter may be fine: align it this way. The point is aligned in the background (orange on the map, progress in the banner, \bouton{Stop aligning} to stop) while you keep reviewing; it then appears in the timeline, to review like the others, and Step~6 always maps it (unless you flag it), whatever its own time tolerance and image selection. The time filter is the widest one of Step~4 ($|dt|<6$~s); a point without an infrared image close enough in time cannot be aligned.

\figureIRD[fig:step4_review_map]{step4_review_map.png}{``Map'' mode of the review: flight, footprint and aligned image of the current point, Step~6 orthophoto over the IGN aerial photos.}{0.9}

\section{Moving around}\label{btn:review_navigation}
\begin{itemize}
	\item \bouton{◀} and \bouton{▶} (keyboard arrows): previous point, next point.
	\item \bouton{Next to review} (key N): next point that needs a look (realigned to confirm, measured as doubtful or not measured). A point measured aligned (green) needs nothing: there is nothing to mark as good.
	\item The timeline at the bottom: hovering shows the image, a click selects the point; mouse wheel to zoom, drag to pan, double-click for the whole survey.
	\item \bouton{Close} closes the review; what has been flagged and confirmed is already saved. \bouton{Continue with mapping} also closes it, then opens Step~6 (Section~\ref{btn:review_continue}).
	\item The tool tip of the \bouton{?} button lists all the keys.
\end{itemize}

\section{Deciding: flagging a point}\label{btn:review_decide}
\begin{description}
	\item[\bouton{Flag: leave out of the map} (key F)] flags the point, \textbf{which leaves it out of the map}: misaligned, but also blurred, or with a black area at the edge of the infrared image (the drone tilted by a gust of wind: the infrared camera does not see all of the visible field). Pressing it again removes the flag. The flag is written immediately to \texttt{ImgIRdrone\textbackslash shoot\_points.json}\footnote{This file does not contain the same information as the file \texttt{AerialPhotography\textbackslash shoot\_points.json}.} and Step~6 ignores this point. A flagged point can be realigned by hand (Section~\ref{btn:review_redo}); if that does not recover it, leave it flagged.
	\item[\bouton{Flag the N probably misaligned}] flags in one go all the points whose measured residual misalignment is above 6~px (red \emph{probably misaligned} points). It is shown only when there are some: the measure does not see a blur or a black area, which you flag yourself.
\end{description}

The example below shows the beginning of the file \texttt{ImgIRdrone\textbackslash shoot\_points.json}, which stores this information. All the images were processed with the loose criterion $|dt|<4\,s$ and the ``nearest'' method. This corresponds to 501 VIS images and 350 NIR images. Here the user judged visually that shoot point 120 was misaligned and flagged it, which switched the value of the \texttt{failed} key from \texttt{false} to \texttt{true}.
\begin{verbatim}
	{
		"mission_path": "C:\\Air-Mission\\FLY-20220524-1715-Chalivoy-Milon",
		"time_shift": -48.116,
		"n_VIS": 501,
		"n_NIR": 350,
		"pairing_method": "nearest_time",
		"created_at": "2026-10-29T20:44:13.203434",
		"schema_version": "1.0",
		"shoot_points": [
		{
			"FileName": "VIS_0119.dng",
			"dt": -0.5390370000000075,
			"failed": false
		},
		{
			"FileName": "VIS_0120.dng",
			"dt": 0.6902460000000019,
			"failed": true
		},
		{
			"FileName": "VIS_0121.dng",
			"dt": -1.3297540000000083,
			"failed": false
		},
		...
		]
	}
\end{verbatim}

What should be done next with the misaligned shoot points? Two options:
\begin{itemize}
	\item ``help'' the IRDrone algorithm align them: \bouton{Realign N flagged by hand} (Section~\ref{btn:review_redo}). Depending on the result, the flag is kept or removed.
	\item exclude them from later processing: this is the default effect of the flag, Step~6 ignores the points with \texttt{"failed": true}.
\end{itemize}

\subsection{Example of a successful alignment}
\figurePaireIRD{fig:goodalignment}
{step4_goodalignment_VIS.jpg}{Image in the visible spectrum.}
{step4_goodalignment_NIR.jpg}{Nearest image $|dt|=0.165~s$ in the infrared spectrum.}
{Successful alignment: the time difference is small and the attitude of the drone stayed within the flight parameter specifications.}

\figureIRD[fig:step4_goodalignment_VIR]{step4_goodalignment_VIR.jpg}{Successful alignment. The two images were taken almost at the same moment and with no wind. }{0.5}

\subsection{Example of a failed alignment}
If the drone leaves its flight parameter envelope, for example in strong gusts of wind that force it to increase its pitch and yaw angles, then the aiming centers of the two cameras can be so far apart that no alignment can be found. The alignment can also fail if the time difference between the VIS and NIR shots is large (it can sometimes exceed 2\,s), because the horizontal movement of the drone is then poorly corrected. The alignment can also be imperfect without failing if there is a time difference and the plants move in the wind. The VIR image then shows ``ghosts''. Figures \ref{fig:misalignment} and \ref{fig:step4_misalignment_VIR} show a failed alignment caused by a strong gust of wind that destabilized the drone.

\figurePaireIRD{fig:misalignment}
{step4_misalignment_VIS.jpg}{Image in the visible spectrum.}
{Step4_misalignment_NIR.jpg}{Nearest image $|dt|=1.256~s$ in the infrared spectrum.}
{Failed alignment caused by a drone attitude outside the flight parameter specifications, due to a gust of wind at the moment of the shot.}

\figureIRD[fig:step4_misalignment_VIR]{step4_misalignment_VIR.jpg}{Failed alignment. The ``hole'' effect is especially typical. The user can easily spot it on the VIR images. }{0.5}

\section{Realignment by hand}\label{btn:review_redo}
The \bouton{Realign N flagged by hand} button recovers the shoot points flagged as misaligned. There is no miracle: if a gust of wind pushed the drone out of its flight envelope, the two cameras do not see the same scene and no alignment is possible. But sometimes the information is there and the coarse alignment parameters only need a ``nudge''. Under normal conditions, at most about ten shoot points need to be redone.

\subsection*{Procedure}
\begin{enumerate}
	\item Flag the misaligned points, then press \bouton{Realign N flagged by hand}. The alignment window opens with all the flagged points (Figure~\ref{fig:step4_align_by_hand}).
	\item For each point, move the infrared image over the visible one until the same features coincide (see below), then \bouton{Placed, next} (Enter key). \textbf{The alignment of this point starts at once, in the background}, while you place the next one: no waiting between two points (the point turns to \emph{⚙ Computing} in the list and can no longer be moved). The images of the next points are loaded in the background. \bouton{Skip this point} (S key) leaves a point out: it stays flagged, hence left out of the map.
	\item \bouton{Compute N alignments} sends the points placed without \bouton{Placed, next} (when everything is sent already, \bouton{Cancel} becomes \bouton{Close}) and closes the window. The alignments are computed in the background, several at a time (as many as the alignment workers of Step~4, gear button). Count about half a minute per point. The banner at the top of the review shows the progress: points done, points running with their stage and elapsed time, points waiting. \textbf{The review stays usable}: you can keep checking the other points, or do something else. The points being computed cannot be flagged or unflagged. \bouton{Stop} starts no new alignment: the running ones finish, the others stay flagged. The review cannot be closed while alignments are being computed.
	\item Each realigned point turns purple, \emph{Realigned, to confirm}: it is still flagged (Step~6 skips it). When everything is computed, the review shows the first one. Check it with the curtain, blink or overlay and answer \bouton{Yes, aligned} (Y key) or \bouton{Still misaligned}: the next point to confirm is then shown. \bouton{All N aligned} confirms all of them at once. The N key (next to review) goes to them first. The points to confirm are remembered if the review is closed.
	\item With \bouton{Yes, aligned}, the \texttt{"failed"} key of the point is set to \texttt{false} in \texttt{shoot\_points.json} and the file \texttt{shoot\_points\_failed\_summary.json} is updated (a copy \texttt{shoot\_points.json.bak} is kept). With \bouton{Still misaligned}, the point stays flagged and can be realigned again; if it cannot be recovered (infrared image too far off, blur), just leave it flagged: \textbf{it is left out of the map} (Step~6).
\end{enumerate}

The outputs of the previous alignments are archived in the folder \texttt{VIS\_XXXX\_align\_traces\_manual\textbackslash original}. The point gets back the images it had (those ticked in Step~4, Section~\ref{btn:step4_options}); the diagnostic images of the alignment (folder \texttt{Tech}, several minutes of computation per point) are only written when \texttt{Log detail} is set to \emph{Debug} in the Step~4 settings (gear button, Section~\ref{btn:step4_settings}), or from the command line (below).

\subsection*{The alignment window}
\figureIRD[fig:step4_align_by_hand]{step4_align_by_hand.png}{Alignment by hand: the flagged points on the left (a tick: placed, a circle: to place), the two images overlaid in the middle (here with the visible in cyan and the infrared in red), the angles on the right. Here the infrared is still about 140~px too high: the edges are doubled.}{0.9}

\begin{itemize}
	\item On the left, the list of the points: a circle for the points to place, a tick for the placed ones, a dash for the skipped ones. A click goes to a point.
	\item In the middle, the two images, shown as in the review: \bouton{Overlay} (both images on one: doubled edges mean misaligned; the default). \bouton{Colours} chooses its colours, kept for the next time: visible in colour and infrared in grey (the default; with \bouton{Edges only}, the visible edges are green), visible in green and infrared in grey, visible in blue and infrared in red, or visible in cyan and infrared in red (as in the review, close to the VIR image). \bouton{Swipe}, \bouton{Blink}, \bouton{Side by side} (keys 1 to 4). \bouton{Edges only} (G key) shows the outlines of both images, which helps on fields and forests where the infrared and the visible look different. Wheel to zoom, drag to pan, double-click to see the whole image.
	\item To move the infrared, \textbf{quickly then finely}: first the \textbf{sliders} of \texttt{Yaw}, \texttt{Pitch} and \texttt{Roll} (right panel) for a large move, then the buttons at their ends to refine, in the direction they show: \bouton{◀}~\bouton{▶} left and right (yaw), \bouton{▲}~\bouton{▼} up and down (pitch), \bouton{↺}~\bouton{↻} rotation (roll); held down, they repeat. The \textbf{arrow keys} do the same as these buttons (0.05\textdegree{} per press, about 2 to 3 pixels of the full image; with Shift, 10 times more), wherever the cursor is in the window, even in a value field; Q and E rotate. You can also \textbf{drag the infrared with the right mouse button} (or Shift + drag) in the image, or type a value in a field.
	\item \bouton{Fit} (F key) measures the offset between the edges of the two images (tile by tile) and moves the infrared by that much (yaw and pitch). It works once the images roughly overlap; if it finds no better position, it says so and changes nothing. \bouton{Reset} goes back to the angles of the previous alignment.
	\item Under these buttons, the offset measured between the two images, in pixels of the full image: below 10~px the point is \emph{well placed}.
	\item Roughly is enough: these angles are only the starting point, the coarse alignment and the pyramidal refinement start again from there (Section~\ref{sec:ann_algorithme}).
	\item A point whose images cannot be read is marked with a warning sign and skipped.
\end{itemize}

\subsection*{From the command line: the slider window}
Each aligned point has a script \texttt{VIS\_XXXX\_REDO.bat} (Windows) or \texttt{VIS\_XXXX\_REDO.sh} in \texttt{ImgIRdrone}, which runs its alignment again with a slider window (matplotlib): this is the former procedure, still usable outside the wizard. After a realignment in the wizard, the script starts from the angles placed by hand.



\figureIRD[fig:redo_info]{step4_redo_info.png}{Information window of the manual alignment help window. }{0.5}

The window shows the VIS image and the NIR image, with five sliders: \texttt{YAW}, \texttt{PITCH}, \texttt{ROLL}, \texttt{Absgrad visualize} and \texttt{Alpha}.
\begin{itemize}
	\item The \texttt{Alpha} slider chooses what is shown. Between $-1$ (far left) and $-0.5$ (small red middle mark), a slightly transparent rectangle of the NIR image is overlaid on the VIS image: the closer \texttt{Alpha} is to $-1$, the larger this rectangle. The initial value is $-0.75$.
	\item The \texttt{YAW}, \texttt{PITCH} and \texttt{ROLL} sliders move the NIR image. The method is to change one angle, then vary \texttt{Alpha} to observe the result, and repeat with the three angles. The key is to spot identical features in both spectral bands (tree, edge of a path, building) and make them coincide, first in the center and then at the edges of the rectangle.
	\item When \texttt{Alpha} is moved to the right, the window no longer shows a NIR rectangle: between $-0.5$ and $0$ it shows the difference between the two images, and between $0$ and $1$ a VIS/NIR cross-fade. This use is faster but intended for an experienced operator.
	\item \texttt{Absgrad visualize} (initial value $0.3$) shows the gradients of the two images in two colors when it is above $0.5$: the edges coincide when the alignment is good.\\
\end{itemize}

When the user considers that the manual alignment is satisfactory, they press the Q key and the semi-automatic alignment procedure starts. \\


As with any tool of this kind, you learn it by trying it out; once the method is understood, realignment is quick and effective.
Here are a few steps of the procedure in pictures
\footnote{The tool is rich and offers several ways of viewing the layers. So there is no single method to reach the desired result. The \texttt{Absgrad visualize} slider changes the representations. For example, moving it to the right gives an edge effect. This is not always easy to interpret, with forests for example, and even less so in flat, bare areas.}.
In the image, a road crosses diagonally.
There is also a large tree at the top right, a small shrub at the bottom left along the road, and a ditch at the top left (``black line''). These are the characteristic features we will try to align.
Note that at initialization the NIR layer is slightly transparent
\footnote{Moving the \texttt{Alpha} slider to the right between 0 and 1 takes the NIR layer from fully transparent to opaque.}.
There is no real order for reaching a satisfactory result.
Do not aim for perfection, because only global linear transformations are applied. The code will refine globally and then locally.
\figurePaireIRD{fig:redo_A}
{step4_redo_1.png}{Starting situation after a pass through the procedure of the \texttt{run alignement} button.}
{step4_redo_2.png}{The rotation was changed from $0$ to $-2.46$ to align the NIR road parallel to the VIS road. Use this slider carefully. When there is no obvious line such as the road here, rotation is rather used last, for the ``finishing touch''.}
{First change, using the rotation slider (\texttt{ROLL}).}
\figurePaireIRD{fig:redo_B}
{step4_redo_3.png}{Vertical translation downwards with the \texttt{PITCH} slider. The value goes from  $-1.14$ to $2.04$. This is a fairly large amplitude and, above all, the sign changes.}
{step4_redo_4.png}{Horizontal translation with the \texttt{YAW} slider. The value goes from $0.96$ to $1.8$.}
{Fine-tuning by playing with the \texttt{Alpha} slider (decreasing then increasing it) and watching the alignment of the characteristic features. The slight blur at the end comes from the transparency and from small local alignment errors that the fine alignment procedure will correct.}

\section{\bouton{Continue with mapping} button}\label{btn:review_continue}
At the bottom right, after \bouton{Close}. It closes the review and opens Step~6. The map, which is long, does not start there by itself: it is started with \bouton{Build the map\ldots}. The flagged points, and the realigned points not confirmed yet, are left out of the map. While alignments are being computed, the review cannot be closed: wait for them (or \bouton{Stop}).



%==========================================================
\chapter{Step 6 -- Mapping and georeferencing}\label{ch:step6}
\section{Principle}
The multispectral images produced in Step~4 are ``flat'' images: they are not positioned on the ground.

\figurePaireIRD{fig:step6_ortho}
{step5_VIS_plan.png}{Image corrected for lens distortion (DJI Mavic Air 2 and Zhang calibration).}
{step5_VIS_ortho.png}{Image orthorectified with a terrain model with a $10\,m$ step.}
{The effect of orthorectification is clearly visible in this image, where the terrain has a steep slope (see Figure~\ref{fig:step6_topo}).}

 Step~6 makes them usable in a geographic information system (GIS) such as QGIS. It offers four processes, one per button:
\begin{itemize}
	\item \bouton{Build the map\ldots} (the main one): stitching all the selected images into a single 4-band \emph{orthophoto} (red, green, blue, near infrared) with the free software OpenDroneMap (ODM), then the terrain model and the contour lines;
	\item \bouton{Terrain model and contours\ldots}: building the DTM and the \emph{contour lines} again, for example with another grid step, from elevations supplied by an online service (IGN);
	\item \bouton{Spectral balance\ldots}: balancing the near infrared with the visible bands before the map (see \ref{btn:step6_balance});
	\item \bouton{Georeference images\ldots}: \emph{georeferencing} and \emph{orthorectification} of individual images. Each image is projected onto a digital terrain model (DTM) and then saved in GeoTIFF format in UTM coordinates.
\end{itemize}
These processes rely on the flight metadata stored in the ``companion'' \texttt{.exif} files (GPS position, altitude, yaw, pitch and roll angles, see Step~1).

\figureIRD{step5_window.png}{Step~6 window.}{0.6}

The window is laid out as Step~3. From top to bottom: the name of the step and of the mission; the stages of the map (\texttt{Image selection}, \texttt{Orthophoto (OpenDroneMap)}, \texttt{Terrain model and contours}: $\circ$ to do, $\bullet$ running, \checkmark{} done, $\times$ failed, with a detail under each); a message area (log); the four buttons, \bouton{Build the map\ldots} first; a progress bar preceded by a label showing the operation in progress; an image preview area; and the \bouton{?} and \bouton{Close} buttons, then \bouton{Continue with archiving} once the map is built.

When the map was built earlier, the stages are already ticked and the orthophoto is shown as soon as the window opens. \bouton{Terrain model and contours\ldots} stays greyed out until the map exists.

\textbf{Prerequisites}
\begin{itemize}
	\item Step~4 must have been run for the images to process (folder \texttt{ImgIRdrone\textbackslash imgMultispectral});
	\item for \bouton{Build the map\ldots}: the \emph{Docker Desktop} software must be installed and running;
	\item for \bouton{Terrain model and contours\ldots} (and for the ground elevation during georeferencing): an Internet connection is needed to query the elevation service.
\end{itemize}

\section{\bouton{Georeference images\ldots} button}\label{btn:step6_georef}
This button georeferences and orthorectifies images chosen by the user.

A message recalls the accepted image types, then a file selection window opens. All the chosen files must be of the same type:
\begin{itemize}
	\item DNG images from the DJI camera (visible spectrum);
	\item multispectral TIFF images in \textbf{groups of four}: \texttt{IRD\_XXXX\_1.tif} to \texttt{IRD\_XXXX\_4.tif} (red, green, blue, near infrared). All four files must be selected; an incomplete group is rejected with a warning.
\end{itemize}

In a visible only mission, the message and the file selection window offer the DJI DNG images only.

\subsection*{Georeferencing offsets}
Once the images are chosen, a window asks for three corrections, all equal to \textbf{0} by default:
\begin{itemize}
	\item the \emph{yaw offset} (in degrees), which corrects a systematic error in the heading of the drone;
	\item the \emph{UTM east offset} (in meters);
	\item the \emph{UTM north offset} (in meters).
\end{itemize}

These values apply to all the images processed in this run. The offsets are useful because the GPS
\footnote{Positioning accuracy could be greatly improved with an RTK beacon, but this kind of equipment is expensive.} and attitude data of the DJI drone have limited accuracy, especially if the GPS calibration and/or the camera orientation calibration were not done before take-off. These ``offset'' values can only be determined experimentally, for example by using QGIS to refine the positioning.

 The \emph{Cancel} button cancels the georeferencing.

\subsection*{What the software does for each image}
\begin{enumerate}
	\item DNG images are converted to TIFF with RawTherapee. For multispectral images, the four bands are normalized between 0 and 1 and then stacked into a 4-channel image.
	\item The flight metadata are read from the companion \texttt{.exif} file (created if needed). They give the position of the drone, its altitude above the ground and its attitude angles. If the elevation of the take-off point is unknown, the user may be asked for it.
	\item The ground footprint of the image is estimated from the camera geometry, then a terrain mesh is built. The elevations of the mesh points are requested from IGN.
	\item Each mesh point is projected into the image. The image is resampled on a regular ground grid (orthorectification).
	\item The result is saved in GeoTIFF format (UTM coordinates).
\end{enumerate}

\subsection*{Output files}
The files are created in the \texttt{geo\_ref} subfolder of the folder containing the selected images. An existing file is never overwritten: a suffix \texttt{\_1}, \texttt{\_2}, \dots~is added to the name (useful for example if the previous file is open in QGIS).
\begin{itemize}
	\item \texttt{*\_georef.tif}: the orthorectified image, to open in QGIS. For a multispectral quadruplet the name is \texttt{IRD\_XXXX\_0\_georef.tif} and the image has 4 bands (Red, Green, Blue, NIR);
	\item \texttt{*\_MNT.png} and \texttt{*\_Grid.png}: check plots of the terrain and of the mesh;
	\item \texttt{*\_dtm\_map.png}: shaded 2D map of the DTM under the image, with contour lines, UTM grid and IGN elevation points 	(same layout as for the mapping).
	\end{itemize}
	During the computation the window stops responding: this is normal. The image being processed is shown in the preview area and the log details the progress.

	The placement accuracy depends on the quality of the flight metadata (GPS, yaw) and of the DTM. A systematic offset can be corrected with the offsets window (see above).

\figureIRD{step5_run_Geo_ref.png}{Window during the georeferencing of an image. The original, non-orthorectified image is shown.}{0.6}


\section{\bouton{Build the map\ldots} button}\label{btn:step6_mapping}
This button stitches the aligned images from Step~4 into a single orthophoto of the whole surveyed area, using the OpenDroneMap (ODM) software. ODM runs in a \emph{Docker} container: the user has nothing to configure, but Docker Desktop must be running (see below).

\subsection{Choice of images}\label{btn:step6_selection}
Clicking the button opens a window for choosing the criteria. A note at the bottom recalls that the computation can be long and use several gigabytes of memory and disk space. \bouton{Build the map} starts it at once; \bouton{Cancel} closes the window and starts nothing. Once the map has been built, the window proposes the settings of the last run (the default values below otherwise; the best mapping chosen again in Step~2 or~4 after that run wins). The images can come from the whole flight, not only from the survey phase of Step~2.

\begin{itemize}
	\item \texttt{Maximum |dt| (NIR rank 1)}: maximum allowed time difference between the VIS image and the nearest NIR image. This is the same criterion as the \texttt{Maximum VIS/NIR time difference} of Step~4. Default value: $0.5\,s$ (range $0.05$ to $6\,s$).
	\item \texttt{Best mapping: select by image overlap (recommended)}: default choice. Keeps only the images whose ground footprints overlap by the value given (\texttt{Overlap between successive images}, by default the value chosen in Step~2 or~4, $75\,\%$ otherwise). This is the same selection as in Steps~2 and~4 (Section~\ref{btn:step2_best_mapping}). If Step~4 aligned only the images the map needs, they are all kept; if it aligned everything, only the images the map needs are given to ODM. An image declared misaligned in the review is left out and the next one takes its place.
	\item \texttt{Select one image every n eligible images}: keeps only one image out of $n$ ($n$ from 1 to 4, default value 3) to limit the computation time and the overlap. This is the default choice when the best mapping was unticked in Step~2 or~4.
	\item \texttt{All the aligned images (every point of the review not flagged)}: ignores the criteria above: every image Step~4 aligned (or aligned from the review) and not flagged is mapped, whatever its time difference and the overlap. The most images, the longest mapping. It is proposed again the next time if it was used.
\end{itemize}

\figureIRD{step5_odm_selection.png}{Window of the image selection criteria for the mapping.}{0.4}

On the right, as in Step~4, the flight map shows the images selected for the map, in green with their ground footprints (their overlap shows), the eligible images left out by the selection in grey, and the others as small dots. It is updated as soon as a setting changes ($|dt|$, overlap, $n$, choice); the count below the map gives the number of images selected among the eligible ones. Wheel to zoom, drag to pan, double-click to see the whole flight again.

\textbf{OpenDroneMap: version and GPU.} Under the choices, the \texttt{OpenDroneMap} box shows, after a check of a few seconds (Docker, Internet), the version of ODM installed on the computer and whether a newer one has been published on Docker Hub:
\begin{itemize}
	\item \texttt{Download the newer version before the map}: shown only when a newer version exists, unticked by default. Ticked, the newest ODM image is downloaded (\texttt{docker pull}, a few hundred megabytes to a few gigabytes) before the computation. A newer version can give a slightly different map, so it is never installed without being asked. Without Internet, the box says that the newest version could not be checked and the installed one is used.
	\item \texttt{Use the NVIDIA GPU (faster image matching)}: available, and ticked, only when the computer has an NVIDIA graphics card \emph{and} Docker can give it to a container (the software starts a test container with \texttt{--gpus all}). ODM then runs its GPU version (\texttt{opendronemap/odm:gpu}, about 3\,GB downloaded the first time), which detects and matches the image features on the graphics card; the rest of the computation is the same. Otherwise the box is greyed out and the reason is given (no NVIDIA card; on Linux the \emph{NVIDIA Container Toolkit} is missing; on Windows, Docker Desktop must use WSL~2 with a recent NVIDIA driver); its tooltip gives Docker's answer. If the GPU cannot be used when the computation starts, ODM runs on the processor and the log says so.
\end{itemize}
The log gives the ODM version used and whether it ran on the GPU or the processor.


The selected images (\emph{eligible images}) are the shoot points that satisfy all of the following:
\begin{enumerate}
	\item they were not flagged as misaligned in the review, Step~5 (key \texttt{"failed": false}, see Section~\ref{btn:review_decide});
	\item the difference $|dt|$ with the nearest NIR image is less than or equal to the tolerance;
	\item the corresponding VIS image belongs to the flight phase of the mission chosen in Step~2 (the take-off and synchronization images are therefore excluded).
\end{enumerate}
Among them, the best mapping keeps the images the map needs (or, with the other criterion, one eligible image out of $n$).

\subsection{Starting the computation}
\bouton{Build the map} starts the computation without another question (the selection is written in the log):
\begin{enumerate}
	\item the software checks that Docker Desktop is available. If not, a message asks you to start it and then choose \texttt{Retry}. ODM is then downloaded if it is not yet on the computer, or updated if asked (see above);
	\item the list of selected images is saved in \texttt{mapping\_MULTI\textbackslash img\_best\_mappping.json};
	\item the resolution of the orthophoto is estimated automatically: it is the ground size of a pixel (GSD, in cm) obtained from the flight altitude above the ground and the focal length of the camera, median over the selected images, rounded up to the next tenth of a cm;
	\item the selected TIFF images are made available to ODM in \texttt{mapping\_MULTI\textbackslash images} through \emph{hard links}: no image is copied, so disk space is not doubled. If the disk does not support hard links (external disk in exFAT or FAT32, some network shares), the images are copied, which uses more disk space; the log says so;
	\item ODM is started. The window log shows the main messages; the full output is saved in a file \texttt{mapping\_MULTI\textbackslash odm\_run\_<date>.log};
	\item at the end, the orthophoto is checked (4 bands, georeferenced), the useful results are archived in \texttt{mapping\_result} and the very large ODM working folders are deleted.
\end{enumerate}
During the computation the progress bar is indeterminate, and the buttons and closing the window are disabled. When the computation ends, its duration is shown in the log.

\textbf{Results}\label{sec:step6_mapping_result}

The results are in \texttt{mapping\_MULTI\textbackslash mapping\_result}:
\begin{itemize}
	\item \texttt{odm\_orthophoto.tif}: the 4-band orthophoto (red, green, blue, near infrared), in UTM;
	\item \texttt{report.pdf}: quality report produced by ODM;
	\item \texttt{odm\_25dmesh.ply}: 2.5D mesh of the terrain;
	\item \texttt{shots.geojson} (and \texttt{shots.qml}): estimated position of each shot, with its shoot point number, to open in QGIS;  \texttt{shots.qml} is a style file.
	\item technical files (\texttt{ortho.png}, georeferencing information).
\end{itemize}
A color preview (red, green, blue bands) of the orthophoto is then shown in the window, then the computation of the DTM and contour lines starts automatically (this is equivalent to the \bouton{Terrain model and contours\ldots} button, see \ref{btn:step6_dtm}).

\figureIRD{step5_orthophoto_preview.png}{Preview of the orthophoto in the Step~6 window.}{0.6}

\subsection{Precautions}
\begin{itemize}
	\item Existing results are \textbf{never overwritten}. When a map is already built, it is kept: its folder is renamed \texttt{mapping\_MULTI\_<date of its run>} and the new map is built in a new \texttt{mapping\_MULTI} (with the same camera configuration); the selection window says so. If the new run fails, the previous map is put back in \texttt{mapping\_MULTI} and the failed run, with its log, is kept in \texttt{mapping\_MULTI\_failed\_<date>}. Step~7 keeps these folders.
	\item Docker must use the ``desktop-linux'' engine of Docker Desktop. The first run downloads the ODM image (several gigabytes): an Internet connection is needed.
	\item If some selected images are incomplete (missing TIFF quadruplet or different sizes), the computation is refused and the log lists the files concerned.
\end{itemize}

\section{\bouton{Terrain model and contours\ldots} button}\label{btn:step6_dtm}
This button builds a digital terrain model (DTM) and contour lines over the footprint of the orthophoto. It runs \emph{automatically} at the end of a successful mapping (with a grid step of 10\,m); this button lets you run it again, for example with another step. It needs the orthophoto \texttt{mapping\_MULTI\textbackslash mapping\_result\textbackslash odm\_orthophoto.tif}.

A window asks for the \texttt{Altitude grid step} value, that is, the spacing (in meters) of the points of the regular grid whose elevations are requested from IGN (elevation service of the French National Institute of Geographic and Forest Information); if it is unavailable, the OpenTopoData service is used, with lower accuracy in France. The smaller the step, the larger the number of requests. The default step is 10\,m (range 1 to 100\,m). The number of points is limited to 5000: above that, the software automatically increases the step. \bouton{Build} starts the computation, \bouton{Cancel} starts nothing.

\subsection*{Computation steps}
\begin{enumerate}
	\item reading the footprint of the orthophoto (UTM);
	\item building a regular grid that covers the footprint;
	\item elevation requests (no missing point is tolerated: otherwise the computation fails and a message says so);
	\item interpolation with a smooth surface (quadratic spline), then sampling on a 1\,m raster;
	\item export of the DTM, the contour lines (1\,m interval, index contours every 10\,m), a shaded 2D map and a 3D view.
\end{enumerate}
If it fails, the orthophoto is not affected.

\subsection{Output files}
The files are created next to the orthophoto, in \texttt{mapping\_result} (never overwritten, a numeric suffix is added if needed):
\begin{itemize}
	\item \texttt{ortho\_dtm.tif}: DTM (elevation in meters, UTM);
	\item \texttt{ortho\_contours.geojson} and \texttt{ortho\_contours.qml}: contour lines and their style for QGIS;
	\item \texttt{ortho\_dtm\_map.png}: shaded 2D map with contour lines, positions of the grid points (black dot) and UTM grid (light blue);
	\item \texttt{ortho\_dtm\_3d.png}: 3D surface;
	\item \texttt{ortho\_dtm\_grid.json}: grid points with their elevation and the source of the elevations.
\end{itemize}

\figurePaireIRD{fig:step6_topo}
{step5_dtm_map.png}{2D map of the DTM with contour lines and UTM grid.}
{step5_dtm_3d.png}{3D surface of the DTM.}
{Topographic representation of the terrain from the digital terrain model derived from IGN elevation data.}

\subsection{Opening in QGIS}
In the QGIS map window, drag in the orthophoto \texttt{odm\_orthophoto.tif}, the contour lines \texttt{ortho\_contours.geojson} and the shot points \texttt{shots.geojson}    (the \texttt{.qml} styles with the same names are applied automatically). \\

A few simple actions with a 4-band orthophoto red, green, blue, nir

\begin{itemize}
	\item Color image in the visible spectrum: right-click on the layer> \underline{P}roperties>Symbology, then map bands 1, 2, 3 to red, green, blue for a natural color view. This is the default display option for a multispectral image.
	\item VIR false color image: right-click>\underline{P}roperties>Symbology, then map bands 4~,1~,2 to red, green, blue.
	\item NDVI index map\footnote{The NDVI (\textit{Normalized Difference Vegetation Index}) measures the vigor and density of vegetation. It is computed with the simple formula: $NDVI = \frac{NIR - RED}{NIR + RED}$} : \underline{R}aster>Raster Calculator> enter the formula
	\begin{equation}
		\frac{(VIS@4 - VIS@1)}{(VIS@4 + VIS@1)}
	\end{equation}
	then double-click on the layer and, in Symbology, choose the render type: Singleband pseudocolor> Color ramp: NDVI. The result is shown in Figure~\ref{fig:step6_mapping_NDVI}
	\item NIR layer. In Symbology, select NIR as the first band, then in Symbology choose the render type: Singleband gray.
\end{itemize}

\figurePaireIRD{fig:step6_mapping_1}
{step5_topo_QGIS.png}{Detailed representation of the topography. Contour interval: $10\,m$.}
{step5_mapping_VIS.png}{Mapping in the visible spectrum.}
{IRdrone data shown in the QGIS software over an OpenStreetMap background.}

\figurePaireIRD{fig:step6_mapping_NDVI}
{step5_mapping_VIR.png}{VIR multispectral image.  $\qquad\qquad\qquad\quad$ Spectral shift $[Nir|R|G] \longrightarrow [R|G|B]$.}
{step5_mapping_NDVI.png}{NDVI index map.}
{IRdrone data shown in the QGIS software over a Google Satellite background.}


\section{\bouton{Spectral balance\ldots} button}\label{btn:step6_balance}
The red, green and blue bands come from the drone camera, whose automatic white balance keeps them well balanced. The near infrared (NIR) comes from the SJCam, with its own exposure: in the multispectral images it is often much darker (or brighter) than the visible bands, and a false colour image (NIR, red, green) is dominated by one colour. This window adjusts the tones of the bands so that they match. The button is greyed out in a visible only mission and until Step~4 has written the multispectral images.

\subsection*{The window}
\begin{itemize}
	\item \textbf{Left}: the layer adjusted, \bouton{RGB} (the three visible bands together), \bouton{R}, \bouton{G}, \bouton{B} (one visible band, after the RGB settings) or \bouton{NIR}; four sliders, \emph{Exposure} (in EV: $\times 2$ per EV), \emph{Gamma} (above 1 lifts the dark and mid tones), \emph{Contrast} and \emph{Brightness}, with a box on the right where the value can be typed; the \emph{tone curve} of the layer (a click adds a handle, drag it, a double-click or a right-click removes it; the curve is a monotone spline, the straight line by default; the grey histogram behind it is what the curve receives); \bouton{Reset this layer}.
	\item \textbf{Centre}: the image, the original against the balanced one, on a black background: \bouton{Swipe} (original on the left of the curtain, balanced on the right), \bouton{Blink} (one then the other) or \bouton{Side by side}; shown in \bouton{Visible}, \bouton{NIR} (grey) or \bouton{VIR} (NIR in red, red in green, green in blue). The wheel zooms, a drag pans, a double-click fits the image; the zoom is written at the bottom right. \bouton{Map} shows the flight on aerial photos or maps (IGN, Esri, OpenStreetMap or none), the images of the map and the ground covered by the selected image.
	\item \textbf{Right}: the selected image and \bouton{Validate}; the histograms of the four bands of the balanced image (the original NIR dotted); \bouton{\(\blacktriangleleft\)}, \bouton{Next to validate}, \bouton{\(\blacktriangleright\)}; the map (a click opens it large); \bouton{Write all the balanced images}.
	\item \textbf{Bottom}: the images of the map (those the next \bouton{Build the map\ldots} uses: aligned, not flagged, within the time difference and the overlap chosen), as false colour thumbnails. An orange star marks an image whose balanced bands are written, a grey one an image whose bands are out of date (the settings or the image changed since). The wheel zooms (the times of the shots appear), a click selects an image.
\end{itemize}

\subsection*{One balance for the mission}
The settings apply to \textbf{every image} of the mission: the map is a mosaic of many images, treated alike they leave no seams. An image may get a balance of its own (\emph{This image has a balance of its own}): it starts from a copy of the mission balance (a small pencil marks it on the strip).

The exposure of the SJCam changes from image to image, a fixed setting cannot balance them all: \emph{Match the NIR to the visible (each image)} computes on each image the exposure and the gamma of the NIR so that its median and its highlights (98th percentile) are those of the visible bands. The values computed for the selected image are shown under the sliders; contrast, brightness and curve still apply after. This is usually the only setting needed.

\subsection*{Files written}
In \texttt{ImgIRdrone\textbackslash img\_Bal\_Multispectral}:
\begin{itemize}
	\item \texttt{spectral\_balance.json}: the settings, saved when the window closes (\bouton{Close}, ``Save and close'' once something changed, or the cross);
	\item \texttt{IRD\_XXXX\_1.tif} to \texttt{IRD\_XXXX\_4.tif}: the balanced bands (red, green, blue, NIR), in 16 bits with the metadata of the originals, written by \bouton{Validate} (key \texttt{V}) or \bouton{Write all the balanced images} (about a second per image, in the background, \bouton{Stop} to interrupt);
	\item \texttt{IRD\_XXXX\_bal.json}: the settings applied to the image and the values computed for it.
\end{itemize}
The originals in \texttt{imgMultispectral} are never changed. Once a balance is set, \bouton{Build the map\ldots} uses the balanced bands; those not written yet, or out of date, are written first (stage \texttt{Spectral balance} of the log). \bouton{\(\circlearrowleft\)} at the bottom left puts every setting back to neutral: the map then uses the original bands again.

Keys: \(\leftarrow\) \(\rightarrow\) previous / next image, \texttt{N} next image to validate, \texttt{V} validate, \texttt{1} to \texttt{3} the comparison modes, \texttt{4} the map.

\section{\bouton{Close} button}\label{btn:step6_close}
Closes the window and returns to the main dialog. Closing is refused while an ODM or DTM computation is running. The \bouton{?} button opens this help at the page of this chapter.

\section{\bouton{Continue with archiving} button}\label{btn:step6_continue}
Appears at the bottom right, after \bouton{Close}, once the map is built. It closes the window and opens Step~7, which shows how much space archiving frees. Nothing is deleted until you confirm there.

%==========================================================
\chapter{Step 7 -- Mission archiving}\label{ch:step7}
Step~7 is the archiving of the mission. It is done \emph{once you have finished working with the images}. It cleans the mission folder \emph{in place}:
the many technical files (intermediate images, working folders, copies of synchronization images, etc.) are deleted,
and only the items essential for processing the mission again, if needed, are kept. On the missions tested, the size is reduced
by a factor of 2 to 5 (for example from 131\,GB to 25\,GB for a mission with 600 DJI VIS images).

\begin{center}
\fbox{\parbox{0.9\linewidth}{\small\textbf{Warning: this operation cannot be undone.} Deleted files do not go to the Windows recycle bin; they are permanently lost.
If you want to keep a copy of the mission, of a particular \texttt{mapping\_MULTI} folder or of any other file, \emph{copy them elsewhere before running the archiving}.
Anything in the mission folder that is not in the list below is deleted, including copies of folders (for example \texttt{mapping\_MULTI - Copie}).}}
\end{center}

As for the other steps, the mission archived is the mission open in the main window.

\section{What is kept}\label{btn:step7_rules}
File names are compared without regard to case (\texttt{.dng} and \texttt{.DNG} are equivalent).
\begin{description}
  \item[Mission root] the folders \texttt{AerialPhotography}, \texttt{cameras}, \texttt{FlightAnalytics}, \texttt{ImgIRdrone}, \texttt{mapping\_MULTI} (and the previous maps \texttt{mapping\_MULTI\_<date>} set aside by Step~6), \texttt{Synchro}
        and the file \texttt{desktop.ini}. Everything else is deleted.
  \item[\texttt{AerialPhotography}] the \texttt{VIS} and \texttt{NIR} folders and the file \texttt{shoot\_points.json}. In \texttt{VIS} and \texttt{NIR}: the original
        \texttt{.dng} images with their companion \texttt{.exif} files, as well as \texttt{time\_line.json} and \texttt{transfer\_info\_VIS\_dng.json} (respectively \texttt{transfer\_info\_NIR\_dng.json}).
        The other files (in particular the \texttt{*\_RawTherapee.tif}) are deleted.
  \item[\texttt{cameras} and \texttt{FlightAnalytics}] left as they are.
  \item[\texttt{ImgIRdrone}] the images \texttt{VIS\_XXXX.jpg} and \texttt{VIR\_XXXX\_local.jpg}, the files \texttt{VIS\_XXXX\_motion\_model.npy}, the images
        \texttt{NDVI\_XXXX\_local.jpg} and \texttt{NIR\_XXXX\_local.jpg} (if they exist), \texttt{shoot\_points.json} and \texttt{shoot\_points\_failed\_summary.json}.
        The large \texttt{imgMultispectral} folder and the individual ``redo'' folders are deleted.
  \item[\texttt{mapping\_MULTI}] left as it is: it contains the items directly usable in QGIS and has already been slimmed down in Step~6.
  \item[\texttt{Synchro}] only \texttt{check\_time\_alignment.png}, \texttt{check\_time\_before\_alignment.png} and \texttt{desktop.ini} are kept; the \texttt{VIS}, \texttt{NIR} and \texttt{Check\_ARUCO} folders are deleted.
\end{description}

\paragraph{Why this choice?} The time shift found in Step~3 and the list of image pairs (\texttt{AerialPhotography\textbackslash shoot\_points.json}) sum up the synchronization:
the \texttt{Synchro} images are no longer needed. The \texttt{VIS\_XXXX\_motion\_model.npy} files store the VIS--NIR alignment of each \emph{shoot point}, including the alignments
recovered by the ``redo'' procedure (Step~4). To resume the mission, you therefore only need to run Step~4 again, without losing these corrections, from the original images that were kept.
The \texttt{.jpg} images that are kept also make it easy to find an image quickly to illustrate a report.

\figureIRD[fig:step7]{step6_window.png}{Archiving window: the space that archiving frees is shown as soon as the window opens.}{0.6}

\section{What archiving frees (dry run)}\label{btn:step7_simulate}
As soon as the window opens, the folder is analyzed \emph{without deleting anything} (a dry run). The card under the warning gives the answer:
\begin{itemize}
  \item \texttt{12.3 GB can be freed}, with, in big, the size of the mission folder before and after archiving (for example \texttt{15.1 GB → 2.8 GB}) and the number of files that will be deleted. The details, folder by folder (files to delete, files kept, folders deleted), are in the area below;
  \item \texttt{Archiving impossible}, with the required files that are missing (see below): nothing can be deleted;
  \item \texttt{Mission already archived}, with the date of the archiving and, in big, the size of the folder;
  \item \texttt{Nothing to delete}, when the folder only holds what archiving keeps.
\end{itemize}
Check that the list matches what you expect before archiving.

\section{\bouton{Archive and free \ldots} button}\label{btn:step7_archive}
The button shows the space that will be freed (for example \bouton{Archive and free 12.3 GB\ldots}). It is greyed out while the folder is analyzed, and when nothing can be deleted. The folder is analyzed again when you click, then the user is asked to confirm, as shown in Figure~\ref{fig:step7_confirm_archving}. The message also shows the disk space that will be freed and states that the deletion will be permanent. Two buttons, \texttt{OK} and \texttt{Cancel}, confirm or cancel. For safety, and to avoid a hasty click on \texttt{Ok}, the user must type exactly the same mission folder name as the one shown (respecting case).

\figureIRD[fig:step7_confirm_archving]{step6_confirm_archiving.png}{Confirmation message for the archiving request.}{0.3}

 If OK is clicked, IRDrone runs the actual archiving:

\begin{enumerate}
  \item analyzes the folder and checks that the required files exist (\texttt{mission\_parameters.json}, \texttt{shoot\_points.json}, \texttt{time\_line.json}, \texttt{transfer\_info\_*\_dng.json}, etc.).
        If any are missing, the archiving is \textbf{refused} and the list of missing files is shown; nothing is deleted. Missing folders are never recreated;
  \item shows the number of files and the space that will be deleted, then asks for confirmation: you must \emph{type the name of the mission folder};
  \item deletes the files. The progress bar shows the progress;
  \item if a file is locked (open in another program, for example QGIS), the error is recorded and the deletion continues. Then close the programs concerned and run the archiving again;
  \item writes the files \texttt{archive\_manifest.json} and \texttt{archive\_log.txt} at the root of the mission.
\end{enumerate}

A mission that has already been archived is detected and is not archived a second time.


 A special icon is given to archived folders so you can see at a glance that a mission is archived (Figure~\ref{fig:Air_Mission_folder}).


\paragraph{Archiving and processing progress.} At the end of each step, the mission folder reflects how far the processing of the mission has progressed.

Archiving should only be run at the very end of the study. Note that the original images of the mission (converted to dng format for the SJCam) are always archived. You therefore keep the field data. So there is no need to archive the original images from the DJI and SJCam SD cards (but you can also archive them for safety).

\section{\bouton{Close} button}\label{btn:step7_close}
Closes the window and returns to the main dialog. The \bouton{?} button opens this help at the page of this chapter.

%==========================================================
\chapter{Visible only missions}\label{ch:visible_only}
\section{Without the infrared camera}\label{sec:visible_only}
A mission can be flown with the drone camera only, without the SJCam: the type \bouton{Visible only (drone DNG)} is chosen at the top of the first page of Step~1 (Section~\ref{btn:step1_mission_type}). It is stored in the mission (\texttt{"mission type": "visible"} in \texttt{FlightAnalytics/mission\_parameters.json}; the missions without this key are multispectral). The steps are the same, minus what concerns the infrared:
\begin{itemize}
	\item \textbf{Step 1}: only the SD card of the drone is needed; the copy has three stages (copy, EXIF, ground altitude). The mission folder has no \texttt{AerialPhotography/NIR} folder (\texttt{Synchro} only holds the thumbnails of the time line of Step~2).
	\item \textbf{Step 2}: the time line shows the visible images only; there is no rotating target to look for and no synchronisation phase: only the survey phase is checked. The best mapping works as for a multispectral mission. \bouton{Continue with development} saves the survey and the visible shoot points (\texttt{AerialPhotography/shoot\_points.json}, one point per visible image, without infrared) and opens Step~4.
	\item \textbf{Step 3} is not needed: its tile reads ``Not needed'' and \bouton{Continue} goes straight from Step~2 to Step~4.
	\item \textbf{Step 4} (``Develop'') develops the visible DNG images of the survey (or the best mapping selection): colour, correction of the lens distortion and the same crop as the multispectral images, so that the camera model of OpenDroneMap is the same. Each image gives \texttt{ImgIRdrone/VIS\_XXXX.jpg} (review) and the three bands \texttt{ImgIRdrone/imgMultispectral/IRD\_XXXX\_1..3.tif} (red, green, blue) for the map.
	\item \textbf{Step 5}: the review shows the developed visible images; there is nothing to realign. \bouton{Good: keep for the map} (key \texttt{Y}) marks the image checked and shows the next one to check; a blurred or badly exposed image is flagged (key \texttt{F}): it is left out of the map. The points that Step~4 did not develop (grey on the map) are developed with a right-click on the map.
	\item \textbf{Step 6}: OpenDroneMap builds a three band orthophoto (red, green, blue) instead of four. There is no infrared false colour layer, no VIR nor NDVI; the terrain model is computed as usual.
	\item \textbf{Step 7}: the archive keeps the same files, without the infrared ones.
\end{itemize}

%==========================================================
\appendix
\chapter{Alignment of the VIS and NIR images (Step 4)}\label{sec:ann_algorithme}
This appendix explains in detail how IRdrone superimposes the near infrared image on the visible image. It is written for a reader with a scientific background (linear algebra, a little signal processing) who is not a specialist of computer vision. The first sections give the geometric and signal processing basics, the following ones describe each stage of the algorithm with the diagnostic traces of real alignments, then the methods that were tried and abandoned, the history of the development, and the known limits. The last section (\S\ref{sec:ann_theorique}) deals with the physical origin of the misalignment between the two cameras (theoretical ``pitch'' and ``yaw'' angles).

The code is in \texttt{automatic\_registration.py} (functions \texttt{align\_raw} and \texttt{coarse\_alignment}), in the \texttt{registration} folder (\texttt{rigid.py}: pyramid, \texttt{cost.py}: representation and cost, \texttt{warp\_flow.py}: resampling) and in \texttt{irdrone/semi\_auto\_registration.py} (rotation of the camera, \texttt{manual\_warp}). The figures of this appendix are produced by \texttt{help/make\_alignment\_figures.py} from two shoot points:
\begin{itemize}
	\item \textbf{point A}: image \texttt{VIS\_0201} of the flight of 25 January 2022 (Blassac, 90\,m above the ground, time difference between the two shots 0.0\,s), a typical case where the alignment succeeds;
	\item \textbf{point B}: the sample pair of the repository (\texttt{samples/\_visible.DNG}, \texttt{samples/\_nir.RAW}, May 2022), a difficult case where the automatic alignment fails (\S\ref{sec:ann_limits}).
\end{itemize}
Both were aligned from the command line with the \texttt{--debug} option (\S\ref{sec:ann_traces}), i.e. with the field extension 1.6 of the command line instead of 1.4 in Step~4.

\section{The problem}\label{sec:ann_problem}
\subsection{What ``aligning'' means}
The drone carries two cameras (Figure~\ref{fig:ann_pair}): the camera of the DJI Mavic Air~2, which records the visible light (VIS, red, green and blue), and an SJCam M20 whose infrared blocking filter was removed and replaced by a filter that only lets the near infrared through (NIR, beyond about 720\,nm). The vegetation indices such as the NDVI combine, \emph{pixel by pixel}, the red of the visible image and the near infrared:
\begin{equation}
	\text{NDVI} = \frac{\text{NIR} - \text{R}}{\text{NIR} + \text{R}}.
\end{equation}
This formula only makes sense if the two values of a pixel come from the same point of the ground. Aligning (one also says \emph{registering}) the two images means finding, for every pixel of the visible image, the position of the same ground point in the infrared image, then resampling the infrared image on the grid of the visible one. An error of a few pixels already creates false index values along every contrasted edge (a hedge, a road, the shadow of a tree).

\figureIRD[fig:ann_pair]{align_pair_raw.jpg}{Point B as recorded by the two cameras. Besides the different fields of view and the strong distortion of the infrared camera, the contrasts differ: the road is dark in both bands but its verges are bright in the visible only, the vegetation is green and darker than the dry grass in the visible, bright in the near infrared.}{1}

\subsection{Why it is difficult}
Three difficulties add up:
\begin{enumerate}
	\item \textbf{The two cameras do not look in the same direction.} The visible camera is held vertical by its gimbal, the infrared camera is fixed to the body of the drone and follows its tilts (\S\ref{sec:ann_theorique}). The difference reaches several degrees, i.e. several hundred pixels.
	\item \textbf{The two optics are very different.} The infrared camera has a very wide field of view and a strong distortion; straight lines appear curved (\S\ref{sec:ann_distortion}).
	\item \textbf{The two bands do not see the same thing.} Vegetation, water, soil and roads do not have the same reflectance in the visible and in the near infrared. A contrast can be inverted from one band to the other (Figure~\ref{fig:ann_pair}, and Figure~\ref{fig:step4_pair_VIS_NIR} where the river is black in the infrared). The classical methods, which compare intensities or match characteristic points (SIFT, ORB), are therefore unreliable (see the surveys of Zitová and Flusser~[\hyperref[ann_zitova]{19}] and Jiang et al.~[\hyperref[ann_jiang]{20}]).
\end{enumerate}
The first two difficulties are solved by \emph{geometry}: a model of each camera, obtained by calibration, and the knowledge that the two images differ mostly by a rotation. The third one is solved by \emph{comparing edges rather than intensities}.

\subsection{How accurate must it be?}\label{sec:ann_accuracy}
A vegetation index is computed pixel by pixel. If the infrared image is shifted by $\delta$ pixels, the index read at a pixel is, to first order,
\begin{equation}
	\text{NDVI}_{\text{measured}}(p) \approx \text{NDVI}(p) + \frac{\partial\, \text{NDVI}}{\partial\, NIR}\; \nabla NIR(p) \cdot \delta,
\end{equation}
where $\delta$ is now the displacement vector: the error is proportional to the misalignment \emph{and} to the local contrast. In a uniform field nothing is visible; along an edge the error is maximal. Consider the edge between a meadow (NDVI 0.8) and a road (NDVI 0.1), spread over about 4 pixels by the optics: the index varies by about 0.17 per pixel across the edge. A misalignment of 1 pixel creates an error of the order of 0.1 to 0.2 along the whole edge, a misalignment of 0.1 pixel an error of the order of 0.02. A misalignment produces a typical signature: two thin fringes of opposite signs on either side of every edge (too high on one side, too low on the other), all oriented in the same direction. These fringes are what the eye detects in the review, and what the false color superimpositions of this appendix show in green and magenta. The target is therefore an accuracy of a fraction of a pixel, at full resolution and over the whole image.

\section{Geometry of a camera}\label{sec:ann_camera}
The reference textbook is \emph{Multiple View Geometry in Computer Vision} by Hartley and Zisserman~[\hyperref[ann_hz]{1}] (chapters 2 and 6 for the camera model, 7.4 for distortion, 8.4 for rotating cameras). Szeliski's book~[\hyperref[ann_szeliski]{2}] is a more accessible introduction (chapter 2).

\subsection{The pinhole model and homogeneous coordinates}
To first order, a camera is a \emph{pinhole}: every ray of light goes through the optical center $C$ before reaching the sensor. Let $(X,Y,Z)$ be the coordinates of a point of the scene in the frame of the camera ($Z$ along the optical axis). Its image is at the pixel
\begin{equation}
	u = f\,\frac{X}{Z} + c_x, \qquad v = f\,\frac{Y}{Z} + c_y,
\end{equation}
where $f$ is the focal length \emph{expressed in pixels} (physical focal length divided by the size of a pixel) and $(c_x, c_y)$ the \emph{principal point}, intersection of the optical axis with the sensor, close to the center of the image. The division by $Z$ is the perspective: an object twice as far appears twice as small.

This division is not linear; it becomes linear by writing the pixel in \emph{homogeneous coordinates}, a vector of three numbers defined up to a factor: $(u, v, 1) \sim (\lambda u, \lambda v, \lambda)$. Then
\begin{equation}\label{eq:ann_pinhole}
	\lambda \begin{pmatrix} u \\ v \\ 1 \end{pmatrix} = K \begin{pmatrix} X \\ Y \\ Z \end{pmatrix},
	\qquad
	K = \begin{pmatrix} f & 0 & c_x \\ 0 & f & c_y \\ 0 & 0 & 1 \end{pmatrix}.
\end{equation}
The matrix $K$ of the \emph{intrinsic parameters} (or calibration matrix) describes the camera alone. Conversely, $K^{-1}(u, v, 1)^T$ is the \emph{direction} of the ray that arrives on the pixel $(u,v)$: an image is a set of directions, and the depth $Z$ is lost.

\begin{table}[htbp]\centering
\begin{tabular}{@{}lll@{}}
\hline
 & \textbf{DJI Mavic Air~2 (VIS)} & \textbf{SJCam M20 (NIR)} \\ \hline
Image size (pixels) & $3992 \times 2992$ & $4600 \times 3448$ \\
Focal length $f$ (pixels) & 2898 & 2311 \\
Principal point $(c_x, c_y)$ & (2032, 1511) & (2348, 1725) \\
Horizontal field of view & $68^\circ$ & more than $110^\circ$ \\
Distortion $k_1, k_2, k_3$ & $0.185,\ -0.455,\ 0.245$ & $-0.337,\ 0.131,\ -0.025$ \\ \hline
\end{tabular}
\caption{Calibrations of the two cameras (\texttt{calibration/DJI\_RAW} and \texttt{calibration/M20\_RAW}, RAW images). At the focal length of the visible camera, $1^\circ$ of rotation moves the image by $f \tan 1^\circ \approx 51$ pixels.}
\label{tab:ann_calib}
\end{table}

\subsection{Distortion}\label{sec:ann_distortion}
Real lenses deviate from the pinhole model: the image of a straight line is curved, more and more towards the edges. IRdrone uses the polynomial model of Brown and Conrady~[\hyperref[ann_brown]{3}], the one of OpenCV. With the normalized coordinates $x = X/Z$, $y = Y/Z$ and $r^2 = x^2 + y^2$, the ideal point $(x, y)$ is observed at
\begin{equation}\label{eq:ann_distortion}
	x_d = x\,(1 + k_1 r^2 + k_2 r^4 + k_3 r^6) + 2p_1 xy + p_2 (r^2 + 2x^2),
\end{equation}
and similarly for $y_d$; the pixel is then $K\,(x_d, y_d, 1)^T$. The coefficients $k_i$ describe the radial distortion (barrel when $k_1 < 0$, pincushion when $k_1 > 0$), $p_1, p_2$ a small decentering of the lenses.

Figure~\ref{fig:ann_distortion} shows the magnitudes. For the visible camera the correction does not exceed 40 pixels. For the infrared camera it reaches 210 pixels at the left and right edges and 375 pixels in the corners of the field of the visible camera, and 800 pixels at $50^\circ$ from the axis. Its strongly negative $k_1$ is a barrel distortion close to that of a ``fish-eye''. The polynomial is only valid in the area where calibration targets were seen: beyond about $57^\circ$ from the axis it decreases again, so the model is no longer invertible, and the edges of the infrared image (more than $110^\circ$ of horizontal field) are outside its validity. These areas are never used, since the visible field only reaches $41^\circ$ in the corners.

\figureIRD[fig:ann_distortion]{align_distortion.png}{Distortion of the two cameras. Left and center: a square grid of straight lines on the ground, as it is recorded on the sensor (black frame) of each camera; in red, the field of the visible camera as seen by the infrared camera. Right: radial shift of a point as a function of its angle to the optical axis.}{1}

\subsection{Calibration}
The intrinsic parameters $f, c_x, c_y$ and the distortion coefficients are measured once for each camera by Zhang's method~[\hyperref[ann_zhang]{4}]: about thirty photographs of a planar checkerboard in various orientations; the corners of the squares are detected, and all the parameters (plus the position of the board in each photo) are estimated by minimizing the distance between the observed corners and the corners predicted by equations \eqref{eq:ann_pinhole} and \eqref{eq:ann_distortion} (\texttt{calibration/cameracalibration.py}, OpenCV). This calibration is fixed: it does not depend on the flight.

\subsection{Vignetting (lens shading)}\label{sec:ann_shading}
A lens does not transmit the same amount of light to every pixel: the image of a uniform white wall is darker towards the edges. Several effects add up: the natural fall-off in $\cos^4\theta$ of the illumination with the angle $\theta$ to the optical axis (the pupil is seen obliquely, the light arrives obliquely on the sensor and travels further), the mechanical vignetting by the lens mounts, and the angular acceptance of the micro-lenses of the sensor~[\hyperref[ann_goldman]{21}]. This multiplicative \emph{shading} is described by a gain map $g(x, y)$: the corrected image is $g \cdot I$.

Both maps were measured in September 2021 on photographs of a uniform white chart (\texttt{calibration/cameravignetting.py}). For the visible camera the map is stored at 1/10 of the resolution, for each color channel (\texttt{calibration/DJI\_RAW/shading\_calibration.npy}). For the infrared camera a radial profile is fitted by hand (cubic curve through control points, on a polar representation of the white image), then turned into a map by revolution around the center (\texttt{get\_polar\_shading\_map}). Figure~\ref{fig:ann_shading} shows them: the visible camera needs a gain of about 1.15 in the corners; the infrared camera a gain of up to 1.3 at 80\,\% of the half diagonal, beyond which its fish-eye image circle ends and the gain is reduced again so as not to amplify the noise of unlit areas. The correction is applied to the linear RAW data, before anything else (\texttt{irdrone/process.py}); since September 2021 the RawTherapee profiles have their own vignetting correction turned off, because it did not match these cameras (``not negligible on DJI'').

For the vegetation indices, the shading is a radiometric error: a gain error of 15\,\% on the infrared band alone changes an NDVI of 0.67 into 0.70. For the alignment, it matters through the gradients. By the product rule, $\nabla(g I) = g\, \nabla I + I\, \nabla g \approx g\, \nabla I$, since $g$ varies slowly: the gradient energies of \S\ref{sec:ann_msr} are multiplied by $g^2$. Without correction, the energy of the same edge would vary by a factor of up to 1.3 (visible) and 1.7 (infrared) between the center and the edges, with different patterns in the two cameras; the criterion \eqref{eq:ann_cost}, which compares absolute energies, would then see differences where there is none.

\figureIRD[fig:ann_shading]{align_shading.png}{Lens shading correction of the two cameras. (a) Gain applied as a function of the distance to the center of the image, for each color channel (solid: infrared camera, the three channels of its sensor all record the near infrared; dashed: visible camera, profile along the diagonal). (b), (c) Gain maps (green channel).}{1}

\section{Two cameras that rotate: the homography}\label{sec:ann_homography}
\subsection{A rotation of the camera is a homography of the image}
Suppose that the two cameras have the \emph{same optical center} and only differ by a rotation $R$ (a $3 \times 3$ orthogonal matrix). The ray that arrives on the pixel $p_N$ of the infrared camera has the direction $K_{NIR}^{-1} p_N$ in the frame of that camera, $R\, K_{NIR}^{-1} p_N$ in the frame of the visible camera, and therefore arrives on the pixel
\begin{equation}\label{eq:ann_rotation_homography}
	p_V \sim K_{VIS}\, R\, K_{NIR}^{-1}\, p_N = H\, p_N
\end{equation}
of the visible camera (after correction of the distortions). The $3 \times 3$ matrix $H$, defined up to a factor, is a \emph{homography}: the most general transformation that maps straight lines to straight lines (Hartley and Zisserman~[\hyperref[ann_hz]{1}], \S2.3 and \S8.4). It has 8 degrees of freedom. A remarkable point: equation \eqref{eq:ann_rotation_homography} does not involve the depth $Z$: whatever the relief, a pure rotation of the camera is a homography of the image. This is the principle of the panoramas made by turning on oneself.

The rotation is parameterized by three angles: the \emph{yaw} (rotation about the vertical axis of the image, which moves it horizontally), the \emph{pitch} (which moves it vertically) and the \emph{roll} (rotation about the optical axis, which turns the image about its center). In the code, $R$ is built from the rotation vector $(-\text{pitch}, \text{yaw}, \text{roll})$ (Rodrigues formula, \texttt{manual\_warp}). For small angles, a yaw $\theta$ moves the center of the image by $f \tan\theta$ pixels, i.e. about 51 pixels per degree; a roll $\varphi$ moves a corner, at about 2500 pixels from the center, by $2500\,\varphi$ (44 pixels per degree).

\subsection{Why the rotation is (almost) enough: parallax}
The two optical centers are not exactly at the same place: they are 46\,mm apart on the drone (\texttt{CNIRCVIS\_0} in \texttt{config.py}), and the drone moves between the two shots. A distance $B$ between the centers (the \emph{baseline}) creates a \emph{parallax}: a point at the distance $Z$ is shifted by
\begin{equation}
	d = f\,\frac{B}{Z} \quad \text{pixels}.
\end{equation}
Only the \emph{variations} of this shift matter, the constant part being absorbed by the rotation. With $B = 46$\,mm, at 90\,m: $d = 2898 \times 0.046 / 90 \approx 1.5$ pixel, and a relief of 10\,m changes it by 0.2 pixel: negligible. If the two shots are $\Delta t = 0.5$\,s apart and the drone flies at 3\,m/s, $B = 1.5$\,m: $d \approx 48$ pixels, and a 10\,m relief (a hedge of trees, a slope) now changes it by $f B \Delta Z / Z^2 \approx 5$ pixels. This is the reason why Step~4 only aligns shoot points whose time difference is small, and why a \emph{local} correction is added to the global model (\S\ref{sec:ann_local}). At very low altitude (take-off, landing), the parallax dominates and the alignment is meaningless.

Two remarks make this picture less pessimistic. First, a \emph{plane} also maps to a homography: if the ground is a plane of normal $n$ at the distance $d$, two cameras related by a rotation $R$ \emph{and} a translation $t$ see it through
\begin{equation}\label{eq:ann_plane_homography}
	H = K_{VIS} \Big( R - \frac{t\, n^T}{d} \Big) K_{NIR}^{-1}
\end{equation}
(Hartley and Zisserman~[\hyperref[ann_hz]{1}], \S13.1). The homography fitted in the pyramidal refinement (\S\ref{sec:ann_pyramid}) therefore absorbs the parallax of a flat ground, even sloping: only the relief that departs from the mean plane remains, which the local correction handles. Second, both sensors are CMOS sensors read line by line (\emph{rolling shutter}): the lines of an image are not exposed at the same instant but over some tens of milliseconds. A drone that moves or vibrates during this time slightly shears or waves the image, differently in the two cameras. This deformation is not a homography either; when it is smooth, the local correction absorbs it as well.

\subsection{The virtual camera}
The visible image is first corrected for its distortion: it becomes the image of an ideal pinhole camera, the \emph{reference}. The infrared image is then resampled directly into this reference camera, through equation \eqref{eq:ann_rotation_homography} and its own distortion model, but on a canvas wider than the visible image (1.4 times by default, 1.6 in the REDO script): this is the \emph{virtual camera} of Figure~\ref{fig:ann_virtual}. If the rotation were the identity, the visible image would correspond exactly to the central rectangle; the search for the rotation consists in finding where it actually lies.

\figureIRD[fig:ann_virtual]{align_virtual_camera.jpg}{Point B. Left: the infrared image corrected for its distortion and resampled in the virtual camera (focal length of the visible camera, field 1.4 times wider). Dashed: position of the visible image if the two cameras were parallel; red: position found by the coarse search (yaw $8.2^\circ$, pitch $1.9^\circ$). Right: the visible reference. Outside the visible field the correction of the distortion is less accurate (the road bends on the left).}{1}

\section{Comparing two spectral bands}\label{sec:ann_similarity}
To find the rotation, we need a \emph{similarity criterion}: a number that says how well two images are superimposed, and that is optimal for the right superposition.

\subsection{Why intensities cannot be compared}
For two images of the same band, the simplest criterion is the sum of squared differences, $\text{SSD} = \sum_p (I_1(p) - I_2(p))^2$, or the correlation. Both assume that the same object has the same (or a proportional) intensity in the two images. Between the visible and the near infrared this is false: the vegetation is dark in the visible and bright in the infrared, the road verges are bright in the visible only (Figures~\ref{fig:ann_pair} and \ref{fig:ann_features}). A classical answer to such ``multimodal'' problems is the mutual information~[\hyperref[ann_viola]{5}], which measures a statistical dependence instead of an equality; Another approach combines characteristic points with a gradient similarity~[\hyperref[ann_vayssade]{18}]. IRdrone uses a simpler and more local idea: the \emph{edges} are shared even when the contrasts are not.

\subsection{A worked example: inverted contrasts}\label{sec:ann_inversion}
Figure~\ref{fig:ann_inversion} shows the difficulty on a synthetic cut across a landscape: vegetation, bare soil, a road with bright verges, a crop. In the visible, the vegetation is darker than the soil; in the near infrared it is brighter (the chlorophyll absorbs the red, the leaf structure scatters the infrared strongly). The road is dark in both bands, its verges bright in the visible only. The sum of squared differences and the correlation of the intensities are both minimal for a shift of 60 pixels, the limit of the search range: shifting the visible profile until its dark vegetation faces the dark road of the infrared ``explains'' the intensities better than the right superposition.

A \emph{global} inversion ($I_{NIR} = a - b\, I_{VIS}$) would still be found: the correlation would be $-1$ at the right position, and the phase correlation (\S\ref{sec:ann_fourier}) would give a negative peak, since a change of sign is a phase shift of $\pi$ at every frequency. But here some materials are inverted and others are not: no function $I_{NIR} = \phi(I_{VIS})$ links the two bands. The mutual information~[\hyperref[ann_viola]{5}] tolerates any one-to-one function $\phi$, but needs many pixels per block to estimate the joint histogram, and fails as soon as two materials of the same visible color have different infrared values. The edges, on the other hand, are at the same place in both bands: this is what the energy maps keep (Figure~\ref{fig:ann_inversion}b), and the criterion on these maps finds the right shift.

\figureIRD[fig:ann_inversion]{align_inversion.png}{Synthetic cut across a landscape, with noise and the blur of each camera. (a) Intensities: the vegetation and the bare soil are inverted between the bands, the road is not, its verges are visible in the visible band only. (b) Squared gradients, after smoothing: the edges are at the same place whatever their sign. (c) Criteria as a function of the shift of the visible profile: the intensity criteria choose $+60$ pixels, the criterion of IRdrone the right position.}{1}

\subsection{Preparing the images}\label{sec:ann_preparation}
The alignment works on images prepared as follows (\texttt{irdrone/process.py}): the RAW file is demosaiced with a neutral linear profile (RawTherapee, color matrix of the camera, no tone curve), the black level is subtracted, the lens shading is corrected (\S\ref{sec:ann_shading}); the image is then stretched linearly between its 2nd and 98th percentiles (green channel) to 8 bits, and converted to luminance ($0.299R + 0.587G + 0.114B$). The stretch absorbs the difference of exposure between the two cameras. No gamma or logarithm is applied: the gradient energies are proportional to the squared contrast in linear units, so that a shadow or the vegetation in the visible, which are dark, contribute less than a bright road. Working on the logarithm of the intensities would make every edge count for its \emph{relative} contrast, and the criterion insensitive to any residual gain (exposure, shading); this has not been tested.

\subsection{The feature maps: directional gradient energies}\label{sec:ann_msr}
Each image is converted to luminance (gray levels), smoothed by a Gaussian filter of standard deviation $\sigma$ (5 pixels for the visible, 3 for the less sharp infrared; 3 and 1 for the coarse search), then replaced by four maps (\texttt{multispectral\_representation}, mode \texttt{LAPLACIAN\_ENERGIES}):
\begin{equation}
	E_x = (\partial_x I)^2, \quad E_y = (\partial_y I)^2, \quad E_{\nearrow} = (\partial_{\nearrow} I)^2, \quad E_{\searrow} = (\partial_{\searrow} I)^2.
\end{equation}
The horizontal and vertical derivatives are computed with Scharr filters~[\hyperref[ann_scharr]{6}] (a $3 \times 3$ derivative with an optimized smoothing, better isotropic than Sobel), the diagonal ones with the $2 \times 2$ Roberts cross~[\hyperref[ann_roberts]{7}]. These maps are called the \emph{multispectral representation} (MSR) in the code. The idea comes from Irani and Anandan~[\hyperref[ann_irani]{9}], who aligned images of different sensors (visible and infrared among others) through such directional derivative energies.

The square is the key: an edge from dark to bright and an edge from bright to dark give the same energy. A contrast inversion between the bands therefore becomes invisible, while the \emph{position} of the edge is kept. The four directions keep the orientation of the edges, which helps to distinguish, for example, a road from the perpendicular edge of a field. The smoothing removes the fine textures (grass, leaves) that are different in the two bands and keeps the large structures. Figure~\ref{fig:ann_features} shows these maps on a detail of point B: the luminances are very different, the energies of the edges of the road are at the same place.

For a reader from signal processing: squaring a band-pass signal, then smoothing it, is an \emph{envelope detector} (the demodulation of an amplitude modulated radio signal). The derivative is a high-pass filter that keeps the edges and textures, the square makes the signal positive, and the smoothing of the pyramid keeps its envelope, i.e. \emph{where} there is structure, regardless of its fine phase. This is why the pyramid is built on the energy maps and not on the images (\texttt{compute\_pyramid} on the MSR in \texttt{pyramidal\_search}): reduced 32 times, an image of a fine texture becomes a uniform gray and its edges vanish, while its energy map still shows where the hedges, roads and field edges are (Figure~\ref{fig:ann_coarse}).

\figureIRD[fig:ann_features]{align_feature_maps.jpg}{Directional gradient energies on a detail of point B ($600 \times 600$ pixels). The luminances of the two bands are not comparable (verges bright in the visible, dark in the infrared), but the edges of the road produce the same lines in the energy maps. Some edges exist in one band only (the four lines of the visible against two in the infrared): the criterion must tolerate them.}{1}

\subsection{The similarity criterion: total gradient of the difference}\label{sec:ann_ntg}
The criterion derives from the \emph{Normalized Total Gradient} (NTG) of Chen et al.~[\hyperref[ann_ntg]{8}]. For two images $f$ and $g$:
\begin{equation}
	\text{NTG}(f, g) = \frac{\sum_p \|\nabla (f - g)(p)\|_1}{\sum_p \|\nabla f(p)\|_1 + \sum_p \|\nabla g(p)\|_1}.
\end{equation}
The idea: when two images are well superimposed, their common edges cancel in the difference $f-g$, which becomes smooth; its total gradient (sum of the absolute values of the derivatives) is then minimal. When they are shifted, each edge appears twice in the difference, with opposite signs, and the total gradient increases. The norm $\|\cdot\|_1$ (absolute values rather than squares) makes the criterion robust to the edges that exist in one image only: they add a roughly constant cost that hardly depends on the shift, instead of dominating it. The denominator makes the criterion independent of the amount of texture.

IRdrone applies it to the energy maps: for a shift $u$ of the reference,
\begin{equation}\label{eq:ann_cost}
	C(u) = \sum_{c=1}^{4} \sum_p \Big( \big|S_x * (E^{NIR}_c - E^{VIS}_c(\cdot + u))\big|(p) + \big|S_y * (E^{NIR}_c - E^{VIS}_c(\cdot + u))\big|(p) \Big),
\end{equation}
where $S_x, S_y$ are the Sobel derivative filters (\texttt{cost\_surface\_NTG} in \texttt{cost.py}). Only the numerator of the NTG is computed. Since the Sobel filter is linear, $S * (a - b) = S*a - S*b$: the gradients of both images are computed once, and each shift is then a mere sum of absolute differences (15 times faster, November 2026). A sum of squared differences on the same maps is also implemented (\texttt{SSD}) but not used.

Figure~\ref{fig:ann_profiles} compares the criteria on point A, as a function of a horizontal shift around the right alignment. The SSD of the luminances has a wide valley: it locates the alignment, but with little precision, and it relies on the intensities being comparable, which is the case on this winter scene but not in summer. The NTG has a narrow and deep well: it is very precise, but only ``sees'' the right position within a few pixels, hence the need for a good starting point and for a coarse-to-fine strategy.

\figureIRD[fig:ann_profiles]{align_cost_profiles.png}{Cost as a function of a horizontal shift (point A, after alignment, computed at 1/4 resolution), normalized between 0 and 1. (a) The visible image against itself. (b) Visible against infrared on a block with trees and a hedge. (c) On a block of meadow only, with few structures: the NTG still has its minimum at the right place but outside the narrow well the curve is no longer informative.}{1}

\section{Coarse search}\label{sec:ann_coarse}
The coarse search finds most of the difference of orientation (function \texttt{coarse\_alignment}):
\begin{enumerate}
	\item the infrared image is resampled in the virtual camera with the initial angles: zero in Step~4, the angles set by hand during a realignment (\S\ref{btn:review_redo});
	\item the energy maps of both images are reduced by a factor of 32 ($125 \times 94$ pixels for the visible), the visible map is placed at the center of a canvas of the size of the infrared map;
	\item an \emph{exhaustive search} computes the cost \eqref{eq:ann_cost} for every translation of the visible map inside the infrared map: with the field extension 1.4 of Step~4, $\pm 25$ pixels at this scale, i.e. about $\pm 800$ pixels at full resolution or $\pm 15^\circ$; with the extension 1.6 of the command line, $\pm 37$ pixels, $\pm 1200$ pixels or $\pm 22^\circ$ (Figure~\ref{fig:ann_coarse});
	\item the best position is located with sub-pixel precision (below);
	\item at small angles a horizontal translation $t_x$ is a yaw and a vertical translation $t_y$ a pitch: $\text{yaw} = \arctan(t_x/f)$, $\text{pitch} = \arctan(t_y/f)$;
	\item the infrared image is resampled once again, at full resolution, with the corrected angles.
\end{enumerate}
The roll is not searched: it keeps its initial value (\S\ref{sec:ann_fourier}).

\paragraph{Choosing the minimum.} The absolute minimum of the cost surface is not always the right one: a textureless area or the black border of the canvas can produce a wide low valley. The right superposition produces a narrow and deep well (Figure~\ref{fig:ann_coarse}). The code therefore looks for the position where the surface is the most sharply curved: at every position, a quadratic surface is fitted on the $3\times3$ neighborhood, and the position with the largest determinant of its Hessian matrix (product of the two principal curvatures) is kept (\texttt{minimum\_cost\_max\_hessian}).

\paragraph{Sub-pixel refinement.} Around the selected position, the cost is approximated by a quadratic form~[\hyperref[ann_hasinoff]{10}]
\begin{equation}
	C(u) \approx \tfrac{1}{2}\, u^T A\, u + b^T u + c,
\end{equation}
whose coefficients are obtained by small convolutions of the $3\times3$ neighborhood. Its minimum is at $u^* = -A^{-1} b$ (one Newton step), kept if it is less than one pixel away. At 1/32 resolution, one pixel is 32 pixels at full resolution: this refinement makes the coarse search accurate to a few pixels. For point A, the translation found is $(21.7, 163.9)$ pixels, i.e. a yaw of $0.43^\circ$ and a pitch of $3.24^\circ$.

\figureIRD[fig:ann_coarse]{align_coarse_search.png}{Coarse search of point A at 1/32 resolution. Left: the cost surface over the whole search range ($\pm 1200$ pixels): a sharp well in a smooth landscape. Center: zoom on the center of the cost. Right: determinant of the Hessian, maximal at the bottom of the well.}{1}

\section{Pyramidal refinement}\label{sec:ann_pyramid}
After the coarse search the residual misalignment is a few pixels to a few tens of pixels, and it is not uniform: residual roll, errors of the calibrations, parallax. The refinement estimates it by blocks, on an image pyramid (\texttt{pyramidal\_search} in \texttt{rigid.py}).

\subsection{Coarse to fine}
A \emph{Gaussian pyramid}~[\hyperref[ann_burt]{11}] is the series of an image smoothed and reduced by 2, 4, 8, 16\ldots A displacement of 100 pixels at full resolution is 6 pixels at 1/16: it can be found with a small search window, quickly and with little risk of confusion. Once found, it is applied, and the next finer level only has to correct what remains. This is the classical strategy of the Lucas--Kanade alignment methods~[\hyperref[ann_lucas]{12}, \hyperref[ann_baker]{13}]. IRdrone uses the following scheme (\texttt{iterative\_scheme} in \texttt{align\_raw}):

\begin{table}[htbp]\centering
\begin{tabular}{@{}cccccc@{}}
\hline
\textbf{Scale} & \textbf{Iterations} & \textbf{Blocks} & \textbf{Block size (full res.)} & \textbf{Cost window} & \textbf{Largest vector (full res.)} \\ \hline
1/16 & 2 & $4 \times 4$ & $998 \times 748$ & $\pm 8$ px & $\pm 80$ px \\
1/16 & 2 & $5 \times 5$ & $798 \times 598$ & $\pm 6$ px & $\pm 80$ px \\
1/4 & 3 & $5 \times 5$ & $798 \times 598$ & $\pm 6$ px & $\pm 20$ px \\ \hline
\end{tabular}
\caption{Scheme of the pyramidal refinement. The cost is computed over the window shown, at the scale of the level, but the minimum is searched within $\pm 5$ pixels of the center (\texttt{brute\_force\_vector\_field\_search}): a block vector cannot exceed 80 pixels at 1/16, 20 pixels at 1/4. The iterations are numbered 1, 2, 4, 5, 7, 8, 9 in the names of the trace files (the counter is incremented once more at each change of level).}
\label{tab:ann_pyramid}
\end{table}

\subsection{One iteration}
\begin{enumerate}
	\item The energy maps of the infrared image are resampled with the current model (identity at the start).
	\item The image is split into a grid of blocks. For each block, the cost \eqref{eq:ann_cost} is computed for every shift of the search window: this gives one \emph{cost surface} per block (Figures~\ref{fig:ann_blocks_first} and \ref{fig:ann_blocks_last}).
	\item The minimum of each surface (within $\pm 5$ pixels of the center), refined to the sub-pixel as above, gives a displacement vector per block: a coarse \emph{vector field}, or optical flow.
	\item A homography is fitted on these vectors: the 16 or 25 pairs (block center, center + vector) give an over-determined system for its 8 unknowns. It is solved by the direct linear transform (DLT) inside a RANSAC loop~[\hyperref[ann_ransac]{14}], which ignores the blocks whose vector disagrees by more than 3 pixels with the majority (\texttt{geometric\_rigid\_transform\_estimation}, OpenCV \texttt{findHomography}; Hartley and Zisserman~[\hyperref[ann_hz]{1}], chapter 4).
	\item The model is updated by composition, $H_k = \Delta H_k\, H_{k-1}$, then the next iteration starts.
\end{enumerate}
A homography rather than three angles: it represents the rotation of a camera \emph{and} the residual errors of the calibrations (for point A, the final homography contains a scaling of 1.9\,\% horizontally and 1.1\,\% vertically, for instance a small error on the focal lengths).

The 8 parameters of a homography are not equally well determined by vectors measured in a field of view of $\pm 34^\circ$: the translation and the rotation are constrained by every block, the two \emph{perspective} terms (third row of $H$) only through small differences between opposite sides of the image. They are the first to absorb erroneous vectors: on point B, a few wrong vectors in the corners produce a homography that moves the corners by several hundred pixels (\S\ref{sec:ann_limits}). OpenCV normalizes the coordinates before the DLT, which makes the system numerically well conditioned (Hartley and Zisserman~[\hyperref[ann_hz]{1}], \S4.4), but cannot make up for the lack of information.

\figureIRD[fig:ann_blocks_first]{align_blocks_first.jpg}{First iteration of the refinement of point A (1/16 resolution, $4\times4$ blocks, search $\pm 8$ pixels). Left: superimposition after the coarse search (visible in green, infrared in magenta; they are gray where they agree), with the measured vectors (blue, magnified 20 times) and the vectors of the fitted homography (red). Right: the cost surface of each block; the black dot is the zero shift, the cross the minimum. The vectors are up to 50 pixels long and vary smoothly: a residual scaling and a residual pitch, which the homography represents well.}{1}

\figureIRD[fig:ann_blocks_last]{align_blocks_last.jpg}{Last iteration (1/4 resolution, $5\times5$ blocks, search $\pm 6$ pixels), on the globally aligned image. All the minima are at the center: the residual vectors (magnified 200 times) are below 2 pixels, the homography no longer changes. These residual vectors form the local correction (\S\ref{sec:ann_local}).}{1}

\figureIRD[fig:ann_convergence]{align_convergence.png}{Length of the block vectors at each iteration, for the two points. Point A (blue) converges: from 23 pixels on average to 0.6 pixel. The steps at the changes of level are due to the change of block size and of resolution (one pixel at 1/16 is 16 pixels). Point B (red) does not converge: the vectors stay at their limit, 80 pixels at 1/16 and 20 pixels at 1/4 (\S\ref{sec:ann_limits}).}{0.7}

\subsection{What a block can measure: the aperture problem}\label{sec:ann_aperture}
A block only constrains the displacement in the directions where it contains structure (Figure~\ref{fig:ann_aperture}). A corner, or two edges of different orientations, give a well with a single minimum: both components of the displacement are measured. A straight edge, a road or a row of trees, gives a \emph{valley}: the displacement across the edge is measured, the one along it is not, since sliding the block along the road does not change the cost. This is the classical \emph{aperture problem} of motion estimation. A block of meadow, whose textures differ between the two bands, gives a surface without a significant minimum: its vector is noise.

The two curvatures of the quadratic fit of \S\ref{sec:ann_coarse} (the eigenvalues of the matrix $A$) quantify this: two large curvatures for a corner, one large and one small for an edge, two small ones for a structureless block. It is the same matrix as the \emph{structure tensor} used by the Harris corner detector~[\hyperref[ann_harris]{22}] and by the Lucas--Kanade method~[\hyperref[ann_lucas]{12}], whose inverse is the covariance of the estimated displacement. The code uses its determinant (product of the curvatures) for the coarse search, but not yet as a confidence for the block vectors: a vector measured along a road counts as much as a vector measured on a crossroads in the homography fit, and only the RANSAC selection rejects the most aberrant ones (issue \#4, \S\ref{sec:ann_limits}).

\figureIRD[fig:ann_aperture]{align_aperture.png}{The aperture problem on synthetic blocks of $120 \times 120$ pixels: the NTG cost on the energy maps as a function of the shift (dark: low cost), normalized between 0 and 1, with the rise of the cost 3 pixels away from the right position (cross) in the best and worst directions. (a) A corner: a well. (b) A straight edge: a valley along the edge, the displacement along it is not measured. (c) Unrelated textures in the two images: the right position is not even a local minimum.}{0.78}

\subsection{Local alignment}\label{sec:ann_local}
A homography is a global model: it cannot represent the parallax of the relief or a local imperfection of the distortion model. The vector field of the last iteration ($5 \times 5$ vectors) is therefore kept as a \emph{local correction}: it is interpolated bilinearly between the block centers (constant beyond them) into a displacement for every pixel, which is composed with the global model (\texttt{warp\_from\_sparse\_vector\_field}). Two results are produced:
\begin{itemize}
	\item \emph{global} alignment: rotation and final homography;
	\item \emph{local} alignment: rotation, homography of the previous iteration and local displacement field. This is generally the most accurate, and the one used for the multispectral \texttt{IRD\_XXXX\_n.tif} images.
\end{itemize}

\subsection{A single resampling}
Each resampling of an image interpolates between its pixels, and therefore blurs it a little. Instead of resampling the infrared image at each stage, the final transformation is composed in a single map: for every pixel of the visible image, the position to read in the \emph{raw} infrared image is obtained by chaining the local displacement, the homography, the rotation of equation \eqref{eq:ann_rotation_homography} and the distortion of the infrared camera; the image is then interpolated only once (bicubic interpolation, \texttt{cv2.remap}). In the code this principle is called ``You Only Warp Once'' (\texttt{FULLRES\_REGISTERED\_REFINED\_WARPED\_ONCE\_ONLY.jpg}).

The resampling is done ``backwards'': one does not push each infrared pixel to its new position, which would leave holes and overlaps, but one asks, for each pixel $p$ of the output grid, which position of the raw infrared image it comes from. The chain is then read from the output towards the input, each transformation being applied in the coordinates produced by the previous one:
\begin{equation}
	p \;\xrightarrow{\;+\,v(p)\;}\; p' \;\xrightarrow{\;H^{-1}\;}\; p'' \;\xrightarrow{\;K_{NIR}\, R^{-1} K_{VIS}^{-1}\;}\; p_N \;\xrightarrow{\;\text{distortion}\;}\; p_{raw},
\end{equation}
where $v$ is the local displacement field. The composed map (two arrays, of the $x$ and $y$ coordinates to read for each output pixel) is computed exactly in floating point, and only the final reading of the image interpolates.

Each interpolation is a low-pass filter. The bilinear interpolation, a convolution by a triangle, has a frequency response in $\text{sinc}^2$: at half the Nyquist frequency it attenuates the amplitude by up to 20\,\% (depending on the sub-pixel position), and the losses of successive resamplings multiply. The bicubic interpolation~[\hyperref[ann_keys]{23}] is closer to the ideal interpolator but slightly overshoots in the middle frequencies, an effect that also accumulates. Figure~\ref{fig:ann_warp_once} shows four successive resamplings against a single one. Moreover, each intermediate image would lose its borders, outside the canvas, whereas the composed map reads the raw image directly.

\figureIRD[fig:ann_warp_once]{align_warp_once.png}{A single resampling against four. A detail ($128 \times 128$ pixels) of the visible image of point A is turned and scaled by four small successive transformations. (a) The composed transformation, applied once (bicubic). (b) The four transformations applied one after the other (bilinear): the texture is blurred. (c) The same in bicubic: sharper, but with exaggerated contrasts. (d) Power spectrum of (b) and (c) relative to (a), as a function of the spatial frequency (0.5 cycle per pixel is the Nyquist frequency).}{1} The motion model (angles, homographies, vector field) is saved in \texttt{VIS\_XXXX\_motion\_model.npy} and its readable version \texttt{.json}; when it exists, running Step~4 again reads it back instead of searching again.

Figure~\ref{fig:ann_stages} shows the progress of the alignment on a detail of point A.

\figureIRD[fig:ann_stages]{align_stages.jpg}{The four stages on a detail of point A (visible in green, infrared in magenta, gray where they agree). (a) Calibrations only: the images are about 160 pixels apart. (b) After the coarse search: about ten pixels remain (colored fringes along the shadows). (c) With the homography. (d) With the local correction: the edges coincide.}{1}

\section{Explored alternatives: phase correlation, Fourier--Mellin and polar matching}\label{sec:ann_fourier}
\subsection{Phase correlation}
The first versions (August--September 2021) estimated the translation by \emph{phase correlation}~[\hyperref[ann_kuglin]{15}]. If $g(x) = f(x - t)$, their Fourier transforms satisfy $\hat g(\omega) = \hat f(\omega)\, e^{-i\omega \cdot t}$ (shift theorem). The normalized cross-power spectrum
\begin{equation}
	\frac{\hat f\, \overline{\hat g}}{|\hat f\, \overline{\hat g}|} = e^{i\omega \cdot t}
\end{equation}
only keeps the phase, and its inverse Fourier transform is a Dirac peak at the position $t$. The method is fast (two FFTs), insensitive to a global change of brightness, and finds large translations. But it assumes that the two images are the \emph{same} signal shifted: a rotation, a scaling, a distortion or a contrast inversion spread the peak. Figure~\ref{fig:ann_phase} shows it: on point A, after the coarse search, the peak is clear; on point B it is drowned in the noise, both on the luminances and on the energy maps.

\figureIRD[fig:ann_phase]{align_phase_correlation.png}{Phase correlation surfaces between the visible image and the infrared image after the coarse search (1/8 resolution, zoom on $\pm 320$ pixels). Above, point A: a clear peak, about 2.5 times higher than the next values. Below, point B (residual rotation, contrasts inverted, large textured areas): the peak barely emerges.}{0.65}

\subsection{Fourier--Mellin: rotation and scale}
The Fourier--Mellin method~[\hyperref[ann_reddy]{16}] extends the phase correlation to rotations and scalings. The modulus of the Fourier transform does not depend on the translation; a rotation of the image rotates it, a scaling scales it. Resampled in \emph{log-polar} coordinates $(\log r, \theta)$, the modulus turns a rotation into a translation along $\theta$ and a scaling into a translation along $\log r$, which a second phase correlation finds. Rabatel and Labbé~[\hyperref[ann_rabatel]{17}] applied it to visible and near infrared drone images. In IRdrone, the scale is known from the calibrations, and only the roll was searched this way: the semi-automatic alignment of 2021 resampled both images in polar coordinates and estimated the roll by phase correlation along the angle (\texttt{mellin\_transform} and \texttt{estimate\_rotation} in \texttt{semi\_auto\_registration.py}, still present for a demonstration but not used by Step~4). Combined with the translation by phase correlation, or with SIFT points, it did not give reliable results on the multispectral images, for the reasons above. From September 2021 it was replaced by the minimization of the cost \eqref{eq:ann_cost} on the energy maps.

\subsection{Polar matching of the energy maps (2022 experiment)}
The roll is the weak point of the current algorithm: the coarse search only finds the yaw and the pitch, and the pyramidal refinement only corrects a small roll. On a flight at Thaumiers (November 2022) the infrared camera was mounted with a constant roll of about $-2.4^\circ$, i.e. displacements of 100 pixels in the corners, too much for the refinement (GitHub issue \#33). An experiment combined the two approaches: the energy maps are resampled in polar coordinates around the image center, where a roll becomes a vertical translation, and the cost \eqref{eq:ann_cost} is computed for every shift of $1^\circ$; a parabola fitted on the three lowest values gives the sub-degree roll (\texttt{coarse\_roll\_estimation}, branch \texttt{local\_alignment\_improvements}, commit \texttt{86353cb}, not merged). Figure~\ref{fig:ann_polar} reproduces it: on point A, rolled by a known $2.4^\circ$, the estimate is $2.26^\circ$. On point B, the road is at about $34^\circ$ in the visible image and at about $28^\circ$ in the infrared image after the coarse search, and the polar matching finds a roll of $-6.3^\circ$, but with a second minimum almost as deep at $+5^\circ$ (Figure~\ref{fig:ann_polar_b}): the profile is much less clear-cut than on point A.

\figureIRD[fig:ann_polar]{align_polar_roll.jpg}{Polar matching on point A, whose infrared image was rolled by $2.4^\circ$ on purpose. Left: energy maps of both images in polar coordinates (angle vertically, distance to the center horizontally); a roll is a vertical shift. Right: cost as a function of the angular shift, and parabola fitted around the minimum ($2.26^\circ$).}{1}

\figureIRD[fig:ann_polar_b]{align_polar_roll_b.jpg}{The same on point B, without artificial roll, after the coarse search: minimum at $-6.3^\circ$, but a second minimum near $+5^\circ$ is almost as deep.}{1}

\section{History of the development}\label{sec:ann_history}
The git history of the project gives the following milestones.
\begin{itemize}
	\item \textbf{August 2021}: semi-automatic alignment. The user turns the infrared image with three sliders (yaw, pitch, roll, equation~\eqref{eq:ann_rotation_homography}); the result is refined by phase correlation, by polar phase correlation for the roll, or by SIFT points or block phase correlation with an affine fit. Comparison on the absolute value of the gradient rather than on the intensities.
	\item \textbf{September 2021}: lens shading calibration of both cameras, applied to the RAW data instead of the RawTherapee correction; cost functions (SSD, NTG) on several representations, Newton search and quadratic fit of the minimum (from the HDR+ supplementary material~[\hyperref[ann_hasinoff]{10}]), then pyramidal iterative search fitting a homography on the block vectors, accelerated with Numba.
	\item \textbf{October 2021}: local warp following the vector field, virtual camera with a wider field for the fish-eye, full pipeline coarse search + pyramid, choice of the minimum by the Hessian, single resampling (``You Only Warp Once''), outputs local and global.
	\item \textbf{March 2022}: GitHub issue \#4, ``fix local alignment issues'' (still open): use the curvature of the cost to decide whether a local vector is reliable, regularize the field with its neighbors.
	\item \textbf{November--December 2022}: issue \#33, roll of the infrared camera; polar matching of the energy maps, parabolic refinement of the minimum (experimental branch).
	\item \textbf{November 2026}: speed-up without change of the results (cost through the linearity of the Sobel filter, pyramids in float32 with OpenCV, reduction by 32 computed on the useful rows and columns only): about 15\,s per shoot point; realignment by hand from the review in a dedicated window (\S\ref{btn:review_redo}).
\end{itemize}

\section{Diagnostic traces}\label{sec:ann_traces}
The figures of this appendix come from the diagnostic traces of the alignment. They are not written by Step~4 by default (about 100\,MB of figures per point). They are obtained:
\begin{itemize}
	\item by running the \texttt{VIS\_XXXX\_REDO.bat} (or \texttt{.sh}) script written next to each aligned image: a ``full recovery run'' with the \texttt{--debug} option, which opens the slider window then writes everything;
	\item from the review, when the log detail of the Step~4 settings (\bouton{⚙}) is ``Debug'': a realignment by hand then makes the full run;
	\item from the command line: \texttt{python automatic\_registration.py --images VIS\_XXXX.dng NIR\_YYYY.dng --outdir OUT --debug}.
\end{itemize}
The traces are written in the folder \texttt{VIS\_XXXX\_align\_traces}:
\begin{itemize}
	\item \texttt{Full Search\_\ldots}: the coarse search (cost surface, \texttt{.npy} array of the costs);
	\item \texttt{\_PADDED\_LOWRES\_*}, \texttt{\_LOWRES\_*}: energy maps at 1/32 before and after the coarse search;
	\item \texttt{ds16\_\ldots\_cost\_surface\_Y\_X\_\_itNN.png}: the inputs and the four cost surfaces of block $(Y, X)$ at iteration NN; \texttt{\ldots overview\_cost\_surfaces\ldots}: all the blocks of an iteration; \texttt{\ldots\_cost\_search\_\_itNN.npy}: the costs, block centers and scale, readable with \texttt{numpy.load(\ldots, allow\_pickle=True)};
	\item \texttt{\_msr\_itNN\_dsMM\_\ldots}, \texttt{\_image\_itNN\_\ldots}: energy maps and images at each iteration;
	\item \texttt{FLOW\_itNN\_\ldots}: vector field and image warped by it;
	\item \texttt{FULLRES\_REF.jpg}, \texttt{FULLRES\_REGISTERED\_COARSE.jpg}, \texttt{FULLRES\_REGISTERED\_REFINED\_WARPED\_ONCE\_ONLY.jpg}: the reference and the infrared image after each stage, to compare by flipping between them;
	\item \texttt{motion\_model.npy}: the motion model.
\end{itemize}

\section{Limits and open problems}\label{sec:ann_limits}
Multispectral alignment is an ill-posed problem, and the algorithm fails in identifiable situations:
\begin{itemize}
	\item the two cameras do not see the same scene: gust of wind, large time difference, low altitude (parallax);
	\item the area has few structures: water, uniform field, forest canopy; the cost is then flat or misleading (Figure~\ref{fig:ann_profiles}c);
	\item the misalignment exceeds what the stages can correct: a roll of several degrees, which the coarse search does not model;
	\item the coarse search falls into a misleading minimum.
\end{itemize}
Point B gathers several of these difficulties: large yaw ($8.2^\circ$), probable roll of several degrees, large vegetation areas whose textures differ between the bands. The block vectors stay at their limit of $\pm 80$ then $\pm 20$ pixels (Figures~\ref{fig:ann_blocks_hard} and \ref{fig:ann_convergence}); the RANSAC fit then selects a homography with strong perspective terms, which moves the corners of the image by several hundred pixels. The result is acceptable near the center and wrong towards the corners (stretched textures at the bottom left). In our tests, starting the alignment from the roll estimated by polar matching ($\pm 6.3^\circ$) did not make it converge either: this point needs a realignment by hand. Such failures are detected in the review (Chapter~\ref{ch:review}), which measures the residual offset of the edges.

\figureIRD[fig:ann_blocks_hard]{align_blocks_hard.jpg}{First iteration of the refinement of point B: the cost surfaces of the vegetation blocks have no clear minimum, the minima are at random positions, often on the edge of the search window (vectors magnified 10 times).}{1}

The realignment by hand (\S\ref{btn:review_redo}) replaces the zero initial angles by angles chosen by eye, including the roll. These angles are only a starting point: the coarse search and the pyramidal refinement start again from there and generally find an accurate alignment.

Possible improvements, partly listed in GitHub issues \#4 and \#33:
\begin{itemize}
	\item use the curvature of each cost surface (matrix $A$ of the quadratic fit) as a confidence: a vector measured on a flat surface should not count, neither in the homography nor in the local correction;
	\item regularize the local vector field (median of the neighbors, smoothness constraint) instead of using it raw;
	\item add the roll to the coarse search (polar matching, or a search over a few roll values);
	\item normalize the cost (denominator of the NTG), so that a block with less texture is not favored;
	\item compute the energy maps on the logarithm of the intensities, to compare relative contrasts (\S\ref{sec:ann_preparation});
	\item constrain the global model when the vectors are unreliable: three angles (and possibly a scale) instead of the 8 parameters of a general homography, or a penalty on its perspective terms;
	\item for the difficult cases, multimodal matching methods learned on data (dense optical flow networks such as RAFT~[\hyperref[ann_raft]{24}], descriptors trained on visible--infrared pairs), surveyed by Jiang et al.~[\hyperref[ann_jiang]{20}];
	\item calibrate the colors of the two cameras (radiometric calibration) to compute physically meaningful indices.
\end{itemize}

\section{References}\label{sec:ann_references}
{\renewcommand{\labelenumi}{[\arabic{enumi}]}
\begin{enumerate}
	\item\label{ann_hz} R. Hartley, A. Zisserman, \emph{Multiple View Geometry in Computer Vision}, 2nd edition, Cambridge University Press, 2004. Chapters 2 (projective geometry of the plane, homographies), 4 (estimation of a homography, DLT, RANSAC), 6 (camera models), 7.4 (radial distortion), 8.4 (rotating camera, synthetic rotations, mosaics).
	\item\label{ann_szeliski} R. Szeliski, \emph{Computer Vision: Algorithms and Applications}, 2nd edition, Springer, 2022 (chapter 2: image formation; chapters 8 and 9: alignment and motion estimation). Free online at \texttt{szeliski.org/Book}.
	\item\label{ann_brown} D. C. Brown, \emph{Close-range camera calibration}, Photogrammetric Engineering 37(8), 1971.
	\item\label{ann_zhang} Z. Zhang, \emph{A flexible new technique for camera calibration}, IEEE Transactions on Pattern Analysis and Machine Intelligence 22(11), 2000.
	\item\label{ann_viola} P. Viola, W. M. Wells III, \emph{Alignment by maximization of mutual information}, International Journal of Computer Vision 24(2), 1997.
	\item\label{ann_scharr} H. Scharr, \emph{Optimale Operatoren in der digitalen Bildverarbeitung}, PhD thesis, University of Heidelberg, 2000.
	\item\label{ann_roberts} L. G. Roberts, \emph{Machine perception of three-dimensional solids}, PhD thesis, MIT, 1963.
	\item\label{ann_ntg} S.-J. Chen, H.-L. Shen, C. Li, J. H. Xin, \emph{Normalized Total Gradient: a new measure for multispectral image registration}, IEEE Transactions on Image Processing 27(3), 2018 (\texttt{arXiv:1702.04562}).
	\item\label{ann_irani} M. Irani, P. Anandan, \emph{Robust multi-sensor image alignment}, ICCV 1998.
	\item\label{ann_hasinoff} S. W. Hasinoff et al., \emph{Burst photography for high dynamic range and low-light imaging on mobile cameras}, SIGGRAPH Asia 2016 (supplementary material: pyramidal alignment by blocks, quadratic sub-pixel refinement).
	\item\label{ann_burt} P. J. Burt, E. H. Adelson, \emph{The Laplacian pyramid as a compact image code}, IEEE Transactions on Communications 31(4), 1983.
	\item\label{ann_lucas} B. D. Lucas, T. Kanade, \emph{An iterative image registration technique with an application to stereo vision}, IJCAI 1981.
	\item\label{ann_baker} S. Baker, I. Matthews, \emph{Lucas-Kanade 20 years on: a unifying framework}, International Journal of Computer Vision 56(3), 2004.
	\item\label{ann_ransac} M. A. Fischler, R. C. Bolles, \emph{Random sample consensus: a paradigm for model fitting with applications to image analysis and automated cartography}, Communications of the ACM 24(6), 1981.
	\item\label{ann_kuglin} C. D. Kuglin, D. C. Hines, \emph{The phase correlation image alignment method}, Proceedings of the IEEE International Conference on Cybernetics and Society, 1975.
	\item\label{ann_reddy} B. S. Reddy, B. N. Chatterji, \emph{An FFT-based technique for translation, rotation, and scale-invariant image registration}, IEEE Transactions on Image Processing 5(8), 1996.
	\item\label{ann_rabatel} G. Rabatel, S. Labbé, \emph{Registration of visible and near infrared unmanned aerial vehicle images based on Fourier-Mellin transform}, Precision Agriculture 17, 2016.
	\item\label{ann_vayssade} J.-A. Vayssade, G. Jones, J.-N. Paoli, C. Gée, \emph{Two-step multi-spectral registration via key-point detector and gradient similarity}, 2020.
	\item\label{ann_zitova} B. Zitová, J. Flusser, \emph{Image registration methods: a survey}, Image and Vision Computing 21(11), 2003.
	\item\label{ann_jiang} X. Jiang, J. Ma, G. Xiao, Z. Shao, X. Guo, \emph{A review of multimodal image matching: methods and applications}, Information Fusion 73, 2021.
	\item\label{ann_goldman} D. B. Goldman, \emph{Vignette and exposure calibration and compensation}, IEEE Transactions on Pattern Analysis and Machine Intelligence 32(12), 2010.
	\item\label{ann_harris} C. Harris, M. Stephens, \emph{A combined corner and edge detector}, Alvey Vision Conference, 1988.
	\item\label{ann_keys} R. G. Keys, \emph{Cubic convolution interpolation for digital image processing}, IEEE Transactions on Acoustics, Speech, and Signal Processing 29(6), 1981.
	\item\label{ann_raft} Z. Teed, J. Deng, \emph{RAFT: Recurrent All-Pairs Field Transforms for optical flow}, ECCV 2020.
\end{enumerate}}

\section{Determination of the theoretical ``Pitch'' and ``Yaw''.}\label{sec:ann_theorique}

\subsection{Description of the camera positioning.}

The first step in registering the infrared image onto the visible image consists in applying a translation to the NIR image so as to align its center with that of the VIS image. Indeed, when the pictures are taken, the viewing axes of the two cameras are not parallel. Moreover, the cameras are not triggered simultaneously. The drone camera is mounted on a ``gimbal'' that maintains the viewing direction whatever the movements of the drone
\footnote{Within the limits of the specifications given by the manufacturer DJI and reproduced in Figure~\ref{AnglesDroneGimbal}}.

\figureIRD[AnglesDroneGimbal]{AnglesDroneGimbal.pdf}{Schematic description of the ``gimbal'' and of the allowed angular ranges.}{0.6}

Conversely, the NIR camera is rigidly attached to the drone (Figure~\ref{Definition_Reperes_Cameras}). Its viewing angle therefore follows the movements of the drone. It is triggered continuously with a period $\Delta T^{Nir} = 3s$. This is the highest possible frequency for the SJCam M20.\\


\figureIRD[Definition_Reperes_Cameras]{Definition_Reperes_Cameras.pdf}{Definition of the yaw, pitch and roll angles of the drone, of the gimbal and of the two cameras VIS and NIR. The cameras point towards the ground while the drone is horizontal, which swaps the yaw and roll axes between the drone and the gimbal on the one hand and the cameras on the other hand.}{0.8}

We will take the example of the pitch, assuming that the angles are decoupled. This amounts to the implicit assumption that the yaw and roll angles are zero (or very small). In the vertical plane passing through the axis of the drone, we define the various angles used to locate the position of the drone and of the two cameras at a given time (Figure~\ref{Definition_des_angles_de_Pitch}):
\figureIRD[Definition_des_angles_de_Pitch]{Definition_des_angles_de_Pitch.pdf}{Definition of the angles in the vertical plane.}{0.8}

\begin{itemize}
	\item $\alpha\qquad$  Flight Pitch. Pitch angle of the drone.
	\item $\delta\qquad$  NIR misalignement Pitch. This angle corresponds to a positioning error of the NIR camera mount on the drone.
	\item $\varphi$ Angular wedge. Corresponds to a possible wedge correcting the error $\delta$.
	\item $\zeta\qquad$ Residual deviation Pitch. $\quad \zeta = \delta + \phi$.
	\item $\theta_{Nir}\:\:\;$ Camera NIR Pitch.  $\quad \theta_{Nir} = \alpha + \zeta$.
	\item $\beta\qquad$ Gimbal Pitch. Pitch angle of the gimbal (and of the VIS camera).
	\item $\delta_{Vis}\;\;\;$ Offset VIS. By definition $\delta_{Vis} =90^\circ$.
	\item $\theta_{Vis}\;\;\;$ Camera VIS Pitch.  $\quad \theta_{Vis} = \beta + \delta_{Vis}$.
	\item $\psi\qquad$  Theoritical NIR Pitch. $\quad \psi =\delta_{Vis}+\beta-\alpha$.
	\item $\psi_{r}\quad\;\;$ NIR Pitch. Angle between the viewing axes of the two cameras at the same time $t$. $\qquad \qquad \zeta = \psi_{r} -\psi$. \\
\end{itemize}

The Exif data of the VIS camera images give the GPS coordinates but also the yaw, pitch and roll angles of the drone and of the gimbal. In other words, $\alpha$ and $\beta$ are known at each acquisition of a VIS camera image. Tilting the viewing axis produces a trapezoidal (keystone) distortion of the image, as shown in Figure~\ref{Effet_du_pitch_sur_camera_pinhole}.

\figureIRD[Effet_du_pitch_sur_camera_pinhole]{Effet_du_pitch_sur_camera_pinhole.pdf}{Effect of pitch on the distortion of the image of a horizontal plane formed on the sensor of a ``pinhole'' camera. Note that the image on the sensor is inverted.}{0.6}


For simplicity, we will assume in the following that the NIR camera mounting has no defect, in other words $\delta=0$.
We assume that the clocks of the two cameras VIS and NIR are synchronized.
We therefore know the exact times at which the images of these two cameras were taken.
The (VIS, NIR) image pairs are formed so that the time interval between the shots of these two images is minimal.

\subsection{Computation of the theoretical pitch in the synchronized case.}

\figureIRD[Pitch_cas_synchro]{Pitch_cas_synchro.pdf}{Position of the drone and of the two cameras for images taken at the same time. The pitch has been exaggerated for the purposes of the figure.}{0.5}

If the VIS and NIR images are taken at the same time (synchronized shooting
\footnote{Clock synchronization should not be confused with shot synchronization.}) it is possible to determine, at each shot, the angle $\psi $ between the viewing axes of the two cameras.

Note that the distance between the camera centers $C^{Nir}$ and $C^{Vis}$ (``base line'') is negligible compared with the altitude $H$ of the drone above the ground.
The ratio is less than $4\;10^{-4}$.
We can therefore take $\phi \cong \theta_{Vis}$.

Finally, since we assume that the NIR camera is mounted without defect, $\theta_{Nir} =\alpha_k$. \\

By looking at Figure~\ref{Pitch_cas_synchro} and using elementary geometry formulas, we deduce the formula of the theoretical pitch for the infrared image in the case of synchronized triggering at time $t_k$
\begin{equation}
	\label{Theoritical_Pitch_syncho_2}
	\psi_k  =\phi-\theta_{Nir} =\beta_k + \delta_{Vis} - \alpha_k  \quad .
\end{equation}

The distance between the image centers, computed at ground level, is
\begin{equation}
	L = H \tan \psi_k \quad ,
\end{equation}
and the translation to apply to the NIR image to center it on the VIS image is written, with $\mathrm{f}^{Nir}$ the focal length of the infrared camera,
\begin{equation}
	\mathbf{t}_y = \mathrm{f}^{Nir} \tan \psi_k \quad .
\end{equation}

Of course this situation is an ideal case
\footnote{This is the configuration chosen for professional equipment.}.

\subsection{Estimation of the theoretical pitch in the unsynchronized case.}
In practice the VIS and NIR cameras are not triggered simultaneously.
With a ``hyperlapse'' acquisition
\footnote{The term ``time-lapse'' or chronophotography refers to taking pictures at a regular time interval with the camera at a fixed position. In ``hyperlapse'', on the contrary, the camera moves with a regular motion.}
the interval for the VIS camera is $2s$
\footnote{In practice the interval is rather equal to $2.019s$. In other words, roughly every 50 VIS images the recording date jumps by $1s$. Our algorithm corrects this phenomenon to obtain a ``timeline'' accurate to a tenth of a second.} while it is $3s$ for the NIR camera.
Figure~\ref{Periodicite-du-pitch-NIR} shows an example of a series of hyperlapse images placed on the same time axis once the clocks are synchronized.\\

\figureIRD[Periodicite-du-pitch-NIR]{Periodicite-du-pitch-NIR.pdf}{Hyperlapse images of the VIS and NIR cameras. They are out of sync by a variable positive or negative delay, with a periodicity of three visible images, i.e. $6s$. To check the desynchronization of the shots, we used the presence of a moving vehicle. Flight FLY-20211109-1m/s.}{1}


This time the base line $\mathbf{Bl}(t)$ at the time the NIR image is taken can no longer be neglected. The pitch formulas become
\footnote{ $(\mathbf{e}_1, \mathbf{e}_2, \mathbf{e}_3)$ right-handed orthonormal frame. $\mathbf{e}_2$ is along the axis of the drone and $\mathbf{e}_3$ is on the vertical axis, pointing upwards. }

\begin{equation}
	\label{Theoritical_Pitch_non_syncho_0}
	\begin{aligned}
		\psi_k&=\phi-\alpha(t) \quad,\\
		\phi&= \arctan \left(\frac{\mathbf{Bl}(t) \cdot \mathbf{e}_2}{H} +\tan \theta_{Vis} \right)\quad.\\
	\end{aligned}
\end{equation}\\

Unfortunately the pitch of the drone $\alpha(t)$ is not known at the time $t$ when the NIR image was taken. The same goes for the exact position of the drone at time $t$, and therefore for the ``base line'' $\mathbf{Bl}(t)$. \\

We will assume that between two VIS images spaced $\Delta T^{Vis} =2s$ apart, the drone moves in a straight line at a constant altitude.
We propose to determine $\alpha(t)$ and $\mathbf{Bl}(t)$ by simple linear interpolation.\\

The delay between the shots of the VIS and NIR images is known and denoted $\delta t= t_k -\,t$ .
It can be positive or negative depending on whether the NIR image is taken before or after the VIS image.
The GPS positions of the drone relative to the ground are known.
In the geographic frame
\footnote{$(\mathbf{e}_{\mathrm{W}}, \mathbf{e}_{\mathrm{N}}, \mathbf{e}_3)$ right-handed orthonormal frame.}
we will use the UTM coordinates $(X, Y, Z)$. Of course the displacements of the drone must be determined in its own frame. The orientation of the drone relative to geographic north is given by the yaw angle $\textit{Yaw}_k$. Let\\
\begin{equation}
	\label{Theoritical_Pitch_non_syncho_11}
	\begin{aligned}
		{x}^{EW}_k&=X(t_{k+1})-X(t_{k}) \quad,\\
		{y}^{SN}_k&=Y(t_{k+1})-Y(t_{k}) \quad,\\
		{x_1}_k&=-{x}^{EW}_k\; \cos \textit{Yaw}_k + {y}^{SN}_k\; \sin \textit{Yaw}_k \quad,\\
		{x_2}_k&=\;\;{x}^{EW}_k\:\: \sin \textit{Yaw}_k + {y}^{SN}_k\; \cos \textit{Yaw}_k  \quad.\\
	\end{aligned} \\
\end{equation}


\begin{equation}
	\label{Theoritical_Pitch_non_syncho_1}
	\begin{aligned}
		if \;\delta t <0 \qquad& \qquad\alpha(t)&=&\;\alpha_k \left(\frac{\delta t}{ \Delta T^{Vis} }+1\right) - \alpha_{k+1}\,\frac{\delta t}{ \Delta T^{Vis}} \quad,\\
		&\mathbf{Bl}(t) \cdot \mathbf{e}_2&=&\; {x_2}_k\,\frac{\delta t}{ \Delta T^{Vis}} \quad,&\\
		&&&\\
		if \;\delta t >0 \qquad &\qquad \alpha(t)&=&\;\alpha_{k-1} \,\frac{\delta t}{ \Delta T^{Vis}}- \alpha_k\, \left(\frac{\delta t}{ \Delta T^{Vis} }-1\right)  \quad,\\
		&\mathbf{Bl}(t) \cdot \mathbf{e}_2&=&\; {x_2}_{k-1}\,\frac{\delta t}{ \Delta T^{Vis}} \quad.&\\
	\end{aligned} \\
\end{equation}\\

An example of the positions of the drone and of the cameras is shown in Figure~\ref{Pitch_cas_non_synchro}.\\

\figureIRD[Pitch_cas_non_synchro]{Pitch_cas_non_synchro.pdf}{Case of a ``backward'' desynchronization. The NIR image was taken after the visible image. The ``base-line'' has been exaggerated. In practice the green point is located before the red point for a backward desynchronization.}{0.9}


The desynchronization gap $\delta t$ is periodic and we will denote its period $T$.
In our example $T= 2 \Delta T^{Nir} = 3 \Delta T^{Vis}$.
The pitch of the VIS camera, controlled by the gimbal, is almost constant.
In the absence of wind and at constant speed, the pitch of the drone is also constant.
Under these ideal conditions, the signal $\psi(t)$ will also have period T.
During a flight with wind, the pitch of the drone varies to compensate for turbulence.
This phenomenon can be expressed as a random pitch
\footnote{We will not go into very complex considerations on turbulence modeling. For example ergodicity, homogeneous turbulence, etc. }.
The pitch will therefore have a periodic component and a random component. The Fourier spectrum of such a signal will show a peak at the frequency $0.166\, Hz$, which corresponds to the periodic component of period $T=6\,s$.\\

\figureIRD[Image ViR]{_VIR__local_HYPERLAPSE_0057.jpg}{Multispectral VIR image produced by IRDrone. FLY-20211109-1m/s.}{0.65}

\subsection{Analysis on real data.}

Here we use the data of the mission FLY-20220125, carried out on 25 January 2022 in the late morning (10:59, solar time). The sky was clear, with a light wind (< 3 m/s) from the north-west. After the synchronization phase, the drone climbed to an altitude of 100m and then started its flight to acquire the ``useful'' images. During this phase, the orientation of the drone remained constant ($-82^\circ$), as did its speed $1.5\,m/s$. During the $232 s$ of the acquisition phase, the VIS and NIR cameras took $116$ and $85$ images respectively. The trajectory of the drone breaks down into three straight segments:
\begin{itemize}
	\item Forward flight in a straight line. Direction $\sim 281^\circ$. Duration $166s$.
	\item Flight to the right. Direction  $\sim 8^\circ$ (normal to the previous direction). Duration $26s$.
	\item Backward flight. Direction  $\sim 95^\circ$ (parallel to the initial direction). Duration $40s$.\\
\end{itemize}
\figureIRD[Fly_Blassac_trajectoire]{Fly_Blassac_trajectoire.png}{Diagram of the mission FLY-20220125. The shoot points used for the mapping are shown by the red marks. They have $|dt| \leq 0.5\,s$}{0.8}

After synchronizing the clocks, the algorithm formed 116 (VIS, NIR) image pairs. Note that every other NIR image is associated with two successive VIS images.

\figureIRD[Pitch_Yaw_reel_vs_theorique]{Process-and-Theoretical_Offset_Pitch_vol_reel_01.pdf}{Time evolution of the yaw and pitch angles during the acquisition phase. The angles are determined by the theoretical formulas and by the ``coarse process''. The theoretical angles have been corrected to account for the positioning errors of the NIR camera on the drone.  FLY-20220125.}{0.99}

The VIS camera of the drone is set to ``gimbal lock'' mode
\footnote{In this particular mode, it is not certain that the gimbal yaw information is reliable. The difference between the yaw angles of the drone and of the gimbal is constant and equal to $2.7^\circ$. For some flights this difference can reach $7^\circ$. Here the yaw angle corresponds to the roll of the cameras. However, when comparing the VIS and NIR image pairs, no rotation is observed between the two. This is a camera calibration problem when switching to gimbal lock.}.
The information retrieved from the Exif data of the drone images indicates that the yaw, pitch and roll angles of the gimbal are constant ($-78.9^\circ$, $-90^\circ$  and $0^\circ$ respectively). \\


The yaw of the NIR camera is taken equal to the roll of the drone.
The computation of the theoretical yaw to apply to the NIR image to recenter it on the visible image uses the same kind of formulas as those established for the theoretical pitch (``Theoretical yaw'' curve in Figure~\ref{Pitch_Yaw_reel_vs_theorique}). \\
The evolution of the theoretical values of the yaw and pitch angles broadly follows that obtained by the ``coarse process''.\\


Note that there is a systematic offset ($\sim +0.9^\circ$ ) on the yaw. This offset has been found on all the flights carried out so far. It can be attributed to a positioning error of the NIR camera on the drone. A systematic offset also exists for the ``theoretical'' pitch relative to the ``process'' pitch. We also attribute it to a positioning error of the NIR camera on the drone, which causes a systematic alignment error of the camera viewing axis. This residual error is $\zeta = -1.3^\circ$.
\footnote{This is confirmed by a measurement made with a digital inclinometer measuring angles with an accuracy of $1/10\, ^\circ$. The measurement was taken between the carbon mount and the surface of the NIR filter.} .\\

In some parts, the evolution of the angles looks oscillating, which the spectral analysis (DFFT) confirms.
The Fourier spectrum of the pitch shown in Figure~\ref{Fourier_Pitch_process} shows a clear peak at the frequency $0.164\,Hz$ for the theoretical values and for those obtained by the ``coarse process''.
This corresponds to a periodic component of period $T=6\,s$.
We have shown that this phenomenon is linked to the periodic variation of the ``base line'' (see Figure~\ref{Pitch_cas_non_synchro} and formulas \eqref{Theoritical_Pitch_non_syncho_0}, \eqref{Theoritical_Pitch_non_syncho_1} ).
The rest of the spectrum corresponds to the effect of the wind.
When the drone flies in a straight line forward or backward, the projection of its ``base-line'' onto its axis is periodic.
The pitch therefore has a periodic component during this flight phase.
The projection of its ``base-line'' onto the normal axis is zero.
In this flight phase, the yaw therefore only has the random component associated with the wind.
When the drone moves perpendicular to its axis (still keeping the same orientation), it is the yaw that has a periodic component, while the pitch only has a random component linked to the wind.
This phenomenon is clearly visible in Figure~\ref{Pitch_Yaw_reel_vs_theorique}. The changes in the shape of the curves even make it possible to identify the flight phases.\\


\figureIRD[Fourier_Pitch_process]{Spectre-Pitch-process-vol-reel_test.png}{Fourier transform of the Pitch computed by the ``coarse process'' from the digital images or by the theoretical formulas.  FLY-20222501.}{0.9}

Note that the theoretical pitch computed with the synchronized formula \eqref{Theoritical_Pitch_syncho_2} only involves the angles of the drone and of the gimbal. It shows no periodic behavior, as its DFFT confirms. This is consistent, since the variation of these angles is associated with the ``random'' effects of the wind on the drone (turbulence).


\subsection{Conclusion.}


The proposed method for determining a priori the yaw and pitch angles to apply to the NIR image to recenter it on the VIS image relies on several elements:\\

\begin{itemize}
	\item The VIS and NIR images must be taken in time-lapse mode.
	The recording frequency of the VIS camera (the one with the smallest step, here $2\,s$) must be regular.
	This requires fast SD cards. However, for the DJI the dates are recorded to the nearest second. One-second jumps are observed periodically. It seems that the step is not exactly an integer. For our flight FLY-20211109-1m/s we find $2.01865 \, s$. This phenomenon is found on all the flights carried out so far, with a value of $2.02 \, s$. The IRDrone software automatically determines the exact period for each flight to obtain a constant frequency.
	We use the Python date format, filling in its \texttt{micosecond} attribute.
	Note that the Exif date data provide (for some of them) a decimal value of the seconds
	The situation for the NIR camera is different. There are also jumps, but of the value of the timelapse ($3\,s$). The dates are therefore always multiples of the timelapse. Here it is clearly a failure to record an image
	\footnote{As we were able to check during the manual synchronization of the images, the selected sequence happening to contain such a recording failure.}. The dates have therefore not been corrected. During a jump of this type the ``base-line'' is larger. This remains acceptable if the speed of the drone is low. Note however that the formulas must be adapted to account for it. Again, the code handles this phenomenon automatically.\\

	\item The attitude data (yaw, pitch, roll) must be available for both cameras. This is indeed the case with our setup, since the Exif data of the DJI VIS images (dng format) provide these angles for the drone (hence for the NIR camera, which is rigidly attached to it) and for the gimbal (hence for the VIS camera, which is mounted on it).\\
	\item Computing the ``base-line'' requires knowing the UTM coordinates. Currently the altitude $H$ of the drone above the ground is determined from the altitude of the take-off point, the altitude of the drone relative to the take-off point and the altitude directly below the shot (via an IGN API for France). .\\
	\item Synchronizing the clocks of the two cameras is a critical point. We use a synchronization sequence to be carried out at each take-off. It consists in rotating
	\footnote{The target is placed on a platform rotated by a stepper motor. This motor is driven by an ARDUINO microcontroller.} a QR code (ARUCO marker) under the two cameras. An algorithm then detects this QR code in the images and deduces its orientation from it. A minimization method brings the angle/date curves together. The accuracy of this synchronization is better than $1/10\,s$. We stress that it is essential not to interrupt the hyperlapse shooting during the whole mission (synchronization phase and acquisition phase). Indeed, the synchronization value is only valid (with this accuracy) for a given hyperlapse sequence. Between two hyperlapse sequences, the synchronization value can vary by more than one second (even up to 3s in the worst case).
	Such a clock synchronization error may distort the image pairing (and the computation of the ``base-line'').
	For the example presented here, we also performed this synchronization manually from a moving object present on 25 consecutive images (i.e. 50s). The moving object had been moved specifically for this purpose during our test flight. This technique is accurate (accuracy of about $1/4 s$) but tedious. We found the same image pairing as the one obtained with the synchronization by ARUCO target. \\
	\item The pitch angle of the drone and the ``base-line'' at the time the NIR image is captured were linearly interpolated. This method gives excellent results.\\
\end{itemize}

We have verified the relevance of all these elements on more than ten flights in varied weather conditions, wind in particular.
For example, Figures~\ref{Pitch_Yaw_reel_vs_theorique_FLY2021109} and \ref{Fourier_Pitch_process_FLY2021109} correspond to a flight programmed in hyperlapse (``mode line'') with a speed of $1\,m/s$ and about $600$ images.

\figureIRD[Pitch_Yaw_reel_vs_theorique_FLY2021109]{Process-and-Theoretical_Offset_Pitch_vol_reel_02.png}{Time evolution of the yaw and pitch angles during the acquisition phase. Here the drone moves at a constant heading ($179^\circ$) parallel to its axis. FLY-20211109.}{1}

\figureIRD[Fourier_Pitch_process_FLY2021109]{Spectre-Pitch-process-vol-reel_02.png}{Fourier transform of the Pitch. FLY-20211109.}{.75}




\chapter{Georeferencing and orthorectification (Step 6)}\label{ann:georef}
\todo{Technical text: camera/drone geometry, reference frames, projection onto the DEM, accuracy, intended for advanced users.}

\end{document}
